% Options for packages loaded elsewhere
% Options for packages loaded elsewhere
\PassOptionsToPackage{unicode}{hyperref}
\PassOptionsToPackage{hyphens}{url}
\PassOptionsToPackage{dvipsnames,svgnames,x11names}{xcolor}
%
\documentclass[
  a4paper,
  openany]{ltjsbook}
\usepackage{xcolor}
\usepackage[top=25mm,bottom=25mm,left=25mm,right=25mm]{geometry}
\usepackage{amsmath,amssymb}
\setcounter{secnumdepth}{-\maxdimen} % remove section numbering
\usepackage{iftex}
\ifPDFTeX
  \usepackage[T1]{fontenc}
  \usepackage[utf8]{inputenc}
  \usepackage{textcomp} % provide euro and other symbols
\else % if luatex or xetex
  \usepackage{unicode-math} % this also loads fontspec
  \defaultfontfeatures{Scale=MatchLowercase}
  \defaultfontfeatures[\rmfamily]{Ligatures=TeX,Scale=1}
\fi
\usepackage{lmodern}
\ifPDFTeX\else
  % xetex/luatex font selection
\fi
% Use upquote if available, for straight quotes in verbatim environments
\IfFileExists{upquote.sty}{\usepackage{upquote}}{}
\IfFileExists{microtype.sty}{% use microtype if available
  \usepackage[]{microtype}
  \UseMicrotypeSet[protrusion]{basicmath} % disable protrusion for tt fonts
}{}
\makeatletter
\@ifundefined{KOMAClassName}{% if non-KOMA class
  \IfFileExists{parskip.sty}{%
    \usepackage{parskip}
  }{% else
    \setlength{\parindent}{0pt}
    \setlength{\parskip}{6pt plus 2pt minus 1pt}}
}{% if KOMA class
  \KOMAoptions{parskip=half}}
\makeatother
% Make \paragraph and \subparagraph free-standing
\makeatletter
\ifx\paragraph\undefined\else
  \let\oldparagraph\paragraph
  \renewcommand{\paragraph}{
    \@ifstar
      \xxxParagraphStar
      \xxxParagraphNoStar
  }
  \newcommand{\xxxParagraphStar}[1]{\oldparagraph*{#1}\mbox{}}
  \newcommand{\xxxParagraphNoStar}[1]{\oldparagraph{#1}\mbox{}}
\fi
\ifx\subparagraph\undefined\else
  \let\oldsubparagraph\subparagraph
  \renewcommand{\subparagraph}{
    \@ifstar
      \xxxSubParagraphStar
      \xxxSubParagraphNoStar
  }
  \newcommand{\xxxSubParagraphStar}[1]{\oldsubparagraph*{#1}\mbox{}}
  \newcommand{\xxxSubParagraphNoStar}[1]{\oldsubparagraph{#1}\mbox{}}
\fi
\makeatother

\usepackage{color}
\usepackage{fancyvrb}
\newcommand{\VerbBar}{|}
\newcommand{\VERB}{\Verb[commandchars=\\\{\}]}
\DefineVerbatimEnvironment{Highlighting}{Verbatim}{commandchars=\\\{\}}
% Add ',fontsize=\small' for more characters per line
\usepackage{framed}
\definecolor{shadecolor}{RGB}{241,243,245}
\newenvironment{Shaded}{\begin{snugshade}}{\end{snugshade}}
\newcommand{\AlertTok}[1]{\textcolor[rgb]{0.68,0.00,0.00}{#1}}
\newcommand{\AnnotationTok}[1]{\textcolor[rgb]{0.37,0.37,0.37}{#1}}
\newcommand{\AttributeTok}[1]{\textcolor[rgb]{0.40,0.45,0.13}{#1}}
\newcommand{\BaseNTok}[1]{\textcolor[rgb]{0.68,0.00,0.00}{#1}}
\newcommand{\BuiltInTok}[1]{\textcolor[rgb]{0.00,0.23,0.31}{#1}}
\newcommand{\CharTok}[1]{\textcolor[rgb]{0.13,0.47,0.30}{#1}}
\newcommand{\CommentTok}[1]{\textcolor[rgb]{0.37,0.37,0.37}{#1}}
\newcommand{\CommentVarTok}[1]{\textcolor[rgb]{0.37,0.37,0.37}{\textit{#1}}}
\newcommand{\ConstantTok}[1]{\textcolor[rgb]{0.56,0.35,0.01}{#1}}
\newcommand{\ControlFlowTok}[1]{\textcolor[rgb]{0.00,0.23,0.31}{\textbf{#1}}}
\newcommand{\DataTypeTok}[1]{\textcolor[rgb]{0.68,0.00,0.00}{#1}}
\newcommand{\DecValTok}[1]{\textcolor[rgb]{0.68,0.00,0.00}{#1}}
\newcommand{\DocumentationTok}[1]{\textcolor[rgb]{0.37,0.37,0.37}{\textit{#1}}}
\newcommand{\ErrorTok}[1]{\textcolor[rgb]{0.68,0.00,0.00}{#1}}
\newcommand{\ExtensionTok}[1]{\textcolor[rgb]{0.00,0.23,0.31}{#1}}
\newcommand{\FloatTok}[1]{\textcolor[rgb]{0.68,0.00,0.00}{#1}}
\newcommand{\FunctionTok}[1]{\textcolor[rgb]{0.28,0.35,0.67}{#1}}
\newcommand{\ImportTok}[1]{\textcolor[rgb]{0.00,0.46,0.62}{#1}}
\newcommand{\InformationTok}[1]{\textcolor[rgb]{0.37,0.37,0.37}{#1}}
\newcommand{\KeywordTok}[1]{\textcolor[rgb]{0.00,0.23,0.31}{\textbf{#1}}}
\newcommand{\NormalTok}[1]{\textcolor[rgb]{0.00,0.23,0.31}{#1}}
\newcommand{\OperatorTok}[1]{\textcolor[rgb]{0.37,0.37,0.37}{#1}}
\newcommand{\OtherTok}[1]{\textcolor[rgb]{0.00,0.23,0.31}{#1}}
\newcommand{\PreprocessorTok}[1]{\textcolor[rgb]{0.68,0.00,0.00}{#1}}
\newcommand{\RegionMarkerTok}[1]{\textcolor[rgb]{0.00,0.23,0.31}{#1}}
\newcommand{\SpecialCharTok}[1]{\textcolor[rgb]{0.37,0.37,0.37}{#1}}
\newcommand{\SpecialStringTok}[1]{\textcolor[rgb]{0.13,0.47,0.30}{#1}}
\newcommand{\StringTok}[1]{\textcolor[rgb]{0.13,0.47,0.30}{#1}}
\newcommand{\VariableTok}[1]{\textcolor[rgb]{0.07,0.07,0.07}{#1}}
\newcommand{\VerbatimStringTok}[1]{\textcolor[rgb]{0.13,0.47,0.30}{#1}}
\newcommand{\WarningTok}[1]{\textcolor[rgb]{0.37,0.37,0.37}{\textit{#1}}}

\usepackage{longtable,booktabs,array}
\newcounter{none} % for unnumbered tables
\usepackage{calc} % for calculating minipage widths
% Correct order of tables after \paragraph or \subparagraph
\usepackage{etoolbox}
\makeatletter
\patchcmd\longtable{\par}{\if@noskipsec\mbox{}\fi\par}{}{}
\makeatother
% Allow footnotes in longtable head/foot
\IfFileExists{footnotehyper.sty}{\usepackage{footnotehyper}}{\usepackage{footnote}}
\makesavenoteenv{longtable}
\usepackage{graphicx}
\makeatletter
\newsavebox\pandoc@box
\newcommand*\pandocbounded[1]{% scales image to fit in text height/width
  \sbox\pandoc@box{#1}%
  \Gscale@div\@tempa{\textheight}{\dimexpr\ht\pandoc@box+\dp\pandoc@box\relax}%
  \Gscale@div\@tempb{\linewidth}{\wd\pandoc@box}%
  \ifdim\@tempb\p@<\@tempa\p@\let\@tempa\@tempb\fi% select the smaller of both
  \ifdim\@tempa\p@<\p@\scalebox{\@tempa}{\usebox\pandoc@box}%
  \else\usebox{\pandoc@box}%
  \fi%
}
% Set default figure placement to htbp
\def\fps@figure{htbp}
\makeatother





\setlength{\emergencystretch}{3em} % prevent overfull lines

\providecommand{\tightlist}{%
  \setlength{\itemsep}{0pt}\setlength{\parskip}{0pt}}



 


\usepackage[haranoaji]{luatexja-preset}
\usepackage{titling}

\pretitle{%
  \begin{center}
  \fontsize{30}{36}\selectfont
  \bfseries
}

\posttitle{%
  \par
  \end{center}
  \vskip 1.5em
}

\preauthor{%
  \begin{center}
  \large
}

\postauthor{%
  \par
  \end{center}
}
\makeatletter
\@ifpackageloaded{bookmark}{}{\usepackage{bookmark}}
\makeatother
\makeatletter
\@ifpackageloaded{caption}{}{\usepackage{caption}}
\AtBeginDocument{%
\ifdefined\contentsname
  \renewcommand*\contentsname{Table of contents}
\else
  \newcommand\contentsname{Table of contents}
\fi
\ifdefined\listfigurename
  \renewcommand*\listfigurename{List of Figures}
\else
  \newcommand\listfigurename{List of Figures}
\fi
\ifdefined\listtablename
  \renewcommand*\listtablename{List of Tables}
\else
  \newcommand\listtablename{List of Tables}
\fi
\ifdefined\figurename
  \renewcommand*\figurename{Figure}
\else
  \newcommand\figurename{Figure}
\fi
\ifdefined\tablename
  \renewcommand*\tablename{Table}
\else
  \newcommand\tablename{Table}
\fi
}
\@ifpackageloaded{float}{}{\usepackage{float}}
\floatstyle{ruled}
\@ifundefined{c@chapter}{\newfloat{codelisting}{h}{lop}}{\newfloat{codelisting}{h}{lop}[chapter]}
\floatname{codelisting}{Listing}
\newcommand*\listoflistings{\listof{codelisting}{List of Listings}}
\makeatother
\makeatletter
\makeatother
\makeatletter
\@ifpackageloaded{caption}{}{\usepackage{caption}}
\@ifpackageloaded{subcaption}{}{\usepackage{subcaption}}
\makeatother
\usepackage{bookmark}
\IfFileExists{xurl.sty}{\usepackage{xurl}}{} % add URL line breaks if available
\urlstyle{same}
\hypersetup{
  pdftitle={金融モデル・テクノロジー入門},
  pdfauthor={Eiichi Horigome},
  colorlinks=true,
  linkcolor={blue},
  filecolor={Maroon},
  citecolor={Blue},
  urlcolor={Blue},
  pdfcreator={LaTeX via pandoc}}


\title{金融モデル・テクノロジー入門}
\usepackage{etoolbox}
\makeatletter
\providecommand{\subtitle}[1]{% add subtitle to \maketitle
  \apptocmd{\@title}{\par {\large #1 \par}}{}{}
}
\makeatother
\subtitle{Quant Technology 1st Step}
\author{Eiichi Horigome}
\date{}
\begin{document}
\maketitle

\renewcommand*\contentsname{Table of contents}
{
\setcounter{tocdepth}{2}
\tableofcontents
}

\bookmarksetup{startatroot}

\chapter{はじめに}\label{ux306fux3058ux3081ux306b}

『ラグランジェ・ファイナンス』は、実証会計・ファイナンスと金融工学を、Rおよびtidyverseを用いて学ぶための入門書である。

第I部では、Rの基本操作、財務データ、株式時系列、ファクターモデル、資産価格モデルおよびイベントスタディを扱う。第II部では、確定キャッシュフロー、金利期間構造、OISスワップ、オプション評価モデル、信用リスクおよびエキゾチック・デリバティブを扱う。

各章では、理論だけでなく、再現可能なコードと計算結果を対応させながら金融分析の実装方法を確認する。

\part{第I部 金融工学　均衡モデル}

\chapter{第1章
Rによるデータ分析の基礎}\label{ux7b2c1ux7ae0-rux306bux3088ux308bux30c7ux30fcux30bfux5206ux6790ux306eux57faux790e}

本章では、Rによるデータ分析に必要な基本操作を、Base
Rからtidyverseへ段階的に学ぶ。前半では、オブジェクト、ベクトル、行列、リスト、データフレーム、関数、条件判定、記述統計および確率分布を扱う。後半では、\texttt{dplyr}、\texttt{stringr}、\texttt{tidyr}、\texttt{purrr}、\texttt{ggplot2}を用いて、データの読込み、整形、集計、反復処理、可視化を一つの流れとして組み立てる。さらに、\texttt{nest()}によるリスト列、\texttt{future}と\texttt{furrr}による並列計算、\texttt{dbplyr}とArrowによるデータベース・大規模データ処理へ進む。

Base
Rの項目構成は、Positが公開している\href{https://github.com/rstudio/cheatsheets/blob/main/base-r.pdf}{Base
R
Cheatsheet}を参考にしている。本章のコードでは、反復処理に明示的なループを使わず、ベクトル化または\texttt{purrr::map()}系関数を用いる。また、パイプ演算子はすべて\texttt{\textbar{}\textgreater{}}に統一する。

\section{1.1 Rの基礎}\label{rux306eux57faux790e}

\subsection{使用するパッケージ}\label{ux4f7fux7528ux3059ux308bux30d1ux30c3ux30b1ux30fcux30b8}

\begin{Shaded}
\begin{Highlighting}[]
\FunctionTok{source}\NormalTok{(}\FunctionTok{file.path}\NormalTok{(}\StringTok{"R"}\NormalTok{, }\StringTok{"utility\_LFL.R"}\NormalTok{))}
\FunctionTok{suppressPackageStartupMessages}\NormalTok{(}\FunctionTok{load\_part1\_packages\_LFL}\NormalTok{())}
\end{Highlighting}
\end{Shaded}

\subsection{Rの基本操作}\label{rux306eux57faux672cux64cdux4f5c}

\subsubsection{オブジェクトと代入}\label{ux30aaux30d6ux30b8ux30a7ux30afux30c8ux3068ux4ee3ux5165}

Rでは、計算結果やデータをオブジェクトとして保存する。代入には\texttt{\textless{}-}を用いる。

\begin{Shaded}
\begin{Highlighting}[]
\NormalTok{risk\_free\_rate }\OtherTok{\textless{}{-}} \FloatTok{0.01}
\NormalTok{future\_cash\_flow }\OtherTok{\textless{}{-}} \DecValTok{100}

\NormalTok{present\_value }\OtherTok{\textless{}{-}}\NormalTok{ future\_cash\_flow }\SpecialCharTok{/}\NormalTok{ (}\DecValTok{1} \SpecialCharTok{+}\NormalTok{ risk\_free\_rate)}

\NormalTok{present\_value}
\end{Highlighting}
\end{Shaded}

\begin{verbatim}
[1] 99.0099
\end{verbatim}

オブジェクト名には、内容が分かる英単語を用いる。本書では複数の単語をアンダースコアで結ぶsnake\_caseを基本とする。

\subsubsection{四則演算と代表的な数学関数}\label{ux56dbux5247ux6f14ux7b97ux3068ux4ee3ux8868ux7684ux306aux6570ux5b66ux95a2ux6570}

\begin{Shaded}
\begin{Highlighting}[]
\DecValTok{2} \SpecialCharTok{+} \DecValTok{3}
\end{Highlighting}
\end{Shaded}

\begin{verbatim}
[1] 5
\end{verbatim}

\begin{Shaded}
\begin{Highlighting}[]
\DecValTok{10} \SpecialCharTok{{-}} \DecValTok{4}
\end{Highlighting}
\end{Shaded}

\begin{verbatim}
[1] 6
\end{verbatim}

\begin{Shaded}
\begin{Highlighting}[]
\DecValTok{6} \SpecialCharTok{*} \DecValTok{7}
\end{Highlighting}
\end{Shaded}

\begin{verbatim}
[1] 42
\end{verbatim}

\begin{Shaded}
\begin{Highlighting}[]
\DecValTok{100} \SpecialCharTok{/} \DecValTok{4}
\end{Highlighting}
\end{Shaded}

\begin{verbatim}
[1] 25
\end{verbatim}

\begin{Shaded}
\begin{Highlighting}[]
\DecValTok{2}\SpecialCharTok{\^{}}\DecValTok{5}
\end{Highlighting}
\end{Shaded}

\begin{verbatim}
[1] 32
\end{verbatim}

\begin{Shaded}
\begin{Highlighting}[]
\FunctionTok{sqrt}\NormalTok{(}\DecValTok{16}\NormalTok{)}
\end{Highlighting}
\end{Shaded}

\begin{verbatim}
[1] 4
\end{verbatim}

\begin{Shaded}
\begin{Highlighting}[]
\FunctionTok{log}\NormalTok{(}\FunctionTok{exp}\NormalTok{(}\DecValTok{1}\NormalTok{))}
\end{Highlighting}
\end{Shaded}

\begin{verbatim}
[1] 1
\end{verbatim}

\begin{Shaded}
\begin{Highlighting}[]
\FunctionTok{abs}\NormalTok{(}\SpecialCharTok{{-}}\DecValTok{3}\NormalTok{)}
\end{Highlighting}
\end{Shaded}

\begin{verbatim}
[1] 3
\end{verbatim}

\begin{Shaded}
\begin{Highlighting}[]
\FunctionTok{round}\NormalTok{(pi, }\AttributeTok{digits =} \DecValTok{3}\NormalTok{)}
\end{Highlighting}
\end{Shaded}

\begin{verbatim}
[1] 3.142
\end{verbatim}

\begin{Shaded}
\begin{Highlighting}[]
\FunctionTok{signif}\NormalTok{(pi, }\AttributeTok{digits =} \DecValTok{4}\NormalTok{)}
\end{Highlighting}
\end{Shaded}

\begin{verbatim}
[1] 3.142
\end{verbatim}

\subsubsection{オブジェクトの型と構造}\label{ux30aaux30d6ux30b8ux30a7ux30afux30c8ux306eux578bux3068ux69cbux9020}

\begin{Shaded}
\begin{Highlighting}[]
\NormalTok{numeric\_value }\OtherTok{\textless{}{-}} \FloatTok{10.5}
\NormalTok{integer\_value }\OtherTok{\textless{}{-}} \DecValTok{10}\NormalTok{L}
\NormalTok{character\_value }\OtherTok{\textless{}{-}} \StringTok{"Finance"}
\NormalTok{logical\_value }\OtherTok{\textless{}{-}} \ConstantTok{TRUE}

\FunctionTok{class}\NormalTok{(numeric\_value)}
\end{Highlighting}
\end{Shaded}

\begin{verbatim}
[1] "numeric"
\end{verbatim}

\begin{Shaded}
\begin{Highlighting}[]
\FunctionTok{class}\NormalTok{(integer\_value)}
\end{Highlighting}
\end{Shaded}

\begin{verbatim}
[1] "integer"
\end{verbatim}

\begin{Shaded}
\begin{Highlighting}[]
\FunctionTok{class}\NormalTok{(character\_value)}
\end{Highlighting}
\end{Shaded}

\begin{verbatim}
[1] "character"
\end{verbatim}

\begin{Shaded}
\begin{Highlighting}[]
\FunctionTok{class}\NormalTok{(logical\_value)}
\end{Highlighting}
\end{Shaded}

\begin{verbatim}
[1] "logical"
\end{verbatim}

\begin{Shaded}
\begin{Highlighting}[]
\FunctionTok{typeof}\NormalTok{(numeric\_value)}
\end{Highlighting}
\end{Shaded}

\begin{verbatim}
[1] "double"
\end{verbatim}

\begin{Shaded}
\begin{Highlighting}[]
\FunctionTok{str}\NormalTok{(numeric\_value)}
\end{Highlighting}
\end{Shaded}

\begin{verbatim}
 num 10.5
\end{verbatim}

\subsubsection{欠損値}\label{ux6b20ux640dux5024}

欠損値は\texttt{NA}で表す。欠損値の有無は\texttt{is.na()}で確認する。

\begin{Shaded}
\begin{Highlighting}[]
\NormalTok{returns\_with\_na }\OtherTok{\textless{}{-}} \FunctionTok{c}\NormalTok{(}\FloatTok{0.02}\NormalTok{, }\ConstantTok{NA}\NormalTok{, }\SpecialCharTok{{-}}\FloatTok{0.01}\NormalTok{, }\FloatTok{0.03}\NormalTok{)}

\FunctionTok{is.na}\NormalTok{(returns\_with\_na)}
\end{Highlighting}
\end{Shaded}

\begin{verbatim}
[1] FALSE  TRUE FALSE FALSE
\end{verbatim}

\begin{Shaded}
\begin{Highlighting}[]
\FunctionTok{sum}\NormalTok{(}\FunctionTok{is.na}\NormalTok{(returns\_with\_na))}
\end{Highlighting}
\end{Shaded}

\begin{verbatim}
[1] 1
\end{verbatim}

\begin{Shaded}
\begin{Highlighting}[]
\FunctionTok{mean}\NormalTok{(returns\_with\_na, }\AttributeTok{na.rm =} \ConstantTok{TRUE}\NormalTok{)}
\end{Highlighting}
\end{Shaded}

\begin{verbatim}
[1] 0.01333333
\end{verbatim}

\subsection{ベクトル}\label{ux30d9ux30afux30c8ux30eb}

\subsubsection{ベクトルの作成}\label{ux30d9ux30afux30c8ux30ebux306eux4f5cux6210}

\begin{Shaded}
\begin{Highlighting}[]
\NormalTok{interest\_rates }\OtherTok{\textless{}{-}} \FunctionTok{c}\NormalTok{(}\FloatTok{0.01}\NormalTok{, }\FloatTok{0.015}\NormalTok{, }\FloatTok{0.02}\NormalTok{)}
\NormalTok{year\_sequence }\OtherTok{\textless{}{-}} \DecValTok{1}\SpecialCharTok{:}\DecValTok{5}
\NormalTok{fine\_grid }\OtherTok{\textless{}{-}} \FunctionTok{seq}\NormalTok{(}\AttributeTok{from =} \DecValTok{0}\NormalTok{, }\AttributeTok{to =} \DecValTok{1}\NormalTok{, }\AttributeTok{by =} \FloatTok{0.25}\NormalTok{)}
\NormalTok{repeated\_values }\OtherTok{\textless{}{-}} \FunctionTok{rep}\NormalTok{(}\DecValTok{100}\NormalTok{, }\AttributeTok{times =} \DecValTok{5}\NormalTok{)}

\NormalTok{interest\_rates}
\end{Highlighting}
\end{Shaded}

\begin{verbatim}
[1] 0.010 0.015 0.020
\end{verbatim}

\begin{Shaded}
\begin{Highlighting}[]
\NormalTok{year\_sequence}
\end{Highlighting}
\end{Shaded}

\begin{verbatim}
[1] 1 2 3 4 5
\end{verbatim}

\begin{Shaded}
\begin{Highlighting}[]
\NormalTok{fine\_grid}
\end{Highlighting}
\end{Shaded}

\begin{verbatim}
[1] 0.00 0.25 0.50 0.75 1.00
\end{verbatim}

\begin{Shaded}
\begin{Highlighting}[]
\NormalTok{repeated\_values}
\end{Highlighting}
\end{Shaded}

\begin{verbatim}
[1] 100 100 100 100 100
\end{verbatim}

\subsubsection{ベクトルの要素}\label{ux30d9ux30afux30c8ux30ebux306eux8981ux7d20}

Rの添字は1から始まる。

\begin{Shaded}
\begin{Highlighting}[]
\NormalTok{interest\_rates[}\DecValTok{1}\NormalTok{]}
\end{Highlighting}
\end{Shaded}

\begin{verbatim}
[1] 0.01
\end{verbatim}

\begin{Shaded}
\begin{Highlighting}[]
\NormalTok{interest\_rates[}\DecValTok{2}\SpecialCharTok{:}\DecValTok{3}\NormalTok{]}
\end{Highlighting}
\end{Shaded}

\begin{verbatim}
[1] 0.015 0.020
\end{verbatim}

\begin{Shaded}
\begin{Highlighting}[]
\NormalTok{interest\_rates[}\SpecialCharTok{{-}}\DecValTok{1}\NormalTok{]}
\end{Highlighting}
\end{Shaded}

\begin{verbatim}
[1] 0.015 0.020
\end{verbatim}

\begin{Shaded}
\begin{Highlighting}[]
\NormalTok{interest\_rates[interest\_rates }\SpecialCharTok{\textgreater{}=} \FloatTok{0.015}\NormalTok{]}
\end{Highlighting}
\end{Shaded}

\begin{verbatim}
[1] 0.015 0.020
\end{verbatim}

\subsubsection{名前付きベクトル}\label{ux540dux524dux4ed8ux304dux30d9ux30afux30c8ux30eb}

\begin{Shaded}
\begin{Highlighting}[]
\NormalTok{portfolio\_weights }\OtherTok{\textless{}{-}} \FunctionTok{c}\NormalTok{(}\AttributeTok{stock =} \FloatTok{0.6}\NormalTok{, }\AttributeTok{bond =} \FloatTok{0.3}\NormalTok{, }\AttributeTok{cash =} \FloatTok{0.1}\NormalTok{)}

\NormalTok{portfolio\_weights}
\end{Highlighting}
\end{Shaded}

\begin{verbatim}
stock  bond  cash 
  0.6   0.3   0.1 
\end{verbatim}

\begin{Shaded}
\begin{Highlighting}[]
\NormalTok{portfolio\_weights[}\StringTok{"bond"}\NormalTok{]}
\end{Highlighting}
\end{Shaded}

\begin{verbatim}
bond 
 0.3 
\end{verbatim}

\begin{Shaded}
\begin{Highlighting}[]
\FunctionTok{sum}\NormalTok{(portfolio\_weights)}
\end{Highlighting}
\end{Shaded}

\begin{verbatim}
[1] 1
\end{verbatim}

\subsubsection{ベクトル化}\label{ux30d9ux30afux30c8ux30ebux5316}

Rでは、多くの計算をベクトル全体に一度に適用できる。

\begin{Shaded}
\begin{Highlighting}[]
\NormalTok{discount\_rates }\OtherTok{\textless{}{-}} \FunctionTok{seq}\NormalTok{(}\FloatTok{0.01}\NormalTok{, }\FloatTok{0.05}\NormalTok{, }\AttributeTok{by =} \FloatTok{0.01}\NormalTok{)}
\NormalTok{cash\_flow }\OtherTok{\textless{}{-}} \DecValTok{100}

\NormalTok{present\_values }\OtherTok{\textless{}{-}}\NormalTok{ cash\_flow }\SpecialCharTok{/}\NormalTok{ (}\DecValTok{1} \SpecialCharTok{+}\NormalTok{ discount\_rates)}

\FunctionTok{tibble}\NormalTok{(discount\_rates, present\_values)}
\end{Highlighting}
\end{Shaded}

\begin{verbatim}
# A tibble: 5 x 2
  discount_rates present_values
           <dbl>          <dbl>
1           0.01         99.010
2           0.02         98.039
3           0.03         97.087
4           0.04         96.154
5           0.05         95.238
\end{verbatim}

\subsection{行列、リスト、データフレーム}\label{ux884cux5217ux30eaux30b9ux30c8ux30c7ux30fcux30bfux30d5ux30ecux30fcux30e0}

\subsubsection{行列}\label{ux884cux5217}

\begin{Shaded}
\begin{Highlighting}[]
\NormalTok{return\_matrix }\OtherTok{\textless{}{-}} \FunctionTok{matrix}\NormalTok{(}
  \FunctionTok{c}\NormalTok{(}\FloatTok{0.01}\NormalTok{, }\FloatTok{0.02}\NormalTok{, }\SpecialCharTok{{-}}\FloatTok{0.01}\NormalTok{, }\FloatTok{0.03}\NormalTok{, }\FloatTok{0.00}\NormalTok{, }\FloatTok{0.01}\NormalTok{),}
  \AttributeTok{nrow =} \DecValTok{3}\NormalTok{,}
  \AttributeTok{ncol =} \DecValTok{2}\NormalTok{,}
  \AttributeTok{byrow =} \ConstantTok{TRUE}
\NormalTok{)}

\FunctionTok{colnames}\NormalTok{(return\_matrix) }\OtherTok{\textless{}{-}} \FunctionTok{c}\NormalTok{(}\StringTok{"stock\_a"}\NormalTok{, }\StringTok{"stock\_b"}\NormalTok{)}
\FunctionTok{rownames}\NormalTok{(return\_matrix) }\OtherTok{\textless{}{-}} \FunctionTok{c}\NormalTok{(}\StringTok{"month\_1"}\NormalTok{, }\StringTok{"month\_2"}\NormalTok{, }\StringTok{"month\_3"}\NormalTok{)}

\NormalTok{return\_matrix}
\end{Highlighting}
\end{Shaded}

\begin{verbatim}
        stock_a stock_b
month_1    0.01    0.02
month_2   -0.01    0.03
month_3    0.00    0.01
\end{verbatim}

\begin{Shaded}
\begin{Highlighting}[]
\NormalTok{return\_matrix[, }\StringTok{"stock\_a"}\NormalTok{]}
\end{Highlighting}
\end{Shaded}

\begin{verbatim}
month_1 month_2 month_3 
   0.01   -0.01    0.00 
\end{verbatim}

\begin{Shaded}
\begin{Highlighting}[]
\FunctionTok{colMeans}\NormalTok{(return\_matrix)}
\end{Highlighting}
\end{Shaded}

\begin{verbatim}
stock_a stock_b 
   0.00    0.02 
\end{verbatim}

\subsubsection{リスト}\label{ux30eaux30b9ux30c8}

リストには、型や長さの異なるオブジェクトをまとめて保存できる。

\begin{Shaded}
\begin{Highlighting}[]
\NormalTok{analysis\_result }\OtherTok{\textless{}{-}} \FunctionTok{list}\NormalTok{(}
  \AttributeTok{model\_name =} \StringTok{"simple return model"}\NormalTok{,}
  \AttributeTok{coefficients =} \FunctionTok{c}\NormalTok{(}\AttributeTok{alpha =} \FloatTok{0.001}\NormalTok{, }\AttributeTok{beta =} \FloatTok{1.05}\NormalTok{),}
  \AttributeTok{observations =} \DecValTok{60}\NormalTok{L,}
  \AttributeTok{converged =} \ConstantTok{TRUE}
\NormalTok{)}

\NormalTok{analysis\_result}\SpecialCharTok{$}\NormalTok{model\_name}
\end{Highlighting}
\end{Shaded}

\begin{verbatim}
[1] "simple return model"
\end{verbatim}

\begin{Shaded}
\begin{Highlighting}[]
\NormalTok{analysis\_result}\SpecialCharTok{$}\NormalTok{coefficients}
\end{Highlighting}
\end{Shaded}

\begin{verbatim}
alpha  beta 
0.001 1.050 
\end{verbatim}

\begin{Shaded}
\begin{Highlighting}[]
\NormalTok{analysis\_result[[}\StringTok{"observations"}\NormalTok{]]}
\end{Highlighting}
\end{Shaded}

\begin{verbatim}
[1] 60
\end{verbatim}

\subsubsection{データフレーム}\label{ux30c7ux30fcux30bfux30d5ux30ecux30fcux30e0}

\begin{Shaded}
\begin{Highlighting}[]
\NormalTok{firm\_data\_base }\OtherTok{\textless{}{-}} \FunctionTok{data.frame}\NormalTok{(}
  \AttributeTok{firm\_id =} \DecValTok{1}\SpecialCharTok{:}\DecValTok{3}\NormalTok{,}
  \AttributeTok{firm\_name =} \FunctionTok{c}\NormalTok{(}\StringTok{"Firm A"}\NormalTok{, }\StringTok{"Firm B"}\NormalTok{, }\StringTok{"Firm C"}\NormalTok{),}
  \AttributeTok{roe =} \FunctionTok{c}\NormalTok{(}\FloatTok{0.08}\NormalTok{, }\FloatTok{0.12}\NormalTok{, }\FloatTok{0.15}\NormalTok{)}
\NormalTok{)}

\NormalTok{firm\_data\_base}
\end{Highlighting}
\end{Shaded}

\begin{verbatim}
  firm_id firm_name  roe
1       1    Firm A 0.08
2       2    Firm B 0.12
3       3    Firm C 0.15
\end{verbatim}

\begin{Shaded}
\begin{Highlighting}[]
\FunctionTok{str}\NormalTok{(firm\_data\_base)}
\end{Highlighting}
\end{Shaded}

\begin{verbatim}
'data.frame':   3 obs. of  3 variables:
 $ firm_id  : int  1 2 3
 $ firm_name: chr  "Firm A" "Firm B" "Firm C"
 $ roe      : num  0.08 0.12 0.15
\end{verbatim}

tidyverseでは、データフレームを拡張した\texttt{tibble}を用いる。

\begin{Shaded}
\begin{Highlighting}[]
\NormalTok{firm\_data }\OtherTok{\textless{}{-}} \FunctionTok{tibble}\NormalTok{(}
  \AttributeTok{firm\_id =} \DecValTok{1}\SpecialCharTok{:}\DecValTok{3}\NormalTok{,}
  \AttributeTok{firm\_name =} \FunctionTok{c}\NormalTok{(}\StringTok{"Firm A"}\NormalTok{, }\StringTok{"Firm B"}\NormalTok{, }\StringTok{"Firm C"}\NormalTok{),}
  \AttributeTok{roe =} \FunctionTok{c}\NormalTok{(}\FloatTok{0.08}\NormalTok{, }\FloatTok{0.12}\NormalTok{, }\FloatTok{0.15}\NormalTok{)}
\NormalTok{)}

\NormalTok{firm\_data}
\end{Highlighting}
\end{Shaded}

\begin{verbatim}
# A tibble: 3 x 3
  firm_id firm_name   roe
    <int> <chr>     <dbl>
1       1 Firm A     0.08
2       2 Firm B     0.12
3       3 Firm C     0.15
\end{verbatim}

少量のデータを行単位で記述する場合は\texttt{tribble()}が便利である。

\begin{Shaded}
\begin{Highlighting}[]
\NormalTok{project\_data }\OtherTok{\textless{}{-}} \FunctionTok{tribble}\NormalTok{(}
  \SpecialCharTok{\textasciitilde{}}\NormalTok{project, }\SpecialCharTok{\textasciitilde{}}\NormalTok{initial\_cost, }\SpecialCharTok{\textasciitilde{}}\NormalTok{cash\_flow,}
  \StringTok{"A"}\NormalTok{,             }\DecValTok{100}\NormalTok{,         }\DecValTok{120}\NormalTok{,}
  \StringTok{"B"}\NormalTok{,             }\DecValTok{150}\NormalTok{,         }\DecValTok{175}\NormalTok{,}
  \StringTok{"C"}\NormalTok{,             }\DecValTok{200}\NormalTok{,         }\DecValTok{215}
\NormalTok{)}

\NormalTok{project\_data}
\end{Highlighting}
\end{Shaded}

\begin{verbatim}
# A tibble: 3 x 3
  project initial_cost cash_flow
  <chr>          <dbl>     <dbl>
1 A                100       120
2 B                150       175
3 C                200       215
\end{verbatim}

\subsection{条件判定と独自関数}\label{ux6761ux4ef6ux5224ux5b9aux3068ux72ecux81eaux95a2ux6570}

\subsubsection{条件判定}\label{ux6761ux4ef6ux5224ux5b9a}

\begin{Shaded}
\begin{Highlighting}[]
\NormalTok{npv }\OtherTok{\textless{}{-}} \FloatTok{5.2}

\NormalTok{investment\_decision }\OtherTok{\textless{}{-}} \ControlFlowTok{if}\NormalTok{ (npv }\SpecialCharTok{\textgreater{}=} \DecValTok{0}\NormalTok{) \{}
  \StringTok{"投資を実行する"}
\NormalTok{\} }\ControlFlowTok{else}\NormalTok{ \{}
  \StringTok{"投資を見送る"}
\NormalTok{\}}

\NormalTok{investment\_decision}
\end{Highlighting}
\end{Shaded}

\begin{verbatim}
[1] "投資を実行する"
\end{verbatim}

ベクトルに対する複数の条件分岐には\texttt{case\_when()}を用いる。

\begin{Shaded}
\begin{Highlighting}[]
\NormalTok{credit\_scores }\OtherTok{\textless{}{-}} \FunctionTok{tibble}\NormalTok{(}\AttributeTok{score =} \FunctionTok{c}\NormalTok{(}\DecValTok{85}\NormalTok{, }\DecValTok{72}\NormalTok{, }\DecValTok{58}\NormalTok{, }\DecValTok{91}\NormalTok{))}

\NormalTok{credit\_scores }\SpecialCharTok{|\textgreater{}}
  \FunctionTok{mutate}\NormalTok{(}
    \AttributeTok{rating =} \FunctionTok{case\_when}\NormalTok{(}
\NormalTok{      score }\SpecialCharTok{\textgreater{}=} \DecValTok{80} \SpecialCharTok{\textasciitilde{}} \StringTok{"A"}\NormalTok{,}
\NormalTok{      score }\SpecialCharTok{\textgreater{}=} \DecValTok{65} \SpecialCharTok{\textasciitilde{}} \StringTok{"B"}\NormalTok{,}
      \ConstantTok{TRUE} \SpecialCharTok{\textasciitilde{}} \StringTok{"C"}
\NormalTok{    )}
\NormalTok{  )}
\end{Highlighting}
\end{Shaded}

\begin{verbatim}
# A tibble: 4 x 2
  score rating
  <dbl> <chr> 
1    85 A     
2    72 B     
3    58 C     
4    91 A     
\end{verbatim}

\subsubsection{独自関数}\label{ux72ecux81eaux95a2ux6570}

\begin{Shaded}
\begin{Highlighting}[]
\NormalTok{calculate\_present\_value }\OtherTok{\textless{}{-}} \ControlFlowTok{function}\NormalTok{(cash\_flow, discount\_rate, }\AttributeTok{years =} \DecValTok{1}\NormalTok{) \{}
\NormalTok{  cash\_flow }\SpecialCharTok{/}\NormalTok{ (}\DecValTok{1} \SpecialCharTok{+}\NormalTok{ discount\_rate)}\SpecialCharTok{\^{}}\NormalTok{years}
\NormalTok{\}}

\FunctionTok{calculate\_present\_value}\NormalTok{(}\AttributeTok{cash\_flow =} \DecValTok{100}\NormalTok{, }\AttributeTok{discount\_rate =} \FloatTok{0.02}\NormalTok{, }\AttributeTok{years =} \DecValTok{3}\NormalTok{)}
\end{Highlighting}
\end{Shaded}

\begin{verbatim}
[1] 94.23223
\end{verbatim}

関数はベクトルをそのまま受け取れるように設計すると再利用しやすい。

\begin{Shaded}
\begin{Highlighting}[]
\FunctionTok{calculate\_present\_value}\NormalTok{(}
  \AttributeTok{cash\_flow =} \FunctionTok{c}\NormalTok{(}\DecValTok{100}\NormalTok{, }\DecValTok{100}\NormalTok{, }\DecValTok{100}\NormalTok{),}
  \AttributeTok{discount\_rate =} \FloatTok{0.02}\NormalTok{,}
  \AttributeTok{years =} \DecValTok{1}\SpecialCharTok{:}\DecValTok{3}
\NormalTok{)}
\end{Highlighting}
\end{Shaded}

\begin{verbatim}
[1] 98.03922 96.11688 94.23223
\end{verbatim}

\subsection{記述統計}\label{ux8a18ux8ff0ux7d71ux8a08}

\subsubsection{代表値と散らばり}\label{ux4ee3ux8868ux5024ux3068ux6563ux3089ux3070ux308a}

\begin{Shaded}
\begin{Highlighting}[]
\NormalTok{sample\_returns }\OtherTok{\textless{}{-}} \FunctionTok{c}\NormalTok{(}\FloatTok{0.01}\NormalTok{, }\FloatTok{0.03}\NormalTok{, }\SpecialCharTok{{-}}\FloatTok{0.02}\NormalTok{, }\FloatTok{0.00}\NormalTok{, }\FloatTok{0.04}\NormalTok{, }\FloatTok{0.02}\NormalTok{)}

\FunctionTok{mean}\NormalTok{(sample\_returns)}
\end{Highlighting}
\end{Shaded}

\begin{verbatim}
[1] 0.01333333
\end{verbatim}

\begin{Shaded}
\begin{Highlighting}[]
\FunctionTok{median}\NormalTok{(sample\_returns)}
\end{Highlighting}
\end{Shaded}

\begin{verbatim}
[1] 0.015
\end{verbatim}

\begin{Shaded}
\begin{Highlighting}[]
\FunctionTok{min}\NormalTok{(sample\_returns)}
\end{Highlighting}
\end{Shaded}

\begin{verbatim}
[1] -0.02
\end{verbatim}

\begin{Shaded}
\begin{Highlighting}[]
\FunctionTok{max}\NormalTok{(sample\_returns)}
\end{Highlighting}
\end{Shaded}

\begin{verbatim}
[1] 0.04
\end{verbatim}

\begin{Shaded}
\begin{Highlighting}[]
\FunctionTok{range}\NormalTok{(sample\_returns)}
\end{Highlighting}
\end{Shaded}

\begin{verbatim}
[1] -0.02  0.04
\end{verbatim}

\begin{Shaded}
\begin{Highlighting}[]
\FunctionTok{var}\NormalTok{(sample\_returns)}
\end{Highlighting}
\end{Shaded}

\begin{verbatim}
[1] 0.0004666667
\end{verbatim}

\begin{Shaded}
\begin{Highlighting}[]
\FunctionTok{sd}\NormalTok{(sample\_returns)}
\end{Highlighting}
\end{Shaded}

\begin{verbatim}
[1] 0.02160247
\end{verbatim}

\begin{Shaded}
\begin{Highlighting}[]
\FunctionTok{quantile}\NormalTok{(sample\_returns)}
\end{Highlighting}
\end{Shaded}

\begin{verbatim}
     0%     25%     50%     75%    100% 
-0.0200  0.0025  0.0150  0.0275  0.0400 
\end{verbatim}

\subsubsection{共分散と相関}\label{ux5171ux5206ux6563ux3068ux76f8ux95a2}

\begin{Shaded}
\begin{Highlighting}[]
\NormalTok{stock\_a }\OtherTok{\textless{}{-}} \FunctionTok{c}\NormalTok{(}\FloatTok{0.01}\NormalTok{, }\FloatTok{0.03}\NormalTok{, }\SpecialCharTok{{-}}\FloatTok{0.02}\NormalTok{, }\FloatTok{0.02}\NormalTok{, }\FloatTok{0.04}\NormalTok{)}
\NormalTok{stock\_b }\OtherTok{\textless{}{-}} \FunctionTok{c}\NormalTok{(}\FloatTok{0.00}\NormalTok{, }\FloatTok{0.02}\NormalTok{, }\SpecialCharTok{{-}}\FloatTok{0.01}\NormalTok{, }\FloatTok{0.01}\NormalTok{, }\FloatTok{0.03}\NormalTok{)}

\FunctionTok{cov}\NormalTok{(stock\_a, stock\_b)}
\end{Highlighting}
\end{Shaded}

\begin{verbatim}
[1] 0.00035
\end{verbatim}

\begin{Shaded}
\begin{Highlighting}[]
\FunctionTok{cor}\NormalTok{(stock\_a, stock\_b)}
\end{Highlighting}
\end{Shaded}

\begin{verbatim}
[1] 0.9615239
\end{verbatim}

\subsubsection{順位}\label{ux9806ux4f4d}

\begin{Shaded}
\begin{Highlighting}[]
\NormalTok{firm\_returns }\OtherTok{\textless{}{-}} \FunctionTok{c}\NormalTok{(}\FloatTok{0.08}\NormalTok{, }\FloatTok{0.12}\NormalTok{, }\SpecialCharTok{{-}}\FloatTok{0.03}\NormalTok{, }\FloatTok{0.05}\NormalTok{)}

\FunctionTok{rank}\NormalTok{(firm\_returns)}
\end{Highlighting}
\end{Shaded}

\begin{verbatim}
[1] 3 4 1 2
\end{verbatim}

\begin{Shaded}
\begin{Highlighting}[]
\FunctionTok{rank}\NormalTok{(}\SpecialCharTok{{-}}\NormalTok{firm\_returns)}
\end{Highlighting}
\end{Shaded}

\begin{verbatim}
[1] 2 1 4 3
\end{verbatim}

\begin{Shaded}
\begin{Highlighting}[]
\FunctionTok{order}\NormalTok{(firm\_returns, }\AttributeTok{decreasing =} \ConstantTok{TRUE}\NormalTok{)}
\end{Highlighting}
\end{Shaded}

\begin{verbatim}
[1] 2 1 4 3
\end{verbatim}

\subsection{確率分布}\label{ux78baux7387ux5206ux5e03}

Rの確率分布関数には共通した命名規則がある。

{\def\LTcaptype{none} % do not increment counter
\begin{longtable}[]{@{}ll@{}}
\toprule\noalign{}
接頭辞 & 内容 \\
\midrule\noalign{}
\endhead
\bottomrule\noalign{}
\endlastfoot
\texttt{r} & 乱数を生成する \\
\texttt{d} & 確率密度または確率質量を求める \\
\texttt{p} & 累積確率を求める \\
\texttt{q} & 累積確率に対応する分位点を求める \\
\end{longtable}
}

\subsubsection{一様分布}\label{ux4e00ux69d8ux5206ux5e03}

\begin{Shaded}
\begin{Highlighting}[]
\FunctionTok{set\_random\_seed\_LFL}\NormalTok{(}\DecValTok{1234}\NormalTok{)}

\NormalTok{uniform\_random }\OtherTok{\textless{}{-}} \FunctionTok{runif}\NormalTok{(}\AttributeTok{n =} \DecValTok{10}\NormalTok{, }\AttributeTok{min =} \DecValTok{0}\NormalTok{, }\AttributeTok{max =} \DecValTok{1}\NormalTok{)}

\NormalTok{uniform\_random}
\end{Highlighting}
\end{Shaded}

\begin{verbatim}
 [1] 0.113703411 0.622299405 0.609274733 0.623379442 0.860915384 0.640310605
 [7] 0.009495756 0.232550506 0.666083758 0.514251141
\end{verbatim}

\begin{Shaded}
\begin{Highlighting}[]
\FunctionTok{dunif}\NormalTok{(}\FloatTok{0.5}\NormalTok{, }\AttributeTok{min =} \DecValTok{0}\NormalTok{, }\AttributeTok{max =} \DecValTok{1}\NormalTok{)}
\end{Highlighting}
\end{Shaded}

\begin{verbatim}
[1] 1
\end{verbatim}

\begin{Shaded}
\begin{Highlighting}[]
\FunctionTok{punif}\NormalTok{(}\FloatTok{0.5}\NormalTok{, }\AttributeTok{min =} \DecValTok{0}\NormalTok{, }\AttributeTok{max =} \DecValTok{1}\NormalTok{)}
\end{Highlighting}
\end{Shaded}

\begin{verbatim}
[1] 0.5
\end{verbatim}

\begin{Shaded}
\begin{Highlighting}[]
\FunctionTok{qunif}\NormalTok{(}\FloatTok{0.95}\NormalTok{, }\AttributeTok{min =} \DecValTok{0}\NormalTok{, }\AttributeTok{max =} \DecValTok{1}\NormalTok{)}
\end{Highlighting}
\end{Shaded}

\begin{verbatim}
[1] 0.95
\end{verbatim}

\subsubsection{正規分布}\label{ux6b63ux898fux5206ux5e03}

\begin{Shaded}
\begin{Highlighting}[]
\FunctionTok{set\_random\_seed\_LFL}\NormalTok{(}\DecValTok{1234}\NormalTok{)}

\NormalTok{normal\_random }\OtherTok{\textless{}{-}} \FunctionTok{rnorm}\NormalTok{(}\AttributeTok{n =} \DecValTok{10}\NormalTok{, }\AttributeTok{mean =} \DecValTok{0}\NormalTok{, }\AttributeTok{sd =} \DecValTok{1}\NormalTok{)}

\NormalTok{normal\_random}
\end{Highlighting}
\end{Shaded}

\begin{verbatim}
 [1] -1.2070657  0.2774292  1.0844412 -2.3456977  0.4291247  0.5060559
 [7] -0.5747400 -0.5466319 -0.5644520 -0.8900378
\end{verbatim}

\begin{Shaded}
\begin{Highlighting}[]
\FunctionTok{dnorm}\NormalTok{(}\DecValTok{0}\NormalTok{, }\AttributeTok{mean =} \DecValTok{0}\NormalTok{, }\AttributeTok{sd =} \DecValTok{1}\NormalTok{)}
\end{Highlighting}
\end{Shaded}

\begin{verbatim}
[1] 0.3989423
\end{verbatim}

\begin{Shaded}
\begin{Highlighting}[]
\FunctionTok{pnorm}\NormalTok{(}\FloatTok{1.96}\NormalTok{, }\AttributeTok{mean =} \DecValTok{0}\NormalTok{, }\AttributeTok{sd =} \DecValTok{1}\NormalTok{)}
\end{Highlighting}
\end{Shaded}

\begin{verbatim}
[1] 0.9750021
\end{verbatim}

\begin{Shaded}
\begin{Highlighting}[]
\FunctionTok{qnorm}\NormalTok{(}\FloatTok{0.975}\NormalTok{, }\AttributeTok{mean =} \DecValTok{0}\NormalTok{, }\AttributeTok{sd =} \DecValTok{1}\NormalTok{)}
\end{Highlighting}
\end{Shaded}

\begin{verbatim}
[1] 1.959964
\end{verbatim}

\texttt{dnorm()}は正規分布の密度、\texttt{pnorm()}は指定した値以下となる確率、\texttt{qnorm()}は指定した累積確率に対応する分位点を返す。

\begin{Shaded}
\begin{Highlighting}[]
\NormalTok{probability\_below }\OtherTok{\textless{}{-}} \FunctionTok{pnorm}\NormalTok{(}\SpecialCharTok{{-}}\FloatTok{0.02}\NormalTok{, }\AttributeTok{mean =} \FloatTok{0.01}\NormalTok{, }\AttributeTok{sd =} \FloatTok{0.03}\NormalTok{)}
\NormalTok{value\_at\_risk\_95 }\OtherTok{\textless{}{-}} \FunctionTok{qnorm}\NormalTok{(}\FloatTok{0.05}\NormalTok{, }\AttributeTok{mean =} \FloatTok{0.01}\NormalTok{, }\AttributeTok{sd =} \FloatTok{0.03}\NormalTok{)}

\NormalTok{probability\_below}
\end{Highlighting}
\end{Shaded}

\begin{verbatim}
[1] 0.1586553
\end{verbatim}

\begin{Shaded}
\begin{Highlighting}[]
\NormalTok{value\_at\_risk\_95}
\end{Highlighting}
\end{Shaded}

\begin{verbatim}
[1] -0.03934561
\end{verbatim}

\subsubsection{二項分布}\label{ux4e8cux9805ux5206ux5e03}

\begin{Shaded}
\begin{Highlighting}[]
\FunctionTok{set\_random\_seed\_LFL}\NormalTok{(}\DecValTok{1234}\NormalTok{)}

\NormalTok{binomial\_random }\OtherTok{\textless{}{-}} \FunctionTok{rbinom}\NormalTok{(}\AttributeTok{n =} \DecValTok{10}\NormalTok{, }\AttributeTok{size =} \DecValTok{20}\NormalTok{, }\AttributeTok{prob =} \FloatTok{0.6}\NormalTok{)}

\NormalTok{binomial\_random}
\end{Highlighting}
\end{Shaded}

\begin{verbatim}
 [1] 15 11 11 11 10 11 17 14 11 12
\end{verbatim}

\begin{Shaded}
\begin{Highlighting}[]
\FunctionTok{dbinom}\NormalTok{(}\DecValTok{12}\NormalTok{, }\AttributeTok{size =} \DecValTok{20}\NormalTok{, }\AttributeTok{prob =} \FloatTok{0.6}\NormalTok{)}
\end{Highlighting}
\end{Shaded}

\begin{verbatim}
[1] 0.1797058
\end{verbatim}

\begin{Shaded}
\begin{Highlighting}[]
\FunctionTok{pbinom}\NormalTok{(}\DecValTok{12}\NormalTok{, }\AttributeTok{size =} \DecValTok{20}\NormalTok{, }\AttributeTok{prob =} \FloatTok{0.6}\NormalTok{)}
\end{Highlighting}
\end{Shaded}

\begin{verbatim}
[1] 0.5841071
\end{verbatim}

\begin{Shaded}
\begin{Highlighting}[]
\FunctionTok{qbinom}\NormalTok{(}\FloatTok{0.95}\NormalTok{, }\AttributeTok{size =} \DecValTok{20}\NormalTok{, }\AttributeTok{prob =} \FloatTok{0.6}\NormalTok{)}
\end{Highlighting}
\end{Shaded}

\begin{verbatim}
[1] 16
\end{verbatim}

\subsubsection{ポアソン分布}\label{ux30ddux30a2ux30bdux30f3ux5206ux5e03}

\begin{Shaded}
\begin{Highlighting}[]
\FunctionTok{set\_random\_seed\_LFL}\NormalTok{(}\DecValTok{1234}\NormalTok{)}

\NormalTok{poisson\_random }\OtherTok{\textless{}{-}} \FunctionTok{rpois}\NormalTok{(}\AttributeTok{n =} \DecValTok{10}\NormalTok{, }\AttributeTok{lambda =} \DecValTok{3}\NormalTok{)}

\NormalTok{poisson\_random}
\end{Highlighting}
\end{Shaded}

\begin{verbatim}
 [1] 1 3 3 3 5 3 0 2 4 3
\end{verbatim}

\begin{Shaded}
\begin{Highlighting}[]
\FunctionTok{dpois}\NormalTok{(}\DecValTok{2}\NormalTok{, }\AttributeTok{lambda =} \DecValTok{3}\NormalTok{)}
\end{Highlighting}
\end{Shaded}

\begin{verbatim}
[1] 0.2240418
\end{verbatim}

\begin{Shaded}
\begin{Highlighting}[]
\FunctionTok{ppois}\NormalTok{(}\DecValTok{2}\NormalTok{, }\AttributeTok{lambda =} \DecValTok{3}\NormalTok{)}
\end{Highlighting}
\end{Shaded}

\begin{verbatim}
[1] 0.4231901
\end{verbatim}

\begin{Shaded}
\begin{Highlighting}[]
\FunctionTok{qpois}\NormalTok{(}\FloatTok{0.95}\NormalTok{, }\AttributeTok{lambda =} \DecValTok{3}\NormalTok{)}
\end{Highlighting}
\end{Shaded}

\begin{verbatim}
[1] 6
\end{verbatim}

\subsubsection{分布関数の比較}\label{ux5206ux5e03ux95a2ux6570ux306eux6bd4ux8f03}

\begin{Shaded}
\begin{Highlighting}[]
\NormalTok{distribution\_functions }\OtherTok{\textless{}{-}} \FunctionTok{tribble}\NormalTok{(}
  \SpecialCharTok{\textasciitilde{}}\NormalTok{distribution, }\SpecialCharTok{\textasciitilde{}}\NormalTok{random\_function, }\SpecialCharTok{\textasciitilde{}}\NormalTok{density\_function, }\SpecialCharTok{\textasciitilde{}}\NormalTok{cumulative\_function, }\SpecialCharTok{\textasciitilde{}}\NormalTok{quantile\_function,}
  \StringTok{"一様分布"}\NormalTok{,    }\StringTok{"runif()"}\NormalTok{,        }\StringTok{"dunif()"}\NormalTok{,         }\StringTok{"punif()"}\NormalTok{,            }\StringTok{"qunif()"}\NormalTok{,}
  \StringTok{"正規分布"}\NormalTok{,    }\StringTok{"rnorm()"}\NormalTok{,        }\StringTok{"dnorm()"}\NormalTok{,         }\StringTok{"pnorm()"}\NormalTok{,            }\StringTok{"qnorm()"}\NormalTok{,}
  \StringTok{"二項分布"}\NormalTok{,    }\StringTok{"rbinom()"}\NormalTok{,       }\StringTok{"dbinom()"}\NormalTok{,        }\StringTok{"pbinom()"}\NormalTok{,           }\StringTok{"qbinom()"}\NormalTok{,}
  \StringTok{"ポアソン分布"}\NormalTok{,}\StringTok{"rpois()"}\NormalTok{,        }\StringTok{"dpois()"}\NormalTok{,         }\StringTok{"ppois()"}\NormalTok{,            }\StringTok{"qpois()"}
\NormalTok{)}

\NormalTok{distribution\_functions}
\end{Highlighting}
\end{Shaded}

\begin{verbatim}
# A tibble: 4 x 5
  distribution random_function density_function cumulative_function
  <chr>        <chr>           <chr>            <chr>              
1 一様分布     runif()         dunif()          punif()            
2 正規分布     rnorm()         dnorm()          pnorm()            
3 二項分布     rbinom()        dbinom()         pbinom()           
4 ポアソン分布 rpois()         dpois()          ppois()            
  quantile_function
  <chr>            
1 qunif()          
2 qnorm()          
3 qbinom()         
4 qpois()          
\end{verbatim}

\section{1.2
tidyverseによるデータ整形}\label{tidyverseux306bux3088ux308bux30c7ux30fcux30bfux6574ux5f62}

\subsection{プロジェクト内のデータを読み込む}\label{ux30d7ux30edux30b8ux30a7ux30afux30c8ux5185ux306eux30c7ux30fcux30bfux3092ux8aadux307fux8fbcux3080}

本書で使用するデータは、プロジェクトルートの\texttt{simulation\_data}ディレクトリに置かれている。プロジェクトルートの探索やファイル検索は\texttt{utility.R}へ集約し、各章では分析対象となるファイル名だけを指定する。

\begin{Shaded}
\begin{Highlighting}[]
\NormalTok{tbl\_daily\_stock\_return }\OtherTok{\textless{}{-}} \FunctionTok{read\_csv\_file\_LFL}\NormalTok{(}\StringTok{"ch03\_daily\_stock\_return.csv"}\NormalTok{) }\SpecialCharTok{|\textgreater{}}
  \FunctionTok{mutate}\NormalTok{(}\AttributeTok{date =} \FunctionTok{as.Date}\NormalTok{(date))}

\NormalTok{vec\_required\_columns }\OtherTok{\textless{}{-}} \FunctionTok{c}\NormalTok{(}\StringTok{"date"}\NormalTok{, }\StringTok{"firm1"}\NormalTok{, }\StringTok{"firm2"}\NormalTok{)}
\NormalTok{vec\_missing\_columns }\OtherTok{\textless{}{-}} \FunctionTok{setdiff}\NormalTok{(vec\_required\_columns, }\FunctionTok{names}\NormalTok{(tbl\_daily\_stock\_return))}

\ControlFlowTok{if}\NormalTok{ (}\FunctionTok{length}\NormalTok{(vec\_missing\_columns) }\SpecialCharTok{\textgreater{}} \DecValTok{0}\NormalTok{) \{}
  \FunctionTok{stop}\NormalTok{(}\StringTok{"必要な列が不足しています: "}\NormalTok{, }\FunctionTok{paste}\NormalTok{(vec\_missing\_columns, }\AttributeTok{collapse =} \StringTok{", "}\NormalTok{))}
\NormalTok{\}}

\FunctionTok{show\_table\_LFL}\NormalTok{(tbl\_daily\_stock\_return)}
\end{Highlighting}
\end{Shaded}

\begin{verbatim}
# A tibble: 21 x 3
   date            firm1      firm2
   <date>          <dbl>      <dbl>
 1 2020-04-01 -0.039482   0.076963 
 2 2020-04-02  0.0059787 -0.0072488
 3 2020-04-03  0.055787  -0.017304 
 4 2020-04-06  0.041935   0.0021690
 5 2020-04-07 -0.020193   0.075545 
 6 2020-04-08 -0.0022947 -0.096212 
 7 2020-04-09  0.081707   0.078843 
 8 2020-04-10  0.034561   0.021536 
 9 2020-04-13  0.032601   0.13520  
10 2020-04-14 -0.037663   0.037979 
# i 11 more rows
\end{verbatim}

\subsubsection{データの構造を確認する}\label{ux30c7ux30fcux30bfux306eux69cbux9020ux3092ux78baux8a8dux3059ux308b}

\begin{Shaded}
\begin{Highlighting}[]
\FunctionTok{show\_data\_structure\_LFL}\NormalTok{(tbl\_daily\_stock\_return)}
\end{Highlighting}
\end{Shaded}

\begin{verbatim}
Rows: 21
Columns: 3
$ date  <date> 2020-04-01, 2020-04-02, 2020-04-03, 2020-04-06, 2020-04-07, 202~
$ firm1 <dbl> -0.039482142, 0.005978689, 0.055786605, 0.041934669, -0.02019279~
$ firm2 <dbl> 0.0769625974, -0.0072488484, -0.0173040800, 0.0021689573, 0.0755~
\end{verbatim}

\begin{Shaded}
\begin{Highlighting}[]
\FunctionTok{summary}\NormalTok{(tbl\_daily\_stock\_return)}
\end{Highlighting}
\end{Shaded}

\begin{verbatim}
      date                firm1               firm2           
 Min.   :2020-04-01   Min.   :-0.072071   Min.   :-0.0962117  
 1st Qu.:2020-04-08   1st Qu.:-0.027029   1st Qu.:-0.0196147  
 Median :2020-04-15   Median : 0.005979   Median :-0.0003056  
 Mean   :2020-04-15   Mean   : 0.014161   Mean   : 0.0080253  
 3rd Qu.:2020-04-22   3rd Qu.: 0.044812   3rd Qu.: 0.0379793  
 Max.   :2020-04-30   Max.   : 0.171297   Max.   : 0.1352008  
\end{verbatim}

\begin{Shaded}
\begin{Highlighting}[]
\FunctionTok{nrow}\NormalTok{(tbl\_daily\_stock\_return)}
\end{Highlighting}
\end{Shaded}

\begin{verbatim}
[1] 21
\end{verbatim}

\begin{Shaded}
\begin{Highlighting}[]
\FunctionTok{ncol}\NormalTok{(tbl\_daily\_stock\_return)}
\end{Highlighting}
\end{Shaded}

\begin{verbatim}
[1] 3
\end{verbatim}

\begin{Shaded}
\begin{Highlighting}[]
\FunctionTok{names}\NormalTok{(tbl\_daily\_stock\_return)}
\end{Highlighting}
\end{Shaded}

\begin{verbatim}
[1] "date"  "firm1" "firm2"
\end{verbatim}

\subsection{dplyrによるデータ操作}\label{dplyrux306bux3088ux308bux30c7ux30fcux30bfux64cdux4f5c}

\subsubsection{列を選ぶ}\label{ux5217ux3092ux9078ux3076}

\begin{Shaded}
\begin{Highlighting}[]
\NormalTok{tbl\_daily\_stock\_return }\SpecialCharTok{|\textgreater{}}
  \FunctionTok{select}\NormalTok{(date, firm1)}
\end{Highlighting}
\end{Shaded}

\begin{verbatim}
# A tibble: 21 x 2
   date            firm1
   <date>          <dbl>
 1 2020-04-01 -0.039482 
 2 2020-04-02  0.0059787
 3 2020-04-03  0.055787 
 4 2020-04-06  0.041935 
 5 2020-04-07 -0.020193 
 6 2020-04-08 -0.0022947
 7 2020-04-09  0.081707 
 8 2020-04-10  0.034561 
 9 2020-04-13  0.032601 
10 2020-04-14 -0.037663 
# i 11 more rows
\end{verbatim}

\subsubsection{行を抽出する}\label{ux884cux3092ux62bdux51faux3059ux308b}

\begin{Shaded}
\begin{Highlighting}[]
\NormalTok{tbl\_daily\_stock\_return }\SpecialCharTok{|\textgreater{}}
  \FunctionTok{filter}\NormalTok{(firm1 }\SpecialCharTok{\textgreater{}} \DecValTok{0}\NormalTok{)}
\end{Highlighting}
\end{Shaded}

\begin{verbatim}
# A tibble: 12 x 3
   date           firm1      firm2
   <date>         <dbl>      <dbl>
 1 2020-04-02 0.0059787 -0.0072488
 2 2020-04-03 0.055787  -0.017304 
 3 2020-04-06 0.041935   0.0021690
 4 2020-04-09 0.081707   0.078843 
 5 2020-04-10 0.034561   0.021536 
 6 2020-04-13 0.032601   0.13520  
 7 2020-04-15 0.0098798  0.011427 
 8 2020-04-17 0.17130    0.039463 
 9 2020-04-20 0.046325   0.037274 
10 2020-04-22 0.051543  -0.053818 
11 2020-04-27 0.0045375 -0.019615 
12 2020-04-30 0.044812  -0.013410 
\end{verbatim}

複数の条件を同時に指定できる。

\begin{Shaded}
\begin{Highlighting}[]
\NormalTok{tbl\_daily\_stock\_return }\SpecialCharTok{|\textgreater{}}
  \FunctionTok{filter}\NormalTok{(firm1 }\SpecialCharTok{\textgreater{}} \DecValTok{0}\NormalTok{, firm2 }\SpecialCharTok{\textgreater{}} \DecValTok{0}\NormalTok{)}
\end{Highlighting}
\end{Shaded}

\begin{verbatim}
# A tibble: 7 x 3
  date           firm1     firm2
  <date>         <dbl>     <dbl>
1 2020-04-06 0.041935  0.0021690
2 2020-04-09 0.081707  0.078843 
3 2020-04-10 0.034561  0.021536 
4 2020-04-13 0.032601  0.13520  
5 2020-04-15 0.0098798 0.011427 
6 2020-04-17 0.17130   0.039463 
7 2020-04-20 0.046325  0.037274 
\end{verbatim}

\subsubsection{列を追加する}\label{ux5217ux3092ux8ffdux52a0ux3059ux308b}

\begin{Shaded}
\begin{Highlighting}[]
\NormalTok{return\_data }\OtherTok{\textless{}{-}}\NormalTok{ tbl\_daily\_stock\_return }\SpecialCharTok{|\textgreater{}}
  \FunctionTok{mutate}\NormalTok{(}
    \AttributeTok{spread =}\NormalTok{ firm1 }\SpecialCharTok{{-}}\NormalTok{ firm2,}
    \AttributeTok{firm1\_gross\_return =} \DecValTok{1} \SpecialCharTok{+}\NormalTok{ firm1,}
    \AttributeTok{firm2\_gross\_return =} \DecValTok{1} \SpecialCharTok{+}\NormalTok{ firm2,}
    \AttributeTok{market\_direction =} \FunctionTok{case\_when}\NormalTok{(}
\NormalTok{      firm1 }\SpecialCharTok{\textgreater{}} \DecValTok{0} \SpecialCharTok{\&}\NormalTok{ firm2 }\SpecialCharTok{\textgreater{}} \DecValTok{0} \SpecialCharTok{\textasciitilde{}} \StringTok{"both\_positive"}\NormalTok{,}
\NormalTok{      firm1 }\SpecialCharTok{\textless{}} \DecValTok{0} \SpecialCharTok{\&}\NormalTok{ firm2 }\SpecialCharTok{\textless{}} \DecValTok{0} \SpecialCharTok{\textasciitilde{}} \StringTok{"both\_negative"}\NormalTok{,}
      \ConstantTok{TRUE} \SpecialCharTok{\textasciitilde{}} \StringTok{"mixed"}
\NormalTok{    )}
\NormalTok{  )}

\NormalTok{return\_data}
\end{Highlighting}
\end{Shaded}

\begin{verbatim}
# A tibble: 21 x 7
   date            firm1      firm2     spread firm1_gross_return
   <date>          <dbl>      <dbl>      <dbl>              <dbl>
 1 2020-04-01 -0.039482   0.076963  -0.11644              0.96052
 2 2020-04-02  0.0059787 -0.0072488  0.013228             1.0060 
 3 2020-04-03  0.055787  -0.017304   0.073091             1.0558 
 4 2020-04-06  0.041935   0.0021690  0.039766             1.0419 
 5 2020-04-07 -0.020193   0.075545  -0.095738             0.97981
 6 2020-04-08 -0.0022947 -0.096212   0.093917             0.99771
 7 2020-04-09  0.081707   0.078843   0.0028641            1.0817 
 8 2020-04-10  0.034561   0.021536   0.013025             1.0346 
 9 2020-04-13  0.032601   0.13520   -0.10260              1.0326 
10 2020-04-14 -0.037663   0.037979  -0.075642             0.96234
   firm2_gross_return market_direction
                <dbl> <chr>           
 1            1.0770  mixed           
 2            0.99275 mixed           
 3            0.98270 mixed           
 4            1.0022  both_positive   
 5            1.0755  mixed           
 6            0.90379 both_negative   
 7            1.0788  both_positive   
 8            1.0215  both_positive   
 9            1.1352  both_positive   
10            1.0380  mixed           
# i 11 more rows
\end{verbatim}

\subsubsection{並べ替える}\label{ux4e26ux3079ux66ffux3048ux308b}

\begin{Shaded}
\begin{Highlighting}[]
\NormalTok{return\_data }\SpecialCharTok{|\textgreater{}}
  \FunctionTok{arrange}\NormalTok{(firm1)}
\end{Highlighting}
\end{Shaded}

\begin{verbatim}
# A tibble: 21 x 7
   date            firm1       firm2     spread firm1_gross_return
   <date>          <dbl>       <dbl>      <dbl>              <dbl>
 1 2020-04-28 -0.072071  -0.046873   -0.025197             0.92793
 2 2020-04-24 -0.041378  -0.0041105  -0.037268             0.95862
 3 2020-04-01 -0.039482   0.076963   -0.11644              0.96052
 4 2020-04-14 -0.037663   0.037979   -0.075642             0.96234
 5 2020-04-23 -0.028875  -0.035484    0.0066092            0.97112
 6 2020-04-16 -0.027029  -0.053488    0.026460             0.97297
 7 2020-04-07 -0.020193   0.075545   -0.095738             0.97981
 8 2020-04-21 -0.014593  -0.00030564 -0.014288             0.98541
 9 2020-04-08 -0.0022947 -0.096212    0.093917             0.99771
10 2020-04-27  0.0045375 -0.019615    0.024152             1.0045 
   firm2_gross_return market_direction
                <dbl> <chr>           
 1            0.95313 both_negative   
 2            0.99589 both_negative   
 3            1.0770  mixed           
 4            1.0380  mixed           
 5            0.96452 both_negative   
 6            0.94651 both_negative   
 7            1.0755  mixed           
 8            0.99969 both_negative   
 9            0.90379 both_negative   
10            0.98039 mixed           
# i 11 more rows
\end{verbatim}

\begin{Shaded}
\begin{Highlighting}[]
\NormalTok{return\_data }\SpecialCharTok{|\textgreater{}}
  \FunctionTok{arrange}\NormalTok{(}\FunctionTok{desc}\NormalTok{(firm1))}
\end{Highlighting}
\end{Shaded}

\begin{verbatim}
# A tibble: 21 x 7
   date           firm1      firm2     spread firm1_gross_return
   <date>         <dbl>      <dbl>      <dbl>              <dbl>
 1 2020-04-17 0.17130    0.039463   0.13183               1.1713
 2 2020-04-09 0.081707   0.078843   0.0028641             1.0817
 3 2020-04-03 0.055787  -0.017304   0.073091              1.0558
 4 2020-04-22 0.051543  -0.053818   0.10536               1.0515
 5 2020-04-20 0.046325   0.037274   0.0090516             1.0463
 6 2020-04-30 0.044812  -0.013410   0.058222              1.0448
 7 2020-04-06 0.041935   0.0021690  0.039766              1.0419
 8 2020-04-10 0.034561   0.021536   0.013025              1.0346
 9 2020-04-13 0.032601   0.13520   -0.10260               1.0326
10 2020-04-15 0.0098798  0.011427  -0.0015473             1.0099
   firm2_gross_return market_direction
                <dbl> <chr>           
 1            1.0395  both_positive   
 2            1.0788  both_positive   
 3            0.98270 mixed           
 4            0.94618 mixed           
 5            1.0373  both_positive   
 6            0.98659 mixed           
 7            1.0022  both_positive   
 8            1.0215  both_positive   
 9            1.1352  both_positive   
10            1.0114  both_positive   
# i 11 more rows
\end{verbatim}

\subsubsection{要約する}\label{ux8981ux7d04ux3059ux308b}

\begin{Shaded}
\begin{Highlighting}[]
\NormalTok{return\_data }\SpecialCharTok{|\textgreater{}}
  \FunctionTok{summarise}\NormalTok{(}
    \AttributeTok{observations =} \FunctionTok{n}\NormalTok{(),}
    \AttributeTok{mean\_firm1 =} \FunctionTok{mean}\NormalTok{(firm1),}
    \AttributeTok{sd\_firm1 =} \FunctionTok{sd}\NormalTok{(firm1),}
    \AttributeTok{mean\_firm2 =} \FunctionTok{mean}\NormalTok{(firm2),}
    \AttributeTok{sd\_firm2 =} \FunctionTok{sd}\NormalTok{(firm2),}
    \AttributeTok{correlation =} \FunctionTok{cor}\NormalTok{(firm1, firm2)}
\NormalTok{  )}
\end{Highlighting}
\end{Shaded}

\begin{verbatim}
# A tibble: 1 x 6
  observations mean_firm1 sd_firm1 mean_firm2 sd_firm2 correlation
         <int>      <dbl>    <dbl>      <dbl>    <dbl>       <dbl>
1           21   0.014161 0.053840  0.0080253 0.054390     0.23349
\end{verbatim}

\subsubsection{グループごとに集計する}\label{ux30b0ux30ebux30fcux30d7ux3054ux3068ux306bux96c6ux8a08ux3059ux308b}

\begin{Shaded}
\begin{Highlighting}[]
\NormalTok{return\_data }\SpecialCharTok{|\textgreater{}}
  \FunctionTok{group\_by}\NormalTok{(market\_direction) }\SpecialCharTok{|\textgreater{}}
  \FunctionTok{summarise}\NormalTok{(}
    \AttributeTok{observations =} \FunctionTok{n}\NormalTok{(),}
    \AttributeTok{mean\_firm1 =} \FunctionTok{mean}\NormalTok{(firm1),}
    \AttributeTok{mean\_firm2 =} \FunctionTok{mean}\NormalTok{(firm2),}
    \AttributeTok{.groups =} \StringTok{"drop"}
\NormalTok{  ) }\SpecialCharTok{|\textgreater{}}
  \FunctionTok{arrange}\NormalTok{(}\FunctionTok{desc}\NormalTok{(observations))}
\end{Highlighting}
\end{Shaded}

\begin{verbatim}
# A tibble: 3 x 4
  market_direction observations mean_firm1 mean_firm2
  <chr>                   <int>      <dbl>      <dbl>
1 mixed                       8  0.0081649  0.0098864
2 both_positive               7  0.059758   0.046559 
3 both_negative               6 -0.031040  -0.039412 
\end{verbatim}

\subsubsection{件数を数える}\label{ux4ef6ux6570ux3092ux6570ux3048ux308b}

\begin{Shaded}
\begin{Highlighting}[]
\NormalTok{return\_data }\SpecialCharTok{|\textgreater{}}
  \FunctionTok{count}\NormalTok{(market\_direction, }\AttributeTok{sort =} \ConstantTok{TRUE}\NormalTok{)}
\end{Highlighting}
\end{Shaded}

\begin{verbatim}
# A tibble: 3 x 2
  market_direction     n
  <chr>            <int>
1 mixed                8
2 both_positive        7
3 both_negative        6
\end{verbatim}

\subsubsection{データを結合する}\label{ux30c7ux30fcux30bfux3092ux7d50ux5408ux3059ux308b}

\begin{Shaded}
\begin{Highlighting}[]
\NormalTok{firm\_master }\OtherTok{\textless{}{-}} \FunctionTok{tribble}\NormalTok{(}
  \SpecialCharTok{\textasciitilde{}}\NormalTok{firm,  }\SpecialCharTok{\textasciitilde{}}\NormalTok{firm\_name, }\SpecialCharTok{\textasciitilde{}}\NormalTok{industry,}
  \StringTok{"firm1"}\NormalTok{, }\StringTok{"Firm A"}\NormalTok{,  }\StringTok{"Machinery"}\NormalTok{,}
  \StringTok{"firm2"}\NormalTok{, }\StringTok{"Firm B"}\NormalTok{,  }\StringTok{"Chemicals"}
\NormalTok{)}

\NormalTok{return\_summary }\OtherTok{\textless{}{-}}\NormalTok{ return\_data }\SpecialCharTok{|\textgreater{}}
  \FunctionTok{summarise}\NormalTok{(}
    \AttributeTok{firm1 =} \FunctionTok{mean}\NormalTok{(firm1),}
    \AttributeTok{firm2 =} \FunctionTok{mean}\NormalTok{(firm2)}
\NormalTok{  ) }\SpecialCharTok{|\textgreater{}}
  \FunctionTok{pivot\_longer}\NormalTok{(}
    \AttributeTok{cols =} \FunctionTok{everything}\NormalTok{(),}
    \AttributeTok{names\_to =} \StringTok{"firm"}\NormalTok{,}
    \AttributeTok{values\_to =} \StringTok{"mean\_return"}
\NormalTok{  )}

\NormalTok{return\_summary }\SpecialCharTok{|\textgreater{}}
  \FunctionTok{left\_join}\NormalTok{(firm\_master, }\AttributeTok{by =} \StringTok{"firm"}\NormalTok{)}
\end{Highlighting}
\end{Shaded}

\begin{verbatim}
# A tibble: 2 x 4
  firm  mean_return firm_name industry 
  <chr>       <dbl> <chr>     <chr>    
1 firm1   0.014161  Firm A    Machinery
2 firm2   0.0080253 Firm B    Chemicals
\end{verbatim}

\subsection{stringrによる文字列処理}\label{stringrux306bux3088ux308bux6587ux5b57ux5217ux51e6ux7406}

文字列処理では、関数名が\texttt{str\_}で始まる\texttt{stringr}を使用する。

\begin{Shaded}
\begin{Highlighting}[]
\NormalTok{company\_names }\OtherTok{\textless{}{-}} \FunctionTok{tibble}\NormalTok{(}
  \AttributeTok{raw\_name =} \FunctionTok{c}\NormalTok{(}
    \StringTok{"  Alpha Holdings Co., Ltd. "}\NormalTok{,}
    \StringTok{"Beta Financial Group"}\NormalTok{,}
    \StringTok{"Gamma Industries"}
\NormalTok{  )}
\NormalTok{)}

\NormalTok{company\_names }\OtherTok{\textless{}{-}}\NormalTok{ company\_names }\SpecialCharTok{|\textgreater{}}
  \FunctionTok{mutate}\NormalTok{(}
    \AttributeTok{clean\_name =} \FunctionTok{str\_squish}\NormalTok{(raw\_name),}
    \AttributeTok{lower\_name =} \FunctionTok{str\_to\_lower}\NormalTok{(clean\_name),}
    \AttributeTok{name\_length =} \FunctionTok{str\_length}\NormalTok{(clean\_name),}
    \AttributeTok{contains\_group =} \FunctionTok{str\_detect}\NormalTok{(clean\_name, }\FunctionTok{regex}\NormalTok{(}\StringTok{"group"}\NormalTok{, }\AttributeTok{ignore\_case =} \ConstantTok{TRUE}\NormalTok{)),}
    \AttributeTok{first\_word =} \FunctionTok{str\_extract}\NormalTok{(clean\_name, }\StringTok{"\^{}[A{-}Za{-}z]+"}\NormalTok{)}
\NormalTok{  )}

\NormalTok{company\_names}
\end{Highlighting}
\end{Shaded}

\begin{verbatim}
# A tibble: 3 x 6
  raw_name                      clean_name              
  <chr>                         <chr>                   
1 "  Alpha Holdings Co., Ltd. " Alpha Holdings Co., Ltd.
2 "Beta Financial Group"        Beta Financial Group    
3 "Gamma Industries"            Gamma Industries        
  lower_name               name_length contains_group first_word
  <chr>                          <int> <lgl>          <chr>     
1 alpha holdings co., ltd.          24 FALSE          Alpha     
2 beta financial group              20 TRUE           Beta      
3 gamma industries                  16 FALSE          Gamma     
\end{verbatim}

\subsubsection{置換と結合}\label{ux7f6eux63dbux3068ux7d50ux5408}

\begin{Shaded}
\begin{Highlighting}[]
\NormalTok{company\_names }\SpecialCharTok{|\textgreater{}}
  \FunctionTok{mutate}\NormalTok{(}
    \AttributeTok{short\_name =} \FunctionTok{str\_remove}\NormalTok{(clean\_name, }\FunctionTok{regex}\NormalTok{(}\StringTok{" co}\SpecialCharTok{\textbackslash{}\textbackslash{}}\StringTok{., ltd}\SpecialCharTok{\textbackslash{}\textbackslash{}}\StringTok{.$"}\NormalTok{, }\AttributeTok{ignore\_case =} \ConstantTok{TRUE}\NormalTok{)),}
    \AttributeTok{display\_name =} \FunctionTok{str\_c}\NormalTok{(first\_word, }\StringTok{" / "}\NormalTok{, name\_length, }\StringTok{" characters"}\NormalTok{)}
\NormalTok{  ) }\SpecialCharTok{|\textgreater{}}
  \FunctionTok{select}\NormalTok{(clean\_name, short\_name, display\_name)}
\end{Highlighting}
\end{Shaded}

\begin{verbatim}
# A tibble: 3 x 3
  clean_name               short_name           display_name         
  <chr>                    <chr>                <chr>                
1 Alpha Holdings Co., Ltd. Alpha Holdings       Alpha / 24 characters
2 Beta Financial Group     Beta Financial Group Beta / 20 characters 
3 Gamma Industries         Gamma Industries     Gamma / 16 characters
\end{verbatim}

\subsection{tidyrによるデータ整形}\label{tidyrux306bux3088ux308bux30c7ux30fcux30bfux6574ux5f62}

\subsubsection{wide形式からlong形式へ変換する}\label{wideux5f62ux5f0fux304bux3089longux5f62ux5f0fux3078ux5909ux63dbux3059ux308b}

元のデータでは、\texttt{firm1}と\texttt{firm2}が別々の列に格納されている。\texttt{pivot\_longer()}を使うと、企業名とリターンをそれぞれ一つの列へまとめられる。

\begin{Shaded}
\begin{Highlighting}[]
\NormalTok{return\_long }\OtherTok{\textless{}{-}}\NormalTok{ tbl\_daily\_stock\_return }\SpecialCharTok{|\textgreater{}}
  \FunctionTok{pivot\_longer}\NormalTok{(}
    \AttributeTok{cols =} \FunctionTok{c}\NormalTok{(firm1, firm2),}
    \AttributeTok{names\_to =} \StringTok{"firm"}\NormalTok{,}
    \AttributeTok{values\_to =} \StringTok{"return"}
\NormalTok{  )}

\NormalTok{return\_long}
\end{Highlighting}
\end{Shaded}

\begin{verbatim}
# A tibble: 42 x 3
   date       firm      return
   <date>     <chr>      <dbl>
 1 2020-04-01 firm1 -0.039482 
 2 2020-04-01 firm2  0.076963 
 3 2020-04-02 firm1  0.0059787
 4 2020-04-02 firm2 -0.0072488
 5 2020-04-03 firm1  0.055787 
 6 2020-04-03 firm2 -0.017304 
 7 2020-04-06 firm1  0.041935 
 8 2020-04-06 firm2  0.0021690
 9 2020-04-07 firm1 -0.020193 
10 2020-04-07 firm2  0.075545 
# i 32 more rows
\end{verbatim}

\subsubsection{long形式からwide形式へ戻す}\label{longux5f62ux5f0fux304bux3089wideux5f62ux5f0fux3078ux623bux3059}

\begin{Shaded}
\begin{Highlighting}[]
\NormalTok{return\_long }\SpecialCharTok{|\textgreater{}}
  \FunctionTok{pivot\_wider}\NormalTok{(}
    \AttributeTok{names\_from =}\NormalTok{ firm,}
    \AttributeTok{values\_from =}\NormalTok{ return}
\NormalTok{  )}
\end{Highlighting}
\end{Shaded}

\begin{verbatim}
# A tibble: 21 x 3
   date            firm1      firm2
   <date>          <dbl>      <dbl>
 1 2020-04-01 -0.039482   0.076963 
 2 2020-04-02  0.0059787 -0.0072488
 3 2020-04-03  0.055787  -0.017304 
 4 2020-04-06  0.041935   0.0021690
 5 2020-04-07 -0.020193   0.075545 
 6 2020-04-08 -0.0022947 -0.096212 
 7 2020-04-09  0.081707   0.078843 
 8 2020-04-10  0.034561   0.021536 
 9 2020-04-13  0.032601   0.13520  
10 2020-04-14 -0.037663   0.037979 
# i 11 more rows
\end{verbatim}

\subsubsection{文字列を複数列へ分ける}\label{ux6587ux5b57ux5217ux3092ux8907ux6570ux5217ux3078ux5206ux3051ux308b}

\begin{Shaded}
\begin{Highlighting}[]
\NormalTok{security\_codes }\OtherTok{\textless{}{-}} \FunctionTok{tibble}\NormalTok{(}
  \AttributeTok{security\_id =} \FunctionTok{c}\NormalTok{(}\StringTok{"JP{-}1001"}\NormalTok{, }\StringTok{"JP{-}1002"}\NormalTok{, }\StringTok{"US{-}2001"}\NormalTok{)}
\NormalTok{)}

\NormalTok{security\_codes }\SpecialCharTok{|\textgreater{}}
  \FunctionTok{separate\_wider\_delim}\NormalTok{(}
\NormalTok{    security\_id,}
    \AttributeTok{delim =} \StringTok{"{-}"}\NormalTok{,}
    \AttributeTok{names =} \FunctionTok{c}\NormalTok{(}\StringTok{"country"}\NormalTok{, }\StringTok{"code"}\NormalTok{)}
\NormalTok{  )}
\end{Highlighting}
\end{Shaded}

\begin{verbatim}
# A tibble: 3 x 2
  country code 
  <chr>   <chr>
1 JP      1001 
2 JP      1002 
3 US      2001 
\end{verbatim}

\subsubsection{欠損値を処理する}\label{ux6b20ux640dux5024ux3092ux51e6ux7406ux3059ux308b}

\begin{Shaded}
\begin{Highlighting}[]
\NormalTok{missing\_data }\OtherTok{\textless{}{-}} \FunctionTok{tibble}\NormalTok{(}
  \AttributeTok{firm =} \FunctionTok{c}\NormalTok{(}\StringTok{"A"}\NormalTok{, }\StringTok{"B"}\NormalTok{, }\StringTok{"C"}\NormalTok{),}
  \AttributeTok{return =} \FunctionTok{c}\NormalTok{(}\FloatTok{0.02}\NormalTok{, }\ConstantTok{NA}\NormalTok{, }\SpecialCharTok{{-}}\FloatTok{0.01}\NormalTok{)}
\NormalTok{)}

\NormalTok{missing\_data }\SpecialCharTok{|\textgreater{}}
  \FunctionTok{replace\_na}\NormalTok{(}\FunctionTok{list}\NormalTok{(}\AttributeTok{return =} \DecValTok{0}\NormalTok{))}
\end{Highlighting}
\end{Shaded}

\begin{verbatim}
# A tibble: 3 x 2
  firm  return
  <chr>  <dbl>
1 A       0.02
2 B       0   
3 C      -0.01
\end{verbatim}

\begin{Shaded}
\begin{Highlighting}[]
\NormalTok{missing\_data }\SpecialCharTok{|\textgreater{}}
  \FunctionTok{drop\_na}\NormalTok{()}
\end{Highlighting}
\end{Shaded}

\begin{verbatim}
# A tibble: 2 x 2
  firm  return
  <chr>  <dbl>
1 A       0.02
2 C      -0.01
\end{verbatim}

\subsection{ggplot2による可視化}\label{ggplot2ux306bux3088ux308bux53efux8996ux5316}

\texttt{ggplot2}では、データを\texttt{ggplot()}へ渡し、\texttt{aes()}で変数と視覚要素の対応を指定し、\texttt{geom\_*()}で図形を追加する。

\subsubsection{折れ線グラフ}\label{ux6298ux308cux7ddaux30b0ux30e9ux30d5}

\begin{Shaded}
\begin{Highlighting}[]
\NormalTok{(return\_long }\SpecialCharTok{|\textgreater{}}
  \FunctionTok{ggplot}\NormalTok{(}\FunctionTok{aes}\NormalTok{(}\AttributeTok{x =}\NormalTok{ date, }\AttributeTok{y =}\NormalTok{ return, }\AttributeTok{linetype =}\NormalTok{ firm)) }\SpecialCharTok{+}
  \FunctionTok{geom\_line}\NormalTok{(}\AttributeTok{linewidth =} \FloatTok{0.7}\NormalTok{) }\SpecialCharTok{+}
  \FunctionTok{labs}\NormalTok{(}
    \AttributeTok{x =} \StringTok{"Date"}\NormalTok{,}
    \AttributeTok{y =} \StringTok{"Daily Return"}\NormalTok{,}
    \AttributeTok{linetype =} \StringTok{"Firm"}
\NormalTok{  )) }\SpecialCharTok{|\textgreater{}} \FunctionTok{add\_lagrange\_theme\_LFL}\NormalTok{()}
\end{Highlighting}
\end{Shaded}

\pandocbounded{\includegraphics[keepaspectratio]{part1_empirical/ch01_r_basics_tidyverse_files/figure-pdf/ggplot-line-1.pdf}}

\subsubsection{ヒストグラム}\label{ux30d2ux30b9ux30c8ux30b0ux30e9ux30e0}

\begin{Shaded}
\begin{Highlighting}[]
\NormalTok{(return\_long }\SpecialCharTok{|\textgreater{}}
  \FunctionTok{ggplot}\NormalTok{(}\FunctionTok{aes}\NormalTok{(}\AttributeTok{x =}\NormalTok{ return)) }\SpecialCharTok{+}
  \FunctionTok{geom\_histogram}\NormalTok{(}\AttributeTok{bins =} \DecValTok{10}\NormalTok{, }\AttributeTok{boundary =} \DecValTok{0}\NormalTok{) }\SpecialCharTok{+}
  \FunctionTok{facet\_wrap}\NormalTok{(}\FunctionTok{vars}\NormalTok{(firm), }\AttributeTok{ncol =} \DecValTok{1}\NormalTok{) }\SpecialCharTok{+}
  \FunctionTok{labs}\NormalTok{(}
    \AttributeTok{x =} \StringTok{"Daily Return"}\NormalTok{,}
    \AttributeTok{y =} \StringTok{"Count"}
\NormalTok{  )) }\SpecialCharTok{|\textgreater{}}
  \FunctionTok{add\_lagrange\_theme\_LFL}\NormalTok{()}
\end{Highlighting}
\end{Shaded}

\pandocbounded{\includegraphics[keepaspectratio]{part1_empirical/ch01_r_basics_tidyverse_files/figure-pdf/ggplot-histogram-1.pdf}}

\subsubsection{密度関数}\label{ux5bc6ux5ea6ux95a2ux6570}

\begin{Shaded}
\begin{Highlighting}[]
\NormalTok{(return\_long }\SpecialCharTok{|\textgreater{}}
  \FunctionTok{ggplot}\NormalTok{(}\FunctionTok{aes}\NormalTok{(}\AttributeTok{x =}\NormalTok{ return, }\AttributeTok{linetype =}\NormalTok{ firm)) }\SpecialCharTok{+}
  \FunctionTok{geom\_density}\NormalTok{(}\AttributeTok{linewidth =} \FloatTok{0.8}\NormalTok{) }\SpecialCharTok{+}
  \FunctionTok{labs}\NormalTok{(}
    \AttributeTok{x =} \StringTok{"Daily Return"}\NormalTok{,}
    \AttributeTok{y =} \StringTok{"Density"}\NormalTok{,}
    \AttributeTok{linetype =} \StringTok{"Firm"}
\NormalTok{  )) }\SpecialCharTok{|\textgreater{}}
  \FunctionTok{add\_lagrange\_theme\_LFL}\NormalTok{()}
\end{Highlighting}
\end{Shaded}

\pandocbounded{\includegraphics[keepaspectratio]{part1_empirical/ch01_r_basics_tidyverse_files/figure-pdf/ggplot-density-1.pdf}}

\subsubsection{正規分布の密度と累積分布}\label{ux6b63ux898fux5206ux5e03ux306eux5bc6ux5ea6ux3068ux7d2fux7a4dux5206ux5e03}

\begin{Shaded}
\begin{Highlighting}[]
\NormalTok{normal\_curve }\OtherTok{\textless{}{-}} \FunctionTok{tibble}\NormalTok{(}
  \AttributeTok{x =} \FunctionTok{seq}\NormalTok{(}\SpecialCharTok{{-}}\DecValTok{4}\NormalTok{, }\DecValTok{4}\NormalTok{, }\AttributeTok{length.out =} \DecValTok{401}\NormalTok{),}
  \AttributeTok{density =} \FunctionTok{dnorm}\NormalTok{(x),}
  \AttributeTok{cumulative\_probability =} \FunctionTok{pnorm}\NormalTok{(x)}
\NormalTok{)}

\NormalTok{(normal\_curve }\SpecialCharTok{|\textgreater{}}
  \FunctionTok{pivot\_longer}\NormalTok{(}
    \AttributeTok{cols =} \FunctionTok{c}\NormalTok{(density, cumulative\_probability),}
    \AttributeTok{names\_to =} \StringTok{"function\_type"}\NormalTok{,}
    \AttributeTok{values\_to =} \StringTok{"value"}
\NormalTok{  ) }\SpecialCharTok{|\textgreater{}}
  \FunctionTok{ggplot}\NormalTok{(}\FunctionTok{aes}\NormalTok{(}\AttributeTok{x =}\NormalTok{ x, }\AttributeTok{y =}\NormalTok{ value)) }\SpecialCharTok{+}
  \FunctionTok{geom\_line}\NormalTok{(}\AttributeTok{linewidth =} \FloatTok{0.8}\NormalTok{) }\SpecialCharTok{+}
  \FunctionTok{facet\_wrap}\NormalTok{(}\FunctionTok{vars}\NormalTok{(function\_type), }\AttributeTok{scales =} \StringTok{"free\_y"}\NormalTok{, }\AttributeTok{ncol =} \DecValTok{1}\NormalTok{) }\SpecialCharTok{+}
  \FunctionTok{labs}\NormalTok{(}
    \AttributeTok{x =} \StringTok{"x"}\NormalTok{,}
    \AttributeTok{y =} \ConstantTok{NULL}
\NormalTok{  )) }\SpecialCharTok{|\textgreater{}}
  \FunctionTok{add\_lagrange\_theme\_LFL}\NormalTok{()}
\end{Highlighting}
\end{Shaded}

\pandocbounded{\includegraphics[keepaspectratio]{part1_empirical/ch01_r_basics_tidyverse_files/figure-pdf/ggplot-normal-functions-1.pdf}}

\subsubsection{シミュレーション結果を比較する}\label{ux30b7ux30dfux30e5ux30ecux30fcux30b7ux30e7ux30f3ux7d50ux679cux3092ux6bd4ux8f03ux3059ux308b}

\begin{Shaded}
\begin{Highlighting}[]
\NormalTok{simulated\_distributions }\OtherTok{\textless{}{-}} \FunctionTok{tibble}\NormalTok{(}
  \AttributeTok{distribution =} \FunctionTok{c}\NormalTok{(}\StringTok{"normal"}\NormalTok{, }\StringTok{"uniform"}\NormalTok{, }\StringTok{"exponential"}\NormalTok{),}
  \AttributeTok{values =} \FunctionTok{list}\NormalTok{(}\FunctionTok{rnorm}\NormalTok{(}\DecValTok{10000}\NormalTok{), }\FunctionTok{runif}\NormalTok{(}\DecValTok{10000}\NormalTok{), }\FunctionTok{rexp}\NormalTok{(}\DecValTok{10000}\NormalTok{))}
\NormalTok{) }\SpecialCharTok{|\textgreater{}}
\NormalTok{  tidyr}\SpecialCharTok{::}\FunctionTok{unnest}\NormalTok{(values) }

\NormalTok{p }\OtherTok{\textless{}{-}}\NormalTok{ simulated\_distributions }\SpecialCharTok{|\textgreater{}}
  \FunctionTok{ggplot}\NormalTok{(}\FunctionTok{aes}\NormalTok{(}\AttributeTok{x =}\NormalTok{ values)) }\SpecialCharTok{+}
  \FunctionTok{geom\_histogram}\NormalTok{(}\AttributeTok{bins =} \DecValTok{30}\NormalTok{) }\SpecialCharTok{+}
  \FunctionTok{facet\_wrap}\NormalTok{(}\FunctionTok{vars}\NormalTok{(distribution), }\AttributeTok{scales =} \StringTok{"free"}\NormalTok{) }\SpecialCharTok{+}
  \FunctionTok{labs}\NormalTok{(}
    \AttributeTok{x =} \StringTok{"Simulated Value"}\NormalTok{,}
    \AttributeTok{y =} \StringTok{"Count"}
\NormalTok{  )}

\FunctionTok{add\_lagrange\_theme\_LFL}\NormalTok{(p)}
\end{Highlighting}
\end{Shaded}

\pandocbounded{\includegraphics[keepaspectratio]{part1_empirical/ch01_r_basics_tidyverse_files/figure-pdf/ggplot-simulated-distributions-1.pdf}}

\section{1.3
関数型プログラミングと並列計算}\label{ux95a2ux6570ux578bux30d7ux30edux30b0ux30e9ux30dfux30f3ux30b0ux3068ux4e26ux5217ux8a08ux7b97}

\subsection{purrrによる反復処理}\label{purrrux306bux3088ux308bux53cdux5fa9ux51e6ux7406}

\subsubsection{map()と型を指定するmap系関数}\label{mapux3068ux578bux3092ux6307ux5b9aux3059ux308bmapux7cfbux95a2ux6570}

\texttt{map()}は結果をリストとして返す。数値ベクトルが必要な場合は\texttt{map\_dbl()}を用いる。

\begin{Shaded}
\begin{Highlighting}[]
\NormalTok{discount\_rates }\OtherTok{\textless{}{-}} \FunctionTok{c}\NormalTok{(}\FloatTok{0.01}\NormalTok{, }\FloatTok{0.02}\NormalTok{, }\FloatTok{0.03}\NormalTok{, }\FloatTok{0.04}\NormalTok{)}

\NormalTok{discount\_rates }\SpecialCharTok{|\textgreater{}}
  \FunctionTok{map}\NormalTok{(\textbackslash{}(rate) }\FunctionTok{calculate\_present\_value}\NormalTok{(}\DecValTok{100}\NormalTok{, rate))}
\end{Highlighting}
\end{Shaded}

\begin{verbatim}
[[1]]
[1] 99.0099

[[2]]
[1] 98.03922

[[3]]
[1] 97.08738

[[4]]
[1] 96.15385
\end{verbatim}

\begin{Shaded}
\begin{Highlighting}[]
\NormalTok{present\_value\_table }\OtherTok{\textless{}{-}} \FunctionTok{tibble}\NormalTok{(}
  \AttributeTok{discount\_rate =}\NormalTok{ discount\_rates,}
  \AttributeTok{present\_value =}\NormalTok{ discount\_rates }\SpecialCharTok{|\textgreater{}}
    \FunctionTok{map\_dbl}\NormalTok{(\textbackslash{}(rate) }\FunctionTok{calculate\_present\_value}\NormalTok{(}\DecValTok{100}\NormalTok{, rate))}
\NormalTok{)}

\NormalTok{present\_value\_table}
\end{Highlighting}
\end{Shaded}

\begin{verbatim}
# A tibble: 4 x 2
  discount_rate present_value
          <dbl>         <dbl>
1          0.01        99.010
2          0.02        98.039
3          0.03        97.087
4          0.04        96.154
\end{verbatim}

\subsubsection{map2()で二つのベクトルを同時に処理する}\label{map2ux3067ux4e8cux3064ux306eux30d9ux30afux30c8ux30ebux3092ux540cux6642ux306bux51e6ux7406ux3059ux308b}

\begin{Shaded}
\begin{Highlighting}[]
\NormalTok{cash\_flows }\OtherTok{\textless{}{-}} \FunctionTok{c}\NormalTok{(}\DecValTok{100}\NormalTok{, }\DecValTok{120}\NormalTok{, }\DecValTok{150}\NormalTok{)}
\NormalTok{years }\OtherTok{\textless{}{-}} \FunctionTok{c}\NormalTok{(}\DecValTok{1}\NormalTok{, }\DecValTok{2}\NormalTok{, }\DecValTok{3}\NormalTok{)}

\FunctionTok{map2\_dbl}\NormalTok{(}
\NormalTok{  cash\_flows,}
\NormalTok{  years,}
\NormalTok{  \textbackslash{}(cash\_flow, year) }\FunctionTok{calculate\_present\_value}\NormalTok{(cash\_flow, }\FloatTok{0.02}\NormalTok{, year)}
\NormalTok{)}
\end{Highlighting}
\end{Shaded}

\begin{verbatim}
[1]  98.03922 115.34025 141.34835
\end{verbatim}

\subsubsection{pmap()で複数の列を処理する}\label{pmapux3067ux8907ux6570ux306eux5217ux3092ux51e6ux7406ux3059ux308b}

\begin{Shaded}
\begin{Highlighting}[]
\NormalTok{cash\_flow\_scenarios }\OtherTok{\textless{}{-}} \FunctionTok{tribble}\NormalTok{(}
  \SpecialCharTok{\textasciitilde{}}\NormalTok{cash\_flow, }\SpecialCharTok{\textasciitilde{}}\NormalTok{discount\_rate, }\SpecialCharTok{\textasciitilde{}}\NormalTok{years,}
         \DecValTok{100}\NormalTok{,           }\FloatTok{0.01}\NormalTok{,      }\DecValTok{1}\NormalTok{,}
         \DecValTok{120}\NormalTok{,           }\FloatTok{0.02}\NormalTok{,      }\DecValTok{2}\NormalTok{,}
         \DecValTok{150}\NormalTok{,           }\FloatTok{0.03}\NormalTok{,      }\DecValTok{3}
\NormalTok{)}

\NormalTok{cash\_flow\_scenarios }\SpecialCharTok{|\textgreater{}}
  \FunctionTok{mutate}\NormalTok{(}
    \AttributeTok{present\_value =} \FunctionTok{pmap\_dbl}\NormalTok{(}
      \FunctionTok{list}\NormalTok{(cash\_flow, discount\_rate, years),}
\NormalTok{      \textbackslash{}(cash\_flow, discount\_rate, years) \{}
        \FunctionTok{calculate\_present\_value}\NormalTok{(cash\_flow, discount\_rate, years)}
\NormalTok{      \}}
\NormalTok{    )}
\NormalTok{  )}
\end{Highlighting}
\end{Shaded}

\begin{verbatim}
# A tibble: 3 x 4
  cash_flow discount_rate years present_value
      <dbl>         <dbl> <dbl>         <dbl>
1       100          0.01     1        99.010
2       120          0.02     2       115.34 
3       150          0.03     3       137.27 
\end{verbatim}

\subsubsection{複数列へ同じ集計を適用する}\label{ux8907ux6570ux5217ux3078ux540cux3058ux96c6ux8a08ux3092ux9069ux7528ux3059ux308b}

\begin{Shaded}
\begin{Highlighting}[]
\NormalTok{return\_columns }\OtherTok{\textless{}{-}} \FunctionTok{c}\NormalTok{(}\StringTok{"firm1"}\NormalTok{, }\StringTok{"firm2"}\NormalTok{)}

\NormalTok{return\_columns }\SpecialCharTok{|\textgreater{}}
  \FunctionTok{map\_dfr}\NormalTok{(}
\NormalTok{    \textbackslash{}(column\_name) \{}
\NormalTok{      values }\OtherTok{\textless{}{-}}\NormalTok{ tbl\_daily\_stock\_return[[column\_name]]}

      \FunctionTok{tibble}\NormalTok{(}
        \AttributeTok{firm =}\NormalTok{ column\_name,}
        \AttributeTok{observations =} \FunctionTok{length}\NormalTok{(values),}
        \AttributeTok{mean =} \FunctionTok{mean}\NormalTok{(values),}
        \AttributeTok{standard\_deviation =} \FunctionTok{sd}\NormalTok{(values),}
        \AttributeTok{minimum =} \FunctionTok{min}\NormalTok{(values),}
        \AttributeTok{maximum =} \FunctionTok{max}\NormalTok{(values)}
\NormalTok{      )}
\NormalTok{    \}}
\NormalTok{  )}
\end{Highlighting}
\end{Shaded}

\begin{verbatim}
# A tibble: 2 x 6
  firm  observations      mean standard_deviation   minimum maximum
  <chr>        <int>     <dbl>              <dbl>     <dbl>   <dbl>
1 firm1           21 0.014161            0.053840 -0.072071 0.17130
2 firm2           21 0.0080253           0.054390 -0.096212 0.13520
\end{verbatim}

\subsubsection{複数の確率分布を同じ形式で生成する}\label{ux8907ux6570ux306eux78baux7387ux5206ux5e03ux3092ux540cux3058ux5f62ux5f0fux3067ux751fux6210ux3059ux308b}

\begin{Shaded}
\begin{Highlighting}[]
\FunctionTok{set\_random\_seed\_LFL}\NormalTok{(}\DecValTok{1234}\NormalTok{)}

\NormalTok{distribution\_specs }\OtherTok{\textless{}{-}} \FunctionTok{tribble}\NormalTok{(}
  \SpecialCharTok{\textasciitilde{}}\NormalTok{distribution, }\SpecialCharTok{\textasciitilde{}}\NormalTok{generator,}
  \StringTok{"uniform"}\NormalTok{,      \textbackslash{}(n) }\FunctionTok{runif}\NormalTok{(n, }\AttributeTok{min =} \SpecialCharTok{{-}}\DecValTok{1}\NormalTok{, }\AttributeTok{max =} \DecValTok{1}\NormalTok{),}
  \StringTok{"normal"}\NormalTok{,       \textbackslash{}(n) }\FunctionTok{rnorm}\NormalTok{(n, }\AttributeTok{mean =} \DecValTok{0}\NormalTok{, }\AttributeTok{sd =} \DecValTok{1}\NormalTok{),}
  \StringTok{"poisson"}\NormalTok{,      \textbackslash{}(n) }\FunctionTok{rpois}\NormalTok{(n, }\AttributeTok{lambda =} \DecValTok{3}\NormalTok{)}
\NormalTok{)}

\NormalTok{simulated\_distributions }\OtherTok{\textless{}{-}}\NormalTok{ distribution\_specs }\SpecialCharTok{|\textgreater{}}
  \FunctionTok{mutate}\NormalTok{(}\AttributeTok{values =} \FunctionTok{map}\NormalTok{(generator, \textbackslash{}(generator) }\FunctionTok{generator}\NormalTok{(}\DecValTok{1000}\NormalTok{))) }\SpecialCharTok{|\textgreater{}}
  \FunctionTok{select}\NormalTok{(}\SpecialCharTok{{-}}\NormalTok{generator) }\SpecialCharTok{|\textgreater{}}
  \FunctionTok{unnest\_longer}\NormalTok{(values)}

\NormalTok{simulated\_distributions}
\end{Highlighting}
\end{Shaded}

\begin{verbatim}
# A tibble: 3,000 x 2
   distribution    values
   <chr>            <dbl>
 1 uniform      -0.77259 
 2 uniform       0.24460 
 3 uniform       0.21855 
 4 uniform       0.24676 
 5 uniform       0.72183 
 6 uniform       0.28062 
 7 uniform      -0.98101 
 8 uniform      -0.53490 
 9 uniform       0.33217 
10 uniform       0.028502
# i 2,990 more rows
\end{verbatim}

\subsubsection{walk()で副作用だけを実行する}\label{walkux3067ux526fux4f5cux7528ux3060ux3051ux3092ux5b9fux884cux3059ux308b}

\texttt{walk()}は、返り値を保存するのではなく、表示などの副作用もしくはファイル出力などに実行するときに使う。

\begin{Shaded}
\begin{Highlighting}[]
\FunctionTok{c}\NormalTok{(}\StringTok{"firm1"}\NormalTok{, }\StringTok{"firm2"}\NormalTok{) }\SpecialCharTok{|\textgreater{}}
  \FunctionTok{walk}\NormalTok{(\textbackslash{}(firm) }\FunctionTok{print}\NormalTok{(}\FunctionTok{str\_c}\NormalTok{(}\StringTok{"集計対象: "}\NormalTok{, firm)))}
\end{Highlighting}
\end{Shaded}

\begin{verbatim}
[1] "集計対象: firm1"
[1] "集計対象: firm2"
\end{verbatim}

\subsection{nestによるグループ別データ管理}\label{nestux306bux3088ux308bux30b0ux30ebux30fcux30d7ux5225ux30c7ux30fcux30bfux7ba1ux7406}

\texttt{tidyr::nest()}は、グループごとのデータを一つのリスト列へ格納する。\texttt{group\_by()}が主に集計単位を指定するのに対し、\texttt{nest()}は銘柄別、業種別、ポートフォリオ別などのデータそのものを保持する。これにより、各グループへ同じ回帰分析、リスク計算、バックテストを適用しやすくなる。

\begin{Shaded}
\begin{Highlighting}[]
\NormalTok{tbl\_return\_long }\OtherTok{\textless{}{-}}\NormalTok{ tbl\_daily\_stock\_return }\SpecialCharTok{|\textgreater{}}
  \FunctionTok{pivot\_longer}\NormalTok{(}
    \AttributeTok{cols =} \FunctionTok{c}\NormalTok{(firm1, firm2),}
    \AttributeTok{names\_to =} \StringTok{"firm"}\NormalTok{,}
    \AttributeTok{values\_to =} \StringTok{"return"}
\NormalTok{  )}

\NormalTok{tbl\_nested }\OtherTok{\textless{}{-}}\NormalTok{ tbl\_return\_long }\SpecialCharTok{|\textgreater{}}
  \FunctionTok{nest}\NormalTok{(}\AttributeTok{tbl\_data =} \FunctionTok{c}\NormalTok{(date, return), }\AttributeTok{.by =}\NormalTok{ firm)}

\NormalTok{tbl\_nested}
\end{Highlighting}
\end{Shaded}

\begin{verbatim}
# A tibble: 2 x 2
  firm  tbl_data         
  <chr> <list>           
1 firm1 <tibble [21 x 2]>
2 firm2 <tibble [21 x 2]>
\end{verbatim}

リスト列\texttt{tbl\_data}の各要素はtibbleである。最初の企業のデータは次のように確認できる。

\begin{Shaded}
\begin{Highlighting}[]
\NormalTok{tbl\_nested}\SpecialCharTok{$}\NormalTok{tbl\_data[[}\DecValTok{1}\NormalTok{]]}
\end{Highlighting}
\end{Shaded}

\begin{verbatim}
# A tibble: 21 x 2
   date           return
   <date>          <dbl>
 1 2020-04-01 -0.039482 
 2 2020-04-02  0.0059787
 3 2020-04-03  0.055787 
 4 2020-04-06  0.041935 
 5 2020-04-07 -0.020193 
 6 2020-04-08 -0.0022947
 7 2020-04-09  0.081707 
 8 2020-04-10  0.034561 
 9 2020-04-13  0.032601 
10 2020-04-14 -0.037663 
# i 11 more rows
\end{verbatim}

\texttt{nest()}と\texttt{map()}を組み合わせると、グループ別の処理を一つのパイプラインとして記述できる。

\begin{Shaded}
\begin{Highlighting}[]
\NormalTok{tbl\_nested\_summary }\OtherTok{\textless{}{-}}\NormalTok{ tbl\_nested }\SpecialCharTok{|\textgreater{}}
  \FunctionTok{mutate}\NormalTok{(}
    \AttributeTok{tbl\_summary =} \FunctionTok{map}\NormalTok{(}
\NormalTok{      tbl\_data,}
\NormalTok{      \textbackslash{}(tbl\_data) }\FunctionTok{tibble}\NormalTok{(}
        \AttributeTok{num\_mean\_return =} \FunctionTok{mean}\NormalTok{(tbl\_data}\SpecialCharTok{$}\NormalTok{return),}
        \AttributeTok{num\_volatility =} \FunctionTok{sd}\NormalTok{(tbl\_data}\SpecialCharTok{$}\NormalTok{return)}
\NormalTok{      )}
\NormalTok{    )}
\NormalTok{  ) }\SpecialCharTok{|\textgreater{}}
  \FunctionTok{select}\NormalTok{(firm, tbl\_summary) }\SpecialCharTok{|\textgreater{}}
  \FunctionTok{unnest}\NormalTok{(tbl\_summary)}

\NormalTok{tbl\_nested\_summary}
\end{Highlighting}
\end{Shaded}

\begin{verbatim}
# A tibble: 2 x 3
  firm  num_mean_return num_volatility
  <chr>           <dbl>          <dbl>
1 firm1       0.014161        0.053840
2 firm2       0.0080253       0.054390
\end{verbatim}

\subsection{futureとfurrrによる並列計算}\label{futureux3068furrrux306bux3088ux308bux4e26ux5217ux8a08ux7b97}

\texttt{purrr::map()}は要素を順番に処理する。互いに独立した計算が多数ある場合、\texttt{future}と\texttt{furrr}を使うと、\texttt{map()}に近い記法のまま複数のRセッションへ処理を分配できる。Windowsでは\texttt{plan(multisession)}を使用する。

ここでは、モンテカルロ法による円周率の推定を1000回繰り返し、逐次処理と並列処理の実行時間を比較する。各試行では、正方形内に発生させた点のうち単位円内に入った割合を4倍する。

\begin{Shaded}
\begin{Highlighting}[]
\NormalTok{compute\_pi\_path }\OtherTok{\textless{}{-}} \ControlFlowTok{function}\NormalTok{(int\_path\_id, }\AttributeTok{int\_points =} \DecValTok{100000}\NormalTok{L) \{}
  \FunctionTok{set\_random\_seed\_LFL}\NormalTok{(}\DecValTok{100000}\NormalTok{L }\SpecialCharTok{+}\NormalTok{ int\_path\_id)}

\NormalTok{  vec\_x }\OtherTok{\textless{}{-}} \FunctionTok{runif}\NormalTok{(int\_points, }\AttributeTok{min =} \SpecialCharTok{{-}}\DecValTok{1}\NormalTok{, }\AttributeTok{max =} \DecValTok{1}\NormalTok{)}
\NormalTok{  vec\_y }\OtherTok{\textless{}{-}} \FunctionTok{runif}\NormalTok{(int\_points, }\AttributeTok{min =} \SpecialCharTok{{-}}\DecValTok{1}\NormalTok{, }\AttributeTok{max =} \DecValTok{1}\NormalTok{)}

  \FunctionTok{tibble}\NormalTok{(}
    \AttributeTok{int\_path\_id =}\NormalTok{ int\_path\_id,}
    \AttributeTok{num\_pi =} \DecValTok{4} \SpecialCharTok{*} \FunctionTok{mean}\NormalTok{(vec\_x}\SpecialCharTok{\^{}}\DecValTok{2} \SpecialCharTok{+}\NormalTok{ vec\_y}\SpecialCharTok{\^{}}\DecValTok{2} \SpecialCharTok{\textless{}=} \DecValTok{1}\NormalTok{)}
\NormalTok{  )}
\NormalTok{\}}
\end{Highlighting}
\end{Shaded}

1000本の試行を、利用可能なCPU数を上限として最大8個のチャンクへ分割する。

\begin{Shaded}
\begin{Highlighting}[]
\NormalTok{int\_paths }\OtherTok{\textless{}{-}} \DecValTok{1000}\NormalTok{L}
\NormalTok{int\_workers }\OtherTok{\textless{}{-}} \FunctionTok{min}\NormalTok{(}\DecValTok{8}\NormalTok{L, future}\SpecialCharTok{::}\FunctionTok{availableCores}\NormalTok{())}

\NormalTok{tbl\_chunks }\OtherTok{\textless{}{-}} \FunctionTok{tibble}\NormalTok{(}\AttributeTok{int\_path\_id =} \FunctionTok{seq\_len}\NormalTok{(int\_paths)) }\SpecialCharTok{|\textgreater{}}
  \FunctionTok{mutate}\NormalTok{(}
    \AttributeTok{int\_chunk\_id =} \FunctionTok{rep}\NormalTok{(}\FunctionTok{seq\_len}\NormalTok{(int\_workers), }\AttributeTok{length.out =}\NormalTok{ int\_paths)}
\NormalTok{  ) }\SpecialCharTok{|\textgreater{}}
  \FunctionTok{nest}\NormalTok{(}\AttributeTok{tbl\_path\_ids =}\NormalTok{ int\_path\_id, }\AttributeTok{.by =}\NormalTok{ int\_chunk\_id)}

\NormalTok{tbl\_chunks}
\end{Highlighting}
\end{Shaded}

\begin{verbatim}
# A tibble: 8 x 2
  int_chunk_id tbl_path_ids      
         <int> <list>            
1            1 <tibble [125 x 1]>
2            2 <tibble [125 x 1]>
3            3 <tibble [125 x 1]>
4            4 <tibble [125 x 1]>
5            5 <tibble [125 x 1]>
6            6 <tibble [125 x 1]>
7            7 <tibble [125 x 1]>
8            8 <tibble [125 x 1]>
\end{verbatim}

\begin{Shaded}
\begin{Highlighting}[]
\NormalTok{compute\_pi\_chunk }\OtherTok{\textless{}{-}} \ControlFlowTok{function}\NormalTok{(tbl\_path\_ids, }\AttributeTok{int\_points =} \DecValTok{100000}\NormalTok{L) \{}
\NormalTok{  tbl\_path\_ids}\SpecialCharTok{$}\NormalTok{int\_path\_id }\SpecialCharTok{|\textgreater{}}
    \FunctionTok{map\_dfr}\NormalTok{(compute\_pi\_path, }\AttributeTok{int\_points =}\NormalTok{ int\_points)}
\NormalTok{\}}
\end{Highlighting}
\end{Shaded}

まず\texttt{map()}による逐次処理を実行する。

\begin{Shaded}
\begin{Highlighting}[]
\NormalTok{future}\SpecialCharTok{::}\FunctionTok{plan}\NormalTok{(future}\SpecialCharTok{::}\NormalTok{sequential)}

\NormalTok{num\_time\_map }\OtherTok{\textless{}{-}} \FunctionTok{system.time}\NormalTok{(\{}
\NormalTok{  tbl\_result\_map }\OtherTok{\textless{}{-}}\NormalTok{ tbl\_chunks }\SpecialCharTok{|\textgreater{}}
    \FunctionTok{mutate}\NormalTok{(}
      \AttributeTok{tbl\_result =} \FunctionTok{map}\NormalTok{(tbl\_path\_ids, compute\_pi\_chunk)}
\NormalTok{    ) }\SpecialCharTok{|\textgreater{}}
    \FunctionTok{select}\NormalTok{(int\_chunk\_id, tbl\_result) }\SpecialCharTok{|\textgreater{}}
    \FunctionTok{unnest}\NormalTok{(tbl\_result)}
\NormalTok{\})[[}\StringTok{"elapsed"}\NormalTok{]]}
\end{Highlighting}
\end{Shaded}

次に\texttt{future\_map()}による並列処理を実行する。

\begin{Shaded}
\begin{Highlighting}[]
\NormalTok{future}\SpecialCharTok{::}\FunctionTok{plan}\NormalTok{(future}\SpecialCharTok{::}\NormalTok{multisession, }\AttributeTok{workers =}\NormalTok{ int\_workers)}

\NormalTok{num\_time\_future\_map }\OtherTok{\textless{}{-}} \FunctionTok{system.time}\NormalTok{(\{}
\NormalTok{  tbl\_result\_future }\OtherTok{\textless{}{-}}\NormalTok{ tbl\_chunks }\SpecialCharTok{|\textgreater{}}
    \FunctionTok{mutate}\NormalTok{(}
      \AttributeTok{tbl\_result =} \FunctionTok{future\_map}\NormalTok{(}
\NormalTok{        tbl\_path\_ids,}
\NormalTok{        compute\_pi\_chunk,}
        \AttributeTok{.options =}\NormalTok{ furrr}\SpecialCharTok{::}\FunctionTok{furrr\_options}\NormalTok{(}\AttributeTok{seed =} \ConstantTok{TRUE}\NormalTok{)}
\NormalTok{      )}
\NormalTok{    ) }\SpecialCharTok{|\textgreater{}}
    \FunctionTok{select}\NormalTok{(int\_chunk\_id, tbl\_result) }\SpecialCharTok{|\textgreater{}}
    \FunctionTok{unnest}\NormalTok{(tbl\_result)}
\NormalTok{\})[[}\StringTok{"elapsed"}\NormalTok{]]}

\NormalTok{future}\SpecialCharTok{::}\FunctionTok{plan}\NormalTok{(future}\SpecialCharTok{::}\NormalTok{sequential)}
\end{Highlighting}
\end{Shaded}

推定値と実行時間を比較する。

\begin{Shaded}
\begin{Highlighting}[]
\NormalTok{tbl\_pi\_benchmark }\OtherTok{\textless{}{-}} \FunctionTok{tribble}\NormalTok{(}
  \SpecialCharTok{\textasciitilde{}}\NormalTok{chr\_method, }\SpecialCharTok{\textasciitilde{}}\NormalTok{num\_pi, }\SpecialCharTok{\textasciitilde{}}\NormalTok{num\_elapsed\_seconds,}
  \StringTok{"map"}\NormalTok{, }\FunctionTok{mean}\NormalTok{(tbl\_result\_map}\SpecialCharTok{$}\NormalTok{num\_pi), num\_time\_map,}
  \FunctionTok{str\_c}\NormalTok{(}\StringTok{"future\_map ("}\NormalTok{, int\_workers, }\StringTok{" workers)"}\NormalTok{), }\FunctionTok{mean}\NormalTok{(tbl\_result\_future}\SpecialCharTok{$}\NormalTok{num\_pi), num\_time\_future\_map}
\NormalTok{) }\SpecialCharTok{|\textgreater{}}
  \FunctionTok{mutate}\NormalTok{(}
    \AttributeTok{num\_absolute\_error =} \FunctionTok{abs}\NormalTok{(num\_pi }\SpecialCharTok{{-}}\NormalTok{ pi),}
    \AttributeTok{num\_speedup =} \FunctionTok{first}\NormalTok{(num\_elapsed\_seconds) }\SpecialCharTok{/}\NormalTok{ num\_elapsed\_seconds}
\NormalTok{  )}

\NormalTok{tbl\_pi\_benchmark}
\end{Highlighting}
\end{Shaded}

\begin{verbatim}
# A tibble: 2 x 5
  chr_method             num_pi num_elapsed_seconds num_absolute_error
  <chr>                   <dbl>               <dbl>              <dbl>
1 map                    3.1414                6.31        0.00016693 
2 future_map (8 workers) 3.1415                2.55        0.000094774
  num_speedup
        <dbl>
1      1     
2      2.4745
\end{verbatim}

\begin{Shaded}
\begin{Highlighting}[]
\NormalTok{p }\OtherTok{\textless{}{-}}\NormalTok{ tbl\_pi\_benchmark }\SpecialCharTok{|\textgreater{}}
  \FunctionTok{ggplot}\NormalTok{(}\FunctionTok{aes}\NormalTok{(}\AttributeTok{x =} \FunctionTok{reorder}\NormalTok{(chr\_method, num\_elapsed\_seconds), }\AttributeTok{y =}\NormalTok{ num\_elapsed\_seconds)) }\SpecialCharTok{+}
  \FunctionTok{geom\_col}\NormalTok{() }\SpecialCharTok{+}
  \FunctionTok{coord\_flip}\NormalTok{() }\SpecialCharTok{+}
  \FunctionTok{labs}\NormalTok{(}\AttributeTok{x =} \ConstantTok{NULL}\NormalTok{, }\AttributeTok{y =} \StringTok{"Elapsed time (seconds)"}\NormalTok{) }
  \FunctionTok{add\_lagrange\_theme\_LFL}\NormalTok{(p)}
\end{Highlighting}
\end{Shaded}

\begin{figure}[H]

{\centering \pandocbounded{\includegraphics[keepaspectratio]{part1_empirical/ch01_r_basics_tidyverse_files/figure-pdf/pi-benchmark-plot-1.pdf}}

}

\caption{mapとfuture\_mapの実行時間比較}

\end{figure}%

並列計算には、ワーカーの起動、データ転送、結果回収のオーバーヘッドがある。そのため、処理が軽い場合は\texttt{map()}の方が速いこともある。計算量が大きいモンテカルロ・シミュレーション、銘柄別モデル推定、シナリオ分析では、\texttt{future\_map()}の効果が現れやすい。

\section{1.4
データベースと大規模データ処理}\label{ux30c7ux30fcux30bfux30d9ux30fcux30b9ux3068ux5927ux898fux6a21ux30c7ux30fcux30bfux51e6ux7406}

金融データが大きくなると、すべての行と列をRのメモリへ読み込んでから処理する方法には限界がある。本節では、データベースを\texttt{dplyr}の文法で操作する\texttt{dbplyr}と、ParquetファイルやDatasetを効率よく扱うArrowを確認する。

\subsection{dbplyrによるデータベース操作}\label{dbplyrux306bux3088ux308bux30c7ux30fcux30bfux30d9ux30fcux30b9ux64cdux4f5c}

\texttt{dbplyr}は、\texttt{dplyr}の処理をSQLへ変換する。データベース上で\texttt{filter()}、\texttt{select()}、\texttt{mutate()}、\texttt{summarise()}などを実行し、必要な結果だけを\texttt{collect()}でRへ取り込める。

ここでは、再現可能な例として一時的なSQLiteデータベースを使用する。

\texttt{dbplyr}では、処理は直ちにR上で実行されず、SQLクエリとして組み立てられる。\texttt{show\_query()}を使うと、生成されたSQLを確認できる。

重要なのは、\texttt{collect()}を最後まで遅らせることである。先にデータベース側で行と列を絞り込めば、Rへ転送するデータ量を小さくできる。

\subsection{ArrowによるParquetデータ処理}\label{arrowux306bux3088ux308bparquetux30c7ux30fcux30bfux51e6ux7406}

Apache
Arrowは、列指向データを効率よく扱うための仕組みである。Rでは、Parquetファイルの読書き、複数ファイルから成るDatasetの検索、PythonやC++など他言語との効率的なデータ交換に利用できる。

\begin{Shaded}
\begin{Highlighting}[]
\FunctionTok{set\_random\_seed\_LFL}\NormalTok{(}\DecValTok{20260726}\NormalTok{)}

\NormalTok{tbl\_arrow\_price }\OtherTok{\textless{}{-}} \FunctionTok{crossing}\NormalTok{(}
  \AttributeTok{date =} \FunctionTok{seq.Date}\NormalTok{(}\FunctionTok{as.Date}\NormalTok{(}\StringTok{"2025{-}01{-}01"}\NormalTok{), }\AttributeTok{by =} \StringTok{"day"}\NormalTok{, }\AttributeTok{length.out =} \DecValTok{365}\NormalTok{),}
  \AttributeTok{chr\_symbol =} \FunctionTok{str\_c}\NormalTok{(}\StringTok{"STOCK\_"}\NormalTok{, }\FunctionTok{str\_pad}\NormalTok{(}\DecValTok{1}\SpecialCharTok{:}\DecValTok{20}\NormalTok{, }\AttributeTok{width =} \DecValTok{2}\NormalTok{, }\AttributeTok{pad =} \StringTok{"0"}\NormalTok{))}
\NormalTok{) }\SpecialCharTok{|\textgreater{}}
  \FunctionTok{arrange}\NormalTok{(chr\_symbol, date) }\SpecialCharTok{|\textgreater{}}
  \FunctionTok{mutate}\NormalTok{(}
    \AttributeTok{num\_return =} \FunctionTok{rnorm}\NormalTok{(}\FunctionTok{n}\NormalTok{(), }\AttributeTok{mean =} \FloatTok{0.0002}\NormalTok{, }\AttributeTok{sd =} \FloatTok{0.015}\NormalTok{),}
    \AttributeTok{num\_price =} \DecValTok{100} \SpecialCharTok{*} \FunctionTok{exp}\NormalTok{(}\FunctionTok{cumsum}\NormalTok{(num\_return)),}
    \AttributeTok{.by =}\NormalTok{ chr\_symbol}
\NormalTok{  )}

\NormalTok{tbl\_arrow\_price}
\end{Highlighting}
\end{Shaded}

\begin{verbatim}
# A tibble: 7,300 x 4
   date       chr_symbol  num_return num_price
   <date>     <chr>            <dbl>     <dbl>
 1 2025-01-01 STOCK_01   -0.014847      98.526
 2 2025-01-02 STOCK_01    0.00017318    98.543
 3 2025-01-03 STOCK_01    0.0027754     98.817
 4 2025-01-04 STOCK_01    0.010886      99.899
 5 2025-01-05 STOCK_01    0.016885     101.60 
 6 2025-01-06 STOCK_01    0.017499     103.39 
 7 2025-01-07 STOCK_01    0.046333     108.30 
 8 2025-01-08 STOCK_01   -0.017287     106.44 
 9 2025-01-09 STOCK_01    0.0033907    106.80 
10 2025-01-10 STOCK_01   -0.0078836    105.96 
# i 7,290 more rows
\end{verbatim}

\begin{Shaded}
\begin{Highlighting}[]
\NormalTok{path\_parquet }\OtherTok{\textless{}{-}} \FunctionTok{file.path}\NormalTok{(}\FunctionTok{tempdir}\NormalTok{(), }\StringTok{"financial\_prices.parquet"}\NormalTok{)}

\NormalTok{arrow}\SpecialCharTok{::}\FunctionTok{write\_parquet}\NormalTok{(tbl\_arrow\_price, }\AttributeTok{sink =}\NormalTok{ path\_parquet)}

\NormalTok{path\_parquet}
\end{Highlighting}
\end{Shaded}

\begin{verbatim}
[1] "C:\\Users\\besuk\\AppData\\Local\\Temp\\RtmpIzOhS7/financial_prices.parquet"
\end{verbatim}

\begin{Shaded}
\begin{Highlighting}[]
\NormalTok{arrow}\SpecialCharTok{::}\FunctionTok{read\_parquet}\NormalTok{(path\_parquet) }\SpecialCharTok{|\textgreater{}}
  \FunctionTok{slice\_head}\NormalTok{(}\AttributeTok{n =} \DecValTok{10}\NormalTok{)}
\end{Highlighting}
\end{Shaded}

\begin{verbatim}
# A tibble: 10 x 4
   date       chr_symbol  num_return num_price
   <date>     <chr>            <dbl>     <dbl>
 1 2025-01-01 STOCK_01   -0.014847      98.526
 2 2025-01-02 STOCK_01    0.00017318    98.543
 3 2025-01-03 STOCK_01    0.0027754     98.817
 4 2025-01-04 STOCK_01    0.010886      99.899
 5 2025-01-05 STOCK_01    0.016885     101.60 
 6 2025-01-06 STOCK_01    0.017499     103.39 
 7 2025-01-07 STOCK_01    0.046333     108.30 
 8 2025-01-08 STOCK_01   -0.017287     106.44 
 9 2025-01-09 STOCK_01    0.0033907    106.80 
10 2025-01-10 STOCK_01   -0.0078836    105.96 
\end{verbatim}

\subsection{Arrow Dataset}\label{arrow-dataset}

複数のParquetファイルをまとめて扱う場合はDatasetを使用する。ここでは銘柄をパーティションとして保存する。

\begin{Shaded}
\begin{Highlighting}[]
\NormalTok{path\_dataset }\OtherTok{\textless{}{-}} \FunctionTok{file.path}\NormalTok{(}\FunctionTok{tempdir}\NormalTok{(), }\StringTok{"financial\_dataset"}\NormalTok{)}

\NormalTok{arrow}\SpecialCharTok{::}\FunctionTok{write\_dataset}\NormalTok{(}
\NormalTok{  tbl\_arrow\_price,}
  \AttributeTok{path =}\NormalTok{ path\_dataset,}
  \AttributeTok{format =} \StringTok{"parquet"}\NormalTok{,}
  \AttributeTok{partitioning =} \StringTok{"chr\_symbol"}\NormalTok{,}
  \AttributeTok{existing\_data\_behavior =} \StringTok{"overwrite"}
\NormalTok{)}

\NormalTok{dataset\_finance }\OtherTok{\textless{}{-}}\NormalTok{ arrow}\SpecialCharTok{::}\FunctionTok{open\_dataset}\NormalTok{(path\_dataset)}

\NormalTok{dataset\_finance}
\end{Highlighting}
\end{Shaded}

\begin{verbatim}
FileSystemDataset with 20 Parquet files
4 columns
date: date32[day]
num_return: double
num_price: double
chr_symbol: string
\end{verbatim}

Datasetに対しても\texttt{dplyr}に近い文法を使用できる。必要な銘柄、期間、列を指定してから\texttt{collect()}する。

\begin{Shaded}
\begin{Highlighting}[]
\NormalTok{tbl\_arrow\_selected }\OtherTok{\textless{}{-}}\NormalTok{ dataset\_finance }\SpecialCharTok{|\textgreater{}}
  \FunctionTok{filter}\NormalTok{(}
\NormalTok{    chr\_symbol }\SpecialCharTok{==} \StringTok{"STOCK\_01"}\NormalTok{,}
\NormalTok{    date }\SpecialCharTok{\textgreater{}=} \FunctionTok{as.Date}\NormalTok{(}\StringTok{"2025{-}07{-}01"}\NormalTok{)}
\NormalTok{  ) }\SpecialCharTok{|\textgreater{}}
  \FunctionTok{select}\NormalTok{(date, chr\_symbol, num\_price, num\_return) }\SpecialCharTok{|\textgreater{}}
  \FunctionTok{collect}\NormalTok{()}

\NormalTok{tbl\_arrow\_selected}
\end{Highlighting}
\end{Shaded}

\begin{verbatim}
# A tibble: 184 x 4
   date       chr_symbol num_price num_return
   <date>     <chr>          <dbl>      <dbl>
 1 2025-07-01 STOCK_01      92.811  0.0061451
 2 2025-07-02 STOCK_01      95.183  0.025240 
 3 2025-07-03 STOCK_01      97.579  0.024861 
 4 2025-07-04 STOCK_01      97.833  0.0025990
 5 2025-07-05 STOCK_01      97.612 -0.0022665
 6 2025-07-06 STOCK_01      98.638  0.010456 
 7 2025-07-07 STOCK_01      98.926  0.0029172
 8 2025-07-08 STOCK_01      98.611 -0.0031869
 9 2025-07-09 STOCK_01     100.28   0.016826 
10 2025-07-10 STOCK_01      99.207 -0.010806 
# i 174 more rows
\end{verbatim}

\subsection{Sparkとの関係}\label{sparkux3068ux306eux95a2ux4fc2}

ArrowとSparkは役割が異なる。Arrowは、列指向形式、Parquet、Dataset、異なる言語間の高速なデータ交換を中心とし、1台のPCでも利用しやすい。Sparkは、複数のコンピュータから成るクラスタへデータと計算を分散するための処理基盤である。

データが1台のPCで扱える範囲では、Arrowと\texttt{dbplyr}が導入しやすい。データ量や計算量が1台のメモリとCPUを大きく超える場合にはSparkが候補になる。本書では、実行環境を複雑にしないため、Arrowと\texttt{dbplyr}を実装対象とし、Sparkは位置づけの説明にとどめる。
\#\# 総合演習

最後に、\texttt{simulation\_data}から読み込んだ日次リターンを、読込み、整形、集計、反復処理、可視化の順に分析する。

\subsection{分析用データを作る}\label{ux5206ux6790ux7528ux30c7ux30fcux30bfux3092ux4f5cux308b}

\begin{Shaded}
\begin{Highlighting}[]
\NormalTok{analysis\_data }\OtherTok{\textless{}{-}}\NormalTok{ tbl\_daily\_stock\_return }\SpecialCharTok{|\textgreater{}}
  \FunctionTok{pivot\_longer}\NormalTok{(}
    \AttributeTok{cols =} \FunctionTok{c}\NormalTok{(firm1, firm2),}
    \AttributeTok{names\_to =} \StringTok{"firm"}\NormalTok{,}
    \AttributeTok{values\_to =} \StringTok{"return"}
\NormalTok{  ) }\SpecialCharTok{|\textgreater{}}
  \FunctionTok{mutate}\NormalTok{(}
    \AttributeTok{gross\_return =} \DecValTok{1} \SpecialCharTok{+}\NormalTok{ return,}
    \AttributeTok{direction =} \FunctionTok{case\_when}\NormalTok{(}
\NormalTok{      return }\SpecialCharTok{\textgreater{}} \DecValTok{0} \SpecialCharTok{\textasciitilde{}} \StringTok{"positive"}\NormalTok{,}
\NormalTok{      return }\SpecialCharTok{\textless{}} \DecValTok{0} \SpecialCharTok{\textasciitilde{}} \StringTok{"negative"}\NormalTok{,}
      \ConstantTok{TRUE} \SpecialCharTok{\textasciitilde{}} \StringTok{"zero"}
\NormalTok{    )}
\NormalTok{  ) }\SpecialCharTok{|\textgreater{}}
  \FunctionTok{arrange}\NormalTok{(firm, date)}

\NormalTok{analysis\_data}
\end{Highlighting}
\end{Shaded}

\begin{verbatim}
# A tibble: 42 x 5
   date       firm      return gross_return direction
   <date>     <chr>      <dbl>        <dbl> <chr>    
 1 2020-04-01 firm1 -0.039482       0.96052 negative 
 2 2020-04-02 firm1  0.0059787      1.0060  positive 
 3 2020-04-03 firm1  0.055787       1.0558  positive 
 4 2020-04-06 firm1  0.041935       1.0419  positive 
 5 2020-04-07 firm1 -0.020193       0.97981 negative 
 6 2020-04-08 firm1 -0.0022947      0.99771 negative 
 7 2020-04-09 firm1  0.081707       1.0817  positive 
 8 2020-04-10 firm1  0.034561       1.0346  positive 
 9 2020-04-13 firm1  0.032601       1.0326  positive 
10 2020-04-14 firm1 -0.037663       0.96234 negative 
# i 32 more rows
\end{verbatim}

\subsection{企業別の要約統計量を計算する}\label{ux4f01ux696dux5225ux306eux8981ux7d04ux7d71ux8a08ux91cfux3092ux8a08ux7b97ux3059ux308b}

\begin{Shaded}
\begin{Highlighting}[]
\NormalTok{firm\_summary }\OtherTok{\textless{}{-}}\NormalTok{ analysis\_data }\SpecialCharTok{|\textgreater{}}
  \FunctionTok{group\_by}\NormalTok{(firm) }\SpecialCharTok{|\textgreater{}}
  \FunctionTok{summarise}\NormalTok{(}
    \AttributeTok{observations =} \FunctionTok{n}\NormalTok{(),}
    \AttributeTok{mean\_return =} \FunctionTok{mean}\NormalTok{(return),}
    \AttributeTok{standard\_deviation =} \FunctionTok{sd}\NormalTok{(return),}
    \AttributeTok{minimum\_return =} \FunctionTok{min}\NormalTok{(return),}
    \AttributeTok{maximum\_return =} \FunctionTok{max}\NormalTok{(return),}
    \AttributeTok{positive\_ratio =} \FunctionTok{mean}\NormalTok{(return }\SpecialCharTok{\textgreater{}} \DecValTok{0}\NormalTok{),}
    \AttributeTok{.groups =} \StringTok{"drop"}
\NormalTok{  )}

\NormalTok{firm\_summary}
\end{Highlighting}
\end{Shaded}

\begin{verbatim}
# A tibble: 2 x 7
  firm  observations mean_return standard_deviation minimum_return
  <chr>        <int>       <dbl>              <dbl>          <dbl>
1 firm1           21   0.014161            0.053840      -0.072071
2 firm2           21   0.0080253           0.054390      -0.096212
  maximum_return positive_ratio
           <dbl>          <dbl>
1        0.17130        0.57143
2        0.13520        0.47619
\end{verbatim}

\subsection{累積リターンを計算する}\label{ux7d2fux7a4dux30eaux30bfux30fcux30f3ux3092ux8a08ux7b97ux3059ux308b}

\begin{Shaded}
\begin{Highlighting}[]
\NormalTok{cumulative\_return\_data }\OtherTok{\textless{}{-}}\NormalTok{ analysis\_data }\SpecialCharTok{|\textgreater{}}
  \FunctionTok{group\_by}\NormalTok{(firm) }\SpecialCharTok{|\textgreater{}}
  \FunctionTok{mutate}\NormalTok{(}
    \AttributeTok{cumulative\_gross\_return =} \FunctionTok{cumprod}\NormalTok{(gross\_return),}
    \AttributeTok{cumulative\_return =}\NormalTok{ cumulative\_gross\_return }\SpecialCharTok{{-}} \DecValTok{1}
\NormalTok{  ) }\SpecialCharTok{|\textgreater{}}
  \FunctionTok{ungroup}\NormalTok{()}

\NormalTok{cumulative\_return\_data}
\end{Highlighting}
\end{Shaded}

\begin{verbatim}
# A tibble: 42 x 7
   date       firm      return gross_return direction cumulative_gross_return
   <date>     <chr>      <dbl>        <dbl> <chr>                       <dbl>
 1 2020-04-01 firm1 -0.039482       0.96052 negative                  0.96052
 2 2020-04-02 firm1  0.0059787      1.0060  positive                  0.96626
 3 2020-04-03 firm1  0.055787       1.0558  positive                  1.0202 
 4 2020-04-06 firm1  0.041935       1.0419  positive                  1.0629 
 5 2020-04-07 firm1 -0.020193       0.97981 negative                  1.0415 
 6 2020-04-08 firm1 -0.0022947      0.99771 negative                  1.0391 
 7 2020-04-09 firm1  0.081707       1.0817  positive                  1.1240 
 8 2020-04-10 firm1  0.034561       1.0346  positive                  1.1628 
 9 2020-04-13 firm1  0.032601       1.0326  positive                  1.2007 
10 2020-04-14 firm1 -0.037663       0.96234 negative                  1.1555 
   cumulative_return
               <dbl>
 1         -0.039482
 2         -0.033740
 3          0.020165
 4          0.062945
 5          0.041481
 6          0.039091
 7          0.12399 
 8          0.16284 
 9          0.20075 
10          0.15553 
# i 32 more rows
\end{verbatim}

\subsection{purrrで企業別の統計量を再計算する}\label{purrrux3067ux4f01ux696dux5225ux306eux7d71ux8a08ux91cfux3092ux518dux8a08ux7b97ux3059ux308b}

\begin{Shaded}
\begin{Highlighting}[]
\NormalTok{firm\_names }\OtherTok{\textless{}{-}} \FunctionTok{unique}\NormalTok{(analysis\_data}\SpecialCharTok{$}\NormalTok{firm)}

\NormalTok{firm\_names }\SpecialCharTok{|\textgreater{}}
  \FunctionTok{map\_dfr}\NormalTok{(}
\NormalTok{    \textbackslash{}(firm\_name) \{}
\NormalTok{      firm\_returns }\OtherTok{\textless{}{-}}\NormalTok{ analysis\_data }\SpecialCharTok{|\textgreater{}}
        \FunctionTok{filter}\NormalTok{(firm }\SpecialCharTok{==}\NormalTok{ firm\_name) }\SpecialCharTok{|\textgreater{}}
        \FunctionTok{pull}\NormalTok{(return)}

      \FunctionTok{tibble}\NormalTok{(}
        \AttributeTok{firm =}\NormalTok{ firm\_name,}
        \AttributeTok{mean\_return =} \FunctionTok{mean}\NormalTok{(firm\_returns),}
        \AttributeTok{volatility =} \FunctionTok{sd}\NormalTok{(firm\_returns),}
        \AttributeTok{lower\_5\_percent =} \FunctionTok{quantile}\NormalTok{(firm\_returns, }\FloatTok{0.05}\NormalTok{),}
        \AttributeTok{upper\_95\_percent =} \FunctionTok{quantile}\NormalTok{(firm\_returns, }\FloatTok{0.95}\NormalTok{)}
\NormalTok{      )}
\NormalTok{    \}}
\NormalTok{  )}
\end{Highlighting}
\end{Shaded}

\begin{verbatim}
# A tibble: 2 x 5
  firm  mean_return volatility lower_5_percent upper_95_percent
  <chr>       <dbl>      <dbl>           <dbl>            <dbl>
1 firm1   0.014161    0.053840       -0.041378         0.081707
2 firm2   0.0080253   0.054390       -0.053818         0.078843
\end{verbatim}

\subsection{累積リターンを可視化する}\label{ux7d2fux7a4dux30eaux30bfux30fcux30f3ux3092ux53efux8996ux5316ux3059ux308b}

\begin{Shaded}
\begin{Highlighting}[]
\NormalTok{p }\OtherTok{\textless{}{-}}\NormalTok{ cumulative\_return\_data }\SpecialCharTok{|\textgreater{}}
  \FunctionTok{ggplot}\NormalTok{(}\FunctionTok{aes}\NormalTok{(}\AttributeTok{x =}\NormalTok{ date, }\AttributeTok{y =}\NormalTok{ cumulative\_return, }\AttributeTok{linetype =}\NormalTok{ firm)) }\SpecialCharTok{+}
  \FunctionTok{geom\_line}\NormalTok{(}\AttributeTok{linewidth =} \FloatTok{0.8}\NormalTok{) }\SpecialCharTok{+}
  \FunctionTok{geom\_hline}\NormalTok{(}\AttributeTok{yintercept =} \DecValTok{0}\NormalTok{, }\AttributeTok{linetype =} \StringTok{"dotted"}\NormalTok{) }\SpecialCharTok{+}
  \FunctionTok{labs}\NormalTok{(}
    \AttributeTok{x =} \StringTok{"Date"}\NormalTok{,}
    \AttributeTok{y =} \StringTok{"Cumulative Return"}\NormalTok{,}
    \AttributeTok{linetype =} \StringTok{"Firm"}
\NormalTok{  )}
\NormalTok{p }\SpecialCharTok{|\textgreater{}} \FunctionTok{add\_lagrange\_theme\_LFL}\NormalTok{()}
\end{Highlighting}
\end{Shaded}

\pandocbounded{\includegraphics[keepaspectratio]{part1_empirical/ch01_r_basics_tidyverse_files/figure-pdf/integrated-analysis-plot-1.pdf}}

\section{まとめ}\label{ux307eux3068ux3081}

本章では、Base
Rの基本的なデータ構造と関数から出発し、記述統計、確率分布、tidyverseによるデータ分析までを一通り確認した。重要な点は次のとおりである。

\begin{itemize}
\tightlist
\item
  オブジェクト名には内容を表すsnake\_caseを用いる。
\item
  同じ計算を複数の値へ適用するときは、まずベクトル化を検討する。
\item
  反復処理が必要な場合は\texttt{purrr::map()}系関数を用いる。
\item
  データ操作は\texttt{dplyr}、文字列処理は\texttt{stringr}、整形は\texttt{tidyr}を用いる。
\item
  パイプ演算子は\texttt{\textbar{}\textgreater{}}に統一する。
\item
  グラフは\texttt{ggplot2}でデータ、変数、図形を分けて組み立てる。
\item
  データファイルは絶対パスではなく、プロジェクトルートを基準に読み込む。
\end{itemize}

\subsection{金融データ分析パイプライン}\label{ux91d1ux878dux30c7ux30fcux30bfux5206ux6790ux30d1ux30a4ux30d7ux30e9ux30a4ux30f3}

本章で扱った技術は、次の流れとしてつながる。

\begin{Shaded}
\begin{Highlighting}[]
\NormalTok{CSV・データベース・Parquet}
\NormalTok{            ↓}
\NormalTok{      dbplyr・Arrow}
\NormalTok{            ↓}
\NormalTok{     filter・select}
\NormalTok{            ↓}
\NormalTok{         collect}
\NormalTok{            ↓}
\NormalTok{   tibble・stringr・tidyr}
\NormalTok{            ↓}
\NormalTok{          nest}
\NormalTok{            ↓}
\NormalTok{           map}
\NormalTok{            ↓}
\NormalTok{       future\_map}
\NormalTok{            ↓}
\NormalTok{   集計・モデル推定・可視化}
\end{Highlighting}
\end{Shaded}

\texttt{dbplyr}とArrowは、必要なデータを絞り込む入口を担う。\texttt{collect()}後は、tibbleとして文字列処理、整形、集計を行い、\texttt{nest()}で銘柄別・業種別・シナリオ別のデータを保持する。同じ分析は\texttt{map()}で反復し、計算量が大きい場合は\texttt{future\_map()}へ置き換える。可視化には\texttt{ggplot2}を用い、本書では\texttt{add\_lagrange\_theme\_LFL()}を標準テーマとする。

次章以降では、本章で学んだ操作を用いて、財務データと株式データの分析を進める。

\chapter{第2章
財務データの整理と分析}\label{ux7b2c2ux7ae0-ux8ca1ux52d9ux30c7ux30fcux30bfux306eux6574ux7406ux3068ux5206ux6790}

本章では、企業の財務データを読み込み、データ型、欠損値、年度構成および産業構成を確認する。次に、株主資本、自己資本利益率（ROE）および上級デュポン・モデルの構成要素を計算する。第1章で導入した\texttt{dplyr}、\texttt{tidyr}、\texttt{purrr}および\texttt{ggplot2}を、企業パネルデータの分析へ応用する。

\section{2.1 分析環境}\label{ux5206ux6790ux74b0ux5883}

\begin{Shaded}
\begin{Highlighting}[]
\FunctionTok{source}\NormalTok{(}\FunctionTok{file.path}\NormalTok{(}\StringTok{"R"}\NormalTok{, }\StringTok{"utility\_LFL.R"}\NormalTok{))}

\FunctionTok{load\_part1\_packages\_LFL}\NormalTok{()}
\end{Highlighting}
\end{Shaded}

\section{2.2
財務データの読込みと整備}\label{ux8ca1ux52d9ux30c7ux30fcux30bfux306eux8aadux8fbcux307fux3068ux6574ux5099}

\subsection{2.2.1
財務データの読込み}\label{ux8ca1ux52d9ux30c7ux30fcux30bfux306eux8aadux8fbcux307f}

\begin{Shaded}
\begin{Highlighting}[]
\NormalTok{financial\_data\_raw }\OtherTok{\textless{}{-}}
  \FunctionTok{read\_csv\_file\_LFL}\NormalTok{(}\StringTok{"ch04\_financial\_data.csv"}\NormalTok{)}

\FunctionTok{show\_table\_LFL}\NormalTok{(financial\_data\_raw, }\AttributeTok{n =} \DecValTok{10}\NormalTok{)}
\end{Highlighting}
\end{Shaded}

\begin{verbatim}
# A tibble: 7,920 x 11
    year firm_ID industry_ID  sales     OX      NFE      X      OA      FA
   <dbl>   <dbl>       <dbl>  <dbl>  <dbl>    <dbl>  <dbl>   <dbl>   <dbl>
 1  2015       1           1 5261.4 437.49  NA      286.64 13006.   3543.4
 2  2016       1           1 5949.0 564.14  50.667  513.48 13866.   4642.2
 3  2017       1           1 6505.1 691.18  29.543  661.64 13953.   7744.0
 4  2018       1           1 6846.4 751.29  86.487  664.8  18818.   7284.7
 5  2019       1           1 7572.2 958.53 298.05   660.48 18190    9735.1
 6  2020       1           1 7537.6 778.37 -65.459  843.83 20463.  10274. 
 7  2015       2           1 3505.8  45.82   5.7511  40.07  2977.8  2258.3
 8  2016       2           1 3491.3  51.25   1.8765  49.37  3184.8  1881.8
 9  2017       2           1 3945.7  83.43   7.5279  75.9   3392.2  2425.2
10  2018       2           1 4139.3  93.4    6.8166  86.58  3569.3  2331.0
       OL      FO
    <dbl>   <dbl>
 1 4373.0  2480.7
 2 4534.2  3959.7
 3 5111.2  6159.0
 4 5137.3 10124. 
 5 5488.0 11362. 
 6 5371.4 13772. 
 7 1840.4  2340.9
 8 1769.2  2215.9
 9 1955.3  2727.4
10 2022.1  2694.5
# i 7,910 more rows
\end{verbatim}

\subsection{2.2.2
データ構造の確認}\label{ux30c7ux30fcux30bfux69cbux9020ux306eux78baux8a8d}

\begin{Shaded}
\begin{Highlighting}[]
\FunctionTok{show\_data\_structure\_LFL}\NormalTok{(financial\_data\_raw)}
\end{Highlighting}
\end{Shaded}

\begin{verbatim}
Rows: 7,920
Columns: 11
$ year        <dbl> 2015, 2016, 2017, 2018, 2019, 2020, 2015, 2016, 2017, 2018~
$ firm_ID     <dbl> 1, 1, 1, 1, 1, 1, 2, 2, 2, 2, 2, 2, 3, 3, 3, 3, 3, 3, 4, 4~
$ industry_ID <dbl> 1, 1, 1, 1, 1, 1, 1, 1, 1, 1, 1, 1, 1, 1, 1, 1, 1, 1, 1, 1~
$ sales       <dbl> 5261.40, 5948.96, 6505.06, 6846.38, 7572.24, 7537.63, 3505~
$ OX          <dbl> 437.49, 564.14, 691.18, 751.29, 958.53, 778.37, 45.82, 51.~
$ NFE         <dbl> NA, 50.667498, 29.543157, 86.486500, 298.049774, -65.45877~
$ X           <dbl> 286.64, 513.48, 661.64, 664.80, 660.48, 843.83, 40.07, 49.~
$ OA          <dbl> 13005.55, 13865.58, 13952.58, 18818.48, 18190.00, 20462.86~
$ FA          <dbl> 3543.43, 4642.16, 7743.99, 7284.72, 9735.13, 10274.25, 225~
$ OL          <dbl> 4372.96, 4534.22, 5111.22, 5137.28, 5487.96, 5371.38, 1840~
$ FO          <dbl> 2480.72, 3959.70, 6159.02, 10123.91, 11362.22, 13772.15, 2~
\end{verbatim}

\begin{Shaded}
\begin{Highlighting}[]
\NormalTok{financial\_data\_overview }\OtherTok{\textless{}{-}}\NormalTok{ financial\_data\_raw }\SpecialCharTok{|\textgreater{}}
  \FunctionTok{summarise}\NormalTok{(}
    \AttributeTok{observations =} \FunctionTok{n}\NormalTok{(),}
    \AttributeTok{variables =} \FunctionTok{ncol}\NormalTok{(financial\_data\_raw),}
    \AttributeTok{first\_year =} \FunctionTok{min}\NormalTok{(year, }\AttributeTok{na.rm =} \ConstantTok{TRUE}\NormalTok{),}
    \AttributeTok{last\_year =} \FunctionTok{max}\NormalTok{(year, }\AttributeTok{na.rm =} \ConstantTok{TRUE}\NormalTok{),}
    \AttributeTok{firms =} \FunctionTok{n\_distinct}\NormalTok{(firm\_ID),}
    \AttributeTok{industries =} \FunctionTok{n\_distinct}\NormalTok{(industry\_ID)}
\NormalTok{  )}

\FunctionTok{show\_table\_LFL}\NormalTok{(financial\_data\_overview)}
\end{Highlighting}
\end{Shaded}

\begin{verbatim}
# A tibble: 1 x 6
  observations variables first_year last_year firms industries
         <int>     <int>      <dbl>     <dbl> <int>      <int>
1         7920        11       2015      2020  1515         10
\end{verbatim}

列ごとの欠損数と欠損率を確認する。

\begin{Shaded}
\begin{Highlighting}[]
\NormalTok{missing\_summary }\OtherTok{\textless{}{-}}\NormalTok{ financial\_data\_raw }\SpecialCharTok{|\textgreater{}}
  \FunctionTok{summarise}\NormalTok{(}
    \FunctionTok{across}\NormalTok{(}
      \FunctionTok{everything}\NormalTok{(),}
      \FunctionTok{list}\NormalTok{(}
        \AttributeTok{missing =}\NormalTok{ \textbackslash{}(x) }\FunctionTok{sum}\NormalTok{(}\FunctionTok{is.na}\NormalTok{(x)),}
        \AttributeTok{missing\_rate =}\NormalTok{ \textbackslash{}(x) }\FunctionTok{mean}\NormalTok{(}\FunctionTok{is.na}\NormalTok{(x))}
\NormalTok{      ),}
      \AttributeTok{.names =} \StringTok{"\{.col\}\_\_\{.fn\}"}
\NormalTok{    )}
\NormalTok{  ) }\SpecialCharTok{|\textgreater{}}
  \FunctionTok{pivot\_longer}\NormalTok{(}
    \AttributeTok{cols =} \FunctionTok{everything}\NormalTok{(),}
    \AttributeTok{names\_to =} \FunctionTok{c}\NormalTok{(}\StringTok{"variable"}\NormalTok{, }\StringTok{".value"}\NormalTok{),}
    \AttributeTok{names\_sep =} \StringTok{"\_\_"}
\NormalTok{  ) }\SpecialCharTok{|\textgreater{}}
  \FunctionTok{arrange}\NormalTok{(}\FunctionTok{desc}\NormalTok{(missing\_rate), variable)}

\FunctionTok{show\_table\_LFL}\NormalTok{(missing\_summary, }\AttributeTok{n =} \FunctionTok{nrow}\NormalTok{(missing\_summary))}
\end{Highlighting}
\end{Shaded}

\begin{verbatim}
# A tibble: 11 x 3
   variable    missing missing_rate
   <chr>         <int>        <dbl>
 1 NFE               1   0.00012626
 2 FA                0   0         
 3 FO                0   0         
 4 OA                0   0         
 5 OL                0   0         
 6 OX                0   0         
 7 X                 0   0         
 8 firm_ID           0   0         
 9 industry_ID       0   0         
10 sales             0   0         
11 year              0   0         
\end{verbatim}

\subsection{2.2.3
データ型と並び順の整理}\label{ux30c7ux30fcux30bfux578bux3068ux4e26ux3073ux9806ux306eux6574ux7406}

企業IDと産業IDは分類変数であるためファクター型へ変換する。年度は整数型へ変換し、前年度の値を正しく取得できるように企業IDと年度の順に並べる。

\begin{Shaded}
\begin{Highlighting}[]
\NormalTok{financial\_data }\OtherTok{\textless{}{-}}\NormalTok{ financial\_data\_raw }\SpecialCharTok{|\textgreater{}}
  \FunctionTok{mutate}\NormalTok{(}
    \AttributeTok{firm\_ID =} \FunctionTok{as.factor}\NormalTok{(firm\_ID),}
    \AttributeTok{industry\_ID =} \FunctionTok{as.factor}\NormalTok{(industry\_ID),}
    \AttributeTok{year =} \FunctionTok{as.integer}\NormalTok{(year)}
\NormalTok{  ) }\SpecialCharTok{|\textgreater{}}
  \FunctionTok{arrange}\NormalTok{(firm\_ID, year)}

\FunctionTok{show\_data\_structure\_LFL}\NormalTok{(financial\_data)}
\end{Highlighting}
\end{Shaded}

\begin{verbatim}
Rows: 7,920
Columns: 11
$ year        <int> 2015, 2016, 2017, 2018, 2019, 2020, 2015, 2016, 2017, 2018~
$ firm_ID     <fct> 1, 1, 1, 1, 1, 1, 2, 2, 2, 2, 2, 2, 3, 3, 3, 3, 3, 3, 4, 4~
$ industry_ID <fct> 1, 1, 1, 1, 1, 1, 1, 1, 1, 1, 1, 1, 1, 1, 1, 1, 1, 1, 1, 1~
$ sales       <dbl> 5261.40, 5948.96, 6505.06, 6846.38, 7572.24, 7537.63, 3505~
$ OX          <dbl> 437.49, 564.14, 691.18, 751.29, 958.53, 778.37, 45.82, 51.~
$ NFE         <dbl> NA, 50.667498, 29.543157, 86.486500, 298.049774, -65.45877~
$ X           <dbl> 286.64, 513.48, 661.64, 664.80, 660.48, 843.83, 40.07, 49.~
$ OA          <dbl> 13005.55, 13865.58, 13952.58, 18818.48, 18190.00, 20462.86~
$ FA          <dbl> 3543.43, 4642.16, 7743.99, 7284.72, 9735.13, 10274.25, 225~
$ OL          <dbl> 4372.96, 4534.22, 5111.22, 5137.28, 5487.96, 5371.38, 1840~
$ FO          <dbl> 2480.72, 3959.70, 6159.02, 10123.91, 11362.22, 13772.15, 2~
\end{verbatim}

企業と年度の組合せが重複していないことを確認する。

\begin{Shaded}
\begin{Highlighting}[]
\NormalTok{panel\_key\_duplicates }\OtherTok{\textless{}{-}}\NormalTok{ financial\_data }\SpecialCharTok{|\textgreater{}}
  \FunctionTok{count}\NormalTok{(}
\NormalTok{    firm\_ID,}
\NormalTok{    year,}
    \AttributeTok{name =} \StringTok{"observations"}
\NormalTok{  ) }\SpecialCharTok{|\textgreater{}}
  \FunctionTok{filter}\NormalTok{(observations }\SpecialCharTok{\textgreater{}} \DecValTok{1}\NormalTok{)}

\FunctionTok{show\_table\_LFL}\NormalTok{(panel\_key\_duplicates)}
\end{Highlighting}
\end{Shaded}

\begin{verbatim}
# A tibble: 0 x 3
# i 3 variables: firm_ID <fct>, year <int>, observations <int>
\end{verbatim}

\begin{Shaded}
\begin{Highlighting}[]
\ControlFlowTok{if}\NormalTok{ (}\FunctionTok{nrow}\NormalTok{(panel\_key\_duplicates) }\SpecialCharTok{\textgreater{}} \DecValTok{0}\NormalTok{) \{}
  \FunctionTok{stop}\NormalTok{(}
    \StringTok{"firm\_IDとyearの組合せに重複があります。"}\NormalTok{,}
    \AttributeTok{call. =} \ConstantTok{FALSE}
\NormalTok{  )}
\NormalTok{\}}
\end{Highlighting}
\end{Shaded}

\section{2.3
標本構成と売上高分布}\label{ux6a19ux672cux69cbux6210ux3068ux58f2ux4e0aux9ad8ux5206ux5e03}

\subsection{2.3.1
年度別の企業数}\label{ux5e74ux5ea6ux5225ux306eux4f01ux696dux6570}

年度ごとの企業数と観測数を集計する。

\begin{Shaded}
\begin{Highlighting}[]
\NormalTok{firms\_by\_year }\OtherTok{\textless{}{-}}\NormalTok{ financial\_data }\SpecialCharTok{|\textgreater{}}
  \FunctionTok{group\_by}\NormalTok{(year) }\SpecialCharTok{|\textgreater{}}
  \FunctionTok{summarise}\NormalTok{(}
    \AttributeTok{firms =} \FunctionTok{n\_distinct}\NormalTok{(firm\_ID),}
    \AttributeTok{observations =} \FunctionTok{n}\NormalTok{(),}
    \AttributeTok{.groups =} \StringTok{"drop"}
\NormalTok{  ) }\SpecialCharTok{|\textgreater{}}
  \FunctionTok{arrange}\NormalTok{(year)}

\FunctionTok{show\_table\_LFL}\NormalTok{(firms\_by\_year, }\AttributeTok{n =} \FunctionTok{nrow}\NormalTok{(firms\_by\_year))}
\end{Highlighting}
\end{Shaded}

\begin{verbatim}
# A tibble: 6 x 3
   year firms observations
  <int> <int>        <int>
1  2015  1266         1266
2  2016  1293         1293
3  2017  1319         1319
4  2018  1323         1323
5  2019  1356         1356
6  2020  1363         1363
\end{verbatim}

\begin{Shaded}
\begin{Highlighting}[]
\NormalTok{p\_firms\_by\_year }\OtherTok{\textless{}{-}}\NormalTok{ firms\_by\_year }\SpecialCharTok{|\textgreater{}}
  \FunctionTok{ggplot}\NormalTok{(}\FunctionTok{aes}\NormalTok{(}\AttributeTok{x =}\NormalTok{ year, }\AttributeTok{y =}\NormalTok{ firms, }\AttributeTok{colour =} \StringTok{"Firms"}\NormalTok{)) }\SpecialCharTok{+}
  \FunctionTok{geom\_line}\NormalTok{(}\AttributeTok{linewidth =} \FloatTok{0.8}\NormalTok{) }\SpecialCharTok{+}
  \FunctionTok{geom\_point}\NormalTok{(}\AttributeTok{size =} \DecValTok{2}\NormalTok{) }\SpecialCharTok{+}
  \FunctionTok{scale\_x\_continuous}\NormalTok{(}\AttributeTok{breaks =}\NormalTok{ firms\_by\_year}\SpecialCharTok{$}\NormalTok{year) }\SpecialCharTok{+}
  \FunctionTok{guides}\NormalTok{(}\AttributeTok{colour =} \StringTok{"none"}\NormalTok{) }\SpecialCharTok{+}
  \FunctionTok{labs}\NormalTok{(}
    \AttributeTok{x =} \StringTok{"Year"}\NormalTok{,}
    \AttributeTok{y =} \StringTok{"Number of Firms"}\NormalTok{,}
    \AttributeTok{title =} \StringTok{"Number of Firms by Year"}
\NormalTok{  )}

\FunctionTok{add\_lagrange\_theme\_LFL}\NormalTok{(p\_firms\_by\_year)}
\end{Highlighting}
\end{Shaded}

\pandocbounded{\includegraphics[keepaspectratio]{part1_empirical/ch02_financial_data_files/figure-pdf/plot-firms-by-year-1.pdf}}

\subsection{2.3.2
産業別の企業数}\label{ux7523ux696dux5225ux306eux4f01ux696dux6570}

最新年度について、産業ごとの企業数を集計する。

\begin{Shaded}
\begin{Highlighting}[]
\NormalTok{latest\_year }\OtherTok{\textless{}{-}}
  \FunctionTok{max}\NormalTok{(financial\_data}\SpecialCharTok{$}\NormalTok{year, }\AttributeTok{na.rm =} \ConstantTok{TRUE}\NormalTok{)}

\NormalTok{firms\_by\_industry }\OtherTok{\textless{}{-}}\NormalTok{ financial\_data }\SpecialCharTok{|\textgreater{}}
  \FunctionTok{filter}\NormalTok{(year }\SpecialCharTok{==}\NormalTok{ latest\_year) }\SpecialCharTok{|\textgreater{}}
  \FunctionTok{count}\NormalTok{(}
\NormalTok{    industry\_ID,}
    \AttributeTok{name =} \StringTok{"firms"}\NormalTok{,}
    \AttributeTok{sort =} \ConstantTok{TRUE}
\NormalTok{  )}

\FunctionTok{show\_table\_LFL}\NormalTok{(firms\_by\_industry, }\AttributeTok{n =} \DecValTok{20}\NormalTok{)}
\end{Highlighting}
\end{Shaded}

\begin{verbatim}
# A tibble: 10 x 2
   industry_ID firms
   <fct>       <int>
 1 3             309
 2 10            291
 3 7             233
 4 1             192
 5 9             123
 6 8              71
 7 5              59
 8 2              37
 9 4              28
10 6              20
\end{verbatim}

\begin{Shaded}
\begin{Highlighting}[]
\NormalTok{p\_firms\_by\_industry }\OtherTok{\textless{}{-}}\NormalTok{ firms\_by\_industry }\SpecialCharTok{|\textgreater{}}
  \FunctionTok{slice\_head}\NormalTok{(}\AttributeTok{n =} \DecValTok{15}\NormalTok{) }\SpecialCharTok{|\textgreater{}}
  \FunctionTok{ggplot}\NormalTok{(}
    \FunctionTok{aes}\NormalTok{(}
      \AttributeTok{x =} \FunctionTok{reorder}\NormalTok{(industry\_ID, firms),}
      \AttributeTok{y =}\NormalTok{ firms,}
      \AttributeTok{fill =}\NormalTok{ industry\_ID}
\NormalTok{    )}
\NormalTok{  ) }\SpecialCharTok{+}
  \FunctionTok{geom\_col}\NormalTok{(}\AttributeTok{show.legend =} \ConstantTok{FALSE}\NormalTok{) }\SpecialCharTok{+}
  \FunctionTok{coord\_flip}\NormalTok{() }\SpecialCharTok{+}
  \FunctionTok{labs}\NormalTok{(}
    \AttributeTok{x =} \StringTok{"Industry ID"}\NormalTok{,}
    \AttributeTok{y =} \StringTok{"Number of Firms"}\NormalTok{,}
    \AttributeTok{title =} \FunctionTok{paste}\NormalTok{(}\StringTok{"Largest Industries in"}\NormalTok{, latest\_year)}
\NormalTok{  )}

\FunctionTok{add\_lagrange\_theme\_LFL}\NormalTok{(p\_firms\_by\_industry)}
\end{Highlighting}
\end{Shaded}

\pandocbounded{\includegraphics[keepaspectratio]{part1_empirical/ch02_financial_data_files/figure-pdf/plot-firms-by-industry-1.pdf}}

\subsection{2.3.3
売上高の分布}\label{ux58f2ux4e0aux9ad8ux306eux5206ux5e03}

売上高は企業規模によって大きく異なるため、元の値と自然対数の両方を確認する。

\begin{Shaded}
\begin{Highlighting}[]
\NormalTok{sales\_summary }\OtherTok{\textless{}{-}}\NormalTok{ financial\_data }\SpecialCharTok{|\textgreater{}}
  \FunctionTok{summarise}\NormalTok{(}
    \AttributeTok{observations =} \FunctionTok{sum}\NormalTok{(}\SpecialCharTok{!}\FunctionTok{is.na}\NormalTok{(sales)),}
    \AttributeTok{minimum =} \FunctionTok{min}\NormalTok{(sales, }\AttributeTok{na.rm =} \ConstantTok{TRUE}\NormalTok{),}
    \AttributeTok{first\_quartile =} \FunctionTok{quantile}\NormalTok{(sales, }\FloatTok{0.25}\NormalTok{, }\AttributeTok{na.rm =} \ConstantTok{TRUE}\NormalTok{),}
    \AttributeTok{median =} \FunctionTok{median}\NormalTok{(sales, }\AttributeTok{na.rm =} \ConstantTok{TRUE}\NormalTok{),}
    \AttributeTok{mean =} \FunctionTok{mean}\NormalTok{(sales, }\AttributeTok{na.rm =} \ConstantTok{TRUE}\NormalTok{),}
    \AttributeTok{third\_quartile =} \FunctionTok{quantile}\NormalTok{(sales, }\FloatTok{0.75}\NormalTok{, }\AttributeTok{na.rm =} \ConstantTok{TRUE}\NormalTok{),}
    \AttributeTok{maximum =} \FunctionTok{max}\NormalTok{(sales, }\AttributeTok{na.rm =} \ConstantTok{TRUE}\NormalTok{)}
\NormalTok{  )}

\FunctionTok{show\_table\_LFL}\NormalTok{(sales\_summary)}
\end{Highlighting}
\end{Shaded}

\begin{verbatim}
# A tibble: 1 x 7
  observations minimum first_quartile median    mean third_quartile maximum
         <int>   <dbl>          <dbl>  <dbl>   <dbl>          <dbl>   <dbl>
1         7920  205.34         16103. 40431. 166007.        118314. 3496433
\end{verbatim}

\begin{Shaded}
\begin{Highlighting}[]
\NormalTok{p\_sales\_distribution }\OtherTok{\textless{}{-}}\NormalTok{ financial\_data }\SpecialCharTok{|\textgreater{}}
  \FunctionTok{filter}\NormalTok{(}
\NormalTok{    year }\SpecialCharTok{==}\NormalTok{ latest\_year,}
    \SpecialCharTok{!}\FunctionTok{is.na}\NormalTok{(sales),}
\NormalTok{    sales }\SpecialCharTok{\textgreater{}=} \DecValTok{0}
\NormalTok{  ) }\SpecialCharTok{|\textgreater{}}
  \FunctionTok{ggplot}\NormalTok{(}\FunctionTok{aes}\NormalTok{(}\AttributeTok{x =}\NormalTok{ sales, }\AttributeTok{fill =} \StringTok{"Sales"}\NormalTok{)) }\SpecialCharTok{+}
  \FunctionTok{geom\_histogram}\NormalTok{(}
    \AttributeTok{bins =} \DecValTok{40}\NormalTok{,}
    \AttributeTok{colour =} \StringTok{"white"}\NormalTok{,}
    \AttributeTok{show.legend =} \ConstantTok{FALSE}
\NormalTok{  ) }\SpecialCharTok{+}
  \FunctionTok{scale\_x\_continuous}\NormalTok{(}\AttributeTok{labels =}\NormalTok{ scales}\SpecialCharTok{::}\FunctionTok{label\_comma}\NormalTok{()) }\SpecialCharTok{+}
  \FunctionTok{labs}\NormalTok{(}
    \AttributeTok{x =} \StringTok{"Sales"}\NormalTok{,}
    \AttributeTok{y =} \StringTok{"Count"}\NormalTok{,}
    \AttributeTok{title =} \FunctionTok{paste}\NormalTok{(}\StringTok{"Distribution of Sales in"}\NormalTok{, latest\_year)}
\NormalTok{  )}

\FunctionTok{add\_lagrange\_theme\_LFL}\NormalTok{(p\_sales\_distribution)}
\end{Highlighting}
\end{Shaded}

\pandocbounded{\includegraphics[keepaspectratio]{part1_empirical/ch02_financial_data_files/figure-pdf/plot-sales-distribution-1.pdf}}

\begin{Shaded}
\begin{Highlighting}[]
\NormalTok{p\_log\_sales\_distribution }\OtherTok{\textless{}{-}}\NormalTok{ financial\_data }\SpecialCharTok{|\textgreater{}}
  \FunctionTok{filter}\NormalTok{(}
\NormalTok{    year }\SpecialCharTok{==}\NormalTok{ latest\_year,}
    \SpecialCharTok{!}\FunctionTok{is.na}\NormalTok{(sales),}
\NormalTok{    sales }\SpecialCharTok{\textgreater{}} \DecValTok{0}
\NormalTok{  ) }\SpecialCharTok{|\textgreater{}}
  \FunctionTok{ggplot}\NormalTok{(}\FunctionTok{aes}\NormalTok{(}\AttributeTok{x =} \FunctionTok{log}\NormalTok{(sales), }\AttributeTok{fill =} \StringTok{"Log Sales"}\NormalTok{)) }\SpecialCharTok{+}
  \FunctionTok{geom\_histogram}\NormalTok{(}
    \AttributeTok{bins =} \DecValTok{40}\NormalTok{,}
    \AttributeTok{colour =} \StringTok{"white"}\NormalTok{,}
    \AttributeTok{show.legend =} \ConstantTok{FALSE}
\NormalTok{  ) }\SpecialCharTok{+}
  \FunctionTok{labs}\NormalTok{(}
    \AttributeTok{x =} \StringTok{"Log Sales"}\NormalTok{,}
    \AttributeTok{y =} \StringTok{"Count"}\NormalTok{,}
    \AttributeTok{title =} \FunctionTok{paste}\NormalTok{(}\StringTok{"Distribution of Log Sales in"}\NormalTok{, latest\_year)}
\NormalTok{  )}

\FunctionTok{add\_lagrange\_theme\_LFL}\NormalTok{(p\_log\_sales\_distribution)}
\end{Highlighting}
\end{Shaded}

\pandocbounded{\includegraphics[keepaspectratio]{part1_empirical/ch02_financial_data_files/figure-pdf/plot-log-sales-distribution-1.pdf}}

\section{2.4
株主資本とROE分析}\label{ux682aux4e3bux8cc7ux672cux3068roeux5206ux6790}

\subsection{2.4.1
株主資本とROE}\label{ux682aux4e3bux8cc7ux672cux3068roe}

本章では、営業資産から営業負債を控除した純営業資産と、金融負債から金融資産を控除した純金融負債の差として株主資本を計算する。

\[
BE_t = (OA_t - OL_t) - (FO_t - FA_t)
\]

ROEは当期純利益を期首株主資本で割って計算する。

\[
ROE_t = \frac{X_t}{BE_{t-1}}
\]

ゼロ除算や無限大は欠損値へ置き換える。

\begin{Shaded}
\begin{Highlighting}[]
\NormalTok{divide\_safely\_LFL }\OtherTok{\textless{}{-}} \ControlFlowTok{function}\NormalTok{(numerator, denominator) \{}
\NormalTok{  result }\OtherTok{\textless{}{-}}\NormalTok{ numerator }\SpecialCharTok{/}\NormalTok{ denominator}
  \FunctionTok{replace}\NormalTok{(result, }\SpecialCharTok{!}\FunctionTok{is.finite}\NormalTok{(result), }\ConstantTok{NA\_real\_}\NormalTok{)}
\NormalTok{\}}
\end{Highlighting}
\end{Shaded}

\begin{Shaded}
\begin{Highlighting}[]
\NormalTok{financial\_data }\OtherTok{\textless{}{-}}\NormalTok{ financial\_data }\SpecialCharTok{|\textgreater{}}
  \FunctionTok{mutate}\NormalTok{(}
    \AttributeTok{NOA =}\NormalTok{ OA }\SpecialCharTok{{-}}\NormalTok{ OL,}
    \AttributeTok{NFO =}\NormalTok{ FO }\SpecialCharTok{{-}}\NormalTok{ FA,}
    \AttributeTok{BE =}\NormalTok{ NOA }\SpecialCharTok{{-}}\NormalTok{ NFO}
\NormalTok{  ) }\SpecialCharTok{|\textgreater{}}
  \FunctionTok{group\_by}\NormalTok{(firm\_ID) }\SpecialCharTok{|\textgreater{}}
  \FunctionTok{arrange}\NormalTok{(year, }\AttributeTok{.by\_group =} \ConstantTok{TRUE}\NormalTok{) }\SpecialCharTok{|\textgreater{}}
  \FunctionTok{mutate}\NormalTok{(}
    \AttributeTok{lagged\_BE =} \FunctionTok{lag}\NormalTok{(BE),}
    \AttributeTok{ROE =} \FunctionTok{divide\_safely\_LFL}\NormalTok{(X, lagged\_BE)}
\NormalTok{  ) }\SpecialCharTok{|\textgreater{}}
  \FunctionTok{ungroup}\NormalTok{()}

\NormalTok{financial\_data }\SpecialCharTok{|\textgreater{}}
  \FunctionTok{select}\NormalTok{(}
\NormalTok{    firm\_ID,}
\NormalTok{    industry\_ID,}
\NormalTok{    year,}
\NormalTok{    X,}
\NormalTok{    BE,}
\NormalTok{    lagged\_BE,}
\NormalTok{    ROE}
\NormalTok{  ) }\SpecialCharTok{|\textgreater{}}
  \FunctionTok{slice\_head}\NormalTok{(}\AttributeTok{n =} \DecValTok{10}\NormalTok{) }\SpecialCharTok{|\textgreater{}}
  \FunctionTok{show\_table\_LFL}\NormalTok{(}\AttributeTok{n =} \DecValTok{10}\NormalTok{)}
\end{Highlighting}
\end{Shaded}

\begin{verbatim}
# A tibble: 10 x 7
   firm_ID industry_ID  year      X      BE lagged_BE       ROE
   <fct>   <fct>       <int>  <dbl>   <dbl>     <dbl>     <dbl>
 1 1       1            2015 286.64  9695.3      NA   NA       
 2 1       1            2016 513.48 10014.     9695.3  0.052962
 3 1       1            2017 661.64 10426.    10014.   0.066073
 4 1       1            2018 664.8  10842.    10426.   0.063762
 5 1       1            2019 660.48 11075.    10842.   0.060919
 6 1       1            2020 843.83 11594.    11075.   0.076193
 7 2       1            2015  40.07  1054.9      NA   NA       
 8 2       1            2016  49.37  1081.6    1054.9  0.046800
 9 2       1            2017  75.9   1134.8    1081.6  0.070174
10 2       1            2018  86.58  1183.7    1134.8  0.076295
\end{verbatim}

ROEの要約統計量を確認する。

\begin{Shaded}
\begin{Highlighting}[]
\NormalTok{roe\_summary }\OtherTok{\textless{}{-}}\NormalTok{ financial\_data }\SpecialCharTok{|\textgreater{}}
  \FunctionTok{summarise}\NormalTok{(}
    \AttributeTok{observations =} \FunctionTok{sum}\NormalTok{(}\SpecialCharTok{!}\FunctionTok{is.na}\NormalTok{(ROE)),}
    \AttributeTok{mean =} \FunctionTok{mean}\NormalTok{(ROE, }\AttributeTok{na.rm =} \ConstantTok{TRUE}\NormalTok{),}
    \AttributeTok{median =} \FunctionTok{median}\NormalTok{(ROE, }\AttributeTok{na.rm =} \ConstantTok{TRUE}\NormalTok{),}
    \AttributeTok{standard\_deviation =} \FunctionTok{sd}\NormalTok{(ROE, }\AttributeTok{na.rm =} \ConstantTok{TRUE}\NormalTok{),}
    \AttributeTok{first\_percentile =} \FunctionTok{quantile}\NormalTok{(ROE, }\FloatTok{0.01}\NormalTok{, }\AttributeTok{na.rm =} \ConstantTok{TRUE}\NormalTok{),}
    \AttributeTok{ninety\_ninth\_percentile =} \FunctionTok{quantile}\NormalTok{(ROE, }\FloatTok{0.99}\NormalTok{, }\AttributeTok{na.rm =} \ConstantTok{TRUE}\NormalTok{)}
\NormalTok{  )}

\FunctionTok{show\_table\_LFL}\NormalTok{(roe\_summary)}
\end{Highlighting}
\end{Shaded}

\begin{verbatim}
# A tibble: 1 x 6
  observations     mean   median standard_deviation first_percentile
         <int>    <dbl>    <dbl>              <dbl>            <dbl>
1         6405 0.077713 0.072828           0.066274        -0.079541
  ninety_ninth_percentile
                    <dbl>
1                 0.26731
\end{verbatim}

極端値の影響を避けるため、1パーセンタイルから99パーセンタイルまでを表示する。

\begin{Shaded}
\begin{Highlighting}[]
\NormalTok{roe\_limits }\OtherTok{\textless{}{-}} \FunctionTok{quantile}\NormalTok{(}
\NormalTok{  financial\_data}\SpecialCharTok{$}\NormalTok{ROE,}
  \AttributeTok{probs =} \FunctionTok{c}\NormalTok{(}\FloatTok{0.01}\NormalTok{, }\FloatTok{0.99}\NormalTok{),}
  \AttributeTok{na.rm =} \ConstantTok{TRUE}
\NormalTok{)}

\NormalTok{p\_roe\_distribution }\OtherTok{\textless{}{-}}\NormalTok{ financial\_data }\SpecialCharTok{|\textgreater{}}
  \FunctionTok{filter}\NormalTok{(}
    \SpecialCharTok{!}\FunctionTok{is.na}\NormalTok{(ROE),}
    \FunctionTok{between}\NormalTok{(ROE, roe\_limits[[}\DecValTok{1}\NormalTok{]], roe\_limits[[}\DecValTok{2}\NormalTok{]])}
\NormalTok{  ) }\SpecialCharTok{|\textgreater{}}
  \FunctionTok{ggplot}\NormalTok{(}\FunctionTok{aes}\NormalTok{(}\AttributeTok{x =}\NormalTok{ ROE, }\AttributeTok{fill =} \StringTok{"ROE"}\NormalTok{)) }\SpecialCharTok{+}
  \FunctionTok{geom\_histogram}\NormalTok{(}
    \AttributeTok{bins =} \DecValTok{50}\NormalTok{,}
    \AttributeTok{colour =} \StringTok{"white"}\NormalTok{,}
    \AttributeTok{show.legend =} \ConstantTok{FALSE}
\NormalTok{  ) }\SpecialCharTok{+}
  \FunctionTok{scale\_x\_continuous}\NormalTok{(}\AttributeTok{labels =}\NormalTok{ scales}\SpecialCharTok{::}\FunctionTok{label\_percent}\NormalTok{()) }\SpecialCharTok{+}
  \FunctionTok{labs}\NormalTok{(}
    \AttributeTok{x =} \StringTok{"ROE"}\NormalTok{,}
    \AttributeTok{y =} \StringTok{"Count"}\NormalTok{,}
    \AttributeTok{title =} \StringTok{"Distribution of Return on Equity"}
\NormalTok{  )}

\FunctionTok{add\_lagrange\_theme\_LFL}\NormalTok{(p\_roe\_distribution)}
\end{Highlighting}
\end{Shaded}

\pandocbounded{\includegraphics[keepaspectratio]{part1_empirical/ch02_financial_data_files/figure-pdf/plot-roe-distribution-1.pdf}}

\subsection{2.4.2 産業別のROE}\label{ux7523ux696dux5225ux306eroe}

\begin{Shaded}
\begin{Highlighting}[]
\NormalTok{industry\_roe\_summary }\OtherTok{\textless{}{-}}\NormalTok{ financial\_data }\SpecialCharTok{|\textgreater{}}
  \FunctionTok{group\_by}\NormalTok{(industry\_ID) }\SpecialCharTok{|\textgreater{}}
  \FunctionTok{summarise}\NormalTok{(}
    \AttributeTok{firms =} \FunctionTok{n\_distinct}\NormalTok{(firm\_ID),}
    \AttributeTok{observations =} \FunctionTok{sum}\NormalTok{(}\SpecialCharTok{!}\FunctionTok{is.na}\NormalTok{(ROE)),}
    \AttributeTok{mean\_ROE =} \FunctionTok{mean}\NormalTok{(ROE, }\AttributeTok{na.rm =} \ConstantTok{TRUE}\NormalTok{),}
    \AttributeTok{median\_ROE =} \FunctionTok{median}\NormalTok{(ROE, }\AttributeTok{na.rm =} \ConstantTok{TRUE}\NormalTok{),}
    \AttributeTok{.groups =} \StringTok{"drop"}
\NormalTok{  ) }\SpecialCharTok{|\textgreater{}}
  \FunctionTok{filter}\NormalTok{(observations }\SpecialCharTok{\textgreater{}} \DecValTok{0}\NormalTok{) }\SpecialCharTok{|\textgreater{}}
  \FunctionTok{arrange}\NormalTok{(}\FunctionTok{desc}\NormalTok{(mean\_ROE))}

\FunctionTok{show\_table\_LFL}\NormalTok{(industry\_roe\_summary, }\AttributeTok{n =} \DecValTok{20}\NormalTok{)}
\end{Highlighting}
\end{Shaded}

\begin{verbatim}
# A tibble: 10 x 5
   industry_ID firms observations mean_ROE median_ROE
   <fct>       <int>        <int>    <dbl>      <dbl>
 1 2              40          186 0.10762    0.10283 
 2 6              34          127 0.095142   0.091981
 3 8              80          349 0.093711   0.084671
 4 5              67          259 0.085703   0.084851
 5 9             136          531 0.083403   0.079794
 6 1             225          918 0.077285   0.075118
 7 3             327         1433 0.075489   0.070116
 8 10            314         1388 0.075273   0.070490
 9 4              34          138 0.073719   0.072702
10 7             258         1076 0.067555   0.060409
\end{verbatim}

産業ごとの標本数が小さい場合、平均値は一部企業の影響を強く受ける。ここではROEの観測数が50以上の産業だけを比較する。

\begin{Shaded}
\begin{Highlighting}[]
\NormalTok{p\_industry\_roe }\OtherTok{\textless{}{-}}\NormalTok{ industry\_roe\_summary }\SpecialCharTok{|\textgreater{}}
  \FunctionTok{filter}\NormalTok{(observations }\SpecialCharTok{\textgreater{}=} \DecValTok{50}\NormalTok{) }\SpecialCharTok{|\textgreater{}}
  \FunctionTok{slice\_max}\NormalTok{(}
    \AttributeTok{order\_by =}\NormalTok{ mean\_ROE,}
    \AttributeTok{n =} \DecValTok{15}\NormalTok{,}
    \AttributeTok{with\_ties =} \ConstantTok{FALSE}
\NormalTok{  ) }\SpecialCharTok{|\textgreater{}}
  \FunctionTok{ggplot}\NormalTok{(}
    \FunctionTok{aes}\NormalTok{(}
      \AttributeTok{x =} \FunctionTok{reorder}\NormalTok{(industry\_ID, mean\_ROE),}
      \AttributeTok{y =}\NormalTok{ mean\_ROE,}
      \AttributeTok{fill =}\NormalTok{ industry\_ID}
\NormalTok{    )}
\NormalTok{  ) }\SpecialCharTok{+}
  \FunctionTok{geom\_col}\NormalTok{(}\AttributeTok{show.legend =} \ConstantTok{FALSE}\NormalTok{) }\SpecialCharTok{+}
  \FunctionTok{coord\_flip}\NormalTok{() }\SpecialCharTok{+}
  \FunctionTok{geom\_hline}\NormalTok{(}\AttributeTok{yintercept =} \DecValTok{0}\NormalTok{) }\SpecialCharTok{+}
  \FunctionTok{labs}\NormalTok{(}
    \AttributeTok{x =} \StringTok{"Industry ID"}\NormalTok{,}
    \AttributeTok{y =} \StringTok{"Mean ROE"}\NormalTok{,}
    \AttributeTok{title =} \StringTok{"Industries with High Mean ROE"}
\NormalTok{  )}

\FunctionTok{add\_lagrange\_theme\_LFL}\NormalTok{(p\_industry\_roe) }\SpecialCharTok{+}
  \FunctionTok{scale\_y\_continuous}\NormalTok{(}
    \AttributeTok{position =} \StringTok{"right"}\NormalTok{,}
    \AttributeTok{labels =}\NormalTok{ scales}\SpecialCharTok{::}\FunctionTok{label\_percent}\NormalTok{()}
\NormalTok{  )}
\end{Highlighting}
\end{Shaded}

\pandocbounded{\includegraphics[keepaspectratio]{part1_empirical/ch02_financial_data_files/figure-pdf/plot-industry-roe-1.pdf}}

\subsection{2.4.3
企業の順位付け}\label{ux4f01ux696dux306eux9806ux4f4dux4ed8ux3051}

最新年度について、産業内でROEが高い企業から順に順位を付ける。\texttt{dense\_rank()}を使うと、同順位がある場合でも次の順位を飛ばさない。

\begin{Shaded}
\begin{Highlighting}[]
\NormalTok{latest\_roe\_ranking }\OtherTok{\textless{}{-}}\NormalTok{ financial\_data }\SpecialCharTok{|\textgreater{}}
  \FunctionTok{filter}\NormalTok{(}
\NormalTok{    year }\SpecialCharTok{==}\NormalTok{ latest\_year,}
    \SpecialCharTok{!}\FunctionTok{is.na}\NormalTok{(ROE)}
\NormalTok{  ) }\SpecialCharTok{|\textgreater{}}
  \FunctionTok{group\_by}\NormalTok{(industry\_ID) }\SpecialCharTok{|\textgreater{}}
  \FunctionTok{mutate}\NormalTok{(}
    \AttributeTok{ROE\_rank =} \FunctionTok{dense\_rank}\NormalTok{(}\FunctionTok{desc}\NormalTok{(ROE))}
\NormalTok{  ) }\SpecialCharTok{|\textgreater{}}
  \FunctionTok{ungroup}\NormalTok{() }\SpecialCharTok{|\textgreater{}}
  \FunctionTok{arrange}\NormalTok{(industry\_ID, ROE\_rank)}

\NormalTok{latest\_roe\_ranking }\SpecialCharTok{|\textgreater{}}
  \FunctionTok{select}\NormalTok{(}
\NormalTok{    industry\_ID,}
\NormalTok{    ROE\_rank,}
\NormalTok{    firm\_ID,}
\NormalTok{    ROE}
\NormalTok{  ) }\SpecialCharTok{|\textgreater{}}
  \FunctionTok{slice\_head}\NormalTok{(}\AttributeTok{n =} \DecValTok{20}\NormalTok{) }\SpecialCharTok{|\textgreater{}}
  \FunctionTok{show\_table\_LFL}\NormalTok{(}\AttributeTok{n =} \DecValTok{20}\NormalTok{)}
\end{Highlighting}
\end{Shaded}

\begin{verbatim}
# A tibble: 20 x 4
   industry_ID ROE_rank firm_ID     ROE
   <fct>          <int> <fct>     <dbl>
 1 1                  1 8       0.38826
 2 1                  2 90      0.33395
 3 1                  3 225     0.30061
 4 1                  4 81      0.23113
 5 1                  5 171     0.21847
 6 1                  6 138     0.20680
 7 1                  7 66      0.20335
 8 1                  8 153     0.19822
 9 1                  9 181     0.19080
10 1                 10 65      0.18881
11 1                 11 92      0.17455
12 1                 12 195     0.15685
13 1                 13 13      0.15474
14 1                 14 23      0.15316
15 1                 15 148     0.14986
16 1                 16 214     0.14639
17 1                 17 115     0.14620
18 1                 18 201     0.14286
19 1                 19 50      0.14212
20 1                 20 143     0.14211
\end{verbatim}

各産業でROEが最も高い企業を抽出する。

\begin{Shaded}
\begin{Highlighting}[]
\NormalTok{top\_firm\_by\_industry }\OtherTok{\textless{}{-}}\NormalTok{ latest\_roe\_ranking }\SpecialCharTok{|\textgreater{}}
  \FunctionTok{group\_by}\NormalTok{(industry\_ID) }\SpecialCharTok{|\textgreater{}}
  \FunctionTok{slice\_max}\NormalTok{(}
    \AttributeTok{order\_by =}\NormalTok{ ROE,}
    \AttributeTok{n =} \DecValTok{1}\NormalTok{,}
    \AttributeTok{with\_ties =} \ConstantTok{FALSE}
\NormalTok{  ) }\SpecialCharTok{|\textgreater{}}
  \FunctionTok{ungroup}\NormalTok{() }\SpecialCharTok{|\textgreater{}}
  \FunctionTok{select}\NormalTok{(}
\NormalTok{    industry\_ID,}
\NormalTok{    firm\_ID,}
\NormalTok{    ROE}
\NormalTok{  )}

\FunctionTok{show\_table\_LFL}\NormalTok{(top\_firm\_by\_industry, }\AttributeTok{n =} \DecValTok{20}\NormalTok{)}
\end{Highlighting}
\end{Shaded}

\begin{verbatim}
# A tibble: 10 x 3
   industry_ID firm_ID     ROE
   <fct>       <fct>     <dbl>
 1 1           8       0.38826
 2 2           242     0.37500
 3 3           475     0.49754
 4 4           619     0.14913
 5 5           661     0.26731
 6 6           719     0.14220
 7 7           929     0.56418
 8 8           1042    0.25600
 9 9           1167    0.23462
10 10          1380    0.24979
\end{verbatim}

\section{2.5
上級デュポン・モデル}\label{ux4e0aux7d1aux30c7ux30e5ux30ddux30f3ux30e2ux30c7ux30eb}

\subsection{2.5.1
ROEのデュポン分解}\label{roeux306eux30c7ux30e5ux30ddux30f3ux5206ux89e3}

上級デュポン・モデルでは、ROEを純営業資産利益率と財務レバレッジの効果へ分解する。

\[
ROE_t = RNOA_t + FLEV_{t-1}(RNOA_t - NBC_t)
\]

各構成要素を次のように定義する。

\[
RNOA_t = \frac{OX_t}{NOA_{t-1}}
\]

\[
PM_t = \frac{OX_t}{Sales_t}
\]

\[
ATO_t = \frac{Sales_t}{NOA_{t-1}}
\]

\[
FLEV_{t-1} = \frac{NFO_{t-1}}{BE_{t-1}}
\]

\[
NBC_t = \frac{NFE_t}{NFO_{t-1}}
\]

\begin{Shaded}
\begin{Highlighting}[]
\NormalTok{financial\_data\_dupont }\OtherTok{\textless{}{-}}\NormalTok{ financial\_data }\SpecialCharTok{|\textgreater{}}
  \FunctionTok{group\_by}\NormalTok{(firm\_ID) }\SpecialCharTok{|\textgreater{}}
  \FunctionTok{arrange}\NormalTok{(year, }\AttributeTok{.by\_group =} \ConstantTok{TRUE}\NormalTok{) }\SpecialCharTok{|\textgreater{}}
  \FunctionTok{mutate}\NormalTok{(}
    \AttributeTok{lagged\_NOA =} \FunctionTok{lag}\NormalTok{(NOA),}
    \AttributeTok{lagged\_NFO =} \FunctionTok{lag}\NormalTok{(NFO),}
    \AttributeTok{RNOA =} \FunctionTok{divide\_safely\_LFL}\NormalTok{(OX, lagged\_NOA),}
    \AttributeTok{PM =} \FunctionTok{divide\_safely\_LFL}\NormalTok{(OX, sales),}
    \AttributeTok{ATO =} \FunctionTok{divide\_safely\_LFL}\NormalTok{(sales, lagged\_NOA),}
    \AttributeTok{lagged\_FLEV =} \FunctionTok{divide\_safely\_LFL}\NormalTok{(lagged\_NFO, lagged\_BE),}
    \AttributeTok{NBC =} \FunctionTok{divide\_safely\_LFL}\NormalTok{(NFE, lagged\_NFO),}
    \AttributeTok{ROE\_DuPont =}\NormalTok{ RNOA }\SpecialCharTok{+}\NormalTok{ lagged\_FLEV }\SpecialCharTok{*}\NormalTok{ (RNOA }\SpecialCharTok{{-}}\NormalTok{ NBC),}
    \AttributeTok{DuPont\_gap =}\NormalTok{ ROE }\SpecialCharTok{{-}}\NormalTok{ ROE\_DuPont}
\NormalTok{  ) }\SpecialCharTok{|\textgreater{}}
  \FunctionTok{ungroup}\NormalTok{()}

\NormalTok{financial\_data\_dupont }\SpecialCharTok{|\textgreater{}}
  \FunctionTok{select}\NormalTok{(}
\NormalTok{    firm\_ID,}
\NormalTok{    year,}
\NormalTok{    ROE,}
\NormalTok{    RNOA,}
\NormalTok{    PM,}
\NormalTok{    ATO,}
\NormalTok{    lagged\_FLEV,}
\NormalTok{    NBC,}
\NormalTok{    ROE\_DuPont,}
\NormalTok{    DuPont\_gap}
\NormalTok{  ) }\SpecialCharTok{|\textgreater{}}
  \FunctionTok{slice\_head}\NormalTok{(}\AttributeTok{n =} \DecValTok{10}\NormalTok{) }\SpecialCharTok{|\textgreater{}}
  \FunctionTok{show\_table\_LFL}\NormalTok{(}\AttributeTok{n =} \DecValTok{10}\NormalTok{)}
\end{Highlighting}
\end{Shaded}

\begin{verbatim}
# A tibble: 10 x 10
   firm_ID  year       ROE      RNOA       PM      ATO lagged_FLEV       NBC
   <fct>   <int>     <dbl>     <dbl>    <dbl>    <dbl>       <dbl>     <dbl>
 1 1        2015 NA        NA        0.083151 NA         NA        NA       
 2 1        2016  0.052962  0.065350 0.094830  0.68913   -0.10961  -0.047678
 3 1        2017  0.066073  0.074071 0.10625   0.69712   -0.068152 -0.043289
 4 1        2018  0.063762  0.084974 0.10974   0.77436   -0.15202  -0.054567
 5 1        2019  0.060919  0.070062 0.12658   0.55348    0.26187   0.10498 
 6 1        2020  0.076193  0.061279 0.10326   0.59342    0.14692  -0.040231
 7 2        2015 NA        NA        0.013070 NA         NA        NA       
 8 2        2016  0.046800  0.045056 0.014679  3.0693     0.078263  0.022729
 9 2        2017  0.070174  0.058935 0.021145  2.7872     0.30884   0.022536
10 2        2018  0.076295  0.064999 0.022564  2.8807     0.26625   0.022561
   ROE_DuPont DuPont_gap
        <dbl>      <dbl>
 1  NA        NA        
 2   0.052961  7.7332e-7
 3   0.066072  3.1527e-7
 4   0.063762 -3.3567e-7
 5   0.060919 -2.0862e-8
 6   0.076193  1.1089e-7
 7  NA        NA        
 8   0.046803 -3.2789e-6
 9   0.070176 -1.9587e-6
10   0.076298 -2.9722e-6
\end{verbatim}

デュポン分解で再構成したROEと、直接計算したROEの差を確認する。

\begin{Shaded}
\begin{Highlighting}[]
\NormalTok{dupont\_validation }\OtherTok{\textless{}{-}}\NormalTok{ financial\_data\_dupont }\SpecialCharTok{|\textgreater{}}
  \FunctionTok{summarise}\NormalTok{(}
    \AttributeTok{comparable\_observations =} \FunctionTok{sum}\NormalTok{(}\SpecialCharTok{!}\FunctionTok{is.na}\NormalTok{(DuPont\_gap)),}
    \AttributeTok{mean\_absolute\_gap =} \FunctionTok{mean}\NormalTok{(}\FunctionTok{abs}\NormalTok{(DuPont\_gap), }\AttributeTok{na.rm =} \ConstantTok{TRUE}\NormalTok{),}
    \AttributeTok{maximum\_absolute\_gap =} \FunctionTok{max}\NormalTok{(}\FunctionTok{abs}\NormalTok{(DuPont\_gap), }\AttributeTok{na.rm =} \ConstantTok{TRUE}\NormalTok{)}
\NormalTok{  )}

\FunctionTok{show\_table\_LFL}\NormalTok{(dupont\_validation)}
\end{Highlighting}
\end{Shaded}

\begin{verbatim}
# A tibble: 1 x 3
  comparable_observations mean_absolute_gap maximum_absolute_gap
                    <int>             <dbl>                <dbl>
1                    6405     0.00000036675          0.000020751
\end{verbatim}

\subsection{2.5.2 PMとATOの関係}\label{pmux3068atoux306eux95a2ux4fc2}

営業利益率PMと純営業資産回転率ATOは、企業の事業モデルを異なる側面から表す。産業別中央値を使って両者の関係を確認する。

\begin{Shaded}
\begin{Highlighting}[]
\NormalTok{industry\_pm\_ato }\OtherTok{\textless{}{-}}\NormalTok{ financial\_data\_dupont }\SpecialCharTok{|\textgreater{}}
  \FunctionTok{group\_by}\NormalTok{(industry\_ID) }\SpecialCharTok{|\textgreater{}}
  \FunctionTok{summarise}\NormalTok{(}
    \AttributeTok{observations =} \FunctionTok{sum}\NormalTok{(}\FunctionTok{complete.cases}\NormalTok{(PM, ATO)),}
    \AttributeTok{median\_PM =} \FunctionTok{median}\NormalTok{(PM, }\AttributeTok{na.rm =} \ConstantTok{TRUE}\NormalTok{),}
    \AttributeTok{median\_ATO =} \FunctionTok{median}\NormalTok{(ATO, }\AttributeTok{na.rm =} \ConstantTok{TRUE}\NormalTok{),}
    \AttributeTok{median\_RNOA =} \FunctionTok{median}\NormalTok{(RNOA, }\AttributeTok{na.rm =} \ConstantTok{TRUE}\NormalTok{),}
    \AttributeTok{.groups =} \StringTok{"drop"}
\NormalTok{  ) }\SpecialCharTok{|\textgreater{}}
  \FunctionTok{filter}\NormalTok{(}
\NormalTok{    observations }\SpecialCharTok{\textgreater{}=} \DecValTok{50}\NormalTok{,}
    \FunctionTok{is.finite}\NormalTok{(median\_PM),}
    \FunctionTok{is.finite}\NormalTok{(median\_ATO)}
\NormalTok{  )}

\FunctionTok{show\_table\_LFL}\NormalTok{(industry\_pm\_ato, }\AttributeTok{n =} \DecValTok{20}\NormalTok{)}
\end{Highlighting}
\end{Shaded}

\begin{verbatim}
# A tibble: 10 x 5
   industry_ID observations median_PM median_ATO median_RNOA
   <fct>              <int>     <dbl>      <dbl>       <dbl>
 1 1                    918  0.025331     3.9378    0.083905
 2 2                    186  0.037769     2.4106    0.10883 
 3 3                   1433  0.048140     1.8342    0.089713
 4 4                    138  0.040086     1.7797    0.075919
 5 5                    259  0.034620     2.5099    0.097564
 6 6                    127  0.048194     1.2994    0.074109
 7 7                   1076  0.031837     2.5163    0.080634
 8 8                    349  0.074662     1.0662    0.063506
 9 9                    531  0.042643     2.1790    0.098533
10 10                  1388  0.044615     1.9604    0.090718
\end{verbatim}

\begin{Shaded}
\begin{Highlighting}[]
\NormalTok{p\_pm\_ato\_tradeoff }\OtherTok{\textless{}{-}}\NormalTok{ industry\_pm\_ato }\SpecialCharTok{|\textgreater{}}
  \FunctionTok{ggplot}\NormalTok{(}
    \FunctionTok{aes}\NormalTok{(}
      \AttributeTok{x =}\NormalTok{ median\_ATO,}
      \AttributeTok{y =}\NormalTok{ median\_PM,}
      \AttributeTok{colour =}\NormalTok{ median\_RNOA}
\NormalTok{    )}
\NormalTok{  ) }\SpecialCharTok{+}
  \FunctionTok{geom\_point}\NormalTok{(}\AttributeTok{size =} \FloatTok{2.5}\NormalTok{) }\SpecialCharTok{+}
  \FunctionTok{labs}\NormalTok{(}
    \AttributeTok{x =} \StringTok{"Industry Median ATO"}\NormalTok{,}
    \AttributeTok{y =} \StringTok{"Industry Median PM"}\NormalTok{,}
    \AttributeTok{colour =} \StringTok{"Median RNOA"}\NormalTok{,}
    \AttributeTok{title =} \StringTok{"Profit Margin and Asset Turnover by Industry"}
\NormalTok{  )}

\FunctionTok{add\_lagrange\_theme\_LFL}\NormalTok{(p\_pm\_ato\_tradeoff) }\SpecialCharTok{+}
  \FunctionTok{scale\_y\_continuous}\NormalTok{(}
    \AttributeTok{position =} \StringTok{"right"}\NormalTok{,}
    \AttributeTok{labels =}\NormalTok{ scales}\SpecialCharTok{::}\FunctionTok{label\_percent}\NormalTok{()}
\NormalTok{  )}
\end{Highlighting}
\end{Shaded}

\pandocbounded{\includegraphics[keepaspectratio]{part1_empirical/ch02_financial_data_files/figure-pdf/plot-pm-ato-tradeoff-1.pdf}}

主要産業についてPMの分布を箱ひげ図で比較する。

\begin{Shaded}
\begin{Highlighting}[]
\NormalTok{major\_industries }\OtherTok{\textless{}{-}}\NormalTok{ firms\_by\_industry }\SpecialCharTok{|\textgreater{}}
  \FunctionTok{slice\_head}\NormalTok{(}\AttributeTok{n =} \DecValTok{8}\NormalTok{) }\SpecialCharTok{|\textgreater{}}
  \FunctionTok{pull}\NormalTok{(industry\_ID)}

\NormalTok{pm\_limits }\OtherTok{\textless{}{-}} \FunctionTok{quantile}\NormalTok{(}
\NormalTok{  financial\_data\_dupont}\SpecialCharTok{$}\NormalTok{PM,}
  \AttributeTok{probs =} \FunctionTok{c}\NormalTok{(}\FloatTok{0.01}\NormalTok{, }\FloatTok{0.99}\NormalTok{),}
  \AttributeTok{na.rm =} \ConstantTok{TRUE}
\NormalTok{)}

\NormalTok{p\_pm\_by\_industry }\OtherTok{\textless{}{-}}\NormalTok{ financial\_data\_dupont }\SpecialCharTok{|\textgreater{}}
  \FunctionTok{filter}\NormalTok{(}
\NormalTok{    year }\SpecialCharTok{==}\NormalTok{ latest\_year,}
\NormalTok{    industry\_ID }\SpecialCharTok{\%in\%}\NormalTok{ major\_industries,}
    \SpecialCharTok{!}\FunctionTok{is.na}\NormalTok{(PM)}
\NormalTok{  ) }\SpecialCharTok{|\textgreater{}}
  \FunctionTok{ggplot}\NormalTok{(}
    \FunctionTok{aes}\NormalTok{(}
      \AttributeTok{x =}\NormalTok{ industry\_ID,}
      \AttributeTok{y =}\NormalTok{ PM,}
      \AttributeTok{fill =}\NormalTok{ industry\_ID}
\NormalTok{    )}
\NormalTok{  ) }\SpecialCharTok{+}
  \FunctionTok{geom\_boxplot}\NormalTok{(}
    \AttributeTok{outlier.alpha =} \FloatTok{0.25}\NormalTok{,}
    \AttributeTok{show.legend =} \ConstantTok{FALSE}
\NormalTok{  ) }\SpecialCharTok{+}
  \FunctionTok{coord\_cartesian}\NormalTok{(}\AttributeTok{ylim =}\NormalTok{ pm\_limits) }\SpecialCharTok{+}
  \FunctionTok{labs}\NormalTok{(}
    \AttributeTok{x =} \StringTok{"Industry ID"}\NormalTok{,}
    \AttributeTok{y =} \StringTok{"Profit Margin"}\NormalTok{,}
    \AttributeTok{title =} \FunctionTok{paste}\NormalTok{(}
      \StringTok{"Profit Margin by Industry in"}\NormalTok{,}
\NormalTok{      latest\_year}
\NormalTok{    )}
\NormalTok{  )}

\FunctionTok{add\_lagrange\_theme\_LFL}\NormalTok{(p\_pm\_by\_industry) }\SpecialCharTok{+}
  \FunctionTok{scale\_y\_continuous}\NormalTok{(}
    \AttributeTok{position =} \StringTok{"right"}\NormalTok{,}
    \AttributeTok{labels =}\NormalTok{ scales}\SpecialCharTok{::}\FunctionTok{label\_percent}\NormalTok{()}
\NormalTok{  )}
\end{Highlighting}
\end{Shaded}

\pandocbounded{\includegraphics[keepaspectratio]{part1_empirical/ch02_financial_data_files/figure-pdf/plot-pm-by-industry-1.pdf}}

\section{2.6
財務指標の要約と中間データ}\label{ux8ca1ux52d9ux6307ux6a19ux306eux8981ux7d04ux3068ux4e2dux9593ux30c7ux30fcux30bf}

\subsection{2.6.1
複数の財務指標の要約}\label{ux8907ux6570ux306eux8ca1ux52d9ux6307ux6a19ux306eux8981ux7d04}

複数列に同じ集計を適用する場合は、列名のベクトルと\texttt{purrr::map\_dfr()}を組み合わせる。

\begin{Shaded}
\begin{Highlighting}[]
\NormalTok{financial\_metrics }\OtherTok{\textless{}{-}} \FunctionTok{c}\NormalTok{(}
  \StringTok{"sales"}\NormalTok{,}
  \StringTok{"BE"}\NormalTok{,}
  \StringTok{"ROE"}\NormalTok{,}
  \StringTok{"RNOA"}\NormalTok{,}
  \StringTok{"PM"}\NormalTok{,}
  \StringTok{"ATO"}
\NormalTok{)}

\NormalTok{metric\_summary }\OtherTok{\textless{}{-}}\NormalTok{ financial\_metrics }\SpecialCharTok{|\textgreater{}}
  \FunctionTok{map\_dfr}\NormalTok{(}
\NormalTok{    \textbackslash{}(metric\_name) \{}
\NormalTok{      values }\OtherTok{\textless{}{-}}\NormalTok{ financial\_data\_dupont[[metric\_name]]}
\NormalTok{      finite\_values }\OtherTok{\textless{}{-}}\NormalTok{ values[}\FunctionTok{is.finite}\NormalTok{(values)]}

      \FunctionTok{tibble}\NormalTok{(}
        \AttributeTok{metric =}\NormalTok{ metric\_name,}
        \AttributeTok{observations =} \FunctionTok{length}\NormalTok{(finite\_values),}
        \AttributeTok{mean =} \FunctionTok{mean}\NormalTok{(finite\_values),}
        \AttributeTok{median =} \FunctionTok{median}\NormalTok{(finite\_values),}
        \AttributeTok{standard\_deviation =} \FunctionTok{sd}\NormalTok{(finite\_values),}
        \AttributeTok{minimum =} \FunctionTok{min}\NormalTok{(finite\_values),}
        \AttributeTok{maximum =} \FunctionTok{max}\NormalTok{(finite\_values)}
\NormalTok{      )}
\NormalTok{    \}}
\NormalTok{  )}

\FunctionTok{show\_table\_LFL}\NormalTok{(metric\_summary, }\AttributeTok{n =} \FunctionTok{nrow}\NormalTok{(metric\_summary))}
\end{Highlighting}
\end{Shaded}

\begin{verbatim}
# A tibble: 6 x 7
  metric observations          mean       median standard_deviation   minimum
  <chr>         <int>         <dbl>        <dbl>              <dbl>     <dbl>
1 sales          7920 166007.       40431.            381980.       205.34   
2 BE             7920 111517.       26882.            420624.       136.93   
3 ROE            6405      0.077713     0.072828           0.066274  -0.70192
4 RNOA           6405      0.11141      0.085868           0.12622   -0.82036
5 PM             7920      0.045534     0.040555           0.042223  -0.20079
6 ATO            6405      2.7051       2.1024             1.8684     0.25825
    maximum
      <dbl>
1 3.4964e+6
2 2.5448e+7
3 1.4055e+0
4 1.8283e+0
5 3.8289e-1
6 9.8170e+0
\end{verbatim}

\subsection{2.6.2
中間データの保存}\label{ux4e2dux9593ux30c7ux30fcux30bfux306eux4fddux5b58}

第3章では、本章で計算した株主資本、前年度株主資本およびROEなどを株式データと結合する。そのため、上級デュポン・モデルの計算結果を含む整理済みデータを\texttt{simulation\_data/derived/ch04\_output.csv}へ保存する。

\begin{Shaded}
\begin{Highlighting}[]
\NormalTok{financial\_output\_path }\OtherTok{\textless{}{-}}
  \FunctionTok{write\_derived\_csv\_file\_LFL}\NormalTok{(}
\NormalTok{    financial\_data\_dupont,}
    \StringTok{"ch04\_output.csv"}
\NormalTok{  )}

\NormalTok{financial\_output\_path}
\end{Highlighting}
\end{Shaded}

\begin{verbatim}
[1] "C:/AnalyticFin/Projects/LagrangeFinance/simulation_data/derived/ch04_output.csv"
\end{verbatim}

\begin{Shaded}
\begin{Highlighting}[]
\FunctionTok{file.info}\NormalTok{(financial\_output\_path)[, }\FunctionTok{c}\NormalTok{(}\StringTok{"size"}\NormalTok{, }\StringTok{"mtime"}\NormalTok{)]}
\end{Highlighting}
\end{Shaded}

\begin{verbatim}
                                                                                   size
C:/AnalyticFin/Projects/LagrangeFinance/simulation_data/derived/ch04_output.csv 2425789
                                                                                              mtime
C:/AnalyticFin/Projects/LagrangeFinance/simulation_data/derived/ch04_output.csv 2026-07-30 01:48:35
\end{verbatim}

保存したCSVを読み直し、行数と列名が一致することを確認する。

\begin{Shaded}
\begin{Highlighting}[]
\NormalTok{financial\_output\_check }\OtherTok{\textless{}{-}}
  \FunctionTok{read\_derived\_csv\_file\_LFL}\NormalTok{(}\StringTok{"ch04\_output.csv"}\NormalTok{)}

\FunctionTok{stopifnot}\NormalTok{(}
  \FunctionTok{nrow}\NormalTok{(financial\_output\_check) }\SpecialCharTok{==}
    \FunctionTok{nrow}\NormalTok{(financial\_data\_dupont)}
\NormalTok{)}

\FunctionTok{stopifnot}\NormalTok{(}
  \FunctionTok{identical}\NormalTok{(}
    \FunctionTok{names}\NormalTok{(financial\_output\_check),}
    \FunctionTok{names}\NormalTok{(financial\_data\_dupont)}
\NormalTok{  )}
\NormalTok{)}

\NormalTok{output\_summary }\OtherTok{\textless{}{-}} \FunctionTok{tibble}\NormalTok{(}
  \AttributeTok{output\_file =}\NormalTok{ financial\_output\_path,}
  \AttributeTok{rows =} \FunctionTok{nrow}\NormalTok{(financial\_output\_check),}
  \AttributeTok{columns =} \FunctionTok{ncol}\NormalTok{(financial\_output\_check)}
\NormalTok{)}

\FunctionTok{show\_table\_LFL}\NormalTok{(output\_summary)}
\end{Highlighting}
\end{Shaded}

\begin{verbatim}
# A tibble: 1 x 3
  output_file                                                                   
  <chr>                                                                         
1 C:/AnalyticFin/Projects/LagrangeFinance/simulation_data/derived/ch04_output.c~
   rows columns
  <int>   <int>
1  7920      25
\end{verbatim}

\section{2.7 まとめ}\label{ux307eux3068ux3081-1}

本章では、財務データの構造、欠損値およびパネルデータのキーを確認した。次に、年度別・産業別の企業構成、売上高の分布、株主資本、ROEおよび産業内順位を計算した。最後に、ROEをRNOA、財務レバレッジおよび純借入コストへ分解し、第3章で利用する整理済みデータを\texttt{simulation\_data/derived/ch04\_output.csv}へ保存した。

第3章では、株式データからリターンを計算し、本章で作成した財務データと結合する。

\chapter{第3章
株式データと時系列分析}\label{ux7b2c3ux7ae0-ux682aux5f0fux30c7ux30fcux30bfux3068ux6642ux7cfbux5217ux5206ux6790}

本章では、月次株式データから株式リターン、超過リターンおよび時価総額を作成し、第2章で整理した年次財務データと結合する。さらに、月次リターンの分布、累積リターン、年次リターンおよび簿価時価比率との関係を確認し、第4章以降で利用する分析用データを保存する。

\section{3.1 分析環境}\label{ux5206ux6790ux74b0ux5883-1}

\begin{Shaded}
\begin{Highlighting}[]
\FunctionTok{source}\NormalTok{(}\FunctionTok{file.path}\NormalTok{(}\StringTok{"R"}\NormalTok{, }\StringTok{"utility\_LFL.R"}\NormalTok{))}
\FunctionTok{load\_part1\_packages\_LFL}\NormalTok{()}
\end{Highlighting}
\end{Shaded}

\section{3.2
株式データの読込みと整備}\label{ux682aux5f0fux30c7ux30fcux30bfux306eux8aadux8fbcux307fux3068ux6574ux5099}

\subsection{3.2.1
株式データの読込み}\label{ux682aux5f0fux30c7ux30fcux30bfux306eux8aadux8fbcux307f}

\begin{Shaded}
\begin{Highlighting}[]
\NormalTok{stock\_data\_raw }\OtherTok{\textless{}{-}}
  \FunctionTok{read\_csv\_file\_LFL}\NormalTok{(}\StringTok{"ch05\_stock\_data.csv"}\NormalTok{)}

\FunctionTok{show\_table\_LFL}\NormalTok{(stock\_data\_raw, }\AttributeTok{n =} \DecValTok{10}\NormalTok{)}
\end{Highlighting}
\end{Shaded}

\begin{verbatim}
# A tibble: 95,040 x 9
    year month month_ID firm_ID stock_price   DPS shares_outstanding
   <dbl> <dbl>    <dbl>   <dbl>       <dbl> <dbl>              <dbl>
 1  2015     1        1       1         954     0            2422000
 2  2015     2        2       1         960     0            2422000
 3  2015     3        3       1        1113     0            2422000
 4  2015     4        4       1        1081     0            2422000
 5  2015     5        5       1        1317     0            2422000
 6  2015     6        6       1        1366    29            2422000
 7  2015     7        7       1        1353     0            2422000
 8  2015     8        8       1        1209     0            2422000
 9  2015     9        9       1        1291     0            2422000
10  2015    10       10       1        1407     0            2422000
   adjustment_coefficient        R_F
                    <dbl>      <dbl>
 1                      1 0.00065068
 2                      1 0.00058341
 3                      1 0.00061144
 4                      1 0.00068482
 5                      1 0.00073677
 6                      1 0.00069540
 7                      1 0.00049354
 8                      1 0.00047106
 9                      1 0.00055242
10                      1 0.00063803
# i 95,030 more rows
\end{verbatim}

本章の計算に必要な列がすべて存在することを確認する。

\begin{Shaded}
\begin{Highlighting}[]
\NormalTok{required\_stock\_columns }\OtherTok{\textless{}{-}} \FunctionTok{c}\NormalTok{(}
  \StringTok{"year"}\NormalTok{,}
  \StringTok{"month"}\NormalTok{,}
  \StringTok{"month\_ID"}\NormalTok{,}
  \StringTok{"firm\_ID"}\NormalTok{,}
  \StringTok{"stock\_price"}\NormalTok{,}
  \StringTok{"shares\_outstanding"}\NormalTok{,}
  \StringTok{"DPS"}\NormalTok{,}
  \StringTok{"adjustment\_coefficient"}\NormalTok{,}
  \StringTok{"R\_F"}
\NormalTok{)}

\NormalTok{missing\_stock\_columns }\OtherTok{\textless{}{-}}
  \FunctionTok{setdiff}\NormalTok{(required\_stock\_columns, }\FunctionTok{names}\NormalTok{(stock\_data\_raw))}

\ControlFlowTok{if}\NormalTok{ (}\FunctionTok{length}\NormalTok{(missing\_stock\_columns) }\SpecialCharTok{\textgreater{}} \DecValTok{0}\NormalTok{) \{}
  \FunctionTok{stop}\NormalTok{(}
    \StringTok{"株式データに必要な列が不足しています: "}\NormalTok{,}
    \FunctionTok{paste}\NormalTok{(missing\_stock\_columns, }\AttributeTok{collapse =} \StringTok{", "}\NormalTok{),}
    \AttributeTok{call. =} \ConstantTok{FALSE}
\NormalTok{  )}
\NormalTok{\}}
\end{Highlighting}
\end{Shaded}

\subsection{3.2.2
データ構造と欠損値}\label{ux30c7ux30fcux30bfux69cbux9020ux3068ux6b20ux640dux5024}

\begin{Shaded}
\begin{Highlighting}[]
\FunctionTok{show\_data\_structure\_LFL}\NormalTok{(stock\_data\_raw)}
\end{Highlighting}
\end{Shaded}

\begin{verbatim}
Rows: 95,040
Columns: 9
$ year                   <dbl> 2015, 2015, 2015, 2015, 2015, 2015, 2015, 2015,~
$ month                  <dbl> 1, 2, 3, 4, 5, 6, 7, 8, 9, 10, 11, 12, 1, 2, 3,~
$ month_ID               <dbl> 1, 2, 3, 4, 5, 6, 7, 8, 9, 10, 11, 12, 13, 14, ~
$ firm_ID                <dbl> 1, 1, 1, 1, 1, 1, 1, 1, 1, 1, 1, 1, 1, 1, 1, 1,~
$ stock_price            <dbl> 954, 960, 1113, 1081, 1317, 1366, 1353, 1209, 1~
$ DPS                    <dbl> 0, 0, 0, 0, 0, 29, 0, 0, 0, 0, 0, 29, 0, 0, 0, ~
$ shares_outstanding     <dbl> 2422000, 2422000, 2422000, 2422000, 2422000, 24~
$ adjustment_coefficient <dbl> 1, 1, 1, 1, 1, 1, 1, 1, 1, 1, 1, 1, 1, 1, 1, 1,~
$ R_F                    <dbl> 0.00065068255, 0.00058340985, 0.00061144232, 0.~
\end{verbatim}

\begin{Shaded}
\begin{Highlighting}[]
\NormalTok{stock\_data\_overview }\OtherTok{\textless{}{-}}\NormalTok{ stock\_data\_raw }\SpecialCharTok{|\textgreater{}}
  \FunctionTok{summarise}\NormalTok{(}
    \AttributeTok{observations =} \FunctionTok{n}\NormalTok{(),}
    \AttributeTok{variables =} \FunctionTok{ncol}\NormalTok{(stock\_data\_raw),}
    \AttributeTok{firms =} \FunctionTok{n\_distinct}\NormalTok{(firm\_ID),}
    \AttributeTok{first\_year =} \FunctionTok{min}\NormalTok{(year, }\AttributeTok{na.rm =} \ConstantTok{TRUE}\NormalTok{),}
    \AttributeTok{last\_year =} \FunctionTok{max}\NormalTok{(year, }\AttributeTok{na.rm =} \ConstantTok{TRUE}\NormalTok{),}
    \AttributeTok{first\_month\_ID =} \FunctionTok{min}\NormalTok{(month\_ID, }\AttributeTok{na.rm =} \ConstantTok{TRUE}\NormalTok{),}
    \AttributeTok{last\_month\_ID =} \FunctionTok{max}\NormalTok{(month\_ID, }\AttributeTok{na.rm =} \ConstantTok{TRUE}\NormalTok{)}
\NormalTok{  )}

\FunctionTok{show\_table\_LFL}\NormalTok{(stock\_data\_overview)}
\end{Highlighting}
\end{Shaded}

\begin{verbatim}
# A tibble: 1 x 7
  observations variables firms first_year last_year first_month_ID last_month_ID
         <int>     <int> <int>      <dbl>     <dbl>          <dbl>         <dbl>
1        95040         9  1515       2015      2020              1            72
\end{verbatim}

\begin{Shaded}
\begin{Highlighting}[]
\NormalTok{stock\_missing\_summary }\OtherTok{\textless{}{-}}\NormalTok{ stock\_data\_raw }\SpecialCharTok{|\textgreater{}}
  \FunctionTok{summarise}\NormalTok{(}
    \FunctionTok{across}\NormalTok{(}
      \FunctionTok{everything}\NormalTok{(),}
      \FunctionTok{list}\NormalTok{(}
        \AttributeTok{missing =}\NormalTok{ \textbackslash{}(x) }\FunctionTok{sum}\NormalTok{(}\FunctionTok{is.na}\NormalTok{(x)),}
        \AttributeTok{missing\_rate =}\NormalTok{ \textbackslash{}(x) }\FunctionTok{mean}\NormalTok{(}\FunctionTok{is.na}\NormalTok{(x))}
\NormalTok{      ),}
      \AttributeTok{.names =} \StringTok{"\{.col\}\_\_\{.fn\}"}
\NormalTok{    )}
\NormalTok{  ) }\SpecialCharTok{|\textgreater{}}
  \FunctionTok{pivot\_longer}\NormalTok{(}
    \AttributeTok{cols =} \FunctionTok{everything}\NormalTok{(),}
    \AttributeTok{names\_to =} \FunctionTok{c}\NormalTok{(}\StringTok{"variable"}\NormalTok{, }\StringTok{".value"}\NormalTok{),}
    \AttributeTok{names\_sep =} \StringTok{"\_\_"}
\NormalTok{  ) }\SpecialCharTok{|\textgreater{}}
  \FunctionTok{arrange}\NormalTok{(}\FunctionTok{desc}\NormalTok{(missing\_rate), variable)}

\FunctionTok{show\_table\_LFL}\NormalTok{(}
\NormalTok{  stock\_missing\_summary,}
  \AttributeTok{n =} \FunctionTok{nrow}\NormalTok{(stock\_missing\_summary)}
\NormalTok{)}
\end{Highlighting}
\end{Shaded}

\begin{verbatim}
# A tibble: 9 x 3
  variable               missing missing_rate
  <chr>                    <int>        <dbl>
1 DPS                          0            0
2 R_F                          0            0
3 adjustment_coefficient       0            0
4 firm_ID                      0            0
5 month                        0            0
6 month_ID                     0            0
7 shares_outstanding           0            0
8 stock_price                  0            0
9 year                         0            0
\end{verbatim}

\subsection{3.2.3
データ型とパネルキー}\label{ux30c7ux30fcux30bfux578bux3068ux30d1ux30cdux30ebux30adux30fc}

企業ごとの時系列計算を正しく行うため、企業IDと月IDの順に並べる。

\begin{Shaded}
\begin{Highlighting}[]
\NormalTok{stock\_data }\OtherTok{\textless{}{-}}\NormalTok{ stock\_data\_raw }\SpecialCharTok{|\textgreater{}}
  \FunctionTok{mutate}\NormalTok{(}
    \AttributeTok{firm\_ID =} \FunctionTok{as.integer}\NormalTok{(firm\_ID),}
    \AttributeTok{year =} \FunctionTok{as.integer}\NormalTok{(year),}
    \AttributeTok{month =} \FunctionTok{as.integer}\NormalTok{(month),}
    \AttributeTok{month\_ID =} \FunctionTok{as.integer}\NormalTok{(month\_ID),}
    \AttributeTok{month\_date =}\NormalTok{ lubridate}\SpecialCharTok{::}\FunctionTok{make\_date}\NormalTok{(year, month, }\DecValTok{1}\NormalTok{)}
\NormalTok{  ) }\SpecialCharTok{|\textgreater{}}
  \FunctionTok{arrange}\NormalTok{(firm\_ID, month\_ID)}

\FunctionTok{show\_data\_structure\_LFL}\NormalTok{(stock\_data)}
\end{Highlighting}
\end{Shaded}

\begin{verbatim}
Rows: 95,040
Columns: 10
$ year                   <int> 2015, 2015, 2015, 2015, 2015, 2015, 2015, 2015,~
$ month                  <int> 1, 2, 3, 4, 5, 6, 7, 8, 9, 10, 11, 12, 1, 2, 3,~
$ month_ID               <int> 1, 2, 3, 4, 5, 6, 7, 8, 9, 10, 11, 12, 13, 14, ~
$ firm_ID                <int> 1, 1, 1, 1, 1, 1, 1, 1, 1, 1, 1, 1, 1, 1, 1, 1,~
$ stock_price            <dbl> 954, 960, 1113, 1081, 1317, 1366, 1353, 1209, 1~
$ DPS                    <dbl> 0, 0, 0, 0, 0, 29, 0, 0, 0, 0, 0, 29, 0, 0, 0, ~
$ shares_outstanding     <dbl> 2422000, 2422000, 2422000, 2422000, 2422000, 24~
$ adjustment_coefficient <dbl> 1, 1, 1, 1, 1, 1, 1, 1, 1, 1, 1, 1, 1, 1, 1, 1,~
$ R_F                    <dbl> 0.00065068255, 0.00058340985, 0.00061144232, 0.~
$ month_date             <date> 2015-01-01, 2015-02-01, 2015-03-01, 2015-04-01~
\end{verbatim}

\texttt{firm\_ID}と\texttt{month\_ID}の組合せが一意であることを確認する。

\begin{Shaded}
\begin{Highlighting}[]
\NormalTok{stock\_key\_duplicates }\OtherTok{\textless{}{-}}\NormalTok{ stock\_data }\SpecialCharTok{|\textgreater{}}
  \FunctionTok{count}\NormalTok{(}
\NormalTok{    firm\_ID,}
\NormalTok{    month\_ID,}
    \AttributeTok{name =} \StringTok{"observations"}
\NormalTok{  ) }\SpecialCharTok{|\textgreater{}}
  \FunctionTok{filter}\NormalTok{(observations }\SpecialCharTok{\textgreater{}} \DecValTok{1}\NormalTok{)}

\FunctionTok{show\_table\_LFL}\NormalTok{(stock\_key\_duplicates)}
\end{Highlighting}
\end{Shaded}

\begin{verbatim}
# A tibble: 0 x 3
# i 3 variables: firm_ID <int>, month_ID <int>, observations <int>
\end{verbatim}

\begin{Shaded}
\begin{Highlighting}[]
\ControlFlowTok{if}\NormalTok{ (}\FunctionTok{nrow}\NormalTok{(stock\_key\_duplicates) }\SpecialCharTok{\textgreater{}} \DecValTok{0}\NormalTok{) \{}
  \FunctionTok{stop}\NormalTok{(}
    \StringTok{"firm\_IDとmonth\_IDの組合せに重複があります。"}\NormalTok{,}
    \AttributeTok{call. =} \ConstantTok{FALSE}
\NormalTok{  )}
\NormalTok{\}}
\end{Highlighting}
\end{Shaded}

同じ月IDに複数の年月が対応していないことも確認する。

\begin{Shaded}
\begin{Highlighting}[]
\NormalTok{month\_calendar\_check }\OtherTok{\textless{}{-}}\NormalTok{ stock\_data }\SpecialCharTok{|\textgreater{}}
  \FunctionTok{distinct}\NormalTok{(}
\NormalTok{    month\_ID,}
\NormalTok{    year,}
\NormalTok{    month,}
\NormalTok{    month\_date}
\NormalTok{  ) }\SpecialCharTok{|\textgreater{}}
  \FunctionTok{count}\NormalTok{(}
\NormalTok{    month\_ID,}
    \AttributeTok{name =} \StringTok{"calendar\_rows"}
\NormalTok{  ) }\SpecialCharTok{|\textgreater{}}
  \FunctionTok{filter}\NormalTok{(calendar\_rows }\SpecialCharTok{\textgreater{}} \DecValTok{1}\NormalTok{)}

\FunctionTok{show\_table\_LFL}\NormalTok{(month\_calendar\_check)}
\end{Highlighting}
\end{Shaded}

\begin{verbatim}
# A tibble: 0 x 2
# i 2 variables: month_ID <int>, calendar_rows <int>
\end{verbatim}

\begin{Shaded}
\begin{Highlighting}[]
\ControlFlowTok{if}\NormalTok{ (}\FunctionTok{nrow}\NormalTok{(month\_calendar\_check) }\SpecialCharTok{\textgreater{}} \DecValTok{0}\NormalTok{) \{}
  \FunctionTok{stop}\NormalTok{(}
    \StringTok{"同じmonth\_IDに複数の年月が対応しています。"}\NormalTok{,}
    \AttributeTok{call. =} \ConstantTok{FALSE}
\NormalTok{  )}
\NormalTok{\}}
\end{Highlighting}
\end{Shaded}

年度別の企業数と観測数を確認する。

\begin{Shaded}
\begin{Highlighting}[]
\NormalTok{stock\_observations\_by\_year }\OtherTok{\textless{}{-}}\NormalTok{ stock\_data }\SpecialCharTok{|\textgreater{}}
  \FunctionTok{group\_by}\NormalTok{(year) }\SpecialCharTok{|\textgreater{}}
  \FunctionTok{summarise}\NormalTok{(}
    \AttributeTok{firms =} \FunctionTok{n\_distinct}\NormalTok{(firm\_ID),}
    \AttributeTok{observations =} \FunctionTok{n}\NormalTok{(),}
    \AttributeTok{first\_month\_ID =} \FunctionTok{min}\NormalTok{(month\_ID),}
    \AttributeTok{last\_month\_ID =} \FunctionTok{max}\NormalTok{(month\_ID),}
    \AttributeTok{.groups =} \StringTok{"drop"}
\NormalTok{  ) }\SpecialCharTok{|\textgreater{}}
  \FunctionTok{arrange}\NormalTok{(year)}

\FunctionTok{show\_table\_LFL}\NormalTok{(}
\NormalTok{  stock\_observations\_by\_year,}
  \AttributeTok{n =} \FunctionTok{nrow}\NormalTok{(stock\_observations\_by\_year)}
\NormalTok{)}
\end{Highlighting}
\end{Shaded}

\begin{verbatim}
# A tibble: 6 x 5
   year firms observations first_month_ID last_month_ID
  <int> <int>        <int>          <int>         <int>
1  2015  1266        15192              1            12
2  2016  1293        15516             13            24
3  2017  1319        15828             25            36
4  2018  1323        15876             37            48
5  2019  1356        16272             49            60
6  2020  1363        16356             61            72
\end{verbatim}

\section{3.3
月次リターンと時価総額}\label{ux6708ux6b21ux30eaux30bfux30fcux30f3ux3068ux6642ux4fa1ux7dcfux984d}

\subsection{3.3.1
配当と調整係数}\label{ux914dux5f53ux3068ux8abfux6574ux4fc2ux6570}

配当が支払われた観測と、調整係数が1と異なる観測を確認する。

\begin{Shaded}
\begin{Highlighting}[]
\NormalTok{dividend\_observations }\OtherTok{\textless{}{-}}\NormalTok{ stock\_data }\SpecialCharTok{|\textgreater{}}
  \FunctionTok{filter}\NormalTok{(}
    \SpecialCharTok{!}\FunctionTok{is.na}\NormalTok{(DPS),}
\NormalTok{    DPS }\SpecialCharTok{!=} \DecValTok{0}
\NormalTok{  ) }\SpecialCharTok{|\textgreater{}}
  \FunctionTok{select}\NormalTok{(}
\NormalTok{    firm\_ID,}
\NormalTok{    year,}
\NormalTok{    month,}
\NormalTok{    month\_ID,}
\NormalTok{    stock\_price,}
\NormalTok{    DPS,}
\NormalTok{    adjustment\_coefficient}
\NormalTok{  ) }\SpecialCharTok{|\textgreater{}}
  \FunctionTok{arrange}\NormalTok{(firm\_ID, month\_ID)}

\NormalTok{adjustment\_observations }\OtherTok{\textless{}{-}}\NormalTok{ stock\_data }\SpecialCharTok{|\textgreater{}}
  \FunctionTok{filter}\NormalTok{(}
    \SpecialCharTok{!}\FunctionTok{is.na}\NormalTok{(adjustment\_coefficient),}
    \SpecialCharTok{!}\FunctionTok{near}\NormalTok{(adjustment\_coefficient, }\DecValTok{1}\NormalTok{)}
\NormalTok{  ) }\SpecialCharTok{|\textgreater{}}
  \FunctionTok{select}\NormalTok{(}
\NormalTok{    firm\_ID,}
\NormalTok{    year,}
\NormalTok{    month,}
\NormalTok{    month\_ID,}
\NormalTok{    stock\_price,}
\NormalTok{    DPS,}
\NormalTok{    adjustment\_coefficient}
\NormalTok{  ) }\SpecialCharTok{|\textgreater{}}
  \FunctionTok{arrange}\NormalTok{(firm\_ID, month\_ID)}

\NormalTok{corporate\_action\_summary }\OtherTok{\textless{}{-}}\NormalTok{ stock\_data }\SpecialCharTok{|\textgreater{}}
  \FunctionTok{summarise}\NormalTok{(}
    \AttributeTok{observations =} \FunctionTok{n}\NormalTok{(),}
    \AttributeTok{dividend\_observations =}
      \FunctionTok{sum}\NormalTok{(}\SpecialCharTok{!}\FunctionTok{is.na}\NormalTok{(DPS) }\SpecialCharTok{\&}\NormalTok{ DPS }\SpecialCharTok{!=} \DecValTok{0}\NormalTok{),}
    \AttributeTok{adjusted\_observations =}
      \FunctionTok{sum}\NormalTok{(}
        \SpecialCharTok{!}\FunctionTok{is.na}\NormalTok{(adjustment\_coefficient) }\SpecialCharTok{\&}
          \SpecialCharTok{!}\FunctionTok{near}\NormalTok{(adjustment\_coefficient, }\DecValTok{1}\NormalTok{)}
\NormalTok{      )}
\NormalTok{  )}

\FunctionTok{show\_table\_LFL}\NormalTok{(dividend\_observations, }\AttributeTok{n =} \DecValTok{20}\NormalTok{)}
\end{Highlighting}
\end{Shaded}

\begin{verbatim}
# A tibble: 14,250 x 7
   firm_ID  year month month_ID stock_price   DPS adjustment_coefficient
     <int> <int> <int>    <int>       <dbl> <dbl>                  <dbl>
 1       1  2015     6        6        1366    29                      1
 2       1  2015    12       12        1477    29                      1
 3       1  2016     6       18        1696    40                      1
 4       1  2016    12       24        2842    40                      1
 5       1  2017     6       30        4317    43                      1
 6       1  2017    12       36        3915    43                      1
 7       1  2018     6       42        2986    43                      1
 8       1  2018    12       48        2992    43                      1
 9       1  2019     6       54        3608    53                      1
10       1  2019    12       60        4803    53                      1
11       1  2020     6       66        4698    56                      1
12       1  2020    12       72        3341    56                      1
13       2  2015     6        6        2167     3                      1
14       2  2015    12       12        2613     3                      1
15       2  2016     6       18        2395     7                      1
16       2  2016    12       24        3576     7                      1
17       2  2017     6       30        4966     7                      1
18       2  2017    12       36        5852     7                      1
19       2  2018     6       42        4415    12                      1
20       2  2018    12       48        4542    12                      1
# i 14,230 more rows
\end{verbatim}

\begin{Shaded}
\begin{Highlighting}[]
\FunctionTok{show\_table\_LFL}\NormalTok{(adjustment\_observations, }\AttributeTok{n =} \DecValTok{20}\NormalTok{)}
\end{Highlighting}
\end{Shaded}

\begin{verbatim}
# A tibble: 267 x 7
   firm_ID  year month month_ID stock_price   DPS adjustment_coefficient
     <int> <int> <int>    <int>       <dbl> <dbl>                  <dbl>
 1       1  2017     4       28        4082     0                    1.2
 2       9  2015     4        4        2655     0                    1.2
 3      10  2015     2        2         757     0                    1.2
 4      21  2019    12       60        3008    19                    1.1
 5      27  2019     9       57       13750     0                    0.1
 6      42  2016     5       17       47324     0                    0.1
 7      48  2019     2       50       13232     0                    0.1
 8      53  2019     4       52        1754     0                    1.1
 9      54  2019     2       50        4015     0                    1.1
10      60  2019    12       60        1749     0                    1.1
11      66  2019    11       59       10140     0                    0.1
12      67  2018     1       37        2020     0                    1.1
13      72  2016     5       17        1876     0                    1.1
14      74  2017     6       30        1402    11                    2  
15      80  2018     3       39        1393     0                    1.2
16      84  2017     4       28       10697     0                    0.1
17      90  2020     2       62      153396     0                    0.1
18      96  2015     4        4         805     0                    1.2
19     103  2020     6       66       44488   468                    0.1
20     104  2016     1       13        9191     0                    0.1
# i 247 more rows
\end{verbatim}

\begin{Shaded}
\begin{Highlighting}[]
\FunctionTok{show\_table\_LFL}\NormalTok{(corporate\_action\_summary)}
\end{Highlighting}
\end{Shaded}

\begin{verbatim}
# A tibble: 1 x 3
  observations dividend_observations adjusted_observations
         <int>                 <int>                 <int>
1        95040                 14250                   267
\end{verbatim}

\subsection{3.3.2 時価総額}\label{ux6642ux4fa1ux7dcfux984d}

企業\(i\)の月\(t\)における時価総額\(ME_{i,t}\)は、株価\(P_{i,t}\)と発行済株式数\(N_{i,t}\)の積である。

\[
ME_{i,t}=P_{i,t}N_{i,t}
\]

\begin{Shaded}
\begin{Highlighting}[]
\NormalTok{stock\_data }\OtherTok{\textless{}{-}}\NormalTok{ stock\_data }\SpecialCharTok{|\textgreater{}}
  \FunctionTok{mutate}\NormalTok{(}
    \AttributeTok{ME =}\NormalTok{ stock\_price }\SpecialCharTok{*}\NormalTok{ shares\_outstanding,}
    \AttributeTok{ME =} \FunctionTok{if\_else}\NormalTok{(}
      \FunctionTok{is.finite}\NormalTok{(ME) }\SpecialCharTok{\&}\NormalTok{ ME }\SpecialCharTok{\textgreater{}} \DecValTok{0}\NormalTok{,}
\NormalTok{      ME,}
      \ConstantTok{NA\_real\_}
\NormalTok{    )}
\NormalTok{  )}

\NormalTok{stock\_data }\SpecialCharTok{|\textgreater{}}
  \FunctionTok{select}\NormalTok{(}
\NormalTok{    firm\_ID,}
\NormalTok{    month\_date,}
\NormalTok{    stock\_price,}
\NormalTok{    shares\_outstanding,}
\NormalTok{    ME}
\NormalTok{  ) }\SpecialCharTok{|\textgreater{}}
  \FunctionTok{slice\_head}\NormalTok{(}\AttributeTok{n =} \DecValTok{10}\NormalTok{) }\SpecialCharTok{|\textgreater{}}
  \FunctionTok{show\_table\_LFL}\NormalTok{(}\AttributeTok{n =} \DecValTok{10}\NormalTok{)}
\end{Highlighting}
\end{Shaded}

\begin{verbatim}
# A tibble: 10 x 5
   firm_ID month_date stock_price shares_outstanding         ME
     <int> <date>           <dbl>              <dbl>      <dbl>
 1       1 2015-01-01         954            2422000 2310588000
 2       1 2015-02-01         960            2422000 2325120000
 3       1 2015-03-01        1113            2422000 2695686000
 4       1 2015-04-01        1081            2422000 2618182000
 5       1 2015-05-01        1317            2422000 3189774000
 6       1 2015-06-01        1366            2422000 3308452000
 7       1 2015-07-01        1353            2422000 3276966000
 8       1 2015-08-01        1209            2422000 2928198000
 9       1 2015-09-01        1291            2422000 3126802000
10       1 2015-10-01        1407            2422000 3407754000
\end{verbatim}

\begin{Shaded}
\begin{Highlighting}[]
\NormalTok{p\_market\_equity }\OtherTok{\textless{}{-}}\NormalTok{ stock\_data }\SpecialCharTok{|\textgreater{}}
  \FunctionTok{filter}\NormalTok{(}
    \FunctionTok{is.finite}\NormalTok{(ME),}
\NormalTok{    ME }\SpecialCharTok{\textgreater{}} \DecValTok{0}
\NormalTok{  ) }\SpecialCharTok{|\textgreater{}}
  \FunctionTok{ggplot}\NormalTok{(}\FunctionTok{aes}\NormalTok{(}\AttributeTok{x =}\NormalTok{ ME, }\AttributeTok{fill =} \StringTok{"Market Equity"}\NormalTok{)) }\SpecialCharTok{+}
  \FunctionTok{geom\_histogram}\NormalTok{(}
    \AttributeTok{bins =} \DecValTok{50}\NormalTok{,}
    \AttributeTok{colour =} \StringTok{"white"}\NormalTok{,}
    \AttributeTok{show.legend =} \ConstantTok{FALSE}
\NormalTok{  ) }\SpecialCharTok{+}
  \FunctionTok{scale\_x\_log10}\NormalTok{(}
    \AttributeTok{labels =}\NormalTok{ scales}\SpecialCharTok{::}\FunctionTok{label\_comma}\NormalTok{()}
\NormalTok{  ) }\SpecialCharTok{+}
  \FunctionTok{labs}\NormalTok{(}
    \AttributeTok{x =} \StringTok{"Market Equity (log scale)"}\NormalTok{,}
    \AttributeTok{y =} \StringTok{"Count"}\NormalTok{,}
    \AttributeTok{title =} \StringTok{"Distribution of Market Equity"}
\NormalTok{  )}

\FunctionTok{add\_lagrange\_theme\_LFL}\NormalTok{(p\_market\_equity)}
\end{Highlighting}
\end{Shaded}

\pandocbounded{\includegraphics[keepaspectratio]{part1_empirical/ch03_stock_data_time_series_files/figure-pdf/plot-market-equity-1.pdf}}

\subsection{3.3.3
月次株式リターン}\label{ux6708ux6b21ux682aux5f0fux30eaux30bfux30fcux30f3}

企業ごとに前月株価を作成し、月IDが連続している場合だけ利用する。

\begin{Shaded}
\begin{Highlighting}[]
\NormalTok{stock\_data }\OtherTok{\textless{}{-}}\NormalTok{ stock\_data }\SpecialCharTok{|\textgreater{}}
  \FunctionTok{group\_by}\NormalTok{(firm\_ID) }\SpecialCharTok{|\textgreater{}}
  \FunctionTok{arrange}\NormalTok{(month\_ID, }\AttributeTok{.by\_group =} \ConstantTok{TRUE}\NormalTok{) }\SpecialCharTok{|\textgreater{}}
  \FunctionTok{mutate}\NormalTok{(}
    \AttributeTok{lagged\_stock\_price =} \FunctionTok{lag}\NormalTok{(stock\_price),}
    \AttributeTok{lagged\_month\_ID =} \FunctionTok{lag}\NormalTok{(month\_ID),}
    \AttributeTok{consecutive\_month =}
\NormalTok{      month\_ID }\SpecialCharTok{{-}}\NormalTok{ lagged\_month\_ID }\SpecialCharTok{==} \DecValTok{1}\NormalTok{L,}
    \AttributeTok{lagged\_stock\_price =} \FunctionTok{if\_else}\NormalTok{(}
\NormalTok{      consecutive\_month,}
\NormalTok{      lagged\_stock\_price,}
      \ConstantTok{NA\_real\_}
\NormalTok{    )}
\NormalTok{  ) }\SpecialCharTok{|\textgreater{}}
  \FunctionTok{ungroup}\NormalTok{()}
\end{Highlighting}
\end{Shaded}

月次株式リターン\(R_{i,t}\)を、株価、配当および調整係数から計算する。

\[
R_{i,t}
=
\frac{
  (P_{i,t}+D_{i,t})A_{i,t}-P_{i,t-1}
}{
  P_{i,t-1}
}
\]

\begin{Shaded}
\begin{Highlighting}[]
\NormalTok{stock\_data }\OtherTok{\textless{}{-}}\NormalTok{ stock\_data }\SpecialCharTok{|\textgreater{}}
  \FunctionTok{mutate}\NormalTok{(}
    \AttributeTok{R =} \FunctionTok{case\_when}\NormalTok{(}
      \FunctionTok{is.na}\NormalTok{(lagged\_stock\_price) }\SpecialCharTok{\textasciitilde{}} \ConstantTok{NA\_real\_}\NormalTok{,}
\NormalTok{      lagged\_stock\_price }\SpecialCharTok{\textless{}=} \DecValTok{0} \SpecialCharTok{\textasciitilde{}} \ConstantTok{NA\_real\_}\NormalTok{,}
      \FunctionTok{is.na}\NormalTok{(stock\_price) }\SpecialCharTok{\textasciitilde{}} \ConstantTok{NA\_real\_}\NormalTok{,}
      \FunctionTok{is.na}\NormalTok{(DPS) }\SpecialCharTok{\textasciitilde{}} \ConstantTok{NA\_real\_}\NormalTok{,}
      \FunctionTok{is.na}\NormalTok{(adjustment\_coefficient) }\SpecialCharTok{\textasciitilde{}} \ConstantTok{NA\_real\_}\NormalTok{,}
      \ConstantTok{TRUE} \SpecialCharTok{\textasciitilde{}}\NormalTok{ (}
\NormalTok{        (stock\_price }\SpecialCharTok{+}\NormalTok{ DPS) }\SpecialCharTok{*}
\NormalTok{          adjustment\_coefficient }\SpecialCharTok{{-}}
\NormalTok{          lagged\_stock\_price}
\NormalTok{      ) }\SpecialCharTok{/}\NormalTok{ lagged\_stock\_price}
\NormalTok{    ),}
    \AttributeTok{R =} \FunctionTok{if\_else}\NormalTok{(}
      \FunctionTok{is.finite}\NormalTok{(R),}
\NormalTok{      R,}
      \ConstantTok{NA\_real\_}
\NormalTok{    )}
\NormalTok{  )}

\NormalTok{stock\_data }\SpecialCharTok{|\textgreater{}}
  \FunctionTok{select}\NormalTok{(}
\NormalTok{    firm\_ID,}
\NormalTok{    month\_date,}
\NormalTok{    lagged\_stock\_price,}
\NormalTok{    stock\_price,}
\NormalTok{    DPS,}
\NormalTok{    adjustment\_coefficient,}
\NormalTok{    R}
\NormalTok{  ) }\SpecialCharTok{|\textgreater{}}
  \FunctionTok{filter}\NormalTok{(}\SpecialCharTok{!}\FunctionTok{is.na}\NormalTok{(R)) }\SpecialCharTok{|\textgreater{}}
  \FunctionTok{slice\_head}\NormalTok{(}\AttributeTok{n =} \DecValTok{15}\NormalTok{) }\SpecialCharTok{|\textgreater{}}
  \FunctionTok{show\_table\_LFL}\NormalTok{(}\AttributeTok{n =} \DecValTok{15}\NormalTok{)}
\end{Highlighting}
\end{Shaded}

\begin{verbatim}
# A tibble: 15 x 7
   firm_ID month_date lagged_stock_price stock_price   DPS
     <int> <date>                  <dbl>       <dbl> <dbl>
 1       1 2015-02-01                954         960     0
 2       1 2015-03-01                960        1113     0
 3       1 2015-04-01               1113        1081     0
 4       1 2015-05-01               1081        1317     0
 5       1 2015-06-01               1317        1366    29
 6       1 2015-07-01               1366        1353     0
 7       1 2015-08-01               1353        1209     0
 8       1 2015-09-01               1209        1291     0
 9       1 2015-10-01               1291        1407     0
10       1 2015-11-01               1407        1562     0
11       1 2015-12-01               1562        1477    29
12       1 2016-01-01               1477        1561     0
13       1 2016-02-01               1561        1487     0
14       1 2016-03-01               1487        1526     0
15       1 2016-04-01               1526        1542     0
   adjustment_coefficient          R
                    <dbl>      <dbl>
 1                      1  0.0062893
 2                      1  0.15937  
 3                      1 -0.028751 
 4                      1  0.21832  
 5                      1  0.059226 
 6                      1 -0.0095168
 7                      1 -0.10643  
 8                      1  0.067825 
 9                      1  0.089853 
10                      1  0.11016  
11                      1 -0.035851 
12                      1  0.056872 
13                      1 -0.047406 
14                      1  0.026227 
15                      1  0.010485 
\end{verbatim}

\subsection{3.3.4
無リスク金利と超過リターン}\label{ux7121ux30eaux30b9ux30afux91d1ux5229ux3068ux8d85ux904eux30eaux30bfux30fcux30f3}

同じ月の無リスク金利はすべての企業に共通であることを確認し、月ごとの中央値で統一する。

\begin{Shaded}
\begin{Highlighting}[]
\NormalTok{risk\_free\_rate\_by\_month }\OtherTok{\textless{}{-}}\NormalTok{ stock\_data }\SpecialCharTok{|\textgreater{}}
  \FunctionTok{group\_by}\NormalTok{(}
\NormalTok{    month\_ID,}
\NormalTok{    year,}
\NormalTok{    month,}
\NormalTok{    month\_date}
\NormalTok{  ) }\SpecialCharTok{|\textgreater{}}
  \FunctionTok{summarise}\NormalTok{(}
    \AttributeTok{observations =} \FunctionTok{n}\NormalTok{(),}
    \AttributeTok{valid\_rates =} \FunctionTok{sum}\NormalTok{(}\FunctionTok{is.finite}\NormalTok{(R\_F)),}
    \AttributeTok{distinct\_rates =}
      \FunctionTok{n\_distinct}\NormalTok{(R\_F[}\FunctionTok{is.finite}\NormalTok{(R\_F)]),}
    \AttributeTok{R\_F\_month =}\NormalTok{ \{}
\NormalTok{      values }\OtherTok{\textless{}{-}}\NormalTok{ R\_F[}\FunctionTok{is.finite}\NormalTok{(R\_F)]}

      \ControlFlowTok{if}\NormalTok{ (}\FunctionTok{length}\NormalTok{(values) }\SpecialCharTok{==} \DecValTok{0}\NormalTok{) \{}
        \ConstantTok{NA\_real\_}
\NormalTok{      \} }\ControlFlowTok{else}\NormalTok{ \{}
        \FunctionTok{median}\NormalTok{(values)}
\NormalTok{      \}}
\NormalTok{    \},}
    \AttributeTok{.groups =} \StringTok{"drop"}
\NormalTok{  ) }\SpecialCharTok{|\textgreater{}}
  \FunctionTok{arrange}\NormalTok{(month\_ID)}

\NormalTok{risk\_free\_rate\_diagnostics }\OtherTok{\textless{}{-}}
\NormalTok{  risk\_free\_rate\_by\_month }\SpecialCharTok{|\textgreater{}}
  \FunctionTok{filter}\NormalTok{(distinct\_rates }\SpecialCharTok{\textgreater{}} \DecValTok{1}\NormalTok{)}

\FunctionTok{show\_table\_LFL}\NormalTok{(}
\NormalTok{  risk\_free\_rate\_diagnostics,}
  \AttributeTok{n =} \FunctionTok{nrow}\NormalTok{(risk\_free\_rate\_diagnostics)}
\NormalTok{)}
\end{Highlighting}
\end{Shaded}

\begin{verbatim}
# A tibble: 72 x 8
   month_ID  year month month_date observations valid_rates distinct_rates
      <int> <int> <int> <date>            <int>       <int>          <int>
 1        1  2015     1 2015-01-01         1266        1266             83
 2        2  2015     2 2015-02-01         1266        1266             83
 3        3  2015     3 2015-03-01         1266        1266             83
 4        4  2015     4 2015-04-01         1266        1266             83
 5        5  2015     5 2015-05-01         1266        1266             83
 6        6  2015     6 2015-06-01         1266        1266             83
 7        7  2015     7 2015-07-01         1266        1266             83
 8        8  2015     8 2015-08-01         1266        1266             83
 9        9  2015     9 2015-09-01         1266        1266             83
10       10  2015    10 2015-10-01         1266        1266             83
11       11  2015    11 2015-11-01         1266        1266             83
12       12  2015    12 2015-12-01         1266        1266             83
13       13  2016     1 2016-01-01         1293        1293             83
14       14  2016     2 2016-02-01         1293        1293             83
15       15  2016     3 2016-03-01         1293        1293             83
16       16  2016     4 2016-04-01         1293        1293             83
17       17  2016     5 2016-05-01         1293        1293             83
18       18  2016     6 2016-06-01         1293        1293             83
19       19  2016     7 2016-07-01         1293        1293             83
20       20  2016     8 2016-08-01         1293        1293             83
21       21  2016     9 2016-09-01         1293        1293             83
22       22  2016    10 2016-10-01         1293        1293             83
23       23  2016    11 2016-11-01         1293        1293             83
24       24  2016    12 2016-12-01         1293        1293             83
25       25  2017     1 2017-01-01         1319        1319             83
26       26  2017     2 2017-02-01         1319        1319             83
27       27  2017     3 2017-03-01         1319        1319             83
28       28  2017     4 2017-04-01         1319        1319             83
29       29  2017     5 2017-05-01         1319        1319             83
30       30  2017     6 2017-06-01         1319        1319             83
31       31  2017     7 2017-07-01         1319        1319             83
32       32  2017     8 2017-08-01         1319        1319             83
33       33  2017     9 2017-09-01         1319        1319             83
34       34  2017    10 2017-10-01         1319        1319             83
35       35  2017    11 2017-11-01         1319        1319             83
36       36  2017    12 2017-12-01         1319        1319             83
37       37  2018     1 2018-01-01         1323        1323             83
38       38  2018     2 2018-02-01         1323        1323             83
39       39  2018     3 2018-03-01         1323        1323             83
40       40  2018     4 2018-04-01         1323        1323             83
41       41  2018     5 2018-05-01         1323        1323             83
42       42  2018     6 2018-06-01         1323        1323             83
43       43  2018     7 2018-07-01         1323        1323             83
44       44  2018     8 2018-08-01         1323        1323             83
45       45  2018     9 2018-09-01         1323        1323             83
46       46  2018    10 2018-10-01         1323        1323             83
47       47  2018    11 2018-11-01         1323        1323             83
48       48  2018    12 2018-12-01         1323        1323             83
49       49  2019     1 2019-01-01         1356        1356             83
50       50  2019     2 2019-02-01         1356        1356             83
51       51  2019     3 2019-03-01         1356        1356             83
52       52  2019     4 2019-04-01         1356        1356             83
53       53  2019     5 2019-05-01         1356        1356             83
54       54  2019     6 2019-06-01         1356        1356             83
55       55  2019     7 2019-07-01         1356        1356             83
56       56  2019     8 2019-08-01         1356        1356             83
57       57  2019     9 2019-09-01         1356        1356             83
58       58  2019    10 2019-10-01         1356        1356             83
59       59  2019    11 2019-11-01         1356        1356             83
60       60  2019    12 2019-12-01         1356        1356             83
61       61  2020     1 2020-01-01         1363        1363             83
62       62  2020     2 2020-02-01         1363        1363             83
63       63  2020     3 2020-03-01         1363        1363             83
64       64  2020     4 2020-04-01         1363        1363             83
65       65  2020     5 2020-05-01         1363        1363             83
66       66  2020     6 2020-06-01         1363        1363             83
67       67  2020     7 2020-07-01         1363        1363             83
68       68  2020     8 2020-08-01         1363        1363             83
69       69  2020     9 2020-09-01         1363        1363             83
70       70  2020    10 2020-10-01         1363        1363             83
71       71  2020    11 2020-11-01         1363        1363             83
72       72  2020    12 2020-12-01         1363        1363             83
     R_F_month
         <dbl>
 1 0.000082029
 2 0.000082029
 3 0.000082029
 4 0.000082029
 5 0.000082029
 6 0.000082029
 7 0.000082029
 8 0.000082029
 9 0.000082029
10 0.000082029
11 0.000082029
12 0.000082029
13 0.000082029
14 0.000075959
15 0.000075959
16 0.000082029
17 0.000082029
18 0.000082029
19 0.000082029
20 0.000082029
21 0.000082029
22 0.000075959
23 0.000075959
24 0.000075959
25 0.000075959
26 0.000075959
27 0.000075959
28 0.000082029
29 0.000075959
30 0.000075959
31 0.000075959
32 0.000075959
33 0.000075959
34 0.000075959
35 0.000075959
36 0.000075959
37 0.000075959
38 0.000075959
39 0.000082029
40 0.000082029
41 0.000082029
42 0.000075959
43 0.000075959
44 0.000082029
45 0.000082029
46 0.000082029
47 0.000082029
48 0.000082029
49 0.000075959
50 0.000078994
51 0.000082029
52 0.000082029
53 0.000082029
54 0.000075959
55 0.000082029
56 0.000082029
57 0.000082029
58 0.000075959
59 0.000075959
60 0.000075959
61 0.000082029
62 0.000082029
63 0.000082029
64 0.000075959
65 0.000075959
66 0.000075959
67 0.000075959
68 0.000075959
69 0.000075959
70 0.000082029
71 0.000082029
72 0.000082029
\end{verbatim}

\begin{Shaded}
\begin{Highlighting}[]
\NormalTok{stock\_data }\OtherTok{\textless{}{-}}\NormalTok{ stock\_data }\SpecialCharTok{|\textgreater{}}
  \FunctionTok{select}\NormalTok{(}\SpecialCharTok{{-}}\NormalTok{R\_F) }\SpecialCharTok{|\textgreater{}}
  \FunctionTok{left\_join}\NormalTok{(}
\NormalTok{    risk\_free\_rate\_by\_month }\SpecialCharTok{|\textgreater{}}
      \FunctionTok{select}\NormalTok{(}
\NormalTok{        month\_ID,}
        \AttributeTok{R\_F =}\NormalTok{ R\_F\_month}
\NormalTok{      ),}
    \AttributeTok{by =} \StringTok{"month\_ID"}
\NormalTok{  )}

\NormalTok{missing\_monthly\_risk\_free\_rate }\OtherTok{\textless{}{-}}\NormalTok{ stock\_data }\SpecialCharTok{|\textgreater{}}
  \FunctionTok{distinct}\NormalTok{(month\_ID, R\_F) }\SpecialCharTok{|\textgreater{}}
  \FunctionTok{filter}\NormalTok{(}\SpecialCharTok{!}\FunctionTok{is.finite}\NormalTok{(R\_F))}

\ControlFlowTok{if}\NormalTok{ (}\FunctionTok{nrow}\NormalTok{(missing\_monthly\_risk\_free\_rate) }\SpecialCharTok{\textgreater{}} \DecValTok{0}\NormalTok{) \{}
  \FunctionTok{stop}\NormalTok{(}
    \StringTok{"月次無リスク金利を作成できないmonth\_IDがあります。"}\NormalTok{,}
    \AttributeTok{call. =} \ConstantTok{FALSE}
\NormalTok{  )}
\NormalTok{\}}
\end{Highlighting}
\end{Shaded}

月次超過リターン\(R^e_{i,t}\)は、株式リターンから無リスク金利を差し引いて計算する。

\[
R^e_{i,t}=R_{i,t}-R_{F,t}
\]

\begin{Shaded}
\begin{Highlighting}[]
\NormalTok{stock\_data }\OtherTok{\textless{}{-}}\NormalTok{ stock\_data }\SpecialCharTok{|\textgreater{}}
  \FunctionTok{mutate}\NormalTok{(}
    \AttributeTok{Re =}\NormalTok{ R }\SpecialCharTok{{-}}\NormalTok{ R\_F,}
    \AttributeTok{Re =} \FunctionTok{if\_else}\NormalTok{(}
      \FunctionTok{is.finite}\NormalTok{(Re),}
\NormalTok{      Re,}
      \ConstantTok{NA\_real\_}
\NormalTok{    )}
\NormalTok{  )}
\end{Highlighting}
\end{Shaded}

\section{3.4
月次リターンの分析}\label{ux6708ux6b21ux30eaux30bfux30fcux30f3ux306eux5206ux6790}

\subsection{3.4.1 要約統計量}\label{ux8981ux7d04ux7d71ux8a08ux91cf}

\begin{Shaded}
\begin{Highlighting}[]
\NormalTok{calculate\_skewness\_LFL }\OtherTok{\textless{}{-}} \ControlFlowTok{function}\NormalTok{(x) \{}
\NormalTok{  values }\OtherTok{\textless{}{-}}\NormalTok{ x[}\FunctionTok{is.finite}\NormalTok{(x)]}

  \ControlFlowTok{if}\NormalTok{ (}
    \FunctionTok{length}\NormalTok{(values) }\SpecialCharTok{\textless{}} \DecValTok{3} \SpecialCharTok{||}
      \FunctionTok{near}\NormalTok{(}\FunctionTok{sd}\NormalTok{(values), }\DecValTok{0}\NormalTok{)}
\NormalTok{  ) \{}
    \FunctionTok{return}\NormalTok{(}\ConstantTok{NA\_real\_}\NormalTok{)}
\NormalTok{  \}}

  \FunctionTok{mean}\NormalTok{(}
\NormalTok{    ((values }\SpecialCharTok{{-}} \FunctionTok{mean}\NormalTok{(values)) }\SpecialCharTok{/}
      \FunctionTok{sd}\NormalTok{(values))}\SpecialCharTok{\^{}}\DecValTok{3}
\NormalTok{  )}
\NormalTok{\}}

\NormalTok{calculate\_kurtosis\_LFL }\OtherTok{\textless{}{-}} \ControlFlowTok{function}\NormalTok{(x) \{}
\NormalTok{  values }\OtherTok{\textless{}{-}}\NormalTok{ x[}\FunctionTok{is.finite}\NormalTok{(x)]}

  \ControlFlowTok{if}\NormalTok{ (}
    \FunctionTok{length}\NormalTok{(values) }\SpecialCharTok{\textless{}} \DecValTok{4} \SpecialCharTok{||}
      \FunctionTok{near}\NormalTok{(}\FunctionTok{sd}\NormalTok{(values), }\DecValTok{0}\NormalTok{)}
\NormalTok{  ) \{}
    \FunctionTok{return}\NormalTok{(}\ConstantTok{NA\_real\_}\NormalTok{)}
\NormalTok{  \}}

  \FunctionTok{mean}\NormalTok{(}
\NormalTok{    ((values }\SpecialCharTok{{-}} \FunctionTok{mean}\NormalTok{(values)) }\SpecialCharTok{/}
      \FunctionTok{sd}\NormalTok{(values))}\SpecialCharTok{\^{}}\DecValTok{4}
\NormalTok{  )}
\NormalTok{\}}
\end{Highlighting}
\end{Shaded}

\begin{Shaded}
\begin{Highlighting}[]
\NormalTok{monthly\_return\_summary }\OtherTok{\textless{}{-}}\NormalTok{ stock\_data }\SpecialCharTok{|\textgreater{}}
  \FunctionTok{select}\NormalTok{(R, Re) }\SpecialCharTok{|\textgreater{}}
  \FunctionTok{pivot\_longer}\NormalTok{(}
    \AttributeTok{cols =} \FunctionTok{everything}\NormalTok{(),}
    \AttributeTok{names\_to =} \StringTok{"return\_type"}\NormalTok{,}
    \AttributeTok{values\_to =} \StringTok{"return"}
\NormalTok{  ) }\SpecialCharTok{|\textgreater{}}
  \FunctionTok{filter}\NormalTok{(}\FunctionTok{is.finite}\NormalTok{(return)) }\SpecialCharTok{|\textgreater{}}
  \FunctionTok{group\_by}\NormalTok{(return\_type) }\SpecialCharTok{|\textgreater{}}
  \FunctionTok{summarise}\NormalTok{(}
    \AttributeTok{observations =} \FunctionTok{n}\NormalTok{(),}
    \AttributeTok{mean =} \FunctionTok{mean}\NormalTok{(return),}
    \AttributeTok{median =} \FunctionTok{median}\NormalTok{(return),}
    \AttributeTok{standard\_deviation =} \FunctionTok{sd}\NormalTok{(return),}
    \AttributeTok{minimum =} \FunctionTok{min}\NormalTok{(return),}
    \AttributeTok{first\_quartile =} \FunctionTok{quantile}\NormalTok{(return, }\FloatTok{0.25}\NormalTok{),}
    \AttributeTok{third\_quartile =} \FunctionTok{quantile}\NormalTok{(return, }\FloatTok{0.75}\NormalTok{),}
    \AttributeTok{maximum =} \FunctionTok{max}\NormalTok{(return),}
    \AttributeTok{skewness =} \FunctionTok{calculate\_skewness\_LFL}\NormalTok{(return),}
    \AttributeTok{kurtosis =} \FunctionTok{calculate\_kurtosis\_LFL}\NormalTok{(return),}
    \AttributeTok{.groups =} \StringTok{"drop"}
\NormalTok{  )}

\FunctionTok{show\_table\_LFL}\NormalTok{(monthly\_return\_summary)}
\end{Highlighting}
\end{Shaded}

\begin{verbatim}
# A tibble: 2 x 11
  return_type observations     mean   median standard_deviation  minimum
  <chr>              <int>    <dbl>    <dbl>              <dbl>    <dbl>
1 R                  93525 0.015859 0.010332           0.091131 -0.38028
2 Re                 93525 0.015779 0.010252           0.091131 -0.38036
  first_quartile third_quartile maximum skewness kurtosis
           <dbl>          <dbl>   <dbl>    <dbl>    <dbl>
1      -0.040573       0.064808 0.51498  0.50657   4.2680
2      -0.040652       0.064730 0.51490  0.50656   4.2680
\end{verbatim}

\subsection{3.4.2
リターン分布}\label{ux30eaux30bfux30fcux30f3ux5206ux5e03}

極端値の影響を抑えるため、1パーセンタイルから99パーセンタイルまでを表示する。

\begin{Shaded}
\begin{Highlighting}[]
\NormalTok{monthly\_return\_limits }\OtherTok{\textless{}{-}} \FunctionTok{quantile}\NormalTok{(}
\NormalTok{  stock\_data}\SpecialCharTok{$}\NormalTok{R,}
  \AttributeTok{probs =} \FunctionTok{c}\NormalTok{(}\FloatTok{0.01}\NormalTok{, }\FloatTok{0.99}\NormalTok{),}
  \AttributeTok{na.rm =} \ConstantTok{TRUE}
\NormalTok{)}

\NormalTok{p\_monthly\_return\_distribution }\OtherTok{\textless{}{-}}\NormalTok{ stock\_data }\SpecialCharTok{|\textgreater{}}
  \FunctionTok{filter}\NormalTok{(}
    \FunctionTok{is.finite}\NormalTok{(R),}
    \FunctionTok{between}\NormalTok{(}
\NormalTok{      R,}
\NormalTok{      monthly\_return\_limits[[}\DecValTok{1}\NormalTok{]],}
\NormalTok{      monthly\_return\_limits[[}\DecValTok{2}\NormalTok{]]}
\NormalTok{    )}
\NormalTok{  ) }\SpecialCharTok{|\textgreater{}}
  \FunctionTok{ggplot}\NormalTok{(}\FunctionTok{aes}\NormalTok{(}\AttributeTok{x =}\NormalTok{ R, }\AttributeTok{fill =} \StringTok{"Monthly Return"}\NormalTok{)) }\SpecialCharTok{+}
  \FunctionTok{geom\_histogram}\NormalTok{(}
    \AttributeTok{bins =} \DecValTok{60}\NormalTok{,}
    \AttributeTok{colour =} \StringTok{"white"}\NormalTok{,}
    \AttributeTok{show.legend =} \ConstantTok{FALSE}
\NormalTok{  ) }\SpecialCharTok{+}
  \FunctionTok{scale\_x\_continuous}\NormalTok{(}
    \AttributeTok{labels =}\NormalTok{ scales}\SpecialCharTok{::}\FunctionTok{label\_percent}\NormalTok{()}
\NormalTok{  ) }\SpecialCharTok{+}
  \FunctionTok{labs}\NormalTok{(}
    \AttributeTok{x =} \StringTok{"Monthly Return"}\NormalTok{,}
    \AttributeTok{y =} \StringTok{"Count"}\NormalTok{,}
    \AttributeTok{title =} \StringTok{"Distribution of Monthly Stock Returns"}
\NormalTok{  )}

\FunctionTok{add\_lagrange\_theme\_LFL}\NormalTok{(}
\NormalTok{  p\_monthly\_return\_distribution}
\NormalTok{)}
\end{Highlighting}
\end{Shaded}

\pandocbounded{\includegraphics[keepaspectratio]{part1_empirical/ch03_stock_data_time_series_files/figure-pdf/plot-monthly-return-distribution-1.pdf}}

\subsection{3.4.3
企業別の累積リターン}\label{ux4f01ux696dux5225ux306eux7d2fux7a4dux30eaux30bfux30fcux30f3}

観測数が多い企業から6社を選び、月次リターンを複利で累積する。

\begin{Shaded}
\begin{Highlighting}[]
\NormalTok{sample\_firm\_ids }\OtherTok{\textless{}{-}}\NormalTok{ stock\_data }\SpecialCharTok{|\textgreater{}}
  \FunctionTok{filter}\NormalTok{(}\FunctionTok{is.finite}\NormalTok{(R)) }\SpecialCharTok{|\textgreater{}}
  \FunctionTok{count}\NormalTok{(}
\NormalTok{    firm\_ID,}
    \AttributeTok{sort =} \ConstantTok{TRUE}
\NormalTok{  ) }\SpecialCharTok{|\textgreater{}}
  \FunctionTok{slice\_head}\NormalTok{(}\AttributeTok{n =} \DecValTok{6}\NormalTok{) }\SpecialCharTok{|\textgreater{}}
  \FunctionTok{pull}\NormalTok{(firm\_ID)}

\NormalTok{cumulative\_return\_data }\OtherTok{\textless{}{-}}\NormalTok{ stock\_data }\SpecialCharTok{|\textgreater{}}
  \FunctionTok{filter}\NormalTok{(}
\NormalTok{    firm\_ID }\SpecialCharTok{\%in\%}\NormalTok{ sample\_firm\_ids,}
    \FunctionTok{is.finite}\NormalTok{(R)}
\NormalTok{  ) }\SpecialCharTok{|\textgreater{}}
  \FunctionTok{group\_by}\NormalTok{(firm\_ID) }\SpecialCharTok{|\textgreater{}}
  \FunctionTok{arrange}\NormalTok{(month\_ID, }\AttributeTok{.by\_group =} \ConstantTok{TRUE}\NormalTok{) }\SpecialCharTok{|\textgreater{}}
  \FunctionTok{mutate}\NormalTok{(}
    \AttributeTok{cumulative\_return =} \FunctionTok{cumprod}\NormalTok{(}\DecValTok{1} \SpecialCharTok{+}\NormalTok{ R) }\SpecialCharTok{{-}} \DecValTok{1}
\NormalTok{  ) }\SpecialCharTok{|\textgreater{}}
  \FunctionTok{ungroup}\NormalTok{()}
\end{Highlighting}
\end{Shaded}

\begin{Shaded}
\begin{Highlighting}[]
\NormalTok{p\_cumulative\_return }\OtherTok{\textless{}{-}}\NormalTok{ cumulative\_return\_data }\SpecialCharTok{|\textgreater{}}
  \FunctionTok{ggplot}\NormalTok{(}
    \FunctionTok{aes}\NormalTok{(}
      \AttributeTok{x =}\NormalTok{ month\_date,}
      \AttributeTok{y =}\NormalTok{ cumulative\_return,}
      \AttributeTok{colour =} \FunctionTok{factor}\NormalTok{(firm\_ID)}
\NormalTok{    )}
\NormalTok{  ) }\SpecialCharTok{+}
  \FunctionTok{geom\_line}\NormalTok{(}\AttributeTok{linewidth =} \FloatTok{0.7}\NormalTok{) }\SpecialCharTok{+}
  \FunctionTok{scale\_y\_continuous}\NormalTok{(}
    \AttributeTok{labels =}\NormalTok{ scales}\SpecialCharTok{::}\FunctionTok{label\_percent}\NormalTok{()}
\NormalTok{  ) }\SpecialCharTok{+}
  \FunctionTok{labs}\NormalTok{(}
    \AttributeTok{x =} \StringTok{"Month"}\NormalTok{,}
    \AttributeTok{y =} \StringTok{"Cumulative Return"}\NormalTok{,}
    \AttributeTok{colour =} \StringTok{"Firm ID"}\NormalTok{,}
    \AttributeTok{title =} \StringTok{"Cumulative Returns of Selected Firms"}
\NormalTok{  )}

\FunctionTok{add\_lagrange\_theme\_LFL}\NormalTok{(p\_cumulative\_return) }\SpecialCharTok{+}
\NormalTok{  ggthemes}\SpecialCharTok{::}\FunctionTok{scale\_colour\_economist}\NormalTok{()}
\end{Highlighting}
\end{Shaded}

\pandocbounded{\includegraphics[keepaspectratio]{part1_empirical/ch03_stock_data_time_series_files/figure-pdf/plot-cumulative-return-1.pdf}}

\subsection{3.4.4
企業別の要約統計量}\label{ux4f01ux696dux5225ux306eux8981ux7d04ux7d71ux8a08ux91cf}

\begin{Shaded}
\begin{Highlighting}[]
\NormalTok{firm\_return\_summary }\OtherTok{\textless{}{-}}\NormalTok{ sample\_firm\_ids }\SpecialCharTok{|\textgreater{}}
  \FunctionTok{map\_dfr}\NormalTok{(}
\NormalTok{    \textbackslash{}(firm\_id) \{}
\NormalTok{      firm\_returns }\OtherTok{\textless{}{-}}\NormalTok{ stock\_data }\SpecialCharTok{|\textgreater{}}
        \FunctionTok{filter}\NormalTok{(}
\NormalTok{          firm\_ID }\SpecialCharTok{==}\NormalTok{ firm\_id,}
          \FunctionTok{is.finite}\NormalTok{(R)}
\NormalTok{        ) }\SpecialCharTok{|\textgreater{}}
        \FunctionTok{pull}\NormalTok{(R)}

      \FunctionTok{tibble}\NormalTok{(}
        \AttributeTok{firm\_ID =}\NormalTok{ firm\_id,}
        \AttributeTok{observations =} \FunctionTok{length}\NormalTok{(firm\_returns),}
        \AttributeTok{mean\_return =} \FunctionTok{mean}\NormalTok{(firm\_returns),}
        \AttributeTok{volatility =} \FunctionTok{sd}\NormalTok{(firm\_returns),}
        \AttributeTok{minimum\_return =} \FunctionTok{min}\NormalTok{(firm\_returns),}
        \AttributeTok{maximum\_return =} \FunctionTok{max}\NormalTok{(firm\_returns)}
\NormalTok{      )}
\NormalTok{    \}}
\NormalTok{  )}

\FunctionTok{show\_table\_LFL}\NormalTok{(firm\_return\_summary)}
\end{Highlighting}
\end{Shaded}

\begin{verbatim}
# A tibble: 6 x 6
  firm_ID observations mean_return volatility minimum_return maximum_return
    <int>        <int>       <dbl>      <dbl>          <dbl>          <dbl>
1       1           71   0.029240    0.11544       -0.18157         0.33603
2       2           71   0.018125    0.089135      -0.18825         0.24481
3       3           71   0.023028    0.075369      -0.14460         0.16440
4       4           71  -0.0030159   0.054766      -0.13816         0.11475
5       5           71   0.028390    0.061602      -0.097372        0.25720
6       7           71   0.0077876   0.053950      -0.096949        0.14549
\end{verbatim}

\section{3.5
年次データと財務データの結合}\label{ux5e74ux6b21ux30c7ux30fcux30bfux3068ux8ca1ux52d9ux30c7ux30fcux30bfux306eux7d50ux5408}

\subsection{3.5.1
月次リターンの年次集約}\label{ux6708ux6b21ux30eaux30bfux30fcux30f3ux306eux5e74ux6b21ux96c6ux7d04}

企業\(i\)の年\(t\)における年次リターンは、月次リターンを複利で集約する。

\[
R_{i,t}^{annual}
=
\prod_{m=1}^{M_{i,t}}
(1+R_{i,t,m})-1
\]

\begin{Shaded}
\begin{Highlighting}[]
\NormalTok{annual\_stock\_data }\OtherTok{\textless{}{-}}\NormalTok{ stock\_data }\SpecialCharTok{|\textgreater{}}
  \FunctionTok{group\_by}\NormalTok{(}
\NormalTok{    firm\_ID,}
\NormalTok{    year}
\NormalTok{  ) }\SpecialCharTok{|\textgreater{}}
  \FunctionTok{summarise}\NormalTok{(}
    \AttributeTok{months =} \FunctionTok{sum}\NormalTok{(}\FunctionTok{is.finite}\NormalTok{(R)),}
    \AttributeTok{R =} \ControlFlowTok{if}\NormalTok{ (}
      \FunctionTok{sum}\NormalTok{(}\FunctionTok{is.finite}\NormalTok{(R)) }\SpecialCharTok{==} \DecValTok{0}
\NormalTok{    ) \{}
      \ConstantTok{NA\_real\_}
\NormalTok{    \} }\ControlFlowTok{else}\NormalTok{ \{}
      \FunctionTok{prod}\NormalTok{(}\DecValTok{1} \SpecialCharTok{+}\NormalTok{ R[}\FunctionTok{is.finite}\NormalTok{(R)]) }\SpecialCharTok{{-}} \DecValTok{1}
\NormalTok{    \},}
    \AttributeTok{Re =} \ControlFlowTok{if}\NormalTok{ (}
      \FunctionTok{sum}\NormalTok{(}\FunctionTok{is.finite}\NormalTok{(Re)) }\SpecialCharTok{==} \DecValTok{0}
\NormalTok{    ) \{}
      \ConstantTok{NA\_real\_}
\NormalTok{    \} }\ControlFlowTok{else}\NormalTok{ \{}
      \FunctionTok{prod}\NormalTok{(}\DecValTok{1} \SpecialCharTok{+}\NormalTok{ Re[}\FunctionTok{is.finite}\NormalTok{(Re)]) }\SpecialCharTok{{-}} \DecValTok{1}
\NormalTok{    \},}
    \AttributeTok{ME =}\NormalTok{ \{}
\NormalTok{      values }\OtherTok{\textless{}{-}}\NormalTok{ ME[}\FunctionTok{is.finite}\NormalTok{(ME)]}

      \ControlFlowTok{if}\NormalTok{ (}\FunctionTok{length}\NormalTok{(values) }\SpecialCharTok{==} \DecValTok{0}\NormalTok{) \{}
        \ConstantTok{NA\_real\_}
\NormalTok{      \} }\ControlFlowTok{else}\NormalTok{ \{}
\NormalTok{        dplyr}\SpecialCharTok{::}\FunctionTok{last}\NormalTok{(values)}
\NormalTok{      \}}
\NormalTok{    \},}
    \AttributeTok{.groups =} \StringTok{"drop"}
\NormalTok{  ) }\SpecialCharTok{|\textgreater{}}
  \FunctionTok{arrange}\NormalTok{(firm\_ID, year)}

\FunctionTok{show\_table\_LFL}\NormalTok{(annual\_stock\_data, }\AttributeTok{n =} \DecValTok{20}\NormalTok{)}
\end{Highlighting}
\end{Shaded}

\begin{verbatim}
# A tibble: 7,920 x 6
   firm_ID  year months        R       Re           ME
     <int> <int>  <int>    <dbl>    <dbl>        <dbl>
 1       1  2015     11  0.61213  0.61073   3577294000
 2       1  2016     12  0.99727  0.99547   6883324000
 3       1  2017     12  0.68786  0.68637  11376990000
 4       1  2018     12 -0.21361 -0.21439   8694752000
 5       1  2019     12  0.64684  0.64533  13957518000
 6       1  2020     12 -0.28430 -0.28501   9708946000
 7       2  2015     11  1.1092   1.1074    4086732000
 8       2  2016     12  0.37523  0.37395   5592864000
 9       2  2017     12  0.64073  0.63928   9152528000
10       2  2018     12 -0.21969 -0.22046   7103688000
11       2  2019     12  0.40819  0.40688   9962680000
12       2  2020     12 -0.47547 -0.47599   5194044000
13       3  2015     11  0.27011  0.26898  70434234000
14       3  2016     12  0.37526  0.37398  96064839000
15       3  2017     12  0.38682  0.38558 132546843000
16       3  2018     12 -0.17493 -0.17573 107515395000
17       3  2019     12  0.65869  0.65717 174487833000
18       3  2020     12  0.25382  0.25265 214032195000
19       4  2015     11  0.12885  0.12784  14518674000
20       4  2016     12 -0.18552 -0.18631  11592135000
# i 7,900 more rows
\end{verbatim}

\subsection{3.5.2
年次財務データの読込み}\label{ux5e74ux6b21ux8ca1ux52d9ux30c7ux30fcux30bfux306eux8aadux8fbcux307f}

第2章で作成した\texttt{ch04\_output.csv}を読み込む。

\begin{Shaded}
\begin{Highlighting}[]
\NormalTok{financial\_data }\OtherTok{\textless{}{-}}
  \FunctionTok{read\_derived\_csv\_file\_LFL}\NormalTok{(}\StringTok{"ch04\_output.csv"}\NormalTok{) }\SpecialCharTok{|\textgreater{}}
  \FunctionTok{mutate}\NormalTok{(}
    \AttributeTok{firm\_ID =} \FunctionTok{as.integer}\NormalTok{(firm\_ID),}
    \AttributeTok{year =} \FunctionTok{as.integer}\NormalTok{(year)}
\NormalTok{  )}

\FunctionTok{show\_table\_LFL}\NormalTok{(financial\_data, }\AttributeTok{n =} \DecValTok{10}\NormalTok{)}
\end{Highlighting}
\end{Shaded}

\begin{verbatim}
# A tibble: 7,920 x 25
    year firm_ID industry_ID  sales     OX      NFE      X      OA      FA
   <int>   <int>       <dbl>  <dbl>  <dbl>    <dbl>  <dbl>   <dbl>   <dbl>
 1  2015       1           1 5261.4 437.49  NA      286.64 13006.   3543.4
 2  2016       1           1 5949.0 564.14  50.667  513.48 13866.   4642.2
 3  2017       1           1 6505.1 691.18  29.543  661.64 13953.   7744.0
 4  2018       1           1 6846.4 751.29  86.487  664.8  18818.   7284.7
 5  2019       1           1 7572.2 958.53 298.05   660.48 18190    9735.1
 6  2020       1           1 7537.6 778.37 -65.459  843.83 20463.  10274. 
 7  2015       2           1 3505.8  45.82   5.7511  40.07  2977.8  2258.3
 8  2016       2           1 3491.3  51.25   1.8765  49.37  3184.8  1881.8
 9  2017       2           1 3945.7  83.43   7.5279  75.9   3392.2  2425.2
10  2018       2           1 4139.3  93.4    6.8166  86.58  3569.3  2331.0
       OL      FO     NOA       NFO      BE lagged_BE       ROE lagged_NOA
    <dbl>   <dbl>   <dbl>     <dbl>   <dbl>     <dbl>     <dbl>      <dbl>
 1 4373.0  2480.7  8632.6 -1062.7    9695.3      NA   NA              NA  
 2 4534.2  3959.7  9331.4  -682.46  10014.     9695.3  0.052962     8632.6
 3 5111.2  6159.0  8841.4 -1585.0   10426.    10014.   0.066073     9331.4
 4 5137.3 10124.  13681.   2839.2   10842.    10426.   0.063762     8841.4
 5 5488.0 11362.  12702.   1627.1   11075.    10842.   0.060919    13681. 
 6 5371.4 13772.  15091.   3497.9   11594.    11075.   0.076193    12702. 
 7 1840.4  2340.9  1137.5    82.560  1054.9      NA   NA              NA  
 8 1769.2  2215.9  1415.6   334.04   1081.6    1054.9  0.046800     1137.5
 9 1955.3  2727.4  1436.9   302.14   1134.8    1081.6  0.070174     1415.6
10 2022.1  2694.5  1547.2   363.48   1183.7    1134.8  0.076295     1436.9
   lagged_NFO      RNOA       PM      ATO lagged_FLEV       NBC ROE_DuPont
        <dbl>     <dbl>    <dbl>    <dbl>       <dbl>     <dbl>      <dbl>
 1     NA     NA        0.083151 NA         NA        NA         NA       
 2  -1062.7    0.065350 0.094830  0.68913   -0.10961  -0.047678   0.052961
 3   -682.46   0.074071 0.10625   0.69712   -0.068152 -0.043289   0.066072
 4  -1585.0    0.084974 0.10974   0.77436   -0.15202  -0.054567   0.063762
 5   2839.2    0.070062 0.12658   0.55348    0.26187   0.10498    0.060919
 6   1627.1    0.061279 0.10326   0.59342    0.14692  -0.040231   0.076193
 7     NA     NA        0.013070 NA         NA        NA         NA       
 8     82.560  0.045056 0.014679  3.0693     0.078263  0.022729   0.046803
 9    334.04   0.058935 0.021145  2.7872     0.30884   0.022536   0.070176
10    302.14   0.064999 0.022564  2.8807     0.26625   0.022561   0.076298
   DuPont_gap
        <dbl>
 1 NA        
 2  7.7332e-7
 3  3.1527e-7
 4 -3.3567e-7
 5 -2.0862e-8
 6  1.1089e-7
 7 NA        
 8 -3.2789e-6
 9 -1.9587e-6
10 -2.9722e-6
# i 7,910 more rows
\end{verbatim}

財務データの企業・年度キーが一意であることを確認する。

\begin{Shaded}
\begin{Highlighting}[]
\NormalTok{financial\_key\_duplicates }\OtherTok{\textless{}{-}}\NormalTok{ financial\_data }\SpecialCharTok{|\textgreater{}}
  \FunctionTok{count}\NormalTok{(}
\NormalTok{    firm\_ID,}
\NormalTok{    year,}
    \AttributeTok{name =} \StringTok{"observations"}
\NormalTok{  ) }\SpecialCharTok{|\textgreater{}}
  \FunctionTok{filter}\NormalTok{(observations }\SpecialCharTok{\textgreater{}} \DecValTok{1}\NormalTok{)}

\FunctionTok{show\_table\_LFL}\NormalTok{(financial\_key\_duplicates)}
\end{Highlighting}
\end{Shaded}

\begin{verbatim}
# A tibble: 0 x 3
# i 3 variables: firm_ID <int>, year <int>, observations <int>
\end{verbatim}

\begin{Shaded}
\begin{Highlighting}[]
\ControlFlowTok{if}\NormalTok{ (}\FunctionTok{nrow}\NormalTok{(financial\_key\_duplicates) }\SpecialCharTok{\textgreater{}} \DecValTok{0}\NormalTok{) \{}
  \FunctionTok{stop}\NormalTok{(}
    \StringTok{"財務データのfirm\_IDとyearに重複があります。"}\NormalTok{,}
    \AttributeTok{call. =} \ConstantTok{FALSE}
\NormalTok{  )}
\NormalTok{\}}
\end{Highlighting}
\end{Shaded}

\subsection{3.5.3
年次データの結合}\label{ux5e74ux6b21ux30c7ux30fcux30bfux306eux7d50ux5408}

\begin{Shaded}
\begin{Highlighting}[]
\NormalTok{annual\_data }\OtherTok{\textless{}{-}}\NormalTok{ annual\_stock\_data }\SpecialCharTok{|\textgreater{}}
  \FunctionTok{left\_join}\NormalTok{(}
\NormalTok{    financial\_data,}
    \AttributeTok{by =} \FunctionTok{c}\NormalTok{(}\StringTok{"firm\_ID"}\NormalTok{, }\StringTok{"year"}\NormalTok{)}
\NormalTok{  ) }\SpecialCharTok{|\textgreater{}}
  \FunctionTok{arrange}\NormalTok{(firm\_ID, year)}

\FunctionTok{stopifnot}\NormalTok{(}
  \FunctionTok{nrow}\NormalTok{(annual\_data) }\SpecialCharTok{==}
    \FunctionTok{nrow}\NormalTok{(annual\_stock\_data)}
\NormalTok{)}

\NormalTok{annual\_data\_overview }\OtherTok{\textless{}{-}}\NormalTok{ annual\_data }\SpecialCharTok{|\textgreater{}}
  \FunctionTok{summarise}\NormalTok{(}
    \AttributeTok{observations =} \FunctionTok{n}\NormalTok{(),}
    \AttributeTok{firms =} \FunctionTok{n\_distinct}\NormalTok{(firm\_ID),}
    \AttributeTok{first\_year =} \FunctionTok{min}\NormalTok{(year, }\AttributeTok{na.rm =} \ConstantTok{TRUE}\NormalTok{),}
    \AttributeTok{last\_year =} \FunctionTok{max}\NormalTok{(year, }\AttributeTok{na.rm =} \ConstantTok{TRUE}\NormalTok{),}
    \AttributeTok{financial\_matches =}
      \FunctionTok{sum}\NormalTok{(}\SpecialCharTok{!}\FunctionTok{is.na}\NormalTok{(BE))}
\NormalTok{  )}

\FunctionTok{show\_table\_LFL}\NormalTok{(annual\_data\_overview)}
\end{Highlighting}
\end{Shaded}

\begin{verbatim}
# A tibble: 1 x 5
  observations firms first_year last_year financial_matches
         <int> <int>      <int>     <int>             <int>
1         7920  1515       2015      2020              7920
\end{verbatim}

\subsection{3.5.4
月次データへの財務情報の付加}\label{ux6708ux6b21ux30c7ux30fcux30bfux3078ux306eux8ca1ux52d9ux60c5ux5831ux306eux4ed8ux52a0}

\begin{Shaded}
\begin{Highlighting}[]
\NormalTok{monthly\_data }\OtherTok{\textless{}{-}}\NormalTok{ stock\_data }\SpecialCharTok{|\textgreater{}}
  \FunctionTok{left\_join}\NormalTok{(}
\NormalTok{    financial\_data,}
    \AttributeTok{by =} \FunctionTok{c}\NormalTok{(}\StringTok{"firm\_ID"}\NormalTok{, }\StringTok{"year"}\NormalTok{)}
\NormalTok{  ) }\SpecialCharTok{|\textgreater{}}
  \FunctionTok{arrange}\NormalTok{(firm\_ID, month\_ID)}

\FunctionTok{stopifnot}\NormalTok{(}
  \FunctionTok{nrow}\NormalTok{(monthly\_data) }\SpecialCharTok{==}
    \FunctionTok{nrow}\NormalTok{(stock\_data)}
\NormalTok{)}

\NormalTok{row\_count\_check }\OtherTok{\textless{}{-}} \FunctionTok{tibble}\NormalTok{(}
  \AttributeTok{stock\_rows =} \FunctionTok{nrow}\NormalTok{(stock\_data),}
  \AttributeTok{monthly\_rows =} \FunctionTok{nrow}\NormalTok{(monthly\_data),}
  \AttributeTok{annual\_rows =} \FunctionTok{nrow}\NormalTok{(annual\_data)}
\NormalTok{)}

\FunctionTok{show\_table\_LFL}\NormalTok{(row\_count\_check)}
\end{Highlighting}
\end{Shaded}

\begin{verbatim}
# A tibble: 1 x 3
  stock_rows monthly_rows annual_rows
       <int>        <int>       <int>
1      95040        95040        7920
\end{verbatim}

\subsection{3.5.5
売上高・利益・時価総額}\label{ux58f2ux4e0aux9ad8ux5229ux76caux6642ux4fa1ux7dcfux984d}

最新年度について、売上高、当期純利益および時価総額の関係を確認する。

\begin{Shaded}
\begin{Highlighting}[]
\NormalTok{latest\_analysis\_year }\OtherTok{\textless{}{-}}\NormalTok{ annual\_data }\SpecialCharTok{|\textgreater{}}
  \FunctionTok{filter}\NormalTok{(}
    \FunctionTok{is.finite}\NormalTok{(sales),}
\NormalTok{    sales }\SpecialCharTok{\textgreater{}} \DecValTok{0}\NormalTok{,}
    \FunctionTok{is.finite}\NormalTok{(X),}
\NormalTok{    X }\SpecialCharTok{\textgreater{}} \DecValTok{0}\NormalTok{,}
    \FunctionTok{is.finite}\NormalTok{(ME),}
\NormalTok{    ME }\SpecialCharTok{\textgreater{}} \DecValTok{0}
\NormalTok{  ) }\SpecialCharTok{|\textgreater{}}
  \FunctionTok{summarise}\NormalTok{(}
    \AttributeTok{latest\_year =} \FunctionTok{max}\NormalTok{(year)}
\NormalTok{  ) }\SpecialCharTok{|\textgreater{}}
  \FunctionTok{pull}\NormalTok{(latest\_year)}

\NormalTok{bubble\_chart\_data }\OtherTok{\textless{}{-}}\NormalTok{ annual\_data }\SpecialCharTok{|\textgreater{}}
  \FunctionTok{filter}\NormalTok{(}
\NormalTok{    year }\SpecialCharTok{==}\NormalTok{ latest\_analysis\_year,}
    \FunctionTok{is.finite}\NormalTok{(sales),}
\NormalTok{    sales }\SpecialCharTok{\textgreater{}} \DecValTok{0}\NormalTok{,}
    \FunctionTok{is.finite}\NormalTok{(X),}
\NormalTok{    X }\SpecialCharTok{\textgreater{}} \DecValTok{0}\NormalTok{,}
    \FunctionTok{is.finite}\NormalTok{(ME),}
\NormalTok{    ME }\SpecialCharTok{\textgreater{}} \DecValTok{0}
\NormalTok{  )}
\end{Highlighting}
\end{Shaded}

\begin{Shaded}
\begin{Highlighting}[]
\NormalTok{p\_sales\_earnings\_market\_equity }\OtherTok{\textless{}{-}}\NormalTok{ bubble\_chart\_data }\SpecialCharTok{|\textgreater{}}
  \FunctionTok{ggplot}\NormalTok{(}
    \FunctionTok{aes}\NormalTok{(}
      \AttributeTok{x =} \FunctionTok{log}\NormalTok{(sales),}
      \AttributeTok{y =} \FunctionTok{log}\NormalTok{(X),}
      \AttributeTok{size =}\NormalTok{ ME}
\NormalTok{    )}
\NormalTok{  ) }\SpecialCharTok{+}
  \FunctionTok{geom\_point}\NormalTok{(}\AttributeTok{alpha =} \FloatTok{0.35}\NormalTok{) }\SpecialCharTok{+}
  \FunctionTok{scale\_size}\NormalTok{(}
    \AttributeTok{range =} \FunctionTok{c}\NormalTok{(}\DecValTok{1}\NormalTok{, }\DecValTok{12}\NormalTok{),}
    \AttributeTok{name =} \StringTok{"Market Equity"}
\NormalTok{  ) }\SpecialCharTok{+}
  \FunctionTok{labs}\NormalTok{(}
    \AttributeTok{x =} \StringTok{"Log Sales"}\NormalTok{,}
    \AttributeTok{y =} \StringTok{"Log Net Income"}\NormalTok{,}
    \AttributeTok{title =} \FunctionTok{paste}\NormalTok{(}
      \StringTok{"Sales, Earnings and Market Equity in"}\NormalTok{,}
\NormalTok{      latest\_analysis\_year}
\NormalTok{    )}
\NormalTok{  )}

\FunctionTok{add\_lagrange\_theme\_LFL}\NormalTok{(}
\NormalTok{  p\_sales\_earnings\_market\_equity}
\NormalTok{)}
\end{Highlighting}
\end{Shaded}

\pandocbounded{\includegraphics[keepaspectratio]{part1_empirical/ch03_stock_data_time_series_files/figure-pdf/plot-sales-earnings-market-equity-1.pdf}}

\section{3.6 統計分析}\label{ux7d71ux8a08ux5206ux6790}

\subsection{3.6.1
月次超過リターンの一標本t検定}\label{ux6708ux6b21ux8d85ux904eux30eaux30bfux30fcux30f3ux306eux4e00ux6a19ux672ctux691cux5b9a}

観測数が最も多い企業を選び、月次超過リターンの平均が0と異なるかを検定する。

\begin{Shaded}
\begin{Highlighting}[]
\NormalTok{test\_firm\_id }\OtherTok{\textless{}{-}}\NormalTok{ stock\_data }\SpecialCharTok{|\textgreater{}}
  \FunctionTok{filter}\NormalTok{(}\FunctionTok{is.finite}\NormalTok{(Re)) }\SpecialCharTok{|\textgreater{}}
  \FunctionTok{count}\NormalTok{(}
\NormalTok{    firm\_ID,}
    \AttributeTok{sort =} \ConstantTok{TRUE}
\NormalTok{  ) }\SpecialCharTok{|\textgreater{}}
  \FunctionTok{slice\_head}\NormalTok{(}\AttributeTok{n =} \DecValTok{1}\NormalTok{) }\SpecialCharTok{|\textgreater{}}
  \FunctionTok{pull}\NormalTok{(firm\_ID)}

\NormalTok{test\_firm\_returns }\OtherTok{\textless{}{-}}\NormalTok{ stock\_data }\SpecialCharTok{|\textgreater{}}
  \FunctionTok{filter}\NormalTok{(}
\NormalTok{    firm\_ID }\SpecialCharTok{==}\NormalTok{ test\_firm\_id,}
    \FunctionTok{is.finite}\NormalTok{(Re)}
\NormalTok{  ) }\SpecialCharTok{|\textgreater{}}
  \FunctionTok{pull}\NormalTok{(Re)}

\NormalTok{manual\_t\_value }\OtherTok{\textless{}{-}} \FunctionTok{mean}\NormalTok{(test\_firm\_returns) }\SpecialCharTok{/}
\NormalTok{  (}
    \FunctionTok{sd}\NormalTok{(test\_firm\_returns) }\SpecialCharTok{/}
      \FunctionTok{sqrt}\NormalTok{(}\FunctionTok{length}\NormalTok{(test\_firm\_returns))}
\NormalTok{  )}

\NormalTok{t\_test\_result }\OtherTok{\textless{}{-}} \FunctionTok{t.test}\NormalTok{(}
\NormalTok{  test\_firm\_returns,}
  \AttributeTok{mu =} \DecValTok{0}
\NormalTok{)}

\NormalTok{t\_test\_summary }\OtherTok{\textless{}{-}}\NormalTok{ broom}\SpecialCharTok{::}\FunctionTok{tidy}\NormalTok{(t\_test\_result) }\SpecialCharTok{|\textgreater{}}
  \FunctionTok{mutate}\NormalTok{(}
    \AttributeTok{firm\_ID =}\NormalTok{ test\_firm\_id,}
    \AttributeTok{manual\_t\_value =}\NormalTok{ manual\_t\_value,}
    \AttributeTok{.before =} \DecValTok{1}
\NormalTok{  )}

\FunctionTok{show\_table\_LFL}\NormalTok{(t\_test\_summary)}
\end{Highlighting}
\end{Shaded}

\begin{verbatim}
# A tibble: 1 x 10
  firm_ID manual_t_value estimate statistic  p.value parameter  conf.low
    <int>          <dbl>    <dbl>     <dbl>    <dbl>     <dbl>     <dbl>
1       1         2.1285 0.029161    2.1285 0.036817        70 0.0018367
  conf.high method            alternative
      <dbl> <chr>             <chr>      
1  0.056485 One Sample t-test two.sided  
\end{verbatim}

\subsection{3.6.2
簿価時価比率と年次超過リターン}\label{ux7c3fux4fa1ux6642ux4fa1ux6bd4ux7387ux3068ux5e74ux6b21ux8d85ux904eux30eaux30bfux30fcux30f3}

前年度末の時価総額を作成し、期首株主資本を前年度末時価総額で割って簿価時価比率を計算する。

\begin{Shaded}
\begin{Highlighting}[]
\NormalTok{annual\_data }\OtherTok{\textless{}{-}}\NormalTok{ annual\_data }\SpecialCharTok{|\textgreater{}}
  \FunctionTok{group\_by}\NormalTok{(firm\_ID) }\SpecialCharTok{|\textgreater{}}
  \FunctionTok{arrange}\NormalTok{(year, }\AttributeTok{.by\_group =} \ConstantTok{TRUE}\NormalTok{) }\SpecialCharTok{|\textgreater{}}
  \FunctionTok{mutate}\NormalTok{(}
    \AttributeTok{lagged\_ME =} \FunctionTok{lag}\NormalTok{(ME),}
    \AttributeTok{lagged\_BEME =}\NormalTok{ lagged\_BE }\SpecialCharTok{/}\NormalTok{ lagged\_ME,}
    \AttributeTok{lagged\_BEME =} \FunctionTok{if\_else}\NormalTok{(}
      \FunctionTok{is.finite}\NormalTok{(lagged\_BEME) }\SpecialCharTok{\&}
\NormalTok{        lagged\_BEME }\SpecialCharTok{\textgreater{}} \DecValTok{0}\NormalTok{,}
\NormalTok{      lagged\_BEME,}
      \ConstantTok{NA\_real\_}
\NormalTok{    )}
\NormalTok{  ) }\SpecialCharTok{|\textgreater{}}
  \FunctionTok{ungroup}\NormalTok{()}
\end{Highlighting}
\end{Shaded}

回帰に使用できる年度別観測数を確認する。

\begin{Shaded}
\begin{Highlighting}[]
\NormalTok{regression\_observations\_by\_year }\OtherTok{\textless{}{-}}\NormalTok{ annual\_data }\SpecialCharTok{|\textgreater{}}
  \FunctionTok{filter}\NormalTok{(}
    \FunctionTok{is.finite}\NormalTok{(Re),}
    \FunctionTok{is.finite}\NormalTok{(lagged\_BEME)}
\NormalTok{  ) }\SpecialCharTok{|\textgreater{}}
  \FunctionTok{count}\NormalTok{(}
\NormalTok{    year,}
    \AttributeTok{name =} \StringTok{"observations"}
\NormalTok{  ) }\SpecialCharTok{|\textgreater{}}
  \FunctionTok{arrange}\NormalTok{(year)}

\FunctionTok{show\_table\_LFL}\NormalTok{(}
\NormalTok{  regression\_observations\_by\_year,}
  \AttributeTok{n =} \FunctionTok{nrow}\NormalTok{(regression\_observations\_by\_year)}
\NormalTok{)}
\end{Highlighting}
\end{Shaded}

\begin{verbatim}
# A tibble: 5 x 2
   year observations
  <int>        <int>
1  2016         1241
2  2017         1270
3  2018         1268
4  2019         1304
5  2020         1322
\end{verbatim}

\begin{Shaded}
\begin{Highlighting}[]
\NormalTok{eligible\_regression\_years }\OtherTok{\textless{}{-}}
\NormalTok{  regression\_observations\_by\_year }\SpecialCharTok{|\textgreater{}}
  \FunctionTok{filter}\NormalTok{(observations }\SpecialCharTok{\textgreater{}=} \DecValTok{30}\NormalTok{)}

\ControlFlowTok{if}\NormalTok{ (}\FunctionTok{nrow}\NormalTok{(eligible\_regression\_years) }\SpecialCharTok{==} \DecValTok{0}\NormalTok{) \{}
  \FunctionTok{stop}\NormalTok{(}
    \StringTok{"回帰分析に必要な観測数を満たす年度がありません。"}\NormalTok{,}
    \AttributeTok{call. =} \ConstantTok{FALSE}
\NormalTok{  )}
\NormalTok{\}}

\NormalTok{regression\_year }\OtherTok{\textless{}{-}}\NormalTok{ eligible\_regression\_years }\SpecialCharTok{|\textgreater{}}
  \FunctionTok{summarise}\NormalTok{(}\AttributeTok{year =} \FunctionTok{max}\NormalTok{(year)) }\SpecialCharTok{|\textgreater{}}
  \FunctionTok{pull}\NormalTok{(year)}

\NormalTok{lm\_sample\_data }\OtherTok{\textless{}{-}}\NormalTok{ annual\_data }\SpecialCharTok{|\textgreater{}}
  \FunctionTok{filter}\NormalTok{(}
\NormalTok{    year }\SpecialCharTok{==}\NormalTok{ regression\_year,}
    \FunctionTok{is.finite}\NormalTok{(Re),}
    \FunctionTok{is.finite}\NormalTok{(lagged\_BEME)}
\NormalTok{  ) }\SpecialCharTok{|\textgreater{}}
  \FunctionTok{select}\NormalTok{(}
\NormalTok{    firm\_ID,}
\NormalTok{    year,}
\NormalTok{    Re,}
\NormalTok{    lagged\_BEME}
\NormalTok{  )}
\end{Highlighting}
\end{Shaded}

\begin{Shaded}
\begin{Highlighting}[]
\NormalTok{beme\_plot\_limit }\OtherTok{\textless{}{-}} \FunctionTok{quantile}\NormalTok{(}
\NormalTok{  lm\_sample\_data}\SpecialCharTok{$}\NormalTok{lagged\_BEME,}
  \FloatTok{0.99}\NormalTok{,}
  \AttributeTok{na.rm =} \ConstantTok{TRUE}
\NormalTok{)}

\NormalTok{p\_beme\_return\_regression }\OtherTok{\textless{}{-}}\NormalTok{ lm\_sample\_data }\SpecialCharTok{|\textgreater{}}
  \FunctionTok{ggplot}\NormalTok{(}
    \FunctionTok{aes}\NormalTok{(}
      \AttributeTok{x =}\NormalTok{ lagged\_BEME,}
      \AttributeTok{y =}\NormalTok{ Re}
\NormalTok{    )}
\NormalTok{  ) }\SpecialCharTok{+}
  \FunctionTok{geom\_point}\NormalTok{(}\AttributeTok{alpha =} \FloatTok{0.35}\NormalTok{) }\SpecialCharTok{+}
  \FunctionTok{geom\_smooth}\NormalTok{(}
    \AttributeTok{method =} \StringTok{"lm"}\NormalTok{,}
    \AttributeTok{se =} \ConstantTok{FALSE}
\NormalTok{  ) }\SpecialCharTok{+}
  \FunctionTok{coord\_cartesian}\NormalTok{(}
    \AttributeTok{xlim =} \FunctionTok{c}\NormalTok{(}\DecValTok{0}\NormalTok{, beme\_plot\_limit)}
\NormalTok{  ) }\SpecialCharTok{+}
  \FunctionTok{scale\_y\_continuous}\NormalTok{(}
    \AttributeTok{labels =}\NormalTok{ scales}\SpecialCharTok{::}\FunctionTok{label\_percent}\NormalTok{()}
\NormalTok{  ) }\SpecialCharTok{+}
  \FunctionTok{labs}\NormalTok{(}
    \AttributeTok{x =} \StringTok{"Lagged Book{-}to{-}Market Ratio"}\NormalTok{,}
    \AttributeTok{y =} \StringTok{"Annual Excess Return"}\NormalTok{,}
    \AttributeTok{title =} \FunctionTok{paste}\NormalTok{(}
      \StringTok{"Book{-}to{-}Market Ratio and Excess Return in"}\NormalTok{,}
\NormalTok{      regression\_year}
\NormalTok{    )}
\NormalTok{  )}

\FunctionTok{add\_lagrange\_theme\_LFL}\NormalTok{(}
\NormalTok{  p\_beme\_return\_regression}
\NormalTok{)}
\end{Highlighting}
\end{Shaded}

\pandocbounded{\includegraphics[keepaspectratio]{part1_empirical/ch03_stock_data_time_series_files/figure-pdf/plot-beme-return-regression-1.pdf}}

\begin{Shaded}
\begin{Highlighting}[]
\NormalTok{lm\_result }\OtherTok{\textless{}{-}} \FunctionTok{lm}\NormalTok{(}
\NormalTok{  Re }\SpecialCharTok{\textasciitilde{}}\NormalTok{ lagged\_BEME,}
  \AttributeTok{data =}\NormalTok{ lm\_sample\_data}
\NormalTok{)}

\FunctionTok{show\_table\_LFL}\NormalTok{(}
\NormalTok{  broom}\SpecialCharTok{::}\FunctionTok{tidy}\NormalTok{(lm\_result)}
\NormalTok{)}
\end{Highlighting}
\end{Shaded}

\begin{verbatim}
# A tibble: 2 x 5
  term           estimate    std.error statistic    p.value
  <chr>             <dbl>        <dbl>     <dbl>      <dbl>
1 (Intercept)    -0.10237     0.012394   -8.2593 3.5232e-16
2 lagged_BEME 44359.      10973.          4.0426 5.5911e- 5
\end{verbatim}

\begin{Shaded}
\begin{Highlighting}[]
\FunctionTok{show\_table\_LFL}\NormalTok{(}
\NormalTok{  broom}\SpecialCharTok{::}\FunctionTok{glance}\NormalTok{(lm\_result)}
\NormalTok{)}
\end{Highlighting}
\end{Shaded}

\begin{verbatim}
# A tibble: 1 x 12
  r.squared adj.r.squared   sigma statistic     p.value    df  logLik    AIC
      <dbl>         <dbl>   <dbl>     <dbl>       <dbl> <dbl>   <dbl>  <dbl>
1  0.012230      0.011481 0.24523    16.343 0.000055911     1 -16.667 39.334
     BIC deviance df.residual  nobs
   <dbl>    <dbl>       <int> <int>
1 54.895   79.379        1320  1322
\end{verbatim}

\subsection{3.6.3 年度別回帰}\label{ux5e74ux5ea6ux5225ux56deux5e30}

年度ごとの回帰を\texttt{nest()}と\texttt{purrr::map()}で実行する。

\begin{Shaded}
\begin{Highlighting}[]
\NormalTok{yearly\_regression\_models }\OtherTok{\textless{}{-}}\NormalTok{ annual\_data }\SpecialCharTok{|\textgreater{}}
  \FunctionTok{filter}\NormalTok{(}
    \FunctionTok{is.finite}\NormalTok{(Re),}
    \FunctionTok{is.finite}\NormalTok{(lagged\_BEME)}
\NormalTok{  ) }\SpecialCharTok{|\textgreater{}}
  \FunctionTok{group\_by}\NormalTok{(year) }\SpecialCharTok{|\textgreater{}}
  \FunctionTok{filter}\NormalTok{(}\FunctionTok{n}\NormalTok{() }\SpecialCharTok{\textgreater{}=} \DecValTok{30}\NormalTok{) }\SpecialCharTok{|\textgreater{}}
  \FunctionTok{ungroup}\NormalTok{() }\SpecialCharTok{|\textgreater{}}
  \FunctionTok{nest}\NormalTok{(}\AttributeTok{data =} \SpecialCharTok{{-}}\NormalTok{year) }\SpecialCharTok{|\textgreater{}}
  \FunctionTok{mutate}\NormalTok{(}
    \AttributeTok{model =} \FunctionTok{map}\NormalTok{(}
\NormalTok{      data,}
\NormalTok{      \textbackslash{}(data) }\FunctionTok{lm}\NormalTok{(}
\NormalTok{        Re }\SpecialCharTok{\textasciitilde{}}\NormalTok{ lagged\_BEME,}
        \AttributeTok{data =}\NormalTok{ data}
\NormalTok{      )}
\NormalTok{    ),}
    \AttributeTok{coefficients =} \FunctionTok{map}\NormalTok{(}
\NormalTok{      model,}
\NormalTok{      broom}\SpecialCharTok{::}\NormalTok{tidy}
\NormalTok{    ),}
    \AttributeTok{fit =} \FunctionTok{map}\NormalTok{(}
\NormalTok{      model,}
\NormalTok{      broom}\SpecialCharTok{::}\NormalTok{glance}
\NormalTok{    )}
\NormalTok{  )}
\end{Highlighting}
\end{Shaded}

\begin{Shaded}
\begin{Highlighting}[]
\NormalTok{yearly\_regression\_coefficients }\OtherTok{\textless{}{-}}
\NormalTok{  yearly\_regression\_models }\SpecialCharTok{|\textgreater{}}
  \FunctionTok{select}\NormalTok{(year, coefficients) }\SpecialCharTok{|\textgreater{}}
  \FunctionTok{unnest}\NormalTok{(coefficients) }\SpecialCharTok{|\textgreater{}}
  \FunctionTok{arrange}\NormalTok{(year, term)}

\NormalTok{yearly\_regression\_fit }\OtherTok{\textless{}{-}}
\NormalTok{  yearly\_regression\_models }\SpecialCharTok{|\textgreater{}}
  \FunctionTok{select}\NormalTok{(year, fit) }\SpecialCharTok{|\textgreater{}}
  \FunctionTok{unnest}\NormalTok{(fit) }\SpecialCharTok{|\textgreater{}}
  \FunctionTok{select}\NormalTok{(}
\NormalTok{    year,}
\NormalTok{    r.squared,}
\NormalTok{    adj.r.squared,}
\NormalTok{    sigma,}
\NormalTok{    statistic,}
\NormalTok{    p.value,}
\NormalTok{    nobs}
\NormalTok{  ) }\SpecialCharTok{|\textgreater{}}
  \FunctionTok{arrange}\NormalTok{(year)}

\FunctionTok{show\_table\_LFL}\NormalTok{(}
\NormalTok{  yearly\_regression\_coefficients,}
  \AttributeTok{n =} \FunctionTok{nrow}\NormalTok{(yearly\_regression\_coefficients)}
\NormalTok{)}
\end{Highlighting}
\end{Shaded}

\begin{verbatim}
# A tibble: 10 x 6
    year term             estimate    std.error statistic    p.value
   <int> <chr>               <dbl>        <dbl>     <dbl>      <dbl>
 1  2016 (Intercept)      0.056439     0.016563    3.4076 6.7644e- 4
 2  2016 lagged_BEME  72103.       12275.          5.8738 5.4661e- 9
 3  2017 (Intercept)      0.17989      0.015673   11.478  4.4094e-29
 4  2017 lagged_BEME  32467.       11890.          2.7306 6.4094e- 3
 5  2018 (Intercept)     -0.12679      0.015203   -8.3395 1.9259e-16
 6  2018 lagged_BEME  65283.       13067.          4.9959 6.6754e- 7
 7  2019 (Intercept)      0.28895      0.026406   10.942  1.0121e-26
 8  2019 lagged_BEME 100567.       19057.          5.2771 1.5363e- 7
 9  2020 (Intercept)     -0.10237      0.012394   -8.2593 3.5232e-16
10  2020 lagged_BEME  44359.       10973.          4.0426 5.5911e- 5
\end{verbatim}

\begin{Shaded}
\begin{Highlighting}[]
\FunctionTok{show\_table\_LFL}\NormalTok{(}
\NormalTok{  yearly\_regression\_fit,}
  \AttributeTok{n =} \FunctionTok{nrow}\NormalTok{(yearly\_regression\_fit)}
\NormalTok{)}
\end{Highlighting}
\end{Shaded}

\begin{verbatim}
# A tibble: 5 x 7
   year r.squared adj.r.squared   sigma statistic   p.value  nobs
  <int>     <dbl>         <dbl>   <dbl>     <dbl>     <dbl> <int>
1  2016 0.027092      0.026306  0.28691   34.501  5.4661e-9  1241
2  2017 0.0058459     0.0050619 0.28492    7.4562 6.4094e-3  1270
3  2018 0.019334      0.018559  0.27721   24.959  6.6754e-7  1268
4  2019 0.020940      0.020188  0.49819   27.847  1.5363e-7  1304
5  2020 0.012230      0.011481  0.24523   16.343  5.5911e-5  1322
\end{verbatim}

\section{3.7
分析用データの保存}\label{ux5206ux6790ux7528ux30c7ux30fcux30bfux306eux4fddux5b58}

第4章以降で使用するため、月次データと年次データを\texttt{simulation\_data/derived}へ保存する。

\begin{Shaded}
\begin{Highlighting}[]
\NormalTok{output\_data }\OtherTok{\textless{}{-}} \FunctionTok{list}\NormalTok{(}
  \AttributeTok{monthly =}\NormalTok{ monthly\_data,}
  \AttributeTok{annual =}\NormalTok{ annual\_data}
\NormalTok{)}

\NormalTok{output\_files }\OtherTok{\textless{}{-}} \FunctionTok{c}\NormalTok{(}
  \AttributeTok{monthly =} \StringTok{"ch05\_output1.csv"}\NormalTok{,}
  \AttributeTok{annual =} \StringTok{"ch05\_output2.csv"}
\NormalTok{)}

\NormalTok{output\_paths }\OtherTok{\textless{}{-}}\NormalTok{ output\_data }\SpecialCharTok{|\textgreater{}}
  \FunctionTok{imap\_chr}\NormalTok{(}
\NormalTok{    \textbackslash{}(data, data\_name)}
      \FunctionTok{write\_derived\_csv\_file\_LFL}\NormalTok{(}
\NormalTok{        data,}
\NormalTok{        output\_files[[data\_name]]}
\NormalTok{      )}
\NormalTok{  )}

\NormalTok{output\_paths}
\end{Highlighting}
\end{Shaded}

\begin{verbatim}
                                                                           monthly 
"C:/AnalyticFin/Projects/LagrangeFinance/simulation_data/derived/ch05_output1.csv" 
                                                                            annual 
"C:/AnalyticFin/Projects/LagrangeFinance/simulation_data/derived/ch05_output2.csv" 
\end{verbatim}

保存したファイルを読み直し、行数と列名を検証する。

\begin{Shaded}
\begin{Highlighting}[]
\NormalTok{output\_verification }\OtherTok{\textless{}{-}}\NormalTok{ output\_data }\SpecialCharTok{|\textgreater{}}
  \FunctionTok{imap\_dfr}\NormalTok{(}
\NormalTok{    \textbackslash{}(expected\_data, data\_name) \{}
\NormalTok{      output\_check }\OtherTok{\textless{}{-}}
        \FunctionTok{read\_derived\_csv\_file\_LFL}\NormalTok{(}
\NormalTok{          output\_files[[data\_name]]}
\NormalTok{        )}

      \FunctionTok{stopifnot}\NormalTok{(}
        \FunctionTok{nrow}\NormalTok{(output\_check) }\SpecialCharTok{==}
          \FunctionTok{nrow}\NormalTok{(expected\_data)}
\NormalTok{      )}

      \FunctionTok{stopifnot}\NormalTok{(}
        \FunctionTok{identical}\NormalTok{(}
          \FunctionTok{names}\NormalTok{(output\_check),}
          \FunctionTok{names}\NormalTok{(expected\_data)}
\NormalTok{        )}
\NormalTok{      )}

      \FunctionTok{tibble}\NormalTok{(}
        \AttributeTok{data =}\NormalTok{ data\_name,}
        \AttributeTok{output\_file =}\NormalTok{ output\_paths[[data\_name]],}
        \AttributeTok{rows =} \FunctionTok{nrow}\NormalTok{(output\_check),}
        \AttributeTok{columns =} \FunctionTok{ncol}\NormalTok{(output\_check),}
        \AttributeTok{size\_bytes =}
          \FunctionTok{file.info}\NormalTok{(}
\NormalTok{            output\_paths[[data\_name]]}
\NormalTok{          )}\SpecialCharTok{$}\NormalTok{size}
\NormalTok{      )}
\NormalTok{    \}}
\NormalTok{  )}

\FunctionTok{show\_table\_LFL}\NormalTok{(output\_verification)}
\end{Highlighting}
\end{Shaded}

\begin{verbatim}
# A tibble: 2 x 5
  data   
  <chr>  
1 monthly
2 annual 
  output_file                                                                   
  <chr>                                                                         
1 C:/AnalyticFin/Projects/LagrangeFinance/simulation_data/derived/ch05_output1.~
2 C:/AnalyticFin/Projects/LagrangeFinance/simulation_data/derived/ch05_output2.~
   rows columns size_bytes
  <int>   <int>      <dbl>
1 95040      39   39813837
2  7920      31    3022818
\end{verbatim}

\section{3.8 まとめ}\label{ux307eux3068ux3081-2}

本章では、株価、配当および調整係数から月次リターンを作成し、無リスク金利を差し引いて超過リターンを計算した。さらに、月次リターンを年次リターンへ集約し、第2章で作成した財務データと結合した。最後に、簿価時価比率と年次超過リターンの関係を回帰分析で確認し、月次データを\texttt{ch05\_output1.csv}、年次データを\texttt{ch05\_output2.csv}として保存した。

第4章では、本章で作成した分析用データを利用して、ポートフォリオ分析へ進む。

\chapter{第4章
ファクターモデル}\label{ux7b2c4ux7ae0-ux30d5ux30a1ux30afux30bfux30fcux30e2ux30c7ux30eb}

本章では、第3章で作成した月次株式データと年次財務データを利用し、市場ポートフォリオ、企業規模ポートフォリオ、簿価時価比率ポートフォリオ、SMB、HMLを構築する。さらに、企業規模10分位ポートフォリオを対象として、CAPMとFama--Frenchの3ファクター・モデルを推定する。

本章の処理は、次の順序で進める。

\begin{enumerate}
\def\labelenumi{\arabic{enumi}.}
\tightlist
\item
  第3章の月次・年次データを読み込む
\item
  前年度末の情報からポートフォリオを形成する
\item
  市場ポートフォリオと市場超過リターンを計算する
\item
  企業規模と簿価時価比率による単変量ソートを行う
\item
  2×3ポートフォリオからSMBとHMLを構築する
\item
  CAPMとFF3モデルを推定して比較する
\item
  第5章以降で使用するファクター・データを保存する
\end{enumerate}

反復処理には明示的なループを使わず、\texttt{purrr::map()}系を利用する。

\section{4.1 分析環境}\label{ux5206ux6790ux74b0ux5883-2}

\begin{Shaded}
\begin{Highlighting}[]
\FunctionTok{source}\NormalTok{(}\StringTok{\textquotesingle{}C:/AnalyticFin/Projects/LagrangeFinance/R/utility\_LFL.R\textquotesingle{}}\NormalTok{)}
\FunctionTok{load\_part1\_packages\_LFL}\NormalTok{()}
\end{Highlighting}
\end{Shaded}

本章では、第1章から第3章と同じ共通ユーティリティを利用する。ファイル探索、CSV読込み、表の表示、グラフのテーマ設定を各章で重複して定義しない。

\section{4.2 入力データ}\label{ux5165ux529bux30c7ux30fcux30bf}

\subsection{4.2.1
第3章の出力の読込み}\label{ux7b2c3ux7ae0ux306eux51faux529bux306eux8aadux8fbcux307f}

第3章で作成した次の2ファイルを読み込む。

\begin{Shaded}
\begin{Highlighting}[]
\NormalTok{simulation\_data/derived/ch05\_output1.csv}
\NormalTok{simulation\_data/derived/ch05\_output2.csv}
\end{Highlighting}
\end{Shaded}

\texttt{ch05\_output1.csv}は月次データ、\texttt{ch05\_output2.csv}は年次データである。ファイル名は旧教材の章番号に由来する。

\begin{Shaded}
\begin{Highlighting}[]
\NormalTok{tbl\_monthly }\OtherTok{\textless{}{-}} \FunctionTok{read\_derived\_csv\_LFL}\NormalTok{(}\StringTok{"ch05\_output1.csv"}\NormalTok{) }\SpecialCharTok{|\textgreater{}}
  \FunctionTok{mutate}\NormalTok{(}
    \AttributeTok{firm\_ID =} \FunctionTok{as.integer}\NormalTok{(firm\_ID),}
    \AttributeTok{year =} \FunctionTok{as.integer}\NormalTok{(year),}
    \AttributeTok{month =} \FunctionTok{as.integer}\NormalTok{(month),}
    \AttributeTok{month\_ID =} \FunctionTok{as.integer}\NormalTok{(month\_ID),}
    \AttributeTok{month\_date =} \FunctionTok{as.Date}\NormalTok{(month\_date)}
\NormalTok{  ) }\SpecialCharTok{|\textgreater{}}
  \FunctionTok{arrange}\NormalTok{(firm\_ID, month\_ID)}

\NormalTok{tbl\_annual }\OtherTok{\textless{}{-}} \FunctionTok{read\_derived\_csv\_LFL}\NormalTok{(}\StringTok{"ch05\_output2.csv"}\NormalTok{) }\SpecialCharTok{|\textgreater{}}
  \FunctionTok{mutate}\NormalTok{(}
    \AttributeTok{firm\_ID =} \FunctionTok{as.integer}\NormalTok{(firm\_ID),}
    \AttributeTok{year =} \FunctionTok{as.integer}\NormalTok{(year)}
\NormalTok{  ) }\SpecialCharTok{|\textgreater{}}
  \FunctionTok{arrange}\NormalTok{(firm\_ID, year)}

\FunctionTok{show\_data\_structure\_LFL}\NormalTok{(tbl\_monthly)}
\end{Highlighting}
\end{Shaded}

\begin{verbatim}
Rows: 95,040
Columns: 39
$ year                   <int> 2015, 2015, 2015, 2015, 2015, 2015, 2015, 2015,~
$ month                  <int> 1, 2, 3, 4, 5, 6, 7, 8, 9, 10, 11, 12, 1, 2, 3,~
$ month_ID               <int> 1, 2, 3, 4, 5, 6, 7, 8, 9, 10, 11, 12, 13, 14, ~
$ firm_ID                <int> 1, 1, 1, 1, 1, 1, 1, 1, 1, 1, 1, 1, 1, 1, 1, 1,~
$ stock_price            <dbl> 954, 960, 1113, 1081, 1317, 1366, 1353, 1209, 1~
$ DPS                    <dbl> 0, 0, 0, 0, 0, 29, 0, 0, 0, 0, 0, 29, 0, 0, 0, ~
$ shares_outstanding     <dbl> 2422000, 2422000, 2422000, 2422000, 2422000, 24~
$ adjustment_coefficient <dbl> 1, 1, 1, 1, 1, 1, 1, 1, 1, 1, 1, 1, 1, 1, 1, 1,~
$ month_date             <date> 2015-01-01, 2015-02-01, 2015-03-01, 2015-04-01~
$ ME                     <dbl> 2310588000, 2325120000, 2695686000, 2618182000,~
$ lagged_stock_price     <dbl> NA, 954, 960, 1113, 1081, 1317, 1366, 1353, 120~
$ lagged_month_ID        <dbl> NA, 1, 2, 3, 4, 5, 6, 7, 8, 9, 10, 11, 12, 13, ~
$ consecutive_month      <lgl> NA, TRUE, TRUE, TRUE, TRUE, TRUE, TRUE, TRUE, T~
$ R                      <dbl> NA, 0.006289308, 0.159375000, -0.028751123, 0.2~
$ R_F                    <dbl> 0.00008202927, 0.00008202927, 0.00008202927, 0.~
$ Re                     <dbl> NA, 0.006207279, 0.159292971, -0.028833152, 0.2~
$ industry_ID            <dbl> 1, 1, 1, 1, 1, 1, 1, 1, 1, 1, 1, 1, 1, 1, 1, 1,~
$ sales                  <dbl> 5261.40, 5261.40, 5261.40, 5261.40, 5261.40, 52~
$ OX                     <dbl> 437.49, 437.49, 437.49, 437.49, 437.49, 437.49,~
$ NFE                    <dbl> NA, NA, NA, NA, NA, NA, NA, NA, NA, NA, NA, NA,~
$ X                      <dbl> 286.64, 286.64, 286.64, 286.64, 286.64, 286.64,~
$ OA                     <dbl> 13005.55, 13005.55, 13005.55, 13005.55, 13005.5~
$ FA                     <dbl> 3543.43, 3543.43, 3543.43, 3543.43, 3543.43, 35~
$ OL                     <dbl> 4372.96, 4372.96, 4372.96, 4372.96, 4372.96, 43~
$ FO                     <dbl> 2480.72, 2480.72, 2480.72, 2480.72, 2480.72, 24~
$ NOA                    <dbl> 8632.59, 8632.59, 8632.59, 8632.59, 8632.59, 86~
$ NFO                    <dbl> -1062.71, -1062.71, -1062.71, -1062.71, -1062.7~
$ BE                     <dbl> 9695.30, 9695.30, 9695.30, 9695.30, 9695.30, 96~
$ lagged_BE              <dbl> NA, NA, NA, NA, NA, NA, NA, NA, NA, NA, NA, NA,~
$ ROE                    <dbl> NA, NA, NA, NA, NA, NA, NA, NA, NA, NA, NA, NA,~
$ lagged_NOA             <dbl> NA, NA, NA, NA, NA, NA, NA, NA, NA, NA, NA, NA,~
$ lagged_NFO             <dbl> NA, NA, NA, NA, NA, NA, NA, NA, NA, NA, NA, NA,~
$ RNOA                   <dbl> NA, NA, NA, NA, NA, NA, NA, NA, NA, NA, NA, NA,~
$ PM                     <dbl> 0.08315087, 0.08315087, 0.08315087, 0.08315087,~
$ ATO                    <dbl> NA, NA, NA, NA, NA, NA, NA, NA, NA, NA, NA, NA,~
$ lagged_FLEV            <dbl> NA, NA, NA, NA, NA, NA, NA, NA, NA, NA, NA, NA,~
$ NBC                    <dbl> NA, NA, NA, NA, NA, NA, NA, NA, NA, NA, NA, NA,~
$ ROE_DuPont             <dbl> NA, NA, NA, NA, NA, NA, NA, NA, NA, NA, NA, NA,~
$ DuPont_gap             <dbl> NA, NA, NA, NA, NA, NA, NA, NA, NA, NA, NA, NA,~
\end{verbatim}

\begin{Shaded}
\begin{Highlighting}[]
\FunctionTok{show\_data\_structure\_LFL}\NormalTok{(tbl\_annual)}
\end{Highlighting}
\end{Shaded}

\begin{verbatim}
Rows: 7,920
Columns: 31
$ firm_ID     <int> 1, 1, 1, 1, 1, 1, 2, 2, 2, 2, 2, 2, 3, 3, 3, 3, 3, 3, 4, 4~
$ year        <int> 2015, 2016, 2017, 2018, 2019, 2020, 2015, 2016, 2017, 2018~
$ months      <dbl> 11, 12, 12, 12, 12, 12, 11, 12, 12, 12, 12, 12, 11, 12, 12~
$ R           <dbl> 0.61213017, 0.99727265, 0.68786382, -0.21361287, 0.6468357~
$ Re          <dbl> 0.61073273, 0.99547052, 0.68636684, -0.21438527, 0.6453315~
$ ME          <dbl> 3577294000, 6883324000, 11376990000, 8694752000, 139575180~
$ industry_ID <dbl> 1, 1, 1, 1, 1, 1, 1, 1, 1, 1, 1, 1, 1, 1, 1, 1, 1, 1, 1, 1~
$ sales       <dbl> 5261.40, 5948.96, 6505.06, 6846.38, 7572.24, 7537.63, 3505~
$ OX          <dbl> 437.49, 564.14, 691.18, 751.29, 958.53, 778.37, 45.82, 51.~
$ NFE         <dbl> NA, 50.667498, 29.543157, 86.486500, 298.049774, -65.45877~
$ X           <dbl> 286.64, 513.48, 661.64, 664.80, 660.48, 843.83, 40.07, 49.~
$ OA          <dbl> 13005.55, 13865.58, 13952.58, 18818.48, 18190.00, 20462.86~
$ FA          <dbl> 3543.43, 4642.16, 7743.99, 7284.72, 9735.13, 10274.25, 225~
$ OL          <dbl> 4372.96, 4534.22, 5111.22, 5137.28, 5487.96, 5371.38, 1840~
$ FO          <dbl> 2480.72, 3959.70, 6159.02, 10123.91, 11362.22, 13772.15, 2~
$ NOA         <dbl> 8632.59, 9331.36, 8841.36, 13681.20, 12702.04, 15091.48, 1~
$ NFO         <dbl> -1062.71, -682.46, -1584.97, 2839.19, 1627.09, 3497.90, 82~
$ BE          <dbl> 9695.30, 10013.82, 10426.33, 10842.01, 11074.95, 11593.58,~
$ lagged_BE   <dbl> NA, 9695.30, 10013.82, 10426.33, 10842.01, 11074.95, NA, 1~
$ ROE         <dbl> NA, 0.05296174, 0.06607269, 0.06376165, 0.06091859, 0.0761~
$ lagged_NOA  <dbl> NA, 8632.59, 9331.36, 8841.36, 13681.20, 12702.04, NA, 113~
$ lagged_NFO  <dbl> NA, -1062.71, -682.46, -1584.97, 2839.19, 1627.09, NA, 82.~
$ RNOA        <dbl> NA, 0.06535003, 0.07407066, 0.08497448, 0.07006184, 0.0612~
$ PM          <dbl> 0.083150872, 0.094830021, 0.106252671, 0.109735364, 0.1265~
$ ATO         <dbl> NA, 0.6891281, 0.6971181, 0.7743582, 0.5534778, 0.5934189,~
$ lagged_FLEV <dbl> NA, -0.10961084, -0.06815181, -0.15201610, 0.26186934, 0.1~
$ NBC         <dbl> NA, -0.047677633, -0.043289214, -0.054566648, 0.104977044,~
$ ROE_DuPont  <dbl> NA, 0.05296097, 0.06607237, 0.06376199, 0.06091861, 0.0761~
$ DuPont_gap  <dbl> NA, 0.00000077332412, 0.00000031526680, -0.00000033567247,~
$ lagged_ME   <dbl> NA, 3577294000, 6883324000, 11376990000, 8694752000, 13957~
$ lagged_BEME <dbl> NA, 0.0000027102329, 0.0000014547942, 0.0000009164401, 0.0~
\end{verbatim}

\begin{Shaded}
\begin{Highlighting}[]
\NormalTok{tbl\_input\_overview }\OtherTok{\textless{}{-}} \FunctionTok{tribble}\NormalTok{(}
  \SpecialCharTok{\textasciitilde{}}\NormalTok{data, }\SpecialCharTok{\textasciitilde{}}\NormalTok{observations, }\SpecialCharTok{\textasciitilde{}}\NormalTok{variables, }\SpecialCharTok{\textasciitilde{}}\NormalTok{firms, }\SpecialCharTok{\textasciitilde{}}\NormalTok{first\_year, }\SpecialCharTok{\textasciitilde{}}\NormalTok{last\_year,}
  \StringTok{"monthly"}\NormalTok{, }\FunctionTok{nrow}\NormalTok{(tbl\_monthly), }\FunctionTok{ncol}\NormalTok{(tbl\_monthly),}
  \FunctionTok{n\_distinct}\NormalTok{(tbl\_monthly}\SpecialCharTok{$}\NormalTok{firm\_ID),}
  \FunctionTok{min}\NormalTok{(tbl\_monthly}\SpecialCharTok{$}\NormalTok{year, }\AttributeTok{na.rm =} \ConstantTok{TRUE}\NormalTok{),}
  \FunctionTok{max}\NormalTok{(tbl\_monthly}\SpecialCharTok{$}\NormalTok{year, }\AttributeTok{na.rm =} \ConstantTok{TRUE}\NormalTok{),}
  \StringTok{"annual"}\NormalTok{, }\FunctionTok{nrow}\NormalTok{(tbl\_annual), }\FunctionTok{ncol}\NormalTok{(tbl\_annual),}
  \FunctionTok{n\_distinct}\NormalTok{(tbl\_annual}\SpecialCharTok{$}\NormalTok{firm\_ID),}
  \FunctionTok{min}\NormalTok{(tbl\_annual}\SpecialCharTok{$}\NormalTok{year, }\AttributeTok{na.rm =} \ConstantTok{TRUE}\NormalTok{),}
  \FunctionTok{max}\NormalTok{(tbl\_annual}\SpecialCharTok{$}\NormalTok{year, }\AttributeTok{na.rm =} \ConstantTok{TRUE}\NormalTok{)}
\NormalTok{)}

\FunctionTok{show\_table\_LFL}\NormalTok{(tbl\_input\_overview)}
\end{Highlighting}
\end{Shaded}

\begin{verbatim}
# A tibble: 2 x 6
  data    observations variables firms first_year last_year
  <chr>          <int>     <int> <int>      <int>     <int>
1 monthly        95040        39  1515       2015      2020
2 annual          7920        31  1515       2015      2020
\end{verbatim}

\subsection{4.2.2
必須列とパネルキー}\label{ux5fc5ux9808ux5217ux3068ux30d1ux30cdux30ebux30adux30fc}

分析を開始する前に、必要な列とパネルキーの一意性を確認する。

\begin{Shaded}
\begin{Highlighting}[]
\NormalTok{vec\_required\_monthly\_columns }\OtherTok{\textless{}{-}} \FunctionTok{c}\NormalTok{(}
  \StringTok{"year"}\NormalTok{, }\StringTok{"month"}\NormalTok{, }\StringTok{"month\_ID"}\NormalTok{, }\StringTok{"month\_date"}\NormalTok{,}
  \StringTok{"firm\_ID"}\NormalTok{, }\StringTok{"R"}\NormalTok{, }\StringTok{"R\_F"}\NormalTok{, }\StringTok{"Re"}\NormalTok{, }\StringTok{"ME"}
\NormalTok{)}

\NormalTok{vec\_required\_annual\_columns }\OtherTok{\textless{}{-}} \FunctionTok{c}\NormalTok{(}
  \StringTok{"year"}\NormalTok{, }\StringTok{"firm\_ID"}\NormalTok{, }\StringTok{"ME"}\NormalTok{, }\StringTok{"BE"}\NormalTok{, }\StringTok{"lagged\_BE"}
\NormalTok{)}

\NormalTok{vec\_missing\_monthly\_columns }\OtherTok{\textless{}{-}} \FunctionTok{setdiff}\NormalTok{(}
\NormalTok{  vec\_required\_monthly\_columns,}
  \FunctionTok{names}\NormalTok{(tbl\_monthly)}
\NormalTok{)}

\NormalTok{vec\_missing\_annual\_columns }\OtherTok{\textless{}{-}} \FunctionTok{setdiff}\NormalTok{(}
\NormalTok{  vec\_required\_annual\_columns,}
  \FunctionTok{names}\NormalTok{(tbl\_annual)}
\NormalTok{)}

\ControlFlowTok{if}\NormalTok{ (}\FunctionTok{length}\NormalTok{(vec\_missing\_monthly\_columns) }\SpecialCharTok{\textgreater{}} \DecValTok{0}\NormalTok{) \{}
  \FunctionTok{stop}\NormalTok{(}
    \StringTok{"月次データに必要な列が不足しています: "}\NormalTok{,}
    \FunctionTok{paste}\NormalTok{(vec\_missing\_monthly\_columns, }\AttributeTok{collapse =} \StringTok{", "}\NormalTok{)}
\NormalTok{  )}
\NormalTok{\}}

\ControlFlowTok{if}\NormalTok{ (}\FunctionTok{length}\NormalTok{(vec\_missing\_annual\_columns) }\SpecialCharTok{\textgreater{}} \DecValTok{0}\NormalTok{) \{}
  \FunctionTok{stop}\NormalTok{(}
    \StringTok{"年次データに必要な列が不足しています: "}\NormalTok{,}
    \FunctionTok{paste}\NormalTok{(vec\_missing\_annual\_columns, }\AttributeTok{collapse =} \StringTok{", "}\NormalTok{)}
\NormalTok{  )}
\NormalTok{\}}

\NormalTok{tbl\_monthly\_key\_duplicates }\OtherTok{\textless{}{-}}\NormalTok{ tbl\_monthly }\SpecialCharTok{|\textgreater{}}
  \FunctionTok{count}\NormalTok{(firm\_ID, month\_ID, }\AttributeTok{name =} \StringTok{"observations"}\NormalTok{) }\SpecialCharTok{|\textgreater{}}
  \FunctionTok{filter}\NormalTok{(observations }\SpecialCharTok{\textgreater{}} \DecValTok{1}\NormalTok{)}

\NormalTok{tbl\_annual\_key\_duplicates }\OtherTok{\textless{}{-}}\NormalTok{ tbl\_annual }\SpecialCharTok{|\textgreater{}}
  \FunctionTok{count}\NormalTok{(firm\_ID, year, }\AttributeTok{name =} \StringTok{"observations"}\NormalTok{) }\SpecialCharTok{|\textgreater{}}
  \FunctionTok{filter}\NormalTok{(observations }\SpecialCharTok{\textgreater{}} \DecValTok{1}\NormalTok{)}

\ControlFlowTok{if}\NormalTok{ (}\FunctionTok{nrow}\NormalTok{(tbl\_monthly\_key\_duplicates) }\SpecialCharTok{\textgreater{}} \DecValTok{0}\NormalTok{) \{}
  \FunctionTok{stop}\NormalTok{(}\StringTok{"月次データのfirm\_IDとmonth\_IDの組合せに重複があります。"}\NormalTok{)}
\NormalTok{\}}

\ControlFlowTok{if}\NormalTok{ (}\FunctionTok{nrow}\NormalTok{(tbl\_annual\_key\_duplicates) }\SpecialCharTok{\textgreater{}} \DecValTok{0}\NormalTok{) \{}
  \FunctionTok{stop}\NormalTok{(}\StringTok{"年次データのfirm\_IDとyearの組合せに重複があります。"}\NormalTok{)}
\NormalTok{\}}

\FunctionTok{show\_table\_LFL}\NormalTok{(tbl\_monthly\_key\_duplicates)}
\end{Highlighting}
\end{Shaded}

\begin{verbatim}
# A tibble: 0 x 3
# i 3 variables: firm_ID <int>, month_ID <int>, observations <int>
\end{verbatim}

\begin{Shaded}
\begin{Highlighting}[]
\FunctionTok{show\_table\_LFL}\NormalTok{(tbl\_annual\_key\_duplicates)}
\end{Highlighting}
\end{Shaded}

\begin{verbatim}
# A tibble: 0 x 3
# i 3 variables: firm_ID <int>, year <int>, observations <int>
\end{verbatim}

\subsection{4.2.3
分析用補助関数}\label{ux5206ux6790ux7528ux88dcux52a9ux95a2ux6570}

次の関数は、欠損値や非有限値を除外して要約統計量と加重平均を計算する。

\begin{Shaded}
\begin{Highlighting}[]
\NormalTok{compute\_median\_safe }\OtherTok{\textless{}{-}} \ControlFlowTok{function}\NormalTok{(x) \{}
\NormalTok{  values }\OtherTok{\textless{}{-}}\NormalTok{ x[}\FunctionTok{is.finite}\NormalTok{(x)]}
  \ControlFlowTok{if}\NormalTok{ (}\FunctionTok{length}\NormalTok{(values) }\SpecialCharTok{==} \DecValTok{0}\NormalTok{) }\FunctionTok{return}\NormalTok{(}\ConstantTok{NA\_real\_}\NormalTok{)}
  \FunctionTok{median}\NormalTok{(values)}
\NormalTok{\}}

\NormalTok{compute\_mean\_safe }\OtherTok{\textless{}{-}} \ControlFlowTok{function}\NormalTok{(x) \{}
\NormalTok{  values }\OtherTok{\textless{}{-}}\NormalTok{ x[}\FunctionTok{is.finite}\NormalTok{(x)]}
  \ControlFlowTok{if}\NormalTok{ (}\FunctionTok{length}\NormalTok{(values) }\SpecialCharTok{==} \DecValTok{0}\NormalTok{) }\FunctionTok{return}\NormalTok{(}\ConstantTok{NA\_real\_}\NormalTok{)}
  \FunctionTok{mean}\NormalTok{(values)}
\NormalTok{\}}

\NormalTok{compute\_weighted\_mean\_safe }\OtherTok{\textless{}{-}} \ControlFlowTok{function}\NormalTok{(x, weight) \{}
\NormalTok{  valid }\OtherTok{\textless{}{-}} \FunctionTok{is.finite}\NormalTok{(x) }\SpecialCharTok{\&} \FunctionTok{is.finite}\NormalTok{(weight) }\SpecialCharTok{\&}\NormalTok{ weight }\SpecialCharTok{\textgreater{}} \DecValTok{0}
  \ControlFlowTok{if}\NormalTok{ (}\SpecialCharTok{!}\FunctionTok{any}\NormalTok{(valid)) }\FunctionTok{return}\NormalTok{(}\ConstantTok{NA\_real\_}\NormalTok{)}
  \FunctionTok{sum}\NormalTok{(x[valid] }\SpecialCharTok{*}\NormalTok{ weight[valid]) }\SpecialCharTok{/} \FunctionTok{sum}\NormalTok{(weight[valid])}
\NormalTok{\}}

\NormalTok{compute\_significance\_stars }\OtherTok{\textless{}{-}} \ControlFlowTok{function}\NormalTok{(p\_value) \{}
  \FunctionTok{case\_when}\NormalTok{(}
    \FunctionTok{is.na}\NormalTok{(p\_value) }\SpecialCharTok{\textasciitilde{}} \StringTok{""}\NormalTok{,}
\NormalTok{    p\_value }\SpecialCharTok{\textless{}} \FloatTok{0.01} \SpecialCharTok{\textasciitilde{}} \StringTok{"***"}\NormalTok{,}
\NormalTok{    p\_value }\SpecialCharTok{\textless{}} \FloatTok{0.05} \SpecialCharTok{\textasciitilde{}} \StringTok{"**"}\NormalTok{,}
\NormalTok{    p\_value }\SpecialCharTok{\textless{}} \FloatTok{0.10} \SpecialCharTok{\textasciitilde{}} \StringTok{"*"}\NormalTok{,}
    \ConstantTok{TRUE} \SpecialCharTok{\textasciitilde{}} \StringTok{""}
\NormalTok{  )}
\NormalTok{\}}
\end{Highlighting}
\end{Shaded}

\section{4.3
ポートフォリオ形成と市場リターン}\label{ux30ddux30fcux30c8ux30d5ux30a9ux30eaux30aaux5f62ux6210ux3068ux5e02ux5834ux30eaux30bfux30fcux30f3}

\subsection{4.3.1
前年度情報の準備}\label{ux524dux5e74ux5ea6ux60c5ux5831ux306eux6e96ux5099}

年度 \(t\) の月次リターンに対して、年度 \(t-1\)
末に観測された時価総額と株主資本を使用する。これにより、将来情報をポートフォリオ形成に利用することを防ぐ。

\begin{Shaded}
\begin{Highlighting}[]
\ControlFlowTok{if}\NormalTok{ (}\SpecialCharTok{!}\StringTok{"lagged\_ME"} \SpecialCharTok{\%in\%} \FunctionTok{names}\NormalTok{(tbl\_annual)) \{}
\NormalTok{  tbl\_annual }\OtherTok{\textless{}{-}}\NormalTok{ tbl\_annual }\SpecialCharTok{|\textgreater{}}
    \FunctionTok{group\_by}\NormalTok{(firm\_ID) }\SpecialCharTok{|\textgreater{}}
    \FunctionTok{arrange}\NormalTok{(year, }\AttributeTok{.by\_group =} \ConstantTok{TRUE}\NormalTok{) }\SpecialCharTok{|\textgreater{}}
    \FunctionTok{mutate}\NormalTok{(}\AttributeTok{lagged\_ME =} \FunctionTok{lag}\NormalTok{(ME)) }\SpecialCharTok{|\textgreater{}}
    \FunctionTok{ungroup}\NormalTok{()}
\NormalTok{\}}

\ControlFlowTok{if}\NormalTok{ (}\SpecialCharTok{!}\StringTok{"lagged\_BEME"} \SpecialCharTok{\%in\%} \FunctionTok{names}\NormalTok{(tbl\_annual)) \{}
\NormalTok{  tbl\_annual }\OtherTok{\textless{}{-}}\NormalTok{ tbl\_annual }\SpecialCharTok{|\textgreater{}}
    \FunctionTok{mutate}\NormalTok{(}\AttributeTok{lagged\_BEME =}\NormalTok{ lagged\_BE }\SpecialCharTok{/}\NormalTok{ lagged\_ME)}
\NormalTok{\}}

\NormalTok{tbl\_annual }\OtherTok{\textless{}{-}}\NormalTok{ tbl\_annual }\SpecialCharTok{|\textgreater{}}
  \FunctionTok{mutate}\NormalTok{(}
    \AttributeTok{lagged\_ME =} \FunctionTok{if\_else}\NormalTok{(}
      \FunctionTok{is.finite}\NormalTok{(lagged\_ME) }\SpecialCharTok{\&}\NormalTok{ lagged\_ME }\SpecialCharTok{\textgreater{}} \DecValTok{0}\NormalTok{,}
\NormalTok{      lagged\_ME,}
      \ConstantTok{NA\_real\_}
\NormalTok{    ),}
    \AttributeTok{lagged\_BEME =} \FunctionTok{if\_else}\NormalTok{(}
      \FunctionTok{is.finite}\NormalTok{(lagged\_BEME) }\SpecialCharTok{\&}\NormalTok{ lagged\_BEME }\SpecialCharTok{\textgreater{}} \DecValTok{0}\NormalTok{,}
\NormalTok{      lagged\_BEME,}
      \ConstantTok{NA\_real\_}
\NormalTok{    )}
\NormalTok{  ) }\SpecialCharTok{|\textgreater{}}
  \FunctionTok{arrange}\NormalTok{(firm\_ID, year)}

\NormalTok{tbl\_lagged\_variable\_availability }\OtherTok{\textless{}{-}}\NormalTok{ tbl\_annual }\SpecialCharTok{|\textgreater{}}
  \FunctionTok{group\_by}\NormalTok{(year) }\SpecialCharTok{|\textgreater{}}
  \FunctionTok{summarise}\NormalTok{(}
    \AttributeTok{firms =} \FunctionTok{n\_distinct}\NormalTok{(firm\_ID),}
    \AttributeTok{valid\_lagged\_ME =} \FunctionTok{sum}\NormalTok{(}\FunctionTok{is.finite}\NormalTok{(lagged\_ME)),}
    \AttributeTok{valid\_lagged\_BEME =} \FunctionTok{sum}\NormalTok{(}\FunctionTok{is.finite}\NormalTok{(lagged\_BEME)),}
    \AttributeTok{valid\_for\_2x3 =} \FunctionTok{sum}\NormalTok{(}
      \FunctionTok{is.finite}\NormalTok{(lagged\_ME) }\SpecialCharTok{\&} \FunctionTok{is.finite}\NormalTok{(lagged\_BEME)}
\NormalTok{    ),}
    \AttributeTok{.groups =} \StringTok{"drop"}
\NormalTok{  )}

\FunctionTok{show\_table\_LFL}\NormalTok{(tbl\_lagged\_variable\_availability, }\AttributeTok{n =} \ConstantTok{Inf}\NormalTok{)}
\end{Highlighting}
\end{Shaded}

\begin{verbatim}
# A tibble: 6 x 5
   year firms valid_lagged_ME valid_lagged_BEME valid_for_2x3
  <int> <int>           <int>             <int>         <int>
1  2015  1266               0                 0             0
2  2016  1293            1241              1241          1241
3  2017  1319            1270              1270          1270
4  2018  1323            1268              1268          1268
5  2019  1356            1304              1304          1304
6  2020  1363            1322              1322          1322
\end{verbatim}

\subsection{4.3.2
市場ポートフォリオ}\label{ux5e02ux5834ux30ddux30fcux30c8ux30d5ux30a9ux30eaux30aa}

年度 \(t\) の企業 \(i\)
の市場ポートフォリオ・ウェイトは、前年度末時価総額から計算する。

\[
w_{i,t}^{M}
=
\frac{ME_{i,t-1}}
{\sum_j ME_{j,t-1}}
\]

\begin{Shaded}
\begin{Highlighting}[]
\NormalTok{tbl\_market\_membership }\OtherTok{\textless{}{-}}\NormalTok{ tbl\_annual }\SpecialCharTok{|\textgreater{}}
  \FunctionTok{filter}\NormalTok{(}\FunctionTok{is.finite}\NormalTok{(lagged\_ME), lagged\_ME }\SpecialCharTok{\textgreater{}} \DecValTok{0}\NormalTok{) }\SpecialCharTok{|\textgreater{}}
  \FunctionTok{group\_by}\NormalTok{(year) }\SpecialCharTok{|\textgreater{}}
  \FunctionTok{mutate}\NormalTok{(}\AttributeTok{market\_weight =}\NormalTok{ lagged\_ME }\SpecialCharTok{/} \FunctionTok{sum}\NormalTok{(lagged\_ME)) }\SpecialCharTok{|\textgreater{}}
  \FunctionTok{ungroup}\NormalTok{() }\SpecialCharTok{|\textgreater{}}
  \FunctionTok{select}\NormalTok{(year, firm\_ID, lagged\_ME, market\_weight)}

\NormalTok{tbl\_market\_weight\_check }\OtherTok{\textless{}{-}}\NormalTok{ tbl\_market\_membership }\SpecialCharTok{|\textgreater{}}
  \FunctionTok{group\_by}\NormalTok{(year) }\SpecialCharTok{|\textgreater{}}
  \FunctionTok{summarise}\NormalTok{(}
    \AttributeTok{firms =} \FunctionTok{n\_distinct}\NormalTok{(firm\_ID),}
    \AttributeTok{weight\_sum =} \FunctionTok{sum}\NormalTok{(market\_weight),}
    \AttributeTok{minimum\_weight =} \FunctionTok{min}\NormalTok{(market\_weight),}
    \AttributeTok{maximum\_weight =} \FunctionTok{max}\NormalTok{(market\_weight),}
    \AttributeTok{.groups =} \StringTok{"drop"}
\NormalTok{  )}

\ControlFlowTok{if}\NormalTok{ (}\FunctionTok{any}\NormalTok{(}\FunctionTok{abs}\NormalTok{(tbl\_market\_weight\_check}\SpecialCharTok{$}\NormalTok{weight\_sum }\SpecialCharTok{{-}} \DecValTok{1}\NormalTok{) }\SpecialCharTok{\textgreater{}} \FloatTok{1e{-}10}\NormalTok{)) \{}
  \FunctionTok{stop}\NormalTok{(}\StringTok{"市場ポートフォリオの年度別ウェイト合計が1ではありません。"}\NormalTok{)}
\NormalTok{\}}

\FunctionTok{show\_table\_LFL}\NormalTok{(tbl\_market\_weight\_check, }\AttributeTok{n =} \ConstantTok{Inf}\NormalTok{)}
\end{Highlighting}
\end{Shaded}

\begin{verbatim}
# A tibble: 5 x 5
   year firms weight_sum minimum_weight maximum_weight
  <int> <int>      <dbl>          <dbl>          <dbl>
1  2016  1241          1  0.0000019019        0.040134
2  2017  1270          1  0.0000015083        0.032973
3  2018  1268          1  0.0000012609        0.042639
4  2019  1304          1  0.0000021887        0.039396
5  2020  1322          1  0.00000098548       0.056868
\end{verbatim}

市場ポートフォリオの月次リターンは、前年度末時価総額による加重平均である。

\[
R_{M,t}=\sum_i w^M_{i,t}R_{i,t}
\]

\begin{Shaded}
\begin{Highlighting}[]
\NormalTok{tbl\_market\_returns }\OtherTok{\textless{}{-}}\NormalTok{ tbl\_monthly }\SpecialCharTok{|\textgreater{}}
  \FunctionTok{inner\_join}\NormalTok{(}
\NormalTok{    tbl\_market\_membership,}
    \AttributeTok{by =} \FunctionTok{c}\NormalTok{(}\StringTok{"year"}\NormalTok{, }\StringTok{"firm\_ID"}\NormalTok{)}
\NormalTok{  ) }\SpecialCharTok{|\textgreater{}}
  \FunctionTok{group\_by}\NormalTok{(month\_ID, year, month, month\_date) }\SpecialCharTok{|\textgreater{}}
  \FunctionTok{summarise}\NormalTok{(}
    \AttributeTok{R\_F =} \FunctionTok{compute\_median\_safe}\NormalTok{(R\_F),}
    \AttributeTok{R\_M =} \FunctionTok{compute\_weighted\_mean\_safe}\NormalTok{(R, lagged\_ME),}
    \AttributeTok{eligible\_firms =} \FunctionTok{n\_distinct}\NormalTok{(firm\_ID),}
    \AttributeTok{valid\_return\_firms =} \FunctionTok{sum}\NormalTok{(}
      \FunctionTok{is.finite}\NormalTok{(R) }\SpecialCharTok{\&} \FunctionTok{is.finite}\NormalTok{(lagged\_ME) }\SpecialCharTok{\&}\NormalTok{ lagged\_ME }\SpecialCharTok{\textgreater{}} \DecValTok{0}
\NormalTok{    ),}
    \AttributeTok{.groups =} \StringTok{"drop"}
\NormalTok{  ) }\SpecialCharTok{|\textgreater{}}
  \FunctionTok{mutate}\NormalTok{(}
    \AttributeTok{R\_Me =}\NormalTok{ R\_M }\SpecialCharTok{{-}}\NormalTok{ R\_F,}
    \AttributeTok{return\_coverage =}\NormalTok{ valid\_return\_firms }\SpecialCharTok{/}\NormalTok{ eligible\_firms}
\NormalTok{  ) }\SpecialCharTok{|\textgreater{}}
  \FunctionTok{arrange}\NormalTok{(month\_ID)}

\NormalTok{tbl\_invalid\_market\_months }\OtherTok{\textless{}{-}}\NormalTok{ tbl\_market\_returns }\SpecialCharTok{|\textgreater{}}
  \FunctionTok{filter}\NormalTok{(}\SpecialCharTok{!}\FunctionTok{is.finite}\NormalTok{(R\_F) }\SpecialCharTok{|} \SpecialCharTok{!}\FunctionTok{is.finite}\NormalTok{(R\_M) }\SpecialCharTok{|} \SpecialCharTok{!}\FunctionTok{is.finite}\NormalTok{(R\_Me))}

\ControlFlowTok{if}\NormalTok{ (}\FunctionTok{nrow}\NormalTok{(tbl\_invalid\_market\_months) }\SpecialCharTok{\textgreater{}} \DecValTok{0}\NormalTok{) \{}
  \FunctionTok{stop}\NormalTok{(}\StringTok{"市場リターンを作成できない月があります。"}\NormalTok{)}
\NormalTok{\}}

\FunctionTok{show\_table\_LFL}\NormalTok{(tbl\_market\_returns)}
\end{Highlighting}
\end{Shaded}

\begin{verbatim}
# A tibble: 60 x 10
   month_ID  year month month_date         R_F        R_M eligible_firms
      <int> <int> <int> <date>           <dbl>      <dbl>          <int>
 1       13  2016     1 2016-01-01 0.000082029  0.045835            1241
 2       14  2016     2 2016-02-01 0.000075959  0.028423            1241
 3       15  2016     3 2016-03-01 0.000075959 -0.057976            1241
 4       16  2016     4 2016-04-01 0.000082029 -0.0010718           1241
 5       17  2016     5 2016-05-01 0.000082029 -0.0064160           1241
 6       18  2016     6 2016-06-01 0.000082029 -0.037075            1241
 7       19  2016     7 2016-07-01 0.000082029  0.029853            1241
 8       20  2016     8 2016-08-01 0.000082029  0.053757            1241
 9       21  2016     9 2016-09-01 0.000082029  0.023650            1241
10       22  2016    10 2016-10-01 0.000075959 -0.012053            1241
   valid_return_firms       R_Me return_coverage
                <int>      <dbl>           <dbl>
 1               1241  0.045753                1
 2               1241  0.028347                1
 3               1241 -0.058052                1
 4               1241 -0.0011539               1
 5               1241 -0.0064980               1
 6               1241 -0.037157                1
 7               1241  0.029771                1
 8               1241  0.053675                1
 9               1241  0.023568                1
10               1241 -0.012129                1
# i 50 more rows
\end{verbatim}

\begin{Shaded}
\begin{Highlighting}[]
\NormalTok{tbl\_market\_return\_summary }\OtherTok{\textless{}{-}}\NormalTok{ tbl\_market\_returns }\SpecialCharTok{|\textgreater{}}
  \FunctionTok{summarise}\NormalTok{(}
    \AttributeTok{observations =} \FunctionTok{n}\NormalTok{(),}
    \AttributeTok{mean\_market\_return =} \FunctionTok{mean}\NormalTok{(R\_M),}
    \AttributeTok{vec\_mean\_market\_excess\_return =} \FunctionTok{mean}\NormalTok{(R\_Me),}
    \AttributeTok{market\_volatility =} \FunctionTok{sd}\NormalTok{(R\_M),}
    \AttributeTok{minimum\_market\_return =} \FunctionTok{min}\NormalTok{(R\_M),}
    \AttributeTok{maximum\_market\_return =} \FunctionTok{max}\NormalTok{(R\_M),}
    \AttributeTok{minimum\_coverage =} \FunctionTok{min}\NormalTok{(return\_coverage)}
\NormalTok{  )}

\FunctionTok{show\_table\_LFL}\NormalTok{(tbl\_market\_return\_summary)}
\end{Highlighting}
\end{Shaded}

\begin{verbatim}
# A tibble: 1 x 7
  observations mean_market_return vec_mean_market_excess_return
         <int>              <dbl>                         <dbl>
1           60          0.0041003                     0.0040215
  market_volatility minimum_market_return maximum_market_return minimum_coverage
              <dbl>                 <dbl>                 <dbl>            <dbl>
1          0.041498              -0.10244               0.11104                1
\end{verbatim}

\begin{Shaded}
\begin{Highlighting}[]
\NormalTok{plot\_market\_excess\_return }\OtherTok{\textless{}{-}}\NormalTok{ tbl\_market\_returns }\SpecialCharTok{|\textgreater{}}
  \FunctionTok{ggplot}\NormalTok{(}\FunctionTok{aes}\NormalTok{(}\AttributeTok{x =}\NormalTok{ R\_Me)) }\SpecialCharTok{+}
  \FunctionTok{geom\_histogram}\NormalTok{(}\AttributeTok{bins =} \DecValTok{40}\NormalTok{) }\SpecialCharTok{+}
  \FunctionTok{geom\_vline}\NormalTok{(}\AttributeTok{xintercept =} \DecValTok{0}\NormalTok{, }\AttributeTok{linetype =} \StringTok{"dotted"}\NormalTok{) }\SpecialCharTok{+}
  \FunctionTok{labs}\NormalTok{(}
    \AttributeTok{x =} \StringTok{"Monthly Market Excess Return"}\NormalTok{,}
    \AttributeTok{y =} \StringTok{"Count"}\NormalTok{,}
    \AttributeTok{title =} \StringTok{"Distribution of Market Excess Returns"}
\NormalTok{  )}

\FunctionTok{show\_plot\_LFL}\NormalTok{(}\FunctionTok{add\_lagrange\_theme\_LFL}\NormalTok{(plot\_market\_excess\_return))}
\end{Highlighting}
\end{Shaded}

\pandocbounded{\includegraphics[keepaspectratio]{part1_empirical/ch04_factor_models_files/figure-pdf/plot-market-excess-return-1.pdf}}

\begin{Shaded}
\begin{Highlighting}[]
\NormalTok{plot\_cumulative\_market\_return }\OtherTok{\textless{}{-}}\NormalTok{ tbl\_market\_returns }\SpecialCharTok{|\textgreater{}}
  \FunctionTok{arrange}\NormalTok{(month\_ID) }\SpecialCharTok{|\textgreater{}}
  \FunctionTok{mutate}\NormalTok{(}\AttributeTok{cumulative\_gross\_market\_return =} \FunctionTok{cumprod}\NormalTok{(}\DecValTok{1} \SpecialCharTok{+}\NormalTok{ R\_M)) }\SpecialCharTok{|\textgreater{}}
  \FunctionTok{ggplot}\NormalTok{(}\FunctionTok{aes}\NormalTok{(}\AttributeTok{x =}\NormalTok{ month\_date, }\AttributeTok{y =}\NormalTok{ cumulative\_gross\_market\_return)) }\SpecialCharTok{+}
  \FunctionTok{geom\_line}\NormalTok{(}\AttributeTok{linewidth =} \FloatTok{0.8}\NormalTok{) }\SpecialCharTok{+}
  \FunctionTok{geom\_hline}\NormalTok{(}\AttributeTok{yintercept =} \DecValTok{1}\NormalTok{, }\AttributeTok{linetype =} \StringTok{"dotted"}\NormalTok{) }\SpecialCharTok{+}
  \FunctionTok{labs}\NormalTok{(}
    \AttributeTok{x =} \ConstantTok{NULL}\NormalTok{,}
    \AttributeTok{y =} \StringTok{"Cumulative Gross Return"}\NormalTok{,}
    \AttributeTok{title =} \StringTok{"Value{-}Weighted Market Portfolio"}
\NormalTok{  )}

\FunctionTok{show\_plot\_LFL}\NormalTok{(}\FunctionTok{add\_lagrange\_theme\_LFL}\NormalTok{(plot\_cumulative\_market\_return))}
\end{Highlighting}
\end{Shaded}

\pandocbounded{\includegraphics[keepaspectratio]{part1_empirical/ch04_factor_models_files/figure-pdf/plot-cumulative-market-return-1.pdf}}

\subsection{4.3.3
ポートフォリオ形成ユニバース}\label{ux30ddux30fcux30c8ux30d5ux30a9ux30eaux30aaux5f62ux6210ux30e6ux30cbux30d0ux30fcux30b9}

企業規模と簿価時価比率の両方を利用できる企業を共通の形成ユニバースとする。

\begin{Shaded}
\begin{Highlighting}[]
\NormalTok{tbl\_formation\_universe }\OtherTok{\textless{}{-}}\NormalTok{ tbl\_annual }\SpecialCharTok{|\textgreater{}}
  \FunctionTok{filter}\NormalTok{(}
    \FunctionTok{is.finite}\NormalTok{(lagged\_ME),}
\NormalTok{    lagged\_ME }\SpecialCharTok{\textgreater{}} \DecValTok{0}\NormalTok{,}
    \FunctionTok{is.finite}\NormalTok{(lagged\_BEME),}
\NormalTok{    lagged\_BEME }\SpecialCharTok{\textgreater{}} \DecValTok{0}
\NormalTok{  ) }\SpecialCharTok{|\textgreater{}}
  \FunctionTok{select}\NormalTok{(year, firm\_ID, lagged\_ME, lagged\_BEME)}

\NormalTok{tbl\_formation\_breakpoints }\OtherTok{\textless{}{-}}\NormalTok{ tbl\_formation\_universe }\SpecialCharTok{|\textgreater{}}
  \FunctionTok{group\_by}\NormalTok{(year) }\SpecialCharTok{|\textgreater{}}
  \FunctionTok{summarise}\NormalTok{(}
    \AttributeTok{firms =} \FunctionTok{n\_distinct}\NormalTok{(firm\_ID),}
    \AttributeTok{size\_median =} \FunctionTok{median}\NormalTok{(lagged\_ME),}
    \AttributeTok{BEME\_30 =} \FunctionTok{quantile}\NormalTok{(lagged\_BEME, }\AttributeTok{probs =} \FloatTok{0.30}\NormalTok{, }\AttributeTok{names =} \ConstantTok{FALSE}\NormalTok{),}
    \AttributeTok{BEME\_70 =} \FunctionTok{quantile}\NormalTok{(lagged\_BEME, }\AttributeTok{probs =} \FloatTok{0.70}\NormalTok{, }\AttributeTok{names =} \ConstantTok{FALSE}\NormalTok{),}
    \AttributeTok{.groups =} \StringTok{"drop"}
\NormalTok{  )}

\NormalTok{tbl\_formation }\OtherTok{\textless{}{-}}\NormalTok{ tbl\_formation\_universe }\SpecialCharTok{|\textgreater{}}
  \FunctionTok{left\_join}\NormalTok{(tbl\_formation\_breakpoints, }\AttributeTok{by =} \StringTok{"year"}\NormalTok{) }\SpecialCharTok{|\textgreater{}}
  \FunctionTok{group\_by}\NormalTok{(year) }\SpecialCharTok{|\textgreater{}}
  \FunctionTok{mutate}\NormalTok{(}
    \AttributeTok{ME\_rank10 =} \FunctionTok{ntile}\NormalTok{(lagged\_ME, }\DecValTok{10}\NormalTok{),}
    \AttributeTok{BEME\_rank10 =} \FunctionTok{ntile}\NormalTok{(lagged\_BEME, }\DecValTok{10}\NormalTok{)}
\NormalTok{  ) }\SpecialCharTok{|\textgreater{}}
  \FunctionTok{ungroup}\NormalTok{() }\SpecialCharTok{|\textgreater{}}
  \FunctionTok{mutate}\NormalTok{(}
    \AttributeTok{size\_group =} \FunctionTok{if\_else}\NormalTok{(lagged\_ME }\SpecialCharTok{\textless{}=}\NormalTok{ size\_median, }\StringTok{"S"}\NormalTok{, }\StringTok{"B"}\NormalTok{),}
    \AttributeTok{BEME\_group =} \FunctionTok{case\_when}\NormalTok{(}
\NormalTok{      lagged\_BEME }\SpecialCharTok{\textless{}=}\NormalTok{ BEME\_30 }\SpecialCharTok{\textasciitilde{}} \StringTok{"L"}\NormalTok{,}
\NormalTok{      lagged\_BEME }\SpecialCharTok{\textless{}=}\NormalTok{ BEME\_70 }\SpecialCharTok{\textasciitilde{}} \StringTok{"N"}\NormalTok{,}
      \ConstantTok{TRUE} \SpecialCharTok{\textasciitilde{}} \StringTok{"H"}
\NormalTok{    ),}
    \AttributeTok{FF\_portfolio\_type =} \FunctionTok{factor}\NormalTok{(}
      \FunctionTok{paste0}\NormalTok{(size\_group, BEME\_group),}
      \AttributeTok{levels =} \FunctionTok{c}\NormalTok{(}\StringTok{"SL"}\NormalTok{, }\StringTok{"SN"}\NormalTok{, }\StringTok{"SH"}\NormalTok{, }\StringTok{"BL"}\NormalTok{, }\StringTok{"BN"}\NormalTok{, }\StringTok{"BH"}\NormalTok{)}
\NormalTok{    )}
\NormalTok{  ) }\SpecialCharTok{|\textgreater{}}
  \FunctionTok{select}\NormalTok{(}
\NormalTok{    year, firm\_ID, lagged\_ME, lagged\_BEME,}
\NormalTok{    ME\_rank10, BEME\_rank10,}
\NormalTok{    size\_group, BEME\_group, FF\_portfolio\_type}
\NormalTok{  )}

\NormalTok{tbl\_portfolio\_membership\_counts }\OtherTok{\textless{}{-}}\NormalTok{ tbl\_formation }\SpecialCharTok{|\textgreater{}}
  \FunctionTok{count}\NormalTok{(year, FF\_portfolio\_type, }\AttributeTok{name =} \StringTok{"firms"}\NormalTok{) }\SpecialCharTok{|\textgreater{}}
  \FunctionTok{complete}\NormalTok{(}
\NormalTok{    year,}
\NormalTok{    FF\_portfolio\_type,}
    \AttributeTok{fill =} \FunctionTok{list}\NormalTok{(}\AttributeTok{firms =} \DecValTok{0}\NormalTok{)}
\NormalTok{  ) }\SpecialCharTok{|\textgreater{}}
  \FunctionTok{arrange}\NormalTok{(year, FF\_portfolio\_type)}

\NormalTok{tbl\_empty\_portfolios }\OtherTok{\textless{}{-}}\NormalTok{ tbl\_portfolio\_membership\_counts }\SpecialCharTok{|\textgreater{}}
  \FunctionTok{filter}\NormalTok{(firms }\SpecialCharTok{==} \DecValTok{0}\NormalTok{)}

\ControlFlowTok{if}\NormalTok{ (}\FunctionTok{nrow}\NormalTok{(tbl\_empty\_portfolios) }\SpecialCharTok{\textgreater{}} \DecValTok{0}\NormalTok{) \{}
  \FunctionTok{stop}\NormalTok{(}\StringTok{"企業が割り当てられていない2×3ポートフォリオがあります。"}\NormalTok{)}
\NormalTok{\}}

\FunctionTok{show\_table\_LFL}\NormalTok{(tbl\_formation\_breakpoints, }\AttributeTok{n =} \ConstantTok{Inf}\NormalTok{)}
\end{Highlighting}
\end{Shaded}

\begin{verbatim}
# A tibble: 5 x 5
   year firms size_median       BEME_30      BEME_70
  <int> <int>       <dbl>         <dbl>        <dbl>
1  2016  1241 24564654000 0.00000076171 0.0000014304
2  2017  1270 25133528500 0.00000070949 0.0000013738
3  2018  1268 28268839000 0.00000061518 0.0000012038
4  2019  1304 23853596500 0.00000073364 0.0000014045
5  2020  1322 31878381000 0.00000057440 0.0000011304
\end{verbatim}

\begin{Shaded}
\begin{Highlighting}[]
\FunctionTok{show\_table\_LFL}\NormalTok{(tbl\_portfolio\_membership\_counts, }\AttributeTok{n =} \ConstantTok{Inf}\NormalTok{)}
\end{Highlighting}
\end{Shaded}

\begin{verbatim}
# A tibble: 30 x 3
    year FF_portfolio_type firms
   <int> <fct>             <int>
 1  2016 SL                  140
 2  2016 SN                  210
 3  2016 SH                  271
 4  2016 BL                  233
 5  2016 BN                  286
 6  2016 BH                  101
 7  2017 SL                  141
 8  2017 SN                  224
 9  2017 SH                  270
10  2017 BL                  240
11  2017 BN                  284
12  2017 BH                  111
13  2018 SL                  152
14  2018 SN                  224
15  2018 SH                  258
16  2018 BL                  229
17  2018 BN                  282
18  2018 BH                  123
19  2019 SL                  167
20  2019 SN                  225
21  2019 SH                  260
22  2019 BL                  224
23  2019 BN                  297
24  2019 BH                  131
25  2020 SL                  173
26  2020 SN                  248
27  2020 SH                  240
28  2020 BL                  224
29  2020 BN                  280
30  2020 BH                  157
\end{verbatim}

\section{4.4
単変量ポートフォリオ・ソート}\label{ux5358ux5909ux91cfux30ddux30fcux30c8ux30d5ux30a9ux30eaux30aaux30bdux30fcux30c8}

\subsection{4.4.1
企業規模10分位ポートフォリオ}\label{ux4f01ux696dux898fux6a2110ux5206ux4f4dux30ddux30fcux30c8ux30d5ux30a9ux30eaux30aa}

企業規模10分位ごとに、等加重リターンと前年度末時価総額加重リターンを計算する。

\begin{Shaded}
\begin{Highlighting}[]
\NormalTok{tbl\_size\_decile\_returns }\OtherTok{\textless{}{-}}\NormalTok{ tbl\_monthly }\SpecialCharTok{|\textgreater{}}
  \FunctionTok{inner\_join}\NormalTok{(}
\NormalTok{    tbl\_formation }\SpecialCharTok{|\textgreater{}}
      \FunctionTok{select}\NormalTok{(year, firm\_ID, lagged\_ME, ME\_rank10),}
    \AttributeTok{by =} \FunctionTok{c}\NormalTok{(}\StringTok{"year"}\NormalTok{, }\StringTok{"firm\_ID"}\NormalTok{)}
\NormalTok{  ) }\SpecialCharTok{|\textgreater{}}
  \FunctionTok{group\_by}\NormalTok{(}
\NormalTok{    month\_ID, year, month, month\_date, ME\_rank10}
\NormalTok{  ) }\SpecialCharTok{|\textgreater{}}
  \FunctionTok{summarise}\NormalTok{(}
    \AttributeTok{R\_F =} \FunctionTok{compute\_median\_safe}\NormalTok{(R\_F),}
    \AttributeTok{R\_equal =} \FunctionTok{compute\_mean\_safe}\NormalTok{(R),}
    \AttributeTok{R\_value =} \FunctionTok{compute\_weighted\_mean\_safe}\NormalTok{(R, lagged\_ME),}
    \AttributeTok{firms =} \FunctionTok{n\_distinct}\NormalTok{(firm\_ID),}
    \AttributeTok{valid\_return\_firms =} \FunctionTok{sum}\NormalTok{(}\FunctionTok{is.finite}\NormalTok{(R)),}
    \AttributeTok{.groups =} \StringTok{"drop"}
\NormalTok{  ) }\SpecialCharTok{|\textgreater{}}
  \FunctionTok{mutate}\NormalTok{(}
    \AttributeTok{Re\_equal =}\NormalTok{ R\_equal }\SpecialCharTok{{-}}\NormalTok{ R\_F,}
    \AttributeTok{Re\_value =}\NormalTok{ R\_value }\SpecialCharTok{{-}}\NormalTok{ R\_F}
\NormalTok{  ) }\SpecialCharTok{|\textgreater{}}
  \FunctionTok{arrange}\NormalTok{(month\_ID, ME\_rank10)}

\NormalTok{tbl\_size\_decile\_coverage }\OtherTok{\textless{}{-}}\NormalTok{ tbl\_size\_decile\_returns }\SpecialCharTok{|\textgreater{}}
  \FunctionTok{group\_by}\NormalTok{(month\_ID) }\SpecialCharTok{|\textgreater{}}
  \FunctionTok{summarise}\NormalTok{(}
    \AttributeTok{portfolios =} \FunctionTok{n\_distinct}\NormalTok{(ME\_rank10),}
    \AttributeTok{.groups =} \StringTok{"drop"}
\NormalTok{  ) }\SpecialCharTok{|\textgreater{}}
  \FunctionTok{filter}\NormalTok{(portfolios }\SpecialCharTok{!=} \DecValTok{10}\NormalTok{)}

\ControlFlowTok{if}\NormalTok{ (}\FunctionTok{nrow}\NormalTok{(tbl\_size\_decile\_coverage) }\SpecialCharTok{\textgreater{}} \DecValTok{0}\NormalTok{) \{}
  \FunctionTok{stop}\NormalTok{(}\StringTok{"10個の企業規模ポートフォリオがそろわない月があります。"}\NormalTok{)}
\NormalTok{\}}

\NormalTok{tbl\_size\_decile\_summary }\OtherTok{\textless{}{-}}\NormalTok{ tbl\_size\_decile\_returns }\SpecialCharTok{|\textgreater{}}
  \FunctionTok{group\_by}\NormalTok{(ME\_rank10) }\SpecialCharTok{|\textgreater{}}
  \FunctionTok{summarise}\NormalTok{(}
    \AttributeTok{months =} \FunctionTok{n}\NormalTok{(),}
    \AttributeTok{mean\_equal\_excess\_return =} \FunctionTok{mean}\NormalTok{(Re\_equal, }\AttributeTok{na.rm =} \ConstantTok{TRUE}\NormalTok{),}
    \AttributeTok{mean\_value\_excess\_return =} \FunctionTok{mean}\NormalTok{(Re\_value, }\AttributeTok{na.rm =} \ConstantTok{TRUE}\NormalTok{),}
    \AttributeTok{value\_return\_volatility =} \FunctionTok{sd}\NormalTok{(Re\_value, }\AttributeTok{na.rm =} \ConstantTok{TRUE}\NormalTok{),}
    \AttributeTok{.groups =} \StringTok{"drop"}
\NormalTok{  ) }\SpecialCharTok{|\textgreater{}}
  \FunctionTok{arrange}\NormalTok{(ME\_rank10)}

\FunctionTok{show\_table\_LFL}\NormalTok{(tbl\_size\_decile\_summary, }\AttributeTok{n =} \ConstantTok{Inf}\NormalTok{)}
\end{Highlighting}
\end{Shaded}

\begin{verbatim}
# A tibble: 10 x 5
   ME_rank10 months mean_equal_excess_return mean_value_excess_return
       <int>  <int>                    <dbl>                    <dbl>
 1         1     60                0.014931                 0.015119 
 2         2     60                0.013581                 0.013512 
 3         3     60                0.015247                 0.015316 
 4         4     60                0.013073                 0.012958 
 5         5     60                0.010989                 0.010910 
 6         6     60                0.010277                 0.010388 
 7         7     60                0.0067439                0.0065837
 8         8     60                0.0051647                0.0054430
 9         9     60                0.0046434                0.0045278
10        10     60                0.0037215                0.0029811
   value_return_volatility
                     <dbl>
 1                0.042085
 2                0.040856
 3                0.040102
 4                0.041626
 5                0.041165
 6                0.040414
 7                0.041338
 8                0.041428
 9                0.044503
10                0.042314
\end{verbatim}

\begin{Shaded}
\begin{Highlighting}[]
\NormalTok{plot\_size\_decile }\OtherTok{\textless{}{-}}\NormalTok{ tbl\_size\_decile\_summary }\SpecialCharTok{|\textgreater{}}
  \FunctionTok{select}\NormalTok{(}
\NormalTok{    ME\_rank10,}
    \AttributeTok{equal\_weighted =}\NormalTok{ mean\_equal\_excess\_return,}
    \AttributeTok{value\_weighted =}\NormalTok{ mean\_value\_excess\_return}
\NormalTok{  ) }\SpecialCharTok{|\textgreater{}}
  \FunctionTok{pivot\_longer}\NormalTok{(}
    \AttributeTok{cols =} \FunctionTok{c}\NormalTok{(equal\_weighted, value\_weighted),}
    \AttributeTok{names\_to =} \StringTok{"weighting"}\NormalTok{,}
    \AttributeTok{values\_to =} \StringTok{"mean\_excess\_return"}
\NormalTok{  ) }\SpecialCharTok{|\textgreater{}}
  \FunctionTok{mutate}\NormalTok{(}\AttributeTok{ME\_rank10 =} \FunctionTok{factor}\NormalTok{(ME\_rank10, }\AttributeTok{levels =} \DecValTok{1}\SpecialCharTok{:}\DecValTok{10}\NormalTok{)) }\SpecialCharTok{|\textgreater{}}
  \FunctionTok{ggplot}\NormalTok{(}
    \FunctionTok{aes}\NormalTok{(}
      \AttributeTok{x =}\NormalTok{ ME\_rank10,}
      \AttributeTok{y =}\NormalTok{ mean\_excess\_return,}
      \AttributeTok{fill =}\NormalTok{ weighting}
\NormalTok{    )}
\NormalTok{  ) }\SpecialCharTok{+}
  \FunctionTok{geom\_col}\NormalTok{(}\AttributeTok{position =} \StringTok{"dodge"}\NormalTok{) }\SpecialCharTok{+}
  \FunctionTok{geom\_hline}\NormalTok{(}\AttributeTok{yintercept =} \DecValTok{0}\NormalTok{, }\AttributeTok{linetype =} \StringTok{"dotted"}\NormalTok{) }\SpecialCharTok{+}
  \FunctionTok{scale\_y\_continuous}\NormalTok{(}\AttributeTok{labels =}\NormalTok{ scales}\SpecialCharTok{::}\FunctionTok{label\_percent}\NormalTok{()) }\SpecialCharTok{+}
  \FunctionTok{labs}\NormalTok{(}
    \AttributeTok{x =} \StringTok{"Size Decile"}\NormalTok{,}
    \AttributeTok{y =} \StringTok{"Mean Monthly Excess Return"}\NormalTok{,}
    \AttributeTok{fill =} \StringTok{"Weighting"}\NormalTok{,}
    \AttributeTok{title =} \StringTok{"Size{-}Sorted Portfolio Returns"}
\NormalTok{  )}

\FunctionTok{show\_plot\_LFL}\NormalTok{(}\FunctionTok{add\_lagrange\_theme\_LFL}\NormalTok{(plot\_size\_decile))}
\end{Highlighting}
\end{Shaded}

\pandocbounded{\includegraphics[keepaspectratio]{part1_empirical/ch04_factor_models_files/figure-pdf/plot-size-decile-returns-1.pdf}}

\subsection{4.4.2
簿価時価比率10分位ポートフォリオ}\label{ux7c3fux4fa1ux6642ux4fa1ux6bd4ux738710ux5206ux4f4dux30ddux30fcux30c8ux30d5ux30a9ux30eaux30aa}

\begin{Shaded}
\begin{Highlighting}[]
\NormalTok{tbl\_beme\_decile\_returns }\OtherTok{\textless{}{-}}\NormalTok{ tbl\_monthly }\SpecialCharTok{|\textgreater{}}
  \FunctionTok{inner\_join}\NormalTok{(}
\NormalTok{    tbl\_formation }\SpecialCharTok{|\textgreater{}}
      \FunctionTok{select}\NormalTok{(year, firm\_ID, lagged\_ME, BEME\_rank10),}
    \AttributeTok{by =} \FunctionTok{c}\NormalTok{(}\StringTok{"year"}\NormalTok{, }\StringTok{"firm\_ID"}\NormalTok{)}
\NormalTok{  ) }\SpecialCharTok{|\textgreater{}}
  \FunctionTok{group\_by}\NormalTok{(}
\NormalTok{    month\_ID, year, month, month\_date, BEME\_rank10}
\NormalTok{  ) }\SpecialCharTok{|\textgreater{}}
  \FunctionTok{summarise}\NormalTok{(}
    \AttributeTok{R\_F =} \FunctionTok{compute\_median\_safe}\NormalTok{(R\_F),}
    \AttributeTok{R\_equal =} \FunctionTok{compute\_mean\_safe}\NormalTok{(R),}
    \AttributeTok{R\_value =} \FunctionTok{compute\_weighted\_mean\_safe}\NormalTok{(R, lagged\_ME),}
    \AttributeTok{firms =} \FunctionTok{n\_distinct}\NormalTok{(firm\_ID),}
    \AttributeTok{valid\_return\_firms =} \FunctionTok{sum}\NormalTok{(}\FunctionTok{is.finite}\NormalTok{(R)),}
    \AttributeTok{.groups =} \StringTok{"drop"}
\NormalTok{  ) }\SpecialCharTok{|\textgreater{}}
  \FunctionTok{mutate}\NormalTok{(}
    \AttributeTok{Re\_equal =}\NormalTok{ R\_equal }\SpecialCharTok{{-}}\NormalTok{ R\_F,}
    \AttributeTok{Re\_value =}\NormalTok{ R\_value }\SpecialCharTok{{-}}\NormalTok{ R\_F}
\NormalTok{  ) }\SpecialCharTok{|\textgreater{}}
  \FunctionTok{arrange}\NormalTok{(month\_ID, BEME\_rank10)}

\NormalTok{tbl\_beme\_decile\_summary }\OtherTok{\textless{}{-}}\NormalTok{ tbl\_beme\_decile\_returns }\SpecialCharTok{|\textgreater{}}
  \FunctionTok{group\_by}\NormalTok{(BEME\_rank10) }\SpecialCharTok{|\textgreater{}}
  \FunctionTok{summarise}\NormalTok{(}
    \AttributeTok{months =} \FunctionTok{n}\NormalTok{(),}
    \AttributeTok{mean\_equal\_excess\_return =} \FunctionTok{mean}\NormalTok{(Re\_equal, }\AttributeTok{na.rm =} \ConstantTok{TRUE}\NormalTok{),}
    \AttributeTok{mean\_value\_excess\_return =} \FunctionTok{mean}\NormalTok{(Re\_value, }\AttributeTok{na.rm =} \ConstantTok{TRUE}\NormalTok{),}
    \AttributeTok{value\_return\_volatility =} \FunctionTok{sd}\NormalTok{(Re\_value, }\AttributeTok{na.rm =} \ConstantTok{TRUE}\NormalTok{),}
    \AttributeTok{.groups =} \StringTok{"drop"}
\NormalTok{  ) }\SpecialCharTok{|\textgreater{}}
  \FunctionTok{arrange}\NormalTok{(BEME\_rank10)}

\FunctionTok{show\_table\_LFL}\NormalTok{(tbl\_beme\_decile\_summary, }\AttributeTok{n =} \ConstantTok{Inf}\NormalTok{)}
\end{Highlighting}
\end{Shaded}

\begin{verbatim}
# A tibble: 10 x 5
   BEME_rank10 months mean_equal_excess_return mean_value_excess_return
         <int>  <int>                    <dbl>                    <dbl>
 1           1     60                0.0069080               -0.0012722
 2           2     60                0.0057059                0.0024925
 3           3     60                0.0043601                0.0021706
 4           4     60                0.0090970                0.0034693
 5           5     60                0.0099654                0.0085384
 6           6     60                0.010433                 0.0055909
 7           7     60                0.0086709                0.0085607
 8           8     60                0.012979                 0.0088343
 9           9     60                0.016017                 0.014657 
10          10     60                0.014389                 0.0097273
   value_return_volatility
                     <dbl>
 1                0.052048
 2                0.039802
 3                0.045255
 4                0.045192
 5                0.047651
 6                0.042627
 7                0.044756
 8                0.044844
 9                0.049615
10                0.043590
\end{verbatim}

\begin{Shaded}
\begin{Highlighting}[]
\NormalTok{plot\_beme\_decile }\OtherTok{\textless{}{-}}\NormalTok{ tbl\_beme\_decile\_summary }\SpecialCharTok{|\textgreater{}}
  \FunctionTok{select}\NormalTok{(}
\NormalTok{    BEME\_rank10,}
    \AttributeTok{equal\_weighted =}\NormalTok{ mean\_equal\_excess\_return,}
    \AttributeTok{value\_weighted =}\NormalTok{ mean\_value\_excess\_return}
\NormalTok{  ) }\SpecialCharTok{|\textgreater{}}
  \FunctionTok{pivot\_longer}\NormalTok{(}
    \AttributeTok{cols =} \FunctionTok{c}\NormalTok{(equal\_weighted, value\_weighted),}
    \AttributeTok{names\_to =} \StringTok{"weighting"}\NormalTok{,}
    \AttributeTok{values\_to =} \StringTok{"mean\_excess\_return"}
\NormalTok{  ) }\SpecialCharTok{|\textgreater{}}
  \FunctionTok{mutate}\NormalTok{(}\AttributeTok{BEME\_rank10 =} \FunctionTok{factor}\NormalTok{(BEME\_rank10, }\AttributeTok{levels =} \DecValTok{1}\SpecialCharTok{:}\DecValTok{10}\NormalTok{)) }\SpecialCharTok{|\textgreater{}}
  \FunctionTok{ggplot}\NormalTok{(}
    \FunctionTok{aes}\NormalTok{(}
      \AttributeTok{x =}\NormalTok{ BEME\_rank10,}
      \AttributeTok{y =}\NormalTok{ mean\_excess\_return,}
      \AttributeTok{fill =}\NormalTok{ weighting}
\NormalTok{    )}
\NormalTok{  ) }\SpecialCharTok{+}
  \FunctionTok{geom\_col}\NormalTok{(}\AttributeTok{position =} \StringTok{"dodge"}\NormalTok{) }\SpecialCharTok{+}
  \FunctionTok{geom\_hline}\NormalTok{(}\AttributeTok{yintercept =} \DecValTok{0}\NormalTok{, }\AttributeTok{linetype =} \StringTok{"dotted"}\NormalTok{) }\SpecialCharTok{+}
  \FunctionTok{scale\_y\_continuous}\NormalTok{(}\AttributeTok{labels =}\NormalTok{ scales}\SpecialCharTok{::}\FunctionTok{label\_percent}\NormalTok{()) }\SpecialCharTok{+}
  \FunctionTok{labs}\NormalTok{(}
    \AttributeTok{x =} \StringTok{"Book{-}to{-}Market Decile"}\NormalTok{,}
    \AttributeTok{y =} \StringTok{"Mean Monthly Excess Return"}\NormalTok{,}
    \AttributeTok{fill =} \StringTok{"Weighting"}\NormalTok{,}
    \AttributeTok{title =} \StringTok{"Book{-}to{-}Market{-}Sorted Portfolio Returns"}
\NormalTok{  )}

\FunctionTok{show\_plot\_LFL}\NormalTok{(}\FunctionTok{add\_lagrange\_theme\_LFL}\NormalTok{(plot\_beme\_decile))}
\end{Highlighting}
\end{Shaded}

\pandocbounded{\includegraphics[keepaspectratio]{part1_empirical/ch04_factor_models_files/figure-pdf/plot-beme-decile-returns-1.pdf}}

\section{4.5 CAPM}\label{capm}

\subsection{4.5.1
モデルの推定}\label{ux30e2ux30c7ux30ebux306eux63a8ux5b9a}

CAPMでは、ポートフォリオ \(p\)
の超過リターンを市場超過リターンで説明する。

\[
R^e_{p,t}
=
\alpha_p+\beta_pR^e_{M,t}+\varepsilon_{p,t}
\]

\begin{Shaded}
\begin{Highlighting}[]
\NormalTok{tbl\_capm\_data }\OtherTok{\textless{}{-}}\NormalTok{ tbl\_size\_decile\_returns }\SpecialCharTok{|\textgreater{}}
  \FunctionTok{select}\NormalTok{(month\_ID, ME\_rank10, }\AttributeTok{Re =}\NormalTok{ Re\_value) }\SpecialCharTok{|\textgreater{}}
  \FunctionTok{inner\_join}\NormalTok{(}
\NormalTok{    tbl\_market\_returns }\SpecialCharTok{|\textgreater{}}
      \FunctionTok{select}\NormalTok{(month\_ID, R\_Me),}
    \AttributeTok{by =} \StringTok{"month\_ID"}
\NormalTok{  ) }\SpecialCharTok{|\textgreater{}}
  \FunctionTok{filter}\NormalTok{(}\FunctionTok{is.finite}\NormalTok{(Re), }\FunctionTok{is.finite}\NormalTok{(R\_Me)) }\SpecialCharTok{|\textgreater{}}
  \FunctionTok{arrange}\NormalTok{(ME\_rank10, month\_ID)}

\NormalTok{tbl\_capm\_models }\OtherTok{\textless{}{-}}\NormalTok{ tbl\_capm\_data }\SpecialCharTok{|\textgreater{}}
  \FunctionTok{group\_by}\NormalTok{(ME\_rank10) }\SpecialCharTok{|\textgreater{}}
  \FunctionTok{nest}\NormalTok{() }\SpecialCharTok{|\textgreater{}}
  \FunctionTok{mutate}\NormalTok{(}
    \AttributeTok{model =} \FunctionTok{map}\NormalTok{(data, \textbackslash{}(data) }\FunctionTok{lm}\NormalTok{(Re }\SpecialCharTok{\textasciitilde{}}\NormalTok{ R\_Me, }\AttributeTok{data =}\NormalTok{ data)),}
    \AttributeTok{coefficients =} \FunctionTok{map}\NormalTok{(model, broom}\SpecialCharTok{::}\NormalTok{tidy),}
    \AttributeTok{fit =} \FunctionTok{map}\NormalTok{(model, broom}\SpecialCharTok{::}\NormalTok{glance)}
\NormalTok{  )}

\NormalTok{tbl\_capm\_results }\OtherTok{\textless{}{-}}\NormalTok{ tbl\_capm\_models }\SpecialCharTok{|\textgreater{}}
  \FunctionTok{select}\NormalTok{(ME\_rank10, coefficients) }\SpecialCharTok{|\textgreater{}}
  \FunctionTok{unnest}\NormalTok{(coefficients) }\SpecialCharTok{|\textgreater{}}
  \FunctionTok{mutate}\NormalTok{(}\AttributeTok{significance =} \FunctionTok{compute\_significance\_stars}\NormalTok{(p.value)) }\SpecialCharTok{|\textgreater{}}
  \FunctionTok{arrange}\NormalTok{(ME\_rank10, term)}

\NormalTok{tbl\_capm\_fit }\OtherTok{\textless{}{-}}\NormalTok{ tbl\_capm\_models }\SpecialCharTok{|\textgreater{}}
  \FunctionTok{select}\NormalTok{(ME\_rank10, fit) }\SpecialCharTok{|\textgreater{}}
  \FunctionTok{unnest}\NormalTok{(fit) }\SpecialCharTok{|\textgreater{}}
  \FunctionTok{select}\NormalTok{(}
\NormalTok{    ME\_rank10, r.squared, adj.r.squared,}
\NormalTok{    sigma, statistic, p.value, nobs}
\NormalTok{  ) }\SpecialCharTok{|\textgreater{}}
  \FunctionTok{arrange}\NormalTok{(ME\_rank10)}

\NormalTok{tbl\_capm\_alpha }\OtherTok{\textless{}{-}}\NormalTok{ tbl\_capm\_results }\SpecialCharTok{|\textgreater{}}
  \FunctionTok{filter}\NormalTok{(term }\SpecialCharTok{==} \StringTok{"(Intercept)"}\NormalTok{) }\SpecialCharTok{|\textgreater{}}
  \FunctionTok{transmute}\NormalTok{(}
\NormalTok{    ME\_rank10,}
    \AttributeTok{tbl\_capm\_alpha =}\NormalTok{ estimate,}
    \AttributeTok{standard\_error =}\NormalTok{ std.error,}
\NormalTok{    statistic,}
    \AttributeTok{p\_value =}\NormalTok{ p.value,}
\NormalTok{    significance}
\NormalTok{  )}

\FunctionTok{show\_table\_LFL}\NormalTok{(tbl\_capm\_results, }\AttributeTok{n =} \ConstantTok{Inf}\NormalTok{)}
\end{Highlighting}
\end{Shaded}

\begin{verbatim}
# A tibble: 20 x 7
# Groups:   ME_rank10 [10]
   ME_rank10 term           estimate  std.error statistic    p.value
       <int> <chr>             <dbl>      <dbl>     <dbl>      <dbl>
 1         1 (Intercept)  0.012412   0.0041175    3.0144  3.8151e- 3
 2         1 R_Me         0.67325    0.099584     6.7607  7.3884e- 9
 3         2 (Intercept)  0.010641   0.0036808    2.8910  5.3962e- 3
 4         2 R_Me         0.71390    0.089021     8.0195  5.6813e-11
 5         3 (Intercept)  0.012195   0.0031256    3.9016  2.5116e- 4
 6         3 R_Me         0.77614    0.075595    10.267   1.1583e-14
 7         4 (Intercept)  0.0095370  0.0028845    3.3062  1.6262e- 3
 8         4 R_Me         0.85081    0.069764    12.196   1.2013e-17
 9         5 (Intercept)  0.0073101  0.0023212    3.1493  2.5865e- 3
10         5 R_Me         0.89511    0.056138    15.945   7.2596e-23
11         6 (Intercept)  0.0067526  0.0019672    3.4325  1.1093e- 3
12         6 R_Me         0.90394    0.047579    18.999   1.3994e-26
13         7 (Intercept)  0.0027911  0.0017412    1.6030  1.1437e- 1
14         7 R_Me         0.94309    0.042112    22.395   3.1166e-30
15         8 (Intercept)  0.0016328  0.0017077    0.95613 3.4298e- 1
16         8 R_Me         0.94747    0.041302    22.940   8.8761e-31
17         9 (Intercept)  0.00035036 0.0014465    0.24221 8.0947e- 1
18         9 R_Me         1.0388     0.034985    29.693   9.0177e-37
19        10 (Intercept) -0.0010957  0.00059445  -1.8432  7.0415e- 2
20        10 R_Me         1.0138     0.014377    70.512   6.5289e-58
   significance
   <chr>       
 1 "***"       
 2 "***"       
 3 "***"       
 4 "***"       
 5 "***"       
 6 "***"       
 7 "***"       
 8 "***"       
 9 "***"       
10 "***"       
11 "***"       
12 "***"       
13 ""          
14 "***"       
15 ""          
16 "***"       
17 ""          
18 "***"       
19 "*"         
20 "***"       
\end{verbatim}

\begin{Shaded}
\begin{Highlighting}[]
\FunctionTok{show\_table\_LFL}\NormalTok{(tbl\_capm\_fit, }\AttributeTok{n =} \ConstantTok{Inf}\NormalTok{)}
\end{Highlighting}
\end{Shaded}

\begin{verbatim}
# A tibble: 10 x 7
# Groups:   ME_rank10 [10]
   ME_rank10 r.squared adj.r.squared     sigma statistic    p.value  nobs
       <int>     <dbl>         <dbl>     <dbl>     <dbl>      <dbl> <int>
 1         1   0.44073       0.43109 0.031743     45.706 7.3884e- 9    60
 2         2   0.52580       0.51763 0.028376     64.312 5.6813e-11    60
 3         3   0.64507       0.63895 0.024096    105.41  1.1583e-14    60
 4         4   0.71944       0.71461 0.022238    148.73  1.2013e-17    60
 5         5   0.81424       0.81104 0.017894    254.23  7.2596e-23    60
 6         6   0.86156       0.85918 0.015166    360.96  1.3994e-26    60
 7         7   0.89634       0.89456 0.013423    501.54  3.1166e-30    60
 8         8   0.90073       0.89902 0.013165    526.25  8.8761e-31    60
 9         9   0.93827       0.93721 0.011152    881.65  9.0177e-37    60
10        10   0.98847       0.98827 0.0045827  4972.0   6.5289e-58    60
\end{verbatim}

\begin{Shaded}
\begin{Highlighting}[]
\NormalTok{plot\_capm\_alpha }\OtherTok{\textless{}{-}}\NormalTok{ tbl\_capm\_alpha }\SpecialCharTok{|\textgreater{}}
  \FunctionTok{mutate}\NormalTok{(}\AttributeTok{ME\_rank10 =} \FunctionTok{factor}\NormalTok{(ME\_rank10, }\AttributeTok{levels =} \DecValTok{1}\SpecialCharTok{:}\DecValTok{10}\NormalTok{)) }\SpecialCharTok{|\textgreater{}}
  \FunctionTok{ggplot}\NormalTok{(}\FunctionTok{aes}\NormalTok{(}\AttributeTok{x =}\NormalTok{ ME\_rank10, }\AttributeTok{y =}\NormalTok{ tbl\_capm\_alpha)) }\SpecialCharTok{+}
  \FunctionTok{geom\_col}\NormalTok{() }\SpecialCharTok{+}
  \FunctionTok{geom\_hline}\NormalTok{(}\AttributeTok{yintercept =} \DecValTok{0}\NormalTok{, }\AttributeTok{linetype =} \StringTok{"dotted"}\NormalTok{) }\SpecialCharTok{+}
  \FunctionTok{scale\_y\_continuous}\NormalTok{(}\AttributeTok{labels =}\NormalTok{ scales}\SpecialCharTok{::}\FunctionTok{label\_percent}\NormalTok{()) }\SpecialCharTok{+}
  \FunctionTok{labs}\NormalTok{(}
    \AttributeTok{x =} \StringTok{"Size Decile"}\NormalTok{,}
    \AttributeTok{y =} \StringTok{"CAPM Alpha"}\NormalTok{,}
    \AttributeTok{title =} \StringTok{"CAPM Alphas of Size Portfolios"}
\NormalTok{  )}

\FunctionTok{show\_plot\_LFL}\NormalTok{(}\FunctionTok{add\_lagrange\_theme\_LFL}\NormalTok{(plot\_capm\_alpha))}
\end{Highlighting}
\end{Shaded}

\pandocbounded{\includegraphics[keepaspectratio]{part1_empirical/ch04_factor_models_files/figure-pdf/plot-capm-alpha-1.pdf}}

\subsection{4.5.2 証券市場線}\label{ux8a3cux5238ux5e02ux5834ux7dda}

\begin{Shaded}
\begin{Highlighting}[]
\NormalTok{tbl\_capm\_beta }\OtherTok{\textless{}{-}}\NormalTok{ tbl\_capm\_results }\SpecialCharTok{|\textgreater{}}
  \FunctionTok{filter}\NormalTok{(term }\SpecialCharTok{==} \StringTok{"R\_Me"}\NormalTok{) }\SpecialCharTok{|\textgreater{}}
  \FunctionTok{transmute}\NormalTok{(}
\NormalTok{    ME\_rank10,}
    \AttributeTok{tbl\_capm\_beta =}\NormalTok{ estimate,}
    \AttributeTok{beta\_p\_value =}\NormalTok{ p.value}
\NormalTok{  )}

\NormalTok{tbl\_security\_market\_line }\OtherTok{\textless{}{-}}\NormalTok{ tbl\_size\_decile\_summary }\SpecialCharTok{|\textgreater{}}
  \FunctionTok{select}\NormalTok{(}
\NormalTok{    ME\_rank10,}
    \AttributeTok{mean\_excess\_return =}\NormalTok{ mean\_value\_excess\_return}
\NormalTok{  ) }\SpecialCharTok{|\textgreater{}}
  \FunctionTok{left\_join}\NormalTok{(tbl\_capm\_beta, }\AttributeTok{by =} \StringTok{"ME\_rank10"}\NormalTok{)}

\NormalTok{vec\_mean\_market\_excess\_return }\OtherTok{\textless{}{-}} \FunctionTok{mean}\NormalTok{(}
\NormalTok{  tbl\_market\_returns}\SpecialCharTok{$}\NormalTok{R\_Me,}
  \AttributeTok{na.rm =} \ConstantTok{TRUE}
\NormalTok{)}

\NormalTok{plot\_security\_market\_line }\OtherTok{\textless{}{-}}\NormalTok{ tbl\_security\_market\_line }\SpecialCharTok{|\textgreater{}}
  \FunctionTok{ggplot}\NormalTok{(}
    \FunctionTok{aes}\NormalTok{(}
      \AttributeTok{x =}\NormalTok{ tbl\_capm\_beta,}
      \AttributeTok{y =}\NormalTok{ mean\_excess\_return}
\NormalTok{    )}
\NormalTok{  ) }\SpecialCharTok{+}
  \FunctionTok{geom\_point}\NormalTok{() }\SpecialCharTok{+}
  \FunctionTok{geom\_text}\NormalTok{(}\FunctionTok{aes}\NormalTok{(}\AttributeTok{label =}\NormalTok{ ME\_rank10), }\AttributeTok{nudge\_y =} \FloatTok{0.0005}\NormalTok{) }\SpecialCharTok{+}
  \FunctionTok{geom\_abline}\NormalTok{(}
    \AttributeTok{intercept =} \DecValTok{0}\NormalTok{,}
    \AttributeTok{slope =}\NormalTok{ vec\_mean\_market\_excess\_return,}
    \AttributeTok{linetype =} \StringTok{"dashed"}
\NormalTok{  ) }\SpecialCharTok{+}
  \FunctionTok{scale\_y\_continuous}\NormalTok{(}\AttributeTok{labels =}\NormalTok{ scales}\SpecialCharTok{::}\FunctionTok{label\_percent}\NormalTok{()) }\SpecialCharTok{+}
  \FunctionTok{labs}\NormalTok{(}
    \AttributeTok{x =} \StringTok{"Market Beta"}\NormalTok{,}
    \AttributeTok{y =} \StringTok{"Mean Monthly Excess Return"}\NormalTok{,}
    \AttributeTok{title =} \StringTok{"Security Market Line"}
\NormalTok{  )}

\FunctionTok{show\_plot\_LFL}\NormalTok{(}\FunctionTok{add\_lagrange\_theme\_LFL}\NormalTok{(plot\_security\_market\_line))}
\end{Highlighting}
\end{Shaded}

\pandocbounded{\includegraphics[keepaspectratio]{part1_empirical/ch04_factor_models_files/figure-pdf/plot-security-market-line-1.pdf}}

\section{4.6
2×3ポートフォリオとSMB・HML}\label{ux30ddux30fcux30c8ux30d5ux30a9ux30eaux30aaux3068smbhml}

\subsection{4.6.1
Size--BE/MEポートフォリオ}\label{sizebemeux30ddux30fcux30c8ux30d5ux30a9ux30eaux30aa}

企業規模をSmallとBigの2群、簿価時価比率をLow、Neutral、Highの3群に分ける。

{\def\LTcaptype{none} % do not increment counter
\begin{longtable}[]{@{}lll@{}}
\toprule\noalign{}
ポートフォリオ & 企業規模 & 簿価時価比率 \\
\midrule\noalign{}
\endhead
\bottomrule\noalign{}
\endlastfoot
SL & Small & Low \\
SN & Small & Neutral \\
SH & Small & High \\
BL & Big & Low \\
BN & Big & Neutral \\
BH & Big & High \\
\end{longtable}
}

\begin{Shaded}
\begin{Highlighting}[]
\NormalTok{tbl\_ff\_portfolio\_returns }\OtherTok{\textless{}{-}}\NormalTok{ tbl\_monthly }\SpecialCharTok{|\textgreater{}}
  \FunctionTok{inner\_join}\NormalTok{(}
\NormalTok{    tbl\_formation }\SpecialCharTok{|\textgreater{}}
      \FunctionTok{select}\NormalTok{(}
\NormalTok{        year, firm\_ID, lagged\_ME,}
\NormalTok{        size\_group, BEME\_group, FF\_portfolio\_type}
\NormalTok{      ),}
    \AttributeTok{by =} \FunctionTok{c}\NormalTok{(}\StringTok{"year"}\NormalTok{, }\StringTok{"firm\_ID"}\NormalTok{)}
\NormalTok{  ) }\SpecialCharTok{|\textgreater{}}
  \FunctionTok{filter}\NormalTok{(}\SpecialCharTok{!}\FunctionTok{is.na}\NormalTok{(FF\_portfolio\_type)) }\SpecialCharTok{|\textgreater{}}
  \FunctionTok{group\_by}\NormalTok{(}
\NormalTok{    month\_ID, year, month, month\_date,}
\NormalTok{    FF\_portfolio\_type, size\_group, BEME\_group}
\NormalTok{  ) }\SpecialCharTok{|\textgreater{}}
  \FunctionTok{summarise}\NormalTok{(}
    \AttributeTok{R\_F =} \FunctionTok{compute\_median\_safe}\NormalTok{(R\_F),}
    \AttributeTok{R =} \FunctionTok{compute\_weighted\_mean\_safe}\NormalTok{(R, lagged\_ME),}
    \AttributeTok{firms =} \FunctionTok{n\_distinct}\NormalTok{(firm\_ID),}
    \AttributeTok{valid\_return\_firms =} \FunctionTok{sum}\NormalTok{(}\FunctionTok{is.finite}\NormalTok{(R)),}
    \AttributeTok{.groups =} \StringTok{"drop"}
\NormalTok{  ) }\SpecialCharTok{|\textgreater{}}
  \FunctionTok{mutate}\NormalTok{(}\AttributeTok{Re =}\NormalTok{ R }\SpecialCharTok{{-}}\NormalTok{ R\_F) }\SpecialCharTok{|\textgreater{}}
  \FunctionTok{arrange}\NormalTok{(month\_ID, FF\_portfolio\_type)}

\NormalTok{tbl\_ff\_portfolio\_coverage }\OtherTok{\textless{}{-}}\NormalTok{ tbl\_ff\_portfolio\_returns }\SpecialCharTok{|\textgreater{}}
  \FunctionTok{group\_by}\NormalTok{(month\_ID) }\SpecialCharTok{|\textgreater{}}
  \FunctionTok{summarise}\NormalTok{(}
    \AttributeTok{portfolios =} \FunctionTok{n\_distinct}\NormalTok{(FF\_portfolio\_type),}
    \AttributeTok{.groups =} \StringTok{"drop"}
\NormalTok{  ) }\SpecialCharTok{|\textgreater{}}
  \FunctionTok{filter}\NormalTok{(portfolios }\SpecialCharTok{!=} \DecValTok{6}\NormalTok{)}

\ControlFlowTok{if}\NormalTok{ (}\FunctionTok{nrow}\NormalTok{(tbl\_ff\_portfolio\_coverage) }\SpecialCharTok{\textgreater{}} \DecValTok{0}\NormalTok{) \{}
  \FunctionTok{stop}\NormalTok{(}\StringTok{"6個のSize{-}{-}BE/MEポートフォリオがそろわない月があります。"}\NormalTok{)}
\NormalTok{\}}

\NormalTok{tbl\_ff\_portfolio\_summary }\OtherTok{\textless{}{-}}\NormalTok{ tbl\_ff\_portfolio\_returns }\SpecialCharTok{|\textgreater{}}
  \FunctionTok{group\_by}\NormalTok{(}
\NormalTok{    FF\_portfolio\_type,}
\NormalTok{    size\_group,}
\NormalTok{    BEME\_group}
\NormalTok{  ) }\SpecialCharTok{|\textgreater{}}
  \FunctionTok{summarise}\NormalTok{(}
    \AttributeTok{months =} \FunctionTok{n}\NormalTok{(),}
    \AttributeTok{mean\_return =} \FunctionTok{mean}\NormalTok{(R),}
    \AttributeTok{mean\_excess\_return =} \FunctionTok{mean}\NormalTok{(Re),}
    \AttributeTok{volatility =} \FunctionTok{sd}\NormalTok{(R),}
    \AttributeTok{.groups =} \StringTok{"drop"}
\NormalTok{  ) }\SpecialCharTok{|\textgreater{}}
  \FunctionTok{arrange}\NormalTok{(FF\_portfolio\_type)}

\FunctionTok{show\_table\_LFL}\NormalTok{(tbl\_ff\_portfolio\_summary, }\AttributeTok{n =} \ConstantTok{Inf}\NormalTok{)}
\end{Highlighting}
\end{Shaded}

\begin{verbatim}
# A tibble: 6 x 7
  FF_portfolio_type size_group BEME_group months mean_return mean_excess_return
  <fct>             <chr>      <chr>       <int>       <dbl>              <dbl>
1 SL                S          L              60  0.010056           0.0099773 
2 SN                S          N              60  0.013474           0.013395  
3 SH                S          H              60  0.014002           0.013923  
4 BL                B          L              60  0.00081235         0.00073351
5 BN                B          N              60  0.0063080          0.0062291 
6 BH                B          H              60  0.010766           0.010687  
  volatility
       <dbl>
1   0.044267
2   0.041252
3   0.040505
4   0.043049
5   0.042621
6   0.046846
\end{verbatim}

\begin{Shaded}
\begin{Highlighting}[]
\NormalTok{plot\_ff\_portfolio }\OtherTok{\textless{}{-}}\NormalTok{ tbl\_ff\_portfolio\_summary }\SpecialCharTok{|\textgreater{}}
  \FunctionTok{ggplot}\NormalTok{(}
    \FunctionTok{aes}\NormalTok{(}
      \AttributeTok{x =}\NormalTok{ BEME\_group,}
      \AttributeTok{y =}\NormalTok{ mean\_excess\_return,}
      \AttributeTok{fill =}\NormalTok{ size\_group}
\NormalTok{    )}
\NormalTok{  ) }\SpecialCharTok{+}
  \FunctionTok{geom\_col}\NormalTok{(}\AttributeTok{position =} \StringTok{"dodge"}\NormalTok{) }\SpecialCharTok{+}
  \FunctionTok{geom\_text}\NormalTok{(}
    \FunctionTok{aes}\NormalTok{(}\AttributeTok{label =}\NormalTok{ FF\_portfolio\_type, }\AttributeTok{group =}\NormalTok{ size\_group),}
    \AttributeTok{position =} \FunctionTok{position\_dodge}\NormalTok{(}\AttributeTok{width =} \FloatTok{0.9}\NormalTok{),}
    \AttributeTok{vjust =} \SpecialCharTok{{-}}\FloatTok{0.5}
\NormalTok{  ) }\SpecialCharTok{+}
  \FunctionTok{geom\_hline}\NormalTok{(}\AttributeTok{yintercept =} \DecValTok{0}\NormalTok{, }\AttributeTok{linetype =} \StringTok{"dotted"}\NormalTok{) }\SpecialCharTok{+}
  \FunctionTok{scale\_y\_continuous}\NormalTok{(}
    \AttributeTok{labels =}\NormalTok{ scales}\SpecialCharTok{::}\FunctionTok{label\_percent}\NormalTok{(),}
    \AttributeTok{expand =} \FunctionTok{expansion}\NormalTok{(}\AttributeTok{mult =} \FunctionTok{c}\NormalTok{(}\FloatTok{0.05}\NormalTok{, }\FloatTok{0.15}\NormalTok{))}
\NormalTok{  ) }\SpecialCharTok{+}
  \FunctionTok{labs}\NormalTok{(}
    \AttributeTok{x =} \StringTok{"Book{-}to{-}Market Group"}\NormalTok{,}
    \AttributeTok{y =} \StringTok{"Mean Monthly Excess Return"}\NormalTok{,}
    \AttributeTok{fill =} \StringTok{"Size Group"}\NormalTok{,}
    \AttributeTok{title =} \StringTok{"Six Size{-}{-}Book{-}to{-}Market Portfolios"}
\NormalTok{  )}

\FunctionTok{show\_plot\_LFL}\NormalTok{(}\FunctionTok{add\_lagrange\_theme\_LFL}\NormalTok{(plot\_ff\_portfolio))}
\end{Highlighting}
\end{Shaded}

\pandocbounded{\includegraphics[keepaspectratio]{part1_empirical/ch04_factor_models_files/figure-pdf/plot-ff-portfolio-returns-1.pdf}}

\begin{Shaded}
\begin{Highlighting}[]
\NormalTok{plot\_ff\_cumulative }\OtherTok{\textless{}{-}}\NormalTok{ tbl\_ff\_portfolio\_returns }\SpecialCharTok{|\textgreater{}}
  \FunctionTok{group\_by}\NormalTok{(FF\_portfolio\_type) }\SpecialCharTok{|\textgreater{}}
  \FunctionTok{arrange}\NormalTok{(month\_ID, }\AttributeTok{.by\_group =} \ConstantTok{TRUE}\NormalTok{) }\SpecialCharTok{|\textgreater{}}
  \FunctionTok{mutate}\NormalTok{(}\AttributeTok{cumulative\_gross\_return =} \FunctionTok{cumprod}\NormalTok{(}\DecValTok{1} \SpecialCharTok{+}\NormalTok{ R)) }\SpecialCharTok{|\textgreater{}}
  \FunctionTok{ungroup}\NormalTok{() }\SpecialCharTok{|\textgreater{}}
  \FunctionTok{ggplot}\NormalTok{(}
    \FunctionTok{aes}\NormalTok{(}
      \AttributeTok{x =}\NormalTok{ month\_date,}
      \AttributeTok{y =}\NormalTok{ cumulative\_gross\_return,}
      \AttributeTok{linetype =}\NormalTok{ FF\_portfolio\_type}
\NormalTok{    )}
\NormalTok{  ) }\SpecialCharTok{+}
  \FunctionTok{geom\_line}\NormalTok{(}\AttributeTok{linewidth =} \FloatTok{0.7}\NormalTok{) }\SpecialCharTok{+}
  \FunctionTok{geom\_hline}\NormalTok{(}\AttributeTok{yintercept =} \DecValTok{1}\NormalTok{, }\AttributeTok{linetype =} \StringTok{"dotted"}\NormalTok{) }\SpecialCharTok{+}
  \FunctionTok{labs}\NormalTok{(}
    \AttributeTok{x =} \ConstantTok{NULL}\NormalTok{,}
    \AttributeTok{y =} \StringTok{"Cumulative Gross Return"}\NormalTok{,}
    \AttributeTok{linetype =} \StringTok{"Portfolio"}\NormalTok{,}
    \AttributeTok{title =} \StringTok{"Cumulative Returns of Six Portfolios"}
\NormalTok{  )}

\FunctionTok{show\_plot\_LFL}\NormalTok{(}\FunctionTok{add\_lagrange\_theme\_LFL}\NormalTok{(plot\_ff\_cumulative))}
\end{Highlighting}
\end{Shaded}

\pandocbounded{\includegraphics[keepaspectratio]{part1_empirical/ch04_factor_models_files/figure-pdf/plot-ff-cumulative-return-1.pdf}}

\subsection{4.6.2 SMBとHML}\label{smbux3068hml}

SMBはSmallポートフォリオの平均リターンからBigポートフォリオの平均リターンを引いて計算する。

\[
SMB_t
=
\frac{SL_t+SN_t+SH_t}{3}
-
\frac{BL_t+BN_t+BH_t}{3}
\]

HMLはHighポートフォリオの平均リターンからLowポートフォリオの平均リターンを引いて計算する。

\[
HML_t
=
\frac{SH_t+BH_t}{2}
-
\frac{SL_t+BL_t}{2}
\]

\begin{Shaded}
\begin{Highlighting}[]
\NormalTok{tbl\_ff\_portfolio\_wide }\OtherTok{\textless{}{-}}\NormalTok{ tbl\_ff\_portfolio\_returns }\SpecialCharTok{|\textgreater{}}
  \FunctionTok{select}\NormalTok{(}
\NormalTok{    month\_ID, year, month, month\_date,}
\NormalTok{    FF\_portfolio\_type, R}
\NormalTok{  ) }\SpecialCharTok{|\textgreater{}}
  \FunctionTok{pivot\_wider}\NormalTok{(}
    \AttributeTok{names\_from =}\NormalTok{ FF\_portfolio\_type,}
    \AttributeTok{values\_from =}\NormalTok{ R}
\NormalTok{  ) }\SpecialCharTok{|\textgreater{}}
  \FunctionTok{arrange}\NormalTok{(month\_ID)}

\NormalTok{tbl\_smb\_hml }\OtherTok{\textless{}{-}}\NormalTok{ tbl\_ff\_portfolio\_wide }\SpecialCharTok{|\textgreater{}}
  \FunctionTok{mutate}\NormalTok{(}
    \AttributeTok{SMB =}\NormalTok{ (SL }\SpecialCharTok{+}\NormalTok{ SN }\SpecialCharTok{+}\NormalTok{ SH) }\SpecialCharTok{/} \DecValTok{3} \SpecialCharTok{{-}}\NormalTok{ (BL }\SpecialCharTok{+}\NormalTok{ BN }\SpecialCharTok{+}\NormalTok{ BH) }\SpecialCharTok{/} \DecValTok{3}\NormalTok{,}
    \AttributeTok{HML =}\NormalTok{ (SH }\SpecialCharTok{+}\NormalTok{ BH) }\SpecialCharTok{/} \DecValTok{2} \SpecialCharTok{{-}}\NormalTok{ (SL }\SpecialCharTok{+}\NormalTok{ BL) }\SpecialCharTok{/} \DecValTok{2}
\NormalTok{  ) }\SpecialCharTok{|\textgreater{}}
  \FunctionTok{select}\NormalTok{(}
\NormalTok{    month\_ID, year, month, month\_date,}
\NormalTok{    SL, SN, SH, BL, BN, BH, SMB, HML}
\NormalTok{  )}

\NormalTok{tbl\_factor }\OtherTok{\textless{}{-}}\NormalTok{ tbl\_market\_returns }\SpecialCharTok{|\textgreater{}}
  \FunctionTok{select}\NormalTok{(}
\NormalTok{    month\_ID, year, month, month\_date,}
\NormalTok{    R\_F, R\_M, R\_Me,}
    \AttributeTok{market\_return\_coverage =}\NormalTok{ return\_coverage}
\NormalTok{  ) }\SpecialCharTok{|\textgreater{}}
  \FunctionTok{inner\_join}\NormalTok{(}
\NormalTok{    tbl\_smb\_hml,}
    \AttributeTok{by =} \FunctionTok{c}\NormalTok{(}\StringTok{"month\_ID"}\NormalTok{, }\StringTok{"year"}\NormalTok{, }\StringTok{"month"}\NormalTok{, }\StringTok{"month\_date"}\NormalTok{)}
\NormalTok{  ) }\SpecialCharTok{|\textgreater{}}
  \FunctionTok{arrange}\NormalTok{(month\_ID)}

\NormalTok{tbl\_invalid\_factor\_rows }\OtherTok{\textless{}{-}}\NormalTok{ tbl\_factor }\SpecialCharTok{|\textgreater{}}
  \FunctionTok{filter}\NormalTok{(}
    \SpecialCharTok{!}\FunctionTok{if\_all}\NormalTok{(}
      \FunctionTok{c}\NormalTok{(R\_F, R\_M, R\_Me, SMB, HML),}
\NormalTok{      is.finite}
\NormalTok{    )}
\NormalTok{  )}

\ControlFlowTok{if}\NormalTok{ (}\FunctionTok{nrow}\NormalTok{(tbl\_invalid\_factor\_rows) }\SpecialCharTok{\textgreater{}} \DecValTok{0}\NormalTok{) \{}
  \FunctionTok{stop}\NormalTok{(}\StringTok{"ファクター・データに欠損値または非有限値があります。"}\NormalTok{)}
\NormalTok{\}}

\FunctionTok{show\_table\_LFL}\NormalTok{(tbl\_factor)}
\end{Highlighting}
\end{Shaded}

\begin{verbatim}
# A tibble: 60 x 16
   month_ID  year month month_date         R_F        R_M       R_Me
      <int> <int> <int> <date>           <dbl>      <dbl>      <dbl>
 1       13  2016     1 2016-01-01 0.000082029  0.045835   0.045753 
 2       14  2016     2 2016-02-01 0.000075959  0.028423   0.028347 
 3       15  2016     3 2016-03-01 0.000075959 -0.057976  -0.058052 
 4       16  2016     4 2016-04-01 0.000082029 -0.0010718 -0.0011539
 5       17  2016     5 2016-05-01 0.000082029 -0.0064160 -0.0064980
 6       18  2016     6 2016-06-01 0.000082029 -0.037075  -0.037157 
 7       19  2016     7 2016-07-01 0.000082029  0.029853   0.029771 
 8       20  2016     8 2016-08-01 0.000082029  0.053757   0.053675 
 9       21  2016     9 2016-09-01 0.000082029  0.023650   0.023568 
10       22  2016    10 2016-10-01 0.000075959 -0.012053  -0.012129 
   market_return_coverage        SL         SN          SH         BL         BN
                    <dbl>     <dbl>      <dbl>       <dbl>      <dbl>      <dbl>
 1                      1  0.034424  0.040738   0.049300    0.048792   0.043035 
 2                      1  0.016158  0.039043   0.026503    0.018793   0.043267 
 3                      1 -0.011443 -0.0075343 -0.00082137 -0.065067  -0.058172 
 4                      1 -0.050894 -0.029672  -0.0084329  -0.0096066  0.013125 
 5                      1 -0.050419 -0.020829  -0.013882   -0.0064587 -0.0031242
 6                      1 -0.010094 -0.021100  -0.019579   -0.045272  -0.025807 
 7                      1  0.016136  0.0030885  0.0024072   0.037354   0.027623 
 8                      1  0.061823  0.043624   0.040794    0.060068   0.044383 
 9                      1  0.076890  0.056521   0.059240    0.023654   0.018573 
10                      1  0.014770  0.011308   0.0017272  -0.013614  -0.010340 
           BH        SMB        HML
        <dbl>      <dbl>      <dbl>
 1  0.037868  -0.0017441  0.0019763
 2  0.039974  -0.0067772  0.015763 
 3 -0.035989   0.046477   0.019850 
 4  0.016129  -0.036215   0.034098 
 5 -0.0090124 -0.022178   0.016992 
 6 -0.030949   0.017085   0.0024194
 7  0.0011262 -0.014824  -0.024978 
 8  0.048092  -0.0021008 -0.016503 
 9  0.023533   0.042297  -0.0088858
10 -0.016441   0.022733  -0.0079348
# i 50 more rows
\end{verbatim}

\subsection{4.6.3
ファクターの要約統計量}\label{ux30d5ux30a1ux30afux30bfux30fcux306eux8981ux7d04ux7d71ux8a08ux91cf}

\begin{Shaded}
\begin{Highlighting}[]
\NormalTok{vec\_factor\_names }\OtherTok{\textless{}{-}} \FunctionTok{c}\NormalTok{(}\StringTok{"R\_Me"}\NormalTok{, }\StringTok{"SMB"}\NormalTok{, }\StringTok{"HML"}\NormalTok{)}

\NormalTok{tbl\_factor\_summary }\OtherTok{\textless{}{-}}\NormalTok{ vec\_factor\_names }\SpecialCharTok{|\textgreater{}}
  \FunctionTok{map\_dfr}\NormalTok{(}
\NormalTok{    \textbackslash{}(factor\_name) \{}
\NormalTok{      values }\OtherTok{\textless{}{-}}\NormalTok{ tbl\_factor[[factor\_name]]}
\NormalTok{      finite\_values }\OtherTok{\textless{}{-}}\NormalTok{ values[}\FunctionTok{is.finite}\NormalTok{(values)]}

      \FunctionTok{tibble}\NormalTok{(}
        \AttributeTok{factor =}\NormalTok{ factor\_name,}
        \AttributeTok{observations =} \FunctionTok{length}\NormalTok{(finite\_values),}
        \AttributeTok{mean =} \FunctionTok{mean}\NormalTok{(finite\_values),}
        \AttributeTok{standard\_deviation =} \FunctionTok{sd}\NormalTok{(finite\_values),}
        \AttributeTok{minimum =} \FunctionTok{min}\NormalTok{(finite\_values),}
        \AttributeTok{median =} \FunctionTok{median}\NormalTok{(finite\_values),}
        \AttributeTok{maximum =} \FunctionTok{max}\NormalTok{(finite\_values),}
        \AttributeTok{t\_value =} \FunctionTok{mean}\NormalTok{(finite\_values) }\SpecialCharTok{/}
\NormalTok{          (}\FunctionTok{sd}\NormalTok{(finite\_values) }\SpecialCharTok{/} \FunctionTok{sqrt}\NormalTok{(}\FunctionTok{length}\NormalTok{(finite\_values)))}
\NormalTok{      )}
\NormalTok{    \}}
\NormalTok{  )}

\NormalTok{tbl\_factor\_correlation }\OtherTok{\textless{}{-}}\NormalTok{ tbl\_factor }\SpecialCharTok{|\textgreater{}}
  \FunctionTok{select}\NormalTok{(R\_Me, SMB, HML) }\SpecialCharTok{|\textgreater{}}
  \FunctionTok{cor}\NormalTok{(}\AttributeTok{use =} \StringTok{"complete.obs"}\NormalTok{)}

\FunctionTok{show\_table\_LFL}\NormalTok{(tbl\_factor\_summary, }\AttributeTok{n =} \ConstantTok{Inf}\NormalTok{)}
\end{Highlighting}
\end{Shaded}

\begin{verbatim}
# A tibble: 3 x 8
  factor observations      mean standard_deviation   minimum    median  maximum
  <chr>         <int>     <dbl>              <dbl>     <dbl>     <dbl>    <dbl>
1 R_Me             60 0.0040215           0.041498 -0.10252  0.0060055 0.11097 
2 SMB              60 0.0065488           0.022857 -0.038447 0.0055499 0.069390
3 HML              60 0.0069496           0.018058 -0.043572 0.0089560 0.035442
  t_value
    <dbl>
1 0.75064
2 2.2193 
3 2.9811 
\end{verbatim}

\begin{Shaded}
\begin{Highlighting}[]
\NormalTok{tbl\_factor\_correlation}
\end{Highlighting}
\end{Shaded}

\begin{verbatim}
           R_Me        SMB        HML
R_Me  1.0000000 -0.3431883  0.1702499
SMB  -0.3431883  1.0000000 -0.3806883
HML   0.1702499 -0.3806883  1.0000000
\end{verbatim}

\begin{Shaded}
\begin{Highlighting}[]
\NormalTok{plot\_factor\_time\_series }\OtherTok{\textless{}{-}}\NormalTok{ tbl\_factor }\SpecialCharTok{|\textgreater{}}
  \FunctionTok{select}\NormalTok{(month\_date, R\_Me, SMB, HML) }\SpecialCharTok{|\textgreater{}}
  \FunctionTok{pivot\_longer}\NormalTok{(}
    \AttributeTok{cols =} \FunctionTok{c}\NormalTok{(R\_Me, SMB, HML),}
    \AttributeTok{names\_to =} \StringTok{"factor"}\NormalTok{,}
    \AttributeTok{values\_to =} \StringTok{"return"}
\NormalTok{  ) }\SpecialCharTok{|\textgreater{}}
  \FunctionTok{ggplot}\NormalTok{(}\FunctionTok{aes}\NormalTok{(}\AttributeTok{x =}\NormalTok{ month\_date, }\AttributeTok{y =}\NormalTok{ return)) }\SpecialCharTok{+}
  \FunctionTok{geom\_line}\NormalTok{(}\AttributeTok{linewidth =} \FloatTok{0.6}\NormalTok{) }\SpecialCharTok{+}
  \FunctionTok{geom\_hline}\NormalTok{(}\AttributeTok{yintercept =} \DecValTok{0}\NormalTok{, }\AttributeTok{linetype =} \StringTok{"dotted"}\NormalTok{) }\SpecialCharTok{+}
  \FunctionTok{facet\_wrap}\NormalTok{(}
    \FunctionTok{vars}\NormalTok{(factor),}
    \AttributeTok{ncol =} \DecValTok{1}\NormalTok{,}
    \AttributeTok{scales =} \StringTok{"free\_y"}
\NormalTok{  ) }\SpecialCharTok{+}
  \FunctionTok{scale\_y\_continuous}\NormalTok{(}\AttributeTok{labels =}\NormalTok{ scales}\SpecialCharTok{::}\FunctionTok{label\_percent}\NormalTok{()) }\SpecialCharTok{+}
  \FunctionTok{labs}\NormalTok{(}
    \AttributeTok{x =} \ConstantTok{NULL}\NormalTok{,}
    \AttributeTok{y =} \StringTok{"Monthly Return"}\NormalTok{,}
    \AttributeTok{title =} \StringTok{"Market, SMB and HML Factors"}
\NormalTok{  )}

\FunctionTok{show\_plot\_LFL}\NormalTok{(}\FunctionTok{add\_lagrange\_theme\_LFL}\NormalTok{(plot\_factor\_time\_series))}
\end{Highlighting}
\end{Shaded}

\pandocbounded{\includegraphics[keepaspectratio]{part1_empirical/ch04_factor_models_files/figure-pdf/plot-factor-time-series-1.pdf}}

\section{4.7
FF3モデルとCAPMの比較}\label{ff3ux30e2ux30c7ux30ebux3068capmux306eux6bd4ux8f03}

\subsection{4.7.1
FF3モデルの推定}\label{ff3ux30e2ux30c7ux30ebux306eux63a8ux5b9a}

Fama--Frenchの3ファクター・モデルは、ポートフォリオの超過リターンを市場、SMB、HMLで説明する。

\[
R^e_{p,t}
=
\alpha_p
+\beta_{M,p}R^e_{M,t}
+\beta_{SMB,p}SMB_t
+\beta_{HML,p}HML_t
+\varepsilon_{p,t}
\]

\begin{Shaded}
\begin{Highlighting}[]
\NormalTok{tbl\_ff3\_data }\OtherTok{\textless{}{-}}\NormalTok{ tbl\_size\_decile\_returns }\SpecialCharTok{|\textgreater{}}
  \FunctionTok{select}\NormalTok{(month\_ID, ME\_rank10, }\AttributeTok{Re =}\NormalTok{ Re\_value) }\SpecialCharTok{|\textgreater{}}
  \FunctionTok{inner\_join}\NormalTok{(}
\NormalTok{    tbl\_factor }\SpecialCharTok{|\textgreater{}}
      \FunctionTok{select}\NormalTok{(month\_ID, R\_Me, SMB, HML),}
    \AttributeTok{by =} \StringTok{"month\_ID"}
\NormalTok{  ) }\SpecialCharTok{|\textgreater{}}
  \FunctionTok{filter}\NormalTok{(}
    \FunctionTok{if\_all}\NormalTok{(}
      \FunctionTok{c}\NormalTok{(Re, R\_Me, SMB, HML),}
\NormalTok{      is.finite}
\NormalTok{    )}
\NormalTok{  ) }\SpecialCharTok{|\textgreater{}}
  \FunctionTok{arrange}\NormalTok{(ME\_rank10, month\_ID)}

\NormalTok{tbl\_ff3\_models }\OtherTok{\textless{}{-}}\NormalTok{ tbl\_ff3\_data }\SpecialCharTok{|\textgreater{}}
  \FunctionTok{group\_by}\NormalTok{(ME\_rank10) }\SpecialCharTok{|\textgreater{}}
  \FunctionTok{nest}\NormalTok{() }\SpecialCharTok{|\textgreater{}}
  \FunctionTok{mutate}\NormalTok{(}
    \AttributeTok{model =} \FunctionTok{map}\NormalTok{(}
\NormalTok{      data,}
\NormalTok{      \textbackslash{}(data) }\FunctionTok{lm}\NormalTok{(Re }\SpecialCharTok{\textasciitilde{}}\NormalTok{ R\_Me }\SpecialCharTok{+}\NormalTok{ SMB }\SpecialCharTok{+}\NormalTok{ HML, }\AttributeTok{data =}\NormalTok{ data)}
\NormalTok{    ),}
    \AttributeTok{coefficients =} \FunctionTok{map}\NormalTok{(model, broom}\SpecialCharTok{::}\NormalTok{tidy),}
    \AttributeTok{fit =} \FunctionTok{map}\NormalTok{(model, broom}\SpecialCharTok{::}\NormalTok{glance)}
\NormalTok{  )}

\NormalTok{tbl\_ff3\_results }\OtherTok{\textless{}{-}}\NormalTok{ tbl\_ff3\_models }\SpecialCharTok{|\textgreater{}}
  \FunctionTok{select}\NormalTok{(ME\_rank10, coefficients) }\SpecialCharTok{|\textgreater{}}
  \FunctionTok{unnest}\NormalTok{(coefficients) }\SpecialCharTok{|\textgreater{}}
  \FunctionTok{mutate}\NormalTok{(}\AttributeTok{significance =} \FunctionTok{compute\_significance\_stars}\NormalTok{(p.value)) }\SpecialCharTok{|\textgreater{}}
  \FunctionTok{arrange}\NormalTok{(ME\_rank10, term)}

\NormalTok{tbl\_ff3\_fit }\OtherTok{\textless{}{-}}\NormalTok{ tbl\_ff3\_models }\SpecialCharTok{|\textgreater{}}
  \FunctionTok{select}\NormalTok{(ME\_rank10, fit) }\SpecialCharTok{|\textgreater{}}
  \FunctionTok{unnest}\NormalTok{(fit) }\SpecialCharTok{|\textgreater{}}
  \FunctionTok{select}\NormalTok{(}
\NormalTok{    ME\_rank10, r.squared, adj.r.squared,}
\NormalTok{    sigma, statistic, p.value, nobs}
\NormalTok{  ) }\SpecialCharTok{|\textgreater{}}
  \FunctionTok{arrange}\NormalTok{(ME\_rank10)}

\NormalTok{tbl\_ff3\_alpha }\OtherTok{\textless{}{-}}\NormalTok{ tbl\_ff3\_results }\SpecialCharTok{|\textgreater{}}
  \FunctionTok{filter}\NormalTok{(term }\SpecialCharTok{==} \StringTok{"(Intercept)"}\NormalTok{) }\SpecialCharTok{|\textgreater{}}
  \FunctionTok{transmute}\NormalTok{(}
\NormalTok{    ME\_rank10,}
    \AttributeTok{tbl\_ff3\_alpha =}\NormalTok{ estimate,}
    \AttributeTok{standard\_error =}\NormalTok{ std.error,}
\NormalTok{    statistic,}
    \AttributeTok{p\_value =}\NormalTok{ p.value,}
\NormalTok{    significance}
\NormalTok{  )}

\FunctionTok{show\_table\_LFL}\NormalTok{(tbl\_ff3\_results, }\AttributeTok{n =} \ConstantTok{Inf}\NormalTok{)}
\end{Highlighting}
\end{Shaded}

\begin{verbatim}
# A tibble: 40 x 7
# Groups:   ME_rank10 [10]
   ME_rank10 term           estimate  std.error statistic    p.value
       <int> <chr>             <dbl>      <dbl>     <dbl>      <dbl>
 1         1 (Intercept)  0.0022160  0.0019246    1.1514  2.5446e- 1
 2         1 HML          0.037314   0.098061     0.38052 7.0500e- 1
 3         1 R_Me         0.92776    0.042009    22.085   2.5525e-29
 4         1 SMB          1.3610     0.081276    16.746   1.8226e-23
 5         2 (Intercept)  0.0012670  0.0015741    0.80490 4.2428e- 1
 6         2 HML          0.048745   0.080200     0.60779 5.4578e- 1
 7         2 R_Me         0.94434    0.034357    27.486   3.3917e-34
 8         2 SMB          1.2382     0.066472    18.627   1.1450e-25
 9         3 (Intercept)  0.0024754  0.0014256    1.7364  8.7992e- 2
10         3 HML          0.26825    0.072634     3.6931  5.0419e- 4
11         3 R_Me         0.96149    0.031116    30.900   7.1830e-37
12         3 SMB          1.0857     0.060201    18.034   5.4337e-25
13         4 (Intercept) -0.00088852 0.00085856  -1.0349  3.0517e- 1
14         4 HML          0.40182    0.043744     9.1858  9.0449e-13
15         4 R_Me         1.0215     0.018740    54.513   3.2219e-50
16         4 SMB          1.0607     0.036256    29.256   1.2872e-35
17         5 (Intercept) -0.00088220 0.00091786  -0.96115 3.4061e- 1
18         5 HML          0.31739    0.046765     6.7869  7.7807e- 9
19         5 R_Me         1.0289     0.020034    51.357   8.5258e-49
20         5 SMB          0.83200    0.038760    21.466   1.0638e-28
21         6 (Intercept)  0.0012087  0.0013152    0.91901 3.6203e- 1
22         6 HML          0.15817    0.067009     2.3605  2.1759e- 2
23         6 R_Me         1.0084     0.028706    35.128   7.7326e-40
24         6 SMB          0.61456    0.055539    11.065   1.0274e-15
25         7 (Intercept) -0.0019158  0.0011588   -1.6532  1.0388e- 1
26         7 HML          0.11287    0.059041     1.9117  6.1037e- 2
27         7 R_Me         1.0370     0.025293    41.001   1.8493e-43
28         7 SMB          0.54127    0.048935    11.061   1.0435e-15
29         8 (Intercept)  0.0011876  0.0014388    0.82540 4.1264e- 1
30         8 HML         -0.26686    0.073308    -3.6403  5.9543e- 4
31         8 R_Me         1.0247     0.031405    32.627   3.9979e-38
32         8 SMB          0.30377    0.060760     4.9995  5.9999e- 6
33         9 (Intercept) -0.0010787  0.0015851   -0.68056 4.9895e- 1
34         9 HML         -0.013149   0.080760    -0.16281 8.7125e- 1
35         9 R_Me         1.0790     0.034597    31.187   4.4005e-37
36         9 SMB          0.20749    0.066936     3.0999  3.0278e- 3
37        10 (Intercept)  0.00017429 0.00042384   0.41123 6.8247e- 1
38        10 HML         -0.0028300  0.021595    -0.13105 8.9620e- 1
39        10 R_Me         0.98160    0.0092510  106.11    2.9778e-66
40        10 SMB         -0.17118    0.017898    -9.5641  2.2416e-13
   significance
   <chr>       
 1 ""          
 2 ""          
 3 "***"       
 4 "***"       
 5 ""          
 6 ""          
 7 "***"       
 8 "***"       
 9 "*"         
10 "***"       
11 "***"       
12 "***"       
13 ""          
14 "***"       
15 "***"       
16 "***"       
17 ""          
18 "***"       
19 "***"       
20 "***"       
21 ""          
22 "**"        
23 "***"       
24 "***"       
25 ""          
26 "*"         
27 "***"       
28 "***"       
29 ""          
30 "***"       
31 "***"       
32 "***"       
33 ""          
34 ""          
35 "***"       
36 "***"       
37 ""          
38 ""          
39 "***"       
40 "***"       
\end{verbatim}

\begin{Shaded}
\begin{Highlighting}[]
\FunctionTok{show\_table\_LFL}\NormalTok{(tbl\_ff3\_fit, }\AttributeTok{n =} \ConstantTok{Inf}\NormalTok{)}
\end{Highlighting}
\end{Shaded}

\begin{verbatim}
# A tibble: 10 x 7
# Groups:   ME_rank10 [10]
   ME_rank10 r.squared adj.r.squared     sigma statistic    p.value  nobs
       <int>     <dbl>         <dbl>     <dbl>     <dbl>      <dbl> <int>
 1         1   0.91540       0.91087 0.012564     201.99 5.3609e-30    60
 2         2   0.93996       0.93674 0.010276     292.24 3.6748e-34    60
 3         3   0.94888       0.94615 0.0093063    346.51 4.0778e-36    60
 4         4   0.98279       0.98187 0.0056047   1066.2  2.3878e-49    60
 5         5   0.97989       0.97881 0.0059918    909.61 1.8726e-47    60
 6         6   0.95716       0.95487 0.0085857    417.09 2.9086e-38    60
 7         7   0.96821       0.96651 0.0075647    568.60 6.8782e-42    60
 8         8   0.95121       0.94860 0.0093927    363.94 1.1072e-36    60
 9         9   0.94869       0.94594 0.010347     345.12 4.5371e-36    60
10        10   0.99594       0.99572 0.0027668   4581.1  6.5209e-67    60
\end{verbatim}

\begin{Shaded}
\begin{Highlighting}[]
\NormalTok{plot\_ff3\_alpha }\OtherTok{\textless{}{-}}\NormalTok{ tbl\_ff3\_alpha }\SpecialCharTok{|\textgreater{}}
  \FunctionTok{mutate}\NormalTok{(}\AttributeTok{ME\_rank10 =} \FunctionTok{factor}\NormalTok{(ME\_rank10, }\AttributeTok{levels =} \DecValTok{1}\SpecialCharTok{:}\DecValTok{10}\NormalTok{)) }\SpecialCharTok{|\textgreater{}}
  \FunctionTok{ggplot}\NormalTok{(}\FunctionTok{aes}\NormalTok{(}\AttributeTok{x =}\NormalTok{ ME\_rank10, }\AttributeTok{y =}\NormalTok{ tbl\_ff3\_alpha)) }\SpecialCharTok{+}
  \FunctionTok{geom\_col}\NormalTok{() }\SpecialCharTok{+}
  \FunctionTok{geom\_hline}\NormalTok{(}\AttributeTok{yintercept =} \DecValTok{0}\NormalTok{, }\AttributeTok{linetype =} \StringTok{"dotted"}\NormalTok{) }\SpecialCharTok{+}
  \FunctionTok{scale\_y\_continuous}\NormalTok{(}\AttributeTok{labels =}\NormalTok{ scales}\SpecialCharTok{::}\FunctionTok{label\_percent}\NormalTok{()) }\SpecialCharTok{+}
  \FunctionTok{labs}\NormalTok{(}
    \AttributeTok{x =} \StringTok{"Size Decile"}\NormalTok{,}
    \AttributeTok{y =} \StringTok{"FF3 Alpha"}\NormalTok{,}
    \AttributeTok{title =} \StringTok{"FF3 Alphas of Size Portfolios"}
\NormalTok{  )}

\FunctionTok{show\_plot\_LFL}\NormalTok{(}\FunctionTok{add\_lagrange\_theme\_LFL}\NormalTok{(plot\_ff3\_alpha))}
\end{Highlighting}
\end{Shaded}

\pandocbounded{\includegraphics[keepaspectratio]{part1_empirical/ch04_factor_models_files/figure-pdf/plot-ff3-alpha-1.pdf}}

\subsection{4.7.2
アルファとモデル適合度の比較}\label{ux30a2ux30ebux30d5ux30a1ux3068ux30e2ux30c7ux30ebux9069ux5408ux5ea6ux306eux6bd4ux8f03}

\begin{Shaded}
\begin{Highlighting}[]
\NormalTok{tbl\_model\_alpha\_comparison }\OtherTok{\textless{}{-}}\NormalTok{ tbl\_capm\_alpha }\SpecialCharTok{|\textgreater{}}
  \FunctionTok{select}\NormalTok{(ME\_rank10, tbl\_capm\_alpha) }\SpecialCharTok{|\textgreater{}}
  \FunctionTok{left\_join}\NormalTok{(}
\NormalTok{    tbl\_ff3\_alpha }\SpecialCharTok{|\textgreater{}}
      \FunctionTok{select}\NormalTok{(ME\_rank10, tbl\_ff3\_alpha),}
    \AttributeTok{by =} \StringTok{"ME\_rank10"}
\NormalTok{  ) }\SpecialCharTok{|\textgreater{}}
  \FunctionTok{pivot\_longer}\NormalTok{(}
    \AttributeTok{cols =} \FunctionTok{c}\NormalTok{(tbl\_capm\_alpha, tbl\_ff3\_alpha),}
    \AttributeTok{names\_to =} \StringTok{"model"}\NormalTok{,}
    \AttributeTok{values\_to =} \StringTok{"alpha"}
\NormalTok{  )}

\NormalTok{tbl\_model\_fit\_comparison }\OtherTok{\textless{}{-}}\NormalTok{ tbl\_capm\_fit }\SpecialCharTok{|\textgreater{}}
  \FunctionTok{select}\NormalTok{(}
\NormalTok{    ME\_rank10,}
    \AttributeTok{CAPM\_adjusted\_R2 =}\NormalTok{ adj.r.squared}
\NormalTok{  ) }\SpecialCharTok{|\textgreater{}}
  \FunctionTok{left\_join}\NormalTok{(}
\NormalTok{    tbl\_ff3\_fit }\SpecialCharTok{|\textgreater{}}
      \FunctionTok{select}\NormalTok{(}
\NormalTok{        ME\_rank10,}
        \AttributeTok{FF3\_adjusted\_R2 =}\NormalTok{ adj.r.squared}
\NormalTok{      ),}
    \AttributeTok{by =} \StringTok{"ME\_rank10"}
\NormalTok{  )}

\FunctionTok{show\_table\_LFL}\NormalTok{(tbl\_model\_alpha\_comparison, }\AttributeTok{n =} \ConstantTok{Inf}\NormalTok{)}
\end{Highlighting}
\end{Shaded}

\begin{verbatim}
# A tibble: 20 x 3
# Groups:   ME_rank10 [10]
   ME_rank10 model                alpha
       <int> <chr>                <dbl>
 1         1 tbl_capm_alpha  0.012412  
 2         1 tbl_ff3_alpha   0.0022160 
 3         2 tbl_capm_alpha  0.010641  
 4         2 tbl_ff3_alpha   0.0012670 
 5         3 tbl_capm_alpha  0.012195  
 6         3 tbl_ff3_alpha   0.0024754 
 7         4 tbl_capm_alpha  0.0095370 
 8         4 tbl_ff3_alpha  -0.00088852
 9         5 tbl_capm_alpha  0.0073101 
10         5 tbl_ff3_alpha  -0.00088220
11         6 tbl_capm_alpha  0.0067526 
12         6 tbl_ff3_alpha   0.0012087 
13         7 tbl_capm_alpha  0.0027911 
14         7 tbl_ff3_alpha  -0.0019158 
15         8 tbl_capm_alpha  0.0016328 
16         8 tbl_ff3_alpha   0.0011876 
17         9 tbl_capm_alpha  0.00035036
18         9 tbl_ff3_alpha  -0.0010787 
19        10 tbl_capm_alpha -0.0010957 
20        10 tbl_ff3_alpha   0.00017429
\end{verbatim}

\begin{Shaded}
\begin{Highlighting}[]
\FunctionTok{show\_table\_LFL}\NormalTok{(tbl\_model\_fit\_comparison, }\AttributeTok{n =} \ConstantTok{Inf}\NormalTok{)}
\end{Highlighting}
\end{Shaded}

\begin{verbatim}
# A tibble: 10 x 3
# Groups:   ME_rank10 [10]
   ME_rank10 CAPM_adjusted_R2 FF3_adjusted_R2
       <int>            <dbl>           <dbl>
 1         1          0.43109         0.91087
 2         2          0.51763         0.93674
 3         3          0.63895         0.94615
 4         4          0.71461         0.98187
 5         5          0.81104         0.97881
 6         6          0.85918         0.95487
 7         7          0.89456         0.96651
 8         8          0.89902         0.94860
 9         9          0.93721         0.94594
10        10          0.98827         0.99572
\end{verbatim}

\begin{Shaded}
\begin{Highlighting}[]
\NormalTok{plot\_model\_alpha }\OtherTok{\textless{}{-}}\NormalTok{ tbl\_model\_alpha\_comparison }\SpecialCharTok{|\textgreater{}}
  \FunctionTok{mutate}\NormalTok{(}\AttributeTok{ME\_rank10 =} \FunctionTok{factor}\NormalTok{(ME\_rank10, }\AttributeTok{levels =} \DecValTok{1}\SpecialCharTok{:}\DecValTok{10}\NormalTok{)) }\SpecialCharTok{|\textgreater{}}
  \FunctionTok{ggplot}\NormalTok{(}
    \FunctionTok{aes}\NormalTok{(}
      \AttributeTok{x =}\NormalTok{ ME\_rank10,}
      \AttributeTok{y =}\NormalTok{ alpha,}
      \AttributeTok{fill =}\NormalTok{ model}
\NormalTok{    )}
\NormalTok{  ) }\SpecialCharTok{+}
  \FunctionTok{geom\_col}\NormalTok{(}\AttributeTok{position =} \StringTok{"dodge"}\NormalTok{) }\SpecialCharTok{+}
  \FunctionTok{geom\_hline}\NormalTok{(}\AttributeTok{yintercept =} \DecValTok{0}\NormalTok{, }\AttributeTok{linetype =} \StringTok{"dotted"}\NormalTok{) }\SpecialCharTok{+}
  \FunctionTok{scale\_y\_continuous}\NormalTok{(}\AttributeTok{labels =}\NormalTok{ scales}\SpecialCharTok{::}\FunctionTok{label\_percent}\NormalTok{()) }\SpecialCharTok{+}
  \FunctionTok{labs}\NormalTok{(}
    \AttributeTok{x =} \StringTok{"Size Decile"}\NormalTok{,}
    \AttributeTok{y =} \StringTok{"Monthly Alpha"}\NormalTok{,}
    \AttributeTok{fill =} \StringTok{"Model"}\NormalTok{,}
    \AttributeTok{title =} \StringTok{"CAPM and FF3 Alphas"}
\NormalTok{  )}

\FunctionTok{show\_plot\_LFL}\NormalTok{(}\FunctionTok{add\_lagrange\_theme\_LFL}\NormalTok{(plot\_model\_alpha))}
\end{Highlighting}
\end{Shaded}

\pandocbounded{\includegraphics[keepaspectratio]{part1_empirical/ch04_factor_models_files/figure-pdf/plot-model-alpha-comparison-1.pdf}}

\begin{Shaded}
\begin{Highlighting}[]
\NormalTok{plot\_model\_fit }\OtherTok{\textless{}{-}}\NormalTok{ tbl\_model\_fit\_comparison }\SpecialCharTok{|\textgreater{}}
  \FunctionTok{pivot\_longer}\NormalTok{(}
    \AttributeTok{cols =} \FunctionTok{c}\NormalTok{(CAPM\_adjusted\_R2, FF3\_adjusted\_R2),}
    \AttributeTok{names\_to =} \StringTok{"model"}\NormalTok{,}
    \AttributeTok{values\_to =} \StringTok{"adjusted\_R2"}
\NormalTok{  ) }\SpecialCharTok{|\textgreater{}}
  \FunctionTok{mutate}\NormalTok{(}\AttributeTok{ME\_rank10 =} \FunctionTok{factor}\NormalTok{(ME\_rank10, }\AttributeTok{levels =} \DecValTok{1}\SpecialCharTok{:}\DecValTok{10}\NormalTok{)) }\SpecialCharTok{|\textgreater{}}
  \FunctionTok{ggplot}\NormalTok{(}
    \FunctionTok{aes}\NormalTok{(}
      \AttributeTok{x =}\NormalTok{ ME\_rank10,}
      \AttributeTok{y =}\NormalTok{ adjusted\_R2,}
      \AttributeTok{fill =}\NormalTok{ model}
\NormalTok{    )}
\NormalTok{  ) }\SpecialCharTok{+}
  \FunctionTok{geom\_col}\NormalTok{(}\AttributeTok{position =} \StringTok{"dodge"}\NormalTok{) }\SpecialCharTok{+}
  \FunctionTok{labs}\NormalTok{(}
    \AttributeTok{x =} \StringTok{"Size Decile"}\NormalTok{,}
    \AttributeTok{y =} \StringTok{"Adjusted R{-}squared"}\NormalTok{,}
    \AttributeTok{fill =} \StringTok{"Model"}\NormalTok{,}
    \AttributeTok{title =} \StringTok{"Model Fit: CAPM and FF3"}
\NormalTok{  )}

\FunctionTok{show\_plot\_LFL}\NormalTok{(}\FunctionTok{add\_lagrange\_theme\_LFL}\NormalTok{(plot\_model\_fit))}
\end{Highlighting}
\end{Shaded}

\pandocbounded{\includegraphics[keepaspectratio]{part1_empirical/ch04_factor_models_files/figure-pdf/plot-model-fit-comparison-1.pdf}}

\section{4.8
ファクター・データの保存}\label{ux30d5ux30a1ux30afux30bfux30fcux30c7ux30fcux30bfux306eux4fddux5b58}

第5章以降で使用するファクター・データを、\texttt{simulation\_data/derived/ch06\_output.csv}に保存する。

\begin{Shaded}
\begin{Highlighting}[]
\NormalTok{tbl\_factor\_output }\OtherTok{\textless{}{-}}\NormalTok{ tbl\_factor }\SpecialCharTok{|\textgreater{}}
  \FunctionTok{select}\NormalTok{(}
\NormalTok{    month\_ID, year, month, month\_date,}
\NormalTok{    R\_F, R\_M, R\_Me, SMB, HML,}
\NormalTok{    SL, SN, SH, BL, BN, BH,}
\NormalTok{    market\_return\_coverage}
\NormalTok{  )}

\NormalTok{vec\_derived\_data\_dir }\OtherTok{\textless{}{-}} \FunctionTok{file.path}\NormalTok{(}
  \FunctionTok{get\_simulation\_data\_dir\_LFL}\NormalTok{(),}
  \StringTok{"derived"}
\NormalTok{)}

\ControlFlowTok{if}\NormalTok{ (}\SpecialCharTok{!}\FunctionTok{dir.exists}\NormalTok{(vec\_derived\_data\_dir)) \{}
  \FunctionTok{dir.create}\NormalTok{(vec\_derived\_data\_dir, }\AttributeTok{recursive =} \ConstantTok{TRUE}\NormalTok{)}
\NormalTok{\}}

\NormalTok{vec\_factor\_output\_path }\OtherTok{\textless{}{-}} \FunctionTok{file.path}\NormalTok{(}
\NormalTok{  vec\_derived\_data\_dir,}
  \StringTok{"ch06\_output.csv"}
\NormalTok{)}

\NormalTok{readr}\SpecialCharTok{::}\FunctionTok{write\_csv}\NormalTok{(}
\NormalTok{  tbl\_factor\_output,}
\NormalTok{  vec\_factor\_output\_path}
\NormalTok{)}

\NormalTok{tbl\_factor\_output\_check }\OtherTok{\textless{}{-}}\NormalTok{ readr}\SpecialCharTok{::}\FunctionTok{read\_csv}\NormalTok{(}
\NormalTok{  vec\_factor\_output\_path,}
  \AttributeTok{show\_col\_types =} \ConstantTok{FALSE}
\NormalTok{)}

\FunctionTok{stopifnot}\NormalTok{(}
  \FunctionTok{nrow}\NormalTok{(tbl\_factor\_output\_check) }\SpecialCharTok{==} \FunctionTok{nrow}\NormalTok{(tbl\_factor\_output),}
  \FunctionTok{identical}\NormalTok{(}\FunctionTok{names}\NormalTok{(tbl\_factor\_output\_check), }\FunctionTok{names}\NormalTok{(tbl\_factor\_output))}
\NormalTok{)}

\NormalTok{tbl\_factor\_output\_summary }\OtherTok{\textless{}{-}} \FunctionTok{tibble}\NormalTok{(}
  \AttributeTok{output\_file =}\NormalTok{ vec\_factor\_output\_path,}
  \AttributeTok{rows =} \FunctionTok{nrow}\NormalTok{(tbl\_factor\_output\_check),}
  \AttributeTok{columns =} \FunctionTok{ncol}\NormalTok{(tbl\_factor\_output\_check),}
  \AttributeTok{first\_month =} \FunctionTok{min}\NormalTok{(tbl\_factor\_output\_check}\SpecialCharTok{$}\NormalTok{month\_ID),}
  \AttributeTok{last\_month =} \FunctionTok{max}\NormalTok{(tbl\_factor\_output\_check}\SpecialCharTok{$}\NormalTok{month\_ID)}
\NormalTok{)}

\FunctionTok{show\_table\_LFL}\NormalTok{(tbl\_factor\_output\_summary)}
\end{Highlighting}
\end{Shaded}

\begin{verbatim}
# A tibble: 1 x 5
  output_file                                                                   
  <chr>                                                                         
1 C:/AnalyticFin/Projects/LagrangeFinance/simulation_data/derived/ch06_output.c~
   rows columns first_month last_month
  <int>   <int>       <dbl>      <dbl>
1    60      16          13         72
\end{verbatim}

\section{4.9 まとめ}\label{ux307eux3068ux3081-3}

本章では、前年度末時価総額を用いて市場ポートフォリオと企業規模ポートフォリオを構築した。また、企業規模と簿価時価比率による独立2×3ソートから6ポートフォリオを作成し、SMBとHMLを計算した。

企業規模10分位ポートフォリオに対してCAPMとFF3モデルを推定し、アルファと調整済み決定係数を比較した。回帰の反復処理には\texttt{nest()}と\texttt{purrr::map()}を使用しており、明示的なループは使用していない。

作成したファクター・データは次の場所に保存される。

\begin{Shaded}
\begin{Highlighting}[]
\NormalTok{simulation\_data/derived/ch06\_output.csv}
\end{Highlighting}
\end{Shaded}

第5章では、このファクター・データを利用して、個別企業の資本コスト、共分散構造、平均分散ポートフォリオを分析する。

\chapter{第5章
資産価格モデルとポートフォリオ最適化}\label{ux7b2c5ux7ae0-ux8cc7ux7523ux4fa1ux683cux30e2ux30c7ux30ebux3068ux30ddux30fcux30c8ux30d5ux30a9ux30eaux30aaux6700ux9069ux5316}

本章では、Fama--Frenchの3ファクター・モデルを用いて企業別の期待リターンを推定し、その結果を配当成長率の分析とポートフォリオ最適化へ接続する。

前章までに作成した月次株式データ、年次企業データ、ファクターデータを利用し、企業固有リスクと共通ファクターリスクから株式リターンの分散共分散行列を構築する。最後に、平均分散最適化を実行し、目標期待リターンの変化に応じた効率的フロンティアを確認する。

\section{5.1
分析対象とデータの接続}\label{ux5206ux6790ux5bfeux8c61ux3068ux30c7ux30fcux30bfux306eux63a5ux7d9a}

\subsection{5.1.1
共通関数と分析データ}\label{ux5171ux901aux95a2ux6570ux3068ux5206ux6790ux30c7ux30fcux30bf}

本章では、プロジェクト共通のutilityを読み込み、前章までに作成された派生データを利用する。パッケージの導入方法そのものは扱わず、分析に必要な処理から開始する。

\begin{Shaded}
\begin{Highlighting}[]
\FunctionTok{source}\NormalTok{(}\StringTok{"R/utility\_LFL.R"}\NormalTok{)}
\FunctionTok{load\_part1\_packages\_LFL}\NormalTok{()}
\end{Highlighting}
\end{Shaded}

\begin{Shaded}
\begin{Highlighting}[]
\NormalTok{tbl\_monthly\_raw }\OtherTok{\textless{}{-}} \FunctionTok{read\_derived\_csv\_LFL}\NormalTok{(}\StringTok{"ch05\_output1.csv"}\NormalTok{)}
\NormalTok{tbl\_annual\_raw }\OtherTok{\textless{}{-}} \FunctionTok{read\_derived\_csv\_LFL}\NormalTok{(}\StringTok{"ch05\_output2.csv"}\NormalTok{)}
\NormalTok{tbl\_factor }\OtherTok{\textless{}{-}} \FunctionTok{read\_derived\_csv\_LFL}\NormalTok{(}\StringTok{"ch06\_output.csv"}\NormalTok{)}
\end{Highlighting}
\end{Shaded}

\subsection{5.1.2
推定対象企業の選定}\label{ux63a8ux5b9aux5bfeux8c61ux4f01ux696dux306eux9078ux5b9a}

企業別の時系列回帰には一定数の観測値が必要である。本章では、2020年末まで上場が継続し、月次観測値が36か月以上存在する企業を推定対象とする。

\begin{Shaded}
\begin{Highlighting}[]
\NormalTok{tbl\_monthly }\OtherTok{\textless{}{-}}\NormalTok{ tbl\_monthly\_raw }\SpecialCharTok{|\textgreater{}}
  \FunctionTok{group\_by}\NormalTok{(firm\_ID) }\SpecialCharTok{|\textgreater{}}
  \FunctionTok{mutate}\NormalTok{(}
    \AttributeTok{is\_public\_2020 =} \FunctionTok{max}\NormalTok{(month\_ID, }\AttributeTok{na.rm =} \ConstantTok{TRUE}\NormalTok{) }\SpecialCharTok{==} \DecValTok{72}\NormalTok{,}
    \AttributeTok{n\_observations =} \FunctionTok{n}\NormalTok{()}
\NormalTok{  ) }\SpecialCharTok{|\textgreater{}}
  \FunctionTok{ungroup}\NormalTok{() }\SpecialCharTok{|\textgreater{}}
  \FunctionTok{filter}\NormalTok{(}
\NormalTok{    is\_public\_2020,}
\NormalTok{    n\_observations }\SpecialCharTok{\textgreater{}=} \DecValTok{36}
\NormalTok{  ) }\SpecialCharTok{|\textgreater{}}
  \FunctionTok{select}\NormalTok{(}\SpecialCharTok{{-}}\NormalTok{n\_observations)}
\end{Highlighting}
\end{Shaded}

年次データについても、月次データに含まれる企業と年度の組合せに限定する。

\begin{Shaded}
\begin{Highlighting}[]
\NormalTok{tbl\_annual }\OtherTok{\textless{}{-}}\NormalTok{ tbl\_monthly }\SpecialCharTok{|\textgreater{}}
  \FunctionTok{distinct}\NormalTok{(year, firm\_ID) }\SpecialCharTok{|\textgreater{}}
  \FunctionTok{left\_join}\NormalTok{(}
\NormalTok{    tbl\_annual\_raw,}
    \AttributeTok{by =} \FunctionTok{c}\NormalTok{(}\StringTok{"year"}\NormalTok{, }\StringTok{"firm\_ID"}\NormalTok{)}
\NormalTok{  )}
\end{Highlighting}
\end{Shaded}

\subsection{5.1.3
月次株式リターンとファクターの結合}\label{ux6708ux6b21ux682aux5f0fux30eaux30bfux30fcux30f3ux3068ux30d5ux30a1ux30afux30bfux30fcux306eux7d50ux5408}

月次株式データとファクターデータを \texttt{month\_ID}
で結合する。無リスク金利 \texttt{R\_F}
が月次株式データにも含まれる場合は、重複を避けるため、ファクターデータ側の値を利用する。

\begin{Shaded}
\begin{Highlighting}[]
\NormalTok{tbl\_monthly\_factor }\OtherTok{\textless{}{-}}\NormalTok{ tbl\_monthly }\SpecialCharTok{|\textgreater{}}
  \FunctionTok{select}\NormalTok{(}\SpecialCharTok{{-}}\FunctionTok{any\_of}\NormalTok{(}\StringTok{"R\_F"}\NormalTok{)) }\SpecialCharTok{|\textgreater{}}
  \FunctionTok{left\_join}\NormalTok{(}
\NormalTok{    tbl\_factor,}
    \AttributeTok{by =} \StringTok{"month\_ID"}
\NormalTok{  )}
\end{Highlighting}
\end{Shaded}

\section{5.2
FF3モデルによる期待リターンの推定}\label{ff3ux30e2ux30c7ux30ebux306bux3088ux308bux671fux5f85ux30eaux30bfux30fcux30f3ux306eux63a8ux5b9a}

\subsection{5.2.1 FF3モデル}\label{ff3ux30e2ux30c7ux30eb}

企業 \(i\) の時点 \(t\)
における超過リターンを、次の3ファクター・モデルで表す。

\[
R_{i,t} - R_{F,t}
=
\beta_{i,M} R_{M,t}^{e}
+
\beta_{i,SMB} SMB_t
+
\beta_{i,HML} HML_t
+
\varepsilon_{i,t}
\]

ここで、\(R_{M,t}^{e}\) は市場超過リターン、\(SMB_t\)
は企業規模ファクター、\(HML_t\)
は簿価時価比率ファクターである。本章では、入力データに含まれる超過リターン
\texttt{Re} を被説明変数とし、定数項を置かずに企業別回帰を推定する。

\begin{Shaded}
\begin{Highlighting}[]
\NormalTok{tbl\_ff3\_models }\OtherTok{\textless{}{-}}\NormalTok{ tbl\_monthly\_factor }\SpecialCharTok{|\textgreater{}}
  \FunctionTok{select}\NormalTok{(}
\NormalTok{    firm\_ID,}
\NormalTok{    Re,}
\NormalTok{    R\_Me,}
\NormalTok{    SMB,}
\NormalTok{    HML}
\NormalTok{  ) }\SpecialCharTok{|\textgreater{}}
  \FunctionTok{drop\_na}\NormalTok{() }\SpecialCharTok{|\textgreater{}}
  \FunctionTok{group\_by}\NormalTok{(firm\_ID) }\SpecialCharTok{|\textgreater{}}
  \FunctionTok{nest}\NormalTok{(}\AttributeTok{tbl\_firm =} \SpecialCharTok{{-}}\NormalTok{firm\_ID) }\SpecialCharTok{|\textgreater{}}
  \FunctionTok{mutate}\NormalTok{(}
    \AttributeTok{list\_model =} \FunctionTok{map}\NormalTok{(}
\NormalTok{      tbl\_firm,}
\NormalTok{      \textbackslash{}(tbl\_data) \{}
        \FunctionTok{lm}\NormalTok{(}
\NormalTok{          Re }\SpecialCharTok{\textasciitilde{}} \DecValTok{0} \SpecialCharTok{+}\NormalTok{ R\_Me }\SpecialCharTok{+}\NormalTok{ SMB }\SpecialCharTok{+}\NormalTok{ HML,}
          \AttributeTok{data =}\NormalTok{ tbl\_data}
\NormalTok{        )}
\NormalTok{      \}}
\NormalTok{    ),}
    \AttributeTok{tbl\_coefficient =} \FunctionTok{map}\NormalTok{(}
\NormalTok{      list\_model,}
\NormalTok{      broom}\SpecialCharTok{::}\NormalTok{tidy}
\NormalTok{    )}
\NormalTok{  ) }\SpecialCharTok{|\textgreater{}}
  \FunctionTok{ungroup}\NormalTok{()}
\end{Highlighting}
\end{Shaded}

\subsection{5.2.2
企業別ファクター・ローディング}\label{ux4f01ux696dux5225ux30d5ux30a1ux30afux30bfux30fcux30edux30fcux30c7ux30a3ux30f3ux30b0}

企業別回帰から得られた係数を横持ち形式に変換し、市場、企業規模、簿価時価比率に対する感応度として整理する。

\begin{Shaded}
\begin{Highlighting}[]
\NormalTok{tbl\_ff3\_loadings }\OtherTok{\textless{}{-}}\NormalTok{ tbl\_ff3\_models }\SpecialCharTok{|\textgreater{}}
  \FunctionTok{select}\NormalTok{(}
\NormalTok{    firm\_ID,}
\NormalTok{    tbl\_coefficient}
\NormalTok{  ) }\SpecialCharTok{|\textgreater{}}
  \FunctionTok{unnest}\NormalTok{(tbl\_coefficient) }\SpecialCharTok{|\textgreater{}}
  \FunctionTok{select}\NormalTok{(}
\NormalTok{    firm\_ID,}
\NormalTok{    term,}
\NormalTok{    estimate}
\NormalTok{  ) }\SpecialCharTok{|\textgreater{}}
  \FunctionTok{pivot\_wider}\NormalTok{(}
    \AttributeTok{names\_from =}\NormalTok{ term,}
    \AttributeTok{values\_from =}\NormalTok{ estimate}
\NormalTok{  ) }\SpecialCharTok{|\textgreater{}}
  \FunctionTok{rename}\NormalTok{(}
    \AttributeTok{beta\_market =}\NormalTok{ R\_Me,}
    \AttributeTok{beta\_size =}\NormalTok{ SMB,}
    \AttributeTok{beta\_value =}\NormalTok{ HML}
\NormalTok{  )}

\FunctionTok{show\_table\_head\_LFL}\NormalTok{(tbl\_ff3\_loadings)}
\end{Highlighting}
\end{Shaded}

\begin{verbatim}
# A tibble: 6 x 4
  firm_ID beta_market beta_size beta_value
    <dbl>       <dbl>     <dbl>      <dbl>
1       1     1.6191    2.6384     0.48152
2       2     1.2405    0.50565   -0.52252
3       3     0.69623   1.1439     0.33640
4       4     0.37560   0.50193   -0.43245
5       5     0.58264   1.5325     1.1194 
6       7     0.38022   0.92052   -0.29961
\end{verbatim}

\subsection{5.2.3
ファクター・リスクプレミアム}\label{ux30d5ux30a1ux30afux30bfux30fcux30eaux30b9ux30afux30d7ux30ecux30dfux30a2ux30e0}

月次ファクターリターンの標本平均を12倍し、年率換算したファクター・リスクプレミアムを求める。

\begin{Shaded}
\begin{Highlighting}[]
\NormalTok{vec\_factor\_names }\OtherTok{\textless{}{-}} \FunctionTok{c}\NormalTok{(}\StringTok{"R\_Me"}\NormalTok{, }\StringTok{"SMB"}\NormalTok{, }\StringTok{"HML"}\NormalTok{)}

\NormalTok{tbl\_factor\_risk\_premium }\OtherTok{\textless{}{-}}\NormalTok{ tbl\_factor }\SpecialCharTok{|\textgreater{}}
  \FunctionTok{summarise}\NormalTok{(}
    \FunctionTok{across}\NormalTok{(}
      \FunctionTok{all\_of}\NormalTok{(vec\_factor\_names),}
\NormalTok{      \textbackslash{}(vec\_return) \{}
        \FunctionTok{mean}\NormalTok{(vec\_return, }\AttributeTok{na.rm =} \ConstantTok{TRUE}\NormalTok{) }\SpecialCharTok{*} \DecValTok{12}
\NormalTok{      \}}
\NormalTok{    )}
\NormalTok{  ) }\SpecialCharTok{|\textgreater{}}
  \FunctionTok{pivot\_longer}\NormalTok{(}
    \AttributeTok{cols =} \FunctionTok{everything}\NormalTok{(),}
    \AttributeTok{names\_to =} \StringTok{"factor"}\NormalTok{,}
    \AttributeTok{values\_to =} \StringTok{"annual\_risk\_premium"}
\NormalTok{  )}

\NormalTok{vec\_factor\_risk\_premium }\OtherTok{\textless{}{-}}\NormalTok{ tbl\_factor\_risk\_premium }\SpecialCharTok{|\textgreater{}}
  \FunctionTok{arrange}\NormalTok{(}\FunctionTok{match}\NormalTok{(factor, vec\_factor\_names)) }\SpecialCharTok{|\textgreater{}}
  \FunctionTok{pull}\NormalTok{(annual\_risk\_premium)}

\FunctionTok{names}\NormalTok{(vec\_factor\_risk\_premium) }\OtherTok{\textless{}{-}}\NormalTok{ vec\_factor\_names}

\FunctionTok{show\_table\_LFL}\NormalTok{(tbl\_factor\_risk\_premium)}
\end{Highlighting}
\end{Shaded}

\begin{verbatim}
# A tibble: 3 x 2
  factor annual_risk_premium
  <chr>                <dbl>
1 R_Me              0.048258
2 SMB               0.078586
3 HML               0.083395
\end{verbatim}

\subsection{5.2.4
株式資本コスト}\label{ux682aux5f0fux8cc7ux672cux30b3ux30b9ux30c8}

FF3モデルに基づく企業 \(i\)
の期待超過リターンは、ファクター・ローディングとファクター・リスクプレミアムの積和として求められる。

\[
E[R_i-R_F]
=
\beta_{i,M}E[R_M^e]
+
\beta_{i,SMB}E[SMB]
+
\beta_{i,HML}E[HML]
\]

本章では、この期待超過リターンを企業の株式資本コストとして利用する。

\begin{Shaded}
\begin{Highlighting}[]
\NormalTok{tbl\_cost\_of\_equity }\OtherTok{\textless{}{-}}\NormalTok{ tbl\_ff3\_loadings }\SpecialCharTok{|\textgreater{}}
  \FunctionTok{mutate}\NormalTok{(}
    \AttributeTok{cost\_of\_equity =}
\NormalTok{      beta\_market }\SpecialCharTok{*}\NormalTok{ vec\_factor\_risk\_premium[[}\StringTok{"R\_Me"}\NormalTok{]] }\SpecialCharTok{+}
\NormalTok{      beta\_size }\SpecialCharTok{*}\NormalTok{ vec\_factor\_risk\_premium[[}\StringTok{"SMB"}\NormalTok{]] }\SpecialCharTok{+}
\NormalTok{      beta\_value }\SpecialCharTok{*}\NormalTok{ vec\_factor\_risk\_premium[[}\StringTok{"HML"}\NormalTok{]]}
\NormalTok{  )}

\FunctionTok{show\_table\_head\_LFL}\NormalTok{(tbl\_cost\_of\_equity)}
\end{Highlighting}
\end{Shaded}

\begin{verbatim}
# A tibble: 6 x 5
  firm_ID beta_market beta_size beta_value cost_of_equity
    <dbl>       <dbl>     <dbl>      <dbl>          <dbl>
1       1     1.6191    2.6384     0.48152       0.32563 
2       2     1.2405    0.50565   -0.52252       0.056027
3       3     0.69623   1.1439     0.33640       0.15155 
4       4     0.37560   0.50193   -0.43245       0.021505
5       5     0.58264   1.5325     1.1194        0.24190 
6       7     0.38022   0.92052   -0.29961       0.065702
\end{verbatim}

\section{5.3
期待リターンと配当成長率}\label{ux671fux5f85ux30eaux30bfux30fcux30f3ux3068ux914dux5f53ux6210ux9577ux7387}

\subsection{5.3.1
期待リターンの企業間分布}\label{ux671fux5f85ux30eaux30bfux30fcux30f3ux306eux4f01ux696dux9593ux5206ux5e03}

\begin{Shaded}
\begin{Highlighting}[]
\NormalTok{plot\_cost\_of\_equity\_distribution }\OtherTok{\textless{}{-}}\NormalTok{ tbl\_cost\_of\_equity }\SpecialCharTok{|\textgreater{}}
  \FunctionTok{ggplot}\NormalTok{(}
    \FunctionTok{aes}\NormalTok{(}\AttributeTok{x =}\NormalTok{ cost\_of\_equity)}
\NormalTok{  ) }\SpecialCharTok{+}
  \FunctionTok{geom\_histogram}\NormalTok{(}\AttributeTok{bins =} \DecValTok{30}\NormalTok{) }\SpecialCharTok{+}
  \FunctionTok{labs}\NormalTok{(}
    \AttributeTok{x =} \StringTok{"FF3 Cost of Equity"}\NormalTok{,}
    \AttributeTok{y =} \StringTok{"Count"}
\NormalTok{  ) }\SpecialCharTok{+}
  \FunctionTok{scale\_y\_continuous}\NormalTok{(}
    \AttributeTok{expand =} \FunctionTok{expansion}\NormalTok{(}\AttributeTok{mult =} \FunctionTok{c}\NormalTok{(}\DecValTok{0}\NormalTok{, }\FloatTok{0.05}\NormalTok{))}
\NormalTok{  )}

\NormalTok{plot\_cost\_of\_equity\_distribution }\OtherTok{\textless{}{-}} \FunctionTok{add\_lagrange\_theme\_LFL}\NormalTok{(}
\NormalTok{  plot\_cost\_of\_equity\_distribution}
\NormalTok{)}

\FunctionTok{show\_plot\_LFL}\NormalTok{(plot\_cost\_of\_equity\_distribution)}
\end{Highlighting}
\end{Shaded}

\pandocbounded{\includegraphics[keepaspectratio]{part1_empirical/ch05_asset_pricing_files/figure-pdf/plot-cost-of-equity-distribution-1.pdf}}

\subsection{5.3.2
企業規模と期待リターン}\label{ux4f01ux696dux898fux6a21ux3068ux671fux5f85ux30eaux30bfux30fcux30f3}

\begin{Shaded}
\begin{Highlighting}[]
\NormalTok{tbl\_firm\_characteristics\_2020 }\OtherTok{\textless{}{-}}\NormalTok{ tbl\_annual }\SpecialCharTok{|\textgreater{}}
  \FunctionTok{arrange}\NormalTok{(}
\NormalTok{    firm\_ID,}
\NormalTok{    year}
\NormalTok{  ) }\SpecialCharTok{|\textgreater{}}
  \FunctionTok{group\_by}\NormalTok{(firm\_ID) }\SpecialCharTok{|\textgreater{}}
  \FunctionTok{mutate}\NormalTok{(}
    \AttributeTok{lagged\_ME =} \FunctionTok{lag}\NormalTok{(ME)}
\NormalTok{  ) }\SpecialCharTok{|\textgreater{}}
  \FunctionTok{ungroup}\NormalTok{() }\SpecialCharTok{|\textgreater{}}
  \FunctionTok{filter}\NormalTok{(year }\SpecialCharTok{==} \DecValTok{2020}\NormalTok{) }\SpecialCharTok{|\textgreater{}}
  \FunctionTok{mutate}\NormalTok{(}
    \AttributeTok{size\_group =} \FunctionTok{ntile}\NormalTok{(lagged\_ME, }\DecValTok{2}\NormalTok{),}
    \AttributeTok{size\_group =} \FunctionTok{factor}\NormalTok{(}
\NormalTok{      size\_group,}
      \AttributeTok{levels =} \FunctionTok{c}\NormalTok{(}\DecValTok{1}\NormalTok{, }\DecValTok{2}\NormalTok{),}
      \AttributeTok{labels =} \FunctionTok{c}\NormalTok{(}\StringTok{"Small"}\NormalTok{, }\StringTok{"Large"}\NormalTok{)}
\NormalTok{    )}
\NormalTok{  ) }\SpecialCharTok{|\textgreater{}}
  \FunctionTok{inner\_join}\NormalTok{(}
\NormalTok{    tbl\_cost\_of\_equity,}
    \AttributeTok{by =} \StringTok{"firm\_ID"}
\NormalTok{  ) }\SpecialCharTok{|\textgreater{}}
  \FunctionTok{drop\_na}\NormalTok{(}
\NormalTok{    lagged\_ME,}
\NormalTok{    cost\_of\_equity,}
\NormalTok{    size\_group}
\NormalTok{  )}
\end{Highlighting}
\end{Shaded}

\begin{Shaded}
\begin{Highlighting}[]
\NormalTok{plot\_cost\_of\_equity\_by\_size }\OtherTok{\textless{}{-}}\NormalTok{ tbl\_firm\_characteristics\_2020 }\SpecialCharTok{|\textgreater{}}
  \FunctionTok{ggplot}\NormalTok{(}
    \FunctionTok{aes}\NormalTok{(}
      \AttributeTok{x =}\NormalTok{ cost\_of\_equity,}
      \AttributeTok{linetype =}\NormalTok{ size\_group}
\NormalTok{    )}
\NormalTok{  ) }\SpecialCharTok{+}
  \FunctionTok{geom\_density}\NormalTok{() }\SpecialCharTok{+}
  \FunctionTok{labs}\NormalTok{(}
    \AttributeTok{x =} \StringTok{"FF3 Cost of Equity"}\NormalTok{,}
    \AttributeTok{y =} \StringTok{"Estimated Density"}\NormalTok{,}
    \AttributeTok{linetype =} \StringTok{"Firm Size"}
\NormalTok{  ) }\SpecialCharTok{+}
  \FunctionTok{scale\_y\_continuous}\NormalTok{(}
    \AttributeTok{expand =} \FunctionTok{expansion}\NormalTok{(}\AttributeTok{mult =} \FunctionTok{c}\NormalTok{(}\DecValTok{0}\NormalTok{, }\FloatTok{0.05}\NormalTok{)),}
    \AttributeTok{breaks =} \ConstantTok{NULL}
\NormalTok{  )}

\NormalTok{plot\_cost\_of\_equity\_by\_size }\OtherTok{\textless{}{-}} \FunctionTok{add\_lagrange\_theme\_LFL}\NormalTok{(}
\NormalTok{  plot\_cost\_of\_equity\_by\_size}
\NormalTok{)}

\FunctionTok{show\_plot\_LFL}\NormalTok{(plot\_cost\_of\_equity\_by\_size)}
\end{Highlighting}
\end{Shaded}

\pandocbounded{\includegraphics[keepaspectratio]{part1_empirical/ch05_asset_pricing_files/figure-pdf/plot-cost-of-equity-by-size-1.pdf}}

\subsection{5.3.3 配当利回り}\label{ux914dux5f53ux5229ux56deux308a}

\begin{Shaded}
\begin{Highlighting}[]
\NormalTok{tbl\_dividend\_yield\_2020 }\OtherTok{\textless{}{-}}\NormalTok{ tbl\_monthly }\SpecialCharTok{|\textgreater{}}
  \FunctionTok{filter}\NormalTok{(year }\SpecialCharTok{==} \DecValTok{2020}\NormalTok{) }\SpecialCharTok{|\textgreater{}}
  \FunctionTok{group\_by}\NormalTok{(firm\_ID) }\SpecialCharTok{|\textgreater{}}
  \FunctionTok{summarise}\NormalTok{(}
    \AttributeTok{annual\_dividend =} \FunctionTok{sum}\NormalTok{(}
\NormalTok{      DPS }\SpecialCharTok{*}\NormalTok{ shares\_outstanding,}
      \AttributeTok{na.rm =} \ConstantTok{TRUE}
\NormalTok{    ),}
    \AttributeTok{latest\_ME =}\NormalTok{ ME[}\FunctionTok{which.max}\NormalTok{(month\_ID)],}
    \AttributeTok{.groups =} \StringTok{"drop"}
\NormalTok{  ) }\SpecialCharTok{|\textgreater{}}
  \FunctionTok{mutate}\NormalTok{(}
    \AttributeTok{dividend\_yield =}\NormalTok{ annual\_dividend }\SpecialCharTok{/}\NormalTok{ latest\_ME}
\NormalTok{  ) }\SpecialCharTok{|\textgreater{}}
  \FunctionTok{filter}\NormalTok{(}
    \FunctionTok{is.finite}\NormalTok{(dividend\_yield),}
\NormalTok{    dividend\_yield }\SpecialCharTok{\textgreater{}} \DecValTok{0}
\NormalTok{  )}

\FunctionTok{show\_table\_head\_LFL}\NormalTok{(tbl\_dividend\_yield\_2020)}
\end{Highlighting}
\end{Shaded}

\begin{verbatim}
# A tibble: 6 x 4
  firm_ID annual_dividend    latest_ME dividend_yield
    <dbl>           <dbl>        <dbl>          <dbl>
1       1       325472000   9708946000      0.033523 
2       2        37536000   5194044000      0.0072267
3       3      4660110000 214032195000      0.021773 
4       4       101352000   8652927000      0.011713 
5       5       745488000 235170402000      0.0031700
6       7        64848000  15909376000      0.0040761
\end{verbatim}

\subsection{5.3.4
潜在配当成長率}\label{ux6f5cux5728ux914dux5f53ux6210ux9577ux7387}

株式資本コストを \(r_i\)、配当利回りを \(y_i\)
とすると、本章で用いる潜在配当成長率 \(g_i\) は次式から逆算される。

\[
g_i
=
\frac{r_i-y_i}{1+y_i}
\]

\begin{Shaded}
\begin{Highlighting}[]
\NormalTok{tbl\_implied\_dividend\_growth }\OtherTok{\textless{}{-}}\NormalTok{ tbl\_dividend\_yield\_2020 }\SpecialCharTok{|\textgreater{}}
  \FunctionTok{inner\_join}\NormalTok{(}
\NormalTok{    tbl\_cost\_of\_equity }\SpecialCharTok{|\textgreater{}}
      \FunctionTok{select}\NormalTok{(}
\NormalTok{        firm\_ID,}
\NormalTok{        cost\_of\_equity}
\NormalTok{      ),}
    \AttributeTok{by =} \StringTok{"firm\_ID"}
\NormalTok{  ) }\SpecialCharTok{|\textgreater{}}
  \FunctionTok{mutate}\NormalTok{(}
    \AttributeTok{implied\_dividend\_growth =}
\NormalTok{      (cost\_of\_equity }\SpecialCharTok{{-}}\NormalTok{ dividend\_yield) }\SpecialCharTok{/}
\NormalTok{      (}\DecValTok{1} \SpecialCharTok{+}\NormalTok{ dividend\_yield)}
\NormalTok{  )}

\FunctionTok{show\_table\_head\_LFL}\NormalTok{(tbl\_implied\_dividend\_growth)}
\end{Highlighting}
\end{Shaded}

\begin{verbatim}
# A tibble: 6 x 6
  firm_ID annual_dividend    latest_ME dividend_yield cost_of_equity
    <dbl>           <dbl>        <dbl>          <dbl>          <dbl>
1       1       325472000   9708946000      0.033523        0.32563 
2       2        37536000   5194044000      0.0072267       0.056027
3       3      4660110000 214032195000      0.021773        0.15155 
4       4       101352000   8652927000      0.011713        0.021505
5       5       745488000 235170402000      0.0031700       0.24190 
6       7        64848000  15909376000      0.0040761       0.065702
  implied_dividend_growth
                    <dbl>
1               0.28263  
2               0.048450 
3               0.12701  
4               0.0096791
5               0.23797  
6               0.061376 
\end{verbatim}

\section{5.4
ファクターモデルによるリスク推定}\label{ux30d5ux30a1ux30afux30bfux30fcux30e2ux30c7ux30ebux306bux3088ux308bux30eaux30b9ux30afux63a8ux5b9a}

\subsection{5.4.1
ファクターの分散共分散}\label{ux30d5ux30a1ux30afux30bfux30fcux306eux5206ux6563ux5171ux5206ux6563}

\begin{Shaded}
\begin{Highlighting}[]
\NormalTok{tbl\_factor\_covariance }\OtherTok{\textless{}{-}} \FunctionTok{compute\_covariance\_table\_LFL}\NormalTok{(}
  \AttributeTok{tbl\_data =}\NormalTok{ tbl\_factor,}
  \AttributeTok{vec\_variables =}\NormalTok{ vec\_factor\_names,}
  \AttributeTok{annualization\_factor =} \DecValTok{12}
\NormalTok{)}

\NormalTok{matrix\_factor\_covariance }\OtherTok{\textless{}{-}} \FunctionTok{convert\_covariance\_table\_to\_matrix\_LFL}\NormalTok{(}
\NormalTok{  tbl\_factor\_covariance}
\NormalTok{)}

\FunctionTok{show\_table\_LFL}\NormalTok{(tbl\_factor\_covariance)}
\end{Highlighting}
\end{Shaded}

\begin{verbatim}
# A tibble: 9 x 3
  variable_row variable_column covariance
  <chr>        <chr>                <dbl>
1 R_Me         R_Me             0.020665 
2 R_Me         SMB             -0.0039063
3 R_Me         HML              0.0015310
4 SMB          R_Me            -0.0039063
5 SMB          SMB              0.0062694
6 SMB          HML             -0.0018855
7 HML          R_Me             0.0015310
8 HML          SMB             -0.0018855
9 HML          HML              0.0039129
\end{verbatim}

\subsection{5.4.2
投資ユニバース}\label{ux6295ux8cc7ux30e6ux30cbux30d0ux30fcux30b9}

\begin{Shaded}
\begin{Highlighting}[]
\NormalTok{n\_portfolio\_firms }\OtherTok{\textless{}{-}} \DecValTok{100}\NormalTok{L}

\NormalTok{tbl\_investment\_universe }\OtherTok{\textless{}{-}}\NormalTok{ tbl\_annual }\SpecialCharTok{|\textgreater{}}
  \FunctionTok{filter}\NormalTok{(year }\SpecialCharTok{==} \DecValTok{2020}\NormalTok{) }\SpecialCharTok{|\textgreater{}}
  \FunctionTok{inner\_join}\NormalTok{(}
\NormalTok{    tbl\_cost\_of\_equity }\SpecialCharTok{|\textgreater{}}
      \FunctionTok{select}\NormalTok{(}
\NormalTok{        firm\_ID,}
\NormalTok{        cost\_of\_equity}
\NormalTok{      ),}
    \AttributeTok{by =} \StringTok{"firm\_ID"}
\NormalTok{  ) }\SpecialCharTok{|\textgreater{}}
  \FunctionTok{slice\_max}\NormalTok{(}
    \AttributeTok{order\_by =}\NormalTok{ ME,}
    \AttributeTok{n =}\NormalTok{ n\_portfolio\_firms,}
    \AttributeTok{with\_ties =} \ConstantTok{FALSE}
\NormalTok{  ) }\SpecialCharTok{|\textgreater{}}
  \FunctionTok{arrange}\NormalTok{(}\FunctionTok{desc}\NormalTok{(ME)) }\SpecialCharTok{|\textgreater{}}
  \FunctionTok{select}\NormalTok{(}
\NormalTok{    firm\_ID,}
\NormalTok{    ME,}
\NormalTok{    cost\_of\_equity}
\NormalTok{  )}

\NormalTok{vec\_investment\_firm\_ids }\OtherTok{\textless{}{-}}\NormalTok{ tbl\_investment\_universe }\SpecialCharTok{|\textgreater{}}
  \FunctionTok{pull}\NormalTok{(firm\_ID)}

\NormalTok{n\_portfolio\_firms }\OtherTok{\textless{}{-}} \FunctionTok{length}\NormalTok{(vec\_investment\_firm\_ids)}

\FunctionTok{show\_table\_head\_LFL}\NormalTok{(tbl\_investment\_universe)}
\end{Highlighting}
\end{Shaded}

\begin{verbatim}
# A tibble: 6 x 3
  firm_ID        ME cost_of_equity
    <dbl>     <dbl>          <dbl>
1     790 2.7339e13       0.061424
2     987 5.4441e12      -0.071023
3    1237 5.2764e12       0.091188
4    1057 5.1050e12      -0.058108
5     988 4.9127e12      -0.10059 
6    1400 3.7669e12       0.20979 
\end{verbatim}

\subsection{5.4.3
企業固有リスク}\label{ux4f01ux696dux56faux6709ux30eaux30b9ux30af}

\begin{Shaded}
\begin{Highlighting}[]
\NormalTok{tbl\_idiosyncratic\_variance }\OtherTok{\textless{}{-}}\NormalTok{ tbl\_monthly\_factor }\SpecialCharTok{|\textgreater{}}
  \FunctionTok{filter}\NormalTok{(}
\NormalTok{    firm\_ID }\SpecialCharTok{\%in\%}\NormalTok{ vec\_investment\_firm\_ids}
\NormalTok{  ) }\SpecialCharTok{|\textgreater{}}
  \FunctionTok{inner\_join}\NormalTok{(}
\NormalTok{    tbl\_ff3\_loadings,}
    \AttributeTok{by =} \StringTok{"firm\_ID"}
\NormalTok{  ) }\SpecialCharTok{|\textgreater{}}
  \FunctionTok{mutate}\NormalTok{(}
    \AttributeTok{fitted\_return =}
\NormalTok{      R\_F }\SpecialCharTok{+}
\NormalTok{      beta\_market }\SpecialCharTok{*}\NormalTok{ R\_Me }\SpecialCharTok{+}
\NormalTok{      beta\_size }\SpecialCharTok{*}\NormalTok{ SMB }\SpecialCharTok{+}
\NormalTok{      beta\_value }\SpecialCharTok{*}\NormalTok{ HML,}
    \AttributeTok{residual\_return =}\NormalTok{ R }\SpecialCharTok{{-}}\NormalTok{ fitted\_return}
\NormalTok{  ) }\SpecialCharTok{|\textgreater{}}
  \FunctionTok{group\_by}\NormalTok{(firm\_ID) }\SpecialCharTok{|\textgreater{}}
  \FunctionTok{summarise}\NormalTok{(}
    \AttributeTok{idiosyncratic\_variance =}
      \DecValTok{12} \SpecialCharTok{*} \FunctionTok{var}\NormalTok{(}
\NormalTok{        residual\_return,}
        \AttributeTok{na.rm =} \ConstantTok{TRUE}
\NormalTok{      ),}
    \AttributeTok{.groups =} \StringTok{"drop"}
\NormalTok{  ) }\SpecialCharTok{|\textgreater{}}
  \FunctionTok{right\_join}\NormalTok{(}
\NormalTok{    tbl\_investment\_universe }\SpecialCharTok{|\textgreater{}}
      \FunctionTok{select}\NormalTok{(firm\_ID),}
    \AttributeTok{by =} \StringTok{"firm\_ID"}
\NormalTok{  ) }\SpecialCharTok{|\textgreater{}}
  \FunctionTok{arrange}\NormalTok{(}
    \FunctionTok{match}\NormalTok{(}
\NormalTok{      firm\_ID,}
\NormalTok{      vec\_investment\_firm\_ids}
\NormalTok{    )}
\NormalTok{  )}

\FunctionTok{show\_table\_head\_LFL}\NormalTok{(tbl\_idiosyncratic\_variance)}
\end{Highlighting}
\end{Shaded}

\begin{verbatim}
# A tibble: 6 x 2
  firm_ID idiosyncratic_variance
    <dbl>                  <dbl>
1     790               0.055478
2     987               0.036155
3    1237               0.053695
4    1057               0.045481
5     988               0.12455 
6    1400               0.094833
\end{verbatim}

\subsection{5.4.4
株式リターンの分散共分散}\label{ux682aux5f0fux30eaux30bfux30fcux30f3ux306eux5206ux6563ux5171ux5206ux6563}

企業別ファクター・ローディングを並べた行列を
\(\boldsymbol{B}\)、企業固有リスクの対角行列を
\(\boldsymbol{\Sigma}_{\varepsilon}\)
とすると、株式リターンの分散共分散行列は次式で構築できる。

\[
\boldsymbol{\Sigma}
=
\boldsymbol{B}
\boldsymbol{\Sigma}_F
\boldsymbol{B}^{\mathsf{T}}
+
\boldsymbol{\Sigma}_{\varepsilon}
\]

\begin{Shaded}
\begin{Highlighting}[]
\NormalTok{tbl\_portfolio\_loadings }\OtherTok{\textless{}{-}}\NormalTok{ tbl\_investment\_universe }\SpecialCharTok{|\textgreater{}}
  \FunctionTok{select}\NormalTok{(firm\_ID) }\SpecialCharTok{|\textgreater{}}
  \FunctionTok{left\_join}\NormalTok{(}
\NormalTok{    tbl\_ff3\_loadings,}
    \AttributeTok{by =} \StringTok{"firm\_ID"}
\NormalTok{  )}

\NormalTok{matrix\_factor\_loadings }\OtherTok{\textless{}{-}}\NormalTok{ tbl\_portfolio\_loadings }\SpecialCharTok{|\textgreater{}}
  \FunctionTok{select}\NormalTok{(}
\NormalTok{    beta\_market,}
\NormalTok{    beta\_size,}
\NormalTok{    beta\_value}
\NormalTok{  ) }\SpecialCharTok{|\textgreater{}}
  \FunctionTok{as.matrix}\NormalTok{()}

\NormalTok{vec\_idiosyncratic\_variance }\OtherTok{\textless{}{-}}\NormalTok{ tbl\_idiosyncratic\_variance }\SpecialCharTok{|\textgreater{}}
  \FunctionTok{pull}\NormalTok{(idiosyncratic\_variance)}

\NormalTok{matrix\_idiosyncratic\_covariance }\OtherTok{\textless{}{-}} \FunctionTok{diag}\NormalTok{(}
\NormalTok{  vec\_idiosyncratic\_variance}
\NormalTok{)}

\NormalTok{matrix\_return\_covariance }\OtherTok{\textless{}{-}}
\NormalTok{  matrix\_factor\_loadings }\SpecialCharTok{\%*\%}
\NormalTok{  matrix\_factor\_covariance }\SpecialCharTok{\%*\%}
  \FunctionTok{t}\NormalTok{(matrix\_factor\_loadings) }\SpecialCharTok{+}
\NormalTok{  matrix\_idiosyncratic\_covariance}

\FunctionTok{rownames}\NormalTok{(matrix\_return\_covariance) }\OtherTok{\textless{}{-}}\NormalTok{ vec\_investment\_firm\_ids}
\FunctionTok{colnames}\NormalTok{(matrix\_return\_covariance) }\OtherTok{\textless{}{-}}\NormalTok{ vec\_investment\_firm\_ids}
\end{Highlighting}
\end{Shaded}

\subsection{5.4.5
期待リターン}\label{ux671fux5f85ux30eaux30bfux30fcux30f3}

\begin{Shaded}
\begin{Highlighting}[]
\NormalTok{vec\_expected\_returns }\OtherTok{\textless{}{-}}\NormalTok{ tbl\_investment\_universe }\SpecialCharTok{|\textgreater{}}
  \FunctionTok{pull}\NormalTok{(cost\_of\_equity)}

\FunctionTok{names}\NormalTok{(vec\_expected\_returns) }\OtherTok{\textless{}{-}}\NormalTok{ vec\_investment\_firm\_ids}

\FunctionTok{tibble}\NormalTok{(}
  \AttributeTok{firm\_ID =}\NormalTok{ vec\_investment\_firm\_ids,}
  \AttributeTok{expected\_return =}\NormalTok{ vec\_expected\_returns}
\NormalTok{) }\SpecialCharTok{|\textgreater{}}
  \FunctionTok{show\_table\_head\_LFL}\NormalTok{()}
\end{Highlighting}
\end{Shaded}

\begin{verbatim}
# A tibble: 6 x 2
  firm_ID expected_return
    <dbl>           <dbl>
1     790        0.061424
2     987       -0.071023
3    1237        0.091188
4    1057       -0.058108
5     988       -0.10059 
6    1400        0.20979 
\end{verbatim}

\section{5.5
平均分散最適化}\label{ux5e73ux5747ux5206ux6563ux6700ux9069ux5316}

\subsection{5.5.1 最適化問題}\label{ux6700ux9069ux5316ux554fux984c}

目標期待リターンを \(\mu_p^*\) とすると、ポートフォリオ・ウェイト
\(\boldsymbol{w}\) は次の問題から求められる。

\[
\min_{\boldsymbol{w}}
\frac{1}{2}
\boldsymbol{w}^{\mathsf{T}}
\boldsymbol{\Sigma}
\boldsymbol{w}
\]

\[
\boldsymbol{w}^{\mathsf{T}}
\boldsymbol{\mu}
=
\mu_p^*
\]

\[
\boldsymbol{w}^{\mathsf{T}}
\boldsymbol{1}
=
1
\]

本章の基本設定では、空売りを禁止する不等式制約は追加しない。

\subsection{5.5.2
目標期待リターンの設定}\label{ux76eeux6a19ux671fux5f85ux30eaux30bfux30fcux30f3ux306eux8a2dux5b9a}

\begin{Shaded}
\begin{Highlighting}[]
\NormalTok{target\_return }\OtherTok{\textless{}{-}} \FloatTok{0.10}

\NormalTok{list\_mean\_variance\_portfolio }\OtherTok{\textless{}{-}} \FunctionTok{compute\_mean\_variance\_portfolio\_LFL}\NormalTok{(}
  \AttributeTok{matrix\_covariance =}\NormalTok{ matrix\_return\_covariance,}
  \AttributeTok{vec\_expected\_returns =}\NormalTok{ vec\_expected\_returns,}
  \AttributeTok{target\_return =}\NormalTok{ target\_return}
\NormalTok{)}
\end{Highlighting}
\end{Shaded}

\subsection{5.5.3
最適ポートフォリオ}\label{ux6700ux9069ux30ddux30fcux30c8ux30d5ux30a9ux30eaux30aa}

\begin{Shaded}
\begin{Highlighting}[]
\NormalTok{tbl\_optimal\_portfolio }\OtherTok{\textless{}{-}} \FunctionTok{create\_portfolio\_weight\_table\_LFL}\NormalTok{(}
  \AttributeTok{list\_optimization =}\NormalTok{ list\_mean\_variance\_portfolio,}
  \AttributeTok{vec\_asset\_ids =}\NormalTok{ vec\_investment\_firm\_ids,}
  \AttributeTok{target\_return =}\NormalTok{ target\_return}
\NormalTok{) }\SpecialCharTok{|\textgreater{}}
  \FunctionTok{arrange}\NormalTok{(}
    \FunctionTok{desc}\NormalTok{(}\FunctionTok{abs}\NormalTok{(weight))}
\NormalTok{  )}

\FunctionTok{show\_table\_head\_LFL}\NormalTok{(}
\NormalTok{  tbl\_optimal\_portfolio,}
  \AttributeTok{n =} \DecValTok{10}
\NormalTok{)}
\end{Highlighting}
\end{Shaded}

\begin{verbatim}
# A tibble: 10 x 3
   asset_ID    weight target_return
      <dbl>     <dbl>         <dbl>
 1      208 -0.16960            0.1
 2     1473  0.075792           0.1
 3      727  0.075617           0.1
 4     1253  0.073272           0.1
 5      987  0.068849           0.1
 6      829  0.062825           0.1
 7      409 -0.059685           0.1
 8       33  0.057845           0.1
 9      765  0.057677           0.1
10      709  0.057313           0.1
\end{verbatim}

\subsection{5.5.4
最適化結果の検算}\label{ux6700ux9069ux5316ux7d50ux679cux306eux691cux7b97}

\begin{Shaded}
\begin{Highlighting}[]
\NormalTok{tbl\_portfolio\_summary }\OtherTok{\textless{}{-}} \FunctionTok{compute\_portfolio\_summary\_LFL}\NormalTok{(}
  \AttributeTok{list\_optimization =}\NormalTok{ list\_mean\_variance\_portfolio,}
  \AttributeTok{target\_return =}\NormalTok{ target\_return}
\NormalTok{) }\SpecialCharTok{|\textgreater{}}
  \FunctionTok{mutate}\NormalTok{(}
    \AttributeTok{realized\_expected\_return =} \FunctionTok{sum}\NormalTok{(}
\NormalTok{      list\_mean\_variance\_portfolio}\SpecialCharTok{$}\NormalTok{solution }\SpecialCharTok{*}
\NormalTok{        vec\_expected\_returns}
\NormalTok{    ),}
    \AttributeTok{return\_constraint\_error =}
\NormalTok{      realized\_expected\_return }\SpecialCharTok{{-}}
\NormalTok{      target\_return,}
    \AttributeTok{weight\_constraint\_error =}
\NormalTok{      weight\_sum }\SpecialCharTok{{-}}
      \DecValTok{1}
\NormalTok{  )}

\FunctionTok{show\_table\_LFL}\NormalTok{(tbl\_portfolio\_summary)}
\end{Highlighting}
\end{Shaded}

\begin{verbatim}
# A tibble: 1 x 8
  target_return portfolio_return portfolio_variance portfolio_risk weight_sum
          <dbl>            <dbl>              <dbl>          <dbl>      <dbl>
1           0.1            0.100          0.0069899       0.083606          1
  realized_expected_return return_constraint_error weight_constraint_error
                     <dbl>                   <dbl>                   <dbl>
1                    0.100              4.1633e-17              2.2204e-16
\end{verbatim}

\section{5.6
効率的フロンティア}\label{ux52b9ux7387ux7684ux30d5ux30edux30f3ux30c6ux30a3ux30a2}

\subsection{5.6.1
目標期待リターンの範囲}\label{ux76eeux6a19ux671fux5f85ux30eaux30bfux30fcux30f3ux306eux7bc4ux56f2}

\begin{Shaded}
\begin{Highlighting}[]
\NormalTok{vec\_target\_returns }\OtherTok{\textless{}{-}} \FunctionTok{seq}\NormalTok{(}
  \AttributeTok{from =} \SpecialCharTok{{-}}\FloatTok{0.10}\NormalTok{,}
  \AttributeTok{to =} \FloatTok{0.40}\NormalTok{,}
  \AttributeTok{length.out =} \DecValTok{100}
\NormalTok{)}
\end{Highlighting}
\end{Shaded}

\subsection{5.6.2
複数ポートフォリオの推定}\label{ux8907ux6570ux30ddux30fcux30c8ux30d5ux30a9ux30eaux30aaux306eux63a8ux5b9a}

\begin{Shaded}
\begin{Highlighting}[]
\NormalTok{tbl\_efficient\_frontier }\OtherTok{\textless{}{-}} \FunctionTok{compute\_efficient\_frontier\_LFL}\NormalTok{(}
  \AttributeTok{matrix\_covariance =}\NormalTok{ matrix\_return\_covariance,}
  \AttributeTok{vec\_expected\_returns =}\NormalTok{ vec\_expected\_returns,}
  \AttributeTok{vec\_target\_returns =}\NormalTok{ vec\_target\_returns}
\NormalTok{)}

\FunctionTok{show\_table\_head\_LFL}\NormalTok{(}
\NormalTok{  tbl\_efficient\_frontier }\SpecialCharTok{|\textgreater{}}
    \FunctionTok{select}\NormalTok{(}
\NormalTok{      target\_return,}
\NormalTok{      portfolio\_risk}
\NormalTok{    )}
\NormalTok{)}
\end{Highlighting}
\end{Shaded}

\begin{verbatim}
# A tibble: 6 x 2
  target_return portfolio_risk
          <dbl>          <dbl>
1     -0.1            0.079982
2     -0.094949       0.078497
3     -0.089899       0.077068
4     -0.084848       0.075697
5     -0.079798       0.074389
6     -0.074747       0.073145
\end{verbatim}

\subsection{5.6.3
効率的フロンティアの描画}\label{ux52b9ux7387ux7684ux30d5ux30edux30f3ux30c6ux30a3ux30a2ux306eux63cfux753b}

\begin{Shaded}
\begin{Highlighting}[]
\NormalTok{plot\_efficient\_frontier }\OtherTok{\textless{}{-}}\NormalTok{ tbl\_efficient\_frontier }\SpecialCharTok{|\textgreater{}}
  \FunctionTok{ggplot}\NormalTok{(}
    \FunctionTok{aes}\NormalTok{(}
      \AttributeTok{x =}\NormalTok{ portfolio\_risk,}
      \AttributeTok{y =}\NormalTok{ target\_return}
\NormalTok{    )}
\NormalTok{  ) }\SpecialCharTok{+}
  \FunctionTok{geom\_line}\NormalTok{() }\SpecialCharTok{+}
  \FunctionTok{labs}\NormalTok{(}
    \AttributeTok{x =} \StringTok{"Portfolio Risk"}\NormalTok{,}
    \AttributeTok{y =} \StringTok{"Expected Return"}
\NormalTok{  )}

\NormalTok{plot\_efficient\_frontier }\OtherTok{\textless{}{-}} \FunctionTok{add\_lagrange\_theme\_LFL}\NormalTok{(}
\NormalTok{  plot\_efficient\_frontier}
\NormalTok{)}

\FunctionTok{show\_plot\_LFL}\NormalTok{(plot\_efficient\_frontier)}
\end{Highlighting}
\end{Shaded}

\pandocbounded{\includegraphics[keepaspectratio]{part1_empirical/ch05_asset_pricing_files/figure-pdf/plot-efficient-frontier-1.pdf}}

\section{5.7 まとめ}\label{ux307eux3068ux3081-4}

本章では、Fama--Frenchの3ファクター・モデルを企業別に推定し、市場、企業規模、簿価時価比率に対するファクター・ローディングを求めた。ファクター・リスクプレミアムとの積和から株式資本コストを計算し、配当利回りと組み合わせることで潜在配当成長率を推定した。

さらに、共通ファクターの分散共分散行列と企業固有リスクから株式リターンの分散共分散行列を構築した。これを平均分散最適化へ接続し、単一の目標期待リターンに対応する最適ポートフォリオと、複数の目標期待リターンからなる効率的フロンティアを確認した。

\chapter{第6章
イベントスタディ}\label{ux7b2c6ux7ae0-ux30a4ux30d9ux30f3ux30c8ux30b9ux30bfux30c7ux30a3}

本章では、企業の決算発表を対象としてイベントスタディを行う。イベントスタディは、企業イベントの前後に観察された株式リターンから、通常時に予想されるリターンを差し引き、イベント固有の株価反応を測定する方法である。

本章では、次の処理を順番に行う。

\begin{enumerate}
\def\labelenumi{\arabic{enumi}.}
\tightlist
\item
  個別株リターン、市場リターン、イベント情報を読み込む
\item
  決算サプライズを計算し、イベントを5群に分類する
\item
  発表日の翌取引日をイベント日0として設定する
\item
  推定期間を使ってイベントごとのマーケット・モデルを推定する
\item
  異常リターンと累積異常リターンを計算する
\item
  AAR、CAAR、複数のイベント・ウィンドウについて統計的検定を行う
\item
  決算サプライズと短期CARの関係を確認する
\item
  イベント単位と日次集計の結果を保存する
\end{enumerate}

反復処理には明示的なループを使わず、\texttt{purrr::map()}、\texttt{pmap\_dfr()}、\texttt{walk2()}を使用する。パイプ演算子はnative
pipeの\texttt{\textbar{}\textgreater{}}に統一する。

\section{6.1
イベントスタディの基本設計}\label{ux30a4ux30d9ux30f3ux30c8ux30b9ux30bfux30c7ux30a3ux306eux57faux672cux8a2dux8a08}

\subsection{6.1.1
マーケット・モデル}\label{ux30deux30fcux30b1ux30c3ux30c8ux30e2ux30c7ux30eb}

企業 \(i\) の日 \(t\) における実現リターンを \(R_{i,t}\)、市場リターンを
\(R_{M,t}\) とする。推定期間では、次のマーケット・モデルを推定する。

\[
R_{i,t}
=
\alpha_i
+
\beta_i R_{M,t}
+
\varepsilon_{i,t}
\]

推定された通常リターンは次式で表される。

\[
\widehat{R}_{i,t}
=
\widehat{\alpha}_i
+
\widehat{\beta}_i R_{M,t}
\]

異常リターン \(AR_{i,t}\) は、実現リターンと通常リターンの差である。

\[
AR_{i,t}
=
R_{i,t}
-
\widehat{R}_{i,t}
\]

イベント期間の開始日から日 \(t\) までの累積異常リターン \(CAR_{i,t}\)
は、異常リターンを累積して計算する。

\[
CAR_{i,t}
=
\sum_{\tau=T_1}^{t} AR_{i,\tau}
\]

\subsection{6.1.2 分析期間}\label{ux5206ux6790ux671fux9593}

本章では、イベント日0を基準として次の期間を使用する。

{\def\LTcaptype{none} % do not increment counter
\begin{longtable}[]{@{}lrl@{}}
\toprule\noalign{}
期間 & 相対取引日 & 用途 \\
\midrule\noalign{}
\endhead
\bottomrule\noalign{}
\endlastfoot
推定期間 & -130日から-31日 & マーケット・モデルの推定 \\
事前イベント期間 & -30日から-1日 & 発表前の株価反応 \\
イベント日 & 0日 & 発表後最初の取引日 \\
事後イベント期間 & 1日から30日 & 発表後の株価反応 \\
\end{longtable}
}

推定期間は100取引日、イベント期間は61取引日である。マーケット・モデルの推定には、最低80観測を必要とする。

\subsection{6.1.3 分析環境}\label{ux5206ux6790ux74b0ux5883-3}

\begin{Shaded}
\begin{Highlighting}[]
\FunctionTok{source}\NormalTok{(}\StringTok{"R/utility\_LFL.R"}\NormalTok{)}

\FunctionTok{suppressPackageStartupMessages}\NormalTok{(\{}
  \FunctionTok{library}\NormalTok{(tidyverse)}
  \FunctionTok{library}\NormalTok{(broom)}
  \FunctionTok{library}\NormalTok{(ggthemes)}
\NormalTok{\})}

\NormalTok{estimation\_start }\OtherTok{\textless{}{-}} \SpecialCharTok{{-}}\DecValTok{130}\NormalTok{L}
\NormalTok{estimation\_end }\OtherTok{\textless{}{-}} \SpecialCharTok{{-}}\DecValTok{31}\NormalTok{L}
\NormalTok{event\_start }\OtherTok{\textless{}{-}} \SpecialCharTok{{-}}\DecValTok{30}\NormalTok{L}
\NormalTok{event\_end }\OtherTok{\textless{}{-}} \DecValTok{30}\NormalTok{L}
\NormalTok{minimum\_estimation\_observations }\OtherTok{\textless{}{-}} \DecValTok{80}\NormalTok{L}

\NormalTok{tbl\_window\_parameter }\OtherTok{\textless{}{-}} \FunctionTok{tribble}\NormalTok{(}
  \SpecialCharTok{\textasciitilde{}}\NormalTok{parameter, }\SpecialCharTok{\textasciitilde{}}\NormalTok{value,}
  \StringTok{"estimation\_start"}\NormalTok{, estimation\_start,}
  \StringTok{"estimation\_end"}\NormalTok{, estimation\_end,}
  \StringTok{"event\_start"}\NormalTok{, event\_start,}
  \StringTok{"event\_end"}\NormalTok{, event\_end,}
  \StringTok{"minimum\_estimation\_observations"}\NormalTok{, minimum\_estimation\_observations}
\NormalTok{)}

\NormalTok{tbl\_window\_parameter}
\end{Highlighting}
\end{Shaded}

\begin{verbatim}
# A tibble: 5 x 2
  parameter                       value
  <chr>                           <int>
1 estimation_start                 -130
2 estimation_end                    -31
3 event_start                       -30
4 event_end                          30
5 minimum_estimation_observations    80
\end{verbatim}

プロジェクトルートとデータディレクトリを確認する。

\begin{Shaded}
\begin{Highlighting}[]
\NormalTok{vec\_project\_candidate }\OtherTok{\textless{}{-}} \FunctionTok{unique}\NormalTok{(}\FunctionTok{c}\NormalTok{(}
  \FunctionTok{Sys.getenv}\NormalTok{(}\StringTok{"QUARTO\_PROJECT\_DIR"}\NormalTok{),}
  \FunctionTok{getwd}\NormalTok{(),}
  \FunctionTok{dirname}\NormalTok{(}\FunctionTok{getwd}\NormalTok{())}
\NormalTok{))}

\NormalTok{vec\_project\_candidate }\OtherTok{\textless{}{-}}\NormalTok{ vec\_project\_candidate[}
  \FunctionTok{nzchar}\NormalTok{(vec\_project\_candidate)}
\NormalTok{]}

\NormalTok{vec\_project\_match }\OtherTok{\textless{}{-}}\NormalTok{ vec\_project\_candidate[}
  \FunctionTok{file.exists}\NormalTok{(}\FunctionTok{file.path}\NormalTok{(vec\_project\_candidate, }\StringTok{"\_quarto.yml"}\NormalTok{))}
\NormalTok{]}

\ControlFlowTok{if}\NormalTok{ (}\FunctionTok{length}\NormalTok{(vec\_project\_match) }\SpecialCharTok{==} \DecValTok{0}\NormalTok{L) \{}
  \FunctionTok{stop}\NormalTok{(}\StringTok{"LagrangeFinanceのプロジェクトルートを確認できません。"}\NormalTok{)}
\NormalTok{\}}

\NormalTok{project\_dir }\OtherTok{\textless{}{-}} \FunctionTok{normalizePath}\NormalTok{(}
\NormalTok{  vec\_project\_match[[}\DecValTok{1}\NormalTok{]],}
  \AttributeTok{winslash =} \StringTok{"/"}\NormalTok{,}
  \AttributeTok{mustWork =} \ConstantTok{TRUE}
\NormalTok{)}

\NormalTok{simulation\_data\_dir }\OtherTok{\textless{}{-}} \FunctionTok{file.path}\NormalTok{(}
\NormalTok{  project\_dir,}
  \StringTok{"simulation\_data"}
\NormalTok{)}

\NormalTok{derived\_data\_dir }\OtherTok{\textless{}{-}} \FunctionTok{file.path}\NormalTok{(}
\NormalTok{  simulation\_data\_dir,}
  \StringTok{"derived"}
\NormalTok{)}

\FunctionTok{dir.create}\NormalTok{(}
\NormalTok{  derived\_data\_dir,}
  \AttributeTok{recursive =} \ConstantTok{TRUE}\NormalTok{,}
  \AttributeTok{showWarnings =} \ConstantTok{FALSE}
\NormalTok{)}

\FunctionTok{tibble}\NormalTok{(}
  \AttributeTok{directory =} \FunctionTok{c}\NormalTok{(}
    \StringTok{"project"}\NormalTok{,}
    \StringTok{"simulation\_data"}\NormalTok{,}
    \StringTok{"derived"}
\NormalTok{  ),}
  \AttributeTok{path =} \FunctionTok{c}\NormalTok{(}
\NormalTok{    project\_dir,}
\NormalTok{    simulation\_data\_dir,}
\NormalTok{    derived\_data\_dir}
\NormalTok{  )}
\NormalTok{)}
\end{Highlighting}
\end{Shaded}

\begin{verbatim}
# A tibble: 3 x 2
  directory       path                                                          
  <chr>           <chr>                                                         
1 project         C:/AnalyticFin/Projects/LagrangeFinance                       
2 simulation_data C:/AnalyticFin/Projects/LagrangeFinance/simulation_data       
3 derived         C:/AnalyticFin/Projects/LagrangeFinance/simulation_data/deriv~
\end{verbatim}

\section{6.2
入力データの読み込みと検証}\label{ux5165ux529bux30c7ux30fcux30bfux306eux8aadux307fux8fbcux307fux3068ux691cux8a3c}

\subsection{6.2.1
入力ファイル}\label{ux5165ux529bux30d5ux30a1ux30a4ux30eb}

本章では次の3ファイルを使用する。

\begin{Shaded}
\begin{Highlighting}[]
\NormalTok{ch08\_return\_data.csv}
\NormalTok{ch08\_market\_return\_data.csv}
\NormalTok{ch08\_event\_data.csv}
\end{Highlighting}
\end{Shaded}

ファイル名は旧教材の章番号に由来するが、本書では第6章の入力データとして使用する。

\begin{Shaded}
\begin{Highlighting}[]
\NormalTok{vec\_input\_file }\OtherTok{\textless{}{-}} \FunctionTok{c}\NormalTok{(}
  \AttributeTok{return =} \StringTok{"ch08\_return\_data.csv"}\NormalTok{,}
  \AttributeTok{market =} \StringTok{"ch08\_market\_return\_data.csv"}\NormalTok{,}
  \AttributeTok{event =} \StringTok{"ch08\_event\_data.csv"}
\NormalTok{)}

\NormalTok{vec\_input\_path }\OtherTok{\textless{}{-}}\NormalTok{ purrr}\SpecialCharTok{::}\FunctionTok{map\_chr}\NormalTok{(}
\NormalTok{  vec\_input\_file,}
\NormalTok{  find\_data\_file\_LFL,}
  \AttributeTok{search\_dir =}\NormalTok{ simulation\_data\_dir}
\NormalTok{)}

\FunctionTok{tibble}\NormalTok{(}
  \AttributeTok{data\_name =} \FunctionTok{names}\NormalTok{(vec\_input\_path),}
  \AttributeTok{file\_name =} \FunctionTok{basename}\NormalTok{(vec\_input\_path),}
  \AttributeTok{path =} \FunctionTok{unname}\NormalTok{(vec\_input\_path)}
\NormalTok{)}
\end{Highlighting}
\end{Shaded}

\begin{verbatim}
# A tibble: 3 x 3
  data_name file_name                   path                                    
  <chr>     <chr>                       <chr>                                   
1 return    ch08_return_data.csv        C:/AnalyticFin/Projects/LagrangeFinance~
2 market    ch08_market_return_data.csv C:/AnalyticFin/Projects/LagrangeFinance~
3 event     ch08_event_data.csv         C:/AnalyticFin/Projects/LagrangeFinance~
\end{verbatim}

\subsection{6.2.2
データの読み込み}\label{ux30c7ux30fcux30bfux306eux8aadux307fux8fbcux307f}

\begin{Shaded}
\begin{Highlighting}[]
\NormalTok{tbl\_return\_raw }\OtherTok{\textless{}{-}}\NormalTok{ readr}\SpecialCharTok{::}\FunctionTok{read\_csv}\NormalTok{(}
\NormalTok{  vec\_input\_path[[}\StringTok{"return"}\NormalTok{]],}
  \AttributeTok{show\_col\_types =} \ConstantTok{FALSE}
\NormalTok{)}

\NormalTok{tbl\_market\_raw }\OtherTok{\textless{}{-}}\NormalTok{ readr}\SpecialCharTok{::}\FunctionTok{read\_csv}\NormalTok{(}
\NormalTok{  vec\_input\_path[[}\StringTok{"market"}\NormalTok{]],}
  \AttributeTok{show\_col\_types =} \ConstantTok{FALSE}
\NormalTok{)}

\NormalTok{tbl\_event\_raw }\OtherTok{\textless{}{-}}\NormalTok{ readr}\SpecialCharTok{::}\FunctionTok{read\_csv}\NormalTok{(}
\NormalTok{  vec\_input\_path[[}\StringTok{"event"}\NormalTok{]],}
  \AttributeTok{show\_col\_types =} \ConstantTok{FALSE}
\NormalTok{)}
\end{Highlighting}
\end{Shaded}

\subsection{6.2.3
必須列の確認}\label{ux5fc5ux9808ux5217ux306eux78baux8a8d}

\begin{Shaded}
\begin{Highlighting}[]
\FunctionTok{validate\_required\_columns\_LFL}\NormalTok{(}
\NormalTok{  tbl\_return\_raw,}
  \FunctionTok{c}\NormalTok{(}\StringTok{"date"}\NormalTok{, }\StringTok{"firm\_ID"}\NormalTok{, }\StringTok{"R"}\NormalTok{),}
  \StringTok{"tbl\_return\_raw"}
\NormalTok{)}

\FunctionTok{validate\_required\_columns\_LFL}\NormalTok{(}
\NormalTok{  tbl\_market\_raw,}
  \FunctionTok{c}\NormalTok{(}\StringTok{"date"}\NormalTok{, }\StringTok{"R\_M"}\NormalTok{),}
  \StringTok{"tbl\_market\_raw"}
\NormalTok{)}

\FunctionTok{validate\_required\_columns\_LFL}\NormalTok{(}
\NormalTok{  tbl\_event\_raw,}
  \FunctionTok{c}\NormalTok{(}
    \StringTok{"event\_date"}\NormalTok{,}
    \StringTok{"firm\_ID"}\NormalTok{,}
    \StringTok{"earnings\_forecast"}\NormalTok{,}
    \StringTok{"realized\_earnings"}\NormalTok{,}
    \StringTok{"lagged\_ME"}
\NormalTok{  ),}
  \StringTok{"tbl\_event\_raw"}
\NormalTok{)}
\end{Highlighting}
\end{Shaded}

\subsection{6.2.4
データ型の整理}\label{ux30c7ux30fcux30bfux578bux306eux6574ux7406}

\begin{Shaded}
\begin{Highlighting}[]
\NormalTok{tbl\_return }\OtherTok{\textless{}{-}}\NormalTok{ tbl\_return\_raw }\SpecialCharTok{|\textgreater{}}
  \FunctionTok{transmute}\NormalTok{(}
    \AttributeTok{date =} \FunctionTok{as.Date}\NormalTok{(date),}
    \AttributeTok{firm\_ID =} \FunctionTok{as.integer}\NormalTok{(firm\_ID),}
    \AttributeTok{R =} \FunctionTok{as.numeric}\NormalTok{(R)}
\NormalTok{  ) }\SpecialCharTok{|\textgreater{}}
  \FunctionTok{arrange}\NormalTok{(}
\NormalTok{    firm\_ID,}
\NormalTok{    date}
\NormalTok{  )}

\NormalTok{tbl\_market }\OtherTok{\textless{}{-}}\NormalTok{ tbl\_market\_raw }\SpecialCharTok{|\textgreater{}}
  \FunctionTok{transmute}\NormalTok{(}
    \AttributeTok{date =} \FunctionTok{as.Date}\NormalTok{(date),}
    \AttributeTok{R\_M =} \FunctionTok{as.numeric}\NormalTok{(R\_M)}
\NormalTok{  ) }\SpecialCharTok{|\textgreater{}}
  \FunctionTok{arrange}\NormalTok{(date)}

\NormalTok{tbl\_event }\OtherTok{\textless{}{-}}\NormalTok{ tbl\_event\_raw }\SpecialCharTok{|\textgreater{}}
  \FunctionTok{transmute}\NormalTok{(}
    \AttributeTok{event\_date =} \FunctionTok{as.Date}\NormalTok{(event\_date),}
    \AttributeTok{firm\_ID =} \FunctionTok{as.integer}\NormalTok{(firm\_ID),}
    \AttributeTok{earnings\_forecast =} \FunctionTok{as.numeric}\NormalTok{(}
\NormalTok{      earnings\_forecast}
\NormalTok{    ),}
    \AttributeTok{realized\_earnings =} \FunctionTok{as.numeric}\NormalTok{(}
\NormalTok{      realized\_earnings}
\NormalTok{    ),}
    \AttributeTok{lagged\_ME =} \FunctionTok{as.numeric}\NormalTok{(lagged\_ME)}
\NormalTok{  ) }\SpecialCharTok{|\textgreater{}}
  \FunctionTok{arrange}\NormalTok{(}
\NormalTok{    event\_date,}
\NormalTok{    firm\_ID}
\NormalTok{  )}
\end{Highlighting}
\end{Shaded}

\subsection{6.2.5
キーと欠損値の検証}\label{ux30adux30fcux3068ux6b20ux640dux5024ux306eux691cux8a3c}

\begin{Shaded}
\begin{Highlighting}[]
\FunctionTok{validate\_unique\_keys\_LFL}\NormalTok{(}
\NormalTok{  tbl\_return,}
  \FunctionTok{c}\NormalTok{(}\StringTok{"firm\_ID"}\NormalTok{, }\StringTok{"date"}\NormalTok{),}
  \StringTok{"tbl\_return"}
\NormalTok{)}

\FunctionTok{validate\_unique\_keys\_LFL}\NormalTok{(}
\NormalTok{  tbl\_market,}
  \StringTok{"date"}\NormalTok{,}
  \StringTok{"tbl\_market"}
\NormalTok{)}

\FunctionTok{validate\_unique\_keys\_LFL}\NormalTok{(}
\NormalTok{  tbl\_event,}
  \FunctionTok{c}\NormalTok{(}\StringTok{"firm\_ID"}\NormalTok{, }\StringTok{"event\_date"}\NormalTok{),}
  \StringTok{"tbl\_event"}
\NormalTok{)}
\end{Highlighting}
\end{Shaded}

\begin{Shaded}
\begin{Highlighting}[]
\NormalTok{tbl\_missing\_summary }\OtherTok{\textless{}{-}} \FunctionTok{bind\_rows}\NormalTok{(}
  \FunctionTok{summarize\_missing\_values\_LFL}\NormalTok{(}
\NormalTok{    tbl\_return,}
    \StringTok{"return"}
\NormalTok{  ),}
  \FunctionTok{summarize\_missing\_values\_LFL}\NormalTok{(}
\NormalTok{    tbl\_market,}
    \StringTok{"market"}
\NormalTok{  ),}
  \FunctionTok{summarize\_missing\_values\_LFL}\NormalTok{(}
\NormalTok{    tbl\_event,}
    \StringTok{"event"}
\NormalTok{  )}
\NormalTok{)}

\NormalTok{tbl\_missing\_summary}
\end{Highlighting}
\end{Shaded}

\begin{verbatim}
# A tibble: 10 x 5
   data_name observations missing_ratio variable          missing_values
   <chr>            <int>         <dbl> <chr>                      <int>
 1 return         1610400             0 date                           0
 2 return         1610400             0 firm_ID                        0
 3 return         1610400             0 R                              0
 4 market            1220             0 date                           0
 5 market            1220             0 R_M                            0
 6 event             3960             0 event_date                     0
 7 event             3960             0 firm_ID                        0
 8 event             3960             0 earnings_forecast              0
 9 event             3960             0 realized_earnings              0
10 event             3960             0 lagged_ME                      0
\end{verbatim}

\subsection{6.2.6
入力データの概要}\label{ux5165ux529bux30c7ux30fcux30bfux306eux6982ux8981}

\begin{Shaded}
\begin{Highlighting}[]
\NormalTok{tbl\_input\_overview }\OtherTok{\textless{}{-}} \FunctionTok{tribble}\NormalTok{(}
  \SpecialCharTok{\textasciitilde{}}\NormalTok{data\_name, }\SpecialCharTok{\textasciitilde{}}\NormalTok{observations, }\SpecialCharTok{\textasciitilde{}}\NormalTok{firms, }\SpecialCharTok{\textasciitilde{}}\NormalTok{first\_date, }\SpecialCharTok{\textasciitilde{}}\NormalTok{last\_date,}
  \StringTok{"individual returns"}\NormalTok{,}
  \FunctionTok{nrow}\NormalTok{(tbl\_return),}
  \FunctionTok{n\_distinct}\NormalTok{(tbl\_return}\SpecialCharTok{$}\NormalTok{firm\_ID),}
  \FunctionTok{min}\NormalTok{(tbl\_return}\SpecialCharTok{$}\NormalTok{date, }\AttributeTok{na.rm =} \ConstantTok{TRUE}\NormalTok{),}
  \FunctionTok{max}\NormalTok{(tbl\_return}\SpecialCharTok{$}\NormalTok{date, }\AttributeTok{na.rm =} \ConstantTok{TRUE}\NormalTok{),}
  \StringTok{"market returns"}\NormalTok{,}
  \FunctionTok{nrow}\NormalTok{(tbl\_market),}
  \ConstantTok{NA\_integer\_}\NormalTok{,}
  \FunctionTok{min}\NormalTok{(tbl\_market}\SpecialCharTok{$}\NormalTok{date, }\AttributeTok{na.rm =} \ConstantTok{TRUE}\NormalTok{),}
  \FunctionTok{max}\NormalTok{(tbl\_market}\SpecialCharTok{$}\NormalTok{date, }\AttributeTok{na.rm =} \ConstantTok{TRUE}\NormalTok{),}
  \StringTok{"events"}\NormalTok{,}
  \FunctionTok{nrow}\NormalTok{(tbl\_event),}
  \FunctionTok{n\_distinct}\NormalTok{(tbl\_event}\SpecialCharTok{$}\NormalTok{firm\_ID),}
  \FunctionTok{min}\NormalTok{(tbl\_event}\SpecialCharTok{$}\NormalTok{event\_date, }\AttributeTok{na.rm =} \ConstantTok{TRUE}\NormalTok{),}
  \FunctionTok{max}\NormalTok{(tbl\_event}\SpecialCharTok{$}\NormalTok{event\_date, }\AttributeTok{na.rm =} \ConstantTok{TRUE}\NormalTok{)}
\NormalTok{)}

\NormalTok{tbl\_input\_overview}
\end{Highlighting}
\end{Shaded}

\begin{verbatim}
# A tibble: 3 x 5
  data_name          observations firms first_date last_date 
  <chr>                     <int> <int> <date>     <date>    
1 individual returns      1610400  1320 2016-01-04 2020-12-30
2 market returns             1220    NA 2016-01-04 2020-12-30
3 events                     3960  1320 2017-01-04 2020-04-07
\end{verbatim}

\subsection{6.2.7
個別株リターンの分布}\label{ux500bux5225ux682aux30eaux30bfux30fcux30f3ux306eux5206ux5e03}

\begin{Shaded}
\begin{Highlighting}[]
\NormalTok{plot\_return\_distribution }\OtherTok{\textless{}{-}}\NormalTok{ tbl\_return }\SpecialCharTok{|\textgreater{}}
  \FunctionTok{filter}\NormalTok{(}\FunctionTok{is.finite}\NormalTok{(R)) }\SpecialCharTok{|\textgreater{}}
  \FunctionTok{ggplot}\NormalTok{(}\FunctionTok{aes}\NormalTok{(}\AttributeTok{x =}\NormalTok{ R)) }\SpecialCharTok{+}
  \FunctionTok{geom\_histogram}\NormalTok{(}\AttributeTok{bins =} \DecValTok{80}\NormalTok{) }\SpecialCharTok{+}
  \FunctionTok{geom\_vline}\NormalTok{(}
    \AttributeTok{xintercept =} \DecValTok{0}\NormalTok{,}
    \AttributeTok{linetype =} \StringTok{"dotted"}
\NormalTok{  ) }\SpecialCharTok{+}
  \FunctionTok{scale\_x\_continuous}\NormalTok{(}
    \AttributeTok{labels =}\NormalTok{ scales}\SpecialCharTok{::}\FunctionTok{label\_percent}\NormalTok{()}
\NormalTok{  ) }\SpecialCharTok{+}
  \FunctionTok{labs}\NormalTok{(}
    \AttributeTok{x =} \StringTok{"Daily return"}\NormalTok{,}
    \AttributeTok{y =} \StringTok{"Observations"}\NormalTok{,}
    \AttributeTok{title =} \StringTok{"Distribution of Individual Stock Returns"}
\NormalTok{  ) }\SpecialCharTok{+}
  \FunctionTok{theme\_classic}\NormalTok{()}

\NormalTok{plot\_return\_distribution}
\end{Highlighting}
\end{Shaded}

\pandocbounded{\includegraphics[keepaspectratio]{part1_empirical/ch06_event_study_files/figure-pdf/plot-return-distribution-1.pdf}}

\subsection{6.2.8
市場リターンの時系列}\label{ux5e02ux5834ux30eaux30bfux30fcux30f3ux306eux6642ux7cfbux5217}

\begin{Shaded}
\begin{Highlighting}[]
\NormalTok{plot\_market\_return\_series }\OtherTok{\textless{}{-}}\NormalTok{ tbl\_market }\SpecialCharTok{|\textgreater{}}
  \FunctionTok{ggplot}\NormalTok{(}\FunctionTok{aes}\NormalTok{(}\AttributeTok{x =}\NormalTok{ date, }\AttributeTok{y =}\NormalTok{ R\_M)) }\SpecialCharTok{+}
  \FunctionTok{geom\_line}\NormalTok{(}\AttributeTok{linewidth =} \FloatTok{0.4}\NormalTok{) }\SpecialCharTok{+}
  \FunctionTok{geom\_hline}\NormalTok{(}
    \AttributeTok{yintercept =} \DecValTok{0}\NormalTok{,}
    \AttributeTok{linetype =} \StringTok{"dotted"}
\NormalTok{  ) }\SpecialCharTok{+}
  \FunctionTok{scale\_y\_continuous}\NormalTok{(}
    \AttributeTok{labels =}\NormalTok{ scales}\SpecialCharTok{::}\FunctionTok{label\_percent}\NormalTok{()}
\NormalTok{  ) }\SpecialCharTok{+}
  \FunctionTok{labs}\NormalTok{(}
    \AttributeTok{x =} \ConstantTok{NULL}\NormalTok{,}
    \AttributeTok{y =} \StringTok{"Market return"}\NormalTok{,}
    \AttributeTok{title =} \StringTok{"Daily Market Returns"}
\NormalTok{  ) }\SpecialCharTok{+}
  \FunctionTok{theme\_classic}\NormalTok{()}

\NormalTok{plot\_market\_return\_series}
\end{Highlighting}
\end{Shaded}

\pandocbounded{\includegraphics[keepaspectratio]{part1_empirical/ch06_event_study_files/figure-pdf/plot-market-return-series-1.pdf}}

\section{6.3
決算サプライズとイベント分類}\label{ux6c7aux7b97ux30b5ux30d7ux30e9ux30a4ux30baux3068ux30a4ux30d9ux30f3ux30c8ux5206ux985e}

\subsection{6.3.1
決算サプライズ}\label{ux6c7aux7b97ux30b5ux30d7ux30e9ux30a4ux30ba}

決算サプライズは、実現利益と予想利益の差を期首時価総額で標準化して計算する。

\[
Surprise_i
=
\frac{
RealizedEarnings_i
-
ForecastEarnings_i
}{
ME_{i,-1}
}
\]

正の値は実現利益が予想を上回ったことを、負の値は予想を下回ったことを表す。

\begin{Shaded}
\begin{Highlighting}[]
\NormalTok{tbl\_event }\OtherTok{\textless{}{-}} \FunctionTok{compute\_earnings\_surprise\_LFL}\NormalTok{(}
  \AttributeTok{tbl\_event =}\NormalTok{ tbl\_event,}
  \AttributeTok{number\_of\_groups =} \DecValTok{5}\NormalTok{L}
\NormalTok{)}

\NormalTok{tbl\_event }\SpecialCharTok{|\textgreater{}}
  \FunctionTok{select}\NormalTok{(}
\NormalTok{    event\_ID,}
\NormalTok{    firm\_ID,}
\NormalTok{    event\_date,}
\NormalTok{    event\_year,}
\NormalTok{    earnings\_surprise,}
\NormalTok{    event\_strength}
\NormalTok{  ) }\SpecialCharTok{|\textgreater{}}
  \FunctionTok{slice\_head}\NormalTok{(}\AttributeTok{n =} \DecValTok{20}\NormalTok{)}
\end{Highlighting}
\end{Shaded}

\begin{verbatim}
# A tibble: 20 x 6
   event_ID firm_ID event_date event_year earnings_surprise event_strength
      <int>   <int> <date>          <dbl>             <dbl> <ord>         
 1        1      39 2017-01-04       2017         -0.00813  1             
 2        2      57 2017-01-04       2017         -0.0100   1             
 3        3      73 2017-01-04       2017         -0.00202  4             
 4        4     109 2017-01-04       2017         -0.00440  3             
 5        5     111 2017-01-04       2017         -0.00449  3             
 6        6     231 2017-01-04       2017         -0.0144   1             
 7        7     944 2017-01-04       2017         -0.00690  2             
 8        8    1227 2017-01-04       2017         -0.00112  5             
 9        9     422 2017-01-05       2017         -0.00484  3             
10       10    1203 2017-01-05       2017         -0.000866 5             
11       11     346 2017-01-06       2017         -0.00586  2             
12       12     590 2017-01-06       2017         -0.0103   1             
13       13     644 2017-01-06       2017         -0.00649  2             
14       14    1153 2017-01-06       2017         -0.00210  4             
15       15     191 2017-01-10       2017         -0.00347  4             
16       16     353 2017-01-10       2017         -0.0102   1             
17       17     566 2017-01-10       2017         -0.00439  3             
18       18     322 2017-01-11       2017         -0.00469  3             
19       19    1050 2017-01-11       2017         -0.00250  4             
20       20     426 2017-01-12       2017         -0.00101  5             
\end{verbatim}

\subsection{6.3.2
イベント強度別の概要}\label{ux30a4ux30d9ux30f3ux30c8ux5f37ux5ea6ux5225ux306eux6982ux8981}

\begin{Shaded}
\begin{Highlighting}[]
\NormalTok{tbl\_event\_strength\_summary }\OtherTok{\textless{}{-}}\NormalTok{ tbl\_event }\SpecialCharTok{|\textgreater{}}
  \FunctionTok{group\_by}\NormalTok{(}
\NormalTok{    event\_year,}
\NormalTok{    event\_strength}
\NormalTok{  ) }\SpecialCharTok{|\textgreater{}}
  \FunctionTok{summarise}\NormalTok{(}
    \AttributeTok{events =} \FunctionTok{n}\NormalTok{(),}
    \AttributeTok{mean\_surprise =} \FunctionTok{mean}\NormalTok{(}
\NormalTok{      earnings\_surprise,}
      \AttributeTok{na.rm =} \ConstantTok{TRUE}
\NormalTok{    ),}
    \AttributeTok{median\_surprise =} \FunctionTok{median}\NormalTok{(}
\NormalTok{      earnings\_surprise,}
      \AttributeTok{na.rm =} \ConstantTok{TRUE}
\NormalTok{    ),}
    \AttributeTok{.groups =} \StringTok{"drop"}
\NormalTok{  ) }\SpecialCharTok{|\textgreater{}}
  \FunctionTok{arrange}\NormalTok{(}
\NormalTok{    event\_year,}
\NormalTok{    event\_strength}
\NormalTok{  )}

\NormalTok{tbl\_event\_strength\_summary}
\end{Highlighting}
\end{Shaded}

\begin{verbatim}
# A tibble: 20 x 5
   event_year event_strength events mean_surprise median_surprise
        <dbl> <ord>           <int>         <dbl>           <dbl>
 1       2017 1                 214     -0.0103         -0.00983 
 2       2017 2                 214     -0.00643        -0.00639 
 3       2017 3                 214     -0.00442        -0.00437 
 4       2017 4                 213     -0.00287        -0.00289 
 5       2017 5                 213     -0.000622       -0.000818
 6       2018 1                 263     -0.00994        -0.00926 
 7       2018 2                 263     -0.00669        -0.00664 
 8       2018 3                 263     -0.00488        -0.00487 
 9       2018 4                 263     -0.00319        -0.00318 
10       2018 5                 263     -0.000985       -0.00117 
11       2019 1                 262     -0.00980        -0.00940 
12       2019 2                 262     -0.00646        -0.00642 
13       2019 3                 261     -0.00459        -0.00459 
14       2019 4                 261     -0.00296        -0.00296 
15       2019 5                 261     -0.000887       -0.00108 
16       2020 1                  54     -0.0107         -0.0103  
17       2020 2                  54     -0.00651        -0.00652 
18       2020 3                  54     -0.00436        -0.00444 
19       2020 4                  54     -0.00267        -0.00263 
20       2020 5                  54     -0.000511       -0.000696
\end{verbatim}

\subsection{6.3.3
決算サプライズの分布}\label{ux6c7aux7b97ux30b5ux30d7ux30e9ux30a4ux30baux306eux5206ux5e03}

\begin{Shaded}
\begin{Highlighting}[]
\NormalTok{plot\_earnings\_surprise\_distribution }\OtherTok{\textless{}{-}}\NormalTok{ tbl\_event }\SpecialCharTok{|\textgreater{}}
  \FunctionTok{filter}\NormalTok{(}\FunctionTok{is.finite}\NormalTok{(earnings\_surprise)) }\SpecialCharTok{|\textgreater{}}
  \FunctionTok{ggplot}\NormalTok{(}
    \FunctionTok{aes}\NormalTok{(}
      \AttributeTok{x =}\NormalTok{ earnings\_surprise,}
      \AttributeTok{linetype =}\NormalTok{ event\_strength}
\NormalTok{    )}
\NormalTok{  ) }\SpecialCharTok{+}
  \FunctionTok{geom\_density}\NormalTok{(}\AttributeTok{linewidth =} \FloatTok{0.8}\NormalTok{) }\SpecialCharTok{+}
  \FunctionTok{geom\_vline}\NormalTok{(}
    \AttributeTok{xintercept =} \DecValTok{0}\NormalTok{,}
    \AttributeTok{linetype =} \StringTok{"dotted"}
\NormalTok{  ) }\SpecialCharTok{+}
  \FunctionTok{labs}\NormalTok{(}
    \AttributeTok{x =} \StringTok{"Standardized earnings surprise"}\NormalTok{,}
    \AttributeTok{y =} \StringTok{"Density"}\NormalTok{,}
    \AttributeTok{linetype =} \StringTok{"Event strength"}\NormalTok{,}
    \AttributeTok{title =} \StringTok{"Distribution of Earnings Surprise"}
\NormalTok{  ) }\SpecialCharTok{+}
  \FunctionTok{theme\_classic}\NormalTok{()}

\NormalTok{plot\_earnings\_surprise\_distribution}
\end{Highlighting}
\end{Shaded}

\pandocbounded{\includegraphics[keepaspectratio]{part1_empirical/ch06_event_study_files/figure-pdf/plot-earnings-surprise-distribution-1.pdf}}

\subsection{6.3.4
イベント強度別の箱ひげ図}\label{ux30a4ux30d9ux30f3ux30c8ux5f37ux5ea6ux5225ux306eux7bb1ux3072ux3052ux56f3}

\begin{Shaded}
\begin{Highlighting}[]
\NormalTok{plot\_surprise\_by\_strength }\OtherTok{\textless{}{-}}\NormalTok{ tbl\_event }\SpecialCharTok{|\textgreater{}}
  \FunctionTok{filter}\NormalTok{(}\FunctionTok{is.finite}\NormalTok{(earnings\_surprise)) }\SpecialCharTok{|\textgreater{}}
  \FunctionTok{ggplot}\NormalTok{(}
    \FunctionTok{aes}\NormalTok{(}
      \AttributeTok{x =}\NormalTok{ event\_strength,}
      \AttributeTok{y =}\NormalTok{ earnings\_surprise}
\NormalTok{    )}
\NormalTok{  ) }\SpecialCharTok{+}
  \FunctionTok{geom\_boxplot}\NormalTok{(}
    \AttributeTok{outlier.alpha =} \FloatTok{0.25}
\NormalTok{  ) }\SpecialCharTok{+}
  \FunctionTok{geom\_hline}\NormalTok{(}
    \AttributeTok{yintercept =} \DecValTok{0}\NormalTok{,}
    \AttributeTok{linetype =} \StringTok{"dotted"}
\NormalTok{  ) }\SpecialCharTok{+}
  \FunctionTok{labs}\NormalTok{(}
    \AttributeTok{x =} \StringTok{"Event strength"}\NormalTok{,}
    \AttributeTok{y =} \StringTok{"Standardized earnings surprise"}\NormalTok{,}
    \AttributeTok{title =} \StringTok{"Earnings Surprise by Event{-}Strength Group"}
\NormalTok{  ) }\SpecialCharTok{+}
  \FunctionTok{theme\_classic}\NormalTok{()}

\NormalTok{plot\_surprise\_by\_strength}
\end{Highlighting}
\end{Shaded}

\pandocbounded{\includegraphics[keepaspectratio]{part1_empirical/ch06_event_study_files/figure-pdf/plot-surprise-by-strength-1.pdf}}

\subsection{6.3.5
年度別イベント件数}\label{ux5e74ux5ea6ux5225ux30a4ux30d9ux30f3ux30c8ux4ef6ux6570}

\begin{Shaded}
\begin{Highlighting}[]
\NormalTok{plot\_events\_by\_year }\OtherTok{\textless{}{-}}\NormalTok{ tbl\_event }\SpecialCharTok{|\textgreater{}}
  \FunctionTok{count}\NormalTok{(}
\NormalTok{    event\_year,}
\NormalTok{    event\_strength,}
    \AttributeTok{name =} \StringTok{"events"}
\NormalTok{  ) }\SpecialCharTok{|\textgreater{}}
  \FunctionTok{ggplot}\NormalTok{(}
    \FunctionTok{aes}\NormalTok{(}
      \AttributeTok{x =} \FunctionTok{factor}\NormalTok{(event\_year),}
      \AttributeTok{y =}\NormalTok{ events,}
      \AttributeTok{fill =}\NormalTok{ event\_strength}
\NormalTok{    )}
\NormalTok{  ) }\SpecialCharTok{+}
  \FunctionTok{geom\_col}\NormalTok{(}
    \AttributeTok{position =} \StringTok{"dodge"}
\NormalTok{  ) }\SpecialCharTok{+}
  \FunctionTok{labs}\NormalTok{(}
    \AttributeTok{x =} \StringTok{"Event year"}\NormalTok{,}
    \AttributeTok{y =} \StringTok{"Events"}\NormalTok{,}
    \AttributeTok{fill =} \StringTok{"Event strength"}\NormalTok{,}
    \AttributeTok{title =} \StringTok{"Number of Earnings Events by Year"}
\NormalTok{  ) }\SpecialCharTok{+}
  \FunctionTok{theme\_classic}\NormalTok{()}

\NormalTok{plot\_events\_by\_year}
\end{Highlighting}
\end{Shaded}

\pandocbounded{\includegraphics[keepaspectratio]{part1_empirical/ch06_event_study_files/figure-pdf/plot-events-by-year-1.pdf}}

\section{6.4
イベント日とイベント・パネル}\label{ux30a4ux30d9ux30f3ux30c8ux65e5ux3068ux30a4ux30d9ux30f3ux30c8ux30d1ux30cdux30eb}

\subsection{6.4.1
取引日カレンダー}\label{ux53d6ux5f15ux65e5ux30abux30ecux30f3ux30c0ux30fc}

市場リターンの日付を取引日カレンダーとして使用する。

\begin{Shaded}
\begin{Highlighting}[]
\NormalTok{tbl\_trading\_calendar }\OtherTok{\textless{}{-}}\NormalTok{ tbl\_market }\SpecialCharTok{|\textgreater{}}
  \FunctionTok{distinct}\NormalTok{(date) }\SpecialCharTok{|\textgreater{}}
  \FunctionTok{arrange}\NormalTok{(date) }\SpecialCharTok{|\textgreater{}}
  \FunctionTok{mutate}\NormalTok{(}
    \AttributeTok{date\_ID =} \FunctionTok{row\_number}\NormalTok{()}
\NormalTok{  ) }\SpecialCharTok{|\textgreater{}}
  \FunctionTok{select}\NormalTok{(}
\NormalTok{    date\_ID,}
\NormalTok{    date}
\NormalTok{  )}

\NormalTok{tbl\_market\_with\_ID }\OtherTok{\textless{}{-}}\NormalTok{ tbl\_market }\SpecialCharTok{|\textgreater{}}
  \FunctionTok{inner\_join}\NormalTok{(}
\NormalTok{    tbl\_trading\_calendar,}
    \AttributeTok{by =} \StringTok{"date"}
\NormalTok{  ) }\SpecialCharTok{|\textgreater{}}
  \FunctionTok{select}\NormalTok{(}
\NormalTok{    date\_ID,}
\NormalTok{    date,}
\NormalTok{    R\_M}
\NormalTok{  )}

\NormalTok{tbl\_return\_market }\OtherTok{\textless{}{-}}\NormalTok{ tbl\_return }\SpecialCharTok{|\textgreater{}}
  \FunctionTok{inner\_join}\NormalTok{(}
\NormalTok{    tbl\_market\_with\_ID,}
    \AttributeTok{by =} \StringTok{"date"}
\NormalTok{  ) }\SpecialCharTok{|\textgreater{}}
  \FunctionTok{select}\NormalTok{(}
\NormalTok{    firm\_ID,}
\NormalTok{    date\_ID,}
\NormalTok{    R,}
\NormalTok{    R\_M}
\NormalTok{  ) }\SpecialCharTok{|\textgreater{}}
  \FunctionTok{arrange}\NormalTok{(}
\NormalTok{    firm\_ID,}
\NormalTok{    date\_ID}
\NormalTok{  )}
\end{Highlighting}
\end{Shaded}

\subsection{6.4.2 イベント日0}\label{ux30a4ux30d9ux30f3ux30c8ux65e50}

決算発表は取引終了後に行われると仮定する。そのため、発表日が取引日であっても、その次の取引日をイベント日0とする。

\begin{Shaded}
\begin{Highlighting}[]
\NormalTok{tbl\_event }\OtherTok{\textless{}{-}} \FunctionTok{assign\_event\_day\_zero\_LFL}\NormalTok{(}
  \AttributeTok{tbl\_event =}\NormalTok{ tbl\_event,}
  \AttributeTok{tbl\_trading\_calendar =}\NormalTok{ tbl\_trading\_calendar}
\NormalTok{)}

\NormalTok{tbl\_event }\SpecialCharTok{|\textgreater{}}
  \FunctionTok{select}\NormalTok{(}
\NormalTok{    event\_ID,}
\NormalTok{    firm\_ID,}
\NormalTok{    event\_date,}
\NormalTok{    event\_trade\_date,}
\NormalTok{    event\_trade\_date\_ID,}
\NormalTok{    event\_strength,}
\NormalTok{    earnings\_surprise}
\NormalTok{  ) }\SpecialCharTok{|\textgreater{}}
  \FunctionTok{slice\_head}\NormalTok{(}\AttributeTok{n =} \DecValTok{20}\NormalTok{)}
\end{Highlighting}
\end{Shaded}

\begin{verbatim}
# A tibble: 20 x 7
   event_ID firm_ID event_date event_trade_date event_trade_date_ID
      <int>   <int> <date>     <date>                         <int>
 1        1      39 2017-01-04 2017-01-05                       247
 2        2      57 2017-01-04 2017-01-05                       247
 3        3      73 2017-01-04 2017-01-05                       247
 4        4     109 2017-01-04 2017-01-05                       247
 5        5     111 2017-01-04 2017-01-05                       247
 6        6     231 2017-01-04 2017-01-05                       247
 7        7     944 2017-01-04 2017-01-05                       247
 8        8    1227 2017-01-04 2017-01-05                       247
 9        9     422 2017-01-05 2017-01-06                       248
10       10    1203 2017-01-05 2017-01-06                       248
11       11     346 2017-01-06 2017-01-10                       249
12       12     590 2017-01-06 2017-01-10                       249
13       13     644 2017-01-06 2017-01-10                       249
14       14    1153 2017-01-06 2017-01-10                       249
15       15     191 2017-01-10 2017-01-11                       250
16       16     353 2017-01-10 2017-01-11                       250
17       17     566 2017-01-10 2017-01-11                       250
18       18     322 2017-01-11 2017-01-12                       251
19       19    1050 2017-01-11 2017-01-12                       251
20       20     426 2017-01-12 2017-01-13                       252
# i 2 more variables: event_strength <ord>, earnings_surprise <dbl>
\end{verbatim}

\subsection{6.4.3
発表日からイベント日0までの日数}\label{ux767aux8868ux65e5ux304bux3089ux30a4ux30d9ux30f3ux30c8ux65e50ux307eux3067ux306eux65e5ux6570}

\begin{Shaded}
\begin{Highlighting}[]
\NormalTok{tbl\_event\_gap }\OtherTok{\textless{}{-}}\NormalTok{ tbl\_event }\SpecialCharTok{|\textgreater{}}
  \FunctionTok{mutate}\NormalTok{(}
    \AttributeTok{calendar\_day\_gap =} \FunctionTok{as.integer}\NormalTok{(}
\NormalTok{      event\_trade\_date }\SpecialCharTok{{-}}\NormalTok{ event\_date}
\NormalTok{    )}
\NormalTok{  )}

\NormalTok{plot\_event\_date\_gap }\OtherTok{\textless{}{-}}\NormalTok{ tbl\_event\_gap }\SpecialCharTok{|\textgreater{}}
  \FunctionTok{count}\NormalTok{(}
\NormalTok{    calendar\_day\_gap,}
    \AttributeTok{name =} \StringTok{"events"}
\NormalTok{  ) }\SpecialCharTok{|\textgreater{}}
  \FunctionTok{ggplot}\NormalTok{(}
    \FunctionTok{aes}\NormalTok{(}
      \AttributeTok{x =} \FunctionTok{factor}\NormalTok{(calendar\_day\_gap),}
      \AttributeTok{y =}\NormalTok{ events}
\NormalTok{    )}
\NormalTok{  ) }\SpecialCharTok{+}
  \FunctionTok{geom\_col}\NormalTok{() }\SpecialCharTok{+}
  \FunctionTok{labs}\NormalTok{(}
    \AttributeTok{x =} \StringTok{"Calendar days from announcement to event day 0"}\NormalTok{,}
    \AttributeTok{y =} \StringTok{"Events"}\NormalTok{,}
    \AttributeTok{title =} \StringTok{"Delay from Announcement Date to Event Day 0"}
\NormalTok{  ) }\SpecialCharTok{+}
  \FunctionTok{theme\_classic}\NormalTok{()}

\NormalTok{plot\_event\_date\_gap}
\end{Highlighting}
\end{Shaded}

\pandocbounded{\includegraphics[keepaspectratio]{part1_empirical/ch06_event_study_files/figure-pdf/plot-event-date-gap-1.pdf}}

\subsection{6.4.4
イベント・パネル}\label{ux30a4ux30d9ux30f3ux30c8ux30d1ux30cdux30eb}

\begin{Shaded}
\begin{Highlighting}[]
\NormalTok{tbl\_event\_panel }\OtherTok{\textless{}{-}} \FunctionTok{create\_event\_study\_panel\_LFL}\NormalTok{(}
  \AttributeTok{tbl\_event =}\NormalTok{ tbl\_event,}
  \AttributeTok{tbl\_trading\_calendar =}\NormalTok{ tbl\_trading\_calendar,}
  \AttributeTok{tbl\_return\_market =}\NormalTok{ tbl\_return\_market,}
  \AttributeTok{estimation\_start =}\NormalTok{ estimation\_start,}
  \AttributeTok{event\_end =}\NormalTok{ event\_end}
\NormalTok{)}

\NormalTok{tbl\_event\_panel }\SpecialCharTok{|\textgreater{}}
  \FunctionTok{select}\NormalTok{(}
\NormalTok{    event\_ID,}
\NormalTok{    firm\_ID,}
\NormalTok{    relative\_day,}
\NormalTok{    date,}
\NormalTok{    R,}
\NormalTok{    R\_M}
\NormalTok{  ) }\SpecialCharTok{|\textgreater{}}
  \FunctionTok{slice\_head}\NormalTok{(}\AttributeTok{n =} \DecValTok{20}\NormalTok{)}
\end{Highlighting}
\end{Shaded}

\begin{verbatim}
# A tibble: 20 x 6
   event_ID firm_ID relative_day date               R       R_M
      <int>   <int>        <int> <date>         <dbl>     <dbl>
 1        1      39         -130 2016-06-23  0.00985   0.0181  
 2        1      39         -129 2016-06-24  0.0126    0.00929 
 3        1      39         -128 2016-06-27  0.0151    0.00311 
 4        1      39         -127 2016-06-28  0.0158    0.000673
 5        1      39         -126 2016-06-29 -0.00340  -0.00482 
 6        1      39         -125 2016-06-30 -0.00710  -0.00381 
 7        1      39         -124 2016-07-01  0.0147   -0.00324 
 8        1      39         -123 2016-07-04 -0.00698  -0.00125 
 9        1      39         -122 2016-07-05 -0.0237   -0.0185  
10        1      39         -121 2016-07-06  0.0116    0.0170  
11        1      39         -120 2016-07-07 -0.00649  -0.00114 
12        1      39         -119 2016-07-08 -0.00622  -0.00122 
13        1      39         -118 2016-07-11  0.00555   0.0172  
14        1      39         -117 2016-07-12  0.00643   0.00606 
15        1      39         -116 2016-07-13 -0.0270   -0.0224  
16        1      39         -115 2016-07-14 -0.0141   -0.0209  
17        1      39         -114 2016-07-15  0.00646   0.00811 
18        1      39         -113 2016-07-19 -0.0106   -0.00262 
19        1      39         -112 2016-07-20  0.0242    0.0316  
20        1      39         -111 2016-07-21 -0.000444 -0.000248
\end{verbatim}

\subsection{6.4.5
データ利用可能性}\label{ux30c7ux30fcux30bfux5229ux7528ux53efux80fdux6027}

\begin{Shaded}
\begin{Highlighting}[]
\NormalTok{estimation\_window\_length }\OtherTok{\textless{}{-}}
\NormalTok{  estimation\_end }\SpecialCharTok{{-}}
\NormalTok{  estimation\_start }\SpecialCharTok{+} \DecValTok{1}\NormalTok{L}

\NormalTok{event\_window\_length }\OtherTok{\textless{}{-}}
\NormalTok{  event\_end }\SpecialCharTok{{-}}
\NormalTok{  event\_start }\SpecialCharTok{+} \DecValTok{1}\NormalTok{L}

\NormalTok{tbl\_event\_coverage }\OtherTok{\textless{}{-}}\NormalTok{ tbl\_event\_panel }\SpecialCharTok{|\textgreater{}}
  \FunctionTok{group\_by}\NormalTok{(}
\NormalTok{    event\_ID,}
\NormalTok{    firm\_ID,}
\NormalTok{    event\_strength}
\NormalTok{  ) }\SpecialCharTok{|\textgreater{}}
  \FunctionTok{summarise}\NormalTok{(}
    \AttributeTok{valid\_estimation\_days =} \FunctionTok{sum}\NormalTok{(}
      \FunctionTok{between}\NormalTok{(}
\NormalTok{        relative\_day,}
\NormalTok{        estimation\_start,}
\NormalTok{        estimation\_end}
\NormalTok{      ) }\SpecialCharTok{\&}
        \FunctionTok{is.finite}\NormalTok{(R) }\SpecialCharTok{\&}
        \FunctionTok{is.finite}\NormalTok{(R\_M)}
\NormalTok{    ),}
    \AttributeTok{valid\_event\_days =} \FunctionTok{sum}\NormalTok{(}
      \FunctionTok{between}\NormalTok{(}
\NormalTok{        relative\_day,}
\NormalTok{        event\_start,}
\NormalTok{        event\_end}
\NormalTok{      ) }\SpecialCharTok{\&}
        \FunctionTok{is.finite}\NormalTok{(R) }\SpecialCharTok{\&}
        \FunctionTok{is.finite}\NormalTok{(R\_M)}
\NormalTok{    ),}
    \AttributeTok{.groups =} \StringTok{"drop"}
\NormalTok{  )}

\NormalTok{tbl\_event\_coverage }\SpecialCharTok{|\textgreater{}}
  \FunctionTok{summarise}\NormalTok{(}
    \AttributeTok{events =} \FunctionTok{n}\NormalTok{(),}
    \AttributeTok{complete\_estimation\_window =} \FunctionTok{sum}\NormalTok{(}
\NormalTok{      valid\_estimation\_days }\SpecialCharTok{==}
\NormalTok{        estimation\_window\_length}
\NormalTok{    ),}
    \AttributeTok{usable\_estimation\_window =} \FunctionTok{sum}\NormalTok{(}
\NormalTok{      valid\_estimation\_days }\SpecialCharTok{\textgreater{}=}
\NormalTok{        minimum\_estimation\_observations}
\NormalTok{    ),}
    \AttributeTok{complete\_event\_window =} \FunctionTok{sum}\NormalTok{(}
\NormalTok{      valid\_event\_days }\SpecialCharTok{==}
\NormalTok{        event\_window\_length}
\NormalTok{    )}
\NormalTok{  )}
\end{Highlighting}
\end{Shaded}

\begin{verbatim}
# A tibble: 1 x 4
  events complete_estimation_window usable_estimation_wi~1 complete_event_window
   <int>                      <int>                  <int>                 <int>
1   3960                       3960                   3960                  3960
# i abbreviated name: 1: usable_estimation_window
\end{verbatim}

\subsection{6.4.6
推定期間の利用可能日数}\label{ux63a8ux5b9aux671fux9593ux306eux5229ux7528ux53efux80fdux65e5ux6570}

\begin{Shaded}
\begin{Highlighting}[]
\NormalTok{plot\_estimation\_coverage }\OtherTok{\textless{}{-}}\NormalTok{ tbl\_event\_coverage }\SpecialCharTok{|\textgreater{}}
  \FunctionTok{ggplot}\NormalTok{(}
    \FunctionTok{aes}\NormalTok{(}
      \AttributeTok{x =}\NormalTok{ valid\_estimation\_days}
\NormalTok{    )}
\NormalTok{  ) }\SpecialCharTok{+}
  \FunctionTok{geom\_histogram}\NormalTok{(}
    \AttributeTok{binwidth =} \DecValTok{5}\NormalTok{,}
    \AttributeTok{boundary =} \DecValTok{0}
\NormalTok{  ) }\SpecialCharTok{+}
  \FunctionTok{geom\_vline}\NormalTok{(}
    \AttributeTok{xintercept =}
\NormalTok{      minimum\_estimation\_observations,}
    \AttributeTok{linetype =} \StringTok{"dotted"}
\NormalTok{  ) }\SpecialCharTok{+}
  \FunctionTok{labs}\NormalTok{(}
    \AttributeTok{x =} \StringTok{"Valid estimation{-}window observations"}\NormalTok{,}
    \AttributeTok{y =} \StringTok{"Events"}\NormalTok{,}
    \AttributeTok{title =} \StringTok{"Estimation{-}Window Data Coverage"}
\NormalTok{  ) }\SpecialCharTok{+}
  \FunctionTok{theme\_classic}\NormalTok{()}

\NormalTok{plot\_estimation\_coverage}
\end{Highlighting}
\end{Shaded}

\pandocbounded{\includegraphics[keepaspectratio]{part1_empirical/ch06_event_study_files/figure-pdf/plot-estimation-coverage-1.pdf}}

\section{6.5
マーケット・モデルの推定}\label{ux30deux30fcux30b1ux30c3ux30c8ux30e2ux30c7ux30ebux306eux63a8ux5b9a}

\subsection{6.5.1
イベント別回帰}\label{ux30a4ux30d9ux30f3ux30c8ux5225ux56deux5e30}

\begin{Shaded}
\begin{Highlighting}[]
\NormalTok{tbl\_market\_model }\OtherTok{\textless{}{-}} \FunctionTok{estimate\_market\_models\_LFL}\NormalTok{(}
  \AttributeTok{tbl\_event\_panel =}\NormalTok{ tbl\_event\_panel,}
  \AttributeTok{estimation\_start =}\NormalTok{ estimation\_start,}
  \AttributeTok{estimation\_end =}\NormalTok{ estimation\_end,}
  \AttributeTok{minimum\_observations =}
\NormalTok{    minimum\_estimation\_observations}
\NormalTok{)}

\NormalTok{tbl\_market\_model }\SpecialCharTok{|\textgreater{}}
  \FunctionTok{slice\_head}\NormalTok{(}\AttributeTok{n =} \DecValTok{20}\NormalTok{)}
\end{Highlighting}
\end{Shaded}

\begin{verbatim}
# A tibble: 20 x 13
   event_ID firm_ID event_date event_trade_date event_year event_strength
      <int>   <int> <date>     <date>                <dbl> <ord>         
 1        1      39 2017-01-04 2017-01-05             2017 1             
 2        2      57 2017-01-04 2017-01-05             2017 1             
 3        3      73 2017-01-04 2017-01-05             2017 4             
 4        4     109 2017-01-04 2017-01-05             2017 3             
 5        5     111 2017-01-04 2017-01-05             2017 3             
 6        6     231 2017-01-04 2017-01-05             2017 1             
 7        7     944 2017-01-04 2017-01-05             2017 2             
 8        8    1227 2017-01-04 2017-01-05             2017 5             
 9        9     422 2017-01-05 2017-01-06             2017 3             
10       10    1203 2017-01-05 2017-01-06             2017 5             
11       11     346 2017-01-06 2017-01-10             2017 2             
12       12     590 2017-01-06 2017-01-10             2017 1             
13       13     644 2017-01-06 2017-01-10             2017 2             
14       14    1153 2017-01-06 2017-01-10             2017 4             
15       15     191 2017-01-10 2017-01-11             2017 4             
16       16     353 2017-01-10 2017-01-11             2017 1             
17       17     566 2017-01-10 2017-01-11             2017 3             
18       18     322 2017-01-11 2017-01-12             2017 3             
19       19    1050 2017-01-11 2017-01-12             2017 4             
20       20     426 2017-01-12 2017-01-13             2017 5             
# i 7 more variables: earnings_surprise <dbl>, alpha <dbl>, beta <dbl>,
#   sigma_AR <dbl>, r_squared <dbl>, adjusted_r_squared <dbl>,
#   estimation_observations <int>
\end{verbatim}

\subsection{6.5.2
推定結果の概要}\label{ux63a8ux5b9aux7d50ux679cux306eux6982ux8981}

\begin{Shaded}
\begin{Highlighting}[]
\NormalTok{tbl\_market\_model\_summary }\OtherTok{\textless{}{-}}\NormalTok{ tbl\_market\_model }\SpecialCharTok{|\textgreater{}}
  \FunctionTok{summarise}\NormalTok{(}
    \AttributeTok{estimated\_events =} \FunctionTok{n}\NormalTok{(),}
    \AttributeTok{mean\_alpha =} \FunctionTok{mean}\NormalTok{(alpha),}
    \AttributeTok{mean\_beta =} \FunctionTok{mean}\NormalTok{(beta),}
    \AttributeTok{median\_beta =} \FunctionTok{median}\NormalTok{(beta),}
    \AttributeTok{minimum\_beta =} \FunctionTok{min}\NormalTok{(beta),}
    \AttributeTok{maximum\_beta =} \FunctionTok{max}\NormalTok{(beta),}
    \AttributeTok{mean\_r\_squared =} \FunctionTok{mean}\NormalTok{(r\_squared),}
    \AttributeTok{median\_r\_squared =} \FunctionTok{median}\NormalTok{(r\_squared),}
    \AttributeTok{mean\_sigma\_AR =} \FunctionTok{mean}\NormalTok{(sigma\_AR),}
    \AttributeTok{minimum\_observations =}
      \FunctionTok{min}\NormalTok{(estimation\_observations),}
    \AttributeTok{maximum\_observations =}
      \FunctionTok{max}\NormalTok{(estimation\_observations)}
\NormalTok{  )}

\NormalTok{tbl\_market\_model\_summary}
\end{Highlighting}
\end{Shaded}

\begin{verbatim}
# A tibble: 1 x 11
  estimated_events mean_alpha mean_beta median_beta minimum_beta maximum_beta
             <int>      <dbl>     <dbl>       <dbl>        <dbl>        <dbl>
1             3960 -0.0000173     0.998        1.00     -0.00834         2.18
# i 5 more variables: mean_r_squared <dbl>, median_r_squared <dbl>,
#   mean_sigma_AR <dbl>, minimum_observations <int>, maximum_observations <int>
\end{verbatim}

\subsection{6.5.3
ベータの分布}\label{ux30d9ux30fcux30bfux306eux5206ux5e03}

\begin{Shaded}
\begin{Highlighting}[]
\NormalTok{plot\_market\_beta }\OtherTok{\textless{}{-}}\NormalTok{ tbl\_market\_model }\SpecialCharTok{|\textgreater{}}
  \FunctionTok{ggplot}\NormalTok{(}\FunctionTok{aes}\NormalTok{(}\AttributeTok{x =}\NormalTok{ beta)) }\SpecialCharTok{+}
  \FunctionTok{geom\_histogram}\NormalTok{(}\AttributeTok{bins =} \DecValTok{40}\NormalTok{) }\SpecialCharTok{+}
  \FunctionTok{geom\_vline}\NormalTok{(}
    \AttributeTok{xintercept =} \DecValTok{1}\NormalTok{,}
    \AttributeTok{linetype =} \StringTok{"dotted"}
\NormalTok{  ) }\SpecialCharTok{+}
  \FunctionTok{labs}\NormalTok{(}
    \AttributeTok{x =} \StringTok{"Market beta"}\NormalTok{,}
    \AttributeTok{y =} \StringTok{"Events"}\NormalTok{,}
    \AttributeTok{title =} \StringTok{"Distribution of Event{-}Specific Market Betas"}
\NormalTok{  ) }\SpecialCharTok{+}
  \FunctionTok{theme\_classic}\NormalTok{()}

\NormalTok{plot\_market\_beta}
\end{Highlighting}
\end{Shaded}

\pandocbounded{\includegraphics[keepaspectratio]{part1_empirical/ch06_event_study_files/figure-pdf/plot-market-beta-1.pdf}}

\subsection{6.5.4
アルファの分布}\label{ux30a2ux30ebux30d5ux30a1ux306eux5206ux5e03}

\begin{Shaded}
\begin{Highlighting}[]
\NormalTok{plot\_market\_alpha }\OtherTok{\textless{}{-}}\NormalTok{ tbl\_market\_model }\SpecialCharTok{|\textgreater{}}
  \FunctionTok{ggplot}\NormalTok{(}\FunctionTok{aes}\NormalTok{(}\AttributeTok{x =}\NormalTok{ alpha)) }\SpecialCharTok{+}
  \FunctionTok{geom\_histogram}\NormalTok{(}\AttributeTok{bins =} \DecValTok{40}\NormalTok{) }\SpecialCharTok{+}
  \FunctionTok{geom\_vline}\NormalTok{(}
    \AttributeTok{xintercept =} \DecValTok{0}\NormalTok{,}
    \AttributeTok{linetype =} \StringTok{"dotted"}
\NormalTok{  ) }\SpecialCharTok{+}
  \FunctionTok{scale\_x\_continuous}\NormalTok{(}
    \AttributeTok{labels =}\NormalTok{ scales}\SpecialCharTok{::}\FunctionTok{label\_percent}\NormalTok{()}
\NormalTok{  ) }\SpecialCharTok{+}
  \FunctionTok{labs}\NormalTok{(}
    \AttributeTok{x =} \StringTok{"Market{-}model alpha"}\NormalTok{,}
    \AttributeTok{y =} \StringTok{"Events"}\NormalTok{,}
    \AttributeTok{title =} \StringTok{"Distribution of Event{-}Specific Alphas"}
\NormalTok{  ) }\SpecialCharTok{+}
  \FunctionTok{theme\_classic}\NormalTok{()}

\NormalTok{plot\_market\_alpha}
\end{Highlighting}
\end{Shaded}

\pandocbounded{\includegraphics[keepaspectratio]{part1_empirical/ch06_event_study_files/figure-pdf/plot-market-alpha-1.pdf}}

\subsection{6.5.5
決定係数の分布}\label{ux6c7aux5b9aux4fc2ux6570ux306eux5206ux5e03}

\begin{Shaded}
\begin{Highlighting}[]
\NormalTok{plot\_market\_model\_r\_squared }\OtherTok{\textless{}{-}}\NormalTok{ tbl\_market\_model }\SpecialCharTok{|\textgreater{}}
  \FunctionTok{ggplot}\NormalTok{(}\FunctionTok{aes}\NormalTok{(}\AttributeTok{x =}\NormalTok{ r\_squared)) }\SpecialCharTok{+}
  \FunctionTok{geom\_histogram}\NormalTok{(}\AttributeTok{bins =} \DecValTok{40}\NormalTok{) }\SpecialCharTok{+}
  \FunctionTok{labs}\NormalTok{(}
    \AttributeTok{x =} \StringTok{"R{-}squared"}\NormalTok{,}
    \AttributeTok{y =} \StringTok{"Events"}\NormalTok{,}
    \AttributeTok{title =} \StringTok{"Goodness of Fit of Event{-}Specific Market Models"}
\NormalTok{  ) }\SpecialCharTok{+}
  \FunctionTok{theme\_classic}\NormalTok{()}

\NormalTok{plot\_market\_model\_r\_squared}
\end{Highlighting}
\end{Shaded}

\pandocbounded{\includegraphics[keepaspectratio]{part1_empirical/ch06_event_study_files/figure-pdf/plot-market-model-r-squared-1.pdf}}

\subsection{6.5.6
ベータと決定係数}\label{ux30d9ux30fcux30bfux3068ux6c7aux5b9aux4fc2ux6570}

\begin{Shaded}
\begin{Highlighting}[]
\NormalTok{plot\_beta\_r\_squared }\OtherTok{\textless{}{-}}\NormalTok{ tbl\_market\_model }\SpecialCharTok{|\textgreater{}}
  \FunctionTok{ggplot}\NormalTok{(}
    \FunctionTok{aes}\NormalTok{(}
      \AttributeTok{x =}\NormalTok{ beta,}
      \AttributeTok{y =}\NormalTok{ r\_squared,}
      \AttributeTok{shape =}\NormalTok{ event\_strength}
\NormalTok{    )}
\NormalTok{  ) }\SpecialCharTok{+}
  \FunctionTok{geom\_point}\NormalTok{(}\AttributeTok{alpha =} \FloatTok{0.55}\NormalTok{) }\SpecialCharTok{+}
  \FunctionTok{geom\_vline}\NormalTok{(}
    \AttributeTok{xintercept =} \DecValTok{1}\NormalTok{,}
    \AttributeTok{linetype =} \StringTok{"dotted"}
\NormalTok{  ) }\SpecialCharTok{+}
  \FunctionTok{labs}\NormalTok{(}
    \AttributeTok{x =} \StringTok{"Market beta"}\NormalTok{,}
    \AttributeTok{y =} \StringTok{"R{-}squared"}\NormalTok{,}
    \AttributeTok{shape =} \StringTok{"Event strength"}\NormalTok{,}
    \AttributeTok{title =} \StringTok{"Market Beta and Model Fit"}
\NormalTok{  ) }\SpecialCharTok{+}
  \FunctionTok{theme\_classic}\NormalTok{()}

\NormalTok{plot\_beta\_r\_squared}
\end{Highlighting}
\end{Shaded}

\pandocbounded{\includegraphics[keepaspectratio]{part1_empirical/ch06_event_study_files/figure-pdf/plot-beta-r-squared-1.pdf}}

\section{6.6
異常リターンと累積異常リターン}\label{ux7570ux5e38ux30eaux30bfux30fcux30f3ux3068ux7d2fux7a4dux7570ux5e38ux30eaux30bfux30fcux30f3}

\subsection{6.6.1 ARとCAR}\label{arux3068car}

\begin{Shaded}
\begin{Highlighting}[]
\NormalTok{tbl\_abnormal\_return }\OtherTok{\textless{}{-}} \FunctionTok{compute\_abnormal\_returns\_LFL}\NormalTok{(}
  \AttributeTok{tbl\_event\_panel =}\NormalTok{ tbl\_event\_panel,}
  \AttributeTok{tbl\_market\_model =}\NormalTok{ tbl\_market\_model,}
  \AttributeTok{event\_start =}\NormalTok{ event\_start,}
  \AttributeTok{event\_end =}\NormalTok{ event\_end,}
  \AttributeTok{require\_complete\_window =} \ConstantTok{TRUE}
\NormalTok{)}

\NormalTok{tbl\_abnormal\_return }\SpecialCharTok{|\textgreater{}}
  \FunctionTok{select}\NormalTok{(}
\NormalTok{    event\_ID,}
\NormalTok{    firm\_ID,}
\NormalTok{    event\_strength,}
\NormalTok{    relative\_day,}
\NormalTok{    date,}
\NormalTok{    R,}
\NormalTok{    R\_M,}
\NormalTok{    normal\_return,}
\NormalTok{    AR,}
\NormalTok{    CAR}
\NormalTok{  ) }\SpecialCharTok{|\textgreater{}}
  \FunctionTok{slice\_head}\NormalTok{(}\AttributeTok{n =} \DecValTok{20}\NormalTok{)}
\end{Highlighting}
\end{Shaded}

\begin{verbatim}
# A tibble: 20 x 10
   event_ID firm_ID event_strength relative_day date              R       R_M
      <int>   <int> <ord>                 <int> <date>        <dbl>     <dbl>
 1        1      39 1                       -30 2016-11-18 -0.00132 -0.00147 
 2        1      39 1                       -29 2016-11-21  0.00372 -0.00334 
 3        1      39 1                       -28 2016-11-22  0.0210   0.0171  
 4        1      39 1                       -27 2016-11-24  0.00309  0.0157  
 5        1      39 1                       -26 2016-11-25 -0.00626 -0.0233  
 6        1      39 1                       -25 2016-11-28 -0.00406 -0.00654 
 7        1      39 1                       -24 2016-11-29  0.0123  -0.00414 
 8        1      39 1                       -23 2016-11-30 -0.00965  0.00293 
 9        1      39 1                       -22 2016-12-01  0.0140   0.000473
10        1      39 1                       -21 2016-12-02  0.0122   0.0108  
11        1      39 1                       -20 2016-12-05  0.00978 -0.00548 
12        1      39 1                       -19 2016-12-06  0.0261   0.0186  
13        1      39 1                       -18 2016-12-07 -0.00441 -0.00733 
14        1      39 1                       -17 2016-12-08  0.0150   0.0276  
15        1      39 1                       -16 2016-12-09  0.00745 -0.00144 
16        1      39 1                       -15 2016-12-12  0.0114  -0.00642 
17        1      39 1                       -14 2016-12-13 -0.0110  -0.0143  
18        1      39 1                       -13 2016-12-14  0.0195   0.0224  
19        1      39 1                       -12 2016-12-15  0.00550  0.0186  
20        1      39 1                       -11 2016-12-16 -0.00667 -0.00344 
# i 3 more variables: normal_return <dbl>, AR <dbl>, CAR <dbl>
\end{verbatim}

\subsection{6.6.2 AARとCAAR}\label{aarux3068caar}

イベント強度 \(g\) に属するイベント数を \(N_g\)
とすると、平均異常リターンは次式で表される。

\[
AAR_{g,t}
=
\frac{1}{N_g}
\sum_{i \in g}
AR_{i,t}
\]

平均累積異常リターンは次式で表される。

\[
CAAR_{g,t}
=
\frac{1}{N_g}
\sum_{i \in g}
CAR_{i,t}
\]

\begin{Shaded}
\begin{Highlighting}[]
\NormalTok{significance\_stars }\OtherTok{\textless{}{-}} \ControlFlowTok{function}\NormalTok{(p\_value) \{}
  \FunctionTok{case\_when}\NormalTok{(}
    \FunctionTok{is.na}\NormalTok{(p\_value) }\SpecialCharTok{\textasciitilde{}} \StringTok{""}\NormalTok{,}
\NormalTok{    p\_value }\SpecialCharTok{\textless{}} \FloatTok{0.01} \SpecialCharTok{\textasciitilde{}} \StringTok{"***"}\NormalTok{,}
\NormalTok{    p\_value }\SpecialCharTok{\textless{}} \FloatTok{0.05} \SpecialCharTok{\textasciitilde{}} \StringTok{"**"}\NormalTok{,}
\NormalTok{    p\_value }\SpecialCharTok{\textless{}} \FloatTok{0.10} \SpecialCharTok{\textasciitilde{}} \StringTok{"*"}\NormalTok{,}
    \ConstantTok{TRUE} \SpecialCharTok{\textasciitilde{}} \StringTok{""}
\NormalTok{  )}
\NormalTok{\}}
\end{Highlighting}
\end{Shaded}

\begin{Shaded}
\begin{Highlighting}[]
\NormalTok{tbl\_daily\_statistics }\OtherTok{\textless{}{-}}\NormalTok{ tbl\_abnormal\_return }\SpecialCharTok{|\textgreater{}}
  \FunctionTok{group\_by}\NormalTok{(}
\NormalTok{    relative\_day,}
\NormalTok{    event\_strength}
\NormalTok{  ) }\SpecialCharTok{|\textgreater{}}
  \FunctionTok{summarise}\NormalTok{(}
    \AttributeTok{events =} \FunctionTok{n\_distinct}\NormalTok{(event\_ID),}
    \AttributeTok{AAR =} \FunctionTok{mean}\NormalTok{(AR),}
    \AttributeTok{CAAR =} \FunctionTok{mean}\NormalTok{(CAR),}
    \AttributeTok{standard\_deviation\_AR =} \FunctionTok{sd}\NormalTok{(AR),}
    \AttributeTok{standard\_deviation\_CAR =} \FunctionTok{sd}\NormalTok{(CAR),}
    \AttributeTok{.groups =} \StringTok{"drop"}
\NormalTok{  ) }\SpecialCharTok{|\textgreater{}}
  \FunctionTok{mutate}\NormalTok{(}
    \AttributeTok{standard\_error\_AAR =}
\NormalTok{      standard\_deviation\_AR }\SpecialCharTok{/}
      \FunctionTok{sqrt}\NormalTok{(events),}
    \AttributeTok{standard\_error\_CAAR =}
\NormalTok{      standard\_deviation\_CAR }\SpecialCharTok{/}
      \FunctionTok{sqrt}\NormalTok{(events),}
    \AttributeTok{t\_AAR =} \FunctionTok{if\_else}\NormalTok{(}
      \FunctionTok{is.finite}\NormalTok{(standard\_error\_AAR) }\SpecialCharTok{\&}
\NormalTok{        standard\_error\_AAR }\SpecialCharTok{\textgreater{}} \DecValTok{0}\NormalTok{,}
\NormalTok{      AAR }\SpecialCharTok{/}\NormalTok{ standard\_error\_AAR,}
      \ConstantTok{NA\_real\_}
\NormalTok{    ),}
    \AttributeTok{t\_CAAR =} \FunctionTok{if\_else}\NormalTok{(}
      \FunctionTok{is.finite}\NormalTok{(standard\_error\_CAAR) }\SpecialCharTok{\&}
\NormalTok{        standard\_error\_CAAR }\SpecialCharTok{\textgreater{}} \DecValTok{0}\NormalTok{,}
\NormalTok{      CAAR }\SpecialCharTok{/}\NormalTok{ standard\_error\_CAAR,}
      \ConstantTok{NA\_real\_}
\NormalTok{    ),}
    \AttributeTok{p\_AAR =} \FunctionTok{if\_else}\NormalTok{(}
      \FunctionTok{is.finite}\NormalTok{(t\_AAR),}
      \DecValTok{2} \SpecialCharTok{*} \FunctionTok{pt}\NormalTok{(}
        \SpecialCharTok{{-}}\FunctionTok{abs}\NormalTok{(t\_AAR),}
        \AttributeTok{df =}\NormalTok{ events }\SpecialCharTok{{-}} \DecValTok{1}\NormalTok{L}
\NormalTok{      ),}
      \ConstantTok{NA\_real\_}
\NormalTok{    ),}
    \AttributeTok{p\_CAAR =} \FunctionTok{if\_else}\NormalTok{(}
      \FunctionTok{is.finite}\NormalTok{(t\_CAAR),}
      \DecValTok{2} \SpecialCharTok{*} \FunctionTok{pt}\NormalTok{(}
        \SpecialCharTok{{-}}\FunctionTok{abs}\NormalTok{(t\_CAAR),}
        \AttributeTok{df =}\NormalTok{ events }\SpecialCharTok{{-}} \DecValTok{1}\NormalTok{L}
\NormalTok{      ),}
      \ConstantTok{NA\_real\_}
\NormalTok{    ),}
    \AttributeTok{significance\_AAR =}
      \FunctionTok{significance\_stars}\NormalTok{(p\_AAR),}
    \AttributeTok{significance\_CAAR =}
      \FunctionTok{significance\_stars}\NormalTok{(p\_CAAR)}
\NormalTok{  ) }\SpecialCharTok{|\textgreater{}}
  \FunctionTok{arrange}\NormalTok{(}
\NormalTok{    event\_strength,}
\NormalTok{    relative\_day}
\NormalTok{  )}

\NormalTok{tbl\_daily\_statistics}
\end{Highlighting}
\end{Shaded}

\begin{verbatim}
# A tibble: 305 x 15
   relative_day event_strength events        AAR      CAAR standard_deviation_AR
          <int> <ord>           <int>      <dbl>     <dbl>                 <dbl>
 1          -30 1                 793 -0.000333   -3.33e-4               0.00800
 2          -29 1                 793  0.000346    1.32e-5               0.00812
 3          -28 1                 793 -0.0000305  -1.73e-5               0.00768
 4          -27 1                 793  0.000242    2.25e-4               0.00795
 5          -26 1                 793  0.000420    6.45e-4               0.00825
 6          -25 1                 793  0.000208    8.53e-4               0.00808
 7          -24 1                 793  0.000534    1.39e-3               0.00806
 8          -23 1                 793 -0.0000285   1.36e-3               0.00816
 9          -22 1                 793 -0.000142    1.22e-3               0.00797
10          -21 1                 793 -0.000192    1.02e-3               0.00831
# i 295 more rows
# i 9 more variables: standard_deviation_CAR <dbl>, standard_error_AAR <dbl>,
#   standard_error_CAAR <dbl>, t_AAR <dbl>, t_CAAR <dbl>, p_AAR <dbl>,
#   p_CAAR <dbl>, significance_AAR <chr>, significance_CAAR <chr>
\end{verbatim}

\subsection{6.6.3
イベント強度別CAAR}\label{ux30a4ux30d9ux30f3ux30c8ux5f37ux5ea6ux5225caar}

\begin{Shaded}
\begin{Highlighting}[]
\NormalTok{plot\_caar\_by\_strength }\OtherTok{\textless{}{-}}\NormalTok{ tbl\_daily\_statistics }\SpecialCharTok{|\textgreater{}}
  \FunctionTok{ggplot}\NormalTok{(}
    \FunctionTok{aes}\NormalTok{(}
      \AttributeTok{x =}\NormalTok{ relative\_day,}
      \AttributeTok{y =}\NormalTok{ CAAR,}
      \AttributeTok{linetype =}\NormalTok{ event\_strength}
\NormalTok{    )}
\NormalTok{  ) }\SpecialCharTok{+}
  \FunctionTok{annotate}\NormalTok{(}
    \StringTok{"rect"}\NormalTok{,}
    \AttributeTok{xmin =} \SpecialCharTok{{-}}\DecValTok{1}\NormalTok{,}
    \AttributeTok{xmax =} \DecValTok{1}\NormalTok{,}
    \AttributeTok{ymin =} \SpecialCharTok{{-}}\ConstantTok{Inf}\NormalTok{,}
    \AttributeTok{ymax =} \ConstantTok{Inf}\NormalTok{,}
    \AttributeTok{alpha =} \FloatTok{0.08}
\NormalTok{  ) }\SpecialCharTok{+}
  \FunctionTok{geom\_line}\NormalTok{(}\AttributeTok{linewidth =} \FloatTok{0.9}\NormalTok{) }\SpecialCharTok{+}
  \FunctionTok{geom\_hline}\NormalTok{(}
    \AttributeTok{yintercept =} \DecValTok{0}\NormalTok{,}
    \AttributeTok{linetype =} \StringTok{"dotted"}
\NormalTok{  ) }\SpecialCharTok{+}
  \FunctionTok{geom\_vline}\NormalTok{(}
    \AttributeTok{xintercept =} \DecValTok{0}\NormalTok{,}
    \AttributeTok{linetype =} \StringTok{"dotted"}
\NormalTok{  ) }\SpecialCharTok{+}
  \FunctionTok{scale\_y\_continuous}\NormalTok{(}
    \AttributeTok{labels =}\NormalTok{ scales}\SpecialCharTok{::}\FunctionTok{label\_percent}\NormalTok{()}
\NormalTok{  ) }\SpecialCharTok{+}
  \FunctionTok{labs}\NormalTok{(}
    \AttributeTok{x =} \StringTok{"Relative trading day"}\NormalTok{,}
    \AttributeTok{y =} \StringTok{"CAAR"}\NormalTok{,}
    \AttributeTok{linetype =} \StringTok{"Event strength"}\NormalTok{,}
    \AttributeTok{title =} \StringTok{"CAAR by Earnings{-}Surprise Group"}
\NormalTok{  ) }\SpecialCharTok{+}
  \FunctionTok{theme\_classic}\NormalTok{()}

\NormalTok{plot\_caar\_by\_strength}
\end{Highlighting}
\end{Shaded}

\pandocbounded{\includegraphics[keepaspectratio]{part1_empirical/ch06_event_study_files/figure-pdf/plot-caar-by-strength-1.pdf}}

\subsection{6.6.4
イベント日前後のAAR}\label{ux30a4ux30d9ux30f3ux30c8ux65e5ux524dux5f8cux306eaar}

\begin{Shaded}
\begin{Highlighting}[]
\NormalTok{plot\_aar\_around\_event }\OtherTok{\textless{}{-}}\NormalTok{ tbl\_daily\_statistics }\SpecialCharTok{|\textgreater{}}
  \FunctionTok{filter}\NormalTok{(}
    \FunctionTok{between}\NormalTok{(}
\NormalTok{      relative\_day,}
      \SpecialCharTok{{-}}\DecValTok{10}\NormalTok{L,}
      \DecValTok{10}\NormalTok{L}
\NormalTok{    )}
\NormalTok{  ) }\SpecialCharTok{|\textgreater{}}
  \FunctionTok{ggplot}\NormalTok{(}
    \FunctionTok{aes}\NormalTok{(}
      \AttributeTok{x =}\NormalTok{ relative\_day,}
      \AttributeTok{y =}\NormalTok{ AAR,}
      \AttributeTok{linetype =}\NormalTok{ event\_strength}
\NormalTok{    )}
\NormalTok{  ) }\SpecialCharTok{+}
  \FunctionTok{geom\_line}\NormalTok{(}\AttributeTok{linewidth =} \FloatTok{0.8}\NormalTok{) }\SpecialCharTok{+}
  \FunctionTok{geom\_point}\NormalTok{() }\SpecialCharTok{+}
  \FunctionTok{geom\_vline}\NormalTok{(}
    \AttributeTok{xintercept =} \DecValTok{0}\NormalTok{,}
    \AttributeTok{linetype =} \StringTok{"dotted"}
\NormalTok{  ) }\SpecialCharTok{+}
  \FunctionTok{geom\_hline}\NormalTok{(}
    \AttributeTok{yintercept =} \DecValTok{0}\NormalTok{,}
    \AttributeTok{linetype =} \StringTok{"dotted"}
\NormalTok{  ) }\SpecialCharTok{+}
  \FunctionTok{scale\_x\_continuous}\NormalTok{(}
    \AttributeTok{breaks =} \FunctionTok{seq}\NormalTok{(}\SpecialCharTok{{-}}\DecValTok{10}\NormalTok{, }\DecValTok{10}\NormalTok{, }\DecValTok{2}\NormalTok{)}
\NormalTok{  ) }\SpecialCharTok{+}
  \FunctionTok{scale\_y\_continuous}\NormalTok{(}
    \AttributeTok{labels =}\NormalTok{ scales}\SpecialCharTok{::}\FunctionTok{label\_percent}\NormalTok{()}
\NormalTok{  ) }\SpecialCharTok{+}
  \FunctionTok{labs}\NormalTok{(}
    \AttributeTok{x =} \StringTok{"Relative trading day"}\NormalTok{,}
    \AttributeTok{y =} \StringTok{"AAR"}\NormalTok{,}
    \AttributeTok{linetype =} \StringTok{"Event strength"}\NormalTok{,}
    \AttributeTok{title =} \StringTok{"Average Abnormal Returns around the Event"}
\NormalTok{  ) }\SpecialCharTok{+}
  \FunctionTok{theme\_classic}\NormalTok{()}

\NormalTok{plot\_aar\_around\_event}
\end{Highlighting}
\end{Shaded}

\pandocbounded{\includegraphics[keepaspectratio]{part1_empirical/ch06_event_study_files/figure-pdf/plot-aar-around-event-1.pdf}}

\subsection{6.6.5
イベント日0のAAR}\label{ux30a4ux30d9ux30f3ux30c8ux65e50ux306eaar}

\begin{Shaded}
\begin{Highlighting}[]
\NormalTok{plot\_event\_day\_aar }\OtherTok{\textless{}{-}}\NormalTok{ tbl\_daily\_statistics }\SpecialCharTok{|\textgreater{}}
  \FunctionTok{filter}\NormalTok{(relative\_day }\SpecialCharTok{==} \DecValTok{0}\NormalTok{L) }\SpecialCharTok{|\textgreater{}}
  \FunctionTok{ggplot}\NormalTok{(}
    \FunctionTok{aes}\NormalTok{(}
      \AttributeTok{x =}\NormalTok{ event\_strength,}
      \AttributeTok{y =}\NormalTok{ AAR}
\NormalTok{    )}
\NormalTok{  ) }\SpecialCharTok{+}
  \FunctionTok{geom\_col}\NormalTok{() }\SpecialCharTok{+}
  \FunctionTok{geom\_hline}\NormalTok{(}
    \AttributeTok{yintercept =} \DecValTok{0}\NormalTok{,}
    \AttributeTok{linetype =} \StringTok{"dotted"}
\NormalTok{  ) }\SpecialCharTok{+}
  \FunctionTok{scale\_y\_continuous}\NormalTok{(}
    \AttributeTok{labels =}\NormalTok{ scales}\SpecialCharTok{::}\FunctionTok{label\_percent}\NormalTok{()}
\NormalTok{  ) }\SpecialCharTok{+}
  \FunctionTok{labs}\NormalTok{(}
    \AttributeTok{x =} \StringTok{"Event strength"}\NormalTok{,}
    \AttributeTok{y =} \StringTok{"Event{-}day AAR"}\NormalTok{,}
    \AttributeTok{title =} \StringTok{"Abnormal Return on Event Day 0"}
\NormalTok{  ) }\SpecialCharTok{+}
  \FunctionTok{theme\_classic}\NormalTok{()}

\NormalTok{plot\_event\_day\_aar}
\end{Highlighting}
\end{Shaded}

\pandocbounded{\includegraphics[keepaspectratio]{part1_empirical/ch06_event_study_files/figure-pdf/plot-event-day-aar-1.pdf}}

\subsection{6.6.6
イベント日後のCAAR}\label{ux30a4ux30d9ux30f3ux30c8ux65e5ux5f8cux306ecaar}

\begin{Shaded}
\begin{Highlighting}[]
\NormalTok{plot\_post\_event\_caar }\OtherTok{\textless{}{-}}\NormalTok{ tbl\_daily\_statistics }\SpecialCharTok{|\textgreater{}}
  \FunctionTok{filter}\NormalTok{(}
    \FunctionTok{between}\NormalTok{(}
\NormalTok{      relative\_day,}
      \DecValTok{0}\NormalTok{L,}
      \DecValTok{30}\NormalTok{L}
\NormalTok{    )}
\NormalTok{  ) }\SpecialCharTok{|\textgreater{}}
  \FunctionTok{ggplot}\NormalTok{(}
    \FunctionTok{aes}\NormalTok{(}
      \AttributeTok{x =}\NormalTok{ relative\_day,}
      \AttributeTok{y =}\NormalTok{ CAAR,}
      \AttributeTok{linetype =}\NormalTok{ event\_strength}
\NormalTok{    )}
\NormalTok{  ) }\SpecialCharTok{+}
  \FunctionTok{geom\_line}\NormalTok{(}\AttributeTok{linewidth =} \FloatTok{0.9}\NormalTok{) }\SpecialCharTok{+}
  \FunctionTok{geom\_hline}\NormalTok{(}
    \AttributeTok{yintercept =} \DecValTok{0}\NormalTok{,}
    \AttributeTok{linetype =} \StringTok{"dotted"}
\NormalTok{  ) }\SpecialCharTok{+}
  \FunctionTok{scale\_y\_continuous}\NormalTok{(}
    \AttributeTok{labels =}\NormalTok{ scales}\SpecialCharTok{::}\FunctionTok{label\_percent}\NormalTok{()}
\NormalTok{  ) }\SpecialCharTok{+}
  \FunctionTok{labs}\NormalTok{(}
    \AttributeTok{x =} \StringTok{"Trading days after event"}\NormalTok{,}
    \AttributeTok{y =} \StringTok{"CAAR"}\NormalTok{,}
    \AttributeTok{linetype =} \StringTok{"Event strength"}\NormalTok{,}
    \AttributeTok{title =} \StringTok{"Post{-}Event Cumulative Abnormal Returns"}
\NormalTok{  ) }\SpecialCharTok{+}
  \FunctionTok{theme\_classic}\NormalTok{()}

\NormalTok{plot\_post\_event\_caar}
\end{Highlighting}
\end{Shaded}

\pandocbounded{\includegraphics[keepaspectratio]{part1_empirical/ch06_event_study_files/figure-pdf/plot-post-event-caar-1.pdf}}

\subsection{6.6.7
イベント強度別AR分布}\label{ux30a4ux30d9ux30f3ux30c8ux5f37ux5ea6ux5225arux5206ux5e03}

\begin{Shaded}
\begin{Highlighting}[]
\NormalTok{plot\_ar\_distribution\_by\_strength }\OtherTok{\textless{}{-}}\NormalTok{ tbl\_abnormal\_return }\SpecialCharTok{|\textgreater{}}
  \FunctionTok{filter}\NormalTok{(}
    \FunctionTok{between}\NormalTok{(}
\NormalTok{      relative\_day,}
      \SpecialCharTok{{-}}\DecValTok{1}\NormalTok{L,}
      \DecValTok{1}\NormalTok{L}
\NormalTok{    )}
\NormalTok{  ) }\SpecialCharTok{|\textgreater{}}
  \FunctionTok{ggplot}\NormalTok{(}
    \FunctionTok{aes}\NormalTok{(}
      \AttributeTok{x =}\NormalTok{ AR,}
      \AttributeTok{linetype =}\NormalTok{ event\_strength}
\NormalTok{    )}
\NormalTok{  ) }\SpecialCharTok{+}
  \FunctionTok{geom\_density}\NormalTok{(}\AttributeTok{linewidth =} \FloatTok{0.8}\NormalTok{) }\SpecialCharTok{+}
  \FunctionTok{geom\_vline}\NormalTok{(}
    \AttributeTok{xintercept =} \DecValTok{0}\NormalTok{,}
    \AttributeTok{linetype =} \StringTok{"dotted"}
\NormalTok{  ) }\SpecialCharTok{+}
  \FunctionTok{scale\_x\_continuous}\NormalTok{(}
    \AttributeTok{labels =}\NormalTok{ scales}\SpecialCharTok{::}\FunctionTok{label\_percent}\NormalTok{()}
\NormalTok{  ) }\SpecialCharTok{+}
  \FunctionTok{labs}\NormalTok{(}
    \AttributeTok{x =} \StringTok{"Abnormal return"}\NormalTok{,}
    \AttributeTok{y =} \StringTok{"Density"}\NormalTok{,}
    \AttributeTok{linetype =} \StringTok{"Event strength"}\NormalTok{,}
    \AttributeTok{title =} \StringTok{"Distribution of Abnormal Returns in [{-}1, +1]"}
\NormalTok{  ) }\SpecialCharTok{+}
  \FunctionTok{theme\_classic}\NormalTok{()}

\NormalTok{plot\_ar\_distribution\_by\_strength}
\end{Highlighting}
\end{Shaded}

\pandocbounded{\includegraphics[keepaspectratio]{part1_empirical/ch06_event_study_files/figure-pdf/plot-ar-distribution-by-strength-1.pdf}}

\section{6.7
イベント・ウィンドウ分析}\label{ux30a4ux30d9ux30f3ux30c8ux30a6ux30a3ux30f3ux30c9ux30a6ux5206ux6790}

\subsection{6.7.1
ウィンドウ定義}\label{ux30a6ux30a3ux30f3ux30c9ux30a6ux5b9aux7fa9}

\begin{Shaded}
\begin{Highlighting}[]
\NormalTok{tbl\_summary\_window }\OtherTok{\textless{}{-}} \FunctionTok{tribble}\NormalTok{(}
  \SpecialCharTok{\textasciitilde{}}\NormalTok{window\_ID, }\SpecialCharTok{\textasciitilde{}}\NormalTok{window\_label, }\SpecialCharTok{\textasciitilde{}}\NormalTok{start\_day, }\SpecialCharTok{\textasciitilde{}}\NormalTok{end\_day,}
  \StringTok{"m1\_p1"}\NormalTok{, }\StringTok{"[{-}1, +1]"}\NormalTok{, }\SpecialCharTok{{-}}\DecValTok{1}\NormalTok{L, }\DecValTok{1}\NormalTok{L,}
  \StringTok{"0\_p1"}\NormalTok{, }\StringTok{"[0, +1]"}\NormalTok{, }\DecValTok{0}\NormalTok{L, }\DecValTok{1}\NormalTok{L,}
  \StringTok{"m5\_p5"}\NormalTok{, }\StringTok{"[{-}5, +5]"}\NormalTok{, }\SpecialCharTok{{-}}\DecValTok{5}\NormalTok{L, }\DecValTok{5}\NormalTok{L,}
  \StringTok{"m10\_p10"}\NormalTok{, }\StringTok{"[{-}10, +10]"}\NormalTok{, }\SpecialCharTok{{-}}\DecValTok{10}\NormalTok{L, }\DecValTok{10}\NormalTok{L,}
  \StringTok{"m30\_p30"}\NormalTok{, }\StringTok{"[{-}30, +30]"}\NormalTok{, }\SpecialCharTok{{-}}\DecValTok{30}\NormalTok{L, }\DecValTok{30}\NormalTok{L}
\NormalTok{)}

\NormalTok{tbl\_summary\_window}
\end{Highlighting}
\end{Shaded}

\begin{verbatim}
# A tibble: 5 x 4
  window_ID window_label start_day end_day
  <chr>     <chr>            <int>   <int>
1 m1_p1     [-1, +1]            -1       1
2 0_p1      [0, +1]              0       1
3 m5_p5     [-5, +5]            -5       5
4 m10_p10   [-10, +10]         -10      10
5 m30_p30   [-30, +30]         -30      30
\end{verbatim}

\subsection{6.7.2
イベント別CAR}\label{ux30a4ux30d9ux30f3ux30c8ux5225car}

\begin{Shaded}
\begin{Highlighting}[]
\NormalTok{tbl\_event\_window\_return }\OtherTok{\textless{}{-}}
  \FunctionTok{compute\_event\_window\_returns\_LFL}\NormalTok{(}
    \AttributeTok{tbl\_abnormal\_return =}\NormalTok{ tbl\_abnormal\_return,}
    \AttributeTok{tbl\_window =}\NormalTok{ tbl\_summary\_window}
\NormalTok{  )}

\NormalTok{tbl\_event\_window\_return }\SpecialCharTok{|\textgreater{}}
  \FunctionTok{slice\_head}\NormalTok{(}\AttributeTok{n =} \DecValTok{20}\NormalTok{)}
\end{Highlighting}
\end{Shaded}

\begin{verbatim}
# A tibble: 20 x 14
   event_ID firm_ID event_date event_trade_date event_year event_strength
      <int>   <int> <date>     <date>                <dbl> <ord>         
 1        1      39 2017-01-04 2017-01-05             2017 1             
 2        2      57 2017-01-04 2017-01-05             2017 1             
 3        3      73 2017-01-04 2017-01-05             2017 4             
 4        4     109 2017-01-04 2017-01-05             2017 3             
 5        5     111 2017-01-04 2017-01-05             2017 3             
 6        6     231 2017-01-04 2017-01-05             2017 1             
 7        7     944 2017-01-04 2017-01-05             2017 2             
 8        8    1227 2017-01-04 2017-01-05             2017 5             
 9        9     422 2017-01-05 2017-01-06             2017 3             
10       10    1203 2017-01-05 2017-01-06             2017 5             
11       11     346 2017-01-06 2017-01-10             2017 2             
12       12     590 2017-01-06 2017-01-10             2017 1             
13       13     644 2017-01-06 2017-01-10             2017 2             
14       14    1153 2017-01-06 2017-01-10             2017 4             
15       15     191 2017-01-10 2017-01-11             2017 4             
16       16     353 2017-01-10 2017-01-11             2017 1             
17       17     566 2017-01-10 2017-01-11             2017 3             
18       18     322 2017-01-11 2017-01-12             2017 3             
19       19    1050 2017-01-11 2017-01-12             2017 4             
20       20     426 2017-01-12 2017-01-13             2017 5             
# i 8 more variables: earnings_surprise <dbl>, window_ID <chr>,
#   window_label <chr>, start_day <int>, end_day <int>, expected_days <int>,
#   CAR_window <dbl>, observed_days <int>
\end{verbatim}

\subsection{6.7.3
イベント強度別の平均CAR}\label{ux30a4ux30d9ux30f3ux30c8ux5f37ux5ea6ux5225ux306eux5e73ux5747car}

\begin{Shaded}
\begin{Highlighting}[]
\NormalTok{tbl\_event\_window\_summary }\OtherTok{\textless{}{-}}\NormalTok{ tbl\_event\_window\_return }\SpecialCharTok{|\textgreater{}}
  \FunctionTok{group\_by}\NormalTok{(}
\NormalTok{    window\_ID,}
\NormalTok{    window\_label,}
\NormalTok{    start\_day,}
\NormalTok{    end\_day,}
\NormalTok{    event\_strength}
\NormalTok{  ) }\SpecialCharTok{|\textgreater{}}
  \FunctionTok{summarise}\NormalTok{(}
    \AttributeTok{events =} \FunctionTok{n\_distinct}\NormalTok{(event\_ID),}
    \AttributeTok{mean\_CAR =} \FunctionTok{mean}\NormalTok{(CAR\_window),}
    \AttributeTok{median\_CAR =} \FunctionTok{median}\NormalTok{(CAR\_window),}
    \AttributeTok{standard\_deviation\_CAR =}
      \FunctionTok{sd}\NormalTok{(CAR\_window),}
    \AttributeTok{positive\_CAR\_ratio =}
      \FunctionTok{mean}\NormalTok{(CAR\_window }\SpecialCharTok{\textgreater{}} \DecValTok{0}\NormalTok{),}
    \AttributeTok{.groups =} \StringTok{"drop"}
\NormalTok{  ) }\SpecialCharTok{|\textgreater{}}
  \FunctionTok{mutate}\NormalTok{(}
    \AttributeTok{standard\_error\_CAR =}
\NormalTok{      standard\_deviation\_CAR }\SpecialCharTok{/}
      \FunctionTok{sqrt}\NormalTok{(events),}
    \AttributeTok{t\_value =} \FunctionTok{if\_else}\NormalTok{(}
      \FunctionTok{is.finite}\NormalTok{(standard\_error\_CAR) }\SpecialCharTok{\&}
\NormalTok{        standard\_error\_CAR }\SpecialCharTok{\textgreater{}} \DecValTok{0}\NormalTok{,}
\NormalTok{      mean\_CAR }\SpecialCharTok{/}
\NormalTok{        standard\_error\_CAR,}
      \ConstantTok{NA\_real\_}
\NormalTok{    ),}
    \AttributeTok{p\_value =} \FunctionTok{if\_else}\NormalTok{(}
      \FunctionTok{is.finite}\NormalTok{(t\_value),}
      \DecValTok{2} \SpecialCharTok{*} \FunctionTok{pt}\NormalTok{(}
        \SpecialCharTok{{-}}\FunctionTok{abs}\NormalTok{(t\_value),}
        \AttributeTok{df =}\NormalTok{ events }\SpecialCharTok{{-}} \DecValTok{1}\NormalTok{L}
\NormalTok{      ),}
      \ConstantTok{NA\_real\_}
\NormalTok{    ),}
    \AttributeTok{significance =}
      \FunctionTok{significance\_stars}\NormalTok{(p\_value)}
\NormalTok{  ) }\SpecialCharTok{|\textgreater{}}
  \FunctionTok{arrange}\NormalTok{(}
\NormalTok{    window\_ID,}
\NormalTok{    event\_strength}
\NormalTok{  )}

\NormalTok{tbl\_event\_window\_summary}
\end{Highlighting}
\end{Shaded}

\begin{verbatim}
# A tibble: 25 x 14
   window_ID window_label start_day end_day event_strength events  mean_CAR
   <chr>     <chr>            <int>   <int> <ord>           <int>     <dbl>
 1 0_p1      [0, +1]              0       1 1                 793  0.00680 
 2 0_p1      [0, +1]              0       1 2                 793  0.00323 
 3 0_p1      [0, +1]              0       1 3                 792  0.000579
 4 0_p1      [0, +1]              0       1 4                 791 -0.00328 
 5 0_p1      [0, +1]              0       1 5                 791 -0.00629 
 6 m10_p10   [-10, +10]         -10      10 1                 793  0.00936 
 7 m10_p10   [-10, +10]         -10      10 2                 793  0.00476 
 8 m10_p10   [-10, +10]         -10      10 3                 792 -0.000173
 9 m10_p10   [-10, +10]         -10      10 4                 791 -0.00301 
10 m10_p10   [-10, +10]         -10      10 5                 791 -0.00942 
# i 15 more rows
# i 7 more variables: median_CAR <dbl>, standard_deviation_CAR <dbl>,
#   positive_CAR_ratio <dbl>, standard_error_CAR <dbl>, t_value <dbl>,
#   p_value <dbl>, significance <chr>
\end{verbatim}

\subsection{6.7.4
短期CARの分布}\label{ux77edux671fcarux306eux5206ux5e03}

\begin{Shaded}
\begin{Highlighting}[]
\NormalTok{plot\_short\_window\_car }\OtherTok{\textless{}{-}}\NormalTok{ tbl\_event\_window\_return }\SpecialCharTok{|\textgreater{}}
  \FunctionTok{filter}\NormalTok{(window\_ID }\SpecialCharTok{==} \StringTok{"m1\_p1"}\NormalTok{) }\SpecialCharTok{|\textgreater{}}
  \FunctionTok{ggplot}\NormalTok{(}
    \FunctionTok{aes}\NormalTok{(}
      \AttributeTok{x =}\NormalTok{ event\_strength,}
      \AttributeTok{y =}\NormalTok{ CAR\_window}
\NormalTok{    )}
\NormalTok{  ) }\SpecialCharTok{+}
  \FunctionTok{geom\_boxplot}\NormalTok{(}
    \AttributeTok{outlier.alpha =} \FloatTok{0.25}
\NormalTok{  ) }\SpecialCharTok{+}
  \FunctionTok{geom\_hline}\NormalTok{(}
    \AttributeTok{yintercept =} \DecValTok{0}\NormalTok{,}
    \AttributeTok{linetype =} \StringTok{"dotted"}
\NormalTok{  ) }\SpecialCharTok{+}
  \FunctionTok{scale\_y\_continuous}\NormalTok{(}
    \AttributeTok{labels =}\NormalTok{ scales}\SpecialCharTok{::}\FunctionTok{label\_percent}\NormalTok{()}
\NormalTok{  ) }\SpecialCharTok{+}
  \FunctionTok{labs}\NormalTok{(}
    \AttributeTok{x =} \StringTok{"Event strength"}\NormalTok{,}
    \AttributeTok{y =} \StringTok{"CAR [{-}1, +1]"}\NormalTok{,}
    \AttributeTok{title =} \StringTok{"Short{-}Window CAR by Earnings Surprise"}
\NormalTok{  ) }\SpecialCharTok{+}
  \FunctionTok{theme\_classic}\NormalTok{()}

\NormalTok{plot\_short\_window\_car}
\end{Highlighting}
\end{Shaded}

\pandocbounded{\includegraphics[keepaspectratio]{part1_empirical/ch06_event_study_files/figure-pdf/plot-short-window-car-1.pdf}}

\subsection{6.7.5
ウィンドウ別平均CAR}\label{ux30a6ux30a3ux30f3ux30c9ux30a6ux5225ux5e73ux5747car}

\begin{Shaded}
\begin{Highlighting}[]
\NormalTok{plot\_window\_mean\_car }\OtherTok{\textless{}{-}}\NormalTok{ tbl\_event\_window\_summary }\SpecialCharTok{|\textgreater{}}
  \FunctionTok{ggplot}\NormalTok{(}
    \FunctionTok{aes}\NormalTok{(}
      \AttributeTok{x =}\NormalTok{ event\_strength,}
      \AttributeTok{y =}\NormalTok{ mean\_CAR,}
      \AttributeTok{fill =}\NormalTok{ window\_label}
\NormalTok{    )}
\NormalTok{  ) }\SpecialCharTok{+}
  \FunctionTok{geom\_col}\NormalTok{(}
    \AttributeTok{position =} \StringTok{"dodge"}
\NormalTok{  ) }\SpecialCharTok{+}
  \FunctionTok{geom\_hline}\NormalTok{(}
    \AttributeTok{yintercept =} \DecValTok{0}\NormalTok{,}
    \AttributeTok{linetype =} \StringTok{"dotted"}
\NormalTok{  ) }\SpecialCharTok{+}
  \FunctionTok{scale\_y\_continuous}\NormalTok{(}
    \AttributeTok{labels =}\NormalTok{ scales}\SpecialCharTok{::}\FunctionTok{label\_percent}\NormalTok{()}
\NormalTok{  ) }\SpecialCharTok{+}
  \FunctionTok{labs}\NormalTok{(}
    \AttributeTok{x =} \StringTok{"Event strength"}\NormalTok{,}
    \AttributeTok{y =} \StringTok{"Mean CAR"}\NormalTok{,}
    \AttributeTok{fill =} \StringTok{"Window"}\NormalTok{,}
    \AttributeTok{title =} \StringTok{"Mean CAR across Event Windows"}
\NormalTok{  ) }\SpecialCharTok{+}
  \FunctionTok{theme\_classic}\NormalTok{()}

\NormalTok{plot\_window\_mean\_car}
\end{Highlighting}
\end{Shaded}

\pandocbounded{\includegraphics[keepaspectratio]{part1_empirical/ch06_event_study_files/figure-pdf/plot-window-mean-car-1.pdf}}

\subsection{6.7.6
正のCARとなったイベントの割合}\label{ux6b63ux306ecarux3068ux306aux3063ux305fux30a4ux30d9ux30f3ux30c8ux306eux5272ux5408}

\begin{Shaded}
\begin{Highlighting}[]
\NormalTok{plot\_positive\_car\_ratio }\OtherTok{\textless{}{-}}\NormalTok{ tbl\_event\_window\_summary }\SpecialCharTok{|\textgreater{}}
  \FunctionTok{ggplot}\NormalTok{(}
    \FunctionTok{aes}\NormalTok{(}
      \AttributeTok{x =}\NormalTok{ event\_strength,}
      \AttributeTok{y =}\NormalTok{ positive\_CAR\_ratio,}
      \AttributeTok{fill =}\NormalTok{ window\_label}
\NormalTok{    )}
\NormalTok{  ) }\SpecialCharTok{+}
  \FunctionTok{geom\_col}\NormalTok{(}
    \AttributeTok{position =} \StringTok{"dodge"}
\NormalTok{  ) }\SpecialCharTok{+}
  \FunctionTok{geom\_hline}\NormalTok{(}
    \AttributeTok{yintercept =} \FloatTok{0.5}\NormalTok{,}
    \AttributeTok{linetype =} \StringTok{"dotted"}
\NormalTok{  ) }\SpecialCharTok{+}
  \FunctionTok{scale\_y\_continuous}\NormalTok{(}
    \AttributeTok{labels =}\NormalTok{ scales}\SpecialCharTok{::}\FunctionTok{label\_percent}\NormalTok{(),}
    \AttributeTok{limits =} \FunctionTok{c}\NormalTok{(}\DecValTok{0}\NormalTok{, }\DecValTok{1}\NormalTok{)}
\NormalTok{  ) }\SpecialCharTok{+}
  \FunctionTok{labs}\NormalTok{(}
    \AttributeTok{x =} \StringTok{"Event strength"}\NormalTok{,}
    \AttributeTok{y =} \StringTok{"Positive CAR ratio"}\NormalTok{,}
    \AttributeTok{fill =} \StringTok{"Window"}\NormalTok{,}
    \AttributeTok{title =} \StringTok{"Share of Events with Positive CAR"}
\NormalTok{  ) }\SpecialCharTok{+}
  \FunctionTok{theme\_classic}\NormalTok{()}

\NormalTok{plot\_positive\_car\_ratio}
\end{Highlighting}
\end{Shaded}

\pandocbounded{\includegraphics[keepaspectratio]{part1_empirical/ch06_event_study_files/figure-pdf/plot-positive-car-ratio-1.pdf}}

\subsection{6.7.7
平均CARのヒートマップ}\label{ux5e73ux5747carux306eux30d2ux30fcux30c8ux30deux30c3ux30d7}

\begin{Shaded}
\begin{Highlighting}[]
\NormalTok{plot\_window\_car\_heatmap }\OtherTok{\textless{}{-}}\NormalTok{ tbl\_event\_window\_summary }\SpecialCharTok{|\textgreater{}}
  \FunctionTok{ggplot}\NormalTok{(}
    \FunctionTok{aes}\NormalTok{(}
      \AttributeTok{x =}\NormalTok{ event\_strength,}
      \AttributeTok{y =}\NormalTok{ window\_label,}
      \AttributeTok{fill =}\NormalTok{ mean\_CAR}
\NormalTok{    )}
\NormalTok{  ) }\SpecialCharTok{+}
  \FunctionTok{geom\_tile}\NormalTok{() }\SpecialCharTok{+}
  \FunctionTok{geom\_text}\NormalTok{(}
    \FunctionTok{aes}\NormalTok{(}
      \AttributeTok{label =}\NormalTok{ scales}\SpecialCharTok{::}\FunctionTok{percent}\NormalTok{(}
\NormalTok{        mean\_CAR,}
        \AttributeTok{accuracy =} \FloatTok{0.1}
\NormalTok{      )}
\NormalTok{    )}
\NormalTok{  ) }\SpecialCharTok{+}
  \FunctionTok{labs}\NormalTok{(}
    \AttributeTok{x =} \StringTok{"Event strength"}\NormalTok{,}
    \AttributeTok{y =} \StringTok{"Event window"}\NormalTok{,}
    \AttributeTok{fill =} \StringTok{"Mean CAR"}\NormalTok{,}
    \AttributeTok{title =} \StringTok{"Mean CAR by Event Strength and Event Window"}
\NormalTok{  ) }\SpecialCharTok{+}
  \FunctionTok{theme\_classic}\NormalTok{()}

\NormalTok{plot\_window\_car\_heatmap}
\end{Highlighting}
\end{Shaded}

\pandocbounded{\includegraphics[keepaspectratio]{part1_empirical/ch06_event_study_files/figure-pdf/plot-window-car-heatmap-1.pdf}}

\section{6.8
決算サプライズと短期CAR}\label{ux6c7aux7b97ux30b5ux30d7ux30e9ux30a4ux30baux3068ux77edux671fcar}

\subsection{6.8.1 回帰分析}\label{ux56deux5e30ux5206ux6790}

\begin{Shaded}
\begin{Highlighting}[]
\NormalTok{tbl\_surprise\_regression }\OtherTok{\textless{}{-}}\NormalTok{ tbl\_event\_window\_return }\SpecialCharTok{|\textgreater{}}
  \FunctionTok{filter}\NormalTok{(}
\NormalTok{    window\_ID }\SpecialCharTok{==} \StringTok{"m1\_p1"}\NormalTok{,}
    \FunctionTok{is.finite}\NormalTok{(earnings\_surprise),}
    \FunctionTok{is.finite}\NormalTok{(CAR\_window)}
\NormalTok{  )}

\NormalTok{model\_surprise\_car }\OtherTok{\textless{}{-}} \FunctionTok{lm}\NormalTok{(}
\NormalTok{  CAR\_window }\SpecialCharTok{\textasciitilde{}}\NormalTok{ earnings\_surprise,}
  \AttributeTok{data =}\NormalTok{ tbl\_surprise\_regression}
\NormalTok{)}

\NormalTok{broom}\SpecialCharTok{::}\FunctionTok{tidy}\NormalTok{(model\_surprise\_car)}
\end{Highlighting}
\end{Shaded}

\begin{verbatim}
# A tibble: 2 x 5
  term              estimate std.error statistic  p.value
  <chr>                <dbl>     <dbl>     <dbl>    <dbl>
1 (Intercept)       -0.00668  0.000416     -16.1 2.66e-56
2 earnings_surprise -1.37     0.0691       -19.8 1.19e-83
\end{verbatim}

\begin{Shaded}
\begin{Highlighting}[]
\NormalTok{broom}\SpecialCharTok{::}\FunctionTok{glance}\NormalTok{(model\_surprise\_car)}
\end{Highlighting}
\end{Shaded}

\begin{verbatim}
# A tibble: 1 x 12
  r.squared adj.r.squared  sigma statistic  p.value    df logLik     AIC     BIC
      <dbl>         <dbl>  <dbl>     <dbl>    <dbl> <dbl>  <dbl>   <dbl>   <dbl>
1    0.0905        0.0903 0.0145      394. 1.19e-83     1 11144. -22282. -22263.
# i 3 more variables: deviance <dbl>, df.residual <int>, nobs <int>
\end{verbatim}

\subsection{6.8.2
決算サプライズと短期CAR}\label{ux6c7aux7b97ux30b5ux30d7ux30e9ux30a4ux30baux3068ux77edux671fcar-1}

\begin{Shaded}
\begin{Highlighting}[]
\NormalTok{plot\_surprise\_car }\OtherTok{\textless{}{-}}\NormalTok{ tbl\_surprise\_regression }\SpecialCharTok{|\textgreater{}}
  \FunctionTok{ggplot}\NormalTok{(}
    \FunctionTok{aes}\NormalTok{(}
      \AttributeTok{x =}\NormalTok{ earnings\_surprise,}
      \AttributeTok{y =}\NormalTok{ CAR\_window,}
      \AttributeTok{shape =}\NormalTok{ event\_strength}
\NormalTok{    )}
\NormalTok{  ) }\SpecialCharTok{+}
  \FunctionTok{geom\_point}\NormalTok{(}\AttributeTok{alpha =} \FloatTok{0.45}\NormalTok{) }\SpecialCharTok{+}
  \FunctionTok{geom\_smooth}\NormalTok{(}
    \AttributeTok{method =} \StringTok{"lm"}\NormalTok{,}
    \AttributeTok{se =} \ConstantTok{TRUE}
\NormalTok{  ) }\SpecialCharTok{+}
  \FunctionTok{geom\_hline}\NormalTok{(}
    \AttributeTok{yintercept =} \DecValTok{0}\NormalTok{,}
    \AttributeTok{linetype =} \StringTok{"dotted"}
\NormalTok{  ) }\SpecialCharTok{+}
  \FunctionTok{scale\_y\_continuous}\NormalTok{(}
    \AttributeTok{labels =}\NormalTok{ scales}\SpecialCharTok{::}\FunctionTok{label\_percent}\NormalTok{()}
\NormalTok{  ) }\SpecialCharTok{+}
  \FunctionTok{labs}\NormalTok{(}
    \AttributeTok{x =} \StringTok{"Standardized earnings surprise"}\NormalTok{,}
    \AttributeTok{y =} \StringTok{"CAR [{-}1, +1]"}\NormalTok{,}
    \AttributeTok{shape =} \StringTok{"Event strength"}\NormalTok{,}
    \AttributeTok{title =} \StringTok{"Earnings Surprise and Short{-}Window CAR"}
\NormalTok{  ) }\SpecialCharTok{+}
  \FunctionTok{theme\_classic}\NormalTok{()}

\NormalTok{plot\_surprise\_car}
\end{Highlighting}
\end{Shaded}

\pandocbounded{\includegraphics[keepaspectratio]{part1_empirical/ch06_event_study_files/figure-pdf/plot-surprise-car-1.pdf}}

\subsection{6.8.3
年度別の関係}\label{ux5e74ux5ea6ux5225ux306eux95a2ux4fc2}

\begin{Shaded}
\begin{Highlighting}[]
\NormalTok{plot\_surprise\_car\_by\_year }\OtherTok{\textless{}{-}}\NormalTok{ tbl\_surprise\_regression }\SpecialCharTok{|\textgreater{}}
  \FunctionTok{ggplot}\NormalTok{(}
    \FunctionTok{aes}\NormalTok{(}
      \AttributeTok{x =}\NormalTok{ earnings\_surprise,}
      \AttributeTok{y =}\NormalTok{ CAR\_window}
\NormalTok{    )}
\NormalTok{  ) }\SpecialCharTok{+}
  \FunctionTok{geom\_point}\NormalTok{(}\AttributeTok{alpha =} \FloatTok{0.35}\NormalTok{) }\SpecialCharTok{+}
  \FunctionTok{geom\_smooth}\NormalTok{(}
    \AttributeTok{method =} \StringTok{"lm"}\NormalTok{,}
    \AttributeTok{se =} \ConstantTok{FALSE}
\NormalTok{  ) }\SpecialCharTok{+}
  \FunctionTok{facet\_wrap}\NormalTok{(}
    \FunctionTok{vars}\NormalTok{(event\_year),}
    \AttributeTok{scales =} \StringTok{"free\_x"}
\NormalTok{  ) }\SpecialCharTok{+}
  \FunctionTok{geom\_hline}\NormalTok{(}
    \AttributeTok{yintercept =} \DecValTok{0}\NormalTok{,}
    \AttributeTok{linetype =} \StringTok{"dotted"}
\NormalTok{  ) }\SpecialCharTok{+}
  \FunctionTok{scale\_y\_continuous}\NormalTok{(}
    \AttributeTok{labels =}\NormalTok{ scales}\SpecialCharTok{::}\FunctionTok{label\_percent}\NormalTok{()}
\NormalTok{  ) }\SpecialCharTok{+}
  \FunctionTok{labs}\NormalTok{(}
    \AttributeTok{x =} \StringTok{"Standardized earnings surprise"}\NormalTok{,}
    \AttributeTok{y =} \StringTok{"CAR [{-}1, +1]"}\NormalTok{,}
    \AttributeTok{title =} \StringTok{"Earnings Surprise and Short{-}Window CAR by Year"}
\NormalTok{  ) }\SpecialCharTok{+}
  \FunctionTok{theme\_classic}\NormalTok{()}

\NormalTok{plot\_surprise\_car\_by\_year}
\end{Highlighting}
\end{Shaded}

\pandocbounded{\includegraphics[keepaspectratio]{part1_empirical/ch06_event_study_files/figure-pdf/plot-surprise-car-by-year-1.pdf}}

\section{6.9
分析結果の保存}\label{ux5206ux6790ux7d50ux679cux306eux4fddux5b58}

\subsection{6.9.1
イベント単位の結果}\label{ux30a4ux30d9ux30f3ux30c8ux5358ux4f4dux306eux7d50ux679c}

\begin{Shaded}
\begin{Highlighting}[]
\NormalTok{tbl\_event\_window\_wide }\OtherTok{\textless{}{-}}\NormalTok{ tbl\_event\_window\_return }\SpecialCharTok{|\textgreater{}}
  \FunctionTok{select}\NormalTok{(}
\NormalTok{    event\_ID,}
\NormalTok{    window\_ID,}
\NormalTok{    CAR\_window}
\NormalTok{  ) }\SpecialCharTok{|\textgreater{}}
  \FunctionTok{pivot\_wider}\NormalTok{(}
    \AttributeTok{names\_from =}\NormalTok{ window\_ID,}
    \AttributeTok{values\_from =}\NormalTok{ CAR\_window,}
    \AttributeTok{names\_prefix =} \StringTok{"CAR\_"}
\NormalTok{  )}

\NormalTok{tbl\_event\_output }\OtherTok{\textless{}{-}}\NormalTok{ tbl\_market\_model }\SpecialCharTok{|\textgreater{}}
  \FunctionTok{semi\_join}\NormalTok{(}
\NormalTok{    tbl\_abnormal\_return }\SpecialCharTok{|\textgreater{}}
      \FunctionTok{distinct}\NormalTok{(event\_ID),}
    \AttributeTok{by =} \StringTok{"event\_ID"}
\NormalTok{  ) }\SpecialCharTok{|\textgreater{}}
  \FunctionTok{left\_join}\NormalTok{(}
\NormalTok{    tbl\_event\_window\_wide,}
    \AttributeTok{by =} \StringTok{"event\_ID"}
\NormalTok{  ) }\SpecialCharTok{|\textgreater{}}
  \FunctionTok{arrange}\NormalTok{(event\_ID)}
\end{Highlighting}
\end{Shaded}

\subsection{6.9.2 日次結果}\label{ux65e5ux6b21ux7d50ux679c}

\begin{Shaded}
\begin{Highlighting}[]
\NormalTok{tbl\_event\_daily\_output }\OtherTok{\textless{}{-}}\NormalTok{ tbl\_abnormal\_return }\SpecialCharTok{|\textgreater{}}
  \FunctionTok{select}\NormalTok{(}
\NormalTok{    event\_ID,}
\NormalTok{    firm\_ID,}
\NormalTok{    event\_date,}
\NormalTok{    event\_trade\_date,}
\NormalTok{    event\_year,}
\NormalTok{    event\_strength,}
\NormalTok{    earnings\_surprise,}
\NormalTok{    relative\_day,}
\NormalTok{    date,}
\NormalTok{    R,}
\NormalTok{    R\_M,}
\NormalTok{    normal\_return,}
\NormalTok{    AR,}
\NormalTok{    CAR,}
\NormalTok{    alpha,}
\NormalTok{    beta,}
\NormalTok{    sigma\_AR,}
\NormalTok{    r\_squared,}
\NormalTok{    estimation\_observations}
\NormalTok{  )}

\NormalTok{tbl\_event\_summary\_output }\OtherTok{\textless{}{-}}\NormalTok{ tbl\_daily\_statistics }\SpecialCharTok{|\textgreater{}}
  \FunctionTok{arrange}\NormalTok{(}
\NormalTok{    event\_strength,}
\NormalTok{    relative\_day}
\NormalTok{  )}
\end{Highlighting}
\end{Shaded}

\subsection{6.9.3 CSV保存}\label{csvux4fddux5b58}

\begin{Shaded}
\begin{Highlighting}[]
\NormalTok{list\_output }\OtherTok{\textless{}{-}} \FunctionTok{list}\NormalTok{(}
  \AttributeTok{event\_results =}\NormalTok{ tbl\_event\_output,}
  \AttributeTok{event\_daily =}\NormalTok{ tbl\_event\_daily\_output,}
  \AttributeTok{event\_summary =}\NormalTok{ tbl\_event\_summary\_output}
\NormalTok{)}

\NormalTok{vec\_output\_path }\OtherTok{\textless{}{-}} \FunctionTok{c}\NormalTok{(}
  \AttributeTok{event\_results =} \FunctionTok{file.path}\NormalTok{(}
\NormalTok{    derived\_data\_dir,}
    \StringTok{"ch08\_event\_results.csv"}
\NormalTok{  ),}
  \AttributeTok{event\_daily =} \FunctionTok{file.path}\NormalTok{(}
\NormalTok{    derived\_data\_dir,}
    \StringTok{"ch08\_event\_daily.csv"}
\NormalTok{  ),}
  \AttributeTok{event\_summary =} \FunctionTok{file.path}\NormalTok{(}
\NormalTok{    derived\_data\_dir,}
    \StringTok{"ch08\_event\_summary.csv"}
\NormalTok{  )}
\NormalTok{)}

\NormalTok{purrr}\SpecialCharTok{::}\FunctionTok{walk2}\NormalTok{(}
\NormalTok{  list\_output,}
\NormalTok{  vec\_output\_path,}
\NormalTok{  readr}\SpecialCharTok{::}\NormalTok{write\_csv}
\NormalTok{)}

\FunctionTok{tibble}\NormalTok{(}
  \AttributeTok{output\_name =} \FunctionTok{names}\NormalTok{(vec\_output\_path),}
  \AttributeTok{path =} \FunctionTok{unname}\NormalTok{(vec\_output\_path)}
\NormalTok{)}
\end{Highlighting}
\end{Shaded}

\begin{verbatim}
# A tibble: 3 x 2
  output_name   path                                                            
  <chr>         <chr>                                                           
1 event_results C:/AnalyticFin/Projects/LagrangeFinance/simulation_data/derived~
2 event_daily   C:/AnalyticFin/Projects/LagrangeFinance/simulation_data/derived~
3 event_summary C:/AnalyticFin/Projects/LagrangeFinance/simulation_data/derived~
\end{verbatim}

\subsection{6.9.4
保存結果の検証}\label{ux4fddux5b58ux7d50ux679cux306eux691cux8a3c}

\begin{Shaded}
\begin{Highlighting}[]
\NormalTok{tbl\_output\_verification }\OtherTok{\textless{}{-}}\NormalTok{ purrr}\SpecialCharTok{::}\FunctionTok{imap\_dfr}\NormalTok{(}
\NormalTok{  list\_output,}
\NormalTok{  \textbackslash{}(tbl\_expected, output\_name) \{}
\NormalTok{    output\_path }\OtherTok{\textless{}{-}}\NormalTok{ vec\_output\_path[[output\_name]]}

\NormalTok{    tbl\_check }\OtherTok{\textless{}{-}}\NormalTok{ readr}\SpecialCharTok{::}\FunctionTok{read\_csv}\NormalTok{(}
\NormalTok{      output\_path,}
      \AttributeTok{show\_col\_types =} \ConstantTok{FALSE}
\NormalTok{    )}

    \FunctionTok{stopifnot}\NormalTok{(}
      \FunctionTok{nrow}\NormalTok{(tbl\_check) }\SpecialCharTok{==}
        \FunctionTok{nrow}\NormalTok{(tbl\_expected)}
\NormalTok{    )}

    \FunctionTok{stopifnot}\NormalTok{(}
      \FunctionTok{identical}\NormalTok{(}
        \FunctionTok{names}\NormalTok{(tbl\_check),}
        \FunctionTok{names}\NormalTok{(tbl\_expected)}
\NormalTok{      )}
\NormalTok{    )}

    \FunctionTok{tibble}\NormalTok{(}
      \AttributeTok{output\_name =}\NormalTok{ output\_name,}
      \AttributeTok{output\_file =}\NormalTok{ output\_path,}
      \AttributeTok{rows =} \FunctionTok{nrow}\NormalTok{(tbl\_check),}
      \AttributeTok{columns =} \FunctionTok{ncol}\NormalTok{(tbl\_check),}
      \AttributeTok{size\_bytes =}
        \FunctionTok{file.info}\NormalTok{(output\_path)}\SpecialCharTok{$}\NormalTok{size}
\NormalTok{    )}
\NormalTok{  \}}
\NormalTok{)}

\NormalTok{tbl\_output\_verification}
\end{Highlighting}
\end{Shaded}

\begin{verbatim}
# A tibble: 3 x 5
  output_name   output_file                              rows columns size_bytes
  <chr>         <chr>                                   <int>   <int>      <dbl>
1 event_results C:/AnalyticFin/Projects/LagrangeFinan~   3960      18    1067043
2 event_daily   C:/AnalyticFin/Projects/LagrangeFinan~ 241560      19   63708974
3 event_summary C:/AnalyticFin/Projects/LagrangeFinan~    305      15      67027
\end{verbatim}

\section{6.10 まとめ}\label{ux307eux3068ux3081-5}

本章では、決算発表の翌取引日をイベント日0として、マーケット・モデルによるイベントスタディを実施した。推定期間にはイベント日の130取引日前から31取引日前までの100日を使用し、イベント期間には30取引日前から30取引日後までの61日を使用した。

決算サプライズは、実現利益と予想利益の差を期首時価総額で標準化し、年度ごとに5群へ分類した。イベント別にマーケット・モデルを推定し、通常リターン、異常リターン、累積異常リターンを計算した。さらに、日次のAARとCAAR、複数イベント・ウィンドウのCAR、正のCARとなったイベントの割合、決算サプライズと短期CARの回帰関係を確認した。

反復処理には\texttt{purrr}を利用し、明示的なループは使用していない。データ探索、入力検証、イベント日0の設定、イベント・パネルの作成、マーケット・モデルの推定、AR・CARの計算、イベント・ウィンドウ集計は、\texttt{utility\_LFL.R}の共通関数へ分離した。

結果は次の3ファイルに保存される。

\begin{Shaded}
\begin{Highlighting}[]
\NormalTok{simulation\_data/derived/ch08\_event\_results.csv}
\NormalTok{simulation\_data/derived/ch08\_event\_daily.csv}
\NormalTok{simulation\_data/derived/ch08\_event\_summary.csv}
\end{Highlighting}
\end{Shaded}

\part{第II部 金融工学　無裁定モデル}

\chapter{第II-1章
確定キャッシュフロー債券評価}\label{ux7b2cii-1ux7ae0-ux78baux5b9aux30adux30e3ux30c3ux30b7ux30e5ux30d5ux30edux30fcux50b5ux5238ux8a55ux4fa1}

金融では、異なる時点に発生する金額をそのまま比較することはできない。今日受け取る100円と1年後に受け取る100円では、運用可能期間や金利が異なるため、経済的な価値も異なるからである。そこで、時点の異なるキャッシュフローを評価する場合には、将来の金額を現在価値へ割り引くか、現在の金額を将来価値へ換算し、同一時点の価値にそろえる必要がある。金融商品の価格評価では、通常、将来のすべてのキャッシュフローを評価時点の現在価値へ換算する。

\begin{Shaded}
\begin{Highlighting}[]
\NormalTok{knitr}\SpecialCharTok{::}\NormalTok{opts\_chunk}\SpecialCharTok{$}\FunctionTok{set}\NormalTok{(}
  \AttributeTok{collapse =} \ConstantTok{TRUE}\NormalTok{,}
  \AttributeTok{comment =} \StringTok{"\#\textgreater{}"}\NormalTok{,}
  \AttributeTok{warning =} \ConstantTok{FALSE}\NormalTok{,}
  \AttributeTok{message =} \ConstantTok{FALSE}
\NormalTok{)}


\FunctionTok{suppressPackageStartupMessages}\NormalTok{(\{}
  \FunctionTok{library}\NormalTok{(tidyverse)}
  \FunctionTok{library}\NormalTok{(ggthemes)}
  \FunctionTok{library}\NormalTok{(QuantLib)}
  \FunctionTok{library}\NormalTok{(GaussQuant)}
  \FunctionTok{library}\NormalTok{(tidyverse)}
  \FunctionTok{library}\NormalTok{(broom)}
\NormalTok{\})}
\end{Highlighting}
\end{Shaded}

\section{1.1
金利期間構造の確認}\label{ux91d1ux5229ux671fux9593ux69cbux9020ux306eux78baux8a8d}

債券価格は、将来の各キャッシュフローを支払時点に対応する割引係数で現在価値へ換算し、その合計として求める。個別債券の評価に入る前に、金利期間構造をスポットレート、割引係数およびフォワードレートの三つの側面から確認する。

市場で直接観測される金利は、必ずしも各満期に対応するスポットレートではない。市場では、預金金利、OISレート、スワップレート、国債の価格や複利利回りなどが観測される。これらの市場情報と、日数計算規則、利払頻度および受渡日などのマーケットコンベンションを用いて、各満期のゼロクーポン・スポットレートとディスカウントファクター（以下、DF）を順次求める。この手続きをブートストラップという。

スポットレートは、評価時点から特定の満期までのゼロクーポン金利を表す。DFは、将来の1単位の金額を評価時点の価値へ換算する倍率であり、債券、スワップおよびデリバティブのキャッシュフロー評価に直接用いられる。さらに、異なる満期のDFを比較することで、将来の一定期間に適用されるフォワードレートを求めることができる。このように、市場で観測される複数の金利商品からスポットレート、DFおよびフォワードレートからなる金利期間構造を構築することは、金融工学における価格評価の基礎となる。

以下では、EONIAの市場データから構築した金利曲線を用いて、スポットカーブ、DFカーブおよびフォワードカーブの関係を確認する。さらに、本章では、固定利付債を例として、クリーン価格、ダーティ価格、経過利子、最終利回り、デュレーションおよびコンベクシティを計算する。あわせて、日数計算規則、休日カレンダー、営業日調整、受渡日および利払スケジュールといったマーケットコンベンションを確認する。

\begin{Shaded}
\begin{Highlighting}[]
\NormalTok{eonia\_benchmark }\OtherTok{\textless{}{-}}\NormalTok{ GaussQuant}\SpecialCharTok{::}\FunctionTok{eonia\_curve\_benchmark\_GQL}\NormalTok{()}

\NormalTok{eonia\_curve\_tbl }\OtherTok{\textless{}{-}}\NormalTok{ eonia\_benchmark}\SpecialCharTok{$}\NormalTok{final\_validation}\SpecialCharTok{$}\NormalTok{quote\_repricing }\SpecialCharTok{|\textgreater{}}
\NormalTok{  dplyr}\SpecialCharTok{::}\FunctionTok{select}\NormalTok{(}
\NormalTok{    pillar\_date,}
\NormalTok{    discount\_factor,}
\NormalTok{    zero\_rate,}
\NormalTok{    one\_day\_forward\_rate}
\NormalTok{  ) }\SpecialCharTok{|\textgreater{}}
\NormalTok{  dplyr}\SpecialCharTok{::}\FunctionTok{filter}\NormalTok{(}\SpecialCharTok{!}\FunctionTok{is.na}\NormalTok{(.data}\SpecialCharTok{$}\NormalTok{pillar\_date)) }\SpecialCharTok{|\textgreater{}}
\NormalTok{  dplyr}\SpecialCharTok{::}\FunctionTok{distinct}\NormalTok{(.data}\SpecialCharTok{$}\NormalTok{pillar\_date, }\AttributeTok{.keep\_all =} \ConstantTok{TRUE}\NormalTok{) }\SpecialCharTok{|\textgreater{}}
\NormalTok{  dplyr}\SpecialCharTok{::}\FunctionTok{arrange}\NormalTok{(.data}\SpecialCharTok{$}\NormalTok{pillar\_date)}

\NormalTok{spot\_curve\_plot }\OtherTok{\textless{}{-}}\NormalTok{ eonia\_curve\_tbl }\SpecialCharTok{|\textgreater{}}
\NormalTok{  dplyr}\SpecialCharTok{::}\FunctionTok{filter}\NormalTok{(}\SpecialCharTok{!}\FunctionTok{is.na}\NormalTok{(.data}\SpecialCharTok{$}\NormalTok{zero\_rate)) }\SpecialCharTok{|\textgreater{}}
\NormalTok{  ggplot2}\SpecialCharTok{::}\FunctionTok{ggplot}\NormalTok{(}
\NormalTok{    ggplot2}\SpecialCharTok{::}\FunctionTok{aes}\NormalTok{(}
      \AttributeTok{x =}\NormalTok{ .data}\SpecialCharTok{$}\NormalTok{pillar\_date,}
      \AttributeTok{y =} \DecValTok{100} \SpecialCharTok{*}\NormalTok{ .data}\SpecialCharTok{$}\NormalTok{zero\_rate}
\NormalTok{    )}
\NormalTok{  ) }\SpecialCharTok{+}
\NormalTok{  ggplot2}\SpecialCharTok{::}\FunctionTok{geom\_line}\NormalTok{(}\AttributeTok{linewidth =} \FloatTok{0.7}\NormalTok{) }\SpecialCharTok{+}
\NormalTok{  ggplot2}\SpecialCharTok{::}\FunctionTok{geom\_point}\NormalTok{(}\AttributeTok{size =} \FloatTok{1.5}\NormalTok{) }\SpecialCharTok{+}
\NormalTok{  ggplot2}\SpecialCharTok{::}\FunctionTok{scale\_x\_date}\NormalTok{(}
    \AttributeTok{date\_breaks =} \StringTok{"5 years"}\NormalTok{,}
    \AttributeTok{date\_labels =} \StringTok{"\%Y"}
\NormalTok{  ) }\SpecialCharTok{+}
\NormalTok{  ggplot2}\SpecialCharTok{::}\FunctionTok{labs}\NormalTok{(}
    \AttributeTok{title =} \StringTok{"EONIAスポットカーブ"}\NormalTok{,}
    \AttributeTok{x =} \StringTok{"満期"}\NormalTok{,}
    \AttributeTok{y =} \StringTok{"スポットレート（\%）"}
\NormalTok{  ) }\SpecialCharTok{+}
\NormalTok{  ggplot2}\SpecialCharTok{::}\FunctionTok{theme\_minimal}\NormalTok{()}

\NormalTok{discount\_curve\_plot }\OtherTok{\textless{}{-}}\NormalTok{ eonia\_curve\_tbl }\SpecialCharTok{|\textgreater{}}
\NormalTok{  dplyr}\SpecialCharTok{::}\FunctionTok{filter}\NormalTok{(}\SpecialCharTok{!}\FunctionTok{is.na}\NormalTok{(.data}\SpecialCharTok{$}\NormalTok{discount\_factor)) }\SpecialCharTok{|\textgreater{}}
\NormalTok{  ggplot2}\SpecialCharTok{::}\FunctionTok{ggplot}\NormalTok{(}
\NormalTok{    ggplot2}\SpecialCharTok{::}\FunctionTok{aes}\NormalTok{(}
      \AttributeTok{x =}\NormalTok{ .data}\SpecialCharTok{$}\NormalTok{pillar\_date,}
      \AttributeTok{y =}\NormalTok{ .data}\SpecialCharTok{$}\NormalTok{discount\_factor}
\NormalTok{    )}
\NormalTok{  ) }\SpecialCharTok{+}
\NormalTok{  ggplot2}\SpecialCharTok{::}\FunctionTok{geom\_line}\NormalTok{(}\AttributeTok{linewidth =} \FloatTok{0.7}\NormalTok{) }\SpecialCharTok{+}
\NormalTok{  ggplot2}\SpecialCharTok{::}\FunctionTok{geom\_point}\NormalTok{(}\AttributeTok{size =} \FloatTok{1.5}\NormalTok{) }\SpecialCharTok{+}
\NormalTok{  ggplot2}\SpecialCharTok{::}\FunctionTok{scale\_x\_date}\NormalTok{(}
    \AttributeTok{date\_breaks =} \StringTok{"5 years"}\NormalTok{,}
    \AttributeTok{date\_labels =} \StringTok{"\%Y"}
\NormalTok{  ) }\SpecialCharTok{+}
\NormalTok{  ggplot2}\SpecialCharTok{::}\FunctionTok{labs}\NormalTok{(}
    \AttributeTok{title =} \StringTok{"EONIA割引係数カーブ"}\NormalTok{,}
    \AttributeTok{x =} \StringTok{"満期"}\NormalTok{,}
    \AttributeTok{y =} \StringTok{"割引係数"}
\NormalTok{  ) }\SpecialCharTok{+}
\NormalTok{  ggplot2}\SpecialCharTok{::}\FunctionTok{theme\_minimal}\NormalTok{()}

\NormalTok{forward\_curve\_plot }\OtherTok{\textless{}{-}}\NormalTok{ eonia\_curve\_tbl }\SpecialCharTok{|\textgreater{}}
\NormalTok{  dplyr}\SpecialCharTok{::}\FunctionTok{filter}\NormalTok{(}\SpecialCharTok{!}\FunctionTok{is.na}\NormalTok{(.data}\SpecialCharTok{$}\NormalTok{one\_day\_forward\_rate)) }\SpecialCharTok{|\textgreater{}}
\NormalTok{  ggplot2}\SpecialCharTok{::}\FunctionTok{ggplot}\NormalTok{(}
\NormalTok{    ggplot2}\SpecialCharTok{::}\FunctionTok{aes}\NormalTok{(}
      \AttributeTok{x =}\NormalTok{ .data}\SpecialCharTok{$}\NormalTok{pillar\_date,}
      \AttributeTok{y =} \DecValTok{100} \SpecialCharTok{*}\NormalTok{ .data}\SpecialCharTok{$}\NormalTok{one\_day\_forward\_rate}
\NormalTok{    )}
\NormalTok{  ) }\SpecialCharTok{+}
\NormalTok{  ggplot2}\SpecialCharTok{::}\FunctionTok{geom\_line}\NormalTok{(}\AttributeTok{linewidth =} \FloatTok{0.7}\NormalTok{) }\SpecialCharTok{+}
\NormalTok{  ggplot2}\SpecialCharTok{::}\FunctionTok{geom\_point}\NormalTok{(}\AttributeTok{size =} \FloatTok{1.5}\NormalTok{) }\SpecialCharTok{+}
\NormalTok{  ggplot2}\SpecialCharTok{::}\FunctionTok{scale\_x\_date}\NormalTok{(}
    \AttributeTok{date\_breaks =} \StringTok{"5 years"}\NormalTok{,}
    \AttributeTok{date\_labels =} \StringTok{"\%Y"}
\NormalTok{  ) }\SpecialCharTok{+}
\NormalTok{  ggplot2}\SpecialCharTok{::}\FunctionTok{labs}\NormalTok{(}
    \AttributeTok{title =} \StringTok{"EONIA 1営業日フォワードカーブ"}\NormalTok{,}
    \AttributeTok{x =} \StringTok{"開始日"}\NormalTok{,}
    \AttributeTok{y =} \StringTok{"フォワードレート（\%）"}
\NormalTok{  ) }\SpecialCharTok{+}
\NormalTok{  ggplot2}\SpecialCharTok{::}\FunctionTok{theme\_minimal}\NormalTok{()}

\FunctionTok{print}\NormalTok{(spot\_curve\_plot)}
\end{Highlighting}
\end{Shaded}

\pandocbounded{\includegraphics[keepaspectratio]{part2_quantive_riemann/ch01_deterministic_cashflow_valuation_files/figure-pdf/eonia-spot-discount-forward-curves-1.pdf}}

\begin{Shaded}
\begin{Highlighting}[]
\FunctionTok{print}\NormalTok{(discount\_curve\_plot)}
\end{Highlighting}
\end{Shaded}

\pandocbounded{\includegraphics[keepaspectratio]{part2_quantive_riemann/ch01_deterministic_cashflow_valuation_files/figure-pdf/eonia-spot-discount-forward-curves-2.pdf}}

\begin{Shaded}
\begin{Highlighting}[]
\FunctionTok{print}\NormalTok{(forward\_curve\_plot)}
\end{Highlighting}
\end{Shaded}

\pandocbounded{\includegraphics[keepaspectratio]{part2_quantive_riemann/ch01_deterministic_cashflow_valuation_files/figure-pdf/eonia-spot-discount-forward-curves-3.pdf}}

\section{1.2
マーケットコンベンション}\label{ux30deux30fcux30b1ux30c3ux30c8ux30b3ux30f3ux30d9ux30f3ux30b7ux30e7ux30f3}

債券評価では、価格や利回りの計算方法だけでなく、日付、利払頻度、休日および端数処理に関する市場慣行を確認する必要がある。同じ債券であっても、採用する日数計算規則や利回り表示方式が異なれば、経過利子、残存年数および計算価格も異なる。

\subsection{1.2.1
一般的な市場慣行}\label{ux4e00ux822cux7684ux306aux5e02ux5834ux6163ux884c}

日数計算規則（Day Count
Convention）は、二つの日付の間の期間を年数へ換算する規則である。主な方式は次のとおりである。

\begin{itemize}
\item
  \textbf{Actual/365 Fixed} 二つの日付間の実日数を365で除する。
\item
  \textbf{Actual/360}
  二つの日付間の実日数を360で除する。短期金融市場や金利商品で広く用いられる。
\item
  \textbf{Actual/Actual ICMA}
  実日数を、利払期間の実日数と年間利払回数に基づいて年換算する。固定利付債で用いられる代表的な方式である。
\item
  \textbf{Actual/Actual ISDA}
  対象期間を平年部分と閏年部分に分け、それぞれ365または366で除する。
\item
  \textbf{30/360 Bond Basis}
  1か月を30日、1年を360日として計算し、月末の日付には所定の調整を行う。
\item
  \textbf{30E/360}
  開始日または終了日が31日の場合に30日へ調整し、1か月を30日、1年を360日として計算する。
\end{itemize}

30/360方式は手計算を容易にするために考案された方式であり、半年、四半期および1か月をそれぞれ180日、90日および30日として扱える。一方、Actual方式では暦上の実日数を用いるため、同じ利率でも計算期間によって利息額が異なることがある。日数計算規則の違いは、経過利子だけでなく、割引期間や利回りにも影響する。
詳細は\href{https://en.wikipedia.org/wiki/Day_count_convention}{Day
count convention}を参照されたい。
ただし、日数計算規則はすべての計算に同じものを用いればよいわけではない。Actual/Actual
ICMAや30/360 Bond
Basisは、規則的な利払期間を基準としてクーポン額や経過利子を求める用途に適している。一方、金利期間構造では、評価日から各将来日までの距離を一貫した時間軸へ変換する必要があるため、Actual/360またはActual/365
Fixedのような単純な規則を用いることが多い。

とくにActual/Actual ICMAは、基準となる利払期間（reference
period）を必要とする。不規則な初回または最終利払期間では、同じ二つの日付であっても、基準期間の与え方によって年数が変わり得る。利払スケジュールを持たない金利カーブにこの規則をそのまま使用すると、割引期間や割引係数が不自然になる場合がある。したがって、債券のクーポン計算に用いる日数計算規則と、割引カーブの時間軸に用いる日数計算規則は区別して設定する必要がある。

受渡日や利払日が非営業日に当たる場合には、対象市場の休日カレンダーと営業日調整規則を用いる。代表的な規則には、次営業日へ移すFollowing、翌月へ越える場合には前営業日へ戻すModified
Following、前営業日へ戻すPreceding、および日付を調整しないUnadjustedがある。

\subsection{1.2.2
日本の債券市場における慣行}\label{ux65e5ux672cux306eux50b5ux5238ux5e02ux5834ux306bux304aux3051ux308bux6163ux884c}

日本の公社債市場では、単価と複利利回りに加えて、単利利回りが広く用いられてきた。日本証券業協会の公社債店頭売買参考統計値でも、原則として報告された単利利回りから単価を求め、その単価から複利利回りを計算する手順が採用されている。

利付債の経過利子に用いる経過日数は、直前利払期日から受渡日までの実日数を片端入れで数える。一方、最終利回りに用いる残存日数は、受渡日の翌日から償還日までを片端入れで数えるため、経過日数と残存日数を区別する必要がある。

残存日数を \(N\)、残存年数を \(T\) とすると、残存年数は次式で求める。

\[
T=\frac{N}{365}
\]

日本市場の実務では、残存期間が1年以上の場合には、原則として期間中の2月29日を残存日数から除外する。残存期間が1年未満の場合には2月29日を含めるが、分母は365のままとする。

さらに、計算された残存年数は小数点以下第8位で切り捨て、小数点以下第7位までの値を以後の計算に用いる。四捨五入ではなく切り捨てである点が重要である。

額面100円、年間クーポン額を \(C\)、償還価格を \(R\)、単価を
\(P\)、残存年数を \(T\) とすると、単利の最終利回り \(Y\)
は次式で表される。

\[
Y=
\frac{
C+\dfrac{R-P}{T}
}{
P
}
\times 100
\]

この単利利回りは、将来の各キャッシュフローを支払時点ごとに割り引く複利利回りとは異なる。クーポンと償還差損益を残存年数で年換算する日本市場独自の表示方法であり、キャッシュフローごとの時間価値を直接には反映しない。

利回りから単価を逆算する場合には、小数点以下第7位までに切り捨てた残存年数を使用する。算出された単価も、原則として小数点以下第4位を切り捨て、小数点以下第3位までとする。

単価で約定する取引では、計算途中の端数が約定価格を直接変更する場面は限られる。しかし、単利利回りから単価を逆算する場合には、残存年数の切り捨てが償還差損益の年換算額を変化させ、さらに単価の切り捨てによって価格差が生じる。このため、日本の債券計算では、残存日数、残存年数および単価の端数処理を正確に再現する必要がある。

なお、以下のコード例では一般的な債券評価の確認を目的として、30/360 Bond
Basis、半年複利およびUnadjustedを採用している。したがって、この設定は日本国内債の単利利回り慣行をそのまま再現するものではない。

以下では、日本の休日カレンダー、取引日から2営業日後の受渡日、30/360 Bond
Basis、半年複利および年2回払いを設定する。

\begin{Shaded}
\begin{Highlighting}[]
\NormalTok{jp\_calendar }\OtherTok{\textless{}{-}}\NormalTok{ QuantLib}\SpecialCharTok{::}\FunctionTok{Japan}\NormalTok{()}

\NormalTok{dc\_a365 }\OtherTok{\textless{}{-}}\NormalTok{ QuantLib}\SpecialCharTok{::}\FunctionTok{Actual365Fixed}\NormalTok{()}
\NormalTok{dc\_a360 }\OtherTok{\textless{}{-}}\NormalTok{ QuantLib}\SpecialCharTok{::}\FunctionTok{Actual360}\NormalTok{()}
\NormalTok{dc\_30 }\OtherTok{\textless{}{-}}\NormalTok{ GaussQuant}\SpecialCharTok{::}\FunctionTok{day\_counter\_GQL}\NormalTok{(}
  \StringTok{"Thirty360"}\NormalTok{,}
  \AttributeTok{convention =} \StringTok{"BondBasis"}
\NormalTok{)}

\NormalTok{cmpd\_compounded }\OtherTok{\textless{}{-}}\NormalTok{ QuantLib}\SpecialCharTok{::}\FunctionTok{Compounding\_Compounded\_get}\NormalTok{()}
\NormalTok{cmpd\_simple }\OtherTok{\textless{}{-}}\NormalTok{ QuantLib}\SpecialCharTok{::}\FunctionTok{Compounding\_Simple\_get}\NormalTok{()}

\NormalTok{freq\_annual }\OtherTok{\textless{}{-}}\NormalTok{ QuantLib}\SpecialCharTok{::}\FunctionTok{Frequency\_Annual\_get}\NormalTok{()}
\NormalTok{freq\_semiannual }\OtherTok{\textless{}{-}}\NormalTok{ QuantLib}\SpecialCharTok{::}\FunctionTok{Frequency\_Semiannual\_get}\NormalTok{()}

\NormalTok{trade\_date\_chr }\OtherTok{\textless{}{-}} \StringTok{"2022{-}08{-}19"}
\NormalTok{issue\_date\_chr }\OtherTok{\textless{}{-}} \StringTok{"2022{-}07{-}28"}
\NormalTok{maturity\_date\_chr }\OtherTok{\textless{}{-}} \StringTok{"2025{-}07{-}28"}

\NormalTok{trade\_date }\OtherTok{\textless{}{-}}\NormalTok{ GaussQuant}\SpecialCharTok{::}\FunctionTok{date\_GQL}\NormalTok{(trade\_date\_chr)}
\NormalTok{issue\_date\_ql }\OtherTok{\textless{}{-}}\NormalTok{ GaussQuant}\SpecialCharTok{::}\FunctionTok{date\_GQL}\NormalTok{(issue\_date\_chr)}
\NormalTok{maturity\_date\_ql }\OtherTok{\textless{}{-}}\NormalTok{ GaussQuant}\SpecialCharTok{::}\FunctionTok{date\_GQL}\NormalTok{(maturity\_date\_chr)}

\NormalTok{settlement\_days }\OtherTok{\textless{}{-}} \DecValTok{2}\NormalTok{L}

\NormalTok{GaussQuant}\SpecialCharTok{::}\FunctionTok{set\_eval\_date\_GQL}\NormalTok{(trade\_date)}

\NormalTok{settlement\_date }\OtherTok{\textless{}{-}}\NormalTok{ GaussQuant}\SpecialCharTok{::}\FunctionTok{advance\_days\_GQL}\NormalTok{(}
  \AttributeTok{calendar\_obj =}\NormalTok{ jp\_calendar,}
  \AttributeTok{date\_obj =}\NormalTok{ trade\_date,}
  \AttributeTok{n\_days =}\NormalTok{ settlement\_days}
\NormalTok{)}

\NormalTok{settlement\_date\_chr }\OtherTok{\textless{}{-}} \FunctionTok{as.character}\NormalTok{(settlement\_date)}
\NormalTok{settlement\_date\_ql }\OtherTok{\textless{}{-}}\NormalTok{ GaussQuant}\SpecialCharTok{::}\FunctionTok{date\_GQL}\NormalTok{(settlement\_date)}

\NormalTok{setup\_tbl }\OtherTok{\textless{}{-}}\NormalTok{ tibble}\SpecialCharTok{::}\FunctionTok{tibble}\NormalTok{(}
  \AttributeTok{item =} \FunctionTok{c}\NormalTok{(}
    \StringTok{"trade\_date"}\NormalTok{,}
    \StringTok{"settlement\_date"}\NormalTok{,}
    \StringTok{"issue\_date"}\NormalTok{,}
    \StringTok{"maturity\_date"}\NormalTok{,}
    \StringTok{"settlement\_days"}\NormalTok{,}
    \StringTok{"day\_counter"}\NormalTok{,}
    \StringTok{"compounding"}\NormalTok{,}
    \StringTok{"frequency"}
\NormalTok{  ),}
  \AttributeTok{value =} \FunctionTok{c}\NormalTok{(}
\NormalTok{    trade\_date\_chr,}
\NormalTok{    settlement\_date\_chr,}
\NormalTok{    issue\_date\_chr,}
\NormalTok{    maturity\_date\_chr,}
    \FunctionTok{as.character}\NormalTok{(settlement\_days),}
    \StringTok{"Thirty360(BondBasis)"}\NormalTok{,}
    \StringTok{"Compounded"}\NormalTok{,}
    \StringTok{"Semiannual"}
\NormalTok{  )}
\NormalTok{)}

\NormalTok{GaussQuant}\SpecialCharTok{::}\FunctionTok{show\_tbl\_GQH}\NormalTok{(setup\_tbl, }\StringTok{"Bond setup dates and conventions"}\NormalTok{)}
\CommentTok{\#\textgreater{} \# A tibble: 8 x 2}
\CommentTok{\#\textgreater{}   item            value               }
\CommentTok{\#\textgreater{}   \textless{}chr\textgreater{}           \textless{}chr\textgreater{}               }
\CommentTok{\#\textgreater{} 1 trade\_date      2022{-}08{-}19          }
\CommentTok{\#\textgreater{} 2 settlement\_date 2022{-}08{-}23          }
\CommentTok{\#\textgreater{} 3 issue\_date      2022{-}07{-}28          }
\CommentTok{\#\textgreater{} 4 maturity\_date   2025{-}07{-}28          }
\CommentTok{\#\textgreater{} 5 settlement\_days 2                   }
\CommentTok{\#\textgreater{} 6 day\_counter     Thirty360(BondBasis)}
\CommentTok{\#\textgreater{} 7 compounding     Compounded          }
\CommentTok{\#\textgreater{} 8 frequency       Semiannual}
\end{Highlighting}
\end{Shaded}

表示された表から、取引日、受渡日、発行日、満期日および評価に用いる市場慣行を確認する。受渡日は日本の休日カレンダーを用いて取引日から2営業日後に計算される。

\section{1.3
固定利付債計算}\label{ux56faux5b9aux5229ux4ed8ux50b5ux8a08ux7b97}

額面100、年率0.37\%、半年払い、満期2025年7月28日の固定利付債を作成する。利払スケジュールは満期日から発行日へ向かって6か月ごとに生成し、各利払日はUnadjustedとする。

\begin{Shaded}
\begin{Highlighting}[]
\NormalTok{face\_amount }\OtherTok{\textless{}{-}} \DecValTok{100}
\NormalTok{coupon\_rate }\OtherTok{\textless{}{-}} \FloatTok{0.00370}
\NormalTok{clean\_price\_input }\OtherTok{\textless{}{-}} \FloatTok{97.0}

\NormalTok{bond\_obj }\OtherTok{\textless{}{-}}\NormalTok{ GaussQuant}\SpecialCharTok{::}\FunctionTok{fixed\_rate\_bond\_GQL}\NormalTok{(}
  \AttributeTok{issue\_date =}\NormalTok{ issue\_date\_chr,}
  \AttributeTok{maturity\_date =}\NormalTok{ maturity\_date\_chr,}
  \AttributeTok{face\_amount =}\NormalTok{ face\_amount,}
  \AttributeTok{coupon\_rate =}\NormalTok{ coupon\_rate,}
  \AttributeTok{settlement\_days =}\NormalTok{ settlement\_days,}
  \AttributeTok{calendar =}\NormalTok{ jp\_calendar,}
  \AttributeTok{accrual\_day\_counter =}\NormalTok{ dc\_30,}
  \AttributeTok{payment\_convention =} \StringTok{"Unadjusted"}\NormalTok{,}
  \AttributeTok{schedule\_convention =} \StringTok{"Unadjusted"}\NormalTok{,}
  \AttributeTok{maturity\_convention =} \StringTok{"Unadjusted"}\NormalTok{,}
  \AttributeTok{date\_generation =} \StringTok{"Backward"}\NormalTok{,}
  \AttributeTok{end\_of\_month =} \ConstantTok{FALSE}\NormalTok{,}
  \AttributeTok{schedule\_frequency =}\NormalTok{ GaussQuant}\SpecialCharTok{::}\FunctionTok{period\_months\_GQL}\NormalTok{(}\DecValTok{6}\NormalTok{)}
\NormalTok{)}

\NormalTok{bond\_schedule }\OtherTok{\textless{}{-}}\NormalTok{ QuantLib}\SpecialCharTok{::}\FunctionTok{Schedule}\NormalTok{(}
\NormalTok{  issue\_date\_ql,}
\NormalTok{  maturity\_date\_ql,}
\NormalTok{  GaussQuant}\SpecialCharTok{::}\FunctionTok{period\_months\_GQL}\NormalTok{(}\DecValTok{6}\NormalTok{),}
\NormalTok{  jp\_calendar,}
  \StringTok{"Unadjusted"}\NormalTok{,}
  \StringTok{"Unadjusted"}\NormalTok{,}
  \StringTok{"Backward"}\NormalTok{,}
  \ConstantTok{FALSE}
\NormalTok{)}

\NormalTok{schedule\_tbl }\OtherTok{\textless{}{-}}\NormalTok{ GaussQuant}\SpecialCharTok{::}\FunctionTok{schedule\_dates\_GQL}\NormalTok{(bond\_schedule)}

\NormalTok{bond\_setup\_tbl }\OtherTok{\textless{}{-}}\NormalTok{ tibble}\SpecialCharTok{::}\FunctionTok{tibble}\NormalTok{(}
  \AttributeTok{metric =} \FunctionTok{c}\NormalTok{(}
    \StringTok{"bond\_class"}\NormalTok{,}
    \StringTok{"face\_amount"}\NormalTok{,}
    \StringTok{"coupon\_rate"}\NormalTok{,}
    \StringTok{"clean\_price\_input"}
\NormalTok{  ),}
  \AttributeTok{value =} \FunctionTok{list}\NormalTok{(}
    \FunctionTok{paste}\NormalTok{(}\FunctionTok{class}\NormalTok{(bond\_obj), }\AttributeTok{collapse =} \StringTok{", "}\NormalTok{),}
\NormalTok{    face\_amount,}
\NormalTok{    coupon\_rate,}
\NormalTok{    clean\_price\_input}
\NormalTok{  )}
\NormalTok{)}

\NormalTok{GaussQuant}\SpecialCharTok{::}\FunctionTok{show\_tbl\_GQH}\NormalTok{(GaussQuant}\SpecialCharTok{::}\FunctionTok{metric\_tbl\_display\_GQH}\NormalTok{(bond\_setup\_tbl), }\StringTok{"Fixed{-}rate bond object"}\NormalTok{)}
\CommentTok{\#\textgreater{} \# A tibble: 4 x 2}
\CommentTok{\#\textgreater{}   metric            value           }
\CommentTok{\#\textgreater{}   \textless{}chr\textgreater{}             \textless{}chr\textgreater{}           }
\CommentTok{\#\textgreater{} 1 bond\_class        \_p\_FixedRateBond}
\CommentTok{\#\textgreater{} 2 face\_amount       100.00000000    }
\CommentTok{\#\textgreater{} 3 coupon\_rate       0.00370000      }
\CommentTok{\#\textgreater{} 4 clean\_price\_input 97.00000000}
\NormalTok{GaussQuant}\SpecialCharTok{::}\FunctionTok{show\_tbl\_GQH}\NormalTok{(schedule\_tbl, }\StringTok{"Bond schedule"}\NormalTok{, }\AttributeTok{n =} \DecValTok{20}\NormalTok{)}
\CommentTok{\#\textgreater{} \# A tibble: 7 x 1}
\CommentTok{\#\textgreater{}   schedule\_date}
\CommentTok{\#\textgreater{}   \textless{}date\textgreater{}       }
\CommentTok{\#\textgreater{} 1 2022{-}07{-}28   }
\CommentTok{\#\textgreater{} 2 2023{-}01{-}28   }
\CommentTok{\#\textgreater{} 3 2023{-}07{-}28   }
\CommentTok{\#\textgreater{} 4 2024{-}01{-}28   }
\CommentTok{\#\textgreater{} 5 2024{-}07{-}28   }
\CommentTok{\#\textgreater{} 6 2025{-}01{-}28   }
\CommentTok{\#\textgreater{} 7 2025{-}07{-}28}
\end{Highlighting}
\end{Shaded}

債券オブジェクトの条件と、発行日から満期日までの利払日を確認する。

\subsection{1.3.1
クリーン価格から最終利回り}\label{ux30afux30eaux30fcux30f3ux4fa1ux683cux304bux3089ux6700ux7d42ux5229ux56deux308a}

市場で提示されるクリーン価格97から、将来キャッシュフローの現在価値がこの価格と一致する最終利回りを逆算する。

\begin{Shaded}
\begin{Highlighting}[]
\NormalTok{price\_to\_yield }\OtherTok{\textless{}{-}}\NormalTok{ GaussQuant}\SpecialCharTok{::}\FunctionTok{bond\_yield\_from\_clean\_safe\_GQL}\NormalTok{(}
  \AttributeTok{bond =}\NormalTok{ bond\_obj,}
  \AttributeTok{clean\_price =}\NormalTok{ clean\_price\_input,}
  \AttributeTok{day\_counter =}\NormalTok{ dc\_30,}
  \AttributeTok{compounding =}\NormalTok{ cmpd\_compounded,}
  \AttributeTok{frequency =}\NormalTok{ freq\_semiannual}
\NormalTok{)}

\NormalTok{yield\_tbl }\OtherTok{\textless{}{-}}\NormalTok{ tibble}\SpecialCharTok{::}\FunctionTok{tibble}\NormalTok{(}
  \AttributeTok{metric =} \FunctionTok{c}\NormalTok{(}
    \StringTok{"clean\_price\_input"}\NormalTok{,}
    \StringTok{"yield\_from\_clean\_price"}\NormalTok{,}
    \StringTok{"yield\_percent"}
\NormalTok{  ),}
  \AttributeTok{value =} \FunctionTok{list}\NormalTok{(}
\NormalTok{    clean\_price\_input,}
\NormalTok{    price\_to\_yield,}
    \DecValTok{100} \SpecialCharTok{*}\NormalTok{ price\_to\_yield}
\NormalTok{  )}
\NormalTok{)}

\NormalTok{GaussQuant}\SpecialCharTok{::}\FunctionTok{show\_tbl\_GQH}\NormalTok{(GaussQuant}\SpecialCharTok{::}\FunctionTok{metric\_tbl\_display\_GQH}\NormalTok{(yield\_tbl), }\StringTok{"Clean price {-}\textgreater{} yield"}\NormalTok{)}
\CommentTok{\#\textgreater{} \# A tibble: 3 x 2}
\CommentTok{\#\textgreater{}   metric                 value      }
\CommentTok{\#\textgreater{}   \textless{}chr\textgreater{}                  \textless{}chr\textgreater{}      }
\CommentTok{\#\textgreater{} 1 clean\_price\_input      97.00000000}
\CommentTok{\#\textgreater{} 2 yield\_from\_clean\_price 0.01418724 }
\CommentTok{\#\textgreater{} 3 yield\_percent          1.41872370}
\end{Highlighting}
\end{Shaded}

債券価格が額面を下回っているため、算出される最終利回りはクーポンレートを上回る。

\subsection{1.3.2
最終利回りから価格}\label{ux6700ux7d42ux5229ux56deux308aux304bux3089ux4fa1ux683c}

前節で求めた最終利回りを使って、クリーン価格、ダーティ価格および経過利子を計算する。両価格の関係は次式で表される。

\[
\text{Dirty Price}=\text{Clean Price}+\text{Accrued Interest}
\]

\begin{Shaded}
\begin{Highlighting}[]
\NormalTok{price\_tbl }\OtherTok{\textless{}{-}}\NormalTok{ GaussQuant}\SpecialCharTok{::}\FunctionTok{bond\_price\_measures\_GQL}\NormalTok{(}
  \AttributeTok{bond =}\NormalTok{ bond\_obj,}
  \AttributeTok{yield =}\NormalTok{ price\_to\_yield,}
  \AttributeTok{day\_counter =}\NormalTok{ dc\_30,}
  \AttributeTok{compounding =}\NormalTok{ cmpd\_compounded,}
  \AttributeTok{frequency =}\NormalTok{ freq\_semiannual,}
  \AttributeTok{settlement\_date =}\NormalTok{ settlement\_date\_chr}
\NormalTok{)}

\NormalTok{price\_tbl }\OtherTok{\textless{}{-}}\NormalTok{ GaussQuant}\SpecialCharTok{::}\FunctionTok{complete\_accrual\_metrics\_GQL}\NormalTok{(}
  \AttributeTok{price\_tbl =}\NormalTok{ price\_tbl,}
  \AttributeTok{bond\_schedule =}\NormalTok{ bond\_schedule,}
  \AttributeTok{settlement\_date =}\NormalTok{ settlement\_date,}
  \AttributeTok{day\_counter =}\NormalTok{ dc\_30}
\NormalTok{)}

\NormalTok{price\_tbl\_display }\OtherTok{\textless{}{-}}\NormalTok{ GaussQuant}\SpecialCharTok{::}\FunctionTok{metric\_tbl\_display\_GQH}\NormalTok{(price\_tbl)}

\NormalTok{GaussQuant}\SpecialCharTok{::}\FunctionTok{show\_tbl\_GQH}\NormalTok{(price\_tbl\_display, }\StringTok{"Yield {-}\textgreater{} price summary"}\NormalTok{)}
\CommentTok{\#\textgreater{} \# A tibble: 6 x 2}
\CommentTok{\#\textgreater{}   metric               value      }
\CommentTok{\#\textgreater{}   \textless{}chr\textgreater{}                \textless{}chr\textgreater{}      }
\CommentTok{\#\textgreater{} 1 settlement\_date      2022{-}08{-}23 }
\CommentTok{\#\textgreater{} 2 previous\_coupon\_date 2022{-}07{-}28 }
\CommentTok{\#\textgreater{} 3 accrued\_days         25.00000000}
\CommentTok{\#\textgreater{} 4 accrued\_amount       0.02569444 }
\CommentTok{\#\textgreater{} 5 clean\_price          97.00000012}
\CommentTok{\#\textgreater{} 6 dirty\_price          97.02569456}

\NormalTok{dirty\_price\_from\_yield }\OtherTok{\textless{}{-}}\NormalTok{ GaussQuant}\SpecialCharTok{::}\FunctionTok{get\_metric\_num\_GQH}\NormalTok{(price\_tbl, }\StringTok{"dirty\_price"}\NormalTok{)}
\NormalTok{clean\_price\_from\_yield }\OtherTok{\textless{}{-}}\NormalTok{ GaussQuant}\SpecialCharTok{::}\FunctionTok{get\_metric\_num\_GQH}\NormalTok{(price\_tbl, }\StringTok{"clean\_price"}\NormalTok{)}
\NormalTok{accrued\_amount }\OtherTok{\textless{}{-}}\NormalTok{ GaussQuant}\SpecialCharTok{::}\FunctionTok{get\_metric\_num\_GQH}\NormalTok{(price\_tbl, }\StringTok{"accrued\_amount"}\NormalTok{)}

\FunctionTok{cat}\NormalTok{(}
  \StringTok{"clean price:"}\NormalTok{,}
  \FunctionTok{sprintf}\NormalTok{(}\StringTok{"\%.8f"}\NormalTok{, clean\_price\_from\_yield),}
  \StringTok{"}\SpecialCharTok{\textbackslash{}n}\StringTok{"}\NormalTok{,}
  \StringTok{"dirty price:"}\NormalTok{,}
  \FunctionTok{sprintf}\NormalTok{(}\StringTok{"\%.8f"}\NormalTok{, dirty\_price\_from\_yield),}
  \StringTok{"}\SpecialCharTok{\textbackslash{}n}\StringTok{"}\NormalTok{,}
  \StringTok{"accrued:"}\NormalTok{,}
  \FunctionTok{sprintf}\NormalTok{(}\StringTok{"\%.8f"}\NormalTok{, accrued\_amount),}
  \StringTok{"}\SpecialCharTok{\textbackslash{}n}\StringTok{"}
\NormalTok{)}
\CommentTok{\#\textgreater{} clean price: 97.00000012 }
\CommentTok{\#\textgreater{}  dirty price: 97.02569456 }
\CommentTok{\#\textgreater{}  accrued: 0.02569444}
\end{Highlighting}
\end{Shaded}

再計算されたクリーン価格が入力価格97と一致することを確認する。ダーティ価格との差額は、前回利払日から受渡日までに発生した経過利子である。

\subsection{1.3.3 金利リスク}\label{ux91d1ux5229ux30eaux30b9ux30af}

修正デュレーションは、利回りの小さな変化に対する債券価格の一次感応度を表す。コンベクシティは価格と利回りの非線形な関係を補正する二次感応度である。

\begin{Shaded}
\begin{Highlighting}[]
\NormalTok{risk\_tbl }\OtherTok{\textless{}{-}}\NormalTok{ GaussQuant}\SpecialCharTok{::}\FunctionTok{bond\_risk\_measures\_GQL}\NormalTok{(}
  \AttributeTok{bond =}\NormalTok{ bond\_obj,}
  \AttributeTok{yield =}\NormalTok{ price\_to\_yield,}
  \AttributeTok{day\_counter =}\NormalTok{ dc\_30,}
  \AttributeTok{compounding =}\NormalTok{ cmpd\_compounded,}
  \AttributeTok{frequency =}\NormalTok{ freq\_semiannual}
\NormalTok{)}

\NormalTok{risk\_tbl\_display }\OtherTok{\textless{}{-}}\NormalTok{ GaussQuant}\SpecialCharTok{::}\FunctionTok{metric\_tbl\_display\_GQH}\NormalTok{(risk\_tbl)}

\NormalTok{GaussQuant}\SpecialCharTok{::}\FunctionTok{show\_tbl\_GQH}\NormalTok{(risk\_tbl\_display, }\StringTok{"Bond risk measures"}\NormalTok{)}
\CommentTok{\#\textgreater{} \# A tibble: 4 x 2}
\CommentTok{\#\textgreater{}   metric            value     }
\CommentTok{\#\textgreater{}   \textless{}chr\textgreater{}             \textless{}chr\textgreater{}     }
\CommentTok{\#\textgreater{} 1 modified\_duration 2.89593215}
\CommentTok{\#\textgreater{} 2 bpv               0.02809798}
\CommentTok{\#\textgreater{} 3 dv01              0.02809798}
\CommentTok{\#\textgreater{} 4 convexity         9.84952310}

\NormalTok{modified\_duration }\OtherTok{\textless{}{-}}\NormalTok{ GaussQuant}\SpecialCharTok{::}\FunctionTok{get\_metric\_num\_GQH}\NormalTok{(risk\_tbl, }\StringTok{"modified\_duration"}\NormalTok{)}
\NormalTok{convexity\_value }\OtherTok{\textless{}{-}}\NormalTok{ GaussQuant}\SpecialCharTok{::}\FunctionTok{get\_metric\_num\_GQH}\NormalTok{(risk\_tbl, }\StringTok{"convexity"}\NormalTok{)}

\FunctionTok{cat}\NormalTok{(}
  \StringTok{"modified duration:"}\NormalTok{,}
  \FunctionTok{sprintf}\NormalTok{(}\StringTok{"\%.8f"}\NormalTok{, modified\_duration),}
  \StringTok{"}\SpecialCharTok{\textbackslash{}n}\StringTok{"}\NormalTok{,}
  \StringTok{"convexity:"}\NormalTok{,}
  \FunctionTok{sprintf}\NormalTok{(}\StringTok{"\%.8f"}\NormalTok{, convexity\_value),}
  \StringTok{"}\SpecialCharTok{\textbackslash{}n}\StringTok{"}
\NormalTok{)}
\CommentTok{\#\textgreater{} modified duration: 2.89593215 }
\CommentTok{\#\textgreater{}  convexity: 9.84952310}
\end{Highlighting}
\end{Shaded}

修正デュレーションが大きいほど、同じ利回り変化に対する価格変化も大きくなる。コンベクシティは、利回り変化が大きくなるほど重要になる。

\subsection{1.3.4
デュレーションとコンベクシティによる価格変化の近似}\label{ux30c7ux30e5ux30ecux30fcux30b7ux30e7ux30f3ux3068ux30b3ux30f3ux30d9ux30afux30b7ux30c6ux30a3ux306bux3088ux308bux4fa1ux683cux5909ux5316ux306eux8fd1ux4f3c}

利回り変化を \(\Delta y\)
とすると、債券価格の相対変化は修正デュレーションとコンベクシティを使って近似できる。

\[
\frac{\Delta P}{P}
\approx
-D_{\mathrm{mod}}\Delta y
+\frac{1}{2}C(\Delta y)^2
\]

\begin{Shaded}
\begin{Highlighting}[]
\NormalTok{handcalc\_tbl }\OtherTok{\textless{}{-}}\NormalTok{ GaussQuant}\SpecialCharTok{::}\FunctionTok{bond\_risk\_handcalc\_GQL}\NormalTok{(}
  \AttributeTok{bond =}\NormalTok{ bond\_obj,}
  \AttributeTok{yield =}\NormalTok{ price\_to\_yield,}
  \AttributeTok{dirty\_price =}\NormalTok{ dirty\_price\_from\_yield,}
  \AttributeTok{modified\_duration =}\NormalTok{ modified\_duration,}
  \AttributeTok{convexity =}\NormalTok{ convexity\_value,}
  \AttributeTok{day\_counter =}\NormalTok{ dc\_30,}
  \AttributeTok{compounding =}\NormalTok{ cmpd\_compounded,}
  \AttributeTok{frequency =}\NormalTok{ freq\_semiannual}
\NormalTok{)}

\NormalTok{handcalc\_tbl\_display }\OtherTok{\textless{}{-}}\NormalTok{ GaussQuant}\SpecialCharTok{::}\FunctionTok{metric\_tbl\_display\_GQH}\NormalTok{(handcalc\_tbl)}

\NormalTok{GaussQuant}\SpecialCharTok{::}\FunctionTok{show\_tbl\_GQH}\NormalTok{(handcalc\_tbl\_display, }\StringTok{"Risk hand calculations"}\NormalTok{)}
\CommentTok{\#\textgreater{} \# A tibble: 5 x 2}
\CommentTok{\#\textgreater{}   metric                   value      }
\CommentTok{\#\textgreater{}   \textless{}chr\textgreater{}                    \textless{}chr\textgreater{}      }
\CommentTok{\#\textgreater{} 1 bpv\_hand                 {-}0.02809798}
\CommentTok{\#\textgreater{} 2 modified\_duration\_hand   0.02895932 }
\CommentTok{\#\textgreater{} 3 convexity\_hand           0.00098495 }
\CommentTok{\#\textgreater{} 4 price\_up\_100bp           94.26306473}
\CommentTok{\#\textgreater{} 5 price\_approx\_delta\_gamma 94.26367912}
\end{Highlighting}
\end{Shaded}

一次近似と二次近似を比較し、コンベクシティを加えることで実際の価格変化に近づくことを確認する。

\section{1.4
キャッシュフローによる検証}\label{ux30adux30e3ux30c3ux30b7ux30e5ux30d5ux30edux30fcux306bux3088ux308bux691cux8a3c}

各利払日のクーポンと満期時の元本を最終利回りで割り引き、ダーティ価格を手計算する。そこから経過利子を差し引くことでクリーン価格を求める。

\begin{Shaded}
\begin{Highlighting}[]
\NormalTok{schedule\_dates\_chr }\OtherTok{\textless{}{-}}\NormalTok{ schedule\_tbl[[}\DecValTok{1}\NormalTok{]] }\SpecialCharTok{|\textgreater{}}
  \FunctionTok{as.character}\NormalTok{()}

\NormalTok{schedule\_dates\_chr }\OtherTok{\textless{}{-}}\NormalTok{ schedule\_dates\_chr[}\SpecialCharTok{!}\FunctionTok{is.na}\NormalTok{(schedule\_dates\_chr)]}

\NormalTok{future\_dates\_chr }\OtherTok{\textless{}{-}}\NormalTok{ schedule\_dates\_chr[schedule\_dates\_chr }\SpecialCharTok{\textgreater{}}\NormalTok{ settlement\_date\_chr]}

\NormalTok{coupon\_per\_period }\OtherTok{\textless{}{-}}\NormalTok{ face\_amount }\SpecialCharTok{*}\NormalTok{ coupon\_rate }\SpecialCharTok{/} \DecValTok{2}

\NormalTok{manual\_cf\_tbl }\OtherTok{\textless{}{-}}\NormalTok{ tibble}\SpecialCharTok{::}\FunctionTok{tibble}\NormalTok{(}
  \AttributeTok{payment\_date =}\NormalTok{ future\_dates\_chr,}
  \AttributeTok{period\_no =} \FunctionTok{seq\_along}\NormalTok{(future\_dates\_chr),}
  \AttributeTok{cashflow =}\NormalTok{ coupon\_per\_period}
\NormalTok{) }\SpecialCharTok{|\textgreater{}}
\NormalTok{  dplyr}\SpecialCharTok{::}\FunctionTok{mutate}\NormalTok{(}
    \AttributeTok{cashflow =} \FunctionTok{ifelse}\NormalTok{(}
\NormalTok{      .data}\SpecialCharTok{$}\NormalTok{payment\_date }\SpecialCharTok{==}\NormalTok{ maturity\_date\_chr,}
\NormalTok{      .data}\SpecialCharTok{$}\NormalTok{cashflow }\SpecialCharTok{+}\NormalTok{ face\_amount,}
\NormalTok{      .data}\SpecialCharTok{$}\NormalTok{cashflow}
\NormalTok{    ),}
    \AttributeTok{ql\_date =}\NormalTok{ purrr}\SpecialCharTok{::}\FunctionTok{map}\NormalTok{(payment\_date, GaussQuant}\SpecialCharTok{::}\NormalTok{date\_GQL),}
    \AttributeTok{time =}\NormalTok{ purrr}\SpecialCharTok{::}\FunctionTok{map\_dbl}\NormalTok{(}
\NormalTok{      ql\_date,}
      \SpecialCharTok{\textasciitilde{}}\NormalTok{ dc\_30}\SpecialCharTok{$}\FunctionTok{yearFraction}\NormalTok{(settlement\_date\_ql, .x)}
\NormalTok{    ),}
    \AttributeTok{discount\_factor =} \DecValTok{1} \SpecialCharTok{/}\NormalTok{ (}\DecValTok{1} \SpecialCharTok{+}\NormalTok{ price\_to\_yield }\SpecialCharTok{/} \DecValTok{2}\NormalTok{)}\SpecialCharTok{\^{}}\NormalTok{(}\DecValTok{2} \SpecialCharTok{*}\NormalTok{ time),}
    \AttributeTok{pv =}\NormalTok{ cashflow }\SpecialCharTok{*}\NormalTok{ discount\_factor}
\NormalTok{  ) }\SpecialCharTok{|\textgreater{}}
\NormalTok{  dplyr}\SpecialCharTok{::}\FunctionTok{select}\NormalTok{(}\SpecialCharTok{{-}}\NormalTok{ql\_date)}

\NormalTok{manual\_dirty\_price }\OtherTok{\textless{}{-}} \FunctionTok{sum}\NormalTok{(manual\_cf\_tbl}\SpecialCharTok{$}\NormalTok{pv, }\AttributeTok{na.rm =} \ConstantTok{TRUE}\NormalTok{)}
\NormalTok{manual\_clean\_price }\OtherTok{\textless{}{-}}\NormalTok{ manual\_dirty\_price }\SpecialCharTok{{-}}\NormalTok{ accrued\_amount}

\NormalTok{manual\_summary\_tbl }\OtherTok{\textless{}{-}}\NormalTok{ tibble}\SpecialCharTok{::}\FunctionTok{tibble}\NormalTok{(}
  \AttributeTok{metric =} \FunctionTok{c}\NormalTok{(}
    \StringTok{"manual\_dirty\_price"}\NormalTok{,}
    \StringTok{"manual\_clean\_price"}\NormalTok{,}
    \StringTok{"wrapper\_dirty\_price"}\NormalTok{,}
    \StringTok{"wrapper\_clean\_price"}\NormalTok{,}
    \StringTok{"dirty\_price\_difference"}\NormalTok{,}
    \StringTok{"clean\_price\_difference"}
\NormalTok{  ),}
  \AttributeTok{value =} \FunctionTok{list}\NormalTok{(}
\NormalTok{    manual\_dirty\_price,}
\NormalTok{    manual\_clean\_price,}
\NormalTok{    dirty\_price\_from\_yield,}
\NormalTok{    clean\_price\_from\_yield,}
\NormalTok{    manual\_dirty\_price }\SpecialCharTok{{-}}\NormalTok{ dirty\_price\_from\_yield,}
\NormalTok{    manual\_clean\_price }\SpecialCharTok{{-}}\NormalTok{ clean\_price\_from\_yield}
\NormalTok{  )}
\NormalTok{)}

\NormalTok{GaussQuant}\SpecialCharTok{::}\FunctionTok{show\_tbl\_GQH}\NormalTok{(manual\_cf\_tbl, }\StringTok{"Manual cashflow PV table"}\NormalTok{)}
\CommentTok{\#\textgreater{} \# A tibble: 6 x 6}
\CommentTok{\#\textgreater{}   payment\_date period\_no cashflow  time discount\_factor     pv}
\CommentTok{\#\textgreater{}   \textless{}chr\textgreater{}            \textless{}int\textgreater{}    \textless{}dbl\textgreater{} \textless{}dbl\textgreater{}           \textless{}dbl\textgreater{}  \textless{}dbl\textgreater{}}
\CommentTok{\#\textgreater{} 1 2023{-}01{-}28           1    0.185 0.431           0.994  0.184}
\CommentTok{\#\textgreater{} 2 2023{-}07{-}28           2    0.185 0.931           0.987  0.183}
\CommentTok{\#\textgreater{} 3 2024{-}01{-}28           3    0.185 1.43            0.980  0.181}
\CommentTok{\#\textgreater{} 4 2024{-}07{-}28           4    0.185 1.93            0.973  0.180}
\CommentTok{\#\textgreater{} 5 2025{-}01{-}28           5    0.185 2.43            0.966  0.179}
\CommentTok{\#\textgreater{} 6 2025{-}07{-}28           6  100.    2.93            0.959 96.1}
\NormalTok{GaussQuant}\SpecialCharTok{::}\FunctionTok{show\_tbl\_GQH}\NormalTok{(GaussQuant}\SpecialCharTok{::}\FunctionTok{metric\_tbl\_display\_GQH}\NormalTok{(manual\_summary\_tbl), }\StringTok{"Manual PV check"}\NormalTok{)}
\CommentTok{\#\textgreater{} \# A tibble: 6 x 2}
\CommentTok{\#\textgreater{}   metric                 value      }
\CommentTok{\#\textgreater{}   \textless{}chr\textgreater{}                  \textless{}chr\textgreater{}      }
\CommentTok{\#\textgreater{} 1 manual\_dirty\_price     97.02569456}
\CommentTok{\#\textgreater{} 2 manual\_clean\_price     97.00000012}
\CommentTok{\#\textgreater{} 3 wrapper\_dirty\_price    97.02569456}
\CommentTok{\#\textgreater{} 4 wrapper\_clean\_price    97.00000012}
\CommentTok{\#\textgreater{} 5 dirty\_price\_difference 0.00000000 }
\CommentTok{\#\textgreater{} 6 clean\_price\_difference 0.00000000}
\end{Highlighting}
\end{Shaded}

手計算と関数による価格との差が十分に小さいことを確認する。これにより、キャッシュフロー、割引期間および経過利子の関係を検算できる。

\subsection{1.4.1
フラットな利回りによる価格評価}\label{ux30d5ux30e9ux30c3ux30c8ux306aux5229ux56deux308aux306bux3088ux308bux4fa1ux683cux8a55ux4fa1}

次に、最終利回りを一定とするフラットな割引曲線を構築し、割引エンジンから債券価格を計算する。

\begin{Shaded}
\begin{Highlighting}[]
\NormalTok{flat\_yield\_handle }\OtherTok{\textless{}{-}}\NormalTok{ GaussQuant}\SpecialCharTok{::}\FunctionTok{bond\_flat\_forward\_handle\_GQL}\NormalTok{(}
  \AttributeTok{settlement\_date =}\NormalTok{ settlement\_date\_chr,}
  \AttributeTok{rate =}\NormalTok{ price\_to\_yield,}
  \AttributeTok{day\_counter =}\NormalTok{ dc\_30,}
  \AttributeTok{compounding =}\NormalTok{ cmpd\_compounded,}
  \AttributeTok{frequency =}\NormalTok{ freq\_semiannual}
\NormalTok{)}

\NormalTok{flat\_engine }\OtherTok{\textless{}{-}}\NormalTok{ QuantLib}\SpecialCharTok{::}\FunctionTok{DiscountingBondEngine}\NormalTok{(flat\_yield\_handle)}

\NormalTok{QuantLib}\SpecialCharTok{::}\FunctionTok{Instrument\_setPricingEngine}\NormalTok{(}
\NormalTok{  bond\_obj,}
\NormalTok{  flat\_engine}
\NormalTok{)}
\CommentTok{\#\textgreater{} NULL}

\NormalTok{flat\_dirty\_price }\OtherTok{\textless{}{-}}\NormalTok{ QuantLib}\SpecialCharTok{::}\FunctionTok{Instrument\_NPV}\NormalTok{(bond\_obj)}

\NormalTok{flat\_yield\_tbl }\OtherTok{\textless{}{-}}\NormalTok{ tibble}\SpecialCharTok{::}\FunctionTok{tibble}\NormalTok{(}
  \AttributeTok{metric =} \FunctionTok{c}\NormalTok{(}
    \StringTok{"flat\_yield"}\NormalTok{,}
    \StringTok{"dirty\_price\_from\_flat\_engine"}\NormalTok{,}
    \StringTok{"dirty\_price\_from\_yield"}\NormalTok{,}
    \StringTok{"difference"}
\NormalTok{  ),}
  \AttributeTok{value =} \FunctionTok{list}\NormalTok{(}
\NormalTok{    price\_to\_yield,}
\NormalTok{    flat\_dirty\_price,}
\NormalTok{    dirty\_price\_from\_yield,}
\NormalTok{    flat\_dirty\_price }\SpecialCharTok{{-}}\NormalTok{ dirty\_price\_from\_yield}
\NormalTok{  )}
\NormalTok{)}

\NormalTok{GaussQuant}\SpecialCharTok{::}\FunctionTok{show\_tbl\_GQH}\NormalTok{(GaussQuant}\SpecialCharTok{::}\FunctionTok{metric\_tbl\_display\_GQH}\NormalTok{(flat\_yield\_tbl), }\StringTok{"Flat yield pricing"}\NormalTok{)}
\CommentTok{\#\textgreater{} \# A tibble: 4 x 2}
\CommentTok{\#\textgreater{}   metric                       value      }
\CommentTok{\#\textgreater{}   \textless{}chr\textgreater{}                        \textless{}chr\textgreater{}      }
\CommentTok{\#\textgreater{} 1 flat\_yield                   0.01418724 }
\CommentTok{\#\textgreater{} 2 dirty\_price\_from\_flat\_engine 97.02569456}
\CommentTok{\#\textgreater{} 3 dirty\_price\_from\_yield       97.02569456}
\CommentTok{\#\textgreater{} 4 difference                   {-}0.00000000}
\end{Highlighting}
\end{Shaded}

フラットな利回りを使った割引価格と、最終利回りから直接求めたダーティ価格を比較する。
\#\#\# 1.4.2 Z-spreadによる価格調整

Z-spreadは、基準となる割引カーブの各年限に一律に加えるスプレッドである。ここでは基準カーブに50ベーシスポイントを上乗せし、債券価格への影響を確認する。

\begin{Shaded}
\begin{Highlighting}[]
\NormalTok{z\_spread }\OtherTok{\textless{}{-}} \DecValTok{50} \SpecialCharTok{/} \DecValTok{10000}

\NormalTok{zspread\_result }\OtherTok{\textless{}{-}}\NormalTok{ GaussQuant}\SpecialCharTok{::}\FunctionTok{bond\_npv\_with\_zspread\_GQL}\NormalTok{(}
  \AttributeTok{bond =}\NormalTok{ bond\_obj,}
  \AttributeTok{base\_curve\_handle =}\NormalTok{ flat\_yield\_handle,}
  \AttributeTok{z\_spread =}\NormalTok{ z\_spread,}
  \AttributeTok{compounding =}\NormalTok{ cmpd\_compounded,}
  \AttributeTok{frequency =}\NormalTok{ freq\_semiannual,}
  \AttributeTok{day\_counter =}\NormalTok{ dc\_30}
\NormalTok{)}

\NormalTok{zspread\_dirty\_price }\OtherTok{\textless{}{-}}\NormalTok{ zspread\_result}\SpecialCharTok{$}\NormalTok{npv}

\NormalTok{zspread\_tbl }\OtherTok{\textless{}{-}}\NormalTok{ tibble}\SpecialCharTok{::}\FunctionTok{tibble}\NormalTok{(}
  \AttributeTok{metric =} \FunctionTok{c}\NormalTok{(}
    \StringTok{"base\_yield"}\NormalTok{,}
    \StringTok{"z\_spread\_bp"}\NormalTok{,}
    \StringTok{"dirty\_price\_before\_spread"}\NormalTok{,}
    \StringTok{"dirty\_price\_after\_spread"}\NormalTok{,}
    \StringTok{"price\_difference"}
\NormalTok{  ),}
  \AttributeTok{value =} \FunctionTok{list}\NormalTok{(}
\NormalTok{    price\_to\_yield,}
    \DecValTok{10000} \SpecialCharTok{*}\NormalTok{ z\_spread,}
\NormalTok{    flat\_dirty\_price,}
\NormalTok{    zspread\_dirty\_price,}
\NormalTok{    zspread\_dirty\_price }\SpecialCharTok{{-}}\NormalTok{ flat\_dirty\_price}
\NormalTok{  )}
\NormalTok{)}

\NormalTok{GaussQuant}\SpecialCharTok{::}\FunctionTok{show\_tbl\_GQH}\NormalTok{(}
\NormalTok{  GaussQuant}\SpecialCharTok{::}\FunctionTok{metric\_tbl\_display\_GQH}\NormalTok{(zspread\_tbl),}
  \StringTok{"Z{-}spread pricing"}
\NormalTok{)}
\CommentTok{\#\textgreater{} \# A tibble: 5 x 2}
\CommentTok{\#\textgreater{}   metric                    value      }
\CommentTok{\#\textgreater{}   \textless{}chr\textgreater{}                     \textless{}chr\textgreater{}      }
\CommentTok{\#\textgreater{} 1 base\_yield                0.01418724 }
\CommentTok{\#\textgreater{} 2 z\_spread\_bp               50.00000000}
\CommentTok{\#\textgreater{} 3 dirty\_price\_before\_spread 97.02569456}
\CommentTok{\#\textgreater{} 4 dirty\_price\_after\_spread  95.63266391}
\CommentTok{\#\textgreater{} 5 price\_difference          {-}1.39303065}
\end{Highlighting}
\end{Shaded}

正のZ-spreadを加えると割引率が上昇するため、債券価格は低下する。Z-spreadには発行体の信用リスクだけでなく、流動性、市場需給および基準カーブとの相違も含まれ得るため、純粋なデフォルト確率を直接表す指標ではない。ここでは全期間への一律なスプレッドだけを扱い、年限別の非平行シフトは金利カーブの章で扱う。
\#\#\# 1.4.3 債券キャッシュフロー

債券に含まれる将来のクーポンと元本償還を一覧で確認する。

\begin{Shaded}
\begin{Highlighting}[]
\NormalTok{bond\_cf\_tbl }\OtherTok{\textless{}{-}}\NormalTok{ GaussQuant}\SpecialCharTok{::}\FunctionTok{bond\_cashflow\_table\_GQL}\NormalTok{(}
  \AttributeTok{bond =}\NormalTok{ bond\_obj}
\NormalTok{)}

\NormalTok{GaussQuant}\SpecialCharTok{::}\FunctionTok{show\_tbl\_GQH}\NormalTok{(bond\_cf\_tbl, }\StringTok{"Bond cashflow table"}\NormalTok{, }\AttributeTok{n =} \DecValTok{30}\NormalTok{)}
\CommentTok{\#\textgreater{} \# A tibble: 7 x 2}
\CommentTok{\#\textgreater{}   date        amount}
\CommentTok{\#\textgreater{}   \textless{}chr\textgreater{}        \textless{}dbl\textgreater{}}
\CommentTok{\#\textgreater{} 1 2023{-}01{-}28   0.185}
\CommentTok{\#\textgreater{} 2 2023{-}07{-}28   0.185}
\CommentTok{\#\textgreater{} 3 2024{-}01{-}28   0.185}
\CommentTok{\#\textgreater{} 4 2024{-}07{-}28   0.185}
\CommentTok{\#\textgreater{} 5 2025{-}01{-}28   0.185}
\CommentTok{\#\textgreater{} 6 2025{-}07{-}28   0.185}
\CommentTok{\#\textgreater{} 7 2025{-}07{-}28 100}
\end{Highlighting}
\end{Shaded}

各行は支払日と支払額を表し、満期日のキャッシュフローには最終クーポンと元本償還が含まれる。

\section{1.5
計算結果のまとめ}\label{ux8a08ux7b97ux7d50ux679cux306eux307eux3068ux3081}

これまでに設定した日付、価格、利回りおよび金利リスク指標を一つの表にまとめる。

\begin{Shaded}
\begin{Highlighting}[]
\NormalTok{summary\_tbl }\OtherTok{\textless{}{-}}\NormalTok{ tibble}\SpecialCharTok{::}\FunctionTok{tibble}\NormalTok{(}
  \AttributeTok{metric =} \FunctionTok{c}\NormalTok{(}
    \StringTok{"trade\_date"}\NormalTok{,}
    \StringTok{"settlement\_date"}\NormalTok{,}
    \StringTok{"issue\_date"}\NormalTok{,}
    \StringTok{"maturity\_date"}\NormalTok{,}
    \StringTok{"clean\_price\_input"}\NormalTok{,}
    \StringTok{"yield\_from\_clean\_price"}\NormalTok{,}
    \StringTok{"clean\_price\_from\_yield"}\NormalTok{,}
    \StringTok{"dirty\_price\_from\_yield"}\NormalTok{,}
    \StringTok{"accrued\_amount"}\NormalTok{,}
    \StringTok{"modified\_duration"}\NormalTok{,}
    \StringTok{"convexity"}
\NormalTok{  ),}
  \AttributeTok{value =} \FunctionTok{c}\NormalTok{(}
\NormalTok{    GaussQuant}\SpecialCharTok{::}\FunctionTok{iso\_GQL}\NormalTok{(trade\_date),}
\NormalTok{    settlement\_date\_chr,}
\NormalTok{    issue\_date\_chr,}
\NormalTok{    maturity\_date\_chr,}
    \FunctionTok{sprintf}\NormalTok{(}\StringTok{"\%.8f"}\NormalTok{, clean\_price\_input),}
    \FunctionTok{sprintf}\NormalTok{(}\StringTok{"\%.8f"}\NormalTok{, price\_to\_yield),}
    \FunctionTok{sprintf}\NormalTok{(}\StringTok{"\%.8f"}\NormalTok{, clean\_price\_from\_yield),}
    \FunctionTok{sprintf}\NormalTok{(}\StringTok{"\%.8f"}\NormalTok{, dirty\_price\_from\_yield),}
    \FunctionTok{sprintf}\NormalTok{(}\StringTok{"\%.8f"}\NormalTok{, accrued\_amount),}
    \FunctionTok{sprintf}\NormalTok{(}\StringTok{"\%.8f"}\NormalTok{, modified\_duration),}
    \FunctionTok{sprintf}\NormalTok{(}\StringTok{"\%.8f"}\NormalTok{, convexity\_value)}
\NormalTok{  )}
\NormalTok{)}

\NormalTok{GaussQuant}\SpecialCharTok{::}\FunctionTok{show\_tbl\_GQH}\NormalTok{(summary\_tbl, }\StringTok{"Bond analysis summary"}\NormalTok{, }\AttributeTok{n =} \DecValTok{20}\NormalTok{)}
\CommentTok{\#\textgreater{} \# A tibble: 11 x 2}
\CommentTok{\#\textgreater{}    metric                 value      }
\CommentTok{\#\textgreater{}    \textless{}chr\textgreater{}                  \textless{}chr\textgreater{}      }
\CommentTok{\#\textgreater{}  1 trade\_date             2022{-}08{-}19 }
\CommentTok{\#\textgreater{}  2 settlement\_date        2022{-}08{-}23 }
\CommentTok{\#\textgreater{}  3 issue\_date             2022{-}07{-}28 }
\CommentTok{\#\textgreater{}  4 maturity\_date          2025{-}07{-}28 }
\CommentTok{\#\textgreater{}  5 clean\_price\_input      97.00000000}
\CommentTok{\#\textgreater{}  6 yield\_from\_clean\_price 0.01418724 }
\CommentTok{\#\textgreater{}  7 clean\_price\_from\_yield 97.00000012}
\CommentTok{\#\textgreater{}  8 dirty\_price\_from\_yield 97.02569456}
\CommentTok{\#\textgreater{}  9 accrued\_amount         0.02569444 }
\CommentTok{\#\textgreater{} 10 modified\_duration      2.89593215 }
\CommentTok{\#\textgreater{} 11 convexity              9.84952310}

\FunctionTok{cat}\NormalTok{(}\StringTok{"}\SpecialCharTok{\textbackslash{}n}\StringTok{chapter01 bonds stable rewrite completed successfully.}\SpecialCharTok{\textbackslash{}n}\StringTok{"}\NormalTok{)}
\CommentTok{\#\textgreater{} }
\CommentTok{\#\textgreater{} chapter01 bonds stable rewrite completed successfully.}
\end{Highlighting}
\end{Shaded}

\chapter{第II-2章
金利カーブ構築と債券先物CTD、BPV}\label{ux7b2cii-2ux7ae0-ux91d1ux5229ux30abux30fcux30d6ux69cbux7bc9ux3068ux50b5ux5238ux5148ux7269ctdbpv}

債券、金利スワップおよび金利オプションの評価では、将来キャッシュフローを現在価値へ換算するための金利期間構造が必要となる。金利期間構造は、満期ごとの金利を並べた単なる曲線ではない。市場で観測される預金金利、債券価格、FRAおよびスワップレートなどから、各将来日の割引係数を整合的に推定した結果である。

同じ期間構造は、割引係数、ゼロ金利またはフォワード金利として表現できる。これらは相互に変換可能であるが、表示方法によって曲線の特徴の見え方が異なる。割引係数は現在価値計算に直接用いられ、ゼロ金利は満期ごとの平均的な金利水準を示し、フォワード金利は将来の区間ごとの限界的な金利を示す。

本章では、最初にゼロ金利ノードから単純なカーブを構築し、割引係数、ゼロ金利およびフォワード金利の関係を確認する。次に、補間方式による曲線形状の違いを比較する。その後、預金と債券、ならびに預金とスワップを用いたブートストラップを行い、最後にImplied
Term Structureと平行シフトを確認する。

\begin{Shaded}
\begin{Highlighting}[]
\NormalTok{knitr}\SpecialCharTok{::}\NormalTok{opts\_chunk}\SpecialCharTok{$}\FunctionTok{set}\NormalTok{(}
  \AttributeTok{collapse =} \ConstantTok{TRUE}\NormalTok{,}
  \AttributeTok{comment =} \StringTok{"\#\textgreater{}"}\NormalTok{,}
  \AttributeTok{warning =} \ConstantTok{FALSE}\NormalTok{,}
  \AttributeTok{message =} \ConstantTok{FALSE}
\NormalTok{)}

\FunctionTok{suppressPackageStartupMessages}\NormalTok{(\{}
  \FunctionTok{library}\NormalTok{(tidyverse)}
  \FunctionTok{library}\NormalTok{(QuantLib)}
  \FunctionTok{library}\NormalTok{(GaussQuant)}
\NormalTok{\})}
\end{Highlighting}
\end{Shaded}

\section{2.1
金利カーブの基礎}\label{ux91d1ux5229ux30abux30fcux30d6ux306eux57faux790e}

評価日を0、将来時点を \(t\) とし、時点 \(t\)
に1を受け取るキャッシュフローの現在価値を \(D(0,t)\) とする。\(D(0,t)\)
は割引係数である。連続複利のゼロ金利を \(z(0,t)\)
とすると、両者の関係は次式で表される。

\[
D(0,t)
=
\exp\left[-z(0,t)t\right]
\]

したがって、割引係数からゼロ金利を求めると次式となる。

\[
z(0,t)
=
-\frac{\log D(0,t)}{t}
\]

二つの時点 \(t_1\) と \(t_2\) の間に適用される連続複利のフォワード金利を
\(f(t_1,t_2)\) とすると、次式で求められる。

\[
f(t_1,t_2)
=
-\frac{
\log\left[D(0,t_2)/D(0,t_1)\right]
}{
t_2-t_1
}
\]

割引係数、ゼロ金利およびフォワード金利は同じ期間構造を異なる方法で表している。ただし、補間方法の違いは、とくにフォワード金利の形状に強く表れる。

\subsection{2.1.1
ゼロ金利ノード}\label{ux30bcux30edux91d1ux5229ux30ceux30fcux30c9}

最初に、評価日と複数の将来日について連続複利のゼロ金利を設定する。ここで用いる数値は、カーブの構造を確認するための例示値である。

\begin{Shaded}
\begin{Highlighting}[]
\NormalTok{evaluation\_date\_chr }\OtherTok{\textless{}{-}} \StringTok{"2022{-}08{-}19"}

\NormalTok{GaussQuant}\SpecialCharTok{::}\FunctionTok{set\_eval\_date\_GQL}\NormalTok{(evaluation\_date\_chr)}

\NormalTok{zero\_nodes }\OtherTok{\textless{}{-}}\NormalTok{ tibble}\SpecialCharTok{::}\FunctionTok{tribble}\NormalTok{(}
  \SpecialCharTok{\textasciitilde{}}\NormalTok{tenor, }\SpecialCharTok{\textasciitilde{}}\NormalTok{date, }\SpecialCharTok{\textasciitilde{}}\NormalTok{zero\_rate,}
  \StringTok{"0D"}\NormalTok{,  }\FunctionTok{as.Date}\NormalTok{(}\StringTok{"2022{-}08{-}19"}\NormalTok{), }\FloatTok{0.0005}\NormalTok{,}
  \StringTok{"1M"}\NormalTok{,  }\FunctionTok{as.Date}\NormalTok{(}\StringTok{"2022{-}09{-}19"}\NormalTok{), }\FloatTok{0.0006}\NormalTok{,}
  \StringTok{"3M"}\NormalTok{,  }\FunctionTok{as.Date}\NormalTok{(}\StringTok{"2022{-}11{-}21"}\NormalTok{), }\FloatTok{0.0008}\NormalTok{,}
  \StringTok{"6M"}\NormalTok{,  }\FunctionTok{as.Date}\NormalTok{(}\StringTok{"2023{-}02{-}20"}\NormalTok{), }\FloatTok{0.0011}\NormalTok{,}
  \StringTok{"1Y"}\NormalTok{,  }\FunctionTok{as.Date}\NormalTok{(}\StringTok{"2023{-}08{-}21"}\NormalTok{), }\FloatTok{0.0018}\NormalTok{,}
  \StringTok{"2Y"}\NormalTok{,  }\FunctionTok{as.Date}\NormalTok{(}\StringTok{"2024{-}08{-}19"}\NormalTok{), }\FloatTok{0.0032}\NormalTok{,}
  \StringTok{"3Y"}\NormalTok{,  }\FunctionTok{as.Date}\NormalTok{(}\StringTok{"2025{-}08{-}19"}\NormalTok{), }\FloatTok{0.0048}\NormalTok{,}
  \StringTok{"5Y"}\NormalTok{,  }\FunctionTok{as.Date}\NormalTok{(}\StringTok{"2027{-}08{-}19"}\NormalTok{), }\FloatTok{0.0075}\NormalTok{,}
  \StringTok{"10Y"}\NormalTok{, }\FunctionTok{as.Date}\NormalTok{(}\StringTok{"2032{-}08{-}19"}\NormalTok{), }\FloatTok{0.0115}
\NormalTok{)}

\NormalTok{GaussQuant}\SpecialCharTok{::}\FunctionTok{show\_tbl\_GQH}\NormalTok{(}
\NormalTok{  zero\_nodes,}
  \StringTok{"Zero{-}rate nodes"}\NormalTok{,}
  \AttributeTok{n =} \DecValTok{20}
\NormalTok{)}
\CommentTok{\#\textgreater{} \# A tibble: 9 x 3}
\CommentTok{\#\textgreater{}   tenor date       zero\_rate}
\CommentTok{\#\textgreater{}   \textless{}chr\textgreater{} \textless{}date\textgreater{}         \textless{}dbl\textgreater{}}
\CommentTok{\#\textgreater{} 1 0D    2022{-}08{-}19    0.0005}
\CommentTok{\#\textgreater{} 2 1M    2022{-}09{-}19    0.0006}
\CommentTok{\#\textgreater{} 3 3M    2022{-}11{-}21    0.0008}
\CommentTok{\#\textgreater{} 4 6M    2023{-}02{-}20    0.0011}
\CommentTok{\#\textgreater{} 5 1Y    2023{-}08{-}21    0.0018}
\CommentTok{\#\textgreater{} 6 2Y    2024{-}08{-}19    0.0032}
\CommentTok{\#\textgreater{} 7 3Y    2025{-}08{-}19    0.0048}
\CommentTok{\#\textgreater{} 8 5Y    2027{-}08{-}19    0.0075}
\CommentTok{\#\textgreater{} 9 10Y   2032{-}08{-}19    0.0115}
\end{Highlighting}
\end{Shaded}

\subsection{2.1.2 ゼロカーブ}\label{ux30bcux30edux30abux30fcux30d6}

\texttt{build\_zero\_curve\_GQL()}は、日付とゼロ金利の組からQuantLibのZeroCurveを構築する。ここではカーブの時間軸にActual/365
Fixedを用いる。

\begin{Shaded}
\begin{Highlighting}[]
\NormalTok{zero\_curve }\OtherTok{\textless{}{-}}\NormalTok{ GaussQuant}\SpecialCharTok{::}\FunctionTok{build\_zero\_curve\_GQL}\NormalTok{(}
  \AttributeTok{date\_chr =} \FunctionTok{as.character}\NormalTok{(zero\_nodes}\SpecialCharTok{$}\NormalTok{date),}
  \AttributeTok{zero\_rates =}\NormalTok{ zero\_nodes}\SpecialCharTok{$}\NormalTok{zero\_rate,}
  \AttributeTok{day\_counter =} \StringTok{"Actual365Fixed"}
\NormalTok{)}

\NormalTok{zero\_curve\_handle }\OtherTok{\textless{}{-}}\NormalTok{ GaussQuant}\SpecialCharTok{::}\FunctionTok{yield\_curve\_handle\_GQL}\NormalTok{(}
\NormalTok{  zero\_curve}
\NormalTok{)}

\NormalTok{zero\_curve\_tbl }\OtherTok{\textless{}{-}}\NormalTok{ GaussQuant}\SpecialCharTok{::}\FunctionTok{curve\_tbl\_GQH}\NormalTok{(}
  \AttributeTok{curve =}\NormalTok{ zero\_curve,}
  \AttributeTok{tenors =} \FunctionTok{c}\NormalTok{(}
    \StringTok{"1D"}\NormalTok{, }\StringTok{"1W"}\NormalTok{, }\StringTok{"1M"}\NormalTok{, }\StringTok{"3M"}\NormalTok{, }\StringTok{"6M"}\NormalTok{,}
    \StringTok{"1Y"}\NormalTok{, }\StringTok{"2Y"}\NormalTok{, }\StringTok{"3Y"}\NormalTok{, }\StringTok{"5Y"}\NormalTok{, }\StringTok{"7Y"}\NormalTok{, }\StringTok{"10Y"}
\NormalTok{  )}
\NormalTok{)}

\NormalTok{GaussQuant}\SpecialCharTok{::}\FunctionTok{show\_tbl\_GQH}\NormalTok{(}
\NormalTok{  zero\_curve\_tbl,}
  \StringTok{"Zero curve"}\NormalTok{,}
  \AttributeTok{n =} \DecValTok{20}
\NormalTok{)}
\CommentTok{\#\textgreater{} \# A tibble: 11 x 4}
\CommentTok{\#\textgreater{}    tenor date       discount zero\_rate}
\CommentTok{\#\textgreater{}    \textless{}chr\textgreater{} \textless{}date\textgreater{}        \textless{}dbl\textgreater{}     \textless{}dbl\textgreater{}}
\CommentTok{\#\textgreater{}  1 1D    2022{-}08{-}20    1.000  0.000503}
\CommentTok{\#\textgreater{}  2 1W    2022{-}08{-}26    1.000  0.000523}
\CommentTok{\#\textgreater{}  3 1M    2022{-}09{-}19    1.000  0.000600}
\CommentTok{\#\textgreater{}  4 3M    2022{-}11{-}19    1.000  0.000794}
\CommentTok{\#\textgreater{}  5 6M    2023{-}02{-}19    0.999  0.00110 }
\CommentTok{\#\textgreater{}  6 1Y    2023{-}08{-}19    0.998  0.00179 }
\CommentTok{\#\textgreater{}  7 2Y    2024{-}08{-}19    0.994  0.00320 }
\CommentTok{\#\textgreater{}  8 3Y    2025{-}08{-}19    0.986  0.00480 }
\CommentTok{\#\textgreater{}  9 5Y    2027{-}08{-}19    0.963  0.0075  }
\CommentTok{\#\textgreater{} 10 7Y    2029{-}08{-}19    0.938  0.00910 }
\CommentTok{\#\textgreater{} 11 10Y   2032{-}08{-}19    0.891  0.0115}
\end{Highlighting}
\end{Shaded}

カーブから取得したゼロ金利は連続複利表示である。入力ノードの間ではQuantLibの補間規則に従って金利が補間される。

\subsection{2.1.3
カーブの表現}\label{ux30abux30fcux30d6ux306eux8868ux73fe}

カーブを細かい時間間隔へ展開し、割引係数、ゼロ金利および隣接区間のフォワード金利を計算する。

\begin{Shaded}
\begin{Highlighting}[]
\NormalTok{zero\_curve\_grid }\OtherTok{\textless{}{-}}\NormalTok{ GaussQuant}\SpecialCharTok{::}\FunctionTok{curve\_grid\_tbl\_GQL}\NormalTok{(}
  \AttributeTok{curve =}\NormalTok{ zero\_curve,}
  \AttributeTok{n =} \DecValTok{240}\NormalTok{L,}
  \AttributeTok{extrapolate =} \ConstantTok{TRUE}
\NormalTok{) }\SpecialCharTok{|\textgreater{}}
\NormalTok{  dplyr}\SpecialCharTok{::}\FunctionTok{mutate}\NormalTok{(}
    \AttributeTok{next\_time =}\NormalTok{ dplyr}\SpecialCharTok{::}\FunctionTok{lead}\NormalTok{(.data}\SpecialCharTok{$}\NormalTok{time),}
    \AttributeTok{next\_discount =}\NormalTok{ dplyr}\SpecialCharTok{::}\FunctionTok{lead}\NormalTok{(.data}\SpecialCharTok{$}\NormalTok{discount\_factor),}
    \AttributeTok{forward\_rate =}\NormalTok{ dplyr}\SpecialCharTok{::}\FunctionTok{if\_else}\NormalTok{(}
      \SpecialCharTok{!}\FunctionTok{is.na}\NormalTok{(.data}\SpecialCharTok{$}\NormalTok{next\_time) }\SpecialCharTok{\&}
\NormalTok{        .data}\SpecialCharTok{$}\NormalTok{next\_time }\SpecialCharTok{\textgreater{}}\NormalTok{ .data}\SpecialCharTok{$}\NormalTok{time }\SpecialCharTok{\&}
\NormalTok{        .data}\SpecialCharTok{$}\NormalTok{discount\_factor }\SpecialCharTok{\textgreater{}} \DecValTok{0} \SpecialCharTok{\&}
\NormalTok{        .data}\SpecialCharTok{$}\NormalTok{next\_discount }\SpecialCharTok{\textgreater{}} \DecValTok{0}\NormalTok{,}
      \SpecialCharTok{{-}}\FunctionTok{log}\NormalTok{(.data}\SpecialCharTok{$}\NormalTok{next\_discount }\SpecialCharTok{/}\NormalTok{ .data}\SpecialCharTok{$}\NormalTok{discount\_factor) }\SpecialCharTok{/}
\NormalTok{        (.data}\SpecialCharTok{$}\NormalTok{next\_time }\SpecialCharTok{{-}}\NormalTok{ .data}\SpecialCharTok{$}\NormalTok{time),}
      \ConstantTok{NA\_real\_}
\NormalTok{    )}
\NormalTok{  )}

\NormalTok{grid\_show\_index }\OtherTok{\textless{}{-}} \FunctionTok{unique}\NormalTok{(}
  \FunctionTok{round}\NormalTok{(}
    \FunctionTok{seq}\NormalTok{(}
      \AttributeTok{from =} \DecValTok{1}\NormalTok{,}
      \AttributeTok{to =} \FunctionTok{nrow}\NormalTok{(zero\_curve\_grid),}
      \AttributeTok{length.out =} \DecValTok{12}
\NormalTok{    )}
\NormalTok{  )}
\NormalTok{)}

\NormalTok{GaussQuant}\SpecialCharTok{::}\FunctionTok{show\_tbl\_GQH}\NormalTok{(}
\NormalTok{  zero\_curve\_grid }\SpecialCharTok{|\textgreater{}}
\NormalTok{    dplyr}\SpecialCharTok{::}\FunctionTok{slice}\NormalTok{(grid\_show\_index) }\SpecialCharTok{|\textgreater{}}
\NormalTok{    dplyr}\SpecialCharTok{::}\FunctionTok{select}\NormalTok{(}
\NormalTok{      .data}\SpecialCharTok{$}\NormalTok{curve\_date,}
\NormalTok{      .data}\SpecialCharTok{$}\NormalTok{time,}
\NormalTok{      .data}\SpecialCharTok{$}\NormalTok{discount\_factor,}
\NormalTok{      .data}\SpecialCharTok{$}\NormalTok{zero\_rate,}
\NormalTok{      .data}\SpecialCharTok{$}\NormalTok{forward\_rate}
\NormalTok{    ),}
  \StringTok{"Zero{-}curve grid"}\NormalTok{,}
  \AttributeTok{n =} \DecValTok{20}
\NormalTok{)}
\CommentTok{\#\textgreater{} \# A tibble: 12 x 5}
\CommentTok{\#\textgreater{}    curve\_date   time discount\_factor zero\_rate forward\_rate}
\CommentTok{\#\textgreater{}    \textless{}date\textgreater{}      \textless{}dbl\textgreater{}           \textless{}dbl\textgreater{}     \textless{}dbl\textgreater{}        \textless{}dbl\textgreater{}}
\CommentTok{\#\textgreater{}  1 2022{-}08{-}19  0               1       0           0.000548}
\CommentTok{\#\textgreater{}  2 2023{-}07{-}21  0.921           0.998   0.00168     0.00303 }
\CommentTok{\#\textgreater{}  3 2024{-}06{-}06  1.8             0.995   0.00292     0.00550 }
\CommentTok{\#\textgreater{}  4 2025{-}05{-}08  2.72            0.988   0.00435     0.00877 }
\CommentTok{\#\textgreater{}  5 2026{-}04{-}09  3.64            0.980   0.00566     0.0106  }
\CommentTok{\#\textgreater{}  6 2027{-}03{-}12  4.56            0.969   0.00691     0.0131  }
\CommentTok{\#\textgreater{}  7 2028{-}01{-}26  5.44            0.958   0.00785     0.0122  }
\CommentTok{\#\textgreater{}  8 2028{-}12{-}28  6.36            0.947   0.00859     0.0137  }
\CommentTok{\#\textgreater{}  9 2029{-}11{-}29  7.28            0.934   0.00932     0.0152  }
\CommentTok{\#\textgreater{} 10 2030{-}10{-}31  8.21            0.921   0.0101      0.0167  }
\CommentTok{\#\textgreater{} 11 2031{-}09{-}17  9.08            0.907   0.0108      0.0181  }
\CommentTok{\#\textgreater{} 12 2032{-}08{-}19 10.0             0.891   0.0115     NA}
\end{Highlighting}
\end{Shaded}

\begin{Shaded}
\begin{Highlighting}[]
\NormalTok{discount\_factor\_plot }\OtherTok{\textless{}{-}}\NormalTok{ zero\_curve\_grid }\SpecialCharTok{|\textgreater{}}
\NormalTok{  ggplot2}\SpecialCharTok{::}\FunctionTok{ggplot}\NormalTok{(}
\NormalTok{    ggplot2}\SpecialCharTok{::}\FunctionTok{aes}\NormalTok{(}
      \AttributeTok{x =}\NormalTok{ .data}\SpecialCharTok{$}\NormalTok{curve\_date,}
      \AttributeTok{y =}\NormalTok{ .data}\SpecialCharTok{$}\NormalTok{discount\_factor}
\NormalTok{    )}
\NormalTok{  ) }\SpecialCharTok{+}
\NormalTok{  ggplot2}\SpecialCharTok{::}\FunctionTok{geom\_line}\NormalTok{() }\SpecialCharTok{+}
\NormalTok{  ggplot2}\SpecialCharTok{::}\FunctionTok{labs}\NormalTok{(}
    \AttributeTok{title =} \StringTok{"Discount factor curve"}\NormalTok{,}
    \AttributeTok{x =} \ConstantTok{NULL}\NormalTok{,}
    \AttributeTok{y =} \StringTok{"Discount factor"}
\NormalTok{  ) }\SpecialCharTok{+}
\NormalTok{  ggplot2}\SpecialCharTok{::}\FunctionTok{theme\_minimal}\NormalTok{()}

\NormalTok{zero\_rate\_plot }\OtherTok{\textless{}{-}}\NormalTok{ zero\_curve\_grid }\SpecialCharTok{|\textgreater{}}
\NormalTok{  ggplot2}\SpecialCharTok{::}\FunctionTok{ggplot}\NormalTok{(}
\NormalTok{    ggplot2}\SpecialCharTok{::}\FunctionTok{aes}\NormalTok{(}
      \AttributeTok{x =}\NormalTok{ .data}\SpecialCharTok{$}\NormalTok{curve\_date,}
      \AttributeTok{y =} \DecValTok{100} \SpecialCharTok{*}\NormalTok{ .data}\SpecialCharTok{$}\NormalTok{zero\_rate}
\NormalTok{    )}
\NormalTok{  ) }\SpecialCharTok{+}
\NormalTok{  ggplot2}\SpecialCharTok{::}\FunctionTok{geom\_line}\NormalTok{() }\SpecialCharTok{+}
\NormalTok{  ggplot2}\SpecialCharTok{::}\FunctionTok{labs}\NormalTok{(}
    \AttributeTok{title =} \StringTok{"Continuous zero{-}rate curve"}\NormalTok{,}
    \AttributeTok{x =} \ConstantTok{NULL}\NormalTok{,}
    \AttributeTok{y =} \StringTok{"Zero rate (\%)"}
\NormalTok{  ) }\SpecialCharTok{+}
\NormalTok{  ggplot2}\SpecialCharTok{::}\FunctionTok{theme\_minimal}\NormalTok{()}

\NormalTok{forward\_rate\_plot }\OtherTok{\textless{}{-}}\NormalTok{ zero\_curve\_grid }\SpecialCharTok{|\textgreater{}}
\NormalTok{  dplyr}\SpecialCharTok{::}\FunctionTok{filter}\NormalTok{(}\SpecialCharTok{!}\FunctionTok{is.na}\NormalTok{(.data}\SpecialCharTok{$}\NormalTok{forward\_rate)) }\SpecialCharTok{|\textgreater{}}
\NormalTok{  ggplot2}\SpecialCharTok{::}\FunctionTok{ggplot}\NormalTok{(}
\NormalTok{    ggplot2}\SpecialCharTok{::}\FunctionTok{aes}\NormalTok{(}
      \AttributeTok{x =}\NormalTok{ .data}\SpecialCharTok{$}\NormalTok{curve\_date,}
      \AttributeTok{y =} \DecValTok{100} \SpecialCharTok{*}\NormalTok{ .data}\SpecialCharTok{$}\NormalTok{forward\_rate}
\NormalTok{    )}
\NormalTok{  ) }\SpecialCharTok{+}
\NormalTok{  ggplot2}\SpecialCharTok{::}\FunctionTok{geom\_line}\NormalTok{() }\SpecialCharTok{+}
\NormalTok{  ggplot2}\SpecialCharTok{::}\FunctionTok{labs}\NormalTok{(}
    \AttributeTok{title =} \StringTok{"One{-}step forward{-}rate curve"}\NormalTok{,}
    \AttributeTok{x =} \ConstantTok{NULL}\NormalTok{,}
    \AttributeTok{y =} \StringTok{"Forward rate (\%)"}
\NormalTok{  ) }\SpecialCharTok{+}
\NormalTok{  ggplot2}\SpecialCharTok{::}\FunctionTok{theme\_minimal}\NormalTok{()}

\FunctionTok{print}\NormalTok{(discount\_factor\_plot)}
\end{Highlighting}
\end{Shaded}

\pandocbounded{\includegraphics[keepaspectratio]{part2_quantive_riemann/ch02_interest_rate_curves_files/figure-pdf/zero-curve-plots-1.pdf}}

\begin{Shaded}
\begin{Highlighting}[]
\FunctionTok{print}\NormalTok{(zero\_rate\_plot)}
\end{Highlighting}
\end{Shaded}

\pandocbounded{\includegraphics[keepaspectratio]{part2_quantive_riemann/ch02_interest_rate_curves_files/figure-pdf/zero-curve-plots-2.pdf}}

\begin{Shaded}
\begin{Highlighting}[]
\FunctionTok{print}\NormalTok{(forward\_rate\_plot)}
\end{Highlighting}
\end{Shaded}

\pandocbounded{\includegraphics[keepaspectratio]{part2_quantive_riemann/ch02_interest_rate_curves_files/figure-pdf/zero-curve-plots-3.pdf}}

割引係数は将来の1単位を現在価値へ換算する倍率である。正の金利環境では、満期が長くなるほど一般に小さくなる。ゼロ金利は評価日から各満期までの平均的な金利水準を示す。一方、フォワード金利は隣接する二時点間の金利を表すため、補間方法やノード配置の影響を強く受ける。

\section{2.2 補間と参照日}\label{ux88dcux9593ux3068ux53c2ux7167ux65e5}

\subsection{2.2.1
補間方式とフォワード金利}\label{ux88dcux9593ux65b9ux5f0fux3068ux30d5ux30a9ux30efux30fcux30c9ux91d1ux5229}

市場では限られた満期の金利または価格しか観測できないため、ノード間の値を補間する必要がある。補間対象には、ゼロ金利、割引係数、対数割引係数またはフォワード金利などがある。

ここでは、同じゼロ金利ノードから次の二つのカーブを構築する。

\begin{itemize}
\tightlist
\item
  ZeroCurve：ゼロ金利を基礎とするカーブ
\item
  DiscountCurve：ノードのゼロ金利を割引係数へ変換し、割引係数を基礎とするカーブ
\end{itemize}

\begin{Shaded}
\begin{Highlighting}[]
\NormalTok{discount\_curve }\OtherTok{\textless{}{-}}\NormalTok{ GaussQuant}\SpecialCharTok{::}\FunctionTok{build\_discount\_curve\_GQL}\NormalTok{(}
  \AttributeTok{nodes =}\NormalTok{ zero\_nodes }\SpecialCharTok{|\textgreater{}}
\NormalTok{    dplyr}\SpecialCharTok{::}\FunctionTok{select}\NormalTok{(}
\NormalTok{      .data}\SpecialCharTok{$}\NormalTok{date,}
\NormalTok{      .data}\SpecialCharTok{$}\NormalTok{zero\_rate}
\NormalTok{    ),}
  \AttributeTok{day\_counter =} \StringTok{"Actual365Fixed"}
\NormalTok{)}

\NormalTok{discount\_curve\_handle }\OtherTok{\textless{}{-}}\NormalTok{ GaussQuant}\SpecialCharTok{::}\FunctionTok{yield\_curve\_handle\_GQL}\NormalTok{(}
\NormalTok{  discount\_curve}
\NormalTok{)}

\NormalTok{comparison\_max\_time }\OtherTok{\textless{}{-}} \FunctionTok{min}\NormalTok{(}
  \FunctionTok{as.numeric}\NormalTok{(zero\_curve}\SpecialCharTok{$}\FunctionTok{maxTime}\NormalTok{()),}
  \FunctionTok{as.numeric}\NormalTok{(discount\_curve}\SpecialCharTok{$}\FunctionTok{maxTime}\NormalTok{())}
\NormalTok{)}

\NormalTok{comparison\_time\_grid }\OtherTok{\textless{}{-}} \FunctionTok{seq}\NormalTok{(}
  \AttributeTok{from =} \DecValTok{0}\NormalTok{,}
  \AttributeTok{to =}\NormalTok{ comparison\_max\_time,}
  \AttributeTok{length.out =} \DecValTok{300}
\NormalTok{)}
\NormalTok{comparison\_reference\_date }\OtherTok{\textless{}{-}} \FunctionTok{as.Date}\NormalTok{(}
\NormalTok{  GaussQuant}\SpecialCharTok{::}\FunctionTok{iso\_GQL}\NormalTok{(}
\NormalTok{    zero\_curve}\SpecialCharTok{$}\FunctionTok{referenceDate}\NormalTok{()}
\NormalTok{  )}
\NormalTok{)}

\NormalTok{interpolation\_tbl }\OtherTok{\textless{}{-}}\NormalTok{ tibble}\SpecialCharTok{::}\FunctionTok{tibble}\NormalTok{(}
  \AttributeTok{time =}\NormalTok{ comparison\_time\_grid,}
  \AttributeTok{date =}\NormalTok{ comparison\_reference\_date }\SpecialCharTok{+}
    \FunctionTok{round}\NormalTok{(comparison\_time\_grid }\SpecialCharTok{*} \DecValTok{365}\NormalTok{)}
\NormalTok{) }\SpecialCharTok{|\textgreater{}}
\NormalTok{  dplyr}\SpecialCharTok{::}\FunctionTok{mutate}\NormalTok{(}
    \AttributeTok{ql\_date =}\NormalTok{ purrr}\SpecialCharTok{::}\FunctionTok{map}\NormalTok{(}
      \FunctionTok{as.character}\NormalTok{(.data}\SpecialCharTok{$}\NormalTok{date),}
\NormalTok{      GaussQuant}\SpecialCharTok{::}\NormalTok{date\_GQL}
\NormalTok{    ),}
    \AttributeTok{zero\_curve\_discount =}\NormalTok{ purrr}\SpecialCharTok{::}\FunctionTok{map\_dbl}\NormalTok{(}
\NormalTok{      .data}\SpecialCharTok{$}\NormalTok{ql\_date,}
      \SpecialCharTok{\textasciitilde{}}\NormalTok{ GaussQuant}\SpecialCharTok{::}\FunctionTok{curve\_discount\_safe\_GQL}\NormalTok{(}
\NormalTok{        zero\_curve,}
\NormalTok{        .x}
\NormalTok{      )}
\NormalTok{    ),}
    \AttributeTok{discount\_curve\_discount =}\NormalTok{ purrr}\SpecialCharTok{::}\FunctionTok{map\_dbl}\NormalTok{(}
\NormalTok{      .data}\SpecialCharTok{$}\NormalTok{ql\_date,}
      \SpecialCharTok{\textasciitilde{}}\NormalTok{ GaussQuant}\SpecialCharTok{::}\FunctionTok{curve\_discount\_safe\_GQL}\NormalTok{(}
\NormalTok{        discount\_curve,}
\NormalTok{        .x}
\NormalTok{      )}
\NormalTok{    )}
\NormalTok{  ) }\SpecialCharTok{|\textgreater{}}
\NormalTok{  dplyr}\SpecialCharTok{::}\FunctionTok{select}\NormalTok{(}
    \SpecialCharTok{{-}}\NormalTok{.data}\SpecialCharTok{$}\NormalTok{ql\_date}
\NormalTok{  ) }\SpecialCharTok{|\textgreater{}}
\NormalTok{  tidyr}\SpecialCharTok{::}\FunctionTok{pivot\_longer}\NormalTok{(}
    \AttributeTok{cols =} \FunctionTok{c}\NormalTok{(}
\NormalTok{      zero\_curve\_discount,}
\NormalTok{      discount\_curve\_discount}
\NormalTok{    ),}
    \AttributeTok{names\_to =} \StringTok{"interpolation"}\NormalTok{,}
    \AttributeTok{values\_to =} \StringTok{"discount\_factor"}
\NormalTok{  ) }\SpecialCharTok{|\textgreater{}}
\NormalTok{  dplyr}\SpecialCharTok{::}\FunctionTok{mutate}\NormalTok{(}
    \AttributeTok{interpolation =}\NormalTok{ dplyr}\SpecialCharTok{::}\FunctionTok{recode}\NormalTok{(}
\NormalTok{      .data}\SpecialCharTok{$}\NormalTok{interpolation,}
      \AttributeTok{zero\_curve\_discount =} \StringTok{"ZeroCurve"}\NormalTok{,}
      \AttributeTok{discount\_curve\_discount =} \StringTok{"DiscountCurve"}
\NormalTok{    )}
\NormalTok{  ) }\SpecialCharTok{|\textgreater{}}
\NormalTok{  dplyr}\SpecialCharTok{::}\FunctionTok{group\_by}\NormalTok{(}
\NormalTok{    .data}\SpecialCharTok{$}\NormalTok{interpolation}
\NormalTok{  ) }\SpecialCharTok{|\textgreater{}}
\NormalTok{  dplyr}\SpecialCharTok{::}\FunctionTok{arrange}\NormalTok{(}
\NormalTok{    .data}\SpecialCharTok{$}\NormalTok{interpolation,}
\NormalTok{    .data}\SpecialCharTok{$}\NormalTok{time}
\NormalTok{  ) }\SpecialCharTok{|\textgreater{}}
\NormalTok{  dplyr}\SpecialCharTok{::}\FunctionTok{mutate}\NormalTok{(}
    \AttributeTok{zero\_rate =}\NormalTok{ dplyr}\SpecialCharTok{::}\FunctionTok{if\_else}\NormalTok{(}
\NormalTok{      .data}\SpecialCharTok{$}\NormalTok{time }\SpecialCharTok{\textgreater{}} \DecValTok{0} \SpecialCharTok{\&}
\NormalTok{        .data}\SpecialCharTok{$}\NormalTok{discount\_factor }\SpecialCharTok{\textgreater{}} \DecValTok{0}\NormalTok{,}
      \SpecialCharTok{{-}}\FunctionTok{log}\NormalTok{(.data}\SpecialCharTok{$}\NormalTok{discount\_factor) }\SpecialCharTok{/}
\NormalTok{        .data}\SpecialCharTok{$}\NormalTok{time,}
\NormalTok{      zero\_nodes}\SpecialCharTok{$}\NormalTok{zero\_rate[[}\DecValTok{1}\NormalTok{]]}
\NormalTok{    ),}
    \AttributeTok{next\_time =}\NormalTok{ dplyr}\SpecialCharTok{::}\FunctionTok{lead}\NormalTok{(}
\NormalTok{      .data}\SpecialCharTok{$}\NormalTok{time}
\NormalTok{    ),}
    \AttributeTok{next\_discount =}\NormalTok{ dplyr}\SpecialCharTok{::}\FunctionTok{lead}\NormalTok{(}
\NormalTok{      .data}\SpecialCharTok{$}\NormalTok{discount\_factor}
\NormalTok{    ),}
    \AttributeTok{forward\_rate =}\NormalTok{ dplyr}\SpecialCharTok{::}\FunctionTok{if\_else}\NormalTok{(}
      \SpecialCharTok{!}\FunctionTok{is.na}\NormalTok{(.data}\SpecialCharTok{$}\NormalTok{next\_time) }\SpecialCharTok{\&}
\NormalTok{        .data}\SpecialCharTok{$}\NormalTok{next\_time }\SpecialCharTok{\textgreater{}}\NormalTok{ .data}\SpecialCharTok{$}\NormalTok{time }\SpecialCharTok{\&}
\NormalTok{        .data}\SpecialCharTok{$}\NormalTok{discount\_factor }\SpecialCharTok{\textgreater{}} \DecValTok{0} \SpecialCharTok{\&}
\NormalTok{        .data}\SpecialCharTok{$}\NormalTok{next\_discount }\SpecialCharTok{\textgreater{}} \DecValTok{0}\NormalTok{,}
      \SpecialCharTok{{-}}\FunctionTok{log}\NormalTok{(}
\NormalTok{        .data}\SpecialCharTok{$}\NormalTok{next\_discount }\SpecialCharTok{/}
\NormalTok{          .data}\SpecialCharTok{$}\NormalTok{discount\_factor}
\NormalTok{      ) }\SpecialCharTok{/}\NormalTok{ (}
\NormalTok{        .data}\SpecialCharTok{$}\NormalTok{next\_time }\SpecialCharTok{{-}}
\NormalTok{          .data}\SpecialCharTok{$}\NormalTok{time}
\NormalTok{      ),}
      \ConstantTok{NA\_real\_}
\NormalTok{    )}
\NormalTok{  ) }\SpecialCharTok{|\textgreater{}}
\NormalTok{  dplyr}\SpecialCharTok{::}\FunctionTok{ungroup}\NormalTok{()}
\end{Highlighting}
\end{Shaded}

\begin{Shaded}
\begin{Highlighting}[]
\NormalTok{interpolation\_zero\_plot }\OtherTok{\textless{}{-}}\NormalTok{ interpolation\_tbl }\SpecialCharTok{|\textgreater{}}
\NormalTok{  ggplot2}\SpecialCharTok{::}\FunctionTok{ggplot}\NormalTok{(}
\NormalTok{    ggplot2}\SpecialCharTok{::}\FunctionTok{aes}\NormalTok{(}
      \AttributeTok{x =}\NormalTok{ .data}\SpecialCharTok{$}\NormalTok{time,}
      \AttributeTok{y =} \DecValTok{100} \SpecialCharTok{*}\NormalTok{ .data}\SpecialCharTok{$}\NormalTok{zero\_rate,}
      \AttributeTok{linetype =}\NormalTok{ .data}\SpecialCharTok{$}\NormalTok{interpolation}
\NormalTok{    )}
\NormalTok{  ) }\SpecialCharTok{+}
\NormalTok{  ggplot2}\SpecialCharTok{::}\FunctionTok{geom\_line}\NormalTok{() }\SpecialCharTok{+}
\NormalTok{  ggplot2}\SpecialCharTok{::}\FunctionTok{labs}\NormalTok{(}
    \AttributeTok{title =} \StringTok{"Zero rates under different interpolation targets"}\NormalTok{,}
    \AttributeTok{x =} \StringTok{"Time from reference date"}\NormalTok{,}
    \AttributeTok{y =} \StringTok{"Zero rate (\%)"}\NormalTok{,}
    \AttributeTok{linetype =} \ConstantTok{NULL}
\NormalTok{  ) }\SpecialCharTok{+}
\NormalTok{  ggplot2}\SpecialCharTok{::}\FunctionTok{theme\_minimal}\NormalTok{()}

\NormalTok{interpolation\_forward\_plot }\OtherTok{\textless{}{-}}\NormalTok{ interpolation\_tbl }\SpecialCharTok{|\textgreater{}}
\NormalTok{  dplyr}\SpecialCharTok{::}\FunctionTok{filter}\NormalTok{(}\SpecialCharTok{!}\FunctionTok{is.na}\NormalTok{(.data}\SpecialCharTok{$}\NormalTok{forward\_rate)) }\SpecialCharTok{|\textgreater{}}
\NormalTok{  ggplot2}\SpecialCharTok{::}\FunctionTok{ggplot}\NormalTok{(}
\NormalTok{    ggplot2}\SpecialCharTok{::}\FunctionTok{aes}\NormalTok{(}
      \AttributeTok{x =}\NormalTok{ .data}\SpecialCharTok{$}\NormalTok{time,}
      \AttributeTok{y =} \DecValTok{100} \SpecialCharTok{*}\NormalTok{ .data}\SpecialCharTok{$}\NormalTok{forward\_rate,}
      \AttributeTok{linetype =}\NormalTok{ .data}\SpecialCharTok{$}\NormalTok{interpolation}
\NormalTok{    )}
\NormalTok{  ) }\SpecialCharTok{+}
\NormalTok{  ggplot2}\SpecialCharTok{::}\FunctionTok{geom\_line}\NormalTok{() }\SpecialCharTok{+}
\NormalTok{  ggplot2}\SpecialCharTok{::}\FunctionTok{labs}\NormalTok{(}
    \AttributeTok{title =} \StringTok{"Forward rates under different interpolation targets"}\NormalTok{,}
    \AttributeTok{x =} \StringTok{"Time from reference date"}\NormalTok{,}
    \AttributeTok{y =} \StringTok{"Forward rate (\%)"}\NormalTok{,}
    \AttributeTok{linetype =} \ConstantTok{NULL}
\NormalTok{  ) }\SpecialCharTok{+}
\NormalTok{  ggplot2}\SpecialCharTok{::}\FunctionTok{theme\_minimal}\NormalTok{()}

\FunctionTok{print}\NormalTok{(interpolation\_zero\_plot)}
\end{Highlighting}
\end{Shaded}

\pandocbounded{\includegraphics[keepaspectratio]{part2_quantive_riemann/ch02_interest_rate_curves_files/figure-pdf/interpolation-plots-1.pdf}}

\begin{Shaded}
\begin{Highlighting}[]
\FunctionTok{print}\NormalTok{(interpolation\_forward\_plot)}
\end{Highlighting}
\end{Shaded}

\pandocbounded{\includegraphics[keepaspectratio]{part2_quantive_riemann/ch02_interest_rate_curves_files/figure-pdf/interpolation-plots-2.pdf}}

両カーブは入力ノードでほぼ同じ情報を再現するが、ノード間では異なる形状を示す。ゼロ金利のグラフでは差が小さく見えても、フォワード金利へ変換すると差が拡大することがある。したがって、カーブの妥当性はゼロ金利だけでなく、フォワード金利も確認する必要がある。

\subsection{2.2.2
参照日と評価日}\label{ux53c2ux7167ux65e5ux3068ux8a55ux4fa1ux65e5}

カーブの参照日は、割引係数が1となる時点であり、時間軸の起点である。日付ノードから構築したカーブでは、通常、最初のノード日が固定参照日となる。

次の確認では、QuantLibの評価日を変更した後も、すでに構築済みのZeroCurveの参照日が変わらないことを確認する。

\begin{Shaded}
\begin{Highlighting}[]
\NormalTok{reference\_date\_before }\OtherTok{\textless{}{-}} \FunctionTok{as.Date}\NormalTok{(}
\NormalTok{  GaussQuant}\SpecialCharTok{::}\FunctionTok{iso\_GQL}\NormalTok{(}
\NormalTok{    zero\_curve}\SpecialCharTok{$}\FunctionTok{referenceDate}\NormalTok{()}
\NormalTok{  )}
\NormalTok{)}

\NormalTok{GaussQuant}\SpecialCharTok{::}\FunctionTok{set\_eval\_date\_GQL}\NormalTok{(}
  \StringTok{"2022{-}09{-}19"}
\NormalTok{)}

\NormalTok{reference\_date\_after }\OtherTok{\textless{}{-}} \FunctionTok{as.Date}\NormalTok{(}
\NormalTok{  GaussQuant}\SpecialCharTok{::}\FunctionTok{iso\_GQL}\NormalTok{(}
\NormalTok{    zero\_curve}\SpecialCharTok{$}\FunctionTok{referenceDate}\NormalTok{()}
\NormalTok{  )}
\NormalTok{)}

\NormalTok{reference\_date\_tbl }\OtherTok{\textless{}{-}}\NormalTok{ tibble}\SpecialCharTok{::}\FunctionTok{tribble}\NormalTok{(}
  \SpecialCharTok{\textasciitilde{}}\NormalTok{item, }\SpecialCharTok{\textasciitilde{}}\NormalTok{value,}
  \StringTok{"curve\_reference\_date\_before"}\NormalTok{, }\FunctionTok{as.character}\NormalTok{(reference\_date\_before),}
  \StringTok{"evaluation\_date\_after\_change"}\NormalTok{, }\StringTok{"2022{-}09{-}19"}\NormalTok{,}
  \StringTok{"curve\_reference\_date\_after"}\NormalTok{, }\FunctionTok{as.character}\NormalTok{(reference\_date\_after)}
\NormalTok{)}

\NormalTok{GaussQuant}\SpecialCharTok{::}\FunctionTok{show\_tbl\_GQH}\NormalTok{(}
\NormalTok{  reference\_date\_tbl,}
  \StringTok{"Fixed reference date"}
\NormalTok{)}
\CommentTok{\#\textgreater{} \# A tibble: 3 x 2}
\CommentTok{\#\textgreater{}   item                         value     }
\CommentTok{\#\textgreater{}   \textless{}chr\textgreater{}                        \textless{}chr\textgreater{}     }
\CommentTok{\#\textgreater{} 1 curve\_reference\_date\_before  2022{-}08{-}19}
\CommentTok{\#\textgreater{} 2 evaluation\_date\_after\_change 2022{-}09{-}19}
\CommentTok{\#\textgreater{} 3 curve\_reference\_date\_after   2022{-}08{-}19}

\NormalTok{GaussQuant}\SpecialCharTok{::}\FunctionTok{set\_eval\_date\_GQL}\NormalTok{(}
\NormalTok{  evaluation\_date\_chr}
\NormalTok{)}
\end{Highlighting}
\end{Shaded}

固定参照日のカーブは、評価日を変更しても参照日が変化しない。一方、決済日数と休日カレンダーから参照日を決定するMoving
Reference
Date型のカーブでは、評価日の変更に応じて参照日も移動する。評価日を進めるシナリオ分析では、既存カーブを固定したまま使うのか、市場データから新しいカーブを再構築するのかを区別する必要がある。

\section{2.3
ブートストラップ}\label{ux30d6ux30fcux30c8ux30b9ux30c8ux30e9ux30c3ux30d7}

\subsection{2.3.1 預金と債券}\label{ux9810ux91d1ux3068ux50b5ux5238}

市場ではゼロ金利を直接観測できる満期は限られている。そのため、短期預金、国債価格およびスワップレートなどから、割引係数を順次推定するブートストラップを行う。

短期部分では預金金利から最初の割引係数を求め、その先では固定利付債の市場価格と理論価格が一致するようにカーブを延長する。GaussQuantの\texttt{build\_bond\_discount\_curve\_GQL()}は、短期預金と複数の固定利付債を用いてPiecewiseFlatForwardカーブを構築する。

\begin{Shaded}
\begin{Highlighting}[]
\NormalTok{GaussQuant}\SpecialCharTok{::}\FunctionTok{set\_eval\_date\_GQL}\NormalTok{(}
  \StringTok{"2008{-}09{-}15"}
\NormalTok{)}

\NormalTok{bond\_curve }\OtherTok{\textless{}{-}}\NormalTok{ GaussQuant}\SpecialCharTok{::}\FunctionTok{build\_bond\_discount\_curve\_GQL}\NormalTok{(}
  \AttributeTok{settlement\_date =} \StringTok{"2008{-}09{-}18"}\NormalTok{,}
  \AttributeTok{fixing\_days =} \DecValTok{3}\NormalTok{,}
  \AttributeTok{settlement\_days =} \DecValTok{3}
\NormalTok{)}

\NormalTok{bond\_curve\_tbl }\OtherTok{\textless{}{-}}\NormalTok{ GaussQuant}\SpecialCharTok{::}\FunctionTok{curve\_tbl\_GQH}\NormalTok{(}
  \AttributeTok{curve =}\NormalTok{ bond\_curve,}
  \AttributeTok{tenors =} \FunctionTok{c}\NormalTok{(}
    \StringTok{"1M"}\NormalTok{, }\StringTok{"3M"}\NormalTok{, }\StringTok{"6M"}\NormalTok{, }\StringTok{"1Y"}\NormalTok{,}
    \StringTok{"2Y"}\NormalTok{, }\StringTok{"3Y"}\NormalTok{, }\StringTok{"5Y"}\NormalTok{, }\StringTok{"7Y"}\NormalTok{,}
    \StringTok{"10Y"}\NormalTok{, }\StringTok{"15Y"}\NormalTok{, }\StringTok{"20Y"}
\NormalTok{  )}
\NormalTok{)}

\NormalTok{GaussQuant}\SpecialCharTok{::}\FunctionTok{show\_tbl\_GQH}\NormalTok{(}
\NormalTok{  bond\_curve\_tbl,}
  \StringTok{"Deposit and bond bootstrap curve"}\NormalTok{,}
  \AttributeTok{n =} \DecValTok{30}
\NormalTok{)}
\CommentTok{\#\textgreater{} \# A tibble: 11 x 4}
\CommentTok{\#\textgreater{}    tenor date       discount zero\_rate}
\CommentTok{\#\textgreater{}    \textless{}chr\textgreater{} \textless{}date\textgreater{}        \textless{}dbl\textgreater{}     \textless{}dbl\textgreater{}}
\CommentTok{\#\textgreater{}  1 1M    2008{-}10{-}18    0.999   0.00961}
\CommentTok{\#\textgreater{}  2 3M    2008{-}12{-}18    0.998   0.00961}
\CommentTok{\#\textgreater{}  3 6M    2009{-}03{-}18    0.993   0.0145 }
\CommentTok{\#\textgreater{}  4 1Y    2009{-}09{-}18    0.981   0.0192 }
\CommentTok{\#\textgreater{}  5 2Y    2010{-}09{-}18    0.957   0.0218 }
\CommentTok{\#\textgreater{}  6 3Y    2011{-}09{-}18    0.929   0.0245 }
\CommentTok{\#\textgreater{}  7 5Y    2013{-}09{-}18    0.860   0.0302 }
\CommentTok{\#\textgreater{}  8 7Y    2015{-}09{-}18    0.782   0.0351 }
\CommentTok{\#\textgreater{}  9 10Y   2018{-}09{-}18    0.679   0.0388 }
\CommentTok{\#\textgreater{} 10 15Y   2023{-}09{-}18    0.532   0.0421 }
\CommentTok{\#\textgreater{} 11 20Y   2028{-}09{-}18    0.417   0.0438}
\end{Highlighting}
\end{Shaded}

\begin{Shaded}
\begin{Highlighting}[]
\NormalTok{bond\_curve\_grid }\OtherTok{\textless{}{-}}\NormalTok{ GaussQuant}\SpecialCharTok{::}\FunctionTok{curve\_grid\_tbl\_GQL}\NormalTok{(}
  \AttributeTok{curve =}\NormalTok{ bond\_curve,}
  \AttributeTok{n =} \DecValTok{300}\NormalTok{L,}
  \AttributeTok{extrapolate =} \ConstantTok{TRUE}
\NormalTok{)}

\NormalTok{bond\_curve\_plot }\OtherTok{\textless{}{-}}\NormalTok{ bond\_curve\_grid }\SpecialCharTok{|\textgreater{}}
\NormalTok{  ggplot2}\SpecialCharTok{::}\FunctionTok{ggplot}\NormalTok{(}
\NormalTok{    ggplot2}\SpecialCharTok{::}\FunctionTok{aes}\NormalTok{(}
      \AttributeTok{x =}\NormalTok{ .data}\SpecialCharTok{$}\NormalTok{curve\_date,}
      \AttributeTok{y =} \DecValTok{100} \SpecialCharTok{*}\NormalTok{ .data}\SpecialCharTok{$}\NormalTok{zero\_rate}
\NormalTok{    )}
\NormalTok{  ) }\SpecialCharTok{+}
\NormalTok{  ggplot2}\SpecialCharTok{::}\FunctionTok{geom\_line}\NormalTok{() }\SpecialCharTok{+}
\NormalTok{  ggplot2}\SpecialCharTok{::}\FunctionTok{labs}\NormalTok{(}
    \AttributeTok{title =} \StringTok{"Deposit and bond bootstrap curve"}\NormalTok{,}
    \AttributeTok{x =} \ConstantTok{NULL}\NormalTok{,}
    \AttributeTok{y =} \StringTok{"Continuous zero rate (\%)"}
\NormalTok{  ) }\SpecialCharTok{+}
\NormalTok{  ggplot2}\SpecialCharTok{::}\FunctionTok{theme\_minimal}\NormalTok{()}

\FunctionTok{print}\NormalTok{(bond\_curve\_plot)}
\end{Highlighting}
\end{Shaded}

\pandocbounded{\includegraphics[keepaspectratio]{part2_quantive_riemann/ch02_interest_rate_curves_files/figure-pdf/bond-bootstrap-plot-1.pdf}}

債券価格を用いるブートストラップでは、各債券のクーポン日、満期日、経過利子および日数計算規則がカーブ構築に影響する。債券のクーポン計算には商品固有の日数計算規則を用いる一方、期間構造の時間軸には単純で加法的な日数計算規則を用いる必要がある。

\subsection{2.3.2
預金とスワップ}\label{ux9810ux91d1ux3068ux30b9ux30efux30c3ux30d7}

次に、短期部分を預金金利、中長期部分を固定対変動金利スワップの市場レートから構築する。スワップの市場レートは、固定レグと変動レグの現在価値を等しくするパー・スワップレートである。

GaussQuantの\texttt{build\_swap\_curve\_GQL()}は、短期預金と複数満期のスワップレートを用いてPiecewiseFlatForwardカーブを構築する。

\begin{Shaded}
\begin{Highlighting}[]
\NormalTok{swap\_curve }\OtherTok{\textless{}{-}}\NormalTok{ GaussQuant}\SpecialCharTok{::}\FunctionTok{build\_swap\_curve\_GQL}\NormalTok{(}
  \AttributeTok{settlement\_date =} \StringTok{"2008{-}09{-}18"}\NormalTok{,}
  \AttributeTok{fixing\_days =} \DecValTok{3}
\NormalTok{)}

\NormalTok{swap\_curve\_handle }\OtherTok{\textless{}{-}}\NormalTok{ GaussQuant}\SpecialCharTok{::}\FunctionTok{yield\_curve\_handle\_GQL}\NormalTok{(}
\NormalTok{  swap\_curve}
\NormalTok{)}

\NormalTok{swap\_curve\_tbl }\OtherTok{\textless{}{-}}\NormalTok{ GaussQuant}\SpecialCharTok{::}\FunctionTok{curve\_tbl\_GQH}\NormalTok{(}
  \AttributeTok{curve =}\NormalTok{ swap\_curve,}
  \AttributeTok{tenors =} \FunctionTok{c}\NormalTok{(}
    \StringTok{"1W"}\NormalTok{, }\StringTok{"1M"}\NormalTok{, }\StringTok{"3M"}\NormalTok{, }\StringTok{"6M"}\NormalTok{,}
    \StringTok{"1Y"}\NormalTok{, }\StringTok{"2Y"}\NormalTok{, }\StringTok{"3Y"}\NormalTok{, }\StringTok{"5Y"}\NormalTok{,}
    \StringTok{"7Y"}\NormalTok{, }\StringTok{"10Y"}\NormalTok{, }\StringTok{"15Y"}
\NormalTok{  )}
\NormalTok{)}

\NormalTok{GaussQuant}\SpecialCharTok{::}\FunctionTok{show\_tbl\_GQH}\NormalTok{(}
\NormalTok{  swap\_curve\_tbl,}
  \StringTok{"Deposit and swap bootstrap curve"}\NormalTok{,}
  \AttributeTok{n =} \DecValTok{30}
\NormalTok{)}
\CommentTok{\#\textgreater{} \# A tibble: 11 x 4}
\CommentTok{\#\textgreater{}    tenor date       discount zero\_rate}
\CommentTok{\#\textgreater{}    \textless{}chr\textgreater{} \textless{}date\textgreater{}        \textless{}dbl\textgreater{}     \textless{}dbl\textgreater{}}
\CommentTok{\#\textgreater{}  1 1W    2008{-}09{-}25    0.999    0.0441}
\CommentTok{\#\textgreater{}  2 1M    2008{-}10{-}18    0.997    0.0326}
\CommentTok{\#\textgreater{}  3 3M    2008{-}12{-}18    0.992    0.0324}
\CommentTok{\#\textgreater{}  4 6M    2009{-}03{-}18    0.983    0.0341}
\CommentTok{\#\textgreater{}  5 1Y    2009{-}09{-}18    0.967    0.0334}
\CommentTok{\#\textgreater{}  6 2Y    2010{-}09{-}18    0.944    0.0290}
\CommentTok{\#\textgreater{}  7 3Y    2011{-}09{-}18    0.909    0.0318}
\CommentTok{\#\textgreater{}  8 5Y    2013{-}09{-}18    0.837    0.0355}
\CommentTok{\#\textgreater{}  9 7Y    2015{-}09{-}18    0.763    0.0387}
\CommentTok{\#\textgreater{} 10 10Y   2018{-}09{-}18    0.663    0.0411}
\CommentTok{\#\textgreater{} 11 15Y   2023{-}09{-}18    0.521    0.0434}
\end{Highlighting}
\end{Shaded}

\begin{Shaded}
\begin{Highlighting}[]
\NormalTok{comparison\_tenors }\OtherTok{\textless{}{-}} \FunctionTok{c}\NormalTok{(}
  \StringTok{"1M"}\NormalTok{, }\StringTok{"3M"}\NormalTok{, }\StringTok{"6M"}\NormalTok{, }\StringTok{"1Y"}\NormalTok{,}
  \StringTok{"2Y"}\NormalTok{, }\StringTok{"3Y"}\NormalTok{, }\StringTok{"5Y"}\NormalTok{, }\StringTok{"7Y"}\NormalTok{,}
  \StringTok{"10Y"}\NormalTok{, }\StringTok{"15Y"}
\NormalTok{)}

\NormalTok{bootstrap\_comparison\_tbl }\OtherTok{\textless{}{-}}\NormalTok{ dplyr}\SpecialCharTok{::}\FunctionTok{bind\_rows}\NormalTok{(}
\NormalTok{  GaussQuant}\SpecialCharTok{::}\FunctionTok{curve\_tbl\_GQH}\NormalTok{(}
    \AttributeTok{curve =}\NormalTok{ bond\_curve,}
    \AttributeTok{tenors =}\NormalTok{ comparison\_tenors}
\NormalTok{  ) }\SpecialCharTok{|\textgreater{}}
\NormalTok{    dplyr}\SpecialCharTok{::}\FunctionTok{mutate}\NormalTok{(}
      \AttributeTok{curve =} \StringTok{"Deposit + bond"}
\NormalTok{    ),}
\NormalTok{  GaussQuant}\SpecialCharTok{::}\FunctionTok{curve\_tbl\_GQH}\NormalTok{(}
    \AttributeTok{curve =}\NormalTok{ swap\_curve,}
    \AttributeTok{tenors =}\NormalTok{ comparison\_tenors}
\NormalTok{  ) }\SpecialCharTok{|\textgreater{}}
\NormalTok{    dplyr}\SpecialCharTok{::}\FunctionTok{mutate}\NormalTok{(}
      \AttributeTok{curve =} \StringTok{"Deposit + swap"}
\NormalTok{    )}
\NormalTok{)}

\NormalTok{GaussQuant}\SpecialCharTok{::}\FunctionTok{show\_tbl\_GQH}\NormalTok{(}
\NormalTok{  bootstrap\_comparison\_tbl,}
  \StringTok{"Bootstrap curve comparison"}\NormalTok{,}
  \AttributeTok{n =} \DecValTok{30}
\NormalTok{)}
\CommentTok{\#\textgreater{} \# A tibble: 20 x 5}
\CommentTok{\#\textgreater{}    tenor date       discount zero\_rate curve         }
\CommentTok{\#\textgreater{}    \textless{}chr\textgreater{} \textless{}date\textgreater{}        \textless{}dbl\textgreater{}     \textless{}dbl\textgreater{} \textless{}chr\textgreater{}         }
\CommentTok{\#\textgreater{}  1 1M    2008{-}10{-}18    0.999   0.00961 Deposit + bond}
\CommentTok{\#\textgreater{}  2 3M    2008{-}12{-}18    0.998   0.00961 Deposit + bond}
\CommentTok{\#\textgreater{}  3 6M    2009{-}03{-}18    0.993   0.0145  Deposit + bond}
\CommentTok{\#\textgreater{}  4 1Y    2009{-}09{-}18    0.981   0.0192  Deposit + bond}
\CommentTok{\#\textgreater{}  5 2Y    2010{-}09{-}18    0.957   0.0218  Deposit + bond}
\CommentTok{\#\textgreater{}  6 3Y    2011{-}09{-}18    0.929   0.0245  Deposit + bond}
\CommentTok{\#\textgreater{}  7 5Y    2013{-}09{-}18    0.860   0.0302  Deposit + bond}
\CommentTok{\#\textgreater{}  8 7Y    2015{-}09{-}18    0.782   0.0351  Deposit + bond}
\CommentTok{\#\textgreater{}  9 10Y   2018{-}09{-}18    0.679   0.0388  Deposit + bond}
\CommentTok{\#\textgreater{} 10 15Y   2023{-}09{-}18    0.532   0.0421  Deposit + bond}
\CommentTok{\#\textgreater{} 11 1M    2008{-}10{-}18    0.997   0.0326  Deposit + swap}
\CommentTok{\#\textgreater{} 12 3M    2008{-}12{-}18    0.992   0.0324  Deposit + swap}
\CommentTok{\#\textgreater{} 13 6M    2009{-}03{-}18    0.983   0.0341  Deposit + swap}
\CommentTok{\#\textgreater{} 14 1Y    2009{-}09{-}18    0.967   0.0334  Deposit + swap}
\CommentTok{\#\textgreater{} 15 2Y    2010{-}09{-}18    0.944   0.0290  Deposit + swap}
\CommentTok{\#\textgreater{} 16 3Y    2011{-}09{-}18    0.909   0.0318  Deposit + swap}
\CommentTok{\#\textgreater{} 17 5Y    2013{-}09{-}18    0.837   0.0355  Deposit + swap}
\CommentTok{\#\textgreater{} 18 7Y    2015{-}09{-}18    0.763   0.0387  Deposit + swap}
\CommentTok{\#\textgreater{} 19 10Y   2018{-}09{-}18    0.663   0.0411  Deposit + swap}
\CommentTok{\#\textgreater{} 20 15Y   2023{-}09{-}18    0.521   0.0434  Deposit + swap}

\NormalTok{bootstrap\_zero\_plot }\OtherTok{\textless{}{-}}\NormalTok{ bootstrap\_comparison\_tbl }\SpecialCharTok{|\textgreater{}}
\NormalTok{  ggplot2}\SpecialCharTok{::}\FunctionTok{ggplot}\NormalTok{(}
\NormalTok{    ggplot2}\SpecialCharTok{::}\FunctionTok{aes}\NormalTok{(}
      \AttributeTok{x =}\NormalTok{ .data}\SpecialCharTok{$}\NormalTok{date,}
      \AttributeTok{y =} \DecValTok{100} \SpecialCharTok{*}\NormalTok{ .data}\SpecialCharTok{$}\NormalTok{zero\_rate,}
      \AttributeTok{linetype =}\NormalTok{ .data}\SpecialCharTok{$}\NormalTok{curve}
\NormalTok{    )}
\NormalTok{  ) }\SpecialCharTok{+}
\NormalTok{  ggplot2}\SpecialCharTok{::}\FunctionTok{geom\_line}\NormalTok{() }\SpecialCharTok{+}
\NormalTok{  ggplot2}\SpecialCharTok{::}\FunctionTok{geom\_point}\NormalTok{() }\SpecialCharTok{+}
\NormalTok{  ggplot2}\SpecialCharTok{::}\FunctionTok{labs}\NormalTok{(}
    \AttributeTok{title =} \StringTok{"Bond{-}based and swap{-}based curves"}\NormalTok{,}
    \AttributeTok{x =} \ConstantTok{NULL}\NormalTok{,}
    \AttributeTok{y =} \StringTok{"Continuous zero rate (\%)"}\NormalTok{,}
    \AttributeTok{linetype =} \ConstantTok{NULL}
\NormalTok{  ) }\SpecialCharTok{+}
\NormalTok{  ggplot2}\SpecialCharTok{::}\FunctionTok{theme\_minimal}\NormalTok{()}

\FunctionTok{print}\NormalTok{(bootstrap\_zero\_plot)}
\end{Highlighting}
\end{Shaded}

\pandocbounded{\includegraphics[keepaspectratio]{part2_quantive_riemann/ch02_interest_rate_curves_files/figure-pdf/bootstrap-comparison-1.pdf}}

債券カーブとスワップカーブは、異なる市場商品から構築されるため一致するとは限らない。両者の差には、市場の流動性、商品固有の信用・担保条件、日数計算規則および観測時点の違いが含まれる。

FRAは、将来の一定期間に適用される金利を直接観測する商品として、預金とスワップの間を接続するために用いられる。本章の実装では預金とスワップを用いるが、FRAを追加する場合も、各市場商品の理論価格が市場価格と一致するように同じブートストラップの枠組みへ組み込む。

\section{2.4
カーブ操作と検証}\label{ux30abux30fcux30d6ux64cdux4f5cux3068ux691cux8a3c}

\subsection{2.4.1 Implied Term Structure}\label{implied-term-structure}

Implied Term
Structureは、既存カーブが示す将来の割引関係を保ったまま、参照日を将来へ移した期間構造である。元のカーブの参照日を
\(t_0\)、新しい参照日を \(t_1\)、将来日を \(t_2\) とすると、Implied Term
Structureの割引係数は概念的に次式で表される。

\[
D_{t_1}(t_2)
=
\frac{
D_{t_0}(t_2)
}{
D_{t_0}(t_1)
}
\]

\begin{Shaded}
\begin{Highlighting}[]
\NormalTok{implied\_reference\_date }\OtherTok{\textless{}{-}}\NormalTok{ GaussQuant}\SpecialCharTok{::}\FunctionTok{date\_GQL}\NormalTok{(}
  \StringTok{"2009{-}09{-}18"}
\NormalTok{)}

\NormalTok{implied\_constructor }\OtherTok{\textless{}{-}} \FunctionTok{get0}\NormalTok{(}
  \StringTok{"ImpliedTermStructure"}\NormalTok{,}
  \AttributeTok{envir =} \FunctionTok{asNamespace}\NormalTok{(}\StringTok{"QuantLib"}\NormalTok{),}
  \AttributeTok{inherits =} \ConstantTok{FALSE}
\NormalTok{)}

\NormalTok{implied\_curve }\OtherTok{\textless{}{-}} \ControlFlowTok{if}\NormalTok{ (}\FunctionTok{is.null}\NormalTok{(implied\_constructor)) \{}
  \ConstantTok{NULL}
\NormalTok{\} }\ControlFlowTok{else}\NormalTok{ \{}
  \FunctionTok{tryCatch}\NormalTok{(}
    \FunctionTok{implied\_constructor}\NormalTok{(}
\NormalTok{      swap\_curve\_handle,}
\NormalTok{      implied\_reference\_date}
\NormalTok{    ),}
    \AttributeTok{error =} \ControlFlowTok{function}\NormalTok{(e) }\ConstantTok{NULL}
\NormalTok{  )}
\NormalTok{\}}

\NormalTok{implied\_curve\_tbl }\OtherTok{\textless{}{-}} \ControlFlowTok{if}\NormalTok{ (}\FunctionTok{is.null}\NormalTok{(implied\_curve)) \{}
\NormalTok{  tibble}\SpecialCharTok{::}\FunctionTok{tibble}\NormalTok{(}
    \AttributeTok{status =} \StringTok{"ImpliedTermStructure is not available in the installed QuantLib binding."}
\NormalTok{  )}
\NormalTok{\} }\ControlFlowTok{else}\NormalTok{ \{}
\NormalTok{  GaussQuant}\SpecialCharTok{::}\FunctionTok{curve\_tbl\_GQH}\NormalTok{(}
    \AttributeTok{curve =}\NormalTok{ implied\_curve,}
    \AttributeTok{tenors =} \FunctionTok{c}\NormalTok{(}
      \StringTok{"1M"}\NormalTok{, }\StringTok{"3M"}\NormalTok{, }\StringTok{"6M"}\NormalTok{, }\StringTok{"1Y"}\NormalTok{,}
      \StringTok{"2Y"}\NormalTok{, }\StringTok{"3Y"}\NormalTok{, }\StringTok{"5Y"}\NormalTok{, }\StringTok{"10Y"}
\NormalTok{    )}
\NormalTok{  )}
\NormalTok{\}}

\NormalTok{GaussQuant}\SpecialCharTok{::}\FunctionTok{show\_tbl\_GQH}\NormalTok{(}
\NormalTok{  implied\_curve\_tbl,}
  \StringTok{"Implied term structure"}\NormalTok{,}
  \AttributeTok{n =} \DecValTok{20}
\NormalTok{)}
\CommentTok{\#\textgreater{} \# A tibble: 8 x 4}
\CommentTok{\#\textgreater{}   tenor date       discount zero\_rate}
\CommentTok{\#\textgreater{}   \textless{}chr\textgreater{} \textless{}date\textgreater{}        \textless{}dbl\textgreater{}     \textless{}dbl\textgreater{}}
\CommentTok{\#\textgreater{} 1 1M    2009{-}10{-}18    0.998    0.0245}
\CommentTok{\#\textgreater{} 2 3M    2009{-}12{-}18    0.994    0.0245}
\CommentTok{\#\textgreater{} 3 6M    2010{-}03{-}18    0.988    0.0245}
\CommentTok{\#\textgreater{} 4 1Y    2010{-}09{-}18    0.976    0.0245}
\CommentTok{\#\textgreater{} 5 2Y    2011{-}09{-}18    0.940    0.0310}
\CommentTok{\#\textgreater{} 6 3Y    2012{-}09{-}18    0.902    0.0343}
\CommentTok{\#\textgreater{} 7 5Y    2014{-}09{-}18    0.826    0.0381}
\CommentTok{\#\textgreater{} 8 10Y   2019{-}09{-}18    0.653    0.0426}
\end{Highlighting}
\end{Shaded}

Implied Term
Structureは、将来時点から見た割引係数を現在のカーブから機械的に導くものであり、将来の市場カーブを予測したものではない。評価日を将来へ進めたシナリオで市場リスクを評価する場合には、カーブの単純な時間経過と市場金利の変動を区別する必要がある。

\subsection{2.4.2
カーブの平行シフト}\label{ux30abux30fcux30d6ux306eux5e73ux884cux30b7ux30d5ux30c8}

金利リスク分析では、基準カーブに一定のスプレッドを加えたカーブを用いる。すべての満期へ同じスプレッドを加える変化を平行シフトという。

ここでは、預金・スワップカーブへ50ベーシスポイントを加え、ゼロ金利と割引係数の変化を確認する。

\begin{Shaded}
\begin{Highlighting}[]
\NormalTok{parallel\_spread }\OtherTok{\textless{}{-}} \DecValTok{50} \SpecialCharTok{/} \DecValTok{10000}

\NormalTok{spreaded\_curve\_bundle }\OtherTok{\textless{}{-}}\NormalTok{ GaussQuant}\SpecialCharTok{::}\FunctionTok{make\_zero\_spreaded\_curve\_GQL}\NormalTok{(}
  \AttributeTok{base\_curve\_handle =}\NormalTok{ swap\_curve\_handle,}
  \AttributeTok{spread\_rate =}\NormalTok{ parallel\_spread,}
  \AttributeTok{day\_counter =}\NormalTok{ QuantLib}\SpecialCharTok{::}\FunctionTok{Actual365Fixed}\NormalTok{(),}
  \AttributeTok{extrapolate =} \ConstantTok{TRUE}
\NormalTok{)}

\NormalTok{spreaded\_curve }\OtherTok{\textless{}{-}}\NormalTok{ spreaded\_curve\_bundle}\SpecialCharTok{$}\NormalTok{curve}

\NormalTok{parallel\_shift\_tbl }\OtherTok{\textless{}{-}}\NormalTok{ dplyr}\SpecialCharTok{::}\FunctionTok{bind\_rows}\NormalTok{(}
\NormalTok{  GaussQuant}\SpecialCharTok{::}\FunctionTok{curve\_tbl\_GQH}\NormalTok{(}
    \AttributeTok{curve =}\NormalTok{ swap\_curve,}
    \AttributeTok{tenors =}\NormalTok{ comparison\_tenors}
\NormalTok{  ) }\SpecialCharTok{|\textgreater{}}
\NormalTok{    dplyr}\SpecialCharTok{::}\FunctionTok{mutate}\NormalTok{(}
      \AttributeTok{curve =} \StringTok{"Base curve"}
\NormalTok{    ),}
\NormalTok{  GaussQuant}\SpecialCharTok{::}\FunctionTok{curve\_tbl\_GQH}\NormalTok{(}
    \AttributeTok{curve =}\NormalTok{ spreaded\_curve,}
    \AttributeTok{tenors =}\NormalTok{ comparison\_tenors}
\NormalTok{  ) }\SpecialCharTok{|\textgreater{}}
\NormalTok{    dplyr}\SpecialCharTok{::}\FunctionTok{mutate}\NormalTok{(}
      \AttributeTok{curve =} \StringTok{"Base curve + 50bp"}
\NormalTok{    )}
\NormalTok{)}

\NormalTok{parallel\_shift\_wide\_tbl }\OtherTok{\textless{}{-}}\NormalTok{ parallel\_shift\_tbl }\SpecialCharTok{|\textgreater{}}
\NormalTok{  dplyr}\SpecialCharTok{::}\FunctionTok{select}\NormalTok{(}
\NormalTok{    .data}\SpecialCharTok{$}\NormalTok{tenor,}
\NormalTok{    .data}\SpecialCharTok{$}\NormalTok{date,}
\NormalTok{    .data}\SpecialCharTok{$}\NormalTok{curve,}
\NormalTok{    .data}\SpecialCharTok{$}\NormalTok{discount,}
\NormalTok{    .data}\SpecialCharTok{$}\NormalTok{zero\_rate}
\NormalTok{  ) }\SpecialCharTok{|\textgreater{}}
\NormalTok{  tidyr}\SpecialCharTok{::}\FunctionTok{pivot\_wider}\NormalTok{(}
    \AttributeTok{names\_from =}\NormalTok{ curve,}
    \AttributeTok{values\_from =} \FunctionTok{c}\NormalTok{(}
\NormalTok{      discount,}
\NormalTok{      zero\_rate}
\NormalTok{    )}
\NormalTok{  ) }\SpecialCharTok{|\textgreater{}}
\NormalTok{  dplyr}\SpecialCharTok{::}\FunctionTok{mutate}\NormalTok{(}
    \AttributeTok{observed\_zero\_shift\_bp =}
      \DecValTok{10000} \SpecialCharTok{*}\NormalTok{ (}
\NormalTok{        .data}\SpecialCharTok{$}\StringTok{\textasciigrave{}}\AttributeTok{zero\_rate\_Base curve + 50bp}\StringTok{\textasciigrave{}} \SpecialCharTok{{-}}
\NormalTok{          .data}\SpecialCharTok{$}\StringTok{\textasciigrave{}}\AttributeTok{zero\_rate\_Base curve}\StringTok{\textasciigrave{}}
\NormalTok{      )}
\NormalTok{  )}

\NormalTok{GaussQuant}\SpecialCharTok{::}\FunctionTok{show\_tbl\_GQH}\NormalTok{(}
\NormalTok{  parallel\_shift\_wide\_tbl,}
  \StringTok{"Parallel curve shift"}\NormalTok{,}
  \AttributeTok{n =} \DecValTok{30}
\NormalTok{)}
\CommentTok{\#\textgreater{} \# A tibble: 10 x 7}
\CommentTok{\#\textgreater{}    tenor date       \textasciigrave{}discount\_Base curve\textasciigrave{} \textasciigrave{}discount\_Base curve + 50bp\textasciigrave{}}
\CommentTok{\#\textgreater{}    \textless{}chr\textgreater{} \textless{}date\textgreater{}                     \textless{}dbl\textgreater{}                        \textless{}dbl\textgreater{}}
\CommentTok{\#\textgreater{}  1 1M    2008{-}10{-}18                 0.997                        0.997}
\CommentTok{\#\textgreater{}  2 3M    2008{-}12{-}18                 0.992                        0.991}
\CommentTok{\#\textgreater{}  3 6M    2009{-}03{-}18                 0.983                        0.981}
\CommentTok{\#\textgreater{}  4 1Y    2009{-}09{-}18                 0.967                        0.962}
\CommentTok{\#\textgreater{}  5 2Y    2010{-}09{-}18                 0.944                        0.935}
\CommentTok{\#\textgreater{}  6 3Y    2011{-}09{-}18                 0.909                        0.897}
\CommentTok{\#\textgreater{}  7 5Y    2013{-}09{-}18                 0.837                        0.820}
\CommentTok{\#\textgreater{}  8 7Y    2015{-}09{-}18                 0.763                        0.743}
\CommentTok{\#\textgreater{}  9 10Y   2018{-}09{-}18                 0.663                        0.642}
\CommentTok{\#\textgreater{} 10 15Y   2023{-}09{-}18                 0.521                        0.502}
\CommentTok{\#\textgreater{} \# i 3 more variables: \textasciigrave{}zero\_rate\_Base curve\textasciigrave{} \textless{}dbl\textgreater{},}
\CommentTok{\#\textgreater{} \#   \textasciigrave{}zero\_rate\_Base curve + 50bp\textasciigrave{} \textless{}dbl\textgreater{}, observed\_zero\_shift\_bp \textless{}dbl\textgreater{}}
\end{Highlighting}
\end{Shaded}

\begin{Shaded}
\begin{Highlighting}[]
\NormalTok{parallel\_shift\_plot }\OtherTok{\textless{}{-}}\NormalTok{ parallel\_shift\_tbl }\SpecialCharTok{|\textgreater{}}
\NormalTok{  ggplot2}\SpecialCharTok{::}\FunctionTok{ggplot}\NormalTok{(}
\NormalTok{    ggplot2}\SpecialCharTok{::}\FunctionTok{aes}\NormalTok{(}
      \AttributeTok{x =}\NormalTok{ .data}\SpecialCharTok{$}\NormalTok{date,}
      \AttributeTok{y =} \DecValTok{100} \SpecialCharTok{*}\NormalTok{ .data}\SpecialCharTok{$}\NormalTok{zero\_rate,}
      \AttributeTok{linetype =}\NormalTok{ .data}\SpecialCharTok{$}\NormalTok{curve}
\NormalTok{    )}
\NormalTok{  ) }\SpecialCharTok{+}
\NormalTok{  ggplot2}\SpecialCharTok{::}\FunctionTok{geom\_line}\NormalTok{() }\SpecialCharTok{+}
\NormalTok{  ggplot2}\SpecialCharTok{::}\FunctionTok{geom\_point}\NormalTok{() }\SpecialCharTok{+}
\NormalTok{  ggplot2}\SpecialCharTok{::}\FunctionTok{labs}\NormalTok{(}
    \AttributeTok{title =} \StringTok{"Parallel shift of the zero curve"}\NormalTok{,}
    \AttributeTok{x =} \ConstantTok{NULL}\NormalTok{,}
    \AttributeTok{y =} \StringTok{"Continuous zero rate (\%)"}\NormalTok{,}
    \AttributeTok{linetype =} \ConstantTok{NULL}
\NormalTok{  ) }\SpecialCharTok{+}
\NormalTok{  ggplot2}\SpecialCharTok{::}\FunctionTok{theme\_minimal}\NormalTok{()}

\FunctionTok{print}\NormalTok{(parallel\_shift\_plot)}
\end{Highlighting}
\end{Shaded}

\pandocbounded{\includegraphics[keepaspectratio]{part2_quantive_riemann/ch02_interest_rate_curves_files/figure-pdf/parallel-shift-plot-1.pdf}}

正の平行シフトでは割引率が上昇し、将来キャッシュフローの割引係数は低下する。そのため、通常の固定キャッシュフロー商品の現在価値も低下する。

実際の市場変動は完全な平行移動とは限らない。短期だけが動く変化、長期だけが動く変化、曲線の傾きが変化するスティープニングまたはフラットニングも生じる。年限別のKey
Rate
Shiftでは、特定年限のノードを動かし、周辺年限へ補間しながら再評価する。これらは平行シフトより細かなリスク分解である。

\subsection{2.4.3
カーブの検証}\label{ux30abux30fcux30d6ux306eux691cux8a3c}

構築したカーブについて、少なくとも次の点を確認する必要がある。

\begin{itemize}
\tightlist
\item
  参照日と最終日が想定どおりであること
\item
  割引係数が正であること
\item
  欠損値や計算不能な区間がないこと
\item
  ゼロ金利とフォワード金利に不自然な跳びがないこと
\item
  入力商品の理論価格または理論レートが市場値へ再現されること
\end{itemize}

まず、例示的なゼロ金利ノードが二つのカーブでどの程度再現されるかを確認する。

\begin{Shaded}
\begin{Highlighting}[]
\NormalTok{node\_recovery\_tbl }\OtherTok{\textless{}{-}}\NormalTok{ zero\_nodes }\SpecialCharTok{|\textgreater{}}
\NormalTok{  dplyr}\SpecialCharTok{::}\FunctionTok{mutate}\NormalTok{(}
    \AttributeTok{ql\_date =}\NormalTok{ purrr}\SpecialCharTok{::}\FunctionTok{map}\NormalTok{(}
      \FunctionTok{as.character}\NormalTok{(.data}\SpecialCharTok{$}\NormalTok{date),}
\NormalTok{      GaussQuant}\SpecialCharTok{::}\NormalTok{date\_GQL}
\NormalTok{    ),}
    \AttributeTok{time =} \FunctionTok{as.numeric}\NormalTok{(}
\NormalTok{      .data}\SpecialCharTok{$}\NormalTok{date }\SpecialCharTok{{-}}
        \FunctionTok{as.Date}\NormalTok{(evaluation\_date\_chr)}
\NormalTok{    ) }\SpecialCharTok{/} \DecValTok{365}\NormalTok{,}
    \AttributeTok{zero\_curve\_discount =}\NormalTok{ purrr}\SpecialCharTok{::}\FunctionTok{map\_dbl}\NormalTok{(}
\NormalTok{      .data}\SpecialCharTok{$}\NormalTok{ql\_date,}
      \SpecialCharTok{\textasciitilde{}}\NormalTok{ GaussQuant}\SpecialCharTok{::}\FunctionTok{curve\_discount\_safe\_GQL}\NormalTok{(}
\NormalTok{        zero\_curve,}
\NormalTok{        .x}
\NormalTok{      )}
\NormalTok{    ),}
    \AttributeTok{discount\_curve\_discount =}\NormalTok{ purrr}\SpecialCharTok{::}\FunctionTok{map\_dbl}\NormalTok{(}
\NormalTok{      .data}\SpecialCharTok{$}\NormalTok{ql\_date,}
      \SpecialCharTok{\textasciitilde{}}\NormalTok{ GaussQuant}\SpecialCharTok{::}\FunctionTok{curve\_discount\_safe\_GQL}\NormalTok{(}
\NormalTok{        discount\_curve,}
\NormalTok{        .x}
\NormalTok{      )}
\NormalTok{    ),}
    \AttributeTok{zero\_curve\_implied\_rate =}\NormalTok{ dplyr}\SpecialCharTok{::}\FunctionTok{if\_else}\NormalTok{(}
\NormalTok{      .data}\SpecialCharTok{$}\NormalTok{time }\SpecialCharTok{\textgreater{}} \DecValTok{0}\NormalTok{,}
      \SpecialCharTok{{-}}\FunctionTok{log}\NormalTok{(.data}\SpecialCharTok{$}\NormalTok{zero\_curve\_discount) }\SpecialCharTok{/}
\NormalTok{        .data}\SpecialCharTok{$}\NormalTok{time,}
\NormalTok{      .data}\SpecialCharTok{$}\NormalTok{zero\_rate}
\NormalTok{    ),}
    \AttributeTok{discount\_curve\_implied\_rate =}\NormalTok{ dplyr}\SpecialCharTok{::}\FunctionTok{if\_else}\NormalTok{(}
\NormalTok{      .data}\SpecialCharTok{$}\NormalTok{time }\SpecialCharTok{\textgreater{}} \DecValTok{0}\NormalTok{,}
      \SpecialCharTok{{-}}\FunctionTok{log}\NormalTok{(.data}\SpecialCharTok{$}\NormalTok{discount\_curve\_discount) }\SpecialCharTok{/}
\NormalTok{        .data}\SpecialCharTok{$}\NormalTok{time,}
\NormalTok{      .data}\SpecialCharTok{$}\NormalTok{zero\_rate}
\NormalTok{    ),}
    \AttributeTok{zero\_curve\_error =}
\NormalTok{      .data}\SpecialCharTok{$}\NormalTok{zero\_curve\_implied\_rate }\SpecialCharTok{{-}}
\NormalTok{        .data}\SpecialCharTok{$}\NormalTok{zero\_rate,}
    \AttributeTok{discount\_curve\_error =}
\NormalTok{      .data}\SpecialCharTok{$}\NormalTok{discount\_curve\_implied\_rate }\SpecialCharTok{{-}}
\NormalTok{        .data}\SpecialCharTok{$}\NormalTok{zero\_rate}
\NormalTok{  ) }\SpecialCharTok{|\textgreater{}}
\NormalTok{  dplyr}\SpecialCharTok{::}\FunctionTok{select}\NormalTok{(}
\NormalTok{    .data}\SpecialCharTok{$}\NormalTok{tenor,}
\NormalTok{    .data}\SpecialCharTok{$}\NormalTok{date,}
\NormalTok{    .data}\SpecialCharTok{$}\NormalTok{zero\_rate,}
\NormalTok{    .data}\SpecialCharTok{$}\NormalTok{zero\_curve\_implied\_rate,}
\NormalTok{    .data}\SpecialCharTok{$}\NormalTok{discount\_curve\_implied\_rate,}
\NormalTok{    .data}\SpecialCharTok{$}\NormalTok{zero\_curve\_error,}
\NormalTok{    .data}\SpecialCharTok{$}\NormalTok{discount\_curve\_error}
\NormalTok{  )}

\NormalTok{GaussQuant}\SpecialCharTok{::}\FunctionTok{show\_tbl\_GQH}\NormalTok{(}
\NormalTok{  node\_recovery\_tbl,}
  \StringTok{"Node recovery check"}\NormalTok{,}
  \AttributeTok{n =} \DecValTok{20}
\NormalTok{)}
\CommentTok{\#\textgreater{} \# A tibble: 9 x 7}
\CommentTok{\#\textgreater{}   tenor date       zero\_rate zero\_curve\_implied\_rate discount\_curve\_implied\_rate}
\CommentTok{\#\textgreater{}   \textless{}chr\textgreater{} \textless{}date\textgreater{}         \textless{}dbl\textgreater{}                   \textless{}dbl\textgreater{}                       \textless{}dbl\textgreater{}}
\CommentTok{\#\textgreater{} 1 0D    2022{-}08{-}19    0.0005                0.0005                      0.0005  }
\CommentTok{\#\textgreater{} 2 1M    2022{-}09{-}19    0.0006                0.000600                    0.000600}
\CommentTok{\#\textgreater{} 3 3M    2022{-}11{-}21    0.0008                0.000800                    0.000800}
\CommentTok{\#\textgreater{} 4 6M    2023{-}02{-}20    0.0011                0.00110                     0.00110 }
\CommentTok{\#\textgreater{} 5 1Y    2023{-}08{-}21    0.0018                0.00180                     0.00180 }
\CommentTok{\#\textgreater{} 6 2Y    2024{-}08{-}19    0.0032                0.00320                     0.00320 }
\CommentTok{\#\textgreater{} 7 3Y    2025{-}08{-}19    0.0048                0.00480                     0.00480 }
\CommentTok{\#\textgreater{} 8 5Y    2027{-}08{-}19    0.0075                0.0075                      0.0075  }
\CommentTok{\#\textgreater{} 9 10Y   2032{-}08{-}19    0.0115                0.0115                      0.0115  }
\CommentTok{\#\textgreater{} \# i 2 more variables: zero\_curve\_error \textless{}dbl\textgreater{}, discount\_curve\_error \textless{}dbl\textgreater{}}
\end{Highlighting}
\end{Shaded}

続いて、各カーブを密なグリッドへ展開し、割引係数の範囲と欠損値を確認する。

\begin{Shaded}
\begin{Highlighting}[]
\NormalTok{curve\_validation\_tbl }\OtherTok{\textless{}{-}}\NormalTok{ dplyr}\SpecialCharTok{::}\FunctionTok{bind\_rows}\NormalTok{(}
\NormalTok{  GaussQuant}\SpecialCharTok{::}\FunctionTok{curve\_grid\_tbl\_GQL}\NormalTok{(}
\NormalTok{    zero\_curve,}
    \AttributeTok{n =} \DecValTok{200}\NormalTok{L,}
    \AttributeTok{extrapolate =} \ConstantTok{TRUE}
\NormalTok{  ) }\SpecialCharTok{|\textgreater{}}
\NormalTok{    dplyr}\SpecialCharTok{::}\FunctionTok{mutate}\NormalTok{(}
      \AttributeTok{curve =} \StringTok{"ZeroCurve"}
\NormalTok{    ),}
\NormalTok{  GaussQuant}\SpecialCharTok{::}\FunctionTok{curve\_grid\_tbl\_GQL}\NormalTok{(}
\NormalTok{    discount\_curve,}
    \AttributeTok{n =} \DecValTok{200}\NormalTok{L,}
    \AttributeTok{extrapolate =} \ConstantTok{TRUE}
\NormalTok{  ) }\SpecialCharTok{|\textgreater{}}
\NormalTok{    dplyr}\SpecialCharTok{::}\FunctionTok{mutate}\NormalTok{(}
      \AttributeTok{curve =} \StringTok{"DiscountCurve"}
\NormalTok{    ),}
\NormalTok{  GaussQuant}\SpecialCharTok{::}\FunctionTok{curve\_grid\_tbl\_GQL}\NormalTok{(}
\NormalTok{    bond\_curve,}
    \AttributeTok{n =} \DecValTok{200}\NormalTok{L,}
    \AttributeTok{extrapolate =} \ConstantTok{TRUE}
\NormalTok{  ) }\SpecialCharTok{|\textgreater{}}
\NormalTok{    dplyr}\SpecialCharTok{::}\FunctionTok{mutate}\NormalTok{(}
      \AttributeTok{curve =} \StringTok{"Deposit + bond"}
\NormalTok{    ),}
\NormalTok{  GaussQuant}\SpecialCharTok{::}\FunctionTok{curve\_grid\_tbl\_GQL}\NormalTok{(}
\NormalTok{    swap\_curve,}
    \AttributeTok{n =} \DecValTok{200}\NormalTok{L,}
    \AttributeTok{extrapolate =} \ConstantTok{TRUE}
\NormalTok{  ) }\SpecialCharTok{|\textgreater{}}
\NormalTok{    dplyr}\SpecialCharTok{::}\FunctionTok{mutate}\NormalTok{(}
      \AttributeTok{curve =} \StringTok{"Deposit + swap"}
\NormalTok{    ),}
\NormalTok{  GaussQuant}\SpecialCharTok{::}\FunctionTok{curve\_grid\_tbl\_GQL}\NormalTok{(}
\NormalTok{    spreaded\_curve,}
    \AttributeTok{n =} \DecValTok{200}\NormalTok{L,}
    \AttributeTok{extrapolate =} \ConstantTok{TRUE}
\NormalTok{  ) }\SpecialCharTok{|\textgreater{}}
\NormalTok{    dplyr}\SpecialCharTok{::}\FunctionTok{mutate}\NormalTok{(}
      \AttributeTok{curve =} \StringTok{"Deposit + swap + 50bp"}
\NormalTok{    )}
\NormalTok{) }\SpecialCharTok{|\textgreater{}}
\NormalTok{  dplyr}\SpecialCharTok{::}\FunctionTok{group\_by}\NormalTok{(.data}\SpecialCharTok{$}\NormalTok{curve) }\SpecialCharTok{|\textgreater{}}
\NormalTok{  dplyr}\SpecialCharTok{::}\FunctionTok{summarise}\NormalTok{(}
    \AttributeTok{first\_date =} \FunctionTok{min}\NormalTok{(}
\NormalTok{      .data}\SpecialCharTok{$}\NormalTok{curve\_date,}
      \AttributeTok{na.rm =} \ConstantTok{TRUE}
\NormalTok{    ),}
    \AttributeTok{last\_date =} \FunctionTok{max}\NormalTok{(}
\NormalTok{      .data}\SpecialCharTok{$}\NormalTok{curve\_date,}
      \AttributeTok{na.rm =} \ConstantTok{TRUE}
\NormalTok{    ),}
    \AttributeTok{min\_discount =} \FunctionTok{min}\NormalTok{(}
\NormalTok{      .data}\SpecialCharTok{$}\NormalTok{discount\_factor,}
      \AttributeTok{na.rm =} \ConstantTok{TRUE}
\NormalTok{    ),}
    \AttributeTok{max\_discount =} \FunctionTok{max}\NormalTok{(}
\NormalTok{      .data}\SpecialCharTok{$}\NormalTok{discount\_factor,}
      \AttributeTok{na.rm =} \ConstantTok{TRUE}
\NormalTok{    ),}
    \AttributeTok{nonpositive\_discount\_count =} \FunctionTok{sum}\NormalTok{(}
\NormalTok{      .data}\SpecialCharTok{$}\NormalTok{discount\_factor }\SpecialCharTok{\textless{}=} \DecValTok{0}\NormalTok{,}
      \AttributeTok{na.rm =} \ConstantTok{TRUE}
\NormalTok{    ),}
    \AttributeTok{missing\_discount\_count =} \FunctionTok{sum}\NormalTok{(}
      \FunctionTok{is.na}\NormalTok{(.data}\SpecialCharTok{$}\NormalTok{discount\_factor)}
\NormalTok{    ),}
    \AttributeTok{missing\_zero\_rate\_count =} \FunctionTok{sum}\NormalTok{(}
      \FunctionTok{is.na}\NormalTok{(.data}\SpecialCharTok{$}\NormalTok{zero\_rate)}
\NormalTok{    ),}
    \AttributeTok{.groups =} \StringTok{"drop"}
\NormalTok{  )}

\NormalTok{GaussQuant}\SpecialCharTok{::}\FunctionTok{show\_tbl\_GQH}\NormalTok{(}
\NormalTok{  curve\_validation\_tbl,}
  \StringTok{"Curve validation summary"}\NormalTok{,}
  \AttributeTok{n =} \DecValTok{20}
\NormalTok{)}
\CommentTok{\#\textgreater{} \# A tibble: 5 x 8}
\CommentTok{\#\textgreater{}   curve   first\_date last\_date  min\_discount max\_discount nonpositive\_discount\textasciitilde{}1}
\CommentTok{\#\textgreater{}   \textless{}chr\textgreater{}   \textless{}date\textgreater{}     \textless{}date\textgreater{}            \textless{}dbl\textgreater{}        \textless{}dbl\textgreater{}                  \textless{}int\textgreater{}}
\CommentTok{\#\textgreater{} 1 Deposi\textasciitilde{} 2008{-}09{-}18 2038{-}05{-}15        0.260            1                      0}
\CommentTok{\#\textgreater{} 2 Deposi\textasciitilde{} 2008{-}09{-}18 2023{-}09{-}18        0.521            1                      0}
\CommentTok{\#\textgreater{} 3 Deposi\textasciitilde{} 2008{-}09{-}18 2023{-}09{-}18        0.502            1                      0}
\CommentTok{\#\textgreater{} 4 Discou\textasciitilde{} 2022{-}08{-}19 2032{-}08{-}19        0.891            1                      0}
\CommentTok{\#\textgreater{} 5 ZeroCu\textasciitilde{} 2022{-}08{-}19 2032{-}08{-}19        0.891            1                      0}
\CommentTok{\#\textgreater{} \# i abbreviated name: 1: nonpositive\_discount\_count}
\CommentTok{\#\textgreater{} \# i 2 more variables: missing\_discount\_count \textless{}int\textgreater{},}
\CommentTok{\#\textgreater{} \#   missing\_zero\_rate\_count \textless{}int\textgreater{}}
\end{Highlighting}
\end{Shaded}

割引係数が正であることは、裁定のない期間構造を確認するための基本条件である。ただし、正の割引係数だけでは十分ではない。入力商品の再評価、補間区間のフォワード金利、外挿方法、日数計算規則および参照日の設定も確認する必要がある。

\section{2.5
国債先物の基礎}\label{ux56fdux50b5ux5148ux7269ux306eux57faux790e}

株価指数先物では、指数そのものを受け渡すことができないため、満期時には差金決済が行われる。これに対し、国債先物では、満期まで保有された建玉について、受渡適格銘柄に指定された国債を用いた現物受渡しが行われる。現物受渡しを通じて、先物市場と現物国債市場の価格連関が保たれている。

国債先物の取引対象は、実際に発行された特定の国債ではなく、クーポンや償還期限を標準化した標準物である。この標準物は実在しないため、最終決済では、取引所が定めた複数の受渡適格銘柄のいずれかを用いて受渡しを行う。銘柄ごとにクーポン、残存期間および価格が異なるため、その価値の差は変換係数（Conversion
Factor: CF）によって調整される。

受渡銘柄を選択する権利は、国債を受け取る買方ではなく、国債を引き渡す売方にある。したがって売方は、受渡適格銘柄の中から、現物価格だけでなく、変換係数、経過利子、クーポン収入および受渡日までの資金調達費用を考慮し、最も低い経済的費用で受け渡すことのできる銘柄を選択する。この銘柄を最割安銘柄（Cheapest-to-Deliver:
CTD）という。

CTDは、単純に現物単価が最も低い債券ではない。先物価格に変換係数を乗じて求める受渡価値と、現物債を調達して受渡日まで保有する費用とを比較して決定される。

本章では、受渡適格銘柄ごとに現物価格、グロスベーシス、キャリー、ネットベーシスおよびインプライド・レポを計算する。これらの比較を通じて、どのような経済計算によってCTDが選定されるのかを確認し、最後にCTDを基礎とする国債先物のBPVを求める。

\begin{Shaded}
\begin{Highlighting}[]
\NormalTok{knitr}\SpecialCharTok{::}\NormalTok{opts\_chunk}\SpecialCharTok{$}\FunctionTok{set}\NormalTok{(}
  \AttributeTok{collapse =} \ConstantTok{TRUE}\NormalTok{,}
  \AttributeTok{comment =} \StringTok{"\#\textgreater{}"}\NormalTok{,}
  \AttributeTok{warning =} \ConstantTok{FALSE}\NormalTok{,}
  \AttributeTok{message =} \ConstantTok{FALSE}
\NormalTok{)}

\FunctionTok{suppressPackageStartupMessages}\NormalTok{(\{}
  \FunctionTok{library}\NormalTok{(tidyverse)}
  \FunctionTok{library}\NormalTok{(QuantLib)}
  \FunctionTok{library}\NormalTok{(GaussQuant)}
\NormalTok{\})}
\end{Highlighting}
\end{Shaded}

\subsection{2.5.1
計算条件と市場慣行}\label{ux8a08ux7b97ux6761ux4ef6ux3068ux5e02ux5834ux6163ux884c}

本章では、米国国債先物を例として計算する。現物国債には米国政府債カレンダーとActual/Actual（Bond）を使用し、レポ期間にはActual/360を使用する。現物債の受渡日は取引日の1営業日後とする。

米国国債先物の価格は、通常、額面100に対するポイントと32分の1で表示される。たとえば、コード中の
\texttt{109\ +\ 10\ /\ 32} は、相場表示では109-10に相当する。

\begin{Shaded}
\begin{Highlighting}[]
\NormalTok{cal\_us\_gov }\OtherTok{\textless{}{-}}\NormalTok{ QuantLib}\SpecialCharTok{::}\FunctionTok{UnitedStates}\NormalTok{(}\StringTok{"GovernmentBond"}\NormalTok{)}

\NormalTok{dc\_act\_act\_bond }\OtherTok{\textless{}{-}}\NormalTok{ GaussQuant}\SpecialCharTok{::}\FunctionTok{day\_counter\_GQL}\NormalTok{(}\StringTok{"ActualActual\_Bond"}\NormalTok{)}
\NormalTok{dc\_act\_360 }\OtherTok{\textless{}{-}}\NormalTok{ GaussQuant}\SpecialCharTok{::}\FunctionTok{day\_counter\_GQL}\NormalTok{(}\StringTok{"Actual360"}\NormalTok{)}

\NormalTok{cmpd\_compounded }\OtherTok{\textless{}{-}}\NormalTok{ QuantLib}\SpecialCharTok{::}\FunctionTok{Compounding\_Compounded\_get}\NormalTok{()}
\NormalTok{freq\_semiannual }\OtherTok{\textless{}{-}}\NormalTok{ QuantLib}\SpecialCharTok{::}\FunctionTok{Frequency\_Semiannual\_get}\NormalTok{()}

\NormalTok{settlement\_days\_t1 }\OtherTok{\textless{}{-}} \DecValTok{1}\NormalTok{L}
\NormalTok{trade\_date\_chr }\OtherTok{\textless{}{-}} \StringTok{"2023{-}04{-}20"}
\NormalTok{repo\_end\_date\_chr }\OtherTok{\textless{}{-}} \StringTok{"2023{-}07{-}06"}

\NormalTok{trade\_date }\OtherTok{\textless{}{-}}\NormalTok{ GaussQuant}\SpecialCharTok{::}\FunctionTok{date\_GQL}\NormalTok{(trade\_date\_chr)}
\NormalTok{repo\_end\_date\_ql }\OtherTok{\textless{}{-}}\NormalTok{ GaussQuant}\SpecialCharTok{::}\FunctionTok{date\_GQL}\NormalTok{(repo\_end\_date\_chr)}

\NormalTok{GaussQuant}\SpecialCharTok{::}\FunctionTok{set\_eval\_date\_GQL}\NormalTok{(trade\_date)}

\NormalTok{settle\_date }\OtherTok{\textless{}{-}}\NormalTok{ GaussQuant}\SpecialCharTok{::}\FunctionTok{advance\_days\_GQL}\NormalTok{(}
  \AttributeTok{calendar\_obj =}\NormalTok{ cal\_us\_gov,}
  \AttributeTok{date\_obj =}\NormalTok{ trade\_date,}
  \AttributeTok{n\_days =}\NormalTok{ settlement\_days\_t1}
\NormalTok{)}

\NormalTok{settle\_date\_chr }\OtherTok{\textless{}{-}} \FunctionTok{as.character}\NormalTok{(settle\_date)}
\NormalTok{settle\_date\_ql }\OtherTok{\textless{}{-}}\NormalTok{ GaussQuant}\SpecialCharTok{::}\FunctionTok{date\_GQL}\NormalTok{(settle\_date\_chr)}

\NormalTok{setup\_tbl }\OtherTok{\textless{}{-}}\NormalTok{ tibble}\SpecialCharTok{::}\FunctionTok{tribble}\NormalTok{(}
  \SpecialCharTok{\textasciitilde{}}\NormalTok{item, }\SpecialCharTok{\textasciitilde{}}\NormalTok{value,}
  \StringTok{"trade\_date"}\NormalTok{, trade\_date\_chr,}
  \StringTok{"spot\_settlement\_date"}\NormalTok{, settle\_date\_chr,}
  \StringTok{"repo\_end\_date"}\NormalTok{, repo\_end\_date\_chr,}
  \StringTok{"settlement\_days"}\NormalTok{, }\FunctionTok{as.character}\NormalTok{(settlement\_days\_t1),}
  \StringTok{"calendar"}\NormalTok{, }\StringTok{"UnitedStates(GovernmentBond)"}\NormalTok{,}
  \StringTok{"bond\_day\_counter"}\NormalTok{, }\StringTok{"ActualActual(Bond)"}\NormalTok{,}
  \StringTok{"repo\_day\_counter"}\NormalTok{, }\StringTok{"Actual360"}\NormalTok{,}
  \StringTok{"compounding"}\NormalTok{, }\StringTok{"Compounded"}\NormalTok{,}
  \StringTok{"frequency"}\NormalTok{, }\StringTok{"Semiannual"}
\NormalTok{)}

\NormalTok{GaussQuant}\SpecialCharTok{::}\FunctionTok{show\_tbl\_GQH}\NormalTok{(setup\_tbl, }\StringTok{"Treasury futures setup"}\NormalTok{)}
\CommentTok{\#\textgreater{} \# A tibble: 9 x 2}
\CommentTok{\#\textgreater{}   item                 value                       }
\CommentTok{\#\textgreater{}   \textless{}chr\textgreater{}                \textless{}chr\textgreater{}                       }
\CommentTok{\#\textgreater{} 1 trade\_date           2023{-}04{-}20                  }
\CommentTok{\#\textgreater{} 2 spot\_settlement\_date 2023{-}04{-}21                  }
\CommentTok{\#\textgreater{} 3 repo\_end\_date        2023{-}07{-}06                  }
\CommentTok{\#\textgreater{} 4 settlement\_days      1                           }
\CommentTok{\#\textgreater{} 5 calendar             UnitedStates(GovernmentBond)}
\CommentTok{\#\textgreater{} 6 bond\_day\_counter     ActualActual(Bond)          }
\CommentTok{\#\textgreater{} 7 repo\_day\_counter     Actual360                   }
\CommentTok{\#\textgreater{} 8 compounding          Compounded                  }
\CommentTok{\#\textgreater{} 9 frequency            Semiannual}
\end{Highlighting}
\end{Shaded}

\subsection{2.5.2
受渡適格国債}\label{ux53d7ux6e21ux9069ux683cux56fdux50b5}

最初に、受渡適格銘柄の例となる米国国債を作成する。発行日、償還日、クーポンおよび利払スケジュールは、現物価格、経過利子、キャリーおよびBPVを計算する基礎となる。

次のコードでは、年率4.125\%、年2回払いの米国国債を作成し、利払スケジュールを確認する。あわせて、利回り6\%を仮定した場合のクリーン価格を計算する。

\begin{Shaded}
\begin{Highlighting}[]
\NormalTok{bond\_bundle }\OtherTok{\textless{}{-}}\NormalTok{ GaussQuant}\SpecialCharTok{::}\FunctionTok{us\_treasury\_bond\_GQL}\NormalTok{(}
  \AttributeTok{effective\_date =} \StringTok{"2022{-}09{-}30"}\NormalTok{,}
  \AttributeTok{maturity\_date =} \StringTok{"2027{-}09{-}30"}\NormalTok{,}
  \AttributeTok{coupon\_rate\_pct =} \FloatTok{4.125}\NormalTok{,}
  \AttributeTok{face\_amount =} \DecValTok{100}\NormalTok{,}
  \AttributeTok{settlement\_days =}\NormalTok{ settlement\_days\_t1,}
  \AttributeTok{calendar =}\NormalTok{ cal\_us\_gov,}
  \AttributeTok{day\_counter =}\NormalTok{ dc\_act\_act\_bond}
\NormalTok{)}

\NormalTok{bond\_obj }\OtherTok{\textless{}{-}}\NormalTok{ bond\_bundle}\SpecialCharTok{$}\NormalTok{bond}

\NormalTok{GaussQuant}\SpecialCharTok{::}\FunctionTok{show\_tbl\_GQH}\NormalTok{(}
\NormalTok{  GaussQuant}\SpecialCharTok{::}\FunctionTok{schedule\_dates\_GQL}\NormalTok{(bond\_bundle}\SpecialCharTok{$}\NormalTok{schedule),}
  \StringTok{"US Treasury bond schedule"}\NormalTok{,}
  \AttributeTok{n =} \DecValTok{20}
\NormalTok{)}
\CommentTok{\#\textgreater{} \# A tibble: 11 x 1}
\CommentTok{\#\textgreater{}    schedule\_date}
\CommentTok{\#\textgreater{}    \textless{}date\textgreater{}       }
\CommentTok{\#\textgreater{}  1 2022{-}09{-}30   }
\CommentTok{\#\textgreater{}  2 2023{-}03{-}30   }
\CommentTok{\#\textgreater{}  3 2023{-}09{-}30   }
\CommentTok{\#\textgreater{}  4 2024{-}03{-}30   }
\CommentTok{\#\textgreater{}  5 2024{-}09{-}30   }
\CommentTok{\#\textgreater{}  6 2025{-}03{-}30   }
\CommentTok{\#\textgreater{}  7 2025{-}09{-}30   }
\CommentTok{\#\textgreater{}  8 2026{-}03{-}30   }
\CommentTok{\#\textgreater{}  9 2026{-}09{-}30   }
\CommentTok{\#\textgreater{} 10 2027{-}03{-}30   }
\CommentTok{\#\textgreater{} 11 2027{-}09{-}30}

\NormalTok{clean\_price\_at\_six\_percent }\OtherTok{\textless{}{-}}\NormalTok{ GaussQuant}\SpecialCharTok{::}\FunctionTok{bond\_clean\_price\_from\_yield\_safe\_GQL}\NormalTok{(}
  \AttributeTok{bond =}\NormalTok{ bond\_obj,}
  \AttributeTok{yield\_rate =} \DecValTok{6} \SpecialCharTok{/} \DecValTok{100}\NormalTok{,}
  \AttributeTok{day\_counter =}\NormalTok{ dc\_act\_act\_bond,}
  \AttributeTok{compounding =}\NormalTok{ cmpd\_compounded,}
  \AttributeTok{frequency =}\NormalTok{ freq\_semiannual,}
  \AttributeTok{settlement\_date =}\NormalTok{ GaussQuant}\SpecialCharTok{::}\FunctionTok{date\_GQL}\NormalTok{(}\StringTok{"2023{-}07{-}06"}\NormalTok{)}
\NormalTok{)}

\FunctionTok{cat}\NormalTok{(}
  \StringTok{"Yield 6\%, settle 2023{-}07{-}06 clean price:"}\NormalTok{,}
\NormalTok{  GaussQuant}\SpecialCharTok{::}\FunctionTok{fmt\_num\_GQH}\NormalTok{(clean\_price\_at\_six\_percent, }\DecValTok{4}\NormalTok{),}
  \StringTok{"}\SpecialCharTok{\textbackslash{}n}\StringTok{"}
\NormalTok{)}
\CommentTok{\#\textgreater{} Yield 6\%, settle 2023{-}07{-}06 clean price: 93.07}
\end{Highlighting}
\end{Shaded}

\subsection{2.5.3 現物債評価}\label{ux73feux7269ux50b5ux8a55ux4fa1}

CTD分析では、現物債を取得するために必要な資金を把握する必要がある。市場利回りからクリーン価格を求め、受渡日までの経過利子を加えると、実際の決済金額に対応するダーティ価格が得られる。

\[
\text{ダーティ価格}
=
\text{クリーン価格}
+
\text{経過利子}
\]

先物価格には109-10、現物債の市場利回りには3.70\%を用いる。

\begin{Shaded}
\begin{Highlighting}[]
\NormalTok{bond\_yield }\OtherTok{\textless{}{-}} \FloatTok{3.70} \SpecialCharTok{/} \DecValTok{100}
\NormalTok{futures\_price }\OtherTok{\textless{}{-}} \DecValTok{109} \SpecialCharTok{+} \DecValTok{10} \SpecialCharTok{/} \DecValTok{32}

\NormalTok{accrued\_amount }\OtherTok{\textless{}{-}}\NormalTok{ bond\_obj}\SpecialCharTok{$}\FunctionTok{accruedAmount}\NormalTok{(settle\_date\_ql)}
\NormalTok{clean\_price }\OtherTok{\textless{}{-}}\NormalTok{ GaussQuant}\SpecialCharTok{::}\FunctionTok{bond\_clean\_price\_from\_yield\_safe\_GQL}\NormalTok{(}
  \AttributeTok{bond =}\NormalTok{ bond\_obj,}
  \AttributeTok{yield\_rate =}\NormalTok{ bond\_yield,}
  \AttributeTok{day\_counter =}\NormalTok{ dc\_act\_act\_bond,}
  \AttributeTok{compounding =}\NormalTok{ cmpd\_compounded,}
  \AttributeTok{frequency =}\NormalTok{ freq\_semiannual,}
  \AttributeTok{settlement\_date =}\NormalTok{ settle\_date\_ql}
\NormalTok{)}

\NormalTok{dirty\_price }\OtherTok{\textless{}{-}}\NormalTok{ clean\_price }\SpecialCharTok{+}\NormalTok{ accrued\_amount}

\NormalTok{price\_tbl }\OtherTok{\textless{}{-}}\NormalTok{ tibble}\SpecialCharTok{::}\FunctionTok{tibble}\NormalTok{(}
  \AttributeTok{trade\_date =} \FunctionTok{as.Date}\NormalTok{(trade\_date\_chr),}
  \AttributeTok{settlement\_date =} \FunctionTok{as.Date}\NormalTok{(settle\_date\_chr),}
  \AttributeTok{yield =}\NormalTok{ bond\_yield,}
  \AttributeTok{accrued\_amount =}\NormalTok{ accrued\_amount,}
  \AttributeTok{clean\_price =}\NormalTok{ clean\_price,}
  \AttributeTok{dirty\_price =}\NormalTok{ dirty\_price,}
  \AttributeTok{futures\_price =}\NormalTok{ futures\_price}
\NormalTok{)}

\NormalTok{GaussQuant}\SpecialCharTok{::}\FunctionTok{show\_tbl\_GQH}\NormalTok{(price\_tbl, }\StringTok{"Bond price from yield"}\NormalTok{)}
\CommentTok{\#\textgreater{} \# A tibble: 1 x 7}
\CommentTok{\#\textgreater{}   trade\_date settlement\_date yield accrued\_amount clean\_price dirty\_price}
\CommentTok{\#\textgreater{}   \textless{}date\textgreater{}     \textless{}date\textgreater{}          \textless{}dbl\textgreater{}          \textless{}dbl\textgreater{}       \textless{}dbl\textgreater{}       \textless{}dbl\textgreater{}}
\CommentTok{\#\textgreater{} 1 2023{-}04{-}20 2023{-}04{-}21      0.037          0.247        102.        102.}
\CommentTok{\#\textgreater{} \# i 1 more variable: futures\_price \textless{}dbl\textgreater{}}

\FunctionTok{cat}\NormalTok{(}
  \StringTok{"Yield:"}\NormalTok{,}
  \FunctionTok{sprintf}\NormalTok{(}\StringTok{"\%.8f"}\NormalTok{, bond\_yield),}
  \StringTok{"}\SpecialCharTok{\textbackslash{}n}\StringTok{Clean price:"}\NormalTok{,}
  \FunctionTok{sprintf}\NormalTok{(}\StringTok{"\%.8f"}\NormalTok{, clean\_price),}
  \StringTok{"}\SpecialCharTok{\textbackslash{}n}\StringTok{Dirty price:"}\NormalTok{,}
  \FunctionTok{sprintf}\NormalTok{(}\StringTok{"\%.8f"}\NormalTok{, dirty\_price),}
  \StringTok{"}\SpecialCharTok{\textbackslash{}n}\StringTok{Accrued amount:"}\NormalTok{,}
  \FunctionTok{sprintf}\NormalTok{(}\StringTok{"\%.8f"}\NormalTok{, accrued\_amount),}
  \StringTok{"}\SpecialCharTok{\textbackslash{}n}\StringTok{"}
\NormalTok{)}
\CommentTok{\#\textgreater{} Yield: 0.03700000 }
\CommentTok{\#\textgreater{} Clean price: 101.72260614 }
\CommentTok{\#\textgreater{} Dirty price: 101.96920940 }
\CommentTok{\#\textgreater{} Accrued amount: 0.24660326}
\end{Highlighting}
\end{Shaded}

\subsection{2.5.4
キャッシュフロー検証}\label{ux30adux30e3ux30c3ux30b7ux30e5ux30d5ux30edux30fcux691cux8a3c}

次の計算では、各キャッシュフローへ一般的な割引係数を適用し、現在価値の合計と債券価格を比較する。

ただし、Actual/Actual（Bond）は利払参照期間を用いるため、ここで作成する一般的な割引係数は、QuantLibによる債券利回り計算を完全には再現しない。この表は、債券価格が将来キャッシュフローの現在価値から構成されることを確認するためのものであり、差額がゼロになることを目的とした検算ではない。

\begin{Shaded}
\begin{Highlighting}[]
\NormalTok{interest\_rate\_obj }\OtherTok{\textless{}{-}}\NormalTok{ QuantLib}\SpecialCharTok{::}\FunctionTok{InterestRate}\NormalTok{(}
\NormalTok{  bond\_yield,}
\NormalTok{  dc\_act\_act\_bond,}
\NormalTok{  cmpd\_compounded,}
\NormalTok{  freq\_semiannual}
\NormalTok{)}

\NormalTok{settle\_date\_r }\OtherTok{\textless{}{-}} \FunctionTok{as.Date}\NormalTok{(settle\_date\_chr)}

\NormalTok{bond\_cf\_tbl }\OtherTok{\textless{}{-}}\NormalTok{ GaussQuant}\SpecialCharTok{::}\FunctionTok{bond\_cashflow\_table\_GQL}\NormalTok{(bond\_obj) }\SpecialCharTok{|\textgreater{}}
\NormalTok{  dplyr}\SpecialCharTok{::}\FunctionTok{mutate}\NormalTok{(}
    \AttributeTok{pay\_date =} \FunctionTok{as.Date}\NormalTok{(.data}\SpecialCharTok{$}\NormalTok{date),}
    \AttributeTok{discount\_factor =}\NormalTok{ purrr}\SpecialCharTok{::}\FunctionTok{map\_dbl}\NormalTok{(}
      \FunctionTok{as.character}\NormalTok{(.data}\SpecialCharTok{$}\NormalTok{pay\_date),}
      \ControlFlowTok{function}\NormalTok{(x) \{}
        \ControlFlowTok{if}\NormalTok{ (}\FunctionTok{as.Date}\NormalTok{(x) }\SpecialCharTok{\textless{}}\NormalTok{ settle\_date\_r) \{}
          \FunctionTok{return}\NormalTok{(}\ConstantTok{NA\_real\_}\NormalTok{)}
\NormalTok{        \}}

        \FunctionTok{tryCatch}\NormalTok{(}
\NormalTok{          interest\_rate\_obj}\SpecialCharTok{$}\FunctionTok{discountFactor}\NormalTok{(}
\NormalTok{            settle\_date\_ql,}
\NormalTok{            GaussQuant}\SpecialCharTok{::}\FunctionTok{date\_GQL}\NormalTok{(x)}
\NormalTok{          ),}
          \AttributeTok{error =} \ControlFlowTok{function}\NormalTok{(e) }\ConstantTok{NA\_real\_}
\NormalTok{        )}
\NormalTok{      \}}
\NormalTok{    ),}
    \AttributeTok{pv =}\NormalTok{ .data}\SpecialCharTok{$}\NormalTok{amount }\SpecialCharTok{*}\NormalTok{ .data}\SpecialCharTok{$}\NormalTok{discount\_factor}
\NormalTok{  ) }\SpecialCharTok{|\textgreater{}}
\NormalTok{  dplyr}\SpecialCharTok{::}\FunctionTok{select}\NormalTok{(}
\NormalTok{    .data}\SpecialCharTok{$}\NormalTok{pay\_date,}
\NormalTok{    .data}\SpecialCharTok{$}\NormalTok{amount,}
\NormalTok{    .data}\SpecialCharTok{$}\NormalTok{discount\_factor,}
\NormalTok{    .data}\SpecialCharTok{$}\NormalTok{pv}
\NormalTok{  )}

\NormalTok{cashflow\_pv }\OtherTok{\textless{}{-}} \FunctionTok{sum}\NormalTok{(bond\_cf\_tbl}\SpecialCharTok{$}\NormalTok{pv, }\AttributeTok{na.rm =} \ConstantTok{TRUE}\NormalTok{)}

\NormalTok{cashflow\_check\_tbl }\OtherTok{\textless{}{-}}\NormalTok{ tibble}\SpecialCharTok{::}\FunctionTok{tribble}\NormalTok{(}
  \SpecialCharTok{\textasciitilde{}}\NormalTok{metric, }\SpecialCharTok{\textasciitilde{}}\NormalTok{value,}
  \StringTok{"cashflow\_present\_value"}\NormalTok{, cashflow\_pv,}
  \StringTok{"dirty\_price\_from\_yield"}\NormalTok{, dirty\_price,}
  \StringTok{"illustrative\_difference"}\NormalTok{, cashflow\_pv }\SpecialCharTok{{-}}\NormalTok{ dirty\_price}
\NormalTok{)}

\NormalTok{GaussQuant}\SpecialCharTok{::}\FunctionTok{show\_tbl\_GQH}\NormalTok{(bond\_cf\_tbl, }\StringTok{"Bond cashflow present{-}value table"}\NormalTok{, }\AttributeTok{n =} \DecValTok{20}\NormalTok{)}
\CommentTok{\#\textgreater{} \# A tibble: 11 x 4}
\CommentTok{\#\textgreater{}    pay\_date   amount discount\_factor    pv}
\CommentTok{\#\textgreater{}    \textless{}date\textgreater{}      \textless{}dbl\textgreater{}           \textless{}dbl\textgreater{} \textless{}dbl\textgreater{}}
\CommentTok{\#\textgreater{}  1 2023{-}03{-}30   2.06          NA     NA   }
\CommentTok{\#\textgreater{}  2 2023{-}10{-}02   2.06           0.985  2.03}
\CommentTok{\#\textgreater{}  3 2024{-}04{-}01   2.06           0.967  1.99}
\CommentTok{\#\textgreater{}  4 2024{-}09{-}30   2.06           0.949  1.96}
\CommentTok{\#\textgreater{}  5 2025{-}03{-}31   2.06           0.932  1.92}
\CommentTok{\#\textgreater{}  6 2025{-}09{-}30   2.06           0.915  1.89}
\CommentTok{\#\textgreater{}  7 2026{-}03{-}30   2.06           0.899  1.85}
\CommentTok{\#\textgreater{}  8 2026{-}09{-}30   2.06           0.882  1.82}
\CommentTok{\#\textgreater{}  9 2027{-}03{-}30   2.06           0.866  1.79}
\CommentTok{\#\textgreater{} 10 2027{-}09{-}30   2.06           0.851  1.75}
\CommentTok{\#\textgreater{} 11 2027{-}09{-}30 100              0.851 85.1}
\NormalTok{GaussQuant}\SpecialCharTok{::}\FunctionTok{show\_tbl\_GQH}\NormalTok{(}
\NormalTok{  GaussQuant}\SpecialCharTok{::}\FunctionTok{metric\_tbl\_display\_GQH}\NormalTok{(cashflow\_check\_tbl),}
  \StringTok{"Cashflow PV illustration"}
\NormalTok{)}
\CommentTok{\#\textgreater{} \# A tibble: 3 x 2}
\CommentTok{\#\textgreater{}   metric                  value       }
\CommentTok{\#\textgreater{}   \textless{}chr\textgreater{}                   \textless{}chr\textgreater{}       }
\CommentTok{\#\textgreater{} 1 cashflow\_present\_value  102.05828711}
\CommentTok{\#\textgreater{} 2 dirty\_price\_from\_yield  101.96920940}
\CommentTok{\#\textgreater{} 3 illustrative\_difference 0.08907771}
\end{Highlighting}
\end{Shaded}

\subsection{2.5.5 経過利子}\label{ux7d4cux904eux5229ux5b50}

経過利子は、直前利払日から受渡日までに発生したクーポンを日割りして求める。Actual/Actual（Bond）では、経過日数を利払期間の実日数で除してクーポンを按分する。

\[
\text{経過利子}
=
\text{1回分のクーポン}
\times
\frac{\text{経過日数}}{\text{利払期間の日数}}
\]

手計算結果とQuantLibの結果との差を確認する。

\begin{Shaded}
\begin{Highlighting}[]
\NormalTok{previous\_coupon\_date }\OtherTok{\textless{}{-}}\NormalTok{ GaussQuant}\SpecialCharTok{::}\FunctionTok{date\_GQL}\NormalTok{(}\StringTok{"2023{-}03{-}30"}\NormalTok{)}
\NormalTok{next\_coupon\_date }\OtherTok{\textless{}{-}}\NormalTok{ GaussQuant}\SpecialCharTok{::}\FunctionTok{date\_GQL}\NormalTok{(}\StringTok{"2023{-}09{-}30"}\NormalTok{)}

\NormalTok{accrued\_days\_hand }\OtherTok{\textless{}{-}}\NormalTok{ dc\_act\_act\_bond}\SpecialCharTok{$}\FunctionTok{dayCount}\NormalTok{(}
\NormalTok{  previous\_coupon\_date,}
\NormalTok{  settle\_date\_ql}
\NormalTok{)}

\NormalTok{coupon\_period\_days }\OtherTok{\textless{}{-}}\NormalTok{ dc\_act\_act\_bond}\SpecialCharTok{$}\FunctionTok{dayCount}\NormalTok{(}
\NormalTok{  previous\_coupon\_date,}
\NormalTok{  next\_coupon\_date}
\NormalTok{)}

\NormalTok{coupon\_amount\_semiannual }\OtherTok{\textless{}{-}} \FloatTok{4.125} \SpecialCharTok{/} \DecValTok{2}

\NormalTok{accrued\_interest\_hand }\OtherTok{\textless{}{-}}\NormalTok{ coupon\_amount\_semiannual }\SpecialCharTok{*}
\NormalTok{  accrued\_days\_hand }\SpecialCharTok{/}
\NormalTok{  coupon\_period\_days}

\NormalTok{accrued\_tbl }\OtherTok{\textless{}{-}}\NormalTok{ tibble}\SpecialCharTok{::}\FunctionTok{tibble}\NormalTok{(}
  \AttributeTok{previous\_coupon\_date =} \FunctionTok{as.Date}\NormalTok{(GaussQuant}\SpecialCharTok{::}\FunctionTok{iso\_GQL}\NormalTok{(previous\_coupon\_date)),}
  \AttributeTok{settlement\_date =} \FunctionTok{as.Date}\NormalTok{(settle\_date\_chr),}
  \AttributeTok{next\_coupon\_date =} \FunctionTok{as.Date}\NormalTok{(GaussQuant}\SpecialCharTok{::}\FunctionTok{iso\_GQL}\NormalTok{(next\_coupon\_date)),}
  \AttributeTok{accrued\_days =}\NormalTok{ accrued\_days\_hand,}
  \AttributeTok{coupon\_period\_days =}\NormalTok{ coupon\_period\_days,}
  \AttributeTok{coupon\_amount =}\NormalTok{ coupon\_amount\_semiannual,}
  \AttributeTok{accrued\_interest\_hand =}\NormalTok{ accrued\_interest\_hand,}
  \AttributeTok{accrued\_amount\_quantlib =}\NormalTok{ accrued\_amount,}
  \AttributeTok{difference =}\NormalTok{ accrued\_interest\_hand }\SpecialCharTok{{-}}\NormalTok{ accrued\_amount}
\NormalTok{)}

\NormalTok{GaussQuant}\SpecialCharTok{::}\FunctionTok{show\_tbl\_GQH}\NormalTok{(accrued\_tbl, }\StringTok{"Accrued interest hand calculation"}\NormalTok{)}
\CommentTok{\#\textgreater{} \# A tibble: 1 x 9}
\CommentTok{\#\textgreater{}   previous\_coupon\_date settlement\_date next\_coupon\_date accrued\_days}
\CommentTok{\#\textgreater{}   \textless{}date\textgreater{}               \textless{}date\textgreater{}          \textless{}date\textgreater{}                  \textless{}int\textgreater{}}
\CommentTok{\#\textgreater{} 1 2023{-}03{-}30           2023{-}04{-}21      2023{-}09{-}30                 22}
\CommentTok{\#\textgreater{} \# i 5 more variables: coupon\_period\_days \textless{}int\textgreater{}, coupon\_amount \textless{}dbl\textgreater{},}
\CommentTok{\#\textgreater{} \#   accrued\_interest\_hand \textless{}dbl\textgreater{}, accrued\_amount\_quantlib \textless{}dbl\textgreater{},}
\CommentTok{\#\textgreater{} \#   difference \textless{}dbl\textgreater{}}
\end{Highlighting}
\end{Shaded}

\section{2.6
CTD分析と先物BPV}\label{ctdux5206ux6790ux3068ux5148ux7269bpv}

\subsection{2.6.1
受渡適格銘柄と変換係数}\label{ux53d7ux6e21ux9069ux683cux9298ux67c4ux3068ux5909ux63dbux4fc2ux6570}

国債先物では、複数の受渡適格銘柄が存在する。各銘柄はクーポンや残存期間が異なるため、そのままでは同じ先物価格に対して比較できない。そこで、各銘柄へ変換係数を適用し、標準物に対する相対的な受渡価値を調整する。

受渡時に買方が売方へ支払う金額は、概念的には次式で表される。

\[
\text{インボイス価格}
=
\text{先物価格}
\times
\text{変換係数}
+
\text{受渡時経過利子}
\]

ここでいうCFはConversion
Factorの略であり、キャッシュフローを意味するCFとは区別する。

次の表では、3銘柄の発行日、償還日、クーポン、変換係数および市場利回りを設定する。

\begin{Shaded}
\begin{Highlighting}[]
\NormalTok{deliverable\_tbl }\OtherTok{\textless{}{-}}\NormalTok{ tibble}\SpecialCharTok{::}\FunctionTok{tribble}\NormalTok{(}
  \SpecialCharTok{\textasciitilde{}}\NormalTok{issue\_date, }\SpecialCharTok{\textasciitilde{}}\NormalTok{maturity\_date, }\SpecialCharTok{\textasciitilde{}}\NormalTok{coupon\_rate\_pct, }\SpecialCharTok{\textasciitilde{}}\NormalTok{conversion\_factor, }\SpecialCharTok{\textasciitilde{}}\NormalTok{market\_yield\_pct,}
  \StringTok{"2022{-}09{-}30"}\NormalTok{, }\StringTok{"2027{-}09{-}30"}\NormalTok{, }\FloatTok{4.125}\NormalTok{, }\FloatTok{0.9305}\NormalTok{, }\FloatTok{3.70}\NormalTok{,}
  \StringTok{"2022{-}08{-}31"}\NormalTok{, }\StringTok{"2027{-}08{-}31"}\NormalTok{, }\FloatTok{3.125}\NormalTok{, }\FloatTok{0.8953}\NormalTok{, }\FloatTok{3.69}\NormalTok{,}
  \StringTok{"2023{-}01{-}31"}\NormalTok{, }\StringTok{"2028{-}01{-}31"}\NormalTok{, }\FloatTok{3.500}\NormalTok{, }\FloatTok{0.9011}\NormalTok{, }\FloatTok{3.65}
\NormalTok{)}

\NormalTok{GaussQuant}\SpecialCharTok{::}\FunctionTok{show\_tbl\_GQH}\NormalTok{(deliverable\_tbl, }\StringTok{"Treasury futures deliverable basket"}\NormalTok{)}
\CommentTok{\#\textgreater{} \# A tibble: 3 x 5}
\CommentTok{\#\textgreater{}   issue\_date maturity\_date coupon\_rate\_pct conversion\_factor market\_yield\_pct}
\CommentTok{\#\textgreater{}   \textless{}chr\textgreater{}      \textless{}chr\textgreater{}                   \textless{}dbl\textgreater{}             \textless{}dbl\textgreater{}            \textless{}dbl\textgreater{}}
\CommentTok{\#\textgreater{} 1 2022{-}09{-}30 2027{-}09{-}30               4.12             0.930             3.7 }
\CommentTok{\#\textgreater{} 2 2022{-}08{-}31 2027{-}08{-}31               3.12             0.895             3.69}
\CommentTok{\#\textgreater{} 3 2023{-}01{-}31 2028{-}01{-}31               3.5              0.901             3.65}
\end{Highlighting}
\end{Shaded}

\subsection{2.6.2
グロスベーシス}\label{ux30b0ux30edux30b9ux30d9ux30fcux30b7ux30b9}

グロスベーシスは、現物債のクリーン価格と、変換係数で調整した先物価格との差である。

\[
\text{グロスベーシス}
=
\text{現物クリーン価格}
-
\text{先物価格}
\times
\text{変換係数}
\]

グロスベーシスが小さい銘柄ほど受渡しに有利に見えるが、この段階では受渡日までの資金調達費用、クーポン収入および経過利子の変化を考慮していない。

\subsection{2.6.3
キャリーとネットベーシス}\label{ux30adux30e3ux30eaux30fcux3068ux30cdux30c3ux30c8ux30d9ux30fcux30b7ux30b9}

現物債を購入して受渡日まで保有する場合、レポによる資金調達費用が発生する一方、保有期間中にはクーポン収入や経過利子の増加が得られる。これらをまとめたものがキャリーである。

概念的には、ネットベーシスは次式で表される。

\[
\text{ネットベーシス}
=
\text{グロスベーシス}
-
\text{キャリー}
\]

同一の前提条件で比較した場合、一般にネットベーシスが最も小さい銘柄がCTDとなる。

\subsection{2.6.4
インプライド・レポ}\label{ux30a4ux30f3ux30d7ux30e9ux30a4ux30c9ux30ecux30dd}

インプライド・レポは、現物債を購入し、レポで資金調達して受渡日まで保有し、先物契約へ受け渡す取引から得られる収益率である。

ネットベーシスが受渡コストの大きさを表すのに対し、インプライド・レポは同じ取引を収益率として表す。通常、ネットベーシスが最小となる銘柄と、インプライド・レポが最大となる銘柄は一致する。

次のコードでは、各受渡適格銘柄について現物価格、BPV、グロスベーシス、キャリー、ネットベーシスおよびインプライド・レポを計算し、CTDを選定する。

\begin{Shaded}
\begin{Highlighting}[]
\NormalTok{repo\_rate }\OtherTok{\textless{}{-}} \FloatTok{5.10} \SpecialCharTok{/} \DecValTok{100}

\NormalTok{repo\_days }\OtherTok{\textless{}{-}}\NormalTok{ dc\_act\_360}\SpecialCharTok{$}\FunctionTok{dayCount}\NormalTok{(}
\NormalTok{  settle\_date\_ql,}
\NormalTok{  repo\_end\_date\_ql}
\NormalTok{)}

\NormalTok{repo\_year\_fraction }\OtherTok{\textless{}{-}}\NormalTok{ dc\_act\_360}\SpecialCharTok{$}\FunctionTok{yearFraction}\NormalTok{(}
\NormalTok{  settle\_date\_ql,}
\NormalTok{  repo\_end\_date\_ql}
\NormalTok{)}

\NormalTok{ctd\_result }\OtherTok{\textless{}{-}}\NormalTok{ GaussQuant}\SpecialCharTok{::}\FunctionTok{bond\_futures\_ctd\_table\_GQL}\NormalTok{(}
  \AttributeTok{deliverable\_tbl =}\NormalTok{ deliverable\_tbl,}
  \AttributeTok{settlement\_date =}\NormalTok{ settle\_date\_chr,}
  \AttributeTok{futures\_price =}\NormalTok{ futures\_price,}
  \AttributeTok{repo\_end\_date =}\NormalTok{ repo\_end\_date\_chr,}
  \AttributeTok{repo\_rate =}\NormalTok{ repo\_rate,}
  \AttributeTok{settlement\_days =}\NormalTok{ settlement\_days\_t1,}
  \AttributeTok{calendar =}\NormalTok{ cal\_us\_gov,}
  \AttributeTok{day\_counter =}\NormalTok{ dc\_act\_act\_bond,}
  \AttributeTok{compounding =}\NormalTok{ cmpd\_compounded,}
  \AttributeTok{frequency =}\NormalTok{ freq\_semiannual,}
  \AttributeTok{repo\_day\_counter =}\NormalTok{ dc\_act\_360,}
  \AttributeTok{carry\_day\_counter =}\NormalTok{ dc\_act\_360}
\NormalTok{)}

\NormalTok{basket\_tbl }\OtherTok{\textless{}{-}}\NormalTok{ ctd\_result}\SpecialCharTok{$}\NormalTok{basket}

\NormalTok{gross\_basis\_tbl }\OtherTok{\textless{}{-}}\NormalTok{ basket\_tbl }\SpecialCharTok{|\textgreater{}}
\NormalTok{  dplyr}\SpecialCharTok{::}\FunctionTok{select}\NormalTok{(}
\NormalTok{    .data}\SpecialCharTok{$}\NormalTok{issue\_date,}
\NormalTok{    .data}\SpecialCharTok{$}\NormalTok{maturity,}
\NormalTok{    .data}\SpecialCharTok{$}\NormalTok{coupon,}
\NormalTok{    .data}\SpecialCharTok{$}\NormalTok{yield\_pct,}
\NormalTok{    .data}\SpecialCharTok{$}\NormalTok{bpv,}
\NormalTok{    .data}\SpecialCharTok{$}\NormalTok{clean\_price,}
\NormalTok{    .data}\SpecialCharTok{$}\NormalTok{dirty\_price,}
\NormalTok{    .data}\SpecialCharTok{$}\NormalTok{conversion\_factor,}
\NormalTok{    .data}\SpecialCharTok{$}\NormalTok{gross\_basis}
\NormalTok{  )}

\NormalTok{net\_basis\_tbl }\OtherTok{\textless{}{-}}\NormalTok{ basket\_tbl }\SpecialCharTok{|\textgreater{}}
\NormalTok{  dplyr}\SpecialCharTok{::}\FunctionTok{select}\NormalTok{(}
\NormalTok{    .data}\SpecialCharTok{$}\NormalTok{maturity,}
\NormalTok{    .data}\SpecialCharTok{$}\NormalTok{conversion\_factor,}
\NormalTok{    .data}\SpecialCharTok{$}\NormalTok{gross\_basis,}
\NormalTok{    .data}\SpecialCharTok{$}\NormalTok{accrued\_start,}
\NormalTok{    .data}\SpecialCharTok{$}\NormalTok{accrued\_end,}
\NormalTok{    .data}\SpecialCharTok{$}\NormalTok{coupon\_income,}
\NormalTok{    .data}\SpecialCharTok{$}\NormalTok{repo\_cost,}
\NormalTok{    .data}\SpecialCharTok{$}\NormalTok{carry,}
\NormalTok{    .data}\SpecialCharTok{$}\NormalTok{net\_basis,}
\NormalTok{    .data}\SpecialCharTok{$}\NormalTok{forward\_price,}
\NormalTok{    .data}\SpecialCharTok{$}\NormalTok{implied\_repo,}
\NormalTok{    .data}\SpecialCharTok{$}\NormalTok{ctd}
\NormalTok{  )}

\NormalTok{basket\_display\_tbl }\OtherTok{\textless{}{-}}\NormalTok{ basket\_tbl }\SpecialCharTok{|\textgreater{}}
\NormalTok{  dplyr}\SpecialCharTok{::}\FunctionTok{select}\NormalTok{(}\SpecialCharTok{{-}}\NormalTok{dplyr}\SpecialCharTok{::}\FunctionTok{any\_of}\NormalTok{(}\StringTok{"bond\_obj"}\NormalTok{))}

\NormalTok{ctd\_tbl }\OtherTok{\textless{}{-}}\NormalTok{ ctd\_result}\SpecialCharTok{$}\NormalTok{ctd}
\NormalTok{ctd\_display\_tbl }\OtherTok{\textless{}{-}}\NormalTok{ ctd\_tbl }\SpecialCharTok{|\textgreater{}}
\NormalTok{  dplyr}\SpecialCharTok{::}\FunctionTok{select}\NormalTok{(}\SpecialCharTok{{-}}\NormalTok{dplyr}\SpecialCharTok{::}\FunctionTok{any\_of}\NormalTok{(}\StringTok{"bond\_obj"}\NormalTok{))}

\FunctionTok{stopifnot}\NormalTok{(}
  \FunctionTok{nrow}\NormalTok{(ctd\_tbl) }\SpecialCharTok{==} \DecValTok{1}\NormalTok{L,}
  \FunctionTok{isTRUE}\NormalTok{(ctd\_tbl}\SpecialCharTok{$}\NormalTok{ctd[[}\DecValTok{1}\NormalTok{L]])}
\NormalTok{)}

\FunctionTok{cat}\NormalTok{(}
  \StringTok{"Futures price:"}\NormalTok{,}
  \FunctionTok{sprintf}\NormalTok{(}\StringTok{"\%.6f"}\NormalTok{, futures\_price),}
  \StringTok{"}\SpecialCharTok{\textbackslash{}n}\StringTok{Repo rate:"}\NormalTok{,}
  \FunctionTok{sprintf}\NormalTok{(}\StringTok{"\%.6f"}\NormalTok{, repo\_rate),}
  \StringTok{"}\SpecialCharTok{\textbackslash{}n}\StringTok{Repo days:"}\NormalTok{,}
\NormalTok{  repo\_days,}
  \StringTok{"}\SpecialCharTok{\textbackslash{}n}\StringTok{Repo year fraction:"}\NormalTok{,}
  \FunctionTok{sprintf}\NormalTok{(}\StringTok{"\%.8f"}\NormalTok{, repo\_year\_fraction),}
  \StringTok{"}\SpecialCharTok{\textbackslash{}n}\StringTok{"}
\NormalTok{)}
\CommentTok{\#\textgreater{} Futures price: 109.312500 }
\CommentTok{\#\textgreater{} Repo rate: 0.051000 }
\CommentTok{\#\textgreater{} Repo days: 76 }
\CommentTok{\#\textgreater{} Repo year fraction: 0.21111111}

\NormalTok{GaussQuant}\SpecialCharTok{::}\FunctionTok{show\_tbl\_GQH}\NormalTok{(gross\_basis\_tbl, }\StringTok{"Gross basis table"}\NormalTok{, }\AttributeTok{n =} \DecValTok{20}\NormalTok{)}
\CommentTok{\#\textgreater{} \# A tibble: 3 x 9}
\CommentTok{\#\textgreater{}   issue\_date maturity   coupon yield\_pct    bpv clean\_price dirty\_price}
\CommentTok{\#\textgreater{}   \textless{}chr\textgreater{}      \textless{}chr\textgreater{}       \textless{}dbl\textgreater{}     \textless{}dbl\textgreater{}  \textless{}dbl\textgreater{}       \textless{}dbl\textgreater{}       \textless{}dbl\textgreater{}}
\CommentTok{\#\textgreater{} 1 2022{-}09{-}30 2027{-}09{-}30 0.0412      3.7  0.0410       102.        102. }
\CommentTok{\#\textgreater{} 2 2022{-}08{-}31 2027{-}08{-}31 0.0312      3.69 0.0394        97.7        98.2}
\CommentTok{\#\textgreater{} 3 2023{-}01{-}31 2028{-}01{-}31 0.035       3.65 0.0433        99.3       100. }
\CommentTok{\#\textgreater{} \# i 2 more variables: conversion\_factor \textless{}dbl\textgreater{}, gross\_basis \textless{}dbl\textgreater{}}
\NormalTok{GaussQuant}\SpecialCharTok{::}\FunctionTok{show\_tbl\_GQH}\NormalTok{(net\_basis\_tbl, }\StringTok{"Net basis and implied repo table"}\NormalTok{, }\AttributeTok{n =} \DecValTok{20}\NormalTok{)}
\CommentTok{\#\textgreater{} \# A tibble: 3 x 12}
\CommentTok{\#\textgreater{}   maturity conversion\_factor gross\_basis accrued\_start accrued\_end coupon\_income}
\CommentTok{\#\textgreater{}   \textless{}chr\textgreater{}                \textless{}dbl\textgreater{}       \textless{}dbl\textgreater{}         \textless{}dbl\textgreater{}       \textless{}dbl\textgreater{}         \textless{}dbl\textgreater{}}
\CommentTok{\#\textgreater{} 1 2027{-}09\textasciitilde{}             0.930     0.00732         0.247        1.10         0.852}
\CommentTok{\#\textgreater{} 2 2027{-}08\textasciitilde{}             0.895    {-}0.128           0.442        1.09         0.645}
\CommentTok{\#\textgreater{} 3 2028{-}01\textasciitilde{}             0.901     0.841           0.773        1.51         0.735}
\CommentTok{\#\textgreater{} \# i 6 more variables: repo\_cost \textless{}dbl\textgreater{}, carry \textless{}dbl\textgreater{}, net\_basis \textless{}dbl\textgreater{},}
\CommentTok{\#\textgreater{} \#   forward\_price \textless{}dbl\textgreater{}, implied\_repo \textless{}dbl\textgreater{}, ctd \textless{}lgl\textgreater{}}
\NormalTok{GaussQuant}\SpecialCharTok{::}\FunctionTok{show\_tbl\_GQH}\NormalTok{(basket\_display\_tbl, }\StringTok{"Deliverable basket summary"}\NormalTok{, }\AttributeTok{n =} \DecValTok{20}\NormalTok{)}
\CommentTok{\#\textgreater{} \# A tibble: 3 x 18}
\CommentTok{\#\textgreater{}   issue\_date maturity   coupon yield\_pct    bpv clean\_price dirty\_price}
\CommentTok{\#\textgreater{}   \textless{}chr\textgreater{}      \textless{}chr\textgreater{}       \textless{}dbl\textgreater{}     \textless{}dbl\textgreater{}  \textless{}dbl\textgreater{}       \textless{}dbl\textgreater{}       \textless{}dbl\textgreater{}}
\CommentTok{\#\textgreater{} 1 2022{-}09{-}30 2027{-}09{-}30 0.0412      3.7  0.0410       102.        102. }
\CommentTok{\#\textgreater{} 2 2022{-}08{-}31 2027{-}08{-}31 0.0312      3.69 0.0394        97.7        98.2}
\CommentTok{\#\textgreater{} 3 2023{-}01{-}31 2028{-}01{-}31 0.035       3.65 0.0433        99.3       100. }
\CommentTok{\#\textgreater{} \# i 11 more variables: conversion\_factor \textless{}dbl\textgreater{}, gross\_basis \textless{}dbl\textgreater{},}
\CommentTok{\#\textgreater{} \#   accrued\_start \textless{}dbl\textgreater{}, accrued\_end \textless{}dbl\textgreater{}, coupon\_income \textless{}dbl\textgreater{},}
\CommentTok{\#\textgreater{} \#   repo\_cost \textless{}dbl\textgreater{}, carry \textless{}dbl\textgreater{}, net\_basis \textless{}dbl\textgreater{}, forward\_price \textless{}dbl\textgreater{},}
\CommentTok{\#\textgreater{} \#   implied\_repo \textless{}dbl\textgreater{}, ctd \textless{}lgl\textgreater{}}
\NormalTok{GaussQuant}\SpecialCharTok{::}\FunctionTok{show\_tbl\_GQH}\NormalTok{(ctd\_display\_tbl, }\StringTok{"Cheapest{-}to{-}deliver bond"}\NormalTok{, }\AttributeTok{n =} \DecValTok{20}\NormalTok{)}
\CommentTok{\#\textgreater{} \# A tibble: 1 x 18}
\CommentTok{\#\textgreater{}   issue\_date maturity   coupon yield\_pct    bpv clean\_price dirty\_price}
\CommentTok{\#\textgreater{}   \textless{}chr\textgreater{}      \textless{}chr\textgreater{}       \textless{}dbl\textgreater{}     \textless{}dbl\textgreater{}  \textless{}dbl\textgreater{}       \textless{}dbl\textgreater{}       \textless{}dbl\textgreater{}}
\CommentTok{\#\textgreater{} 1 2022{-}09{-}30 2027{-}09{-}30 0.0412       3.7 0.0410        102.        102.}
\CommentTok{\#\textgreater{} \# i 11 more variables: conversion\_factor \textless{}dbl\textgreater{}, gross\_basis \textless{}dbl\textgreater{},}
\CommentTok{\#\textgreater{} \#   accrued\_start \textless{}dbl\textgreater{}, accrued\_end \textless{}dbl\textgreater{}, coupon\_income \textless{}dbl\textgreater{},}
\CommentTok{\#\textgreater{} \#   repo\_cost \textless{}dbl\textgreater{}, carry \textless{}dbl\textgreater{}, net\_basis \textless{}dbl\textgreater{}, forward\_price \textless{}dbl\textgreater{},}
\CommentTok{\#\textgreater{} \#   implied\_repo \textless{}dbl\textgreater{}, ctd \textless{}lgl\textgreater{}}
\end{Highlighting}
\end{Shaded}

結果では、現物価格だけでなく、変換係数とキャリーを反映したネットベーシスを比較する。\texttt{ctd}
がTRUEとなった銘柄が、本章の条件における最割安銘柄である。インプライド・レポもあわせて確認し、CTD選定と整合しているかを見る。

\subsection{2.6.5
CTD調整後の先物BPV}\label{ctdux8abfux6574ux5f8cux306eux5148ux7269bpv}

国債先物の金利感応度は、CTD債券のBPVを変換係数で調整して概算する。

\[
\text{先物BPV}
\approx
\frac{\text{CTD債券のBPV}}
{\text{変換係数}}
\]

この計算は、金利が小幅に変化しても現在のCTDが変わらないことを前提とする。金利水準やイールドカーブの形状が大きく変化すると、別の受渡適格銘柄がCTDとなる可能性がある。そのため、国債先物のリスクは、現在のCTD債券の感応度だけでなく、CTDが切り替わる可能性にも依存する。

\begin{Shaded}
\begin{Highlighting}[]
\NormalTok{ctd\_bond }\OtherTok{\textless{}{-}}\NormalTok{ ctd\_tbl}\SpecialCharTok{$}\NormalTok{bond\_obj[[}\DecValTok{1}\NormalTok{L]]}
\NormalTok{ctd\_conversion\_factor }\OtherTok{\textless{}{-}}\NormalTok{ ctd\_tbl}\SpecialCharTok{$}\NormalTok{conversion\_factor[[}\DecValTok{1}\NormalTok{L]]}
\NormalTok{ctd\_yield }\OtherTok{\textless{}{-}}\NormalTok{ ctd\_tbl}\SpecialCharTok{$}\NormalTok{yield\_pct[[}\DecValTok{1}\NormalTok{L]] }\SpecialCharTok{/} \DecValTok{100}

\NormalTok{ctd\_futures\_bpv }\OtherTok{\textless{}{-}}\NormalTok{ GaussQuant}\SpecialCharTok{::}\FunctionTok{bond\_futures\_bpv\_GQL}\NormalTok{(}
  \AttributeTok{bond =}\NormalTok{ ctd\_bond,}
  \AttributeTok{conversion\_factor =}\NormalTok{ ctd\_conversion\_factor,}
  \AttributeTok{contracts =} \DecValTok{1}\NormalTok{,}
  \AttributeTok{measure =} \StringTok{"dv01"}\NormalTok{,}
  \AttributeTok{ytm =}\NormalTok{ ctd\_yield,}
  \AttributeTok{day\_counter =}\NormalTok{ dc\_act\_act\_bond,}
  \AttributeTok{compounding =}\NormalTok{ cmpd\_compounded,}
  \AttributeTok{frequency =}\NormalTok{ freq\_semiannual}
\NormalTok{)}

\NormalTok{ctd\_summary\_tbl }\OtherTok{\textless{}{-}}\NormalTok{ tibble}\SpecialCharTok{::}\FunctionTok{tibble}\NormalTok{(}
  \AttributeTok{metric =} \FunctionTok{c}\NormalTok{(}
    \StringTok{"ctd\_maturity"}\NormalTok{,}
    \StringTok{"ctd\_coupon"}\NormalTok{,}
    \StringTok{"ctd\_yield\_pct"}\NormalTok{,}
    \StringTok{"conversion\_factor"}\NormalTok{,}
    \StringTok{"gross\_basis"}\NormalTok{,}
    \StringTok{"net\_basis"}\NormalTok{,}
    \StringTok{"implied\_repo"}\NormalTok{,}
    \StringTok{"ctd\_bond\_bpv"}\NormalTok{,}
    \StringTok{"ctd\_futures\_bpv"}
\NormalTok{  ),}
  \AttributeTok{value =} \FunctionTok{list}\NormalTok{(}
    \FunctionTok{as.character}\NormalTok{(ctd\_tbl}\SpecialCharTok{$}\NormalTok{maturity[[}\DecValTok{1}\NormalTok{L]]),}
\NormalTok{    ctd\_tbl}\SpecialCharTok{$}\NormalTok{coupon[[}\DecValTok{1}\NormalTok{L]],}
\NormalTok{    ctd\_tbl}\SpecialCharTok{$}\NormalTok{yield\_pct[[}\DecValTok{1}\NormalTok{L]],}
\NormalTok{    ctd\_conversion\_factor,}
\NormalTok{    ctd\_tbl}\SpecialCharTok{$}\NormalTok{gross\_basis[[}\DecValTok{1}\NormalTok{L]],}
\NormalTok{    ctd\_tbl}\SpecialCharTok{$}\NormalTok{net\_basis[[}\DecValTok{1}\NormalTok{L]],}
\NormalTok{    ctd\_tbl}\SpecialCharTok{$}\NormalTok{implied\_repo[[}\DecValTok{1}\NormalTok{L]],}
\NormalTok{    ctd\_tbl}\SpecialCharTok{$}\NormalTok{bpv[[}\DecValTok{1}\NormalTok{L]],}
\NormalTok{    ctd\_futures\_bpv}
\NormalTok{  )}
\NormalTok{)}

\NormalTok{GaussQuant}\SpecialCharTok{::}\FunctionTok{show\_tbl\_GQH}\NormalTok{(}
\NormalTok{  GaussQuant}\SpecialCharTok{::}\FunctionTok{metric\_tbl\_display\_GQH}\NormalTok{(ctd\_summary\_tbl),}
  \StringTok{"CTD and futures BPV summary"}
\NormalTok{)}
\CommentTok{\#\textgreater{} \# A tibble: 9 x 2}
\CommentTok{\#\textgreater{}   metric            value     }
\CommentTok{\#\textgreater{}   \textless{}chr\textgreater{}             \textless{}chr\textgreater{}     }
\CommentTok{\#\textgreater{} 1 ctd\_maturity      2027{-}09{-}30}
\CommentTok{\#\textgreater{} 2 ctd\_coupon        0.04125000}
\CommentTok{\#\textgreater{} 3 ctd\_yield\_pct     3.70000000}
\CommentTok{\#\textgreater{} 4 conversion\_factor 0.93050000}
\CommentTok{\#\textgreater{} 5 gross\_basis       0.00732489}
\CommentTok{\#\textgreater{} 6 net\_basis         0.25329120}
\CommentTok{\#\textgreater{} 7 implied\_repo      0.03923370}
\CommentTok{\#\textgreater{} 8 ctd\_bond\_bpv      0.04101840}
\CommentTok{\#\textgreater{} 9 ctd\_futures\_bpv   0.04408210}
\end{Highlighting}
\end{Shaded}

\section{2.7 まとめ}\label{ux307eux3068ux3081-6}

本章では、金利期間構造を割引係数、ゼロ金利およびフォワード金利の三つの表現から確認した。これらは同じカーブを表すが、補間方式の影響はフォワード金利に強く現れる。

また、ゼロ金利ノードから直接構築する方法に加え、預金と債券、ならびに預金とスワップを用いたブートストラップを確認した。市場商品から構築したカーブでは、各商品の価格計算、市場慣行および日数計算規則が期間構造へ組み込まれる。

さらに、固定参照日、Implied Term
Structureおよび平行シフトを確認した。カーブを評価へ使用する際には、単にカーブオブジェクトを作成するだけでなく、参照日、補間方式、外挿、入力商品の再現性およびフォワード金利の形状を検証する必要がある。
国債先物では、売方が複数の受渡適格銘柄から実際に引き渡す銘柄を選択する。銘柄間のクーポンや残存期間の差は変換係数によって調整されるが、変換係数だけで経済価値の差が完全に解消されるわけではない。

そのため、CTDの選定では、現物価格と変換係数から求めるグロスベーシスに加え、レポ費用、クーポン収入および経過利子の変化を反映したキャリーを考慮する。ネットベーシスが最小、またはインプライド・レポが最大となる銘柄が、売方にとって最も有利な受渡銘柄となる。

国債先物のBPVは、CTD債券のBPVと変換係数に依存する。ただし、市場環境の変化によってCTDが切り替わる可能性があるため、国債先物のリスク管理では受渡銘柄選択権も考慮する必要がある。

\chapter{第II-3章
OISスワップとマルチカーブ評価}\label{ux7b2cii-3ux7ae0-oisux30b9ux30efux30c3ux30d7ux3068ux30deux30ebux30c1ux30abux30fcux30d6ux8a55ux4fa1}

\section{3.1
IBOR改革とOISの基礎}\label{iborux6539ux9769ux3068oisux306eux57faux790e}

\subsection{3.1.1
IBOR改革と歴史的背景}\label{iborux6539ux9769ux3068ux6b74ux53f2ux7684ux80ccux666f}

2007年から2009年にかけての世界金融危機、とりわけ2008年9月のリーマン・ブラザーズ破綻は、金利デリバティブの評価方法を大きく変化させた。それ以前の市場では、LIBORは信用リスクをほとんど含まない基準金利として扱われ、同一のLIBORカーブから将来の変動金利を予測すると同時に、将来キャッシュフローの現在価値を計算するシングルカーブ評価が広く用いられていた。

しかし、LIBORは銀行間の無担保資金調達金利である。したがって、LIBORには将来の短期金利に対する市場予想だけでなく、借り手銀行の信用リスク、資金流動性リスクおよび借入期間に応じたプレミアムが含まれる。金融危機時にLIBORとOIS金利の差が急拡大したことにより、LIBORを信用リスクのない割引金利として扱うことができないことが明らかになった。

この認識を背景として、デリバティブ市場では中央清算および担保契約の利用が拡大した。担保付取引では、担保に付される翌日物金利に近いOISカーブを用いて将来キャッシュフローを割り引くことが経済的に整合的である。一方、LIBOR、TIBORまたはEURIBORなどを参照する契約の将来変動金利は、それぞれの参照金利に対応するフォワードカーブから予測する必要がある。その結果、将来金利を予測するカーブと現在価値を求めるディスカウントカーブを分離するマルチカーブ評価が定着した。

ここで、「担保付取引への移行」と「翌日物指標の性質」は区別する必要がある。米ドルのSOFRは米国債レポ取引に基づく有担保翌日物金利であるのに対し、日本円のTONAは無担保コール市場の翌日物金利である。いずれも翌日物のリスク・フリー・レートに近い指標として用いられるが、基礎となる市場の担保性は同一ではない。

さらに、LIBORの不正操作および虚偽報告が明らかになり、銀行の判断に大きく依存する金利指標の信頼性が損なわれた。金融安定理事会は2014年、既存IBORの強化と、実取引に基づく代替的なリスク・フリー・レートの育成を提言した。2017年には、英国の金融行為規制機構であるFCAが、2021年末以降も銀行へLIBORの呈示を求め続けることはないとの方針を示した。

代表性のある円、ポンド、ユーロおよびスイスフランLIBORの多くは2021年末に終了し、米ドルLIBORの銀行パネルも2023年6月30日に終了した。経過措置として残された合成LIBORも2024年9月30日に最終公表され、すべてのLIBOR設定が恒久的に終了した。

LIBORに代わる主要な指標として、日本円ではTONA、米ドルではSOFRが採用された。これらの翌日物金利を参照するOISでは、通常、計算期間中の日次金利を後決めで複利計算する。そのため、LIBORのように期間開始時点で適用金利が確定する取引とは異なり、期間全体の金利と利息額が確定するのは期末近くとなる。日次金利の公表時刻、Fixingが利用可能になるまでの時間差、計算終了日から支払日までの事務処理期間は、新たなオペレーショナル・リスクとして管理する必要がある。

以上のIBOR改革によって、金利デリバティブ評価は、単一のLIBORカーブですべてを処理する方法から、翌日物金利を基礎とするOISカーブを割引に用い、参照金利ごとに将来金利を予測する方法へ移行した。

本章では、この歴史的変化を踏まえ、日本円のTONAおよび米ドルのSOFRを参照するOISを取り上げる。まず市場レートからOISカーブを構築し、固定レグと翌日物変動レグを評価する。次に、変動レグを日次金利へ分解し、Fixingの確定前後で評価額がどのように変化するかを確認する。最後に、将来金利を予測するカーブとキャッシュフローを割り引くカーブを分離し、マルチカーブ評価の基本構造を示す。

\begin{Shaded}
\begin{Highlighting}[]
\NormalTok{knitr}\SpecialCharTok{::}\NormalTok{opts\_chunk}\SpecialCharTok{$}\FunctionTok{set}\NormalTok{(}
  \AttributeTok{collapse =} \ConstantTok{TRUE}\NormalTok{,}
  \AttributeTok{comment =} \StringTok{"\#\textgreater{}"}\NormalTok{,}
  \AttributeTok{warning =} \ConstantTok{FALSE}\NormalTok{,}
  \AttributeTok{message =} \ConstantTok{FALSE}
\NormalTok{)}


\FunctionTok{suppressPackageStartupMessages}\NormalTok{(\{}
  \FunctionTok{library}\NormalTok{(tidyverse)}
  \FunctionTok{library}\NormalTok{(ggthemes)}
  \FunctionTok{library}\NormalTok{(QuantLib)}
  \FunctionTok{library}\NormalTok{(GaussQuant)}
  \FunctionTok{library}\NormalTok{(tidyverse)}
  \FunctionTok{library}\NormalTok{(broom)}
\NormalTok{\})}
\end{Highlighting}
\end{Shaded}

\subsection{3.1.2
OISと翌日物金利}\label{oisux3068ux7fccux65e5ux7269ux91d1ux5229}

OIS（Overnight Index
Swap）は、あらかじめ定めた固定金利と、計算期間中の翌日物金利を複利で積み上げた変動金利を交換する金利スワップである。元本そのものは交換せず、通常は固定レグと変動レグの利息差額を決済する。

日次金利を \(r_i\)、その金利が適用される日数を
\(d_i\)、日数計算の年間基準を \(B\)
とすると、計算期間全体の複利収益率は概念的に次式で表される。

\[
R_{\mathrm{comp}}
=
\prod_{i=1}^{n}
\left(
1+r_i\frac{d_i}{B}
\right)-1
\]

元本を \(N\) とすると、変動レグの利息額は次式となる。

\[
I_{\mathrm{overnight}}
=
N R_{\mathrm{comp}}
\]

固定金利を \(K\)、固定レグの年率換算期間を \(\alpha\)
とすると、固定レグの利息額は次式となる。

\[
I_{\mathrm{fixed}}
=
N K \alpha
\]

固定金利を支払うOISの価値は、変動レグの現在価値から固定レグの現在価値を差し引いて求める。

\[
PV_{\mathrm{OIS}}
=
PV_{\mathrm{overnight}}
-
PV_{\mathrm{fixed}}
\]

契約締結時に両レグの現在価値を等しくする固定金利がフェアOISレートである。

以下では、カーブを表形式で確認するための補助関数を定義する。これは計算ロジックを新たに実装するものではなく、QuantLibのカーブから割引係数と連続複利ゼロ金利を取り出す表示用関数である。

\begin{Shaded}
\begin{Highlighting}[]
\NormalTok{curve\_tbl\_local }\OtherTok{\textless{}{-}} \ControlFlowTok{function}\NormalTok{(curve\_obj, }\AttributeTok{n =} \DecValTok{200}\NormalTok{, }\AttributeTok{extrapolate =} \ConstantTok{TRUE}\NormalTok{) \{}
  \ControlFlowTok{if}\NormalTok{ (extrapolate) \{}
\NormalTok{    QuantLib}\SpecialCharTok{::}\FunctionTok{TermStructure\_enableExtrapolation}\NormalTok{(curve\_obj)}
\NormalTok{  \}}

\NormalTok{  reference\_date }\OtherTok{\textless{}{-}} \FunctionTok{as.Date}\NormalTok{(}
\NormalTok{    GaussQuant}\SpecialCharTok{::}\FunctionTok{iso\_GQL}\NormalTok{(curve\_obj}\SpecialCharTok{$}\FunctionTok{referenceDate}\NormalTok{())}
\NormalTok{  )}

\NormalTok{  curve\_times }\OtherTok{\textless{}{-}} \FunctionTok{seq}\NormalTok{(}
    \AttributeTok{from =} \DecValTok{0}\NormalTok{,}
    \AttributeTok{to =}\NormalTok{ curve\_obj}\SpecialCharTok{$}\FunctionTok{maxTime}\NormalTok{(),}
    \AttributeTok{length.out =}\NormalTok{ n}
\NormalTok{  )}

\NormalTok{  tibble}\SpecialCharTok{::}\FunctionTok{tibble}\NormalTok{(}\AttributeTok{time =}\NormalTok{ curve\_times) }\SpecialCharTok{|\textgreater{}}
\NormalTok{    dplyr}\SpecialCharTok{::}\FunctionTok{mutate}\NormalTok{(}
      \AttributeTok{discount\_factor =}\NormalTok{ purrr}\SpecialCharTok{::}\FunctionTok{map\_dbl}\NormalTok{(}
\NormalTok{        time,}
        \ControlFlowTok{function}\NormalTok{(x) curve\_obj}\SpecialCharTok{$}\FunctionTok{discount}\NormalTok{(x)}
\NormalTok{      ),}
      \AttributeTok{zero\_rate =}\NormalTok{ dplyr}\SpecialCharTok{::}\FunctionTok{if\_else}\NormalTok{(}
\NormalTok{        time }\SpecialCharTok{\textgreater{}} \DecValTok{0}\NormalTok{,}
        \SpecialCharTok{{-}}\FunctionTok{log}\NormalTok{(discount\_factor) }\SpecialCharTok{/}\NormalTok{ time,}
        \DecValTok{0}
\NormalTok{      ),}
      \AttributeTok{date =}\NormalTok{ reference\_date }\SpecialCharTok{+} \FunctionTok{round}\NormalTok{(time }\SpecialCharTok{*} \DecValTok{365}\NormalTok{)}
\NormalTok{    ) }\SpecialCharTok{|\textgreater{}}
\NormalTok{    dplyr}\SpecialCharTok{::}\FunctionTok{select}\NormalTok{(}
\NormalTok{      date,}
\NormalTok{      time,}
\NormalTok{      discount\_factor,}
\NormalTok{      zero\_rate}
\NormalTok{    )}
\NormalTok{\}}
\end{Highlighting}
\end{Shaded}

\section{3.2
TONAカーブと短期OIS評価}\label{tonaux30abux30fcux30d6ux3068ux77edux671foisux8a55ux4fa1}

TONAは、日本の無担保コール市場における翌日物金利を基礎とする円のリスク・フリー・レートである。本節では、2022年8月19日の例示的なTONAおよびOIS市場レートから、円OISカーブを構築する。

短期部分には翌日物金利や短期OISを使用し、中長期部分には各満期のOISスワップレートを使用する。ブートストラップでは、各市場商品の理論価格が市場価格と一致するように、満期ごとの割引係数を順次決定する。

\subsection{3.2.1
TONA市場レート}\label{tonaux5e02ux5834ux30ecux30fcux30c8}

\begin{Shaded}
\begin{Highlighting}[]
\NormalTok{trade\_date\_tona }\OtherTok{\textless{}{-}} \StringTok{"2022{-}08{-}19"}

\NormalTok{GaussQuant}\SpecialCharTok{::}\FunctionTok{set\_eval\_date\_GQL}\NormalTok{(trade\_date\_tona)}

\NormalTok{tona\_quotes }\OtherTok{\textless{}{-}}\NormalTok{ tibble}\SpecialCharTok{::}\FunctionTok{tribble}\NormalTok{(}
  \SpecialCharTok{\textasciitilde{}}\NormalTok{as\_of\_date, }\SpecialCharTok{\textasciitilde{}}\NormalTok{currency, }\SpecialCharTok{\textasciitilde{}}\NormalTok{instrument, }\SpecialCharTok{\textasciitilde{}}\NormalTok{kind, }\SpecialCharTok{\textasciitilde{}}\NormalTok{tenor, }\SpecialCharTok{\textasciitilde{}}\NormalTok{rate,}
  \FunctionTok{as.Date}\NormalTok{(trade\_date\_tona), }\StringTok{"JPY"}\NormalTok{, }\StringTok{"TONA"}\NormalTok{, }\StringTok{"depo"}\NormalTok{, }\StringTok{"1D"}\NormalTok{,  }\SpecialCharTok{{-}}\FloatTok{0.0000900}\NormalTok{,}
  \FunctionTok{as.Date}\NormalTok{(trade\_date\_tona), }\StringTok{"JPY"}\NormalTok{, }\StringTok{"TONA"}\NormalTok{, }\StringTok{"ois"}\NormalTok{,  }\StringTok{"1W"}\NormalTok{,  }\SpecialCharTok{{-}}\FloatTok{0.0001261}\NormalTok{,}
  \FunctionTok{as.Date}\NormalTok{(trade\_date\_tona), }\StringTok{"JPY"}\NormalTok{, }\StringTok{"TONA"}\NormalTok{, }\StringTok{"ois"}\NormalTok{,  }\StringTok{"2W"}\NormalTok{,  }\SpecialCharTok{{-}}\FloatTok{0.0001469}\NormalTok{,}
  \FunctionTok{as.Date}\NormalTok{(trade\_date\_tona), }\StringTok{"JPY"}\NormalTok{, }\StringTok{"TONA"}\NormalTok{, }\StringTok{"ois"}\NormalTok{,  }\StringTok{"1M"}\NormalTok{,  }\SpecialCharTok{{-}}\FloatTok{0.0001807}\NormalTok{,}
  \FunctionTok{as.Date}\NormalTok{(trade\_date\_tona), }\StringTok{"JPY"}\NormalTok{, }\StringTok{"TONA"}\NormalTok{, }\StringTok{"ois"}\NormalTok{,  }\StringTok{"3M"}\NormalTok{,  }\SpecialCharTok{{-}}\FloatTok{0.0001919}\NormalTok{,}
  \FunctionTok{as.Date}\NormalTok{(trade\_date\_tona), }\StringTok{"JPY"}\NormalTok{, }\StringTok{"TONA"}\NormalTok{, }\StringTok{"ois"}\NormalTok{,  }\StringTok{"6M"}\NormalTok{,  }\SpecialCharTok{{-}}\FloatTok{0.0001043}\NormalTok{,}
  \FunctionTok{as.Date}\NormalTok{(trade\_date\_tona), }\StringTok{"JPY"}\NormalTok{, }\StringTok{"TONA"}\NormalTok{, }\StringTok{"ois"}\NormalTok{,  }\StringTok{"9M"}\NormalTok{,   }\FloatTok{0.0000022}\NormalTok{,}
  \FunctionTok{as.Date}\NormalTok{(trade\_date\_tona), }\StringTok{"JPY"}\NormalTok{, }\StringTok{"TONA"}\NormalTok{, }\StringTok{"ois"}\NormalTok{,  }\StringTok{"12M"}\NormalTok{,  }\FloatTok{0.0001250}\NormalTok{,}
  \FunctionTok{as.Date}\NormalTok{(trade\_date\_tona), }\StringTok{"JPY"}\NormalTok{, }\StringTok{"TONA"}\NormalTok{, }\StringTok{"ois"}\NormalTok{,  }\StringTok{"18M"}\NormalTok{,  }\FloatTok{0.0003125}\NormalTok{,}
  \FunctionTok{as.Date}\NormalTok{(trade\_date\_tona), }\StringTok{"JPY"}\NormalTok{, }\StringTok{"TONA"}\NormalTok{, }\StringTok{"ois"}\NormalTok{,  }\StringTok{"2Y"}\NormalTok{,   }\FloatTok{0.0004875}\NormalTok{,}
  \FunctionTok{as.Date}\NormalTok{(trade\_date\_tona), }\StringTok{"JPY"}\NormalTok{, }\StringTok{"TONA"}\NormalTok{, }\StringTok{"ois"}\NormalTok{,  }\StringTok{"3Y"}\NormalTok{,   }\FloatTok{0.0007375}\NormalTok{,}
  \FunctionTok{as.Date}\NormalTok{(trade\_date\_tona), }\StringTok{"JPY"}\NormalTok{, }\StringTok{"TONA"}\NormalTok{, }\StringTok{"ois"}\NormalTok{,  }\StringTok{"4Y"}\NormalTok{,   }\FloatTok{0.0009479}\NormalTok{,}
  \FunctionTok{as.Date}\NormalTok{(trade\_date\_tona), }\StringTok{"JPY"}\NormalTok{, }\StringTok{"TONA"}\NormalTok{, }\StringTok{"ois"}\NormalTok{,  }\StringTok{"5Y"}\NormalTok{,   }\FloatTok{0.0011854}\NormalTok{,}
  \FunctionTok{as.Date}\NormalTok{(trade\_date\_tona), }\StringTok{"JPY"}\NormalTok{, }\StringTok{"TONA"}\NormalTok{, }\StringTok{"ois"}\NormalTok{,  }\StringTok{"7Y"}\NormalTok{,   }\FloatTok{0.0019146}
\NormalTok{)}

\NormalTok{GaussQuant}\SpecialCharTok{::}\FunctionTok{show\_tbl\_GQH}\NormalTok{(}
\NormalTok{  tona\_quotes,}
  \StringTok{"TONA market quotes"}\NormalTok{,}
  \AttributeTok{n =} \DecValTok{20}
\NormalTok{)}
\CommentTok{\#\textgreater{} \# A tibble: 14 x 6}
\CommentTok{\#\textgreater{}    as\_of\_date currency instrument kind  tenor       rate}
\CommentTok{\#\textgreater{}    \textless{}date\textgreater{}     \textless{}chr\textgreater{}    \textless{}chr\textgreater{}      \textless{}chr\textgreater{} \textless{}chr\textgreater{}      \textless{}dbl\textgreater{}}
\CommentTok{\#\textgreater{}  1 2022{-}08{-}19 JPY      TONA       depo  1D    {-}0.00009  }
\CommentTok{\#\textgreater{}  2 2022{-}08{-}19 JPY      TONA       ois   1W    {-}0.000126 }
\CommentTok{\#\textgreater{}  3 2022{-}08{-}19 JPY      TONA       ois   2W    {-}0.000147 }
\CommentTok{\#\textgreater{}  4 2022{-}08{-}19 JPY      TONA       ois   1M    {-}0.000181 }
\CommentTok{\#\textgreater{}  5 2022{-}08{-}19 JPY      TONA       ois   3M    {-}0.000192 }
\CommentTok{\#\textgreater{}  6 2022{-}08{-}19 JPY      TONA       ois   6M    {-}0.000104 }
\CommentTok{\#\textgreater{}  7 2022{-}08{-}19 JPY      TONA       ois   9M     0.0000022}
\CommentTok{\#\textgreater{}  8 2022{-}08{-}19 JPY      TONA       ois   12M    0.000125 }
\CommentTok{\#\textgreater{}  9 2022{-}08{-}19 JPY      TONA       ois   18M    0.000312 }
\CommentTok{\#\textgreater{} 10 2022{-}08{-}19 JPY      TONA       ois   2Y     0.000488 }
\CommentTok{\#\textgreater{} 11 2022{-}08{-}19 JPY      TONA       ois   3Y     0.000738 }
\CommentTok{\#\textgreater{} 12 2022{-}08{-}19 JPY      TONA       ois   4Y     0.000948 }
\CommentTok{\#\textgreater{} 13 2022{-}08{-}19 JPY      TONA       ois   5Y     0.00119  }
\CommentTok{\#\textgreater{} 14 2022{-}08{-}19 JPY      TONA       ois   7Y     0.00191}
\end{Highlighting}
\end{Shaded}

\subsection{3.2.2 TONA
OISカーブの構築}\label{tona-oisux30abux30fcux30d6ux306eux69cbux7bc9}

\begin{Shaded}
\begin{Highlighting}[]
\NormalTok{tona\_curve\_envs }\OtherTok{\textless{}{-}}\NormalTok{ GaussQuant}\SpecialCharTok{::}\FunctionTok{ir\_build\_ois\_curve\_envs\_GQH}\NormalTok{(}
  \AttributeTok{quotes =}\NormalTok{ tona\_quotes,}
  \AttributeTok{trade\_date =}\NormalTok{ trade\_date\_tona,}
  \AttributeTok{verbose =} \ConstantTok{TRUE}
\NormalTok{)}

\NormalTok{tona\_curve\_env }\OtherTok{\textless{}{-}}\NormalTok{ tona\_curve\_envs[[}\StringTok{"JPY::TONA"}\NormalTok{]]}
\NormalTok{tona\_curve }\OtherTok{\textless{}{-}}\NormalTok{ tona\_curve\_env}\SpecialCharTok{$}\NormalTok{curve}

\NormalTok{tona\_curve\_display\_tbl }\OtherTok{\textless{}{-}} \FunctionTok{curve\_tbl\_local}\NormalTok{(}
  \AttributeTok{curve\_obj =}\NormalTok{ tona\_curve,}
  \AttributeTok{n =} \DecValTok{200}
\NormalTok{)}

\NormalTok{GaussQuant}\SpecialCharTok{::}\FunctionTok{show\_tbl\_GQH}\NormalTok{(}
\NormalTok{  tona\_curve\_display\_tbl,}
  \StringTok{"TONA OIS curve"}\NormalTok{,}
  \AttributeTok{n =} \DecValTok{12}
\NormalTok{)}
\CommentTok{\#\textgreater{} \# A tibble: 12 x 4}
\CommentTok{\#\textgreater{}    date         time discount\_factor zero\_rate}
\CommentTok{\#\textgreater{}    \textless{}date\textgreater{}      \textless{}dbl\textgreater{}           \textless{}dbl\textgreater{}     \textless{}dbl\textgreater{}}
\CommentTok{\#\textgreater{}  1 2022{-}08{-}23 0                 1     0       }
\CommentTok{\#\textgreater{}  2 2022{-}09{-}05 0.0352            1.00 {-}0.000145}
\CommentTok{\#\textgreater{}  3 2022{-}09{-}18 0.0704            1.00 {-}0.000173}
\CommentTok{\#\textgreater{}  4 2022{-}10{-}01 0.106             1.00 {-}0.000183}
\CommentTok{\#\textgreater{}  5 2022{-}10{-}13 0.141             1.00 {-}0.000187}
\CommentTok{\#\textgreater{}  6 2022{-}10{-}26 0.176             1.00 {-}0.000189}
\CommentTok{\#\textgreater{}  7 2022{-}11{-}08 0.211             1.00 {-}0.000191}
\CommentTok{\#\textgreater{}  8 2022{-}11{-}21 0.246             1.00 {-}0.000192}
\CommentTok{\#\textgreater{}  9 2022{-}12{-}04 0.282             1.00 {-}0.000175}
\CommentTok{\#\textgreater{} 10 2022{-}12{-}17 0.317             1.00 {-}0.000157}
\CommentTok{\#\textgreater{} 11 2022{-}12{-}29 0.352             1.00 {-}0.000143}
\CommentTok{\#\textgreater{} 12 2023{-}01{-}11 0.387             1.00 {-}0.000132}

\NormalTok{tona\_curve\_check\_tbl }\OtherTok{\textless{}{-}}\NormalTok{ tibble}\SpecialCharTok{::}\FunctionTok{tibble}\NormalTok{(}
  \AttributeTok{reference\_date =}\NormalTok{ GaussQuant}\SpecialCharTok{::}\FunctionTok{iso\_GQL}\NormalTok{(}
\NormalTok{    tona\_curve}\SpecialCharTok{$}\FunctionTok{referenceDate}\NormalTok{()}
\NormalTok{  ),}
  \AttributeTok{discount\_2022\_08\_23 =}\NormalTok{ tona\_curve}\SpecialCharTok{$}\FunctionTok{discount}\NormalTok{(}
\NormalTok{    GaussQuant}\SpecialCharTok{::}\FunctionTok{date\_GQL}\NormalTok{(}\StringTok{"2022{-}08{-}23"}\NormalTok{)}
\NormalTok{  )}
\NormalTok{)}

\NormalTok{GaussQuant}\SpecialCharTok{::}\FunctionTok{show\_tbl\_GQH}\NormalTok{(}
\NormalTok{  tona\_curve\_check\_tbl,}
  \StringTok{"TONA curve checks"}\NormalTok{,}
  \AttributeTok{n =} \DecValTok{20}
\NormalTok{)}
\CommentTok{\#\textgreater{} \# A tibble: 1 x 2}
\CommentTok{\#\textgreater{}   reference\_date discount\_2022\_08\_23}
\CommentTok{\#\textgreater{}   \textless{}chr\textgreater{}                        \textless{}dbl\textgreater{}}
\CommentTok{\#\textgreater{} 1 2022{-}08{-}23                       1}
\end{Highlighting}
\end{Shaded}

割引係数 \(P(0,T)\) は、将来時点 \(T\)
の1円を評価日時点の価値へ換算する係数である。満期が長くなるほど、通常は割引係数が低下する。ゼロ金利は、各満期の割引係数を単一の金利として表したものである。

\subsection{3.2.3
JPY短期OISの評価}\label{jpyux77edux671foisux306eux8a55ux4fa1}

次に、TONAカーブを用いて短期円OISを評価する。評価対象は、2022年8月23日から8月30日までの7日間、元本1,000万円、固定金利マイナス0.015\%の固定金利支払いOISである。

OISの固定レグと翌日物レグを確認し、NPVとフェアレートを表示する。

\begin{Shaded}
\begin{Highlighting}[]
\NormalTok{jpy\_ois\_short }\OtherTok{\textless{}{-}}\NormalTok{ GaussQuant}\SpecialCharTok{::}\FunctionTok{trade\_ois\_swap\_py\_GQH}\NormalTok{(}
  \AttributeTok{curve\_env =}\NormalTok{ tona\_curve\_env,}
  \AttributeTok{effective =} \StringTok{"2022{-}08{-}23"}\NormalTok{,}
  \AttributeTok{maturity =} \StringTok{"2022{-}08{-}30"}\NormalTok{,}
  \AttributeTok{notional =} \DecValTok{10000000}\NormalTok{,}
  \AttributeTok{fixed\_rate =} \SpecialCharTok{{-}}\FloatTok{0.00015}\NormalTok{,}
  \AttributeTok{pay\_receive =} \StringTok{"pay"}\NormalTok{,}
  \AttributeTok{verbose =} \ConstantTok{TRUE}
\NormalTok{)}

\NormalTok{GaussQuant}\SpecialCharTok{::}\FunctionTok{ir\_show\_swap\_summary\_GQH}\NormalTok{(}
  \AttributeTok{trade\_env =}\NormalTok{ jpy\_ois\_short,}
  \AttributeTok{title =} \StringTok{"JPY short OIS summary"}
\NormalTok{)}

\NormalTok{jpy\_ois\_fixed\_cf\_tbl }\OtherTok{\textless{}{-}}\NormalTok{ GaussQuant}\SpecialCharTok{::}\FunctionTok{swap\_fixed\_leg\_table\_GQL}\NormalTok{(}
\NormalTok{  jpy\_ois\_short}\SpecialCharTok{$}\NormalTok{swap}
\NormalTok{)}

\NormalTok{jpy\_ois\_overnight\_cf\_tbl }\OtherTok{\textless{}{-}}\NormalTok{ GaussQuant}\SpecialCharTok{::}\FunctionTok{ois\_overnight\_leg\_table\_GQL}\NormalTok{(}
\NormalTok{  jpy\_ois\_short}\SpecialCharTok{$}\NormalTok{swap}
\NormalTok{)}

\NormalTok{GaussQuant}\SpecialCharTok{::}\FunctionTok{show\_tbl\_GQH}\NormalTok{(}
\NormalTok{  jpy\_ois\_fixed\_cf\_tbl,}
  \StringTok{"TONA fixed leg cashflows"}\NormalTok{,}
  \AttributeTok{n =} \DecValTok{10}
\NormalTok{)}
\CommentTok{\#\textgreater{} \# A tibble: 1 x 2}
\CommentTok{\#\textgreater{}   date       amount}
\CommentTok{\#\textgreater{}   \textless{}chr\textgreater{}       \textless{}dbl\textgreater{}}
\CommentTok{\#\textgreater{} 1 2022{-}08{-}30  {-}28.8}

\NormalTok{GaussQuant}\SpecialCharTok{::}\FunctionTok{show\_tbl\_GQH}\NormalTok{(}
\NormalTok{  jpy\_ois\_overnight\_cf\_tbl,}
  \StringTok{"TONA overnight leg cashflows"}\NormalTok{,}
  \AttributeTok{n =} \DecValTok{10}
\NormalTok{)}
\CommentTok{\#\textgreater{} \# A tibble: 1 x 2}
\CommentTok{\#\textgreater{}   date       amount}
\CommentTok{\#\textgreater{}   \textless{}chr\textgreater{}       \textless{}dbl\textgreater{}}
\CommentTok{\#\textgreater{} 1 2022{-}08{-}30  {-}24.2}
\end{Highlighting}
\end{Shaded}

\subsection{3.2.4
フェアレートの手計算}\label{ux30d5ux30a7ux30a2ux30ecux30fcux30c8ux306eux624bux8a08ux7b97}

単一期間のOISでは、開始日の割引係数を \(P(0,T_0)\)、終了日の割引係数を
\(P(0,T_1)\)、固定レグの支払日の割引係数を
\(P(0,T_p)\)、固定レグの年率換算期間を \(\alpha\)
とすると、フェアレートは概念的に次式で求められる。

\[
K_{\mathrm{fair}}
=
\frac{
P(0,T_0)-P(0,T_1)
}{
\alpha P(0,T_p)
}
\]

\begin{Shaded}
\begin{Highlighting}[]
\NormalTok{jpy\_effective\_df }\OtherTok{\textless{}{-}}\NormalTok{ tona\_curve}\SpecialCharTok{$}\FunctionTok{discount}\NormalTok{(}
\NormalTok{  GaussQuant}\SpecialCharTok{::}\FunctionTok{date\_GQL}\NormalTok{(}\StringTok{"2022{-}08{-}23"}\NormalTok{)}
\NormalTok{)}

\NormalTok{jpy\_maturity\_df }\OtherTok{\textless{}{-}}\NormalTok{ tona\_curve}\SpecialCharTok{$}\FunctionTok{discount}\NormalTok{(}
\NormalTok{  GaussQuant}\SpecialCharTok{::}\FunctionTok{date\_GQL}\NormalTok{(}\StringTok{"2022{-}08{-}30"}\NormalTok{)}
\NormalTok{)}

\NormalTok{jpy\_payment\_df }\OtherTok{\textless{}{-}}\NormalTok{ tona\_curve}\SpecialCharTok{$}\FunctionTok{discount}\NormalTok{(}
\NormalTok{  GaussQuant}\SpecialCharTok{::}\FunctionTok{date\_GQL}\NormalTok{(}\StringTok{"2022{-}09{-}01"}\NormalTok{)}
\NormalTok{)}

\NormalTok{jpy\_annuity }\OtherTok{\textless{}{-}}\NormalTok{ (}\DecValTok{7} \SpecialCharTok{/} \DecValTok{365}\NormalTok{) }\SpecialCharTok{*}\NormalTok{ jpy\_payment\_df}

\NormalTok{jpy\_swap\_rate\_hand }\OtherTok{\textless{}{-}}\NormalTok{ (}
\NormalTok{  jpy\_effective\_df }\SpecialCharTok{{-}}\NormalTok{ jpy\_maturity\_df}
\NormalTok{) }\SpecialCharTok{/}\NormalTok{ jpy\_annuity}

\NormalTok{jpy\_fair\_rate\_check\_tbl }\OtherTok{\textless{}{-}}\NormalTok{ tibble}\SpecialCharTok{::}\FunctionTok{tibble}\NormalTok{(}
  \AttributeTok{effective\_discount\_factor =}\NormalTok{ jpy\_effective\_df,}
  \AttributeTok{maturity\_discount\_factor =}\NormalTok{ jpy\_maturity\_df,}
  \AttributeTok{payment\_discount\_factor =}\NormalTok{ jpy\_payment\_df,}
  \AttributeTok{annuity =}\NormalTok{ jpy\_annuity,}
  \AttributeTok{fair\_swap\_rate\_hand =}\NormalTok{ jpy\_swap\_rate\_hand}
\NormalTok{)}

\NormalTok{GaussQuant}\SpecialCharTok{::}\FunctionTok{show\_tbl\_GQH}\NormalTok{(}
\NormalTok{  jpy\_fair\_rate\_check\_tbl,}
  \StringTok{"TONA fair{-}rate hand check"}\NormalTok{,}
  \AttributeTok{n =} \DecValTok{20}
\NormalTok{)}
\CommentTok{\#\textgreater{} \# A tibble: 1 x 5}
\CommentTok{\#\textgreater{}   effective\_discount\_fac\textasciitilde{}1 maturity\_discount\_fa\textasciitilde{}2 payment\_discount\_fac\textasciitilde{}3 annuity}
\CommentTok{\#\textgreater{}                      \textless{}dbl\textgreater{}                  \textless{}dbl\textgreater{}                  \textless{}dbl\textgreater{}   \textless{}dbl\textgreater{}}
\CommentTok{\#\textgreater{} 1                        1                   1.00                   1.00  0.0192}
\CommentTok{\#\textgreater{} \# i abbreviated names: 1: effective\_discount\_factor,}
\CommentTok{\#\textgreater{} \#   2: maturity\_discount\_factor, 3: payment\_discount\_factor}
\CommentTok{\#\textgreater{} \# i 1 more variable: fair\_swap\_rate\_hand \textless{}dbl\textgreater{}}
\end{Highlighting}
\end{Shaded}

\section{3.3
日次複利とFixing}\label{ux65e5ux6b21ux8907ux5229ux3068fixing}

\subsection{3.3.1
OIS変動レグの日次複利}\label{oisux5909ux52d5ux30ecux30b0ux306eux65e5ux6b21ux8907ux5229}

OISの変動レグでは、計算期間中の翌日物金利を日ごとに複利で積み上げる。休日をまたぐ金利は、次の営業日まで複数日分に適用されるため、各行の適用日数を確認する必要がある。

次のコードでは、変動レグを日次単位に分解し、割引係数から日次フォワード金利を求める。

\begin{Shaded}
\begin{Highlighting}[]
\NormalTok{jpy\_ois\_daily }\OtherTok{\textless{}{-}}\NormalTok{ GaussQuant}\SpecialCharTok{::}\FunctionTok{trade\_ois\_swap\_py\_GQH}\NormalTok{(}
  \AttributeTok{curve\_env =}\NormalTok{ tona\_curve\_env,}
  \AttributeTok{effective =} \StringTok{"2022{-}08{-}23"}\NormalTok{,}
  \AttributeTok{maturity =} \StringTok{"2022{-}08{-}30"}\NormalTok{,}
  \AttributeTok{notional =} \DecValTok{10000000}\NormalTok{,}
  \AttributeTok{fixed\_rate =} \SpecialCharTok{{-}}\FloatTok{0.00015}\NormalTok{,}
  \AttributeTok{pay\_receive =} \StringTok{"pay"}\NormalTok{,}
  \AttributeTok{fixed\_schedule\_tenor =} \StringTok{"1Y"}\NormalTok{,}
  \AttributeTok{floating\_schedule\_tenor =} \StringTok{"1D"}\NormalTok{,}
  \AttributeTok{pay\_lag =} \DecValTok{0}\NormalTok{,}
  \AttributeTok{verbose =} \ConstantTok{TRUE}
\NormalTok{)}

\NormalTok{GaussQuant}\SpecialCharTok{::}\FunctionTok{set\_eval\_date\_GQL}\NormalTok{(trade\_date\_tona)}

\NormalTok{jpy\_daily\_float\_cf\_tbl }\OtherTok{\textless{}{-}}\NormalTok{ GaussQuant}\SpecialCharTok{::}\FunctionTok{ois\_daily\_forward\_table\_GQH}\NormalTok{(}
  \AttributeTok{swap =}\NormalTok{ jpy\_ois\_daily}\SpecialCharTok{$}\NormalTok{swap,}
  \AttributeTok{curve =}\NormalTok{ tona\_curve,}
  \AttributeTok{index =}\NormalTok{ jpy\_ois\_daily}\SpecialCharTok{$}\NormalTok{index,}
  \AttributeTok{eval\_date =}\NormalTok{ trade\_date\_tona,}
  \AttributeTok{notional =} \DecValTok{10000000}
\NormalTok{)}

\NormalTok{GaussQuant}\SpecialCharTok{::}\FunctionTok{show\_tbl\_GQH}\NormalTok{(}
\NormalTok{  jpy\_daily\_float\_cf\_tbl,}
  \StringTok{"TONA daily forward leg"}\NormalTok{,}
  \AttributeTok{n =} \DecValTok{20}
\NormalTok{)}
\CommentTok{\#\textgreater{} \# A tibble: 5 x 9}
\CommentTok{\#\textgreater{}   fixing\_date next\_date   days fixing\_before\_eval fixing\_on\_eval fixing\_value}
\CommentTok{\#\textgreater{}   \textless{}chr\textgreater{}       \textless{}chr\textgreater{}      \textless{}int\textgreater{} \textless{}lgl\textgreater{}              \textless{}lgl\textgreater{}                 \textless{}dbl\textgreater{}}
\CommentTok{\#\textgreater{} 1 2022{-}08{-}23  2022{-}08{-}24     1 FALSE              FALSE            {-}0.0000900}
\CommentTok{\#\textgreater{} 2 2022{-}08{-}24  2022{-}08{-}25     1 FALSE              FALSE            {-}0.000132 }
\CommentTok{\#\textgreater{} 3 2022{-}08{-}25  2022{-}08{-}26     1 FALSE              FALSE            {-}0.000132 }
\CommentTok{\#\textgreater{} 4 2022{-}08{-}26  2022{-}08{-}29     3 FALSE              FALSE            {-}0.000132 }
\CommentTok{\#\textgreater{} 5 2022{-}08{-}29  2022{-}08{-}30     1 FALSE              FALSE            {-}0.000132 }
\CommentTok{\#\textgreater{} \# i 3 more variables: forward\_rate\_from\_df \textless{}dbl\textgreater{}, applied\_rate \textless{}dbl\textgreater{},}
\CommentTok{\#\textgreater{} \#   amount \textless{}dbl\textgreater{}}
\end{Highlighting}
\end{Shaded}

各日の単利フォワード金利は、開始日の割引係数を
\(P(0,t_i)\)、終了日の割引係数を \(P(0,t_{i+1})\)、年率換算期間を
\(\delta_i\) とすると、次式で求められる。

\[
f_i
=
\frac{
P(0,t_i)/P(0,t_{i+1})-1
}{
\delta_i
}
\]

日次フォワード金利を用いて複利収益率と利息額を確認する。

\begin{Shaded}
\begin{Highlighting}[]
\NormalTok{daily\_compounding\_result }\OtherTok{\textless{}{-}} \ControlFlowTok{if}\NormalTok{ (}
  \FunctionTok{nrow}\NormalTok{(jpy\_daily\_float\_cf\_tbl) }\SpecialCharTok{\textgreater{}=} \DecValTok{2} \SpecialCharTok{\&\&}
  \FunctionTok{all}\NormalTok{(}\FunctionTok{c}\NormalTok{(}\StringTok{"days"}\NormalTok{, }\StringTok{"rate"}\NormalTok{) }\SpecialCharTok{\%in\%} \FunctionTok{names}\NormalTok{(jpy\_daily\_float\_cf\_tbl))}
\NormalTok{) \{}
\NormalTok{  daily\_compounded\_return }\OtherTok{\textless{}{-}}\NormalTok{ (}
    \DecValTok{1} \SpecialCharTok{+}
\NormalTok{      jpy\_daily\_float\_cf\_tbl}\SpecialCharTok{$}\NormalTok{days[}\SpecialCharTok{{-}}\DecValTok{1}\NormalTok{] }\SpecialCharTok{/}
      \DecValTok{365} \SpecialCharTok{*}
\NormalTok{      jpy\_daily\_float\_cf\_tbl}\SpecialCharTok{$}\NormalTok{rate[}\SpecialCharTok{{-}}\DecValTok{1}\NormalTok{]}
\NormalTok{  ) }\SpecialCharTok{|\textgreater{}}
    \FunctionTok{prod}\NormalTok{() }\SpecialCharTok{{-}}
    \DecValTok{1}

\NormalTok{  tibble}\SpecialCharTok{::}\FunctionTok{tibble}\NormalTok{(}
    \AttributeTok{compounded\_return =}\NormalTok{ daily\_compounded\_return,}
    \AttributeTok{floating\_interest =}\NormalTok{ daily\_compounded\_return }\SpecialCharTok{*} \DecValTok{10000000}\NormalTok{,}
    \AttributeTok{equivalent\_annual\_rate =}\NormalTok{ daily\_compounded\_return }\SpecialCharTok{*} \DecValTok{365} \SpecialCharTok{/} \DecValTok{7}
\NormalTok{  )}
\NormalTok{\} }\ControlFlowTok{else}\NormalTok{ \{}
\NormalTok{  tibble}\SpecialCharTok{::}\FunctionTok{tibble}\NormalTok{(}
    \AttributeTok{compounded\_return =} \ConstantTok{NA\_real\_}\NormalTok{,}
    \AttributeTok{floating\_interest =} \ConstantTok{NA\_real\_}\NormalTok{,}
    \AttributeTok{equivalent\_annual\_rate =} \ConstantTok{NA\_real\_}
\NormalTok{  )}
\NormalTok{\}}

\NormalTok{GaussQuant}\SpecialCharTok{::}\FunctionTok{show\_tbl\_GQH}\NormalTok{(}
\NormalTok{  daily\_compounding\_result,}
  \StringTok{"TONA daily compounding result"}\NormalTok{,}
  \AttributeTok{n =} \DecValTok{20}
\NormalTok{)}
\CommentTok{\#\textgreater{} \# A tibble: 1 x 3}
\CommentTok{\#\textgreater{}   compounded\_return floating\_interest equivalent\_annual\_rate}
\CommentTok{\#\textgreater{}               \textless{}dbl\textgreater{}             \textless{}dbl\textgreater{}                  \textless{}dbl\textgreater{}}
\CommentTok{\#\textgreater{} 1                NA                NA                     NA}
\end{Highlighting}
\end{Shaded}

\subsection{3.3.2
Fixingの確定と再評価}\label{fixingux306eux78baux5b9aux3068ux518dux8a55ux4fa1}

評価日より前の翌日物金利には、実際に公表されたFixingを使用する。評価日より後の金利は、OISカーブから推定したフォワード金利を使用する。

したがって、評価途中のOIS変動レグは、次の二つから構成される。

\begin{itemize}
\tightlist
\item
  評価日までに確定した日次Fixingの複利部分
\item
  評価日より後のフォワード金利による予測部分
\end{itemize}

後決め複利では、期間終了まで変動レグ全体の利息額が確定しない。Fixingの公表時刻と支払処理の間に必要な時間を確保するため、実務ではLookback、Observation
Shift、LockoutまたはPayment Delayなどの規則が用いられることがある。

次の例では、2022年8月23日および24日のTONA
Fixingを登録し、2022年8月25日時点で再評価する。

\begin{Shaded}
\begin{Highlighting}[]
\NormalTok{jpy\_fixings\_tbl }\OtherTok{\textless{}{-}}\NormalTok{ tibble}\SpecialCharTok{::}\FunctionTok{tribble}\NormalTok{(}
  \SpecialCharTok{\textasciitilde{}}\NormalTok{date, }\SpecialCharTok{\textasciitilde{}}\NormalTok{currency, }\SpecialCharTok{\textasciitilde{}}\NormalTok{instrument, }\SpecialCharTok{\textasciitilde{}}\NormalTok{fixing,}
  \StringTok{"2022{-}08{-}23"}\NormalTok{, }\StringTok{"JPY"}\NormalTok{, }\StringTok{"TONA"}\NormalTok{, }\SpecialCharTok{{-}}\FloatTok{0.000090}\NormalTok{,}
  \StringTok{"2022{-}08{-}24"}\NormalTok{, }\StringTok{"JPY"}\NormalTok{, }\StringTok{"TONA"}\NormalTok{, }\SpecialCharTok{{-}}\FloatTok{0.000085}
\NormalTok{)}

\NormalTok{GaussQuant}\SpecialCharTok{::}\FunctionTok{set\_eval\_date\_GQL}\NormalTok{(trade\_date\_tona)}

\NormalTok{jpy\_overnight\_before\_tbl }\OtherTok{\textless{}{-}}\NormalTok{ GaussQuant}\SpecialCharTok{::}\FunctionTok{ois\_overnight\_leg\_table\_GQL}\NormalTok{(}
\NormalTok{  jpy\_ois\_short}\SpecialCharTok{$}\NormalTok{swap}
\NormalTok{)}

\NormalTok{GaussQuant}\SpecialCharTok{::}\FunctionTok{ir\_apply\_fixings\_GQH}\NormalTok{(}
  \AttributeTok{index\_obj =}\NormalTok{ jpy\_ois\_short}\SpecialCharTok{$}\NormalTok{index,}
  \AttributeTok{fixings\_tbl =}\NormalTok{ jpy\_fixings\_tbl,}
  \AttributeTok{currency =} \StringTok{"JPY"}\NormalTok{,}
  \AttributeTok{instrument =} \StringTok{"TONA"}
\NormalTok{)}

\NormalTok{GaussQuant}\SpecialCharTok{::}\FunctionTok{set\_eval\_date\_GQL}\NormalTok{(}\StringTok{"2022{-}08{-}25"}\NormalTok{)}

\NormalTok{jpy\_overnight\_after\_tbl }\OtherTok{\textless{}{-}}\NormalTok{ GaussQuant}\SpecialCharTok{::}\FunctionTok{ois\_overnight\_leg\_table\_GQL}\NormalTok{(}
\NormalTok{  jpy\_ois\_short}\SpecialCharTok{$}\NormalTok{swap}
\NormalTok{)}

\NormalTok{GaussQuant}\SpecialCharTok{::}\FunctionTok{ir\_show\_fixing\_before\_after\_GQH}\NormalTok{(}
  \AttributeTok{before\_tbl =}\NormalTok{ jpy\_overnight\_before\_tbl,}
  \AttributeTok{after\_tbl =}\NormalTok{ jpy\_overnight\_after\_tbl,}
  \AttributeTok{title =} \StringTok{"TONA fixing before / after"}\NormalTok{,}
  \AttributeTok{n =} \DecValTok{10}
\NormalTok{)}
\CommentTok{\#\textgreater{} \# A tibble: 1 x 3}
\CommentTok{\#\textgreater{}   row\_id amount\_before amount\_after}
\CommentTok{\#\textgreater{}    \textless{}int\textgreater{}         \textless{}dbl\textgreater{}        \textless{}dbl\textgreater{}}
\CommentTok{\#\textgreater{} 1      1         {-}24.2        {-}24.2}

\NormalTok{GaussQuant}\SpecialCharTok{::}\FunctionTok{ir\_show\_swap\_summary\_GQH}\NormalTok{(}
  \AttributeTok{trade\_env =}\NormalTok{ jpy\_ois\_short,}
  \AttributeTok{title =} \StringTok{"TONA short OIS summary after fixings"}
\NormalTok{)}

\NormalTok{jpy\_fixing\_diag\_tbl }\OtherTok{\textless{}{-}}\NormalTok{ GaussQuant}\SpecialCharTok{::}\FunctionTok{ir\_fixing\_table\_GQH}\NormalTok{(}
  \AttributeTok{swap =}\NormalTok{ jpy\_ois\_daily}\SpecialCharTok{$}\NormalTok{swap,}
  \AttributeTok{curve =}\NormalTok{ tona\_curve,}
  \AttributeTok{index =}\NormalTok{ jpy\_ois\_short}\SpecialCharTok{$}\NormalTok{index}
\NormalTok{)}

\NormalTok{GaussQuant}\SpecialCharTok{::}\FunctionTok{show\_tbl\_GQH}\NormalTok{(}
\NormalTok{  jpy\_fixing\_diag\_tbl,}
  \StringTok{"TONA daily fixing decomposition"}\NormalTok{,}
  \AttributeTok{n =} \DecValTok{20}
\NormalTok{)}
\CommentTok{\#\textgreater{} \# A tibble: 1 x 7}
\CommentTok{\#\textgreater{}       i fixing\_date fixing\_value pay\_date   amount    df    pv}
\CommentTok{\#\textgreater{}   \textless{}int\textgreater{} \textless{}chr\textgreater{}              \textless{}dbl\textgreater{} \textless{}chr\textgreater{}       \textless{}dbl\textgreater{} \textless{}dbl\textgreater{} \textless{}dbl\textgreater{}}
\CommentTok{\#\textgreater{} 1     1 2022{-}08{-}29     {-}0.000142 2022{-}08{-}30  {-}24.2  1.00 {-}24.2}
\end{Highlighting}
\end{Shaded}

\section{3.4
SOFRとマルチカーブ評価}\label{sofrux3068ux30deux30ebux30c1ux30abux30fcux30d6ux8a55ux4fa1}

SOFRは、米国債を担保とする翌日物レポ取引に基づく米ドルのリスク・フリー・レートである。TONAが無担保翌日物市場に基づくのに対し、SOFRは有担保取引に基づく点が異なる。

OIS評価の基本構造はTONAと同じであり、通貨、休日カレンダー、日数計算規則、市場クォートおよび翌日物指数が異なる。

\subsection{3.4.1
SOFR市場レートとOISカーブ}\label{sofrux5e02ux5834ux30ecux30fcux30c8ux3068oisux30abux30fcux30d6}

\begin{Shaded}
\begin{Highlighting}[]
\NormalTok{trade\_date\_sofr }\OtherTok{\textless{}{-}} \StringTok{"2023{-}09{-}26"}

\NormalTok{GaussQuant}\SpecialCharTok{::}\FunctionTok{set\_eval\_date\_GQL}\NormalTok{(trade\_date\_sofr)}

\NormalTok{sofr\_quotes }\OtherTok{\textless{}{-}}\NormalTok{ tibble}\SpecialCharTok{::}\FunctionTok{tribble}\NormalTok{(}
  \SpecialCharTok{\textasciitilde{}}\NormalTok{as\_of\_date, }\SpecialCharTok{\textasciitilde{}}\NormalTok{currency, }\SpecialCharTok{\textasciitilde{}}\NormalTok{instrument, }\SpecialCharTok{\textasciitilde{}}\NormalTok{kind, }\SpecialCharTok{\textasciitilde{}}\NormalTok{tenor, }\SpecialCharTok{\textasciitilde{}}\NormalTok{rate,}
  \FunctionTok{as.Date}\NormalTok{(trade\_date\_sofr), }\StringTok{"USD"}\NormalTok{, }\StringTok{"SOFR"}\NormalTok{, }\StringTok{"depo"}\NormalTok{, }\StringTok{"1D"}\NormalTok{, }\FloatTok{0.0531}\NormalTok{,}
  \FunctionTok{as.Date}\NormalTok{(trade\_date\_sofr), }\StringTok{"USD"}\NormalTok{, }\StringTok{"SOFR"}\NormalTok{, }\StringTok{"ois"}\NormalTok{,  }\StringTok{"1M"}\NormalTok{, }\FloatTok{0.0532}\NormalTok{,}
  \FunctionTok{as.Date}\NormalTok{(trade\_date\_sofr), }\StringTok{"USD"}\NormalTok{, }\StringTok{"SOFR"}\NormalTok{, }\StringTok{"ois"}\NormalTok{,  }\StringTok{"3M"}\NormalTok{, }\FloatTok{0.0538}\NormalTok{,}
  \FunctionTok{as.Date}\NormalTok{(trade\_date\_sofr), }\StringTok{"USD"}\NormalTok{, }\StringTok{"SOFR"}\NormalTok{, }\StringTok{"ois"}\NormalTok{,  }\StringTok{"6M"}\NormalTok{, }\FloatTok{0.0546}\NormalTok{,}
  \FunctionTok{as.Date}\NormalTok{(trade\_date\_sofr), }\StringTok{"USD"}\NormalTok{, }\StringTok{"SOFR"}\NormalTok{, }\StringTok{"ois"}\NormalTok{,  }\StringTok{"1Y"}\NormalTok{, }\FloatTok{0.0545}\NormalTok{,}
  \FunctionTok{as.Date}\NormalTok{(trade\_date\_sofr), }\StringTok{"USD"}\NormalTok{, }\StringTok{"SOFR"}\NormalTok{, }\StringTok{"ois"}\NormalTok{,  }\StringTok{"2Y"}\NormalTok{, }\FloatTok{0.0501}\NormalTok{,}
  \FunctionTok{as.Date}\NormalTok{(trade\_date\_sofr), }\StringTok{"USD"}\NormalTok{, }\StringTok{"SOFR"}\NormalTok{, }\StringTok{"ois"}\NormalTok{,  }\StringTok{"3Y"}\NormalTok{, }\FloatTok{0.0467}
\NormalTok{)}

\NormalTok{sofr\_curve\_envs }\OtherTok{\textless{}{-}}\NormalTok{ GaussQuant}\SpecialCharTok{::}\FunctionTok{ir\_build\_ois\_curve\_envs\_GQH}\NormalTok{(}
  \AttributeTok{quotes =}\NormalTok{ sofr\_quotes,}
  \AttributeTok{trade\_date =}\NormalTok{ trade\_date\_sofr,}
  \AttributeTok{verbose =} \ConstantTok{TRUE}
\NormalTok{)}

\NormalTok{sofr\_curve\_env }\OtherTok{\textless{}{-}}\NormalTok{ sofr\_curve\_envs[[}\StringTok{"USD::SOFR"}\NormalTok{]]}
\NormalTok{sofr\_curve }\OtherTok{\textless{}{-}}\NormalTok{ sofr\_curve\_env}\SpecialCharTok{$}\NormalTok{curve}

\NormalTok{GaussQuant}\SpecialCharTok{::}\FunctionTok{show\_tbl\_GQH}\NormalTok{(}
  \FunctionTok{curve\_tbl\_local}\NormalTok{(sofr\_curve),}
  \StringTok{"SOFR OIS curve"}\NormalTok{,}
  \AttributeTok{n =} \DecValTok{12}
\NormalTok{)}
\CommentTok{\#\textgreater{} \# A tibble: 12 x 4}
\CommentTok{\#\textgreater{}    date         time discount\_factor zero\_rate}
\CommentTok{\#\textgreater{}    \textless{}date\textgreater{}      \textless{}dbl\textgreater{}           \textless{}dbl\textgreater{}     \textless{}dbl\textgreater{}}
\CommentTok{\#\textgreater{}  1 2023{-}09{-}28 0                1        0     }
\CommentTok{\#\textgreater{}  2 2023{-}10{-}04 0.0153           0.999    0.0531}
\CommentTok{\#\textgreater{}  3 2023{-}10{-}09 0.0306           0.998    0.0531}
\CommentTok{\#\textgreater{}  4 2023{-}10{-}15 0.0459           0.998    0.0531}
\CommentTok{\#\textgreater{}  5 2023{-}10{-}20 0.0612           0.997    0.0531}
\CommentTok{\#\textgreater{}  6 2023{-}10{-}26 0.0765           0.996    0.0531}
\CommentTok{\#\textgreater{}  7 2023{-}11{-}01 0.0918           0.995    0.0531}
\CommentTok{\#\textgreater{}  8 2023{-}11{-}06 0.107            0.994    0.0532}
\CommentTok{\#\textgreater{}  9 2023{-}11{-}12 0.122            0.994    0.0532}
\CommentTok{\#\textgreater{} 10 2023{-}11{-}17 0.138            0.993    0.0533}
\CommentTok{\#\textgreater{} 11 2023{-}11{-}23 0.153            0.992    0.0533}
\CommentTok{\#\textgreater{} 12 2023{-}11{-}28 0.168            0.991    0.0533}
\end{Highlighting}
\end{Shaded}

\subsection{3.4.2 SOFR OISの評価}\label{sofr-oisux306eux8a55ux4fa1}

\begin{Shaded}
\begin{Highlighting}[]
\NormalTok{sofr\_ois\_trade }\OtherTok{\textless{}{-}}\NormalTok{ GaussQuant}\SpecialCharTok{::}\FunctionTok{trade\_ois\_swap\_py\_GQH}\NormalTok{(}
  \AttributeTok{curve\_env =}\NormalTok{ sofr\_curve\_env,}
  \AttributeTok{effective =} \StringTok{"2023{-}09{-}28"}\NormalTok{,}
  \AttributeTok{maturity =} \StringTok{"2025{-}09{-}28"}\NormalTok{,}
  \AttributeTok{notional =} \DecValTok{10000000}\NormalTok{,}
  \AttributeTok{fixed\_rate =} \FloatTok{0.05}\NormalTok{,}
  \AttributeTok{pay\_receive =} \StringTok{"pay"}\NormalTok{,}
  \AttributeTok{verbose =} \ConstantTok{TRUE}
\NormalTok{)}

\NormalTok{GaussQuant}\SpecialCharTok{::}\FunctionTok{ir\_show\_swap\_summary\_GQH}\NormalTok{(}
  \AttributeTok{trade\_env =}\NormalTok{ sofr\_ois\_trade,}
  \AttributeTok{title =} \StringTok{"SOFR OIS summary"}
\NormalTok{)}

\NormalTok{sofr\_overnight\_cf\_tbl }\OtherTok{\textless{}{-}}\NormalTok{ GaussQuant}\SpecialCharTok{::}\FunctionTok{ois\_overnight\_leg\_table\_GQL}\NormalTok{(}
\NormalTok{  sofr\_ois\_trade}\SpecialCharTok{$}\NormalTok{swap}
\NormalTok{)}

\NormalTok{GaussQuant}\SpecialCharTok{::}\FunctionTok{show\_tbl\_GQH}\NormalTok{(}
\NormalTok{  sofr\_overnight\_cf\_tbl,}
  \StringTok{"SOFR overnight leg cashflows"}\NormalTok{,}
  \AttributeTok{n =} \DecValTok{20}
\NormalTok{)}
\CommentTok{\#\textgreater{} \# A tibble: 2 x 2}
\CommentTok{\#\textgreater{}   date        amount}
\CommentTok{\#\textgreater{}   \textless{}chr\textgreater{}        \textless{}dbl\textgreater{}}
\CommentTok{\#\textgreater{} 1 2024{-}09{-}30 557111.}
\CommentTok{\#\textgreater{} 2 2025{-}09{-}29 459522.}
\end{Highlighting}
\end{Shaded}

\subsection{3.4.3
マルチカーブ評価}\label{ux30deux30ebux30c1ux30abux30fcux30d6ux8a55ux4fa1}

金融危機以前のシングルカーブ評価では、同一の金利カーブを用いて、将来変動金利の予測とキャッシュフローの割引を行っていた。

マルチカーブ評価では、二つの役割を分離する。

\begin{itemize}
\item
  \textbf{フォワードカーブ}
  参照金利の将来Fixingまたは期間フォワード金利を予測する。
\item
  \textbf{ディスカウントカーブ}
  担保契約に対応する金利を用いて、将来キャッシュフローを現在価値へ割り引く。
\end{itemize}

将来変動金利を \(F_j\)、年率換算期間を \(\alpha_j\)、割引係数を
\(D_j\)、元本を \(N\)
とすると、変動レグの現在価値は概念的に次式で表される。

\[
PV_{\mathrm{floating}}
=
N
\sum_{j=1}^{m}
\alpha_j F_j D_j
\]

ここで、\(F_j\) はフォワードカーブから求め、\(D_j\)
はディスカウントカーブから求める。

以下では、マルチカーブの効果だけを明確にするため、EURIBOR
6か月物を参照する一般的なバニラスワップを用いる。同じフォワードカーブを使用しながら、割引カーブを同一にした場合と分離した場合のNPVを比較する。

\begin{Shaded}
\begin{Highlighting}[]
\NormalTok{GaussQuant}\SpecialCharTok{::}\FunctionTok{set\_eval\_date\_GQL}\NormalTok{(}\StringTok{"2023{-}09{-}26"}\NormalTok{)}

\NormalTok{multi\_curve\_dates }\OtherTok{\textless{}{-}} \FunctionTok{c}\NormalTok{(}
  \StringTok{"2023{-}09{-}26"}\NormalTok{,}
  \StringTok{"2024{-}09{-}26"}\NormalTok{,}
  \StringTok{"2025{-}09{-}26"}\NormalTok{,}
  \StringTok{"2026{-}09{-}28"}\NormalTok{,}
  \StringTok{"2028{-}09{-}26"}
\NormalTok{)}

\NormalTok{forecast\_zero\_rates }\OtherTok{\textless{}{-}} \FunctionTok{c}\NormalTok{(}
  \FloatTok{0.0450}\NormalTok{,}
  \FloatTok{0.0470}\NormalTok{,}
  \FloatTok{0.0460}\NormalTok{,}
  \FloatTok{0.0445}\NormalTok{,}
  \FloatTok{0.0420}
\NormalTok{)}

\NormalTok{discount\_zero\_rates }\OtherTok{\textless{}{-}} \FunctionTok{c}\NormalTok{(}
  \FloatTok{0.0400}\NormalTok{,}
  \FloatTok{0.0415}\NormalTok{,}
  \FloatTok{0.0410}\NormalTok{,}
  \FloatTok{0.0400}\NormalTok{,}
  \FloatTok{0.0380}
\NormalTok{)}

\NormalTok{forecast\_curve }\OtherTok{\textless{}{-}}\NormalTok{ GaussQuant}\SpecialCharTok{::}\FunctionTok{build\_zero\_curve\_GQL}\NormalTok{(}
  \AttributeTok{date\_chr =}\NormalTok{ multi\_curve\_dates,}
  \AttributeTok{zero\_rates =}\NormalTok{ forecast\_zero\_rates,}
  \AttributeTok{day\_counter =} \StringTok{"Actual365Fixed"}
\NormalTok{)}

\NormalTok{discount\_curve }\OtherTok{\textless{}{-}}\NormalTok{ GaussQuant}\SpecialCharTok{::}\FunctionTok{build\_zero\_curve\_GQL}\NormalTok{(}
  \AttributeTok{date\_chr =}\NormalTok{ multi\_curve\_dates,}
  \AttributeTok{zero\_rates =}\NormalTok{ discount\_zero\_rates,}
  \AttributeTok{day\_counter =} \StringTok{"Actual365Fixed"}
\NormalTok{)}

\NormalTok{QuantLib}\SpecialCharTok{::}\FunctionTok{TermStructure\_enableExtrapolation}\NormalTok{(forecast\_curve)}
\CommentTok{\#\textgreater{} NULL}
\NormalTok{QuantLib}\SpecialCharTok{::}\FunctionTok{TermStructure\_enableExtrapolation}\NormalTok{(discount\_curve)}
\CommentTok{\#\textgreater{} NULL}

\NormalTok{forecast\_handle }\OtherTok{\textless{}{-}}\NormalTok{ GaussQuant}\SpecialCharTok{::}\FunctionTok{yield\_curve\_handle\_GQL}\NormalTok{(}
\NormalTok{  forecast\_curve}
\NormalTok{)}

\NormalTok{discount\_handle }\OtherTok{\textless{}{-}}\NormalTok{ GaussQuant}\SpecialCharTok{::}\FunctionTok{yield\_curve\_handle\_GQL}\NormalTok{(}
\NormalTok{  discount\_curve}
\NormalTok{)}

\NormalTok{multi\_curve\_comparison\_tbl }\OtherTok{\textless{}{-}}\NormalTok{ tibble}\SpecialCharTok{::}\FunctionTok{tibble}\NormalTok{(}
  \AttributeTok{date =} \FunctionTok{as.Date}\NormalTok{(multi\_curve\_dates)}
\NormalTok{) }\SpecialCharTok{|\textgreater{}}
\NormalTok{  dplyr}\SpecialCharTok{::}\FunctionTok{mutate}\NormalTok{(}
    \AttributeTok{forecast\_discount\_factor =}\NormalTok{ purrr}\SpecialCharTok{::}\FunctionTok{map\_dbl}\NormalTok{(}
      \FunctionTok{as.character}\NormalTok{(date),}
      \ControlFlowTok{function}\NormalTok{(x) \{}
\NormalTok{        forecast\_curve}\SpecialCharTok{$}\FunctionTok{discount}\NormalTok{(}
\NormalTok{          GaussQuant}\SpecialCharTok{::}\FunctionTok{date\_GQL}\NormalTok{(x)}
\NormalTok{        )}
\NormalTok{      \}}
\NormalTok{    ),}
    \AttributeTok{ois\_discount\_factor =}\NormalTok{ purrr}\SpecialCharTok{::}\FunctionTok{map\_dbl}\NormalTok{(}
      \FunctionTok{as.character}\NormalTok{(date),}
      \ControlFlowTok{function}\NormalTok{(x) \{}
\NormalTok{        discount\_curve}\SpecialCharTok{$}\FunctionTok{discount}\NormalTok{(}
\NormalTok{          GaussQuant}\SpecialCharTok{::}\FunctionTok{date\_GQL}\NormalTok{(x)}
\NormalTok{        )}
\NormalTok{      \}}
\NormalTok{    )}
\NormalTok{  )}

\NormalTok{GaussQuant}\SpecialCharTok{::}\FunctionTok{show\_tbl\_GQH}\NormalTok{(}
\NormalTok{  multi\_curve\_comparison\_tbl,}
  \StringTok{"Forecast and discount curves"}\NormalTok{,}
  \AttributeTok{n =} \DecValTok{20}
\NormalTok{)}
\CommentTok{\#\textgreater{} \# A tibble: 5 x 3}
\CommentTok{\#\textgreater{}   date       forecast\_discount\_factor ois\_discount\_factor}
\CommentTok{\#\textgreater{}   \textless{}date\textgreater{}                        \textless{}dbl\textgreater{}               \textless{}dbl\textgreater{}}
\CommentTok{\#\textgreater{} 1 2023{-}09{-}26                    1                   1    }
\CommentTok{\#\textgreater{} 2 2024{-}09{-}26                    0.954               0.959}
\CommentTok{\#\textgreater{} 3 2025{-}09{-}26                    0.912               0.921}
\CommentTok{\#\textgreater{} 4 2026{-}09{-}28                    0.875               0.887}
\CommentTok{\#\textgreater{} 5 2028{-}09{-}26                    0.810               0.827}
\end{Highlighting}
\end{Shaded}

\subsection{3.4.4
シングルカーブ評価との比較}\label{ux30b7ux30f3ux30b0ux30ebux30abux30fcux30d6ux8a55ux4fa1ux3068ux306eux6bd4ux8f03}

\begin{Shaded}
\begin{Highlighting}[]
\NormalTok{single\_curve\_swap }\OtherTok{\textless{}{-}}\NormalTok{ GaussQuant}\SpecialCharTok{::}\FunctionTok{make\_vanilla\_swap\_GQL}\NormalTok{(}
  \AttributeTok{effective\_date =} \StringTok{"2023{-}09{-}28"}\NormalTok{,}
  \AttributeTok{maturity\_date =} \StringTok{"2025{-}09{-}28"}\NormalTok{,}
  \AttributeTok{fixed\_rate =} \FloatTok{0.045}\NormalTok{,}
  \AttributeTok{notional =} \DecValTok{10000000}\NormalTok{,}
  \AttributeTok{forecast\_handle =}\NormalTok{ forecast\_handle,}
  \AttributeTok{discount\_handle =}\NormalTok{ forecast\_handle,}
  \AttributeTok{swap\_type =} \StringTok{"payer"}\NormalTok{,}
  \AttributeTok{index =} \StringTok{"Euribor6M"}\NormalTok{,}
  \AttributeTok{fixed\_tenor =}\NormalTok{ GaussQuant}\SpecialCharTok{::}\FunctionTok{period\_years\_GQL}\NormalTok{(}\DecValTok{1}\NormalTok{),}
  \AttributeTok{floating\_tenor =}\NormalTok{ GaussQuant}\SpecialCharTok{::}\FunctionTok{period\_months\_GQL}\NormalTok{(}\DecValTok{6}\NormalTok{),}
  \AttributeTok{calendar =}\NormalTok{ QuantLib}\SpecialCharTok{::}\FunctionTok{TARGET}\NormalTok{()}
\NormalTok{)}

\NormalTok{multi\_curve\_swap }\OtherTok{\textless{}{-}}\NormalTok{ GaussQuant}\SpecialCharTok{::}\FunctionTok{make\_vanilla\_swap\_GQL}\NormalTok{(}
  \AttributeTok{effective\_date =} \StringTok{"2023{-}09{-}28"}\NormalTok{,}
  \AttributeTok{maturity\_date =} \StringTok{"2025{-}09{-}28"}\NormalTok{,}
  \AttributeTok{fixed\_rate =} \FloatTok{0.045}\NormalTok{,}
  \AttributeTok{notional =} \DecValTok{10000000}\NormalTok{,}
  \AttributeTok{forecast\_handle =}\NormalTok{ forecast\_handle,}
  \AttributeTok{discount\_handle =}\NormalTok{ discount\_handle,}
  \AttributeTok{swap\_type =} \StringTok{"payer"}\NormalTok{,}
  \AttributeTok{index =} \StringTok{"Euribor6M"}\NormalTok{,}
  \AttributeTok{fixed\_tenor =}\NormalTok{ GaussQuant}\SpecialCharTok{::}\FunctionTok{period\_years\_GQL}\NormalTok{(}\DecValTok{1}\NormalTok{),}
  \AttributeTok{floating\_tenor =}\NormalTok{ GaussQuant}\SpecialCharTok{::}\FunctionTok{period\_months\_GQL}\NormalTok{(}\DecValTok{6}\NormalTok{),}
  \AttributeTok{calendar =}\NormalTok{ QuantLib}\SpecialCharTok{::}\FunctionTok{TARGET}\NormalTok{()}
\NormalTok{)}

\NormalTok{single\_curve\_summary }\OtherTok{\textless{}{-}}\NormalTok{ GaussQuant}\SpecialCharTok{::}\FunctionTok{swap\_summary\_GQL}\NormalTok{(}
\NormalTok{  single\_curve\_swap}
\NormalTok{)}

\NormalTok{multi\_curve\_summary }\OtherTok{\textless{}{-}}\NormalTok{ GaussQuant}\SpecialCharTok{::}\FunctionTok{swap\_summary\_GQL}\NormalTok{(}
\NormalTok{  multi\_curve\_swap}
\NormalTok{)}

\NormalTok{curve\_valuation\_comparison\_tbl }\OtherTok{\textless{}{-}}\NormalTok{ dplyr}\SpecialCharTok{::}\FunctionTok{bind\_rows}\NormalTok{(}
\NormalTok{  single\_curve\_summary }\SpecialCharTok{|\textgreater{}}
\NormalTok{    dplyr}\SpecialCharTok{::}\FunctionTok{mutate}\NormalTok{(}\AttributeTok{valuation\_method =} \StringTok{"single\_curve"}\NormalTok{),}
\NormalTok{  multi\_curve\_summary }\SpecialCharTok{|\textgreater{}}
\NormalTok{    dplyr}\SpecialCharTok{::}\FunctionTok{mutate}\NormalTok{(}\AttributeTok{valuation\_method =} \StringTok{"multi\_curve"}\NormalTok{)}
\NormalTok{) }\SpecialCharTok{|\textgreater{}}
\NormalTok{  dplyr}\SpecialCharTok{::}\FunctionTok{relocate}\NormalTok{(valuation\_method)}

\NormalTok{GaussQuant}\SpecialCharTok{::}\FunctionTok{show\_tbl\_GQH}\NormalTok{(}
\NormalTok{  curve\_valuation\_comparison\_tbl,}
  \StringTok{"Single{-}curve and multi{-}curve valuation"}\NormalTok{,}
  \AttributeTok{n =} \DecValTok{20}
\NormalTok{)}
\CommentTok{\#\textgreater{} \# A tibble: 2 x 6}
\CommentTok{\#\textgreater{}   valuation\_method    npv fixed\_leg\_npv floating\_leg\_npv fair\_rate fair\_spread}
\CommentTok{\#\textgreater{}   \textless{}chr\textgreater{}             \textless{}dbl\textgreater{}         \textless{}dbl\textgreater{}            \textless{}dbl\textgreater{}     \textless{}dbl\textgreater{}       \textless{}dbl\textgreater{}}
\CommentTok{\#\textgreater{} 1 single\_curve     40292.      {-}840562.          880855.    0.0472    {-}0.00210}
\CommentTok{\#\textgreater{} 2 multi\_curve      39436.      {-}847103.          886539.    0.0471    {-}0.00204}
\end{Highlighting}
\end{Shaded}

両方のスワップは同じフォワードカーブから将来EURIBORを予測するが、シングルカーブ評価ではそのカーブを割引にも使用し、マルチカーブ評価では別のディスカウントカーブを使用する。割引係数が異なるため、同じ将来キャッシュフローであってもNPVおよびフェアレートが変化する。

\section{3.5
ケーススタディ1：EONIAカーブとOIS評価}\label{ux30b1ux30fcux30b9ux30b9ux30bfux30c7ux30a31eoniaux30abux30fcux30d6ux3068oisux8a55ux4fa1}

本ケーススタディでは、2012年12月11日のEONIA市場データからOISカーブを構築する。短期部分には翌日物預金と短期OISを用い、その先には日付指定OISと長期OISを用いる。

カーブ構築後、年末をまたぐ市場金利に含まれる流動性効果を分離し、年末ジャンプを組み込んだ最終的なEONIAカーブを作成する。さらに、カーブ構築に使用した5年OISの市場レートを契約固定金利としてOISを構築し、NPVとフェアレートを確認する。

\subsection{3.5.1
評価日とEONIA市場データ}\label{ux8a55ux4fa1ux65e5ux3068eoniaux5e02ux5834ux30c7ux30fcux30bf}

評価日は2012年12月11日とする。市場データには、翌日物預金、短期OIS、日付指定OISおよび長期OISを使用する。

\begin{Shaded}
\begin{Highlighting}[]
\NormalTok{case\_evaluation\_date }\OtherTok{\textless{}{-}} \StringTok{"2012{-}12{-}11"}

\NormalTok{GaussQuant}\SpecialCharTok{::}\FunctionTok{set\_eval\_date\_GQL}\NormalTok{(}
\NormalTok{  case\_evaluation\_date}
\NormalTok{)}

\NormalTok{case\_eonia\_quotes }\OtherTok{\textless{}{-}}\NormalTok{ GaussQuant}\SpecialCharTok{::}\FunctionTok{ametrano\_bianchetti\_eonia\_quotes\_GQL}\NormalTok{()}

\NormalTok{GaussQuant}\SpecialCharTok{::}\FunctionTok{show\_tbl\_GQH}\NormalTok{(}
\NormalTok{  case\_eonia\_quotes,}
  \StringTok{"EONIA market quotes"}\NormalTok{,}
  \AttributeTok{n =} \DecValTok{40}
\NormalTok{)}
\CommentTok{\#\textgreater{} \# A tibble: 30 x 9}
\CommentTok{\#\textgreater{}    quote\_id  instrument\_type rate\_pct     rate fixing\_days tenor\_n tenor\_unit}
\CommentTok{\#\textgreater{}    \textless{}chr\textgreater{}     \textless{}chr\textgreater{}              \textless{}dbl\textgreater{}    \textless{}dbl\textgreater{}       \textless{}int\textgreater{}   \textless{}int\textgreater{} \textless{}chr\textgreater{}     }
\CommentTok{\#\textgreater{}  1 DEP{-}ON{-}0D deposit            0.04   0.0004            0       1 Days      }
\CommentTok{\#\textgreater{}  2 DEP{-}ON{-}1D deposit            0.04   0.0004            1       1 Days      }
\CommentTok{\#\textgreater{}  3 DEP{-}ON{-}2D deposit            0.04   0.0004            2       1 Days      }
\CommentTok{\#\textgreater{}  4 OIS{-}1W    spot\_ois           0.07   0.0007            2       1 Weeks     }
\CommentTok{\#\textgreater{}  5 OIS{-}2W    spot\_ois           0.069  0.00069           2       2 Weeks     }
\CommentTok{\#\textgreater{}  6 OIS{-}3W    spot\_ois           0.078  0.00078           2       3 Weeks     }
\CommentTok{\#\textgreater{}  7 OIS{-}1M    spot\_ois           0.074  0.00074           2       1 Months    }
\CommentTok{\#\textgreater{}  8 OIS{-}ECB{-}1 dated\_ois          0.046  0.00046          NA      NA \textless{}NA\textgreater{}      }
\CommentTok{\#\textgreater{}  9 OIS{-}ECB{-}2 dated\_ois          0.016  0.00016          NA      NA \textless{}NA\textgreater{}      }
\CommentTok{\#\textgreater{} 10 OIS{-}ECB{-}3 dated\_ois         {-}0.007 {-}0.00007          NA      NA \textless{}NA\textgreater{}      }
\CommentTok{\#\textgreater{} \# i 20 more rows}
\CommentTok{\#\textgreater{} \# i 2 more variables: start\_date \textless{}chr\textgreater{}, end\_date \textless{}chr\textgreater{}}
\end{Highlighting}
\end{Shaded}

\texttt{rate\_pct}はパーセント表示の市場レートであり、\texttt{rate}はQuantLibへ渡す小数表示の金利である。

預金には1日の期間を設定し、評価日からスポット日までを接続するため、Fixing日数を0日、1日および2日とした。短期OISと長期OISにはスポット日から満期までの期間を指定し、日付指定OISには開始日と終了日を直接指定する。

\subsection{3.5.2
EONIAカーブの構築}\label{eoniaux30abux30fcux30d6ux306eux69cbux7bc9}

次に、市場レートからEONIAカーブを構築する。

最初に対数割引係数を3次補間するカーブと、区分的に一定のフォワード金利を用いるカーブを構築する。後者を用いて年末部分のフォワード金利を確認し、年末ジャンプを推定する。

\begin{Shaded}
\begin{Highlighting}[]
\NormalTok{case\_eonia\_benchmark }\OtherTok{\textless{}{-}}\NormalTok{ GaussQuant}\SpecialCharTok{::}\FunctionTok{eonia\_curve\_benchmark\_GQL}\NormalTok{(}
  \AttributeTok{quotes =}\NormalTok{ case\_eonia\_quotes,}
  \AttributeTok{evaluation\_date =}\NormalTok{ case\_evaluation\_date}
\NormalTok{)}

\NormalTok{case\_eonia\_curve }\OtherTok{\textless{}{-}}\NormalTok{ case\_eonia\_benchmark}\SpecialCharTok{$}\NormalTok{final}\SpecialCharTok{$}\NormalTok{curve}
\NormalTok{case\_eonia\_handle }\OtherTok{\textless{}{-}}\NormalTok{ case\_eonia\_benchmark}\SpecialCharTok{$}\NormalTok{final}\SpecialCharTok{$}\NormalTok{curve\_handle}

\NormalTok{case\_jump\_summary }\OtherTok{\textless{}{-}}\NormalTok{ case\_eonia\_benchmark}\SpecialCharTok{$}\NormalTok{jump\_summary}

\NormalTok{GaussQuant}\SpecialCharTok{::}\FunctionTok{show\_tbl\_GQH}\NormalTok{(}
\NormalTok{  case\_jump\_summary,}
  \StringTok{"EONIA year{-}end jump"}\NormalTok{,}
  \AttributeTok{n =} \DecValTok{20}
\NormalTok{)}
\CommentTok{\#\textgreater{} \# A tibble: 1 x 7}
\CommentTok{\#\textgreater{}   original\_forward clean\_forward    t12 jump\_time jump\_rate jump\_discount\_factor}
\CommentTok{\#\textgreater{}              \textless{}dbl\textgreater{}         \textless{}dbl\textgreater{}  \textless{}dbl\textgreater{}     \textless{}dbl\textgreater{}     \textless{}dbl\textgreater{}                \textless{}dbl\textgreater{}}
\CommentTok{\#\textgreater{} 1         0.000815      0.000671 0.0384   0.00548   0.00101                1.000}
\CommentTok{\#\textgreater{} \# i 1 more variable: jump\_date \textless{}date\textgreater{}}
\end{Highlighting}
\end{Shaded}

年末ジャンプの推定では、2012年12月24日から2013年1月7日までのフォワード金利を確認する。次に、2013年1月3日の区分フォワード金利を前後のノードの平均へ置き換え、年末要因を含まないフォワードカーブを作成する。

元のフォワード金利と調整後のフォワード金利の差を、2012年12月31日から2013年1月2日までの期間へ換算し、年末ジャンプを求める。

\subsection{3.5.3
ベンチマーク値の確認}\label{ux30d9ux30f3ux30c1ux30deux30fcux30afux5024ux306eux78baux8a8d}

年末ジャンプを推定する過程で得られたフォワード金利を確認する。

\begin{Shaded}
\begin{Highlighting}[]
\NormalTok{case\_benchmark\_check\_tbl }\OtherTok{\textless{}{-}}\NormalTok{ tibble}\SpecialCharTok{::}\FunctionTok{tribble}\NormalTok{(}
  \SpecialCharTok{\textasciitilde{}}\NormalTok{metric, }\SpecialCharTok{\textasciitilde{}}\NormalTok{reference\_value\_pct, }\SpecialCharTok{\textasciitilde{}}\NormalTok{r\_value\_pct,}
  \StringTok{"Original two{-}week forward"}\NormalTok{, }\FloatTok{0.082}\NormalTok{,}
  \DecValTok{100} \SpecialCharTok{*}\NormalTok{ case\_jump\_summary}\SpecialCharTok{$}\NormalTok{original\_forward[[}\DecValTok{1}\NormalTok{L]],}
  \StringTok{"Adjusted two{-}week forward"}\NormalTok{, }\FloatTok{0.067}\NormalTok{,}
  \DecValTok{100} \SpecialCharTok{*}\NormalTok{ case\_jump\_summary}\SpecialCharTok{$}\NormalTok{clean\_forward[[}\DecValTok{1}\NormalTok{L]],}
  \StringTok{"Year{-}end jump"}\NormalTok{, }\FloatTok{0.101}\NormalTok{,}
  \DecValTok{100} \SpecialCharTok{*}\NormalTok{ case\_jump\_summary}\SpecialCharTok{$}\NormalTok{jump\_rate[[}\DecValTok{1}\NormalTok{L]]}
\NormalTok{) }\SpecialCharTok{|\textgreater{}}
\NormalTok{  dplyr}\SpecialCharTok{::}\FunctionTok{mutate}\NormalTok{(}
    \AttributeTok{difference\_pct =}
\NormalTok{      r\_value\_pct }\SpecialCharTok{{-}}
\NormalTok{      reference\_value\_pct}
\NormalTok{  )}

\NormalTok{GaussQuant}\SpecialCharTok{::}\FunctionTok{show\_tbl\_GQH}\NormalTok{(}
\NormalTok{  case\_benchmark\_check\_tbl,}
  \StringTok{"EONIA benchmark comparison"}\NormalTok{,}
  \AttributeTok{n =} \DecValTok{20}
\NormalTok{)}
\CommentTok{\#\textgreater{} \# A tibble: 3 x 4}
\CommentTok{\#\textgreater{}   metric                    reference\_value\_pct r\_value\_pct difference\_pct}
\CommentTok{\#\textgreater{}   \textless{}chr\textgreater{}                                   \textless{}dbl\textgreater{}       \textless{}dbl\textgreater{}          \textless{}dbl\textgreater{}}
\CommentTok{\#\textgreater{} 1 Original two{-}week forward               0.082      0.0815      {-}0.000469}
\CommentTok{\#\textgreater{} 2 Adjusted two{-}week forward               0.067      0.0671       0.000122}
\CommentTok{\#\textgreater{} 3 Year{-}end jump                           0.101      0.101       {-}0.000137}

\FunctionTok{stopifnot}\NormalTok{(}
  \FunctionTok{all}\NormalTok{(}
    \FunctionTok{abs}\NormalTok{(case\_benchmark\_check\_tbl}\SpecialCharTok{$}\NormalTok{difference\_pct) }\SpecialCharTok{\textless{}=} \FloatTok{0.0005}
\NormalTok{  )}
\NormalTok{)}
\end{Highlighting}
\end{Shaded}

基準値はパーセント単位で小数第3位まで表示されている。そのため、照合では表示桁の半分に相当する0.0005パーセントポイントを許容範囲とする。

\subsection{3.5.4
市場レートの再現性}\label{ux5e02ux5834ux30ecux30fcux30c8ux306eux518dux73feux6027}

構築した最終カーブを用いて、各Rate
Helperのインプライド・レートを求める。

\begin{Shaded}
\begin{Highlighting}[]
\NormalTok{case\_eonia\_validation }\OtherTok{\textless{}{-}}\NormalTok{ case\_eonia\_benchmark}\SpecialCharTok{$}\NormalTok{final\_validation}

\NormalTok{case\_eonia\_repricing\_tbl }\OtherTok{\textless{}{-}}\NormalTok{ case\_eonia\_validation}\SpecialCharTok{$}\NormalTok{quote\_repricing }\SpecialCharTok{|\textgreater{}}
\NormalTok{  dplyr}\SpecialCharTok{::}\FunctionTok{mutate}\NormalTok{(}
    \AttributeTok{quote\_error\_bp =} \DecValTok{10000} \SpecialCharTok{*}\NormalTok{ quote\_error}
\NormalTok{  )}

\NormalTok{GaussQuant}\SpecialCharTok{::}\FunctionTok{show\_tbl\_GQH}\NormalTok{(}
\NormalTok{  case\_eonia\_repricing\_tbl,}
  \StringTok{"EONIA quote repricing"}\NormalTok{,}
  \AttributeTok{n =} \DecValTok{40}
\NormalTok{)}
\CommentTok{\#\textgreater{} \# A tibble: 30 x 10}
\CommentTok{\#\textgreater{}    quote\_id  instrument\_type market\_quote implied\_quote quote\_error pillar\_date}
\CommentTok{\#\textgreater{}    \textless{}chr\textgreater{}     \textless{}chr\textgreater{}                  \textless{}dbl\textgreater{}         \textless{}dbl\textgreater{}       \textless{}dbl\textgreater{} \textless{}date\textgreater{}     }
\CommentTok{\#\textgreater{}  1 DEP{-}ON{-}0D deposit              0.0004      0.000400     2.62e{-}15 2012{-}12{-}12 }
\CommentTok{\#\textgreater{}  2 DEP{-}ON{-}1D deposit              0.0004      0.000400     2.62e{-}15 2012{-}12{-}13 }
\CommentTok{\#\textgreater{}  3 DEP{-}ON{-}2D deposit              0.0004      0.000400     2.62e{-}15 2012{-}12{-}14 }
\CommentTok{\#\textgreater{}  4 OIS{-}1W    spot\_ois             0.0007      0.000700     4.56e{-}13 2012{-}12{-}20 }
\CommentTok{\#\textgreater{}  5 OIS{-}2W    spot\_ois             0.00069     0.000690     1.08e{-}14 2012{-}12{-}27 }
\CommentTok{\#\textgreater{}  6 OIS{-}3W    spot\_ois             0.00078     0.000780    {-}1.36e{-}15 2013{-}01{-}03 }
\CommentTok{\#\textgreater{}  7 OIS{-}1M    spot\_ois             0.00074     0.000740     1.85e{-}15 2013{-}01{-}14 }
\CommentTok{\#\textgreater{}  8 OIS{-}ECB{-}1 dated\_ois            0.00046     0.000460     6.52e{-}12 2013{-}02{-}13 }
\CommentTok{\#\textgreater{}  9 OIS{-}ECB{-}2 dated\_ois            0.00016     0.000160     8.46e{-}12 2013{-}03{-}13 }
\CommentTok{\#\textgreater{} 10 OIS{-}ECB{-}3 dated\_ois           {-}0.00007    {-}0.0000700    9.13e{-}12 2013{-}04{-}10 }
\CommentTok{\#\textgreater{} \# i 20 more rows}
\CommentTok{\#\textgreater{} \# i 4 more variables: discount\_factor \textless{}dbl\textgreater{}, zero\_rate \textless{}dbl\textgreater{},}
\CommentTok{\#\textgreater{} \#   one\_day\_forward\_rate \textless{}dbl\textgreater{}, quote\_error\_bp \textless{}dbl\textgreater{}}

\FunctionTok{stopifnot}\NormalTok{(}
  \FunctionTok{max}\NormalTok{(}
    \FunctionTok{abs}\NormalTok{(case\_eonia\_repricing\_tbl}\SpecialCharTok{$}\NormalTok{quote\_error),}
    \AttributeTok{na.rm =} \ConstantTok{TRUE}
\NormalTok{  ) }\SpecialCharTok{\textless{}} \FloatTok{1e{-}8}
\NormalTok{)}
\end{Highlighting}
\end{Shaded}

各商品のインプライド・レートと市場レートが一致することにより、預金、短期OIS、日付指定OISおよび長期OISが同じ期間構造上で再現されていることを確認できる。

\subsection{3.5.5 5年OISの構築}\label{ux5e74oisux306eux69cbux7bc9}

市場データに含まれる5年OISを取り上げる。契約固定金利には、カーブ構築に用いた5年OISの市場レートを使用する。

\begin{Shaded}
\begin{Highlighting}[]
\NormalTok{case\_ois\_tenor }\OtherTok{\textless{}{-}}\NormalTok{ QuantLib}\SpecialCharTok{::}\FunctionTok{Period}\NormalTok{(}
  \DecValTok{5}\NormalTok{,}
  \StringTok{"Years"}
\NormalTok{)}

\NormalTok{case\_ois\_fixed\_rate }\OtherTok{\textless{}{-}}\NormalTok{ case\_eonia\_quotes }\SpecialCharTok{|\textgreater{}}
\NormalTok{  dplyr}\SpecialCharTok{::}\FunctionTok{filter}\NormalTok{(}
\NormalTok{    quote\_id }\SpecialCharTok{==} \StringTok{"OIS{-}5Y"}
\NormalTok{  ) }\SpecialCharTok{|\textgreater{}}
\NormalTok{  dplyr}\SpecialCharTok{::}\FunctionTok{pull}\NormalTok{(}
\NormalTok{    rate}
\NormalTok{  )}

\NormalTok{case\_eonia\_index }\OtherTok{\textless{}{-}}\NormalTok{ QuantLib}\SpecialCharTok{::}\FunctionTok{Eonia}\NormalTok{(}
\NormalTok{  case\_eonia\_handle}
\NormalTok{)}

\NormalTok{case\_ois\_contract\_tbl }\OtherTok{\textless{}{-}}\NormalTok{ tibble}\SpecialCharTok{::}\FunctionTok{tibble}\NormalTok{(}
  \AttributeTok{evaluation\_date =} \FunctionTok{as.Date}\NormalTok{(case\_evaluation\_date),}
  \AttributeTok{index =} \StringTok{"EONIA"}\NormalTok{,}
  \AttributeTok{tenor =} \StringTok{"5Y"}\NormalTok{,}
  \AttributeTok{settlement\_days =} \DecValTok{2}\NormalTok{L,}
  \AttributeTok{contract\_fixed\_rate =}\NormalTok{ case\_ois\_fixed\_rate}
\NormalTok{)}

\NormalTok{GaussQuant}\SpecialCharTok{::}\FunctionTok{show\_tbl\_GQH}\NormalTok{(}
\NormalTok{  case\_ois\_contract\_tbl,}
  \StringTok{"Five{-}year EONIA OIS contract"}\NormalTok{,}
  \AttributeTok{n =} \DecValTok{20}
\NormalTok{)}
\CommentTok{\#\textgreater{} \# A tibble: 1 x 5}
\CommentTok{\#\textgreater{}   evaluation\_date index tenor settlement\_days contract\_fixed\_rate}
\CommentTok{\#\textgreater{}   \textless{}date\textgreater{}          \textless{}chr\textgreater{} \textless{}chr\textgreater{}           \textless{}int\textgreater{}               \textless{}dbl\textgreater{}}
\CommentTok{\#\textgreater{} 1 2012{-}12{-}11      EONIA 5Y                  2             0.00456}
\end{Highlighting}
\end{Shaded}

\subsection{3.5.6
OISの構築と計算結果の照合}\label{oisux306eux69cbux7bc9ux3068ux8a08ux7b97ux7d50ux679cux306eux7167ux5408}

最初に、QuantLibの\texttt{MakeOIS}を直接使用してOISを構築する。

\begin{Shaded}
\begin{Highlighting}[]
\NormalTok{case\_ois\_builder\_quantlib }\OtherTok{\textless{}{-}}\NormalTok{ QuantLib}\SpecialCharTok{::}\FunctionTok{MakeOIS}\NormalTok{(}
  \AttributeTok{swapTenor =}\NormalTok{ case\_ois\_tenor,}
  \AttributeTok{overnightIndex =}\NormalTok{ case\_eonia\_index,}
  \AttributeTok{fixedRate =}\NormalTok{ case\_ois\_fixed\_rate}
\NormalTok{)}

\NormalTok{case\_ois\_quantlib }\OtherTok{\textless{}{-}}\NormalTok{ QuantLib}\SpecialCharTok{::}\FunctionTok{MakeOIS\_makeOIS}\NormalTok{(}
\NormalTok{  case\_ois\_builder\_quantlib}
\NormalTok{)}
\end{Highlighting}
\end{Shaded}

次に、同じ期間、同じEONIAインデックスおよび同じ固定金利を用いてOISを構築する。

\begin{Shaded}
\begin{Highlighting}[]
\NormalTok{case\_ois\_gaussquant }\OtherTok{\textless{}{-}}\NormalTok{ GaussQuant}\SpecialCharTok{::}\FunctionTok{build\_ois\_GQL}\NormalTok{(}
  \AttributeTok{swap\_tenor =}\NormalTok{ case\_ois\_tenor,}
  \AttributeTok{overnight\_index =}\NormalTok{ case\_eonia\_index,}
  \AttributeTok{fixed\_rate =}\NormalTok{ case\_ois\_fixed\_rate}
\NormalTok{)}
\end{Highlighting}
\end{Shaded}

両方のOISについて、NPV、固定レグの現在価値、翌日物レグの現在価値、フェアレートおよびフェアスプレッドを確認する。

\begin{Shaded}
\begin{Highlighting}[]
\NormalTok{case\_ois\_quantlib\_summary }\OtherTok{\textless{}{-}}\NormalTok{ GaussQuant}\SpecialCharTok{::}\FunctionTok{swap\_summary\_GQL}\NormalTok{(}
\NormalTok{  case\_ois\_quantlib}
\NormalTok{)}

\NormalTok{case\_ois\_gaussquant\_summary }\OtherTok{\textless{}{-}}\NormalTok{ GaussQuant}\SpecialCharTok{::}\FunctionTok{swap\_summary\_GQL}\NormalTok{(}
\NormalTok{  case\_ois\_gaussquant}
\NormalTok{)}

\NormalTok{case\_ois\_comparison\_tbl }\OtherTok{\textless{}{-}}\NormalTok{ dplyr}\SpecialCharTok{::}\FunctionTok{bind\_rows}\NormalTok{(}
\NormalTok{  case\_ois\_quantlib\_summary }\SpecialCharTok{|\textgreater{}}
\NormalTok{    dplyr}\SpecialCharTok{::}\FunctionTok{mutate}\NormalTok{(}
      \AttributeTok{calculation\_route =} \StringTok{"QuantLib"}
\NormalTok{    ),}
\NormalTok{  case\_ois\_gaussquant\_summary }\SpecialCharTok{|\textgreater{}}
\NormalTok{    dplyr}\SpecialCharTok{::}\FunctionTok{mutate}\NormalTok{(}
      \AttributeTok{calculation\_route =} \StringTok{"GaussQuant"}
\NormalTok{    )}
\NormalTok{) }\SpecialCharTok{|\textgreater{}}
\NormalTok{  dplyr}\SpecialCharTok{::}\FunctionTok{relocate}\NormalTok{(}
\NormalTok{    calculation\_route}
\NormalTok{  )}

\NormalTok{GaussQuant}\SpecialCharTok{::}\FunctionTok{show\_tbl\_GQH}\NormalTok{(}
\NormalTok{  case\_ois\_comparison\_tbl,}
  \StringTok{"Five{-}year EONIA OIS valuation"}\NormalTok{,}
  \AttributeTok{n =} \DecValTok{20}
\NormalTok{)}
\CommentTok{\#\textgreater{} \# A tibble: 2 x 6}
\CommentTok{\#\textgreater{}   calculation\_route       npv fixed\_leg\_npv floating\_leg\_npv fair\_rate}
\CommentTok{\#\textgreater{}   \textless{}chr\textgreater{}                 \textless{}dbl\textgreater{}         \textless{}dbl\textgreater{}            \textless{}dbl\textgreater{}     \textless{}dbl\textgreater{}}
\CommentTok{\#\textgreater{} 1 QuantLib          {-}1.65e{-}13       {-}0.0230           0.0230   0.00456}
\CommentTok{\#\textgreater{} 2 GaussQuant        {-}1.65e{-}13       {-}0.0230           0.0230   0.00456}
\CommentTok{\#\textgreater{} \# i 1 more variable: fair\_spread \textless{}dbl\textgreater{}}

\NormalTok{case\_ois\_metrics }\OtherTok{\textless{}{-}} \FunctionTok{c}\NormalTok{(}
  \StringTok{"npv"}\NormalTok{,}
  \StringTok{"fixed\_leg\_npv"}\NormalTok{,}
  \StringTok{"floating\_leg\_npv"}\NormalTok{,}
  \StringTok{"fair\_rate"}\NormalTok{,}
  \StringTok{"fair\_spread"}
\NormalTok{)}

\NormalTok{case\_quantlib\_values }\OtherTok{\textless{}{-}} \FunctionTok{unlist}\NormalTok{(}
\NormalTok{  case\_ois\_quantlib\_summary }\SpecialCharTok{|\textgreater{}}
\NormalTok{    dplyr}\SpecialCharTok{::}\FunctionTok{select}\NormalTok{(}
\NormalTok{      dplyr}\SpecialCharTok{::}\FunctionTok{all\_of}\NormalTok{(case\_ois\_metrics)}
\NormalTok{    ),}
  \AttributeTok{use.names =} \ConstantTok{FALSE}
\NormalTok{)}

\NormalTok{case\_gaussquant\_values }\OtherTok{\textless{}{-}} \FunctionTok{unlist}\NormalTok{(}
\NormalTok{  case\_ois\_gaussquant\_summary }\SpecialCharTok{|\textgreater{}}
\NormalTok{    dplyr}\SpecialCharTok{::}\FunctionTok{select}\NormalTok{(}
\NormalTok{      dplyr}\SpecialCharTok{::}\FunctionTok{all\_of}\NormalTok{(case\_ois\_metrics)}
\NormalTok{    ),}
  \AttributeTok{use.names =} \ConstantTok{FALSE}
\NormalTok{)}

\FunctionTok{stopifnot}\NormalTok{(}
  \FunctionTok{isTRUE}\NormalTok{(}
    \FunctionTok{all.equal}\NormalTok{(}
\NormalTok{      case\_quantlib\_values,}
\NormalTok{      case\_gaussquant\_values,}
      \AttributeTok{tolerance =} \FloatTok{1e{-}12}
\NormalTok{    )}
\NormalTok{  )}
\NormalTok{)}
\end{Highlighting}
\end{Shaded}

\subsection{3.5.7
市場固定金利とフェアレート}\label{ux5e02ux5834ux56faux5b9aux91d1ux5229ux3068ux30d5ux30a7ux30a2ux30ecux30fcux30c8}

カーブ構築に使用した5年OISの市場レートと、構築後のカーブから求めたフェアレートを比較する。

\begin{Shaded}
\begin{Highlighting}[]
\NormalTok{case\_ois\_rate\_check\_tbl }\OtherTok{\textless{}{-}}\NormalTok{ tibble}\SpecialCharTok{::}\FunctionTok{tibble}\NormalTok{(}
  \AttributeTok{market\_fixed\_rate =}\NormalTok{ case\_ois\_fixed\_rate,}
  \AttributeTok{calculated\_fair\_rate =}
\NormalTok{    case\_ois\_gaussquant\_summary}\SpecialCharTok{$}\NormalTok{fair\_rate[[}\DecValTok{1}\NormalTok{L]],}
  \AttributeTok{difference\_bp =} \DecValTok{10000} \SpecialCharTok{*}\NormalTok{ (}
\NormalTok{    case\_ois\_gaussquant\_summary}\SpecialCharTok{$}\NormalTok{fair\_rate[[}\DecValTok{1}\NormalTok{L]] }\SpecialCharTok{{-}}
\NormalTok{      case\_ois\_fixed\_rate}
\NormalTok{  ),}
  \AttributeTok{npv =}\NormalTok{ case\_ois\_gaussquant\_summary}\SpecialCharTok{$}\NormalTok{npv[[}\DecValTok{1}\NormalTok{L]]}
\NormalTok{)}

\NormalTok{GaussQuant}\SpecialCharTok{::}\FunctionTok{show\_tbl\_GQH}\NormalTok{(}
\NormalTok{  case\_ois\_rate\_check\_tbl,}
  \StringTok{"Five{-}year EONIA OIS rate check"}\NormalTok{,}
  \AttributeTok{n =} \DecValTok{20}
\NormalTok{)}
\CommentTok{\#\textgreater{} \# A tibble: 1 x 4}
\CommentTok{\#\textgreater{}   market\_fixed\_rate calculated\_fair\_rate difference\_bp       npv}
\CommentTok{\#\textgreater{}               \textless{}dbl\textgreater{}                \textless{}dbl\textgreater{}         \textless{}dbl\textgreater{}     \textless{}dbl\textgreater{}}
\CommentTok{\#\textgreater{} 1           0.00456              0.00456     {-}3.29e{-}10 {-}1.65e{-}13}

\FunctionTok{stopifnot}\NormalTok{(}
  \FunctionTok{abs}\NormalTok{(}
\NormalTok{    case\_ois\_gaussquant\_summary}\SpecialCharTok{$}\NormalTok{fair\_rate[[}\DecValTok{1}\NormalTok{L]] }\SpecialCharTok{{-}}
\NormalTok{      case\_ois\_fixed\_rate}
\NormalTok{  ) }\SpecialCharTok{\textless{}} \FloatTok{1e{-}8}
\NormalTok{)}
\end{Highlighting}
\end{Shaded}

5年OISはカーブ構築に使用した商品の一つであるため、市場固定金利とカーブから再計算したフェアレートは一致し、NPVはゼロに近い値となる。

\subsection{3.5.8
固定レグと翌日物レグ}\label{ux56faux5b9aux30ecux30b0ux3068ux7fccux65e5ux7269ux30ecux30b0}

最後に、5年OISの固定レグと翌日物レグのキャッシュフローを確認する。

\begin{Shaded}
\begin{Highlighting}[]
\NormalTok{case\_ois\_fixed\_leg\_tbl }\OtherTok{\textless{}{-}}\NormalTok{ GaussQuant}\SpecialCharTok{::}\FunctionTok{leg\_to\_cashflow\_tbl\_GQL}\NormalTok{(}
\NormalTok{  case\_ois\_gaussquant}\SpecialCharTok{$}\FunctionTok{fixedLeg}\NormalTok{()}
\NormalTok{)}

\NormalTok{case\_ois\_overnight\_leg\_tbl }\OtherTok{\textless{}{-}}\NormalTok{ GaussQuant}\SpecialCharTok{::}\FunctionTok{leg\_to\_cashflow\_tbl\_GQL}\NormalTok{(}
\NormalTok{  case\_ois\_gaussquant}\SpecialCharTok{$}\FunctionTok{overnightLeg}\NormalTok{()}
\NormalTok{)}

\NormalTok{GaussQuant}\SpecialCharTok{::}\FunctionTok{show\_tbl\_GQH}\NormalTok{(}
\NormalTok{  case\_ois\_fixed\_leg\_tbl,}
  \StringTok{"Five{-}year EONIA OIS fixed leg"}\NormalTok{,}
  \AttributeTok{n =} \DecValTok{20}
\NormalTok{)}
\CommentTok{\#\textgreater{} \# A tibble: 5 x 2}
\CommentTok{\#\textgreater{}   date        amount}
\CommentTok{\#\textgreater{}   \textless{}chr\textgreater{}        \textless{}dbl\textgreater{}}
\CommentTok{\#\textgreater{} 1 2013{-}12{-}13 0.00462}
\CommentTok{\#\textgreater{} 2 2014{-}12{-}15 0.00465}
\CommentTok{\#\textgreater{} 3 2015{-}12{-}14 0.00461}
\CommentTok{\#\textgreater{} 4 2016{-}12{-}13 0.00462}
\CommentTok{\#\textgreater{} 5 2017{-}12{-}13 0.00462}

\NormalTok{GaussQuant}\SpecialCharTok{::}\FunctionTok{show\_tbl\_GQH}\NormalTok{(}
\NormalTok{  case\_ois\_overnight\_leg\_tbl,}
  \StringTok{"Five{-}year EONIA OIS overnight leg"}\NormalTok{,}
  \AttributeTok{n =} \DecValTok{20}
\NormalTok{)}
\CommentTok{\#\textgreater{} \# A tibble: 5 x 2}
\CommentTok{\#\textgreater{}   date          amount}
\CommentTok{\#\textgreater{}   \textless{}chr\textgreater{}          \textless{}dbl\textgreater{}}
\CommentTok{\#\textgreater{} 1 2013{-}12{-}13 0.0000146}
\CommentTok{\#\textgreater{} 2 2014{-}12{-}15 0.000718 }
\CommentTok{\#\textgreater{} 3 2015{-}12{-}14 0.00314  }
\CommentTok{\#\textgreater{} 4 2016{-}12{-}13 0.00730  }
\CommentTok{\#\textgreater{} 5 2017{-}12{-}13 0.0122}
\end{Highlighting}
\end{Shaded}

固定レグと翌日物レグは反対方向のキャッシュフローを持ち、市場固定金利がフェアレートと一致する場合には、両レグの現在価値が相殺される。

\section{3.6
ケーススタディ2：市場レート変動と評価損益}\label{ux30b1ux30fcux30b9ux30b9ux30bfux30c7ux30a32ux5e02ux5834ux30ecux30fcux30c8ux5909ux52d5ux3068ux8a55ux4fa1ux640dux76ca}

第9節では、2012年12月11日のEONIA市場データから期間構造を構築し、5年OISの市場固定金利とフェアレートが一致することを確認した。

本ケーススタディでは、評価日とOISの約定条件を変更せず、カーブ構築に使用する市場レートだけを上下へ10ベーシスポイント変化させる。各シナリオについてEONIAカーブを再構築し、第9節と同じ5年OISを再評価する。

比較するシナリオは、次の三つである。

\begin{itemize}
\tightlist
\item
  EONIA市場レートを一律に10ベーシスポイント低下
\item
  第9節と同じ基準市場レート
\item
  EONIA市場レートを一律に10ベーシスポイント上昇
\end{itemize}

契約固定金利には、第9節で設定した5年OISの市場レート0.456\%を引き続き使用する。したがって、各シナリオのNPVの差は、市場レートの変化による評価損益を表す。

\subsection{3.6.1
市場レートシナリオ}\label{ux5e02ux5834ux30ecux30fcux30c8ux30b7ux30caux30eaux30aa}

最初に、基準となるEONIA市場レートへ、マイナス10ベーシスポイント、ゼロおよびプラス10ベーシスポイントのシフトを加える。

ここでは構築済みのゼロカーブを直接動かすのではなく、預金、短期OIS、日付指定OISおよび長期OISの市場レートを変化させる。その後、それぞれの市場データからEONIAカーブを改めて構築する。

\begin{Shaded}
\begin{Highlighting}[]
\NormalTok{case\_rate\_scenario\_tbl }\OtherTok{\textless{}{-}}\NormalTok{ tibble}\SpecialCharTok{::}\FunctionTok{tribble}\NormalTok{(}
  \SpecialCharTok{\textasciitilde{}}\NormalTok{scenario, }\SpecialCharTok{\textasciitilde{}}\NormalTok{shift\_bp,}
  \StringTok{"10bp低下"}\NormalTok{, }\SpecialCharTok{{-}}\DecValTok{10}\NormalTok{,}
  \StringTok{"基準"}\NormalTok{,       }\DecValTok{0}\NormalTok{,}
  \StringTok{"10bp上昇"}\NormalTok{,  }\DecValTok{10}
\NormalTok{) }\SpecialCharTok{|\textgreater{}}
\NormalTok{  dplyr}\SpecialCharTok{::}\FunctionTok{mutate}\NormalTok{(}
    \AttributeTok{quotes =}\NormalTok{ purrr}\SpecialCharTok{::}\FunctionTok{map}\NormalTok{(}
\NormalTok{      shift\_bp,}
      \ControlFlowTok{function}\NormalTok{(shift\_bp) \{}
\NormalTok{        case\_eonia\_quotes }\SpecialCharTok{|\textgreater{}}
\NormalTok{          dplyr}\SpecialCharTok{::}\FunctionTok{mutate}\NormalTok{(}
            \AttributeTok{rate =}\NormalTok{ rate }\SpecialCharTok{+}\NormalTok{ shift\_bp }\SpecialCharTok{/} \DecValTok{10000}\NormalTok{,}
            \AttributeTok{rate\_pct =} \DecValTok{100} \SpecialCharTok{*}\NormalTok{ rate}
\NormalTok{          )}
\NormalTok{      \}}
\NormalTok{    )}
\NormalTok{  )}
\end{Highlighting}
\end{Shaded}

10ベーシスポイントは、小数表示の金利では0.001、パーセント表示では0.1パーセントポイントに相当する。

市場レートの変化を確認するため、翌日物預金、1か月OIS、5年OISおよび30年OISを取り出す。

\begin{Shaded}
\begin{Highlighting}[]
\NormalTok{case\_rate\_scenario\_quote\_tbl }\OtherTok{\textless{}{-}}\NormalTok{ purrr}\SpecialCharTok{::}\FunctionTok{map2\_dfr}\NormalTok{(}
\NormalTok{  case\_rate\_scenario\_tbl}\SpecialCharTok{$}\NormalTok{scenario,}
\NormalTok{  case\_rate\_scenario\_tbl}\SpecialCharTok{$}\NormalTok{quotes,}
  \ControlFlowTok{function}\NormalTok{(scenario, quotes) \{}
\NormalTok{    quotes }\SpecialCharTok{|\textgreater{}}
\NormalTok{      dplyr}\SpecialCharTok{::}\FunctionTok{filter}\NormalTok{(}
\NormalTok{        quote\_id }\SpecialCharTok{\%in\%} \FunctionTok{c}\NormalTok{(}
          \StringTok{"DEP{-}ON{-}0D"}\NormalTok{,}
          \StringTok{"OIS{-}1M"}\NormalTok{,}
          \StringTok{"OIS{-}5Y"}\NormalTok{,}
          \StringTok{"OIS{-}30Y"}
\NormalTok{        )}
\NormalTok{      ) }\SpecialCharTok{|\textgreater{}}
\NormalTok{      dplyr}\SpecialCharTok{::}\FunctionTok{transmute}\NormalTok{(}
        \AttributeTok{scenario =}\NormalTok{ scenario,}
        \AttributeTok{quote\_id =}\NormalTok{ quote\_id,}
        \AttributeTok{instrument\_type =}\NormalTok{ instrument\_type,}
        \AttributeTok{rate\_pct =}\NormalTok{ rate\_pct,}
        \AttributeTok{rate =}\NormalTok{ rate}
\NormalTok{      )}
\NormalTok{  \}}
\NormalTok{)}

\NormalTok{GaussQuant}\SpecialCharTok{::}\FunctionTok{show\_tbl\_GQH}\NormalTok{(}
\NormalTok{  case\_rate\_scenario\_quote\_tbl,}
  \StringTok{"EONIA market{-}rate scenarios"}\NormalTok{,}
  \AttributeTok{n =} \DecValTok{20}
\NormalTok{)}
\CommentTok{\#\textgreater{} \# A tibble: 12 x 5}
\CommentTok{\#\textgreater{}    scenario quote\_id  instrument\_type rate\_pct     rate}
\CommentTok{\#\textgreater{}    \textless{}chr\textgreater{}    \textless{}chr\textgreater{}     \textless{}chr\textgreater{}              \textless{}dbl\textgreater{}    \textless{}dbl\textgreater{}}
\CommentTok{\#\textgreater{}  1 10bp低下 DEP{-}ON{-}0D deposit           {-}0.06  {-}0.0006 }
\CommentTok{\#\textgreater{}  2 10bp低下 OIS{-}1M    spot\_ois          {-}0.026 {-}0.00026}
\CommentTok{\#\textgreater{}  3 10bp低下 OIS{-}5Y    spot\_ois           0.356  0.00356}
\CommentTok{\#\textgreater{}  4 10bp低下 OIS{-}30Y   spot\_ois           1.94   0.0194 }
\CommentTok{\#\textgreater{}  5 基準     DEP{-}ON{-}0D deposit            0.04   0.0004 }
\CommentTok{\#\textgreater{}  6 基準     OIS{-}1M    spot\_ois           0.074  0.00074}
\CommentTok{\#\textgreater{}  7 基準     OIS{-}5Y    spot\_ois           0.456  0.00456}
\CommentTok{\#\textgreater{}  8 基準     OIS{-}30Y   spot\_ois           2.04   0.0204 }
\CommentTok{\#\textgreater{}  9 10bp上昇 DEP{-}ON{-}0D deposit            0.14   0.0014 }
\CommentTok{\#\textgreater{} 10 10bp上昇 OIS{-}1M    spot\_ois           0.174  0.00174}
\CommentTok{\#\textgreater{} 11 10bp上昇 OIS{-}5Y    spot\_ois           0.556  0.00556}
\CommentTok{\#\textgreater{} 12 10bp上昇 OIS{-}30Y   spot\_ois           2.14   0.0214}
\end{Highlighting}
\end{Shaded}

各商品のレートは、基準値から一律に10ベーシスポイント低下または上昇している。評価日、Fixing日数、期間、開始日および終了日などの市場慣行は変更しない。

\subsection{3.6.2
シナリオ別EONIAカーブ}\label{ux30b7ux30caux30eaux30aaux5225eoniaux30abux30fcux30d6}

各シナリオの市場データから、年末ジャンプを含むEONIAカーブを再構築する。

\begin{Shaded}
\begin{Highlighting}[]
\NormalTok{case\_rate\_scenario\_tbl }\OtherTok{\textless{}{-}}\NormalTok{ case\_rate\_scenario\_tbl }\SpecialCharTok{|\textgreater{}}
\NormalTok{  dplyr}\SpecialCharTok{::}\FunctionTok{mutate}\NormalTok{(}
    \AttributeTok{benchmark =}\NormalTok{ purrr}\SpecialCharTok{::}\FunctionTok{map}\NormalTok{(}
\NormalTok{      quotes,}
      \ControlFlowTok{function}\NormalTok{(quotes) \{}
\NormalTok{        GaussQuant}\SpecialCharTok{::}\FunctionTok{eonia\_curve\_benchmark\_GQL}\NormalTok{(}
          \AttributeTok{quotes =}\NormalTok{ quotes,}
          \AttributeTok{evaluation\_date =}\NormalTok{ case\_evaluation\_date}
\NormalTok{        )}
\NormalTok{      \}}
\NormalTok{    ),}
    \AttributeTok{curve =}\NormalTok{ purrr}\SpecialCharTok{::}\FunctionTok{map}\NormalTok{(}
\NormalTok{      benchmark,}
      \SpecialCharTok{\textasciitilde{}}\NormalTok{ .x}\SpecialCharTok{$}\NormalTok{final}\SpecialCharTok{$}\NormalTok{curve}
\NormalTok{    ),}
    \AttributeTok{curve\_handle =}\NormalTok{ purrr}\SpecialCharTok{::}\FunctionTok{map}\NormalTok{(}
\NormalTok{      benchmark,}
      \SpecialCharTok{\textasciitilde{}}\NormalTok{ .x}\SpecialCharTok{$}\NormalTok{final}\SpecialCharTok{$}\NormalTok{curve\_handle}
\NormalTok{    ),}
    \AttributeTok{max\_quote\_error =}\NormalTok{ purrr}\SpecialCharTok{::}\FunctionTok{map\_dbl}\NormalTok{(}
\NormalTok{      benchmark,}
      \ControlFlowTok{function}\NormalTok{(benchmark) \{}
        \FunctionTok{max}\NormalTok{(}
          \FunctionTok{abs}\NormalTok{(}
\NormalTok{            benchmark}\SpecialCharTok{$}\NormalTok{final\_validation}\SpecialCharTok{$}
\NormalTok{              quote\_repricing}\SpecialCharTok{$}
\NormalTok{              quote\_error}
\NormalTok{          ),}
          \AttributeTok{na.rm =} \ConstantTok{TRUE}
\NormalTok{        )}
\NormalTok{      \}}
\NormalTok{    )}
\NormalTok{  )}

\NormalTok{case\_curve\_validation\_tbl }\OtherTok{\textless{}{-}}\NormalTok{ case\_rate\_scenario\_tbl }\SpecialCharTok{|\textgreater{}}
\NormalTok{  dplyr}\SpecialCharTok{::}\FunctionTok{select}\NormalTok{(}
\NormalTok{    scenario,}
\NormalTok{    shift\_bp,}
\NormalTok{    max\_quote\_error}
\NormalTok{  )}

\NormalTok{GaussQuant}\SpecialCharTok{::}\FunctionTok{show\_tbl\_GQH}\NormalTok{(}
\NormalTok{  case\_curve\_validation\_tbl,}
  \StringTok{"EONIA scenario{-}curve validation"}\NormalTok{,}
  \AttributeTok{n =} \DecValTok{20}
\NormalTok{)}
\CommentTok{\#\textgreater{} \# A tibble: 3 x 3}
\CommentTok{\#\textgreater{}   scenario shift\_bp max\_quote\_error}
\CommentTok{\#\textgreater{}   \textless{}chr\textgreater{}       \textless{}dbl\textgreater{}           \textless{}dbl\textgreater{}}
\CommentTok{\#\textgreater{} 1 10bp低下      {-}10        9.84e{-}12}
\CommentTok{\#\textgreater{} 2 基準            0        9.86e{-}12}
\CommentTok{\#\textgreater{} 3 10bp上昇       10        9.87e{-}12}

\FunctionTok{stopifnot}\NormalTok{(}
  \FunctionTok{all}\NormalTok{(}
\NormalTok{    case\_curve\_validation\_tbl}\SpecialCharTok{$}\NormalTok{max\_quote\_error }\SpecialCharTok{\textless{}} \FloatTok{1e{-}8}
\NormalTok{  )}
\NormalTok{)}
\end{Highlighting}
\end{Shaded}

すべてのシナリオについて、Rate
Helperのインプライド・レートが入力した市場レートと一致することを確認する。これにより、評価損益の比較に使用する各カーブが、それぞれの市場データを再現していることが分かる。

\subsection{3.6.3
同一OISの再評価}\label{ux540cux4e00oisux306eux518dux8a55ux4fa1}

第9節と同じ5年OISを、各シナリオのEONIAカーブへ接続する。

契約期間は5年、契約固定金利は0.456\%のままとし、市場カーブだけを変更する。

\begin{Shaded}
\begin{Highlighting}[]
\NormalTok{case\_rate\_scenario\_tbl }\OtherTok{\textless{}{-}}\NormalTok{ case\_rate\_scenario\_tbl }\SpecialCharTok{|\textgreater{}}
\NormalTok{  dplyr}\SpecialCharTok{::}\FunctionTok{mutate}\NormalTok{(}
    \AttributeTok{eonia\_index =}\NormalTok{ purrr}\SpecialCharTok{::}\FunctionTok{map}\NormalTok{(}
\NormalTok{      curve\_handle,}
\NormalTok{      QuantLib}\SpecialCharTok{::}\NormalTok{Eonia}
\NormalTok{    ),}
    \AttributeTok{ois =}\NormalTok{ purrr}\SpecialCharTok{::}\FunctionTok{map}\NormalTok{(}
\NormalTok{      eonia\_index,}
      \ControlFlowTok{function}\NormalTok{(eonia\_index) \{}
\NormalTok{        GaussQuant}\SpecialCharTok{::}\FunctionTok{build\_ois\_GQL}\NormalTok{(}
          \AttributeTok{swap\_tenor =}\NormalTok{ case\_ois\_tenor,}
          \AttributeTok{overnight\_index =}\NormalTok{ eonia\_index,}
          \AttributeTok{fixed\_rate =}\NormalTok{ case\_ois\_fixed\_rate}
\NormalTok{        )}
\NormalTok{      \}}
\NormalTok{    ),}
    \AttributeTok{valuation =}\NormalTok{ purrr}\SpecialCharTok{::}\FunctionTok{map}\NormalTok{(}
\NormalTok{      ois,}
\NormalTok{      GaussQuant}\SpecialCharTok{::}\NormalTok{swap\_summary\_GQL}
\NormalTok{    )}
\NormalTok{  )}
\end{Highlighting}
\end{Shaded}

各シナリオの評価結果を一つの表へまとめる。

\begin{Shaded}
\begin{Highlighting}[]
\NormalTok{case\_ois\_scenario\_result\_tbl }\OtherTok{\textless{}{-}}\NormalTok{ case\_rate\_scenario\_tbl }\SpecialCharTok{|\textgreater{}}
\NormalTok{  dplyr}\SpecialCharTok{::}\FunctionTok{select}\NormalTok{(}
\NormalTok{    scenario,}
\NormalTok{    shift\_bp,}
\NormalTok{    valuation}
\NormalTok{  ) }\SpecialCharTok{|\textgreater{}}
\NormalTok{  tidyr}\SpecialCharTok{::}\FunctionTok{unnest}\NormalTok{(}
    \AttributeTok{cols =}\NormalTok{ valuation}
\NormalTok{  )}

\NormalTok{case\_base\_fair\_rate }\OtherTok{\textless{}{-}}\NormalTok{ case\_ois\_scenario\_result\_tbl }\SpecialCharTok{|\textgreater{}}
\NormalTok{  dplyr}\SpecialCharTok{::}\FunctionTok{filter}\NormalTok{(}
\NormalTok{    scenario }\SpecialCharTok{==} \StringTok{"基準"}
\NormalTok{  ) }\SpecialCharTok{|\textgreater{}}
\NormalTok{  dplyr}\SpecialCharTok{::}\FunctionTok{pull}\NormalTok{(}
\NormalTok{    fair\_rate}
\NormalTok{  )}

\NormalTok{case\_base\_npv }\OtherTok{\textless{}{-}}\NormalTok{ case\_ois\_scenario\_result\_tbl }\SpecialCharTok{|\textgreater{}}
\NormalTok{  dplyr}\SpecialCharTok{::}\FunctionTok{filter}\NormalTok{(}
\NormalTok{    scenario }\SpecialCharTok{==} \StringTok{"基準"}
\NormalTok{  ) }\SpecialCharTok{|\textgreater{}}
\NormalTok{  dplyr}\SpecialCharTok{::}\FunctionTok{pull}\NormalTok{(}
\NormalTok{    npv}
\NormalTok{  )}

\NormalTok{case\_ois\_scenario\_result\_tbl }\OtherTok{\textless{}{-}}\NormalTok{ case\_ois\_scenario\_result\_tbl }\SpecialCharTok{|\textgreater{}}
\NormalTok{  dplyr}\SpecialCharTok{::}\FunctionTok{mutate}\NormalTok{(}
    \AttributeTok{contract\_fixed\_rate =}\NormalTok{ case\_ois\_fixed\_rate,}
    \AttributeTok{fair\_rate\_change\_bp =} \DecValTok{10000} \SpecialCharTok{*}\NormalTok{ (}
\NormalTok{      fair\_rate }\SpecialCharTok{{-}}
\NormalTok{        case\_base\_fair\_rate}
\NormalTok{    ),}
    \AttributeTok{payer\_valuation\_change =}
\NormalTok{      npv }\SpecialCharTok{{-}}
\NormalTok{      case\_base\_npv,}
    \AttributeTok{receiver\_npv =} \SpecialCharTok{{-}}\NormalTok{npv,}
    \AttributeTok{receiver\_valuation\_change =}
      \SpecialCharTok{{-}}\NormalTok{(}
\NormalTok{        npv }\SpecialCharTok{{-}}
\NormalTok{          case\_base\_npv}
\NormalTok{      )}
\NormalTok{  ) }\SpecialCharTok{|\textgreater{}}
\NormalTok{  dplyr}\SpecialCharTok{::}\FunctionTok{select}\NormalTok{(}
\NormalTok{    scenario,}
\NormalTok{    shift\_bp,}
\NormalTok{    contract\_fixed\_rate,}
\NormalTok{    fair\_rate,}
\NormalTok{    fair\_rate\_change\_bp,}
    \AttributeTok{payer\_npv =}\NormalTok{ npv,}
\NormalTok{    payer\_valuation\_change,}
\NormalTok{    receiver\_npv,}
\NormalTok{    receiver\_valuation\_change}
\NormalTok{  )}

\NormalTok{GaussQuant}\SpecialCharTok{::}\FunctionTok{show\_tbl\_GQH}\NormalTok{(}
\NormalTok{  case\_ois\_scenario\_result\_tbl,}
  \StringTok{"Five{-}year EONIA OIS scenario valuation"}\NormalTok{,}
  \AttributeTok{n =} \DecValTok{20}
\NormalTok{)}
\CommentTok{\#\textgreater{} \# A tibble: 3 x 9}
\CommentTok{\#\textgreater{}   scenario shift\_bp contract\_fixed\_rate fair\_rate fair\_rate\_change\_bp payer\_npv}
\CommentTok{\#\textgreater{}   \textless{}chr\textgreater{}       \textless{}dbl\textgreater{}               \textless{}dbl\textgreater{}     \textless{}dbl\textgreater{}               \textless{}dbl\textgreater{}     \textless{}dbl\textgreater{}}
\CommentTok{\#\textgreater{} 1 10bp低下      {-}10             0.00456   0.00356              {-}10.0  {-}5.05e{-} 3}
\CommentTok{\#\textgreater{} 2 基準            0             0.00456   0.00456                0    {-}1.65e{-}13}
\CommentTok{\#\textgreater{} 3 10bp上昇       10             0.00456   0.00556               10.00  5.02e{-} 3}
\CommentTok{\#\textgreater{} \# i 3 more variables: payer\_valuation\_change \textless{}dbl\textgreater{}, receiver\_npv \textless{}dbl\textgreater{},}
\CommentTok{\#\textgreater{} \#   receiver\_valuation\_change \textless{}dbl\textgreater{}}
\end{Highlighting}
\end{Shaded}

\texttt{payer\_npv}は固定金利を支払うポジションのNPVであり、\texttt{receiver\_npv}は固定金利を受け取る反対ポジションのNPVである。

\texttt{payer\_valuation\_change}と\texttt{receiver\_valuation\_change}は、それぞれのシナリオのNPVから基準シナリオのNPVを差し引いた評価損益である。

\subsection{3.6.4
基準シナリオの照合}\label{ux57faux6e96ux30b7ux30caux30eaux30aaux306eux7167ux5408}

基準シナリオは、第9節で使用した市場データと約定条件をそのまま使用している。そのため、基準シナリオのNPVとフェアレートは、第9節の計算結果と一致する。

\begin{Shaded}
\begin{Highlighting}[]
\NormalTok{case\_base\_scenario\_tbl }\OtherTok{\textless{}{-}}\NormalTok{ case\_ois\_scenario\_result\_tbl }\SpecialCharTok{|\textgreater{}}
\NormalTok{  dplyr}\SpecialCharTok{::}\FunctionTok{filter}\NormalTok{(}
\NormalTok{    scenario }\SpecialCharTok{==} \StringTok{"基準"}
\NormalTok{  )}

\NormalTok{case\_base\_check\_tbl }\OtherTok{\textless{}{-}}\NormalTok{ tibble}\SpecialCharTok{::}\FunctionTok{tibble}\NormalTok{(}
  \AttributeTok{metric =} \FunctionTok{c}\NormalTok{(}
    \StringTok{"NPV"}\NormalTok{,}
    \StringTok{"Fair rate"}
\NormalTok{  ),}
  \AttributeTok{chapter9\_value =} \FunctionTok{c}\NormalTok{(}
\NormalTok{    case\_ois\_gaussquant\_summary}\SpecialCharTok{$}\NormalTok{npv[[}\DecValTok{1}\NormalTok{L]],}
\NormalTok{    case\_ois\_gaussquant\_summary}\SpecialCharTok{$}\NormalTok{fair\_rate[[}\DecValTok{1}\NormalTok{L]]}
\NormalTok{  ),}
  \AttributeTok{chapter10\_value =} \FunctionTok{c}\NormalTok{(}
\NormalTok{    case\_base\_scenario\_tbl}\SpecialCharTok{$}\NormalTok{payer\_npv[[}\DecValTok{1}\NormalTok{L]],}
\NormalTok{    case\_base\_scenario\_tbl}\SpecialCharTok{$}\NormalTok{fair\_rate[[}\DecValTok{1}\NormalTok{L]]}
\NormalTok{  )}
\NormalTok{) }\SpecialCharTok{|\textgreater{}}
\NormalTok{  dplyr}\SpecialCharTok{::}\FunctionTok{mutate}\NormalTok{(}
    \AttributeTok{difference =}
\NormalTok{      chapter10\_value }\SpecialCharTok{{-}}
\NormalTok{      chapter9\_value}
\NormalTok{  )}

\NormalTok{GaussQuant}\SpecialCharTok{::}\FunctionTok{show\_tbl\_GQH}\NormalTok{(}
\NormalTok{  case\_base\_check\_tbl,}
  \StringTok{"Base{-}scenario consistency check"}\NormalTok{,}
  \AttributeTok{n =} \DecValTok{20}
\NormalTok{)}
\CommentTok{\#\textgreater{} \# A tibble: 2 x 4}
\CommentTok{\#\textgreater{}   metric    chapter9\_value chapter10\_value difference}
\CommentTok{\#\textgreater{}   \textless{}chr\textgreater{}              \textless{}dbl\textgreater{}           \textless{}dbl\textgreater{}      \textless{}dbl\textgreater{}}
\CommentTok{\#\textgreater{} 1 NPV            {-}1.65e{-}13       {-}1.65e{-}13          0}
\CommentTok{\#\textgreater{} 2 Fair rate       4.56e{-} 3        4.56e{-} 3          0}

\FunctionTok{stopifnot}\NormalTok{(}
  \FunctionTok{isTRUE}\NormalTok{(}
    \FunctionTok{all.equal}\NormalTok{(}
\NormalTok{      case\_base\_check\_tbl}\SpecialCharTok{$}\NormalTok{chapter9\_value,}
\NormalTok{      case\_base\_check\_tbl}\SpecialCharTok{$}\NormalTok{chapter10\_value,}
      \AttributeTok{tolerance =} \FloatTok{1e{-}12}
\NormalTok{    )}
\NormalTok{  )}
\NormalTok{)}
\end{Highlighting}
\end{Shaded}

この照合により、第9節の基準評価と、第10節の基準シナリオが同じ市場データ、評価日および約定条件に基づいていることを確認できる。

\subsection{3.6.5
市場レート変動と評価損益}\label{ux5e02ux5834ux30ecux30fcux30c8ux5909ux52d5ux3068ux8a55ux4fa1ux640dux76ca}

市場レートが低下すると、5年OISのフェアレートも低下する。この場合、固定金利支払側は、低下後の市場水準より高い0.456\%を支払うため、既存契約の価値は低下する。

一方、市場レートが上昇すると、固定金利支払側は市場水準より低い0.456\%を支払う契約を保有することになるため、既存契約の価値は上昇する。

固定金利受取側では、この関係が逆になる。固定金利支払側の利益は固定金利受取側の損失となり、両者のNPVと評価損益は符号が反対になる。

\subsection{3.6.6
1ベーシスポイント当たりの評価変化}\label{ux30d9ux30fcux30b7ux30b9ux30ddux30a4ux30f3ux30c8ux5f53ux305fux308aux306eux8a55ux4fa1ux5909ux5316}

上下10ベーシスポイントの結果から、固定金利支払側について、1ベーシスポイント当たりの平均的なNPV変化を求める。

\begin{Shaded}
\begin{Highlighting}[]
\NormalTok{case\_npv\_down }\OtherTok{\textless{}{-}}\NormalTok{ case\_ois\_scenario\_result\_tbl }\SpecialCharTok{|\textgreater{}}
\NormalTok{  dplyr}\SpecialCharTok{::}\FunctionTok{filter}\NormalTok{(}
\NormalTok{    scenario }\SpecialCharTok{==} \StringTok{"10bp低下"}
\NormalTok{  ) }\SpecialCharTok{|\textgreater{}}
\NormalTok{  dplyr}\SpecialCharTok{::}\FunctionTok{pull}\NormalTok{(}
\NormalTok{    payer\_npv}
\NormalTok{  )}

\NormalTok{case\_npv\_up }\OtherTok{\textless{}{-}}\NormalTok{ case\_ois\_scenario\_result\_tbl }\SpecialCharTok{|\textgreater{}}
\NormalTok{  dplyr}\SpecialCharTok{::}\FunctionTok{filter}\NormalTok{(}
\NormalTok{    scenario }\SpecialCharTok{==} \StringTok{"10bp上昇"}
\NormalTok{  ) }\SpecialCharTok{|\textgreater{}}
\NormalTok{  dplyr}\SpecialCharTok{::}\FunctionTok{pull}\NormalTok{(}
\NormalTok{    payer\_npv}
\NormalTok{  )}

\NormalTok{case\_ois\_bpv\_tbl }\OtherTok{\textless{}{-}}\NormalTok{ tibble}\SpecialCharTok{::}\FunctionTok{tibble}\NormalTok{(}
  \AttributeTok{position =} \StringTok{"Fixed{-}rate payer"}\NormalTok{,}
  \AttributeTok{npv\_10bp\_down =}\NormalTok{ case\_npv\_down,}
  \AttributeTok{npv\_base =}\NormalTok{ case\_base\_npv,}
  \AttributeTok{npv\_10bp\_up =}\NormalTok{ case\_npv\_up,}
  \AttributeTok{average\_npv\_change\_per\_bp =}
\NormalTok{    (}
\NormalTok{      case\_npv\_up }\SpecialCharTok{{-}}
\NormalTok{        case\_npv\_down}
\NormalTok{    ) }\SpecialCharTok{/} \DecValTok{20}
\NormalTok{)}

\NormalTok{GaussQuant}\SpecialCharTok{::}\FunctionTok{show\_tbl\_GQH}\NormalTok{(}
\NormalTok{  case\_ois\_bpv\_tbl,}
  \StringTok{"Five{-}year EONIA OIS rate sensitivity"}\NormalTok{,}
  \AttributeTok{n =} \DecValTok{20}
\NormalTok{)}
\CommentTok{\#\textgreater{} \# A tibble: 1 x 5}
\CommentTok{\#\textgreater{}   position         npv\_10bp\_down  npv\_base npv\_10bp\_up average\_npv\_change\_per\_bp}
\CommentTok{\#\textgreater{}   \textless{}chr\textgreater{}                    \textless{}dbl\textgreater{}     \textless{}dbl\textgreater{}       \textless{}dbl\textgreater{}                     \textless{}dbl\textgreater{}}
\CommentTok{\#\textgreater{} 1 Fixed{-}rate payer      {-}0.00505 {-}1.65e{-}13     0.00502                  0.000503}
\end{Highlighting}
\end{Shaded}

\texttt{average\_npv\_change\_per\_bp}は、マイナス10ベーシスポイントからプラス10ベーシスポイントまでのNPV変化を20で除した値である。

本ケーススタディでは名目元本を1としてOISを構築しているため、この値も単位元本当たりの評価変化を表す。実際の取引では、これに取引元本を乗じることによって、ポジション全体の金利感応度を求めることができる。

\section{3.7 まとめ}\label{ux307eux3068ux3081-7}

本章では、世界金融危機とIBOR改革を背景として、金利デリバティブの評価がシングルカーブ方式からマルチカーブ方式へ移行した経緯を確認した。LIBORなどのIBORには銀行の信用リスクや資金流動性リスクが含まれるため、担保付取引の割引金利としてそのまま使用することは適切ではない。そこで、担保契約に対応するOISカーブを割引に用い、参照金利ごとのフォワードカーブから将来金利を予測する評価方法が定着した。

OISでは、固定金利と、計算期間中の日々の翌日物金利を複利で積み上げた変動金利を交換する。評価日以前の翌日物金利には確定したFixingを使用し、評価日以後の金利には期間構造から求めた予測値を使用する。TONAとSOFRはいずれも翌日物金利であるが、その基礎となる市場や担保性は異なるため、対象通貨の市場慣行を正しく設定する必要がある。

マルチカーブ評価では、将来の変動金利を求めるフォワードカーブと、将来キャッシュフローを現在価値へ変換するディスカウントカーブを分離する。同じ将来キャッシュフローであっても、使用する割引カーブが異なればNPVやフェアレートも変化する。したがって、金利スワップの評価では、参照インデックスがどのカーブへ接続され、価格計算エンジンがどのカーブを割引に使用しているかを明確にする必要がある。

ケーススタディでは、期間構造ハンドルを翌日物金利インデックスへ接続し、\texttt{MakeOIS}を用いてOISを構築した。さらに、NPV、フェアレート、固定レグおよび翌日物レグのキャッシュフローを確認し、OISの契約条件と市場カーブが評価結果へどのように反映されるかを示した。

また、契約固定金利を変更せずに市場カーブを上下へ平行シフトし、OISの評価損益を比較した。契約締結後も固定金利は変化しないが、市場カーブの変動によって将来の翌日物金利予測と割引係数が変化するため、OISのNPVは日々変動する。固定金利支払側と固定金利受取側では、同じ市場変動に対する評価損益の符号が反対になる。

実際のリスク管理では、カーブ全体の平行シフトだけでなく、短期、中期および長期の金利を個別に変化させる必要がある。また、Fixingの確定、休日カレンダー、日数計算規則、支払ラグ、予測カーブおよび割引カーブの設定を総合的に確認することが、OISと金利スワップを正しく評価するための基本となる。

\section{参考資料}\label{ux53c2ux8003ux8cc7ux6599}

\begin{itemize}
\tightlist
\item
  Financial Conduct Authority,
  \href{https://www.fca.org.uk/news/speeches/the-future-of-libor}{The
  future of LIBOR}, 27 July 2017.
\item
  Financial Conduct Authority,
  \href{https://www.fca.org.uk/news/press-releases/end-libor}{The end of
  LIBOR}, 1 October 2024.
\item
  Financial Stability Board,
  \href{https://www.fsb.org/2014/07/r_140722/}{Reforming Major Interest
  Rate Benchmarks}, 22 July 2014.
\item
  Bank of Japan,
  \href{https://www.boj.or.jp/en/statistics/market/short/mutan/index.htm}{Uncollateralized
  Overnight Call Rate}.
\end{itemize}

\chapter{---}\label{section}

\chapter{title: ``第II-4章
株価・為替と幾何ブラウン運動モデル''}\label{title-ux7b2cii-4ux7ae0-ux682aux4fa1ux70baux66ffux3068ux5e7eux4f55ux30d6ux30e9ux30a6ux30f3ux904bux52d5ux30e2ux30c7ux30eb}

\chapter{execute:}\label{execute}

\chapter{eval: true}\label{eval-true}

\chapter{echo: true}\label{echo-true}

\chapter{warning: false}\label{warning-false}

\chapter{message: false}\label{message-false}

\chapter{---}\label{section-1}

1973年に公表されたBlack--Scholes--Mertonモデルは、株式オプションの理論価格を求めるための基本モデルとして、金融工学の発展に大きな影響を与えた。このモデルでは、株価が幾何ブラウン運動に従い、将来の株価が対数正規分布に従うと仮定する。

一方、金利先物、債券先物、商品先物およびフォワード金利を原資産とするオプションでは、現物価格よりも満期に対応する先物価格またはフォワード価格を直接扱う方が自然である。Blackは1976年、Black--Scholes--Mertonモデルの考え方を先物オプションへ適用した。現在、この評価方法はBlack-76モデルまたは単にBlackモデルと呼ばれている。

Black--Scholes--Mertonモデルでは、現物価格、無リスク金利および配当利回りから将来のフォワード価格を求める。これに対しBlack-76モデルでは、観測されたフォワード価格または先物価格を直接入力し、満期時のPayoffを割り引いて現在価値を求める。

したがって、Black-76とBlack--Scholes--Mertonは無関係な別のモデルではない。両者は、原資産価格が正の値を取り、その変化率が正規分布に従うという共通の仮定を持つ。現物価格を出発点とするか、フォワード価格を出発点とするかが異なる。

本章では、まずBlack-76による先物・フォワードオプションの解析評価を行い、オプション価格とGreeksを確認する。次に、同じ確率過程からMonte
Carloパスを生成し、満期Payoffの割引期待値が解析解へ近づくことを確認する。

その後、Black--Scholes--Mertonモデルによる株式オプション評価へ戻り、Europeanオプションの解析解と二項Treeを比較する。さらに、Americanオプションの早期行使価値を、近似解、二項TreeおよびLongstaff--Schwartz法によって評価する。

Blackモデルは、原資産価格を正の値に限定する対数正規モデルである。この性質は株式や先物価格の評価には都合がよい一方、金利がゼロまたはマイナスとなる市場では制約となる。次章では、この制約を持たないNormalモデルを取り上げ、Blackモデルとの違いを確認する。

\begin{Shaded}
\begin{Highlighting}[]
\NormalTok{knitr}\SpecialCharTok{::}\NormalTok{opts\_chunk}\SpecialCharTok{$}\FunctionTok{set}\NormalTok{(}
  \AttributeTok{collapse =} \ConstantTok{TRUE}\NormalTok{,}
  \AttributeTok{comment =} \StringTok{"\#\textgreater{}"}\NormalTok{,}
  \AttributeTok{warning =} \ConstantTok{FALSE}\NormalTok{,}
  \AttributeTok{message =} \ConstantTok{FALSE}
\NormalTok{)}

\FunctionTok{suppressPackageStartupMessages}\NormalTok{(\{}
  \FunctionTok{library}\NormalTok{(tidyverse)}
  \FunctionTok{library}\NormalTok{(QuantLib)}
  \FunctionTok{library}\NormalTok{(GaussQuant)}
\NormalTok{\})}
\end{Highlighting}
\end{Shaded}

\section{4.1
Black--Scholes--MertonとBlack-76}\label{blackscholesmertonux3068black-76}

\subsection{4.1.1
Black--Scholes--Mertonモデル}\label{blackscholesmertonux30e2ux30c7ux30eb}

リスク中立測度の下で、株価 \(S_t\)
は次の確率微分方程式に従うと仮定する。

\[
dS_t
=
(r-q)S_t\,dt
+
\sigma S_t\,dW_t
\]

ここで、\(r\) は無リスク金利、\(q\) は配当利回り、\(\sigma\)
はボラティリティ、\(W_t\) はブラウン運動である。

満期 \(T\) のフォワード価格は次式となる。

\[
F_0
=
S_0 e^{(r-q)T}
\]

したがって、Black--Scholes--Mertonモデルは現物価格を出発点として、金利と配当を通じてフォワード価格を内生的に求めるモデルとみることができる。

\subsection{4.1.2 Black-76モデル}\label{black-76ux30e2ux30c7ux30eb}

Black-76では、フォワード価格 \(F_0\) を直接入力する。European
callの価格は次式で与えられる。

\[
C
=
D(0,T)
\left[
F_0N(d_1)-KN(d_2)
\right]
\]

European putの価格は次式となる。

\[
P
=
D(0,T)
\left[
KN(-d_2)-F_0N(-d_1)
\right]
\]

ここで、

\[
d_1
=
\frac{
\log(F_0/K)+\frac{1}{2}\sigma^2T
}{
\sigma\sqrt{T}
}
\]

\[
d_2
=
d_1-\sigma\sqrt{T}
\]

であり、\(K\) は行使価格、\(D(0,T)\) は満期までの割引係数、\(N(\cdot)\)
は標準正規分布の累積分布関数である。

金利が決定論的である場合、同一満期の先物価格とフォワード価格は一致するため、Black-76は先物オプションとフォワードオプションの双方に用いられる。

\subsection{4.1.3
Black-76の入力条件}\label{black-76ux306eux5165ux529bux6761ux4ef6}

本節では、債券先物オプションを意識した価格表示を使用する。フォワード価格は107と32分の6.75、行使価格は107.5、ボラティリティは5.2\%、無リスク金利は5\%とする。

Black-76では、現物価格ではなく、オプション満期に対応するフォワード価格または先物価格を原資産として使用する。

\begin{Shaded}
\begin{Highlighting}[]
\NormalTok{black\_input\_tbl }\OtherTok{\textless{}{-}}\NormalTok{ tibble}\SpecialCharTok{::}\FunctionTok{tribble}\NormalTok{(}
  \SpecialCharTok{\textasciitilde{}}\NormalTok{input,           }\SpecialCharTok{\textasciitilde{}}\NormalTok{value,}
  \StringTok{"trade\_date"}\NormalTok{,     }\StringTok{"2023{-}07{-}21"}\NormalTok{,}
  \StringTok{"maturity\_date"}\NormalTok{,  }\StringTok{"2023{-}08{-}26"}\NormalTok{,}
  \StringTok{"forward\_price"}\NormalTok{,  }\FunctionTok{as.character}\NormalTok{(}\DecValTok{107} \SpecialCharTok{+} \FloatTok{6.75} \SpecialCharTok{/} \DecValTok{32}\NormalTok{),}
  \StringTok{"strike\_price"}\NormalTok{,   }\FunctionTok{as.character}\NormalTok{(}\FloatTok{107.5}\NormalTok{),}
  \StringTok{"volatility"}\NormalTok{,     }\FunctionTok{as.character}\NormalTok{(}\FloatTok{0.052}\NormalTok{),}
  \StringTok{"risk\_free\_rate"}\NormalTok{, }\FunctionTok{as.character}\NormalTok{(}\FloatTok{0.05}\NormalTok{)}
\NormalTok{)}

\NormalTok{trade\_date }\OtherTok{\textless{}{-}}\NormalTok{ GaussQuant}\SpecialCharTok{::}\FunctionTok{date\_GQL}\NormalTok{(}\StringTok{"2023{-}07{-}21"}\NormalTok{)}
\NormalTok{maturity\_date }\OtherTok{\textless{}{-}}\NormalTok{ GaussQuant}\SpecialCharTok{::}\FunctionTok{date\_GQL}\NormalTok{(}\StringTok{"2023{-}08{-}26"}\NormalTok{)}
\NormalTok{forward\_price }\OtherTok{\textless{}{-}} \DecValTok{107} \SpecialCharTok{+} \FloatTok{6.75} \SpecialCharTok{/} \DecValTok{32}
\NormalTok{strike\_price }\OtherTok{\textless{}{-}} \FloatTok{107.5}
\NormalTok{volatility }\OtherTok{\textless{}{-}} \FloatTok{0.052}
\NormalTok{risk\_free\_rate }\OtherTok{\textless{}{-}} \FloatTok{0.05}
\NormalTok{day\_counter }\OtherTok{\textless{}{-}}\NormalTok{ QuantLib}\SpecialCharTok{::}\FunctionTok{Actual365Fixed}\NormalTok{()}

\NormalTok{call\_payoff }\OtherTok{\textless{}{-}}\NormalTok{ QuantLib}\SpecialCharTok{::}\FunctionTok{PlainVanillaPayoff}\NormalTok{(}
\NormalTok{  QuantLib}\SpecialCharTok{::}\FunctionTok{Option\_Call\_get}\NormalTok{(),}
\NormalTok{  strike\_price}
\NormalTok{)}

\NormalTok{GaussQuant}\SpecialCharTok{::}\FunctionTok{set\_eval\_date\_GQL}\NormalTok{(trade\_date)}

\NormalTok{maturity }\OtherTok{\textless{}{-}}\NormalTok{ day\_counter}\SpecialCharTok{$}\FunctionTok{yearFraction}\NormalTok{(}
\NormalTok{  trade\_date,}
\NormalTok{  maturity\_date}
\NormalTok{)}

\NormalTok{GaussQuant}\SpecialCharTok{::}\FunctionTok{show\_tbl\_GQH}\NormalTok{(}
\NormalTok{  black\_input\_tbl,}
  \StringTok{"Black{-}76 input assumptions"}\NormalTok{,}
  \AttributeTok{n =} \DecValTok{20}
\NormalTok{)}
\CommentTok{\#\textgreater{} \# A tibble: 6 x 2}
\CommentTok{\#\textgreater{}   input          value      }
\CommentTok{\#\textgreater{}   \textless{}chr\textgreater{}          \textless{}chr\textgreater{}      }
\CommentTok{\#\textgreater{} 1 trade\_date     2023{-}07{-}21 }
\CommentTok{\#\textgreater{} 2 maturity\_date  2023{-}08{-}26 }
\CommentTok{\#\textgreater{} 3 forward\_price  107.2109375}
\CommentTok{\#\textgreater{} 4 strike\_price   107.5      }
\CommentTok{\#\textgreater{} 5 volatility     0.052      }
\CommentTok{\#\textgreater{} 6 risk\_free\_rate 0.05}
\end{Highlighting}
\end{Shaded}

\section{4.2
Black-76による解析評価}\label{black-76ux306bux3088ux308bux89e3ux6790ux8a55ux4fa1}

\subsection{4.2.1
Black公式による価格とGreeks}\label{blackux516cux5f0fux306bux3088ux308bux4fa1ux683cux3068greeks}

\texttt{BlackCalculator}は、フォワード価格、行使価格、総標準偏差および割引係数から、オプション価格とGreeksを直接計算する。

QuantLibへ渡される総標準偏差は、単なる年率ボラティリティではなく、

\[
\sigma\sqrt{T}
\]

である。

本節では、次のGreeksを確認する。

\begin{itemize}
\tightlist
\item
  \textbf{Delta}：フォワード価格の変化に対するオプション価格の感応度
\item
  \textbf{Gamma}：フォワード価格の変化に対するDeltaの感応度
\item
  \textbf{Vega}：ボラティリティの変化に対するオプション価格の感応度
\item
  \textbf{Theta}：時間経過に対するオプション価格の感応度
\end{itemize}

Black-76のDeltaはフォワード価格に対するDeltaであり、株式オプションのSpot
Deltaとは区別する必要がある。

\begin{Shaded}
\begin{Highlighting}[]
\NormalTok{black\_bundle }\OtherTok{\textless{}{-}}\NormalTok{ GaussQuant}\SpecialCharTok{::}\FunctionTok{black\_process\_bundle\_GQL}\NormalTok{(}
  \AttributeTok{valuation\_date =}\NormalTok{ trade\_date,}
  \AttributeTok{forward =}\NormalTok{ forward\_price,}
  \AttributeTok{risk\_free\_rate =}\NormalTok{ risk\_free\_rate,}
  \AttributeTok{volatility =}\NormalTok{ volatility,}
  \AttributeTok{day\_counter =}\NormalTok{ day\_counter}
\NormalTok{)}

\NormalTok{discount\_factor }\OtherTok{\textless{}{-}}\NormalTok{ black\_bundle}\SpecialCharTok{$}\NormalTok{discount\_curve}\SpecialCharTok{$}\FunctionTok{discount}\NormalTok{(}
\NormalTok{  maturity\_date}
\NormalTok{)}

\NormalTok{black\_calculator\_tbl }\OtherTok{\textless{}{-}}\NormalTok{ GaussQuant}\SpecialCharTok{::}\FunctionTok{black\_calculator\_table\_GQL}\NormalTok{(}
  \AttributeTok{payoff =}\NormalTok{ call\_payoff,}
  \AttributeTok{forward =}\NormalTok{ forward\_price,}
  \AttributeTok{volatility =}\NormalTok{ volatility,}
  \AttributeTok{maturity =}\NormalTok{ maturity,}
  \AttributeTok{discount\_factor =}\NormalTok{ discount\_factor}
\NormalTok{)}

\NormalTok{GaussQuant}\SpecialCharTok{::}\FunctionTok{show\_tbl\_GQH}\NormalTok{(}
\NormalTok{  black\_calculator\_tbl,}
  \StringTok{"Black calculator summary"}\NormalTok{,}
  \AttributeTok{n =} \DecValTok{20}
\NormalTok{)}
\CommentTok{\#\textgreater{} \# A tibble: 6 x 2}
\CommentTok{\#\textgreater{}   metric           value}
\CommentTok{\#\textgreater{}   \textless{}chr\textgreater{}            \textless{}dbl\textgreater{}}
\CommentTok{\#\textgreater{} 1 npv            0.562  }
\CommentTok{\#\textgreater{} 2 delta          0.436  }
\CommentTok{\#\textgreater{} 3 gamma          0.224  }
\CommentTok{\#\textgreater{} 4 vega          13.2    }
\CommentTok{\#\textgreater{} 5 theta         {-}3.45   }
\CommentTok{\#\textgreater{} 6 theta\_per\_day {-}0.00946}
\end{Highlighting}
\end{Shaded}

\subsection{4.2.2 Black過程によるEuropean
option評価}\label{blackux904eux7a0bux306bux3088ux308beuropean-optionux8a55ux4fa1}

前節ではBlack公式を直接使用した。本節では、同じ入力条件からBlack
processを構築し、European optionへ解析的Pricing Engineを設定する。

数式を直接評価する方法と、確率過程およびPricing
Engineを通じて評価する方法が、同一の理論価格を与えることを確認する。

\begin{Shaded}
\begin{Highlighting}[]
\NormalTok{european\_call }\OtherTok{\textless{}{-}}\NormalTok{ QuantLib}\SpecialCharTok{::}\FunctionTok{VanillaOption}\NormalTok{(}
\NormalTok{  call\_payoff,}
\NormalTok{  QuantLib}\SpecialCharTok{::}\FunctionTok{EuropeanExercise}\NormalTok{(maturity\_date)}
\NormalTok{)}

\NormalTok{analytic\_engine }\OtherTok{\textless{}{-}}\NormalTok{ QuantLib}\SpecialCharTok{::}\FunctionTok{AnalyticEuropeanEngine}\NormalTok{(}
\NormalTok{  black\_bundle}\SpecialCharTok{$}\NormalTok{process}
\NormalTok{)}

\NormalTok{european\_call}\SpecialCharTok{$}\FunctionTok{setPricingEngine}\NormalTok{(}
\NormalTok{  analytic\_engine}
\NormalTok{)}
\CommentTok{\#\textgreater{} NULL}

\NormalTok{analytic\_european\_tbl }\OtherTok{\textless{}{-}}\NormalTok{ GaussQuant}\SpecialCharTok{::}\FunctionTok{option\_metric\_table\_GQL}\NormalTok{(}
  \AttributeTok{option =}\NormalTok{ european\_call,}
  \AttributeTok{process =}\NormalTok{ black\_bundle}\SpecialCharTok{$}\NormalTok{process}
\NormalTok{)}

\NormalTok{analytic\_npv }\OtherTok{\textless{}{-}}\NormalTok{ analytic\_european\_tbl }\SpecialCharTok{|\textgreater{}}
\NormalTok{  dplyr}\SpecialCharTok{::}\FunctionTok{filter}\NormalTok{(.data}\SpecialCharTok{$}\NormalTok{metric }\SpecialCharTok{==} \StringTok{"npv"}\NormalTok{) }\SpecialCharTok{|\textgreater{}}
\NormalTok{  dplyr}\SpecialCharTok{::}\FunctionTok{pull}\NormalTok{(.data}\SpecialCharTok{$}\NormalTok{value)}

\NormalTok{GaussQuant}\SpecialCharTok{::}\FunctionTok{show\_tbl\_GQH}\NormalTok{(}
\NormalTok{  analytic\_european\_tbl,}
  \StringTok{"Analytic European call"}\NormalTok{,}
  \AttributeTok{n =} \DecValTok{20}
\NormalTok{)}
\CommentTok{\#\textgreater{} \# A tibble: 7 x 2}
\CommentTok{\#\textgreater{}   metric                value}
\CommentTok{\#\textgreater{}   \textless{}chr\textgreater{}                 \textless{}dbl\textgreater{}}
\CommentTok{\#\textgreater{} 1 npv                 0.562  }
\CommentTok{\#\textgreater{} 2 delta               0.436  }
\CommentTok{\#\textgreater{} 3 gamma               0.224  }
\CommentTok{\#\textgreater{} 4 vega               13.2    }
\CommentTok{\#\textgreater{} 5 theta              {-}3.45   }
\CommentTok{\#\textgreater{} 6 theta\_per\_day      {-}0.00946}
\CommentTok{\#\textgreater{} 7 implied\_volatility  0.0520}
\end{Highlighting}
\end{Shaded}

\section{4.3 Black--Scholes--MertonとEuropean
option}\label{blackscholesmertonux3068european-option}

\subsection{4.3.1
Black--Scholes--Mertonモデルの入力条件}\label{blackscholesmertonux30e2ux30c7ux30ebux306eux5165ux529bux6761ux4ef6}

ここからは現物株価を原資産とするBlack--Scholes--Mertonモデルを扱う。現物価格100、行使価格100、満期1年、ボラティリティ20\%、無リスク金利5\%、配当利回り0\%とする。

Black--Scholes--Merton
processでは、現物価格に加えて、無リスク金利カーブ、配当利回りカーブおよびボラティリティカーブを設定する。

\begin{Shaded}
\begin{Highlighting}[]
\NormalTok{bsm\_input\_tbl }\OtherTok{\textless{}{-}}\NormalTok{ tibble}\SpecialCharTok{::}\FunctionTok{tribble}\NormalTok{(}
  \SpecialCharTok{\textasciitilde{}}\NormalTok{input,           }\SpecialCharTok{\textasciitilde{}}\NormalTok{value,}
  \StringTok{"trade\_date"}\NormalTok{,     }\StringTok{"2024{-}03{-}01"}\NormalTok{,}
  \StringTok{"maturity"}\NormalTok{,       }\StringTok{"1Y"}\NormalTok{,}
  \StringTok{"spot\_price"}\NormalTok{,     }\FunctionTok{as.character}\NormalTok{(}\DecValTok{100}\NormalTok{),}
  \StringTok{"strike\_price"}\NormalTok{,   }\FunctionTok{as.character}\NormalTok{(}\DecValTok{100}\NormalTok{),}
  \StringTok{"volatility"}\NormalTok{,     }\FunctionTok{as.character}\NormalTok{(}\FloatTok{0.20}\NormalTok{),}
  \StringTok{"risk\_free\_rate"}\NormalTok{, }\FunctionTok{as.character}\NormalTok{(}\FloatTok{0.05}\NormalTok{),}
  \StringTok{"dividend\_rate"}\NormalTok{,  }\FunctionTok{as.character}\NormalTok{(}\DecValTok{0}\NormalTok{)}
\NormalTok{)}

\NormalTok{trade\_date\_bsm }\OtherTok{\textless{}{-}}\NormalTok{ GaussQuant}\SpecialCharTok{::}\FunctionTok{date\_GQL}\NormalTok{(}\StringTok{"2024{-}03{-}01"}\NormalTok{)}
\NormalTok{calendar\_bsm }\OtherTok{\textless{}{-}}\NormalTok{ QuantLib}\SpecialCharTok{::}\FunctionTok{TARGET}\NormalTok{()}

\NormalTok{maturity\_date\_bsm }\OtherTok{\textless{}{-}}\NormalTok{ calendar\_bsm}\SpecialCharTok{$}\FunctionTok{advance}\NormalTok{(}
\NormalTok{  trade\_date\_bsm,}
  \DecValTok{1}\NormalTok{L,}
  \StringTok{"Years"}
\NormalTok{)}

\NormalTok{spot\_price\_bsm }\OtherTok{\textless{}{-}} \DecValTok{100}
\NormalTok{strike\_price\_bsm }\OtherTok{\textless{}{-}} \DecValTok{100}
\NormalTok{volatility\_bsm }\OtherTok{\textless{}{-}} \FloatTok{0.20}
\NormalTok{risk\_free\_rate\_bsm }\OtherTok{\textless{}{-}} \FloatTok{0.05}
\NormalTok{dividend\_rate\_bsm }\OtherTok{\textless{}{-}} \DecValTok{0}

\NormalTok{GaussQuant}\SpecialCharTok{::}\FunctionTok{set\_eval\_date\_GQL}\NormalTok{(}
\NormalTok{  trade\_date\_bsm}
\NormalTok{)}

\NormalTok{bsm\_bundle }\OtherTok{\textless{}{-}}\NormalTok{ GaussQuant}\SpecialCharTok{::}\FunctionTok{bsm\_process\_bundle\_GQL}\NormalTok{(}
  \AttributeTok{valuation\_date =}\NormalTok{ trade\_date\_bsm,}
  \AttributeTok{spot =}\NormalTok{ spot\_price\_bsm,}
  \AttributeTok{risk\_free\_rate =}\NormalTok{ risk\_free\_rate\_bsm,}
  \AttributeTok{dividend\_rate =}\NormalTok{ dividend\_rate\_bsm,}
  \AttributeTok{volatility =}\NormalTok{ volatility\_bsm,}
  \AttributeTok{day\_counter =}\NormalTok{ day\_counter,}
  \AttributeTok{calendar =}\NormalTok{ calendar\_bsm}
\NormalTok{)}

\NormalTok{GaussQuant}\SpecialCharTok{::}\FunctionTok{show\_tbl\_GQH}\NormalTok{(}
\NormalTok{  bsm\_input\_tbl,}
  \StringTok{"Black{-}Scholes{-}Merton assumptions"}\NormalTok{,}
  \AttributeTok{n =} \DecValTok{20}
\NormalTok{)}
\CommentTok{\#\textgreater{} \# A tibble: 7 x 2}
\CommentTok{\#\textgreater{}   input          value     }
\CommentTok{\#\textgreater{}   \textless{}chr\textgreater{}          \textless{}chr\textgreater{}     }
\CommentTok{\#\textgreater{} 1 trade\_date     2024{-}03{-}01}
\CommentTok{\#\textgreater{} 2 maturity       1Y        }
\CommentTok{\#\textgreater{} 3 spot\_price     100       }
\CommentTok{\#\textgreater{} 4 strike\_price   100       }
\CommentTok{\#\textgreater{} 5 volatility     0.2       }
\CommentTok{\#\textgreater{} 6 risk\_free\_rate 0.05      }
\CommentTok{\#\textgreater{} 7 dividend\_rate  0}
\end{Highlighting}
\end{Shaded}

\subsection{4.3.2 European
putの解析解と二項Tree}\label{european-putux306eux89e3ux6790ux89e3ux3068ux4e8cux9805tree}

Europeanオプションは満期時にのみ行使できるため、Black--Scholes--Mertonモデルの解析解を利用できる。

二項Treeでは、連続時間の価格過程を有限個の時間ステップへ分割する。ステップ数を増やすほど、理論上は連続時間モデルの解析値へ近づく。

ここでは、25、100および500ステップのCox--Ross--Rubinstein
Treeを解析解と比較する。

\begin{Shaded}
\begin{Highlighting}[]
\NormalTok{european\_put }\OtherTok{\textless{}{-}}\NormalTok{ GaussQuant}\SpecialCharTok{::}\FunctionTok{make\_european\_option\_GQL}\NormalTok{(}
  \AttributeTok{spot =}\NormalTok{ spot\_price\_bsm,}
  \AttributeTok{strike =}\NormalTok{ strike\_price\_bsm,}
  \AttributeTok{maturity\_date =}\NormalTok{ maturity\_date\_bsm,}
  \AttributeTok{option\_type =} \StringTok{"put"}\NormalTok{,}
  \AttributeTok{valuation\_date =}\NormalTok{ trade\_date\_bsm,}
  \AttributeTok{risk\_free\_rate =}\NormalTok{ risk\_free\_rate\_bsm,}
  \AttributeTok{dividend\_yield =}\NormalTok{ dividend\_rate\_bsm,}
  \AttributeTok{volatility =}\NormalTok{ volatility\_bsm,}
  \AttributeTok{day\_counter =}\NormalTok{ day\_counter,}
  \AttributeTok{calendar =}\NormalTok{ calendar\_bsm}
\NormalTok{)}

\NormalTok{european\_engines }\OtherTok{\textless{}{-}} \FunctionTok{list}\NormalTok{(}
  \AttributeTok{crr\_25 =}\NormalTok{ QuantLib}\SpecialCharTok{::}\FunctionTok{BinomialCRRVanillaEngine}\NormalTok{(}
\NormalTok{    bsm\_bundle}\SpecialCharTok{$}\NormalTok{process,}
    \DecValTok{25}\NormalTok{L}
\NormalTok{  ),}
  \AttributeTok{analytic =}\NormalTok{ QuantLib}\SpecialCharTok{::}\FunctionTok{AnalyticEuropeanEngine}\NormalTok{(}
\NormalTok{    bsm\_bundle}\SpecialCharTok{$}\NormalTok{process}
\NormalTok{  ),}
  \AttributeTok{crr\_100 =}\NormalTok{ QuantLib}\SpecialCharTok{::}\FunctionTok{BinomialCRRVanillaEngine}\NormalTok{(}
\NormalTok{    bsm\_bundle}\SpecialCharTok{$}\NormalTok{process,}
    \DecValTok{100}\NormalTok{L}
\NormalTok{  ),}
  \AttributeTok{crr\_500 =}\NormalTok{ QuantLib}\SpecialCharTok{::}\FunctionTok{BinomialCRRVanillaEngine}\NormalTok{(}
\NormalTok{    bsm\_bundle}\SpecialCharTok{$}\NormalTok{process,}
    \DecValTok{500}\NormalTok{L}
\NormalTok{  )}
\NormalTok{)}

\NormalTok{european\_comparison\_tbl }\OtherTok{\textless{}{-}}\NormalTok{ purrr}\SpecialCharTok{::}\FunctionTok{imap\_dfr}\NormalTok{(}
\NormalTok{  european\_engines,}
  \ControlFlowTok{function}\NormalTok{(engine, method) \{}
\NormalTok{    european\_put}\SpecialCharTok{$}\FunctionTok{setPricingEngine}\NormalTok{(engine)}

\NormalTok{    tibble}\SpecialCharTok{::}\FunctionTok{tibble}\NormalTok{(}
      \AttributeTok{method =}\NormalTok{ method,}
      \AttributeTok{npv =}\NormalTok{ european\_put}\SpecialCharTok{$}\FunctionTok{NPV}\NormalTok{()}
\NormalTok{    )}
\NormalTok{  \}}
\NormalTok{)}

\NormalTok{analytic\_put\_npv }\OtherTok{\textless{}{-}}\NormalTok{ european\_comparison\_tbl }\SpecialCharTok{|\textgreater{}}
\NormalTok{  dplyr}\SpecialCharTok{::}\FunctionTok{filter}\NormalTok{(.data}\SpecialCharTok{$}\NormalTok{method }\SpecialCharTok{==} \StringTok{"analytic"}\NormalTok{) }\SpecialCharTok{|\textgreater{}}
\NormalTok{  dplyr}\SpecialCharTok{::}\FunctionTok{pull}\NormalTok{(.data}\SpecialCharTok{$}\NormalTok{npv)}

\NormalTok{european\_comparison\_tbl }\OtherTok{\textless{}{-}}\NormalTok{ european\_comparison\_tbl }\SpecialCharTok{|\textgreater{}}
\NormalTok{  dplyr}\SpecialCharTok{::}\FunctionTok{mutate}\NormalTok{(}
    \AttributeTok{difference\_from\_analytic =}\NormalTok{ .data}\SpecialCharTok{$}\NormalTok{npv }\SpecialCharTok{{-}}\NormalTok{ analytic\_put\_npv}
\NormalTok{  )}

\NormalTok{GaussQuant}\SpecialCharTok{::}\FunctionTok{show\_tbl\_GQH}\NormalTok{(}
\NormalTok{  european\_comparison\_tbl,}
  \StringTok{"European put engine comparison"}\NormalTok{,}
  \AttributeTok{n =} \DecValTok{20}
\NormalTok{)}
\CommentTok{\#\textgreater{} \# A tibble: 4 x 3}
\CommentTok{\#\textgreater{}   method     npv difference\_from\_analytic}
\CommentTok{\#\textgreater{}   \textless{}chr\textgreater{}    \textless{}dbl\textgreater{}                    \textless{}dbl\textgreater{}}
\CommentTok{\#\textgreater{} 1 crr\_25    5.65                  0.0720 }
\CommentTok{\#\textgreater{} 2 analytic  5.58                  0      }
\CommentTok{\#\textgreater{} 3 crr\_100   5.56                 {-}0.0197 }
\CommentTok{\#\textgreater{} 4 crr\_500   5.58                 {-}0.00394}
\end{Highlighting}
\end{Shaded}

\section{4.4 Monte Carlo評価}\label{monte-carloux8a55ux4fa1}

\subsection{4.4.1 擬似乱数Monte
Carlo評価}\label{ux64ecux4f3cux4e71ux6570monte-carloux8a55ux4fa1}

リスク中立評価では、オプション価格は満期Payoffの割引期待値として表される。

\[
V_0
=
D(0,T)
E^Q
\left[
\max(F_T-K,0)
\right]
\]

Monte
Carlo法では、多数の将来パスを生成し、各パスの満期Payoffを求め、その平均値を割り引く。

European callには解析解があるため、ここでMonte
Carlo法を使う目的は、解析解に代わる高度な方法を示すことではない。確率過程から生成したパスの期待値がBlack公式の価格へ近づくことを確認し、後に解析解を持たない商品へ応用する準備をすることにある。

\begin{Shaded}
\begin{Highlighting}[]
\NormalTok{mc\_setting\_tbl }\OtherTok{\textless{}{-}}\NormalTok{ tibble}\SpecialCharTok{::}\FunctionTok{tribble}\NormalTok{(}
  \SpecialCharTok{\textasciitilde{}}\NormalTok{setting, }\SpecialCharTok{\textasciitilde{}}\NormalTok{value,}
  \StringTok{"steps"}\NormalTok{,  }\DecValTok{12}\NormalTok{,}
  \StringTok{"paths"}\NormalTok{,  }\DecValTok{20000}\NormalTok{,}
  \StringTok{"seed"}\NormalTok{,   }\DecValTok{1}
\NormalTok{)}

\NormalTok{mc\_path\_tbl }\OtherTok{\textless{}{-}}\NormalTok{ GaussQuant}\SpecialCharTok{::}\FunctionTok{generate\_path\_table\_GQL}\NormalTok{(}
  \AttributeTok{process =}\NormalTok{ black\_bundle}\SpecialCharTok{$}\NormalTok{process,}
  \AttributeTok{maturity =}\NormalTok{ maturity,}
  \AttributeTok{n\_steps =} \DecValTok{12}\NormalTok{L,}
  \AttributeTok{n\_paths =} \DecValTok{20000}\NormalTok{L,}
  \AttributeTok{seed =} \DecValTok{1}\NormalTok{L,}
  \AttributeTok{sequence =} \StringTok{"pseudo"}
\NormalTok{)}

\NormalTok{mc\_result }\OtherTok{\textless{}{-}}\NormalTok{ GaussQuant}\SpecialCharTok{::}\FunctionTok{terminal\_option\_mc\_GQL}\NormalTok{(}
  \AttributeTok{path\_tbl =}\NormalTok{ mc\_path\_tbl,}
  \AttributeTok{strike =}\NormalTok{ strike\_price,}
  \AttributeTok{discount\_factor =}\NormalTok{ discount\_factor,}
  \AttributeTok{option\_type =} \StringTok{"call"}
\NormalTok{)}

\NormalTok{mc\_comparison\_tbl }\OtherTok{\textless{}{-}}\NormalTok{ tibble}\SpecialCharTok{::}\FunctionTok{tribble}\NormalTok{(}
  \SpecialCharTok{\textasciitilde{}}\NormalTok{method,                     }\SpecialCharTok{\textasciitilde{}}\NormalTok{npv,                       }\SpecialCharTok{\textasciitilde{}}\NormalTok{standard\_error,}
  \StringTok{"Black analytic"}\NormalTok{,            analytic\_npv,               }\ConstantTok{NA\_real\_}\NormalTok{,}
  \StringTok{"Pseudo{-}random Monte Carlo"}\NormalTok{, mc\_result}\SpecialCharTok{$}\NormalTok{summary}\SpecialCharTok{$}\NormalTok{npv[[}\DecValTok{1}\NormalTok{]], mc\_result}\SpecialCharTok{$}\NormalTok{summary}\SpecialCharTok{$}\NormalTok{standard\_error[[}\DecValTok{1}\NormalTok{]]}
\NormalTok{) }\SpecialCharTok{|\textgreater{}}
\NormalTok{  dplyr}\SpecialCharTok{::}\FunctionTok{mutate}\NormalTok{(}
    \AttributeTok{difference\_from\_analytic =}\NormalTok{ .data}\SpecialCharTok{$}\NormalTok{npv }\SpecialCharTok{{-}}\NormalTok{ analytic\_npv}
\NormalTok{  )}

\NormalTok{GaussQuant}\SpecialCharTok{::}\FunctionTok{show\_tbl\_GQH}\NormalTok{(}
\NormalTok{  mc\_setting\_tbl,}
  \StringTok{"Monte Carlo settings"}\NormalTok{,}
  \AttributeTok{n =} \DecValTok{20}
\NormalTok{)}
\CommentTok{\#\textgreater{} \# A tibble: 3 x 2}
\CommentTok{\#\textgreater{}   setting value}
\CommentTok{\#\textgreater{}   \textless{}chr\textgreater{}   \textless{}dbl\textgreater{}}
\CommentTok{\#\textgreater{} 1 steps      12}
\CommentTok{\#\textgreater{} 2 paths   20000}
\CommentTok{\#\textgreater{} 3 seed        1}

\NormalTok{GaussQuant}\SpecialCharTok{::}\FunctionTok{show\_tbl\_GQH}\NormalTok{(}
\NormalTok{  mc\_comparison\_tbl,}
  \StringTok{"Analytic and Monte Carlo comparison"}\NormalTok{,}
  \AttributeTok{n =} \DecValTok{20}
\NormalTok{)}
\CommentTok{\#\textgreater{} \# A tibble: 2 x 4}
\CommentTok{\#\textgreater{}   method                      npv standard\_error difference\_from\_analytic}
\CommentTok{\#\textgreater{}   \textless{}chr\textgreater{}                     \textless{}dbl\textgreater{}          \textless{}dbl\textgreater{}                    \textless{}dbl\textgreater{}}
\CommentTok{\#\textgreater{} 1 Black analytic            0.562       NA                        0      }
\CommentTok{\#\textgreater{} 2 Pseudo{-}random Monte Carlo 0.560        0.00654                 {-}0.00163}
\end{Highlighting}
\end{Shaded}

標準誤差は、有限本数のパスを用いたことによるMonte
Carlo推定値の不確実性を表す。解析値との差が標準誤差と同程度であれば、シミュレーション結果は統計的に整合的と判断できる。

\subsection{4.4.2 Monte
Carloパスの可視化}\label{monte-carloux30d1ux30b9ux306eux53efux8996ux5316}

Black
processでは、フォワード価格は対数正規分布に従う。したがって、各パスの価格は常に正であり、満期価格の分布は左右対称にはならない。

次のグラフでは、前節と同じオプション満期までの120本のパスと、その満期価格分布を表示する。

\begin{Shaded}
\begin{Highlighting}[]
\NormalTok{path\_plot\_tbl }\OtherTok{\textless{}{-}}\NormalTok{ GaussQuant}\SpecialCharTok{::}\FunctionTok{generate\_path\_table\_GQL}\NormalTok{(}
  \AttributeTok{process =}\NormalTok{ black\_bundle}\SpecialCharTok{$}\NormalTok{process,}
  \AttributeTok{maturity =}\NormalTok{ maturity,}
  \AttributeTok{n\_steps =} \DecValTok{12}\NormalTok{L,}
  \AttributeTok{n\_paths =} \DecValTok{120}\NormalTok{L,}
  \AttributeTok{seed =} \DecValTok{1}\NormalTok{L,}
  \AttributeTok{sequence =} \StringTok{"pseudo"}
\NormalTok{)}

\NormalTok{path\_plot }\OtherTok{\textless{}{-}}\NormalTok{ ggplot2}\SpecialCharTok{::}\FunctionTok{ggplot}\NormalTok{(}
\NormalTok{  path\_plot\_tbl,}
\NormalTok{  ggplot2}\SpecialCharTok{::}\FunctionTok{aes}\NormalTok{(}
    \AttributeTok{x =}\NormalTok{ .data}\SpecialCharTok{$}\NormalTok{time,}
    \AttributeTok{y =}\NormalTok{ .data}\SpecialCharTok{$}\NormalTok{price,}
    \AttributeTok{group =}\NormalTok{ .data}\SpecialCharTok{$}\NormalTok{path\_id}
\NormalTok{  )}
\NormalTok{) }\SpecialCharTok{+}
\NormalTok{  ggplot2}\SpecialCharTok{::}\FunctionTok{geom\_line}\NormalTok{(}
    \AttributeTok{alpha =} \FloatTok{0.45}\NormalTok{,}
    \AttributeTok{linewidth =} \FloatTok{0.4}
\NormalTok{  ) }\SpecialCharTok{+}
\NormalTok{  ggplot2}\SpecialCharTok{::}\FunctionTok{labs}\NormalTok{(}
    \AttributeTok{title =} \StringTok{"Black model Monte Carlo paths"}\NormalTok{,}
    \AttributeTok{x =} \StringTok{"Year"}\NormalTok{,}
    \AttributeTok{y =} \StringTok{"Forward price"}
\NormalTok{  ) }\SpecialCharTok{+}
\NormalTok{  ggplot2}\SpecialCharTok{::}\FunctionTok{theme\_minimal}\NormalTok{()}

\NormalTok{terminal\_plot\_tbl }\OtherTok{\textless{}{-}}\NormalTok{ path\_plot\_tbl }\SpecialCharTok{|\textgreater{}}
\NormalTok{  dplyr}\SpecialCharTok{::}\FunctionTok{group\_by}\NormalTok{(.data}\SpecialCharTok{$}\NormalTok{path\_id) }\SpecialCharTok{|\textgreater{}}
\NormalTok{  dplyr}\SpecialCharTok{::}\FunctionTok{slice\_max}\NormalTok{(}
    \AttributeTok{order\_by =}\NormalTok{ .data}\SpecialCharTok{$}\NormalTok{step,}
    \AttributeTok{n =} \DecValTok{1}\NormalTok{L,}
    \AttributeTok{with\_ties =} \ConstantTok{FALSE}
\NormalTok{  ) }\SpecialCharTok{|\textgreater{}}
\NormalTok{  dplyr}\SpecialCharTok{::}\FunctionTok{ungroup}\NormalTok{()}

\NormalTok{terminal\_plot }\OtherTok{\textless{}{-}}\NormalTok{ ggplot2}\SpecialCharTok{::}\FunctionTok{ggplot}\NormalTok{(}
\NormalTok{  terminal\_plot\_tbl,}
\NormalTok{  ggplot2}\SpecialCharTok{::}\FunctionTok{aes}\NormalTok{(}\AttributeTok{x =}\NormalTok{ .data}\SpecialCharTok{$}\NormalTok{price)}
\NormalTok{) }\SpecialCharTok{+}
\NormalTok{  ggplot2}\SpecialCharTok{::}\FunctionTok{geom\_histogram}\NormalTok{(}
    \AttributeTok{bins =} \DecValTok{24}
\NormalTok{  ) }\SpecialCharTok{+}
\NormalTok{  ggplot2}\SpecialCharTok{::}\FunctionTok{labs}\NormalTok{(}
    \AttributeTok{title =} \StringTok{"Terminal forward{-}price distribution"}\NormalTok{,}
    \AttributeTok{x =} \StringTok{"Terminal price"}\NormalTok{,}
    \AttributeTok{y =} \StringTok{"Count"}
\NormalTok{  ) }\SpecialCharTok{+}
\NormalTok{  ggplot2}\SpecialCharTok{::}\FunctionTok{theme\_minimal}\NormalTok{()}

\FunctionTok{print}\NormalTok{(path\_plot)}
\end{Highlighting}
\end{Shaded}

\pandocbounded{\includegraphics[keepaspectratio]{part2_quantive_riemann/ch04_black_model_files/figure-pdf/chapter04-path-plot-1.pdf}}

\begin{Shaded}
\begin{Highlighting}[]
\FunctionTok{print}\NormalTok{(terminal\_plot)}
\end{Highlighting}
\end{Shaded}

\pandocbounded{\includegraphics[keepaspectratio]{part2_quantive_riemann/ch04_black_model_files/figure-pdf/chapter04-path-plot-2.pdf}}

\subsection{4.4.3
Black過程の1パス}\label{blackux904eux7a0bux306e1ux30d1ux30b9}

Black processのフォワード価格は、時間幅 \(\Delta t\)
ごとに次式で更新できる。

\[
F_{t+\Delta t}
=
F_t
\exp
\left[
-\frac{1}{2}\sigma^2\Delta t
+
\sigma\sqrt{\Delta t}Z
\right]
\]

ここで、\(Z\)
は標準正規乱数である。フォワード測度の下ではフォワード価格のドリフトがゼロであるため、指数部にはIto補正項である
\(-\frac{1}{2}\sigma^2\Delta t\) が現れる。

QuantLibのPathGeneratorが使用するGaussian
incrementを取り出し、同じ乱数から手計算したパスと比較する。

\begin{Shaded}
\begin{Highlighting}[]
\NormalTok{n\_steps\_check }\OtherTok{\textless{}{-}} \DecValTok{3}\NormalTok{L}
\NormalTok{seed\_check }\OtherTok{\textless{}{-}} \DecValTok{1}\NormalTok{L}

\NormalTok{generated\_path\_tbl }\OtherTok{\textless{}{-}}\NormalTok{ GaussQuant}\SpecialCharTok{::}\FunctionTok{generate\_path\_table\_GQL}\NormalTok{(}
  \AttributeTok{process =}\NormalTok{ black\_bundle}\SpecialCharTok{$}\NormalTok{process,}
  \AttributeTok{maturity =}\NormalTok{ maturity,}
  \AttributeTok{n\_steps =}\NormalTok{ n\_steps\_check,}
  \AttributeTok{n\_paths =} \DecValTok{1}\NormalTok{L,}
  \AttributeTok{seed =}\NormalTok{ seed\_check,}
  \AttributeTok{sequence =} \StringTok{"pseudo"}
\NormalTok{)}

\NormalTok{gaussian\_rng\_check }\OtherTok{\textless{}{-}}\NormalTok{ QuantLib}\SpecialCharTok{::}\FunctionTok{GaussianRandomSequenceGenerator}\NormalTok{(}
\NormalTok{  QuantLib}\SpecialCharTok{::}\FunctionTok{UniformRandomSequenceGenerator}\NormalTok{(}
\NormalTok{    n\_steps\_check,}
\NormalTok{    QuantLib}\SpecialCharTok{::}\FunctionTok{UniformRandomGenerator}\NormalTok{(seed\_check)}
\NormalTok{  )}
\NormalTok{)}

\NormalTok{gaussian\_increment }\OtherTok{\textless{}{-}}\NormalTok{ gaussian\_rng\_check}\SpecialCharTok{$}\FunctionTok{nextSequence}\NormalTok{()}\SpecialCharTok{$}\FunctionTok{value}\NormalTok{() }\SpecialCharTok{|\textgreater{}}
  \FunctionTok{as.numeric}\NormalTok{()}

\NormalTok{time\_grid\_check }\OtherTok{\textless{}{-}}\NormalTok{ generated\_path\_tbl }\SpecialCharTok{|\textgreater{}}
\NormalTok{  dplyr}\SpecialCharTok{::}\FunctionTok{arrange}\NormalTok{(.data}\SpecialCharTok{$}\NormalTok{step) }\SpecialCharTok{|\textgreater{}}
\NormalTok{  dplyr}\SpecialCharTok{::}\FunctionTok{pull}\NormalTok{(.data}\SpecialCharTok{$}\NormalTok{time)}

\NormalTok{dt\_check }\OtherTok{\textless{}{-}} \FunctionTok{diff}\NormalTok{(time\_grid\_check)}

\NormalTok{growth\_factor }\OtherTok{\textless{}{-}} \FunctionTok{exp}\NormalTok{(}
  \SpecialCharTok{{-}}\FloatTok{0.5} \SpecialCharTok{*}\NormalTok{ volatility}\SpecialCharTok{\^{}}\DecValTok{2} \SpecialCharTok{*}\NormalTok{ dt\_check }\SpecialCharTok{+}
\NormalTok{    volatility }\SpecialCharTok{*} \FunctionTok{sqrt}\NormalTok{(dt\_check) }\SpecialCharTok{*}\NormalTok{ gaussian\_increment}
\NormalTok{)}

\NormalTok{hand\_price }\OtherTok{\textless{}{-}}\NormalTok{ forward\_price }\SpecialCharTok{*} \FunctionTok{c}\NormalTok{(}
  \DecValTok{1}\NormalTok{,}
  \FunctionTok{cumprod}\NormalTok{(growth\_factor)}
\NormalTok{)}

\NormalTok{one\_path\_check\_tbl }\OtherTok{\textless{}{-}}\NormalTok{ generated\_path\_tbl }\SpecialCharTok{|\textgreater{}}
\NormalTok{  dplyr}\SpecialCharTok{::}\FunctionTok{arrange}\NormalTok{(.data}\SpecialCharTok{$}\NormalTok{step) }\SpecialCharTok{|\textgreater{}}
\NormalTok{  dplyr}\SpecialCharTok{::}\FunctionTok{transmute}\NormalTok{(}
    \AttributeTok{step =}\NormalTok{ .data}\SpecialCharTok{$}\NormalTok{step,}
    \AttributeTok{time =}\NormalTok{ .data}\SpecialCharTok{$}\NormalTok{time,}
    \AttributeTok{generated\_price =}\NormalTok{ .data}\SpecialCharTok{$}\NormalTok{price,}
    \AttributeTok{hand\_price =}\NormalTok{ hand\_price,}
    \AttributeTok{difference =}\NormalTok{ .data}\SpecialCharTok{$}\NormalTok{price }\SpecialCharTok{{-}}\NormalTok{ hand\_price}
\NormalTok{  )}

\NormalTok{GaussQuant}\SpecialCharTok{::}\FunctionTok{show\_tbl\_GQH}\NormalTok{(}
\NormalTok{  one\_path\_check\_tbl,}
  \StringTok{"One{-}path hand calculation"}\NormalTok{,}
  \AttributeTok{n =} \DecValTok{20}
\NormalTok{)}
\CommentTok{\#\textgreater{} \# A tibble: 4 x 5}
\CommentTok{\#\textgreater{}    step   time generated\_price hand\_price difference}
\CommentTok{\#\textgreater{}   \textless{}int\textgreater{}  \textless{}dbl\textgreater{}           \textless{}dbl\textgreater{}      \textless{}dbl\textgreater{}      \textless{}dbl\textgreater{}}
\CommentTok{\#\textgreater{} 1     0 0                 107.       107.   0       }
\CommentTok{\#\textgreater{} 2     1 0.0329            107.       107.   0       }
\CommentTok{\#\textgreater{} 3     2 0.0658            110.       110.   0       }
\CommentTok{\#\textgreater{} 4     3 0.0986            110.       110.  {-}1.42e{-}14}
\end{Highlighting}
\end{Shaded}

差が数値誤差の範囲に収まれば、PathGeneratorがBlack過程の離散化式に従っていることを確認できる。

\section{4.5 American
putの早期行使価値}\label{american-putux306eux65e9ux671fux884cux4f7fux4fa1ux5024}

Americanオプションは満期までの任意の時点で行使できる。したがって、その価値は同じ条件のEuropeanオプション以上となる。

\[
V_{\mathrm{American}}
\geq
V_{\mathrm{European}}
\]

両者の差は早期行使権の価値である。

American
putには一般的な閉形式の解析解がないため、二項Treeまたは近似法を用いる。本節では、次の方法を比較する。

\begin{itemize}
\tightlist
\item
  25ステップCRR Tree
\item
  Barone--Adesi--Whaley近似
\item
  Bjerksund--Stensland近似
\item
  500ステップCRR Tree
\end{itemize}

500ステップCRR Treeを比較基準とする。

\begin{Shaded}
\begin{Highlighting}[]
\NormalTok{american\_put }\OtherTok{\textless{}{-}}\NormalTok{ GaussQuant}\SpecialCharTok{::}\FunctionTok{make\_american\_option\_GQL}\NormalTok{(}
  \AttributeTok{spot =}\NormalTok{ spot\_price\_bsm,}
  \AttributeTok{strike =}\NormalTok{ strike\_price\_bsm,}
  \AttributeTok{maturity\_date =}\NormalTok{ maturity\_date\_bsm,}
  \AttributeTok{option\_type =} \StringTok{"put"}\NormalTok{,}
  \AttributeTok{valuation\_date =}\NormalTok{ trade\_date\_bsm,}
  \AttributeTok{risk\_free\_rate =}\NormalTok{ risk\_free\_rate\_bsm,}
  \AttributeTok{dividend\_yield =}\NormalTok{ dividend\_rate\_bsm,}
  \AttributeTok{volatility =}\NormalTok{ volatility\_bsm,}
  \AttributeTok{day\_counter =}\NormalTok{ day\_counter,}
  \AttributeTok{calendar =}\NormalTok{ calendar\_bsm,}
  \AttributeTok{steps =} \DecValTok{500}\NormalTok{L}
\NormalTok{)}

\NormalTok{american\_engines }\OtherTok{\textless{}{-}} \FunctionTok{list}\NormalTok{(}
  \AttributeTok{crr\_25 =}\NormalTok{ QuantLib}\SpecialCharTok{::}\FunctionTok{BinomialCRRVanillaEngine}\NormalTok{(}
\NormalTok{    bsm\_bundle}\SpecialCharTok{$}\NormalTok{process,}
    \DecValTok{25}\NormalTok{L}
\NormalTok{  ),}
  \AttributeTok{barone\_adesi\_whaley =}\NormalTok{ QuantLib}\SpecialCharTok{::}\FunctionTok{BaroneAdesiWhaleyApproximationEngine}\NormalTok{(}
\NormalTok{    bsm\_bundle}\SpecialCharTok{$}\NormalTok{process}
\NormalTok{  ),}
  \AttributeTok{bjerksund\_stensland =}\NormalTok{ QuantLib}\SpecialCharTok{::}\FunctionTok{BjerksundStenslandApproximationEngine}\NormalTok{(}
\NormalTok{    bsm\_bundle}\SpecialCharTok{$}\NormalTok{process}
\NormalTok{  ),}
  \AttributeTok{crr\_500 =}\NormalTok{ QuantLib}\SpecialCharTok{::}\FunctionTok{BinomialCRRVanillaEngine}\NormalTok{(}
\NormalTok{    bsm\_bundle}\SpecialCharTok{$}\NormalTok{process,}
    \DecValTok{500}\NormalTok{L}
\NormalTok{  )}
\NormalTok{)}

\NormalTok{american\_comparison\_tbl }\OtherTok{\textless{}{-}}\NormalTok{ purrr}\SpecialCharTok{::}\FunctionTok{imap\_dfr}\NormalTok{(}
\NormalTok{  american\_engines,}
  \ControlFlowTok{function}\NormalTok{(engine, method) \{}
\NormalTok{    american\_put}\SpecialCharTok{$}\FunctionTok{setPricingEngine}\NormalTok{(engine)}

\NormalTok{    tibble}\SpecialCharTok{::}\FunctionTok{tibble}\NormalTok{(}
      \AttributeTok{method =}\NormalTok{ method,}
      \AttributeTok{npv =}\NormalTok{ american\_put}\SpecialCharTok{$}\FunctionTok{NPV}\NormalTok{()}
\NormalTok{    )}
\NormalTok{  \}}
\NormalTok{)}

\NormalTok{crr\_500\_american\_npv }\OtherTok{\textless{}{-}}\NormalTok{ american\_comparison\_tbl }\SpecialCharTok{|\textgreater{}}
\NormalTok{  dplyr}\SpecialCharTok{::}\FunctionTok{filter}\NormalTok{(.data}\SpecialCharTok{$}\NormalTok{method }\SpecialCharTok{==} \StringTok{"crr\_500"}\NormalTok{) }\SpecialCharTok{|\textgreater{}}
\NormalTok{  dplyr}\SpecialCharTok{::}\FunctionTok{pull}\NormalTok{(.data}\SpecialCharTok{$}\NormalTok{npv)}

\NormalTok{american\_comparison\_tbl }\OtherTok{\textless{}{-}}\NormalTok{ american\_comparison\_tbl }\SpecialCharTok{|\textgreater{}}
\NormalTok{  dplyr}\SpecialCharTok{::}\FunctionTok{mutate}\NormalTok{(}
    \AttributeTok{difference\_from\_crr\_500 =}\NormalTok{ .data}\SpecialCharTok{$}\NormalTok{npv }\SpecialCharTok{{-}}\NormalTok{ crr\_500\_american\_npv}
\NormalTok{  )}

\NormalTok{GaussQuant}\SpecialCharTok{::}\FunctionTok{show\_tbl\_GQH}\NormalTok{(}
\NormalTok{  american\_comparison\_tbl,}
  \StringTok{"American put engine comparison"}\NormalTok{,}
  \AttributeTok{n =} \DecValTok{20}
\NormalTok{)}
\CommentTok{\#\textgreater{} \# A tibble: 4 x 3}
\CommentTok{\#\textgreater{}   method                npv difference\_from\_crr\_500}
\CommentTok{\#\textgreater{}   \textless{}chr\textgreater{}               \textless{}dbl\textgreater{}                   \textless{}dbl\textgreater{}}
\CommentTok{\#\textgreater{} 1 crr\_25               6.16                 0.0584 }
\CommentTok{\#\textgreater{} 2 barone\_adesi\_whaley  6.11                 0.00892}
\CommentTok{\#\textgreater{} 3 bjerksund\_stensland  6.00                {-}0.106  }
\CommentTok{\#\textgreater{} 4 crr\_500              6.10                 0}
\end{Highlighting}
\end{Shaded}

\section{4.6 準Monte
CarloとLongstaff--Schwartz法}\label{ux6e96monte-carloux3068longstaffschwartzux6cd5}

\subsection{4.6.1
Sobol系列によるパス生成}\label{sobolux7cfbux5217ux306bux3088ux308bux30d1ux30b9ux751fux6210}

擬似乱数は乱数生成器によって確率的な点列を作る。これに対しSobol系列は、積分領域をできるだけ均等に覆うように構成された低乖離列である。

Sobol系列を用いる準Monte
Carlo法では、通常の擬似乱数より少ないパスで安定した数値積分が得られる場合がある。ただし、パス同士は独立な確率標本ではないため、通常の標準誤差の解釈をそのまま適用できない点に注意する。

次のコードでは、Black--Scholes--Merton
processから64本のSobolパスを生成する。

\begin{Shaded}
\begin{Highlighting}[]
\NormalTok{maturity\_bsm }\OtherTok{\textless{}{-}}\NormalTok{ day\_counter}\SpecialCharTok{$}\FunctionTok{yearFraction}\NormalTok{(}
\NormalTok{  trade\_date\_bsm,}
\NormalTok{  maturity\_date\_bsm}
\NormalTok{)}

\NormalTok{sobol\_display\_tbl }\OtherTok{\textless{}{-}}\NormalTok{ GaussQuant}\SpecialCharTok{::}\FunctionTok{generate\_path\_table\_GQL}\NormalTok{(}
  \AttributeTok{process =}\NormalTok{ bsm\_bundle}\SpecialCharTok{$}\NormalTok{process,}
  \AttributeTok{maturity =}\NormalTok{ maturity\_bsm,}
  \AttributeTok{n\_steps =} \DecValTok{12}\NormalTok{L,}
  \AttributeTok{n\_paths =} \DecValTok{64}\NormalTok{L,}
  \AttributeTok{seed =} \DecValTok{1}\NormalTok{L,}
  \AttributeTok{sequence =} \StringTok{"sobol"}
\NormalTok{)}

\NormalTok{sobol\_plot }\OtherTok{\textless{}{-}}\NormalTok{ ggplot2}\SpecialCharTok{::}\FunctionTok{ggplot}\NormalTok{(}
\NormalTok{  sobol\_display\_tbl,}
\NormalTok{  ggplot2}\SpecialCharTok{::}\FunctionTok{aes}\NormalTok{(}
    \AttributeTok{x =}\NormalTok{ .data}\SpecialCharTok{$}\NormalTok{time,}
    \AttributeTok{y =}\NormalTok{ .data}\SpecialCharTok{$}\NormalTok{price,}
    \AttributeTok{group =}\NormalTok{ .data}\SpecialCharTok{$}\NormalTok{path\_id}
\NormalTok{  )}
\NormalTok{) }\SpecialCharTok{+}
\NormalTok{  ggplot2}\SpecialCharTok{::}\FunctionTok{geom\_line}\NormalTok{(}
    \AttributeTok{alpha =} \FloatTok{0.5}\NormalTok{,}
    \AttributeTok{linewidth =} \FloatTok{0.4}
\NormalTok{  ) }\SpecialCharTok{+}
\NormalTok{  ggplot2}\SpecialCharTok{::}\FunctionTok{labs}\NormalTok{(}
    \AttributeTok{title =} \StringTok{"Black{-}Scholes{-}Merton Sobol paths"}\NormalTok{,}
    \AttributeTok{x =} \StringTok{"Year"}\NormalTok{,}
    \AttributeTok{y =} \StringTok{"Stock price"}
\NormalTok{  ) }\SpecialCharTok{+}
\NormalTok{  ggplot2}\SpecialCharTok{::}\FunctionTok{theme\_minimal}\NormalTok{()}

\NormalTok{GaussQuant}\SpecialCharTok{::}\FunctionTok{show\_tbl\_GQH}\NormalTok{(}
\NormalTok{  sobol\_display\_tbl,}
  \StringTok{"Sobol path table"}\NormalTok{,}
  \AttributeTok{n =} \DecValTok{20}
\NormalTok{)}
\CommentTok{\#\textgreater{} \# A tibble: 20 x 4}
\CommentTok{\#\textgreater{}    path\_id  step   time price}
\CommentTok{\#\textgreater{}      \textless{}int\textgreater{} \textless{}int\textgreater{}  \textless{}dbl\textgreater{} \textless{}dbl\textgreater{}}
\CommentTok{\#\textgreater{}  1       1     0 0       100 }
\CommentTok{\#\textgreater{}  2       1     1 0.0838  100.}
\CommentTok{\#\textgreater{}  3       1     2 0.168   101.}
\CommentTok{\#\textgreater{}  4       1     3 0.251   101.}
\CommentTok{\#\textgreater{}  5       1     4 0.335   101.}
\CommentTok{\#\textgreater{}  6       1     5 0.419   101.}
\CommentTok{\#\textgreater{}  7       1     6 0.503   102.}
\CommentTok{\#\textgreater{}  8       1     7 0.587   102.}
\CommentTok{\#\textgreater{}  9       1     8 0.670   102.}
\CommentTok{\#\textgreater{} 10       1     9 0.754   102.}
\CommentTok{\#\textgreater{} 11       1    10 0.838   103.}
\CommentTok{\#\textgreater{} 12       1    11 0.922   103.}
\CommentTok{\#\textgreater{} 13       1    12 1.01    103.}
\CommentTok{\#\textgreater{} 14       2     0 0       100 }
\CommentTok{\#\textgreater{} 15       2     1 0.0838  104.}
\CommentTok{\#\textgreater{} 16       2     2 0.168   101.}
\CommentTok{\#\textgreater{} 17       2     3 0.251   105.}
\CommentTok{\#\textgreater{} 18       2     4 0.335   101.}
\CommentTok{\#\textgreater{} 19       2     5 0.419   105.}
\CommentTok{\#\textgreater{} 20       2     6 0.503   102.}

\FunctionTok{print}\NormalTok{(sobol\_plot)}
\end{Highlighting}
\end{Shaded}

\pandocbounded{\includegraphics[keepaspectratio]{part2_quantive_riemann/ch04_black_model_files/figure-pdf/chapter04-sobol-paths-1.pdf}}

\subsection{4.6.2 Longstaff--Schwartz法}\label{longstaffschwartzux6cd5}

Americanオプションでは、各行使可能時点において、即時行使価値と保有を継続する価値を比較する必要がある。

putの即時行使価値は次式である。

\[
I(S_t)
=
\max(K-S_t,0)
\]

Longstaff--Schwartz法では、将来受け取る割引Cashflowを原資産価格の関数として回帰し、条件付き継続価値を推定する。本章では、次の2次多項式を使用する。

\[
\widehat{C}(S_t)
=
\beta_0
+
\beta_1S_t
+
\beta_2S_t^2
\]

即時行使価値が推定継続価値を上回る場合、その時点で行使する。

本節では、2,048本のSobolパスからAmerican putを評価し、500ステップCRR
Treeと比較する。

\begin{Shaded}
\begin{Highlighting}[]
\NormalTok{lsm\_path\_tbl }\OtherTok{\textless{}{-}}\NormalTok{ GaussQuant}\SpecialCharTok{::}\FunctionTok{generate\_path\_table\_GQL}\NormalTok{(}
  \AttributeTok{process =}\NormalTok{ bsm\_bundle}\SpecialCharTok{$}\NormalTok{process,}
  \AttributeTok{maturity =}\NormalTok{ maturity\_bsm,}
  \AttributeTok{n\_steps =} \DecValTok{12}\NormalTok{L,}
  \AttributeTok{n\_paths =} \DecValTok{2048}\NormalTok{L,}
  \AttributeTok{seed =} \DecValTok{1}\NormalTok{L,}
  \AttributeTok{sequence =} \StringTok{"sobol"}
\NormalTok{)}

\NormalTok{lsm\_result }\OtherTok{\textless{}{-}}\NormalTok{ GaussQuant}\SpecialCharTok{::}\FunctionTok{lsm\_american\_put\_GQL}\NormalTok{(}
  \AttributeTok{path\_tbl =}\NormalTok{ lsm\_path\_tbl,}
  \AttributeTok{strike =}\NormalTok{ strike\_price\_bsm,}
  \AttributeTok{risk\_free\_rate =}\NormalTok{ risk\_free\_rate\_bsm}
\NormalTok{)}

\NormalTok{lsm\_npv }\OtherTok{\textless{}{-}}\NormalTok{ lsm\_result}\SpecialCharTok{$}\NormalTok{summary }\SpecialCharTok{|\textgreater{}}
\NormalTok{  dplyr}\SpecialCharTok{::}\FunctionTok{filter}\NormalTok{(.data}\SpecialCharTok{$}\NormalTok{metric }\SpecialCharTok{==} \StringTok{"lsm\_npv"}\NormalTok{) }\SpecialCharTok{|\textgreater{}}
\NormalTok{  dplyr}\SpecialCharTok{::}\FunctionTok{pull}\NormalTok{(.data}\SpecialCharTok{$}\NormalTok{value)}

\NormalTok{lsm\_comparison\_tbl }\OtherTok{\textless{}{-}}\NormalTok{ tibble}\SpecialCharTok{::}\FunctionTok{tribble}\NormalTok{(}
  \SpecialCharTok{\textasciitilde{}}\NormalTok{method,                    }\SpecialCharTok{\textasciitilde{}}\NormalTok{npv,}
  \StringTok{"Longstaff{-}Schwartz Sobol"}\NormalTok{, lsm\_npv,}
  \StringTok{"CRR Tree 500 steps"}\NormalTok{,       crr\_500\_american\_npv}
\NormalTok{) }\SpecialCharTok{|\textgreater{}}
\NormalTok{  dplyr}\SpecialCharTok{::}\FunctionTok{mutate}\NormalTok{(}
    \AttributeTok{difference\_from\_crr\_500 =}\NormalTok{ .data}\SpecialCharTok{$}\NormalTok{npv }\SpecialCharTok{{-}}\NormalTok{ crr\_500\_american\_npv}
\NormalTok{  )}

\NormalTok{exercise\_summary\_tbl }\OtherTok{\textless{}{-}}\NormalTok{ lsm\_result}\SpecialCharTok{$}\NormalTok{exercise }\SpecialCharTok{|\textgreater{}}
\NormalTok{  dplyr}\SpecialCharTok{::}\FunctionTok{count}\NormalTok{(}
\NormalTok{    .data}\SpecialCharTok{$}\NormalTok{exercise\_step,}
\NormalTok{    .data}\SpecialCharTok{$}\NormalTok{exercise\_time,}
    \AttributeTok{name =} \StringTok{"paths"}
\NormalTok{  ) }\SpecialCharTok{|\textgreater{}}
\NormalTok{  dplyr}\SpecialCharTok{::}\FunctionTok{arrange}\NormalTok{(.data}\SpecialCharTok{$}\NormalTok{exercise\_step)}

\NormalTok{GaussQuant}\SpecialCharTok{::}\FunctionTok{show\_tbl\_GQH}\NormalTok{(}
\NormalTok{  lsm\_result}\SpecialCharTok{$}\NormalTok{summary,}
  \StringTok{"Longstaff{-}Schwartz summary"}\NormalTok{,}
  \AttributeTok{n =} \DecValTok{20}
\NormalTok{)}
\CommentTok{\#\textgreater{} \# A tibble: 4 x 2}
\CommentTok{\#\textgreater{}   metric            value}
\CommentTok{\#\textgreater{}   \textless{}chr\textgreater{}             \textless{}dbl\textgreater{}}
\CommentTok{\#\textgreater{} 1 lsm\_npv           6.29 }
\CommentTok{\#\textgreater{} 2 standard\_error    0.167}
\CommentTok{\#\textgreater{} 3 paths          2048    }
\CommentTok{\#\textgreater{} 4 steps            12}

\NormalTok{GaussQuant}\SpecialCharTok{::}\FunctionTok{show\_tbl\_GQH}\NormalTok{(}
\NormalTok{  lsm\_comparison\_tbl,}
  \StringTok{"Longstaff{-}Schwartz and CRR comparison"}\NormalTok{,}
  \AttributeTok{n =} \DecValTok{20}
\NormalTok{)}
\CommentTok{\#\textgreater{} \# A tibble: 2 x 3}
\CommentTok{\#\textgreater{}   method                     npv difference\_from\_crr\_500}
\CommentTok{\#\textgreater{}   \textless{}chr\textgreater{}                    \textless{}dbl\textgreater{}                   \textless{}dbl\textgreater{}}
\CommentTok{\#\textgreater{} 1 Longstaff{-}Schwartz Sobol  6.29                   0.191}
\CommentTok{\#\textgreater{} 2 CRR Tree 500 steps        6.10                   0}

\NormalTok{GaussQuant}\SpecialCharTok{::}\FunctionTok{show\_tbl\_GQH}\NormalTok{(}
\NormalTok{  exercise\_summary\_tbl,}
  \StringTok{"Longstaff{-}Schwartz exercise timing"}\NormalTok{,}
  \AttributeTok{n =} \DecValTok{20}
\NormalTok{)}
\CommentTok{\#\textgreater{} \# A tibble: 11 x 3}
\CommentTok{\#\textgreater{}    exercise\_step exercise\_time paths}
\CommentTok{\#\textgreater{}            \textless{}int\textgreater{}         \textless{}dbl\textgreater{} \textless{}int\textgreater{}}
\CommentTok{\#\textgreater{}  1             2         0.168    28}
\CommentTok{\#\textgreater{}  2             3         0.251    52}
\CommentTok{\#\textgreater{}  3             4         0.335    57}
\CommentTok{\#\textgreater{}  4             5         0.419   109}
\CommentTok{\#\textgreater{}  5             6         0.503   107}
\CommentTok{\#\textgreater{}  6             7         0.587   100}
\CommentTok{\#\textgreater{}  7             8         0.670    87}
\CommentTok{\#\textgreater{}  8             9         0.754    38}
\CommentTok{\#\textgreater{}  9            10         0.838    78}
\CommentTok{\#\textgreater{} 10            11         0.922    96}
\CommentTok{\#\textgreater{} 11            12         1.01   1296}

\NormalTok{GaussQuant}\SpecialCharTok{::}\FunctionTok{show\_tbl\_GQH}\NormalTok{(}
\NormalTok{  lsm\_result}\SpecialCharTok{$}\NormalTok{diagnostics,}
  \StringTok{"Longstaff{-}Schwartz diagnostics"}\NormalTok{,}
  \AttributeTok{n =} \DecValTok{20}
\NormalTok{)}
\CommentTok{\#\textgreater{} \# A tibble: 11 x 6}
\CommentTok{\#\textgreater{}     step   time in\_the\_money\_paths exercised\_paths mean\_immediate\_value}
\CommentTok{\#\textgreater{}    \textless{}int\textgreater{}  \textless{}dbl\textgreater{}              \textless{}int\textgreater{}           \textless{}int\textgreater{}                \textless{}dbl\textgreater{}}
\CommentTok{\#\textgreater{}  1    11 0.922                 931             575                 5.68}
\CommentTok{\#\textgreater{}  2    10 0.838                 918             459                 5.53}
\CommentTok{\#\textgreater{}  3     9 0.754                 932             355                 5.28}
\CommentTok{\#\textgreater{}  4     8 0.670                 937             392                 5.08}
\CommentTok{\#\textgreater{}  5     7 0.587                 945             338                 4.82}
\CommentTok{\#\textgreater{}  6     6 0.503                 946             282                 4.55}
\CommentTok{\#\textgreater{}  7     5 0.419                 939             193                 4.25}
\CommentTok{\#\textgreater{}  8     4 0.335                 959             105                 3.85}
\CommentTok{\#\textgreater{}  9     3 0.251                 970              69                 3.41}
\CommentTok{\#\textgreater{} 10     2 0.168                 963              28                 2.88}
\CommentTok{\#\textgreater{} 11     1 0.0838                989               0                 2.11}
\CommentTok{\#\textgreater{} \# i 1 more variable: mean\_realized\_continuation \textless{}dbl\textgreater{}}
\end{Highlighting}
\end{Shaded}

Longstaff--Schwartz法の価格は、パス数、時間ステップ、基底関数および回帰誤差に依存する。したがって、二項Treeとの差が完全にゼロになることを目的とするのではなく、早期行使判断をMonte
Carloパス上でどのように近似するかを確認する。

\section{4.7 まとめ}\label{ux307eux3068ux3081-8}

Black-76とBlack--Scholes--Mertonは、原資産が対数正規分布に従うという共通の仮定を持つ。Black--Scholes--Mertonは現物価格から出発し、Black-76はフォワード価格または先物価格を直接入力する。

Black-76では、Black公式による直接計算と、Black
processを用いた解析的Pricing
Engineが同じ価格を与えることを確認した。さらに、Monte
Carlo法によって満期Payoffの割引期待値を求め、解析値と比較した。

Black--Scholes--Mertonモデルでは、European putの解析解とCRR
Treeを比較した。American
putでは、早期行使権を評価するため、近似Engine、CRR
TreeおよびLongstaff--Schwartz法を比較した。

解析解、Tree、Monte
Carloおよび回帰法は、互いに競合する単純な代替手法ではない。Payoffの構造、早期行使の有無、モデルの次元および必要な計算精度に応じて使い分ける必要がある。

Blackモデルの対数正規仮定は、価格が正である市場では扱いやすい。一方、金利がゼロまたはマイナスとなる場合には直接適用できない。この問題が、NormalモデルやShifted
Blackモデルを用いる理由となる。

\begin{Shaded}
\begin{Highlighting}[]
\NormalTok{chapter04\_summary\_tbl }\OtherTok{\textless{}{-}}\NormalTok{ tibble}\SpecialCharTok{::}\FunctionTok{tribble}\NormalTok{(}
  \SpecialCharTok{\textasciitilde{}}\NormalTok{section,                    }\SpecialCharTok{\textasciitilde{}}\NormalTok{main\_result,}
  \StringTok{"Black{-}76"}\NormalTok{,                  }\StringTok{"Black calculator and analytic European pricing"}\NormalTok{,}
  \StringTok{"Pseudo{-}random Monte Carlo"}\NormalTok{, }\StringTok{"Discounted terminal payoff and standard error"}\NormalTok{,}
  \StringTok{"European put"}\NormalTok{,              }\StringTok{"Analytic and CRR Tree comparison"}\NormalTok{,}
  \StringTok{"American put"}\NormalTok{,              }\StringTok{"Approximation engines and CRR Tree comparison"}\NormalTok{,}
  \StringTok{"Sobol"}\NormalTok{,                     }\StringTok{"Low{-}discrepancy path generation"}\NormalTok{,}
  \StringTok{"Longstaff{-}Schwartz"}\NormalTok{,        }\StringTok{"Regression{-}based early{-}exercise approximation"}
\NormalTok{)}

\NormalTok{GaussQuant}\SpecialCharTok{::}\FunctionTok{show\_tbl\_GQH}\NormalTok{(}
\NormalTok{  chapter04\_summary\_tbl,}
  \StringTok{"Chapter II{-}4 summary"}\NormalTok{,}
  \AttributeTok{n =} \DecValTok{20}
\NormalTok{)}
\CommentTok{\#\textgreater{} \# A tibble: 6 x 2}
\CommentTok{\#\textgreater{}   section                   main\_result                                   }
\CommentTok{\#\textgreater{}   \textless{}chr\textgreater{}                     \textless{}chr\textgreater{}                                         }
\CommentTok{\#\textgreater{} 1 Black{-}76                  Black calculator and analytic European pricing}
\CommentTok{\#\textgreater{} 2 Pseudo{-}random Monte Carlo Discounted terminal payoff and standard error }
\CommentTok{\#\textgreater{} 3 European put              Analytic and CRR Tree comparison              }
\CommentTok{\#\textgreater{} 4 American put              Approximation engines and CRR Tree comparison }
\CommentTok{\#\textgreater{} 5 Sobol                     Low{-}discrepancy path generation               }
\CommentTok{\#\textgreater{} 6 Longstaff{-}Schwartz        Regression{-}based early{-}exercise approximation}
\end{Highlighting}
\end{Shaded}

\chapter{第II-5章
金利モデル（期間構造なし）Normaモデル}\label{ux7b2cii-5ux7ae0-ux91d1ux5229ux30e2ux30c7ux30ebux671fux9593ux69cbux9020ux306aux3057normaux30e2ux30c7ux30eb}

前章では、株価や為替レートのように、価格の変化率をウィーナー過程によって表現する対象について、Blackモデルを検討した。Blackモデルでは、原資産の変動幅がその時点の価格水準に比例する幾何ブラウン運動を仮定するため、原資産価格は常に正の値を取る。

これに対して、本章では金利を原資産とするオプションを検討する。金利は株価や為替レートとは異なり、満期ごとに異なる期間構造を持ち、ゼロまたはマイナスの値を取る可能性もある。このため、金利を単一の幾何ブラウン運動によって表現することには限界がある。

本章では、金利期間構造全体を確率モデルとして記述するのではなく、特定の将来期間に対応するフォワード金利を直接モデル化する。適切なフォワード測度の下では、フォワード金利をマルチンゲールとして扱うことができる。

Blackモデルでは、フォワード金利の変化幅がその金利水準に比例すると仮定する。

\[
dF_t
=
\sigma_B F_t\,dW_t
\]

この場合、フォワード金利は理論上常に正の値を取り、ゼロまたはマイナスの金利を直接扱うことができない。

そこで本章では、フォワード金利の変化額がウィーナー過程に従うNormalモデルを検討する。

\[
dF_t
=
\sigma_N\,dW_t
\]

このモデルでは、将来のフォワード金利は次のように表される。

\[
F_T
=
F_0
+
\sigma_N\sqrt{T}Z
\]

ここで、\(Z\)
は標準正規分布に従う確率変数である。フォワード金利の変化幅は現在の金利水準に依存しないため、将来金利はゼロまたはマイナスの値も取り得る。

本章では、あえて金利期間構造全体の動学モデルを導入せず、各商品の評価に必要なフォワード金利を、それぞれに対応する測度の下でマルチンゲールとして直接モデル化する。この枠組みに基づき、金利先物オプション、Caplet、CapおよびSwaptionをNormalモデルによって評価する。

\begin{Shaded}
\begin{Highlighting}[]
\NormalTok{knitr}\SpecialCharTok{::}\NormalTok{opts\_chunk}\SpecialCharTok{$}\FunctionTok{set}\NormalTok{(}
  \AttributeTok{collapse =} \ConstantTok{TRUE}\NormalTok{,}
  \AttributeTok{comment =} \StringTok{"\#\textgreater{}"}\NormalTok{,}
  \AttributeTok{warning =} \ConstantTok{FALSE}\NormalTok{,}
  \AttributeTok{message =} \ConstantTok{FALSE}
\NormalTok{)}

\FunctionTok{suppressPackageStartupMessages}\NormalTok{(\{}
  \FunctionTok{library}\NormalTok{(tidyverse)}
  \FunctionTok{library}\NormalTok{(QuantLib)}
  \FunctionTok{library}\NormalTok{(GaussQuant)}
\NormalTok{\})}
\end{Highlighting}
\end{Shaded}

\section{5.1
Normalモデルの基礎}\label{normalux30e2ux30c7ux30ebux306eux57faux790e}

\subsection{5.1.1
金利モデルと測度}\label{ux91d1ux5229ux30e2ux30c7ux30ebux3068ux6e2cux5ea6}

金利デリバティブでは、すべてのフォワード金利を一つの確率過程で表現するとは限らない。単一のCapletを評価する場合、そのCapletが参照する期間のフォワード金利を、支払日を満期とするゼロクーポン債をNumeraireとしたフォワード測度の下で扱う。

期間 \([T_1,T_2]\) のフォワード金利を \(L_t(T_1,T_2)\)
とすると、\(T_2\)-forward
measureの下で、これをマルチンゲールとして次のように置く。

\[
dL_t(T_1,T_2)
=
\sigma_N\,dW_t^{T_2}
\]

一方、Swaptionでは、原資産は単一期間のフォワード金利ではなく、複数期間から構成されるフォワードスワップレートである。この場合は、固定レグのアニュイティをNumeraireとするswap
measureの下で、フォワードスワップレートをマルチンゲールとして扱う。

したがって本章では、期間構造全体を確率的に生成するHJM、LMMまたは短期金利モデルへ進むのではなく、各商品の自然なフォワード変数に対して局所的なNormalモデルを適用する。

現在価値を計算するためには割引係数が必要であるが、割引カーブ自体の動学はモデル化しない。本章では、与えられた割引係数または簡単な決定論的割引率を使用する。

\subsection{5.1.2
Bachelier価格公式}\label{bachelierux4fa1ux683cux516cux5f0f}

Normalモデルの下で、満期時のフォワード金利は次の正規分布に従う。

\[
F_T
\sim
N\left(
F_0,\,
\sigma_N^2T
\right)
\]

European callの価格は次式で与えられる。

\[
C_N
=
D(0,T)
\left[
(F_0-K)N(d)
+
\sigma_N\sqrt{T}\phi(d)
\right]
\]

European putの価格は次式となる。

\[
P_N
=
D(0,T)
\left[
(K-F_0)N(-d)
+
\sigma_N\sqrt{T}\phi(d)
\right]
\]

ここで、

\[
d
=
\frac{
F_0-K
}{
\sigma_N\sqrt{T}
}
\]

であり、\(D(0,T)\) は割引係数、\(N(\cdot)\)
は標準正規分布の累積分布関数、\(\phi(\cdot)\)
は標準正規分布の密度関数である。

Black公式には \(\log(F_0/K)\)が含まれるが、Normal公式では差額 \(F_0-K\)
だけを使用する。このため、フォワード金利または行使金利がゼロ以下であっても計算できる。

\subsection{5.1.3 Normal
volatilityの単位}\label{normal-volatilityux306eux5358ux4f4d}

Black
volatilityは原資産価格に対する相対的な変化率であり、通常は百分率で表示される。

これに対してNormal
volatilityは、原資産と同じ単位を持つ。金利を小数で表す場合、Normal
volatility 50bpは次の値である。

\[
50\text{ bp}
=
\frac{50}{10000}
=
0.0050
\]

したがって、\texttt{vol\ =\ 0.0050} は「0.5\%のBlack
volatility」ではなく、\textbf{年率50bpのNormal volatility}を意味する。

Normalモデルでは、

\[
F_T-F_0
=
\sigma_N\sqrt{T}Z
\]

であるため、\(\sigma_N\sqrt{T}\)
は満期までの金利変化額の標準偏差となる。

\begin{Shaded}
\begin{Highlighting}[]
\NormalTok{normal\_volatility\_bp }\OtherTok{\textless{}{-}} \DecValTok{50}
\NormalTok{normal\_volatility }\OtherTok{\textless{}{-}}\NormalTok{ normal\_volatility\_bp }\SpecialCharTok{/} \DecValTok{10000}

\NormalTok{normal\_volatility\_tbl }\OtherTok{\textless{}{-}}\NormalTok{ tibble}\SpecialCharTok{::}\FunctionTok{tribble}\NormalTok{(}
  \SpecialCharTok{\textasciitilde{}}\NormalTok{metric,                        }\SpecialCharTok{\textasciitilde{}}\NormalTok{value,}
  \StringTok{"normal\_volatility\_bp"}\NormalTok{,         normal\_volatility\_bp,}
  \StringTok{"normal\_volatility\_decimal"}\NormalTok{,    normal\_volatility,}
  \StringTok{"one\_standard\_deviation\_1y\_bp"}\NormalTok{, normal\_volatility }\SpecialCharTok{*} \DecValTok{10000}
\NormalTok{)}

\NormalTok{GaussQuant}\SpecialCharTok{::}\FunctionTok{show\_tbl\_GQH}\NormalTok{(}
\NormalTok{  normal\_volatility\_tbl,}
  \StringTok{"Normal volatility units"}\NormalTok{,}
  \AttributeTok{n =} \DecValTok{20}
\NormalTok{)}
\CommentTok{\#\textgreater{} \# A tibble: 3 x 2}
\CommentTok{\#\textgreater{}   metric                        value}
\CommentTok{\#\textgreater{}   \textless{}chr\textgreater{}                         \textless{}dbl\textgreater{}}
\CommentTok{\#\textgreater{} 1 normal\_volatility\_bp         50    }
\CommentTok{\#\textgreater{} 2 normal\_volatility\_decimal     0.005}
\CommentTok{\#\textgreater{} 3 one\_standard\_deviation\_1y\_bp 50}
\end{Highlighting}
\end{Shaded}

\section{5.2
SOFR先物オプション}\label{sofrux5148ux7269ux30aaux30d7ux30b7ux30e7ux30f3}

\subsection{5.2.1
先物価格から金利への変換}\label{ux5148ux7269ux4fa1ux683cux304bux3089ux91d1ux5229ux3078ux306eux5909ux63db}

短期金利先物は、一般に金利と反対方向に動く価格表示を使用する。本例では次の関係を用いる。

\[
\text{Futures Price}
=
100
-
100\times\text{Rate}
\]

したがって、先物価格94.54は金利5.46\%に対応する。

\[
F_0
=
\frac{100-94.54}{100}
=
0.0546
\]

行使価格94.50は行使金利5.50\%に対応する。

\[
K
=
\frac{100-94.50}{100}
=
0.0550
\]

先物価格のcallは、金利が低下して先物価格が上昇したときに価値を持つ。金利を原資産として表現すると、これはputに相当する。そのため、Normal公式では
\texttt{option\_sign\ =\ -1} を使用する。

\begin{Shaded}
\begin{Highlighting}[]
\NormalTok{trade\_date }\OtherTok{\textless{}{-}} \FunctionTok{as.Date}\NormalTok{(}\StringTok{"2023{-}09{-}26"}\NormalTok{)}
\NormalTok{maturity\_date }\OtherTok{\textless{}{-}} \FunctionTok{as.Date}\NormalTok{(}\StringTok{"2023{-}12{-}15"}\NormalTok{)}

\NormalTok{futures\_price }\OtherTok{\textless{}{-}} \FloatTok{94.54}
\NormalTok{strike\_futures\_price }\OtherTok{\textless{}{-}} \FloatTok{94.50}

\NormalTok{forward\_rate }\OtherTok{\textless{}{-}}\NormalTok{ (}\DecValTok{100} \SpecialCharTok{{-}}\NormalTok{ futures\_price) }\SpecialCharTok{/} \DecValTok{100}
\NormalTok{strike\_rate }\OtherTok{\textless{}{-}}\NormalTok{ (}\DecValTok{100} \SpecialCharTok{{-}}\NormalTok{ strike\_futures\_price) }\SpecialCharTok{/} \DecValTok{100}

\NormalTok{normal\_volatility }\OtherTok{\textless{}{-}} \DecValTok{50} \SpecialCharTok{/} \DecValTok{10000}
\NormalTok{risk\_free\_rate }\OtherTok{\textless{}{-}} \FloatTok{5.4} \SpecialCharTok{/} \DecValTok{100}

\NormalTok{maturity\_year\_fraction }\OtherTok{\textless{}{-}} \FunctionTok{as.numeric}\NormalTok{(}
\NormalTok{  maturity\_date }\SpecialCharTok{{-}}\NormalTok{ trade\_date}
\NormalTok{) }\SpecialCharTok{/} \DecValTok{365}

\NormalTok{maturity\_discount\_factor }\OtherTok{\textless{}{-}} \FunctionTok{exp}\NormalTok{(}
  \SpecialCharTok{{-}}\NormalTok{risk\_free\_rate }\SpecialCharTok{*}\NormalTok{ maturity\_year\_fraction}
\NormalTok{)}

\NormalTok{option\_sign }\OtherTok{\textless{}{-}} \SpecialCharTok{{-}}\DecValTok{1}

\NormalTok{futures\_option\_input\_tbl }\OtherTok{\textless{}{-}}\NormalTok{ tibble}\SpecialCharTok{::}\FunctionTok{tribble}\NormalTok{(}
  \SpecialCharTok{\textasciitilde{}}\NormalTok{input,                   }\SpecialCharTok{\textasciitilde{}}\NormalTok{value,}
  \StringTok{"trade\_date"}\NormalTok{,             }\FunctionTok{as.character}\NormalTok{(trade\_date),}
  \StringTok{"maturity\_date"}\NormalTok{,          }\FunctionTok{as.character}\NormalTok{(maturity\_date),}
  \StringTok{"futures\_price"}\NormalTok{,          }\FunctionTok{as.character}\NormalTok{(futures\_price),}
  \StringTok{"strike\_futures\_price"}\NormalTok{,   }\FunctionTok{as.character}\NormalTok{(strike\_futures\_price),}
  \StringTok{"forward\_rate"}\NormalTok{,           }\FunctionTok{as.character}\NormalTok{(forward\_rate),}
  \StringTok{"strike\_rate"}\NormalTok{,            }\FunctionTok{as.character}\NormalTok{(strike\_rate),}
  \StringTok{"normal\_volatility"}\NormalTok{,      }\FunctionTok{as.character}\NormalTok{(normal\_volatility),}
  \StringTok{"maturity\_year\_fraction"}\NormalTok{, }\FunctionTok{as.character}\NormalTok{(maturity\_year\_fraction),}
  \StringTok{"discount\_factor"}\NormalTok{,        }\FunctionTok{as.character}\NormalTok{(maturity\_discount\_factor)}
\NormalTok{)}

\NormalTok{GaussQuant}\SpecialCharTok{::}\FunctionTok{show\_tbl\_GQH}\NormalTok{(}
\NormalTok{  futures\_option\_input\_tbl,}
  \StringTok{"Normal futures{-}option assumptions"}\NormalTok{,}
  \AttributeTok{n =} \DecValTok{20}
\NormalTok{)}
\CommentTok{\#\textgreater{} \# A tibble: 9 x 2}
\CommentTok{\#\textgreater{}   input                  value             }
\CommentTok{\#\textgreater{}   \textless{}chr\textgreater{}                  \textless{}chr\textgreater{}             }
\CommentTok{\#\textgreater{} 1 trade\_date             2023{-}09{-}26        }
\CommentTok{\#\textgreater{} 2 maturity\_date          2023{-}12{-}15        }
\CommentTok{\#\textgreater{} 3 futures\_price          94.54             }
\CommentTok{\#\textgreater{} 4 strike\_futures\_price   94.5              }
\CommentTok{\#\textgreater{} 5 forward\_rate           0.0545999999999999}
\CommentTok{\#\textgreater{} 6 strike\_rate            0.055             }
\CommentTok{\#\textgreater{} 7 normal\_volatility      0.005             }
\CommentTok{\#\textgreater{} 8 maturity\_year\_fraction 0.219178082191781 }
\CommentTok{\#\textgreater{} 9 discount\_factor        0.988234148959797}
\end{Highlighting}
\end{Shaded}

\subsection{5.2.2 価格とGreeks}\label{ux4fa1ux683cux3068greeks}

Normalモデルによる価格とGreeksを計算する。

\begin{Shaded}
\begin{Highlighting}[]
\NormalTok{normal\_greeks\_tbl }\OtherTok{\textless{}{-}}\NormalTok{ GaussQuant}\SpecialCharTok{::}\FunctionTok{normal\_option\_greeks\_GQL}\NormalTok{(}
  \AttributeTok{option\_sign =}\NormalTok{ option\_sign,}
  \AttributeTok{forward =}\NormalTok{ forward\_rate,}
  \AttributeTok{strike =}\NormalTok{ strike\_rate,}
  \AttributeTok{vol =}\NormalTok{ normal\_volatility,}
  \AttributeTok{maturity =}\NormalTok{ maturity\_year\_fraction,}
  \AttributeTok{discount\_factor =}\NormalTok{ maturity\_discount\_factor,}
  \AttributeTok{risk\_free\_rate =}\NormalTok{ risk\_free\_rate}
\NormalTok{)}

\NormalTok{normal\_npv }\OtherTok{\textless{}{-}}\NormalTok{ normal\_greeks\_tbl }\SpecialCharTok{|\textgreater{}}
\NormalTok{  dplyr}\SpecialCharTok{::}\FunctionTok{filter}\NormalTok{(.data}\SpecialCharTok{$}\NormalTok{metric }\SpecialCharTok{==} \StringTok{"npv"}\NormalTok{) }\SpecialCharTok{|\textgreater{}}
\NormalTok{  dplyr}\SpecialCharTok{::}\FunctionTok{pull}\NormalTok{(.data}\SpecialCharTok{$}\NormalTok{value)}

\NormalTok{contract\_value\_multiplier }\OtherTok{\textless{}{-}} \DecValTok{100} \SpecialCharTok{*} \DecValTok{2500}

\NormalTok{normal\_summary\_tbl }\OtherTok{\textless{}{-}}\NormalTok{ dplyr}\SpecialCharTok{::}\FunctionTok{bind\_rows}\NormalTok{(}
\NormalTok{  tibble}\SpecialCharTok{::}\FunctionTok{tibble}\NormalTok{(}
    \AttributeTok{metric =} \FunctionTok{c}\NormalTok{(}
      \StringTok{"option\_npv\_rate\_unit"}\NormalTok{,}
      \StringTok{"contract\_value\_multiplier"}\NormalTok{,}
      \StringTok{"delivery\_amount"}
\NormalTok{    ),}
    \AttributeTok{value =} \FunctionTok{c}\NormalTok{(}
\NormalTok{      normal\_npv,}
\NormalTok{      contract\_value\_multiplier,}
\NormalTok{      normal\_npv }\SpecialCharTok{*}\NormalTok{ contract\_value\_multiplier}
\NormalTok{    )}
\NormalTok{  ),}
\NormalTok{  normal\_greeks\_tbl }\SpecialCharTok{|\textgreater{}}
\NormalTok{    dplyr}\SpecialCharTok{::}\FunctionTok{filter}\NormalTok{(.data}\SpecialCharTok{$}\NormalTok{metric }\SpecialCharTok{!=} \StringTok{"npv"}\NormalTok{)}
\NormalTok{)}

\NormalTok{GaussQuant}\SpecialCharTok{::}\FunctionTok{show\_tbl\_GQH}\NormalTok{(}
\NormalTok{  normal\_summary\_tbl,}
  \StringTok{"Normal futures{-}option summary"}\NormalTok{,}
  \AttributeTok{n =} \DecValTok{20}
\NormalTok{)}
\CommentTok{\#\textgreater{} \# A tibble: 7 x 2}
\CommentTok{\#\textgreater{}   metric                           value}
\CommentTok{\#\textgreater{}   \textless{}chr\textgreater{}                            \textless{}dbl\textgreater{}}
\CommentTok{\#\textgreater{} 1 option\_npv\_rate\_unit           0.00113}
\CommentTok{\#\textgreater{} 2 contract\_value\_multiplier 250000      }
\CommentTok{\#\textgreater{} 3 delivery\_amount              283.     }
\CommentTok{\#\textgreater{} 4 delta                         {-}0.561  }
\CommentTok{\#\textgreater{} 5 gamma                        166.     }
\CommentTok{\#\textgreater{} 6 vega                           0.182  }
\CommentTok{\#\textgreater{} 7 theta                         {-}0.00201}
\end{Highlighting}
\end{Shaded}

\texttt{option\_npv\_rate\_unit}は金利を小数表示した単位でのオプション価値である。金額換算は商品ごとの契約乗数に依存する。本例では、原章と同じ換算係数
\texttt{100\ *\ 2500} を使用する。

\section{5.3
マイナス金利とボラティリティ分析}\label{ux30deux30a4ux30caux30b9ux91d1ux5229ux3068ux30dcux30e9ux30c6ux30a3ux30eaux30c6ux30a3ux5206ux6790}

\subsection{5.3.1
ゼロ金利とマイナス金利}\label{ux30bcux30edux91d1ux5229ux3068ux30deux30a4ux30caux30b9ux91d1ux5229}

Normalモデルの重要な特徴は、フォワード金利がゼロまたはマイナスの場合でも価格公式を変更する必要がないことである。

次の例では、フォワード金利をマイナス0.20\%、行使金利を0\%、Normal
volatilityを60bpとする。

\begin{Shaded}
\begin{Highlighting}[]
\NormalTok{negative\_rate\_example\_tbl }\OtherTok{\textless{}{-}}\NormalTok{ tibble}\SpecialCharTok{::}\FunctionTok{tribble}\NormalTok{(}
  \SpecialCharTok{\textasciitilde{}}\NormalTok{option\_type, }\SpecialCharTok{\textasciitilde{}}\NormalTok{option\_sign, }\SpecialCharTok{\textasciitilde{}}\NormalTok{forward, }\SpecialCharTok{\textasciitilde{}}\NormalTok{strike,}
  \StringTok{"call"}\NormalTok{,        }\DecValTok{1}\NormalTok{,           }\SpecialCharTok{{-}}\FloatTok{0.0020}\NormalTok{,   }\FloatTok{0.0000}\NormalTok{,}
  \StringTok{"put"}\NormalTok{,        }\SpecialCharTok{{-}}\DecValTok{1}\NormalTok{,           }\SpecialCharTok{{-}}\FloatTok{0.0020}\NormalTok{,   }\FloatTok{0.0000}
\NormalTok{) }\SpecialCharTok{|\textgreater{}}
\NormalTok{  dplyr}\SpecialCharTok{::}\FunctionTok{mutate}\NormalTok{(}
    \AttributeTok{normal\_volatility =} \DecValTok{60} \SpecialCharTok{/} \DecValTok{10000}\NormalTok{,}
    \AttributeTok{maturity =} \DecValTok{1}\NormalTok{,}
    \AttributeTok{discount\_factor =} \FunctionTok{exp}\NormalTok{(}\SpecialCharTok{{-}}\FloatTok{0.01}\NormalTok{),}
    \AttributeTok{npv =}\NormalTok{ purrr}\SpecialCharTok{::}\FunctionTok{pmap\_dbl}\NormalTok{(}
      \FunctionTok{list}\NormalTok{(}
\NormalTok{        .data}\SpecialCharTok{$}\NormalTok{option\_sign,}
\NormalTok{        .data}\SpecialCharTok{$}\NormalTok{forward,}
\NormalTok{        .data}\SpecialCharTok{$}\NormalTok{strike,}
\NormalTok{        .data}\SpecialCharTok{$}\NormalTok{normal\_volatility,}
\NormalTok{        .data}\SpecialCharTok{$}\NormalTok{maturity,}
\NormalTok{        .data}\SpecialCharTok{$}\NormalTok{discount\_factor}
\NormalTok{      ),}
      \ControlFlowTok{function}\NormalTok{(}
\NormalTok{        option\_sign,}
\NormalTok{        forward,}
\NormalTok{        strike,}
\NormalTok{        normal\_volatility,}
\NormalTok{        maturity,}
\NormalTok{        discount\_factor}
\NormalTok{      ) \{}
\NormalTok{        GaussQuant}\SpecialCharTok{::}\FunctionTok{normal\_option\_price\_GQL}\NormalTok{(}
          \AttributeTok{option\_sign =}\NormalTok{ option\_sign,}
          \AttributeTok{forward =}\NormalTok{ forward,}
          \AttributeTok{strike =}\NormalTok{ strike,}
          \AttributeTok{vol =}\NormalTok{ normal\_volatility,}
          \AttributeTok{maturity =}\NormalTok{ maturity,}
          \AttributeTok{discount\_factor =}\NormalTok{ discount\_factor}
\NormalTok{        )}
\NormalTok{      \}}
\NormalTok{    )}
\NormalTok{  )}

\NormalTok{GaussQuant}\SpecialCharTok{::}\FunctionTok{show\_tbl\_GQH}\NormalTok{(}
\NormalTok{  negative\_rate\_example\_tbl,}
  \StringTok{"Normal pricing with a negative forward rate"}\NormalTok{,}
  \AttributeTok{n =} \DecValTok{20}
\NormalTok{)}
\CommentTok{\#\textgreater{} \# A tibble: 2 x 8}
\CommentTok{\#\textgreater{}   option\_type option\_sign forward strike normal\_volatility maturity}
\CommentTok{\#\textgreater{}   \textless{}chr\textgreater{}             \textless{}dbl\textgreater{}   \textless{}dbl\textgreater{}  \textless{}dbl\textgreater{}             \textless{}dbl\textgreater{}    \textless{}dbl\textgreater{}}
\CommentTok{\#\textgreater{} 1 call                  1  {-}0.002      0             0.006        1}
\CommentTok{\#\textgreater{} 2 put                  {-}1  {-}0.002      0             0.006        1}
\CommentTok{\#\textgreater{} \# i 2 more variables: discount\_factor \textless{}dbl\textgreater{}, npv \textless{}dbl\textgreater{}}
\end{Highlighting}
\end{Shaded}

Normal分布は下方に限界を持たないため、極端なマイナス金利も正の確率で許容する。これは計算上の柔軟性である一方、経済的には常に現実的とは限らない。

\subsection{5.3.2
インプライドNormalボラティリティ}\label{ux30a4ux30f3ux30d7ux30e9ux30a4ux30c9normalux30dcux30e9ux30c6ux30a3ux30eaux30c6ux30a3}

市場で観測されるオプション価格から、Normalモデルと整合するボラティリティを逆算したものがインプライドNormalボラティリティである。

市場価格を \(V_{\mathrm{mkt}}\) とすると、次式を満たす \(\sigma_N\)
を求める。

\[
V_N(
F_0,
K,
T,
D,
\sigma_N
)
=
V_{\mathrm{mkt}}
\]

GaussQuantには、Normal価格からボラティリティを直接逆算する関数が用意されている。

\begin{Shaded}
\begin{Highlighting}[]
\NormalTok{target\_npv }\OtherTok{\textless{}{-}} \FloatTok{0.1133} \SpecialCharTok{/} \DecValTok{100}

\NormalTok{implied\_normal\_volatility }\OtherTok{\textless{}{-}}\NormalTok{ GaussQuant}\SpecialCharTok{::}\FunctionTok{normal\_vol\_from\_price\_GQL}\NormalTok{(}
  \AttributeTok{option\_sign =}\NormalTok{ option\_sign,}
  \AttributeTok{strike =}\NormalTok{ strike\_rate,}
  \AttributeTok{forward =}\NormalTok{ forward\_rate,}
  \AttributeTok{maturity =}\NormalTok{ maturity\_year\_fraction,}
  \AttributeTok{option\_npv =}\NormalTok{ target\_npv,}
  \AttributeTok{discount\_factor =}\NormalTok{ maturity\_discount\_factor}
\NormalTok{)}

\NormalTok{repriced\_target }\OtherTok{\textless{}{-}}\NormalTok{ GaussQuant}\SpecialCharTok{::}\FunctionTok{normal\_option\_price\_GQL}\NormalTok{(}
  \AttributeTok{option\_sign =}\NormalTok{ option\_sign,}
  \AttributeTok{forward =}\NormalTok{ forward\_rate,}
  \AttributeTok{strike =}\NormalTok{ strike\_rate,}
  \AttributeTok{vol =}\NormalTok{ implied\_normal\_volatility,}
  \AttributeTok{maturity =}\NormalTok{ maturity\_year\_fraction,}
  \AttributeTok{discount\_factor =}\NormalTok{ maturity\_discount\_factor}
\NormalTok{)}

\NormalTok{implied\_volatility\_tbl }\OtherTok{\textless{}{-}}\NormalTok{ tibble}\SpecialCharTok{::}\FunctionTok{tribble}\NormalTok{(}
  \SpecialCharTok{\textasciitilde{}}\NormalTok{metric,                        }\SpecialCharTok{\textasciitilde{}}\NormalTok{value,}
  \StringTok{"target\_npv"}\NormalTok{,                   target\_npv,}
  \StringTok{"implied\_normal\_volatility"}\NormalTok{,    implied\_normal\_volatility,}
  \StringTok{"implied\_normal\_volatility\_bp"}\NormalTok{, implied\_normal\_volatility }\SpecialCharTok{*} \DecValTok{10000}\NormalTok{,}
  \StringTok{"repriced\_npv"}\NormalTok{,                 repriced\_target,}
  \StringTok{"repricing\_difference"}\NormalTok{,         repriced\_target }\SpecialCharTok{{-}}\NormalTok{ target\_npv}
\NormalTok{)}

\NormalTok{GaussQuant}\SpecialCharTok{::}\FunctionTok{show\_tbl\_GQH}\NormalTok{(}
\NormalTok{  implied\_volatility\_tbl,}
  \StringTok{"Implied normal volatility"}\NormalTok{,}
  \AttributeTok{n =} \DecValTok{20}
\NormalTok{)}
\CommentTok{\#\textgreater{} \# A tibble: 5 x 2}
\CommentTok{\#\textgreater{}   metric                          value}
\CommentTok{\#\textgreater{}   \textless{}chr\textgreater{}                           \textless{}dbl\textgreater{}}
\CommentTok{\#\textgreater{} 1 target\_npv                    0.00113}
\CommentTok{\#\textgreater{} 2 implied\_normal\_volatility     0.00499}
\CommentTok{\#\textgreater{} 3 implied\_normal\_volatility\_bp 49.9    }
\CommentTok{\#\textgreater{} 4 repriced\_npv                  0.00113}
\CommentTok{\#\textgreater{} 5 repricing\_difference          0}
\end{Highlighting}
\end{Shaded}

再評価値と市場価格の差が数値誤差の範囲に収まれば、逆算結果が正しいことを確認できる。

\subsection{5.3.3
金利・ボラティリティ・時間シナリオ}\label{ux91d1ux5229ux30dcux30e9ux30c6ux30a3ux30eaux30c6ux30a3ux6642ux9593ux30b7ux30caux30eaux30aa}

Normalモデルでは、フォワード金利とNormal
volatilityがいずれも絶対変化量で表される。このため、シナリオも1bp単位の加算として設定する方が分かりやすい。

\begin{Shaded}
\begin{Highlighting}[]
\NormalTok{normal\_calculator\_function }\OtherTok{\textless{}{-}}\NormalTok{ GaussQuant}\SpecialCharTok{::}\FunctionTok{make\_normal\_calculator\_GQL}\NormalTok{(}
  \AttributeTok{option\_sign =}\NormalTok{ option\_sign,}
  \AttributeTok{strike =}\NormalTok{ strike\_rate,}
  \AttributeTok{maturity =}\NormalTok{ maturity\_year\_fraction,}
  \AttributeTok{discount\_factor =}\NormalTok{ maturity\_discount\_factor,}
  \AttributeTok{forward =}\NormalTok{ forward\_rate,}
  \AttributeTok{vol =}\NormalTok{ normal\_volatility}
\NormalTok{)}

\NormalTok{normal\_calculator\_object }\OtherTok{\textless{}{-}}\NormalTok{ GaussQuant}\SpecialCharTok{::}\FunctionTok{normal\_calculator\_GQL}\NormalTok{(}
  \AttributeTok{option\_sign =}\NormalTok{ option\_sign,}
  \AttributeTok{strike =}\NormalTok{ strike\_rate,}
  \AttributeTok{maturity =}\NormalTok{ maturity\_year\_fraction,}
  \AttributeTok{discount\_factor =}\NormalTok{ maturity\_discount\_factor,}
  \AttributeTok{forward =}\NormalTok{ forward\_rate,}
  \AttributeTok{vol =}\NormalTok{ normal\_volatility}
\NormalTok{)}

\NormalTok{forward\_up\_1bp }\OtherTok{\textless{}{-}}\NormalTok{ forward\_rate }\SpecialCharTok{+} \DecValTok{1} \SpecialCharTok{/} \DecValTok{10000}
\NormalTok{volatility\_up\_1bp }\OtherTok{\textless{}{-}}\NormalTok{ normal\_volatility }\SpecialCharTok{+} \DecValTok{1} \SpecialCharTok{/} \DecValTok{10000}

\NormalTok{maturity\_after\_one\_day }\OtherTok{\textless{}{-}} \FunctionTok{max}\NormalTok{(}
\NormalTok{  maturity\_year\_fraction }\SpecialCharTok{{-}} \DecValTok{1} \SpecialCharTok{/} \DecValTok{365}\NormalTok{,}
  \DecValTok{0}
\NormalTok{)}

\NormalTok{discount\_factor\_after\_one\_day }\OtherTok{\textless{}{-}} \FunctionTok{exp}\NormalTok{(}
  \SpecialCharTok{{-}}\NormalTok{risk\_free\_rate }\SpecialCharTok{*}\NormalTok{ maturity\_after\_one\_day}
\NormalTok{)}

\NormalTok{closure\_scenario\_tbl }\OtherTok{\textless{}{-}}\NormalTok{ tibble}\SpecialCharTok{::}\FunctionTok{tribble}\NormalTok{(}
  \SpecialCharTok{\textasciitilde{}}\NormalTok{scenario,          }\SpecialCharTok{\textasciitilde{}}\NormalTok{npv,}
  \StringTok{"base"}\NormalTok{,              }\FunctionTok{normal\_calculator\_function}\NormalTok{(),}
  \StringTok{"forward\_up\_1bp"}\NormalTok{,    }\FunctionTok{normal\_calculator\_function}\NormalTok{(}
    \AttributeTok{forward\_new =}\NormalTok{ forward\_up\_1bp}
\NormalTok{  ),}
  \StringTok{"normal\_vol\_up\_1bp"}\NormalTok{, }\FunctionTok{normal\_calculator\_function}\NormalTok{(}
    \AttributeTok{vol\_new =}\NormalTok{ volatility\_up\_1bp}
\NormalTok{  ),}
  \StringTok{"one\_day\_passed"}\NormalTok{,    }\FunctionTok{normal\_calculator\_function}\NormalTok{(}
    \AttributeTok{maturity\_new =}\NormalTok{ maturity\_after\_one\_day,}
    \AttributeTok{discount\_factor\_new =}\NormalTok{ discount\_factor\_after\_one\_day}
\NormalTok{  )}
\NormalTok{) }\SpecialCharTok{|\textgreater{}}
\NormalTok{  dplyr}\SpecialCharTok{::}\FunctionTok{mutate}\NormalTok{(}
    \AttributeTok{difference\_from\_base =}\NormalTok{ .data}\SpecialCharTok{$}\NormalTok{npv }\SpecialCharTok{{-}}
\NormalTok{      .data}\SpecialCharTok{$}\NormalTok{npv[.data}\SpecialCharTok{$}\NormalTok{scenario }\SpecialCharTok{==} \StringTok{"base"}\NormalTok{]}
\NormalTok{  )}

\NormalTok{class\_style\_tbl }\OtherTok{\textless{}{-}}\NormalTok{ tibble}\SpecialCharTok{::}\FunctionTok{tribble}\NormalTok{(}
  \SpecialCharTok{\textasciitilde{}}\NormalTok{scenario,          }\SpecialCharTok{\textasciitilde{}}\NormalTok{npv,}
  \StringTok{"base"}\NormalTok{,              normal\_calculator\_object}\SpecialCharTok{$}\FunctionTok{npv}\NormalTok{(),}
  \StringTok{"forward\_up\_1bp"}\NormalTok{,    normal\_calculator\_object}\SpecialCharTok{$}\FunctionTok{npv}\NormalTok{(}
    \AttributeTok{forward\_new =}\NormalTok{ forward\_up\_1bp}
\NormalTok{  ),}
  \StringTok{"normal\_vol\_up\_1bp"}\NormalTok{, normal\_calculator\_object}\SpecialCharTok{$}\FunctionTok{npv}\NormalTok{(}
    \AttributeTok{vol\_new =}\NormalTok{ volatility\_up\_1bp}
\NormalTok{  ),}
  \StringTok{"one\_day\_passed"}\NormalTok{,    normal\_calculator\_object}\SpecialCharTok{$}\FunctionTok{npv}\NormalTok{(}
    \AttributeTok{maturity\_new =}\NormalTok{ maturity\_after\_one\_day,}
    \AttributeTok{discount\_factor\_new =}\NormalTok{ discount\_factor\_after\_one\_day}
\NormalTok{  )}
\NormalTok{)}

\NormalTok{GaussQuant}\SpecialCharTok{::}\FunctionTok{show\_tbl\_GQH}\NormalTok{(}
\NormalTok{  closure\_scenario\_tbl,}
  \StringTok{"Normal{-}model scenarios"}\NormalTok{,}
  \AttributeTok{n =} \DecValTok{20}
\NormalTok{)}
\CommentTok{\#\textgreater{} \# A tibble: 4 x 3}
\CommentTok{\#\textgreater{}   scenario              npv difference\_from\_base}
\CommentTok{\#\textgreater{}   \textless{}chr\textgreater{}               \textless{}dbl\textgreater{}                \textless{}dbl\textgreater{}}
\CommentTok{\#\textgreater{} 1 base              0.00113           0         }
\CommentTok{\#\textgreater{} 2 forward\_up\_1bp    0.00108          {-}0.0000553 }
\CommentTok{\#\textgreater{} 3 normal\_vol\_up\_1bp 0.00115           0.0000182 }
\CommentTok{\#\textgreater{} 4 one\_day\_passed    0.00113          {-}0.00000553}

\NormalTok{GaussQuant}\SpecialCharTok{::}\FunctionTok{show\_tbl\_GQH}\NormalTok{(}
\NormalTok{  class\_style\_tbl,}
  \StringTok{"Class{-}like normal calculator"}\NormalTok{,}
  \AttributeTok{n =} \DecValTok{20}
\NormalTok{)}
\CommentTok{\#\textgreater{} \# A tibble: 4 x 2}
\CommentTok{\#\textgreater{}   scenario              npv}
\CommentTok{\#\textgreater{}   \textless{}chr\textgreater{}               \textless{}dbl\textgreater{}}
\CommentTok{\#\textgreater{} 1 base              0.00113}
\CommentTok{\#\textgreater{} 2 forward\_up\_1bp    0.00108}
\CommentTok{\#\textgreater{} 3 normal\_vol\_up\_1bp 0.00115}
\CommentTok{\#\textgreater{} 4 one\_day\_passed    0.00113}
\end{Highlighting}
\end{Shaded}

先物価格のcallを金利putとして評価しているため、金利が1bp上昇するとオプション価値は低下する。Normal
volatilityが上昇すると、選択権の時間価値が増加する。

\section{5.4 CapletとCap}\label{capletux3068cap}

Capletは、将来の単一期間のフォワード金利に対するcall optionである。期間
\([T_{i-1},T_i]\) のフォワード金利を \(F_i\)、行使金利を
\(K\)、年率換算期間を \(\alpha_i\)、支払日の割引係数を \(P(0,T_i)\)
とする。

Capletの価格は次式となる。

\[
V_i^{\mathrm{caplet}}
=
N\alpha_iP(0,T_i)
\left[
(F_i-K)N(d_i)
+
\sigma_{N,i}\sqrt{\tau_i}\phi(d_i)
\right]
\]

ここで、

\[
d_i
=
\frac{
F_i-K
}{
\sigma_{N,i}\sqrt{\tau_i}
}
\]

である。Capは複数のCapletの合計として評価される。

\[
V^{\mathrm{cap}}
=
\sum_i
V_i^{\mathrm{caplet}}
\]

本節では、期間構造モデルを構築せず、各Capletのフォワード金利、Fixingまでの時間および割引係数を外生的に与える。

\begin{Shaded}
\begin{Highlighting}[]
\NormalTok{cap\_notional }\OtherTok{\textless{}{-}} \DecValTok{10000000}
\NormalTok{cap\_strike }\OtherTok{\textless{}{-}} \FloatTok{0.05}
\NormalTok{flat\_discount\_rate }\OtherTok{\textless{}{-}} \FloatTok{0.045}

\NormalTok{caplet\_input\_tbl }\OtherTok{\textless{}{-}}\NormalTok{ tibble}\SpecialCharTok{::}\FunctionTok{tribble}\NormalTok{(}
  \SpecialCharTok{\textasciitilde{}}\NormalTok{caplet, }\SpecialCharTok{\textasciitilde{}}\NormalTok{fixing\_time, }\SpecialCharTok{\textasciitilde{}}\NormalTok{payment\_time, }\SpecialCharTok{\textasciitilde{}}\NormalTok{accrual\_fraction, }\SpecialCharTok{\textasciitilde{}}\NormalTok{forward\_rate, }\SpecialCharTok{\textasciitilde{}}\NormalTok{normal\_vol\_bp,}
  \DecValTok{1}\NormalTok{L,       }\FloatTok{0.25}\NormalTok{,         }\FloatTok{0.50}\NormalTok{,          }\FloatTok{0.25}\NormalTok{,              }\FloatTok{0.0510}\NormalTok{,        }\DecValTok{80}\NormalTok{,}
  \DecValTok{2}\NormalTok{L,       }\FloatTok{0.50}\NormalTok{,         }\FloatTok{0.75}\NormalTok{,          }\FloatTok{0.25}\NormalTok{,              }\FloatTok{0.0505}\NormalTok{,        }\DecValTok{82}\NormalTok{,}
  \DecValTok{3}\NormalTok{L,       }\FloatTok{0.75}\NormalTok{,         }\FloatTok{1.00}\NormalTok{,          }\FloatTok{0.25}\NormalTok{,              }\FloatTok{0.0495}\NormalTok{,        }\DecValTok{84}\NormalTok{,}
  \DecValTok{4}\NormalTok{L,       }\FloatTok{1.00}\NormalTok{,         }\FloatTok{1.25}\NormalTok{,          }\FloatTok{0.25}\NormalTok{,              }\FloatTok{0.0485}\NormalTok{,        }\DecValTok{86}
\NormalTok{) }\SpecialCharTok{|\textgreater{}}
\NormalTok{  dplyr}\SpecialCharTok{::}\FunctionTok{mutate}\NormalTok{(}
    \AttributeTok{normal\_volatility =}\NormalTok{ .data}\SpecialCharTok{$}\NormalTok{normal\_vol\_bp }\SpecialCharTok{/} \DecValTok{10000}\NormalTok{,}
    \AttributeTok{discount\_factor =} \FunctionTok{exp}\NormalTok{(}
      \SpecialCharTok{{-}}\NormalTok{flat\_discount\_rate }\SpecialCharTok{*}\NormalTok{ .data}\SpecialCharTok{$}\NormalTok{payment\_time}
\NormalTok{    ),}
    \AttributeTok{option\_value\_per\_rate\_unit =}\NormalTok{ purrr}\SpecialCharTok{::}\FunctionTok{pmap\_dbl}\NormalTok{(}
      \FunctionTok{list}\NormalTok{(}
\NormalTok{        .data}\SpecialCharTok{$}\NormalTok{forward\_rate,}
\NormalTok{        .data}\SpecialCharTok{$}\NormalTok{normal\_volatility,}
\NormalTok{        .data}\SpecialCharTok{$}\NormalTok{fixing\_time,}
\NormalTok{        .data}\SpecialCharTok{$}\NormalTok{discount\_factor}
\NormalTok{      ),}
      \ControlFlowTok{function}\NormalTok{(}
\NormalTok{        forward\_rate,}
\NormalTok{        normal\_volatility,}
\NormalTok{        fixing\_time,}
\NormalTok{        discount\_factor}
\NormalTok{      ) \{}
\NormalTok{        GaussQuant}\SpecialCharTok{::}\FunctionTok{normal\_option\_price\_GQL}\NormalTok{(}
          \AttributeTok{option\_sign =} \DecValTok{1}\NormalTok{,}
          \AttributeTok{forward =}\NormalTok{ forward\_rate,}
          \AttributeTok{strike =}\NormalTok{ cap\_strike,}
          \AttributeTok{vol =}\NormalTok{ normal\_volatility,}
          \AttributeTok{maturity =}\NormalTok{ fixing\_time,}
          \AttributeTok{discount\_factor =}\NormalTok{ discount\_factor}
\NormalTok{        )}
\NormalTok{      \}}
\NormalTok{    ),}
    \AttributeTok{caplet\_npv =}\NormalTok{ cap\_notional }\SpecialCharTok{*}
\NormalTok{      .data}\SpecialCharTok{$}\NormalTok{accrual\_fraction }\SpecialCharTok{*}
\NormalTok{      .data}\SpecialCharTok{$}\NormalTok{option\_value\_per\_rate\_unit}
\NormalTok{  )}

\NormalTok{cap\_summary\_tbl }\OtherTok{\textless{}{-}}\NormalTok{ tibble}\SpecialCharTok{::}\FunctionTok{tribble}\NormalTok{(}
  \SpecialCharTok{\textasciitilde{}}\NormalTok{metric,       }\SpecialCharTok{\textasciitilde{}}\NormalTok{value,}
  \StringTok{"cap\_notional"}\NormalTok{, cap\_notional,}
  \StringTok{"cap\_strike"}\NormalTok{,   cap\_strike,}
  \StringTok{"cap\_npv"}\NormalTok{,      }\FunctionTok{sum}\NormalTok{(caplet\_input\_tbl}\SpecialCharTok{$}\NormalTok{caplet\_npv)}
\NormalTok{)}

\NormalTok{GaussQuant}\SpecialCharTok{::}\FunctionTok{show\_tbl\_GQH}\NormalTok{(}
\NormalTok{  caplet\_input\_tbl,}
  \StringTok{"Bachelier caplet decomposition"}\NormalTok{,}
  \AttributeTok{n =} \DecValTok{20}
\NormalTok{)}
\CommentTok{\#\textgreater{} \# A tibble: 4 x 10}
\CommentTok{\#\textgreater{}   caplet fixing\_time payment\_time accrual\_fraction forward\_rate normal\_vol\_bp}
\CommentTok{\#\textgreater{}    \textless{}int\textgreater{}       \textless{}dbl\textgreater{}        \textless{}dbl\textgreater{}            \textless{}dbl\textgreater{}        \textless{}dbl\textgreater{}         \textless{}dbl\textgreater{}}
\CommentTok{\#\textgreater{} 1      1        0.25         0.5              0.25       0.051             80}
\CommentTok{\#\textgreater{} 2      2        0.5          0.75             0.25       0.0505            82}
\CommentTok{\#\textgreater{} 3      3        0.75         1                0.25       0.0495            84}
\CommentTok{\#\textgreater{} 4      4        1            1.25             0.25       0.0485            86}
\CommentTok{\#\textgreater{} \# i 4 more variables: normal\_volatility \textless{}dbl\textgreater{}, discount\_factor \textless{}dbl\textgreater{},}
\CommentTok{\#\textgreater{} \#   option\_value\_per\_rate\_unit \textless{}dbl\textgreater{}, caplet\_npv \textless{}dbl\textgreater{}}

\NormalTok{GaussQuant}\SpecialCharTok{::}\FunctionTok{show\_tbl\_GQH}\NormalTok{(}
\NormalTok{  cap\_summary\_tbl,}
  \StringTok{"Bachelier cap summary"}\NormalTok{,}
  \AttributeTok{n =} \DecValTok{20}
\NormalTok{)}
\CommentTok{\#\textgreater{} \# A tibble: 3 x 2}
\CommentTok{\#\textgreater{}   metric             value}
\CommentTok{\#\textgreater{}   \textless{}chr\textgreater{}              \textless{}dbl\textgreater{}}
\CommentTok{\#\textgreater{} 1 cap\_notional 10000000   }
\CommentTok{\#\textgreater{} 2 cap\_strike          0.05}
\CommentTok{\#\textgreater{} 3 cap\_npv         24274.}
\end{Highlighting}
\end{Shaded}

各Capletは異なるFixing時点とフォワード金利を持つ。ここでは期間構造を動学的に生成せず、各フォワード金利をそれぞれの支払日フォワード測度の下で直接評価している。

\section{5.5 Swaption評価}\label{swaptionux8a55ux4fa1}

\subsection{5.5.1
フォワードスワップレートとアニュイティ}\label{ux30d5ux30a9ux30efux30fcux30c9ux30b9ux30efux30c3ux30d7ux30ecux30fcux30c8ux3068ux30a2ux30cbux30e5ux30a4ux30c6ux30a3}

Swaptionの原資産は、将来開始する金利スワップのフォワードスワップレートである。

固定レグ支払日を \(T_1,\ldots,T_n\)、年率換算期間を \(\alpha_i\)
とすると、スワップアニュイティは次式である。

\[
A(0)
=
\sum_{i=1}^{n}
\alpha_iP(0,T_i)
\]

スワップ開始日を \(T_0\)、満期を \(T_n\)
とすると、単一カーブの簡略化されたフォワードスワップレートは次式となる。

\[
S_0
=
\frac{
P(0,T_0)-P(0,T_n)
}{
A(0)
}
\]

本節では決定論的な割引係数を与え、フォワードスワップレートとアニュイティを計算する。

\begin{Shaded}
\begin{Highlighting}[]
\NormalTok{swap\_exercise\_time }\OtherTok{\textless{}{-}} \DecValTok{1}
\NormalTok{swap\_start\_time }\OtherTok{\textless{}{-}} \DecValTok{1}
\NormalTok{fixed\_payment\_times }\OtherTok{\textless{}{-}} \FunctionTok{c}\NormalTok{(}\DecValTok{2}\NormalTok{, }\DecValTok{3}\NormalTok{, }\DecValTok{4}\NormalTok{)}
\NormalTok{fixed\_accrual\_fractions }\OtherTok{\textless{}{-}} \FunctionTok{rep}\NormalTok{(}
  \DecValTok{1}\NormalTok{,}
  \FunctionTok{length}\NormalTok{(fixed\_payment\_times)}
\NormalTok{)}
\NormalTok{swap\_flat\_discount\_rate }\OtherTok{\textless{}{-}} \FloatTok{0.035}

\NormalTok{swap\_discount\_tbl }\OtherTok{\textless{}{-}}\NormalTok{ tibble}\SpecialCharTok{::}\FunctionTok{tibble}\NormalTok{(}
  \AttributeTok{payment\_time =}\NormalTok{ fixed\_payment\_times,}
  \AttributeTok{accrual\_fraction =}\NormalTok{ fixed\_accrual\_fractions,}
  \AttributeTok{discount\_factor =} \FunctionTok{exp}\NormalTok{(}
    \SpecialCharTok{{-}}\NormalTok{swap\_flat\_discount\_rate }\SpecialCharTok{*}\NormalTok{ fixed\_payment\_times}
\NormalTok{  )}
\NormalTok{)}

\NormalTok{swap\_annuity }\OtherTok{\textless{}{-}} \FunctionTok{sum}\NormalTok{(}
\NormalTok{  swap\_discount\_tbl}\SpecialCharTok{$}\NormalTok{accrual\_fraction }\SpecialCharTok{*}
\NormalTok{    swap\_discount\_tbl}\SpecialCharTok{$}\NormalTok{discount\_factor}
\NormalTok{)}

\NormalTok{swap\_start\_discount\_factor }\OtherTok{\textless{}{-}} \FunctionTok{exp}\NormalTok{(}
  \SpecialCharTok{{-}}\NormalTok{swap\_flat\_discount\_rate }\SpecialCharTok{*}\NormalTok{ swap\_start\_time}
\NormalTok{)}

\NormalTok{swap\_end\_discount\_factor }\OtherTok{\textless{}{-}} \FunctionTok{exp}\NormalTok{(}
  \SpecialCharTok{{-}}\NormalTok{swap\_flat\_discount\_rate }\SpecialCharTok{*}
    \FunctionTok{max}\NormalTok{(fixed\_payment\_times)}
\NormalTok{)}

\NormalTok{forward\_swap\_rate }\OtherTok{\textless{}{-}}\NormalTok{ (}
\NormalTok{  swap\_start\_discount\_factor }\SpecialCharTok{{-}}
\NormalTok{    swap\_end\_discount\_factor}
\NormalTok{) }\SpecialCharTok{/}\NormalTok{ swap\_annuity}

\NormalTok{forward\_swap\_tbl }\OtherTok{\textless{}{-}}\NormalTok{ tibble}\SpecialCharTok{::}\FunctionTok{tribble}\NormalTok{(}
  \SpecialCharTok{\textasciitilde{}}\NormalTok{metric,                       }\SpecialCharTok{\textasciitilde{}}\NormalTok{value,}
  \StringTok{"swap\_start\_discount\_factor"}\NormalTok{,   swap\_start\_discount\_factor,}
  \StringTok{"swap\_end\_discount\_factor"}\NormalTok{,     swap\_end\_discount\_factor,}
  \StringTok{"swap\_annuity"}\NormalTok{,                 swap\_annuity,}
  \StringTok{"forward\_swap\_rate"}\NormalTok{,            forward\_swap\_rate}
\NormalTok{)}

\NormalTok{GaussQuant}\SpecialCharTok{::}\FunctionTok{show\_tbl\_GQH}\NormalTok{(}
\NormalTok{  swap\_discount\_tbl,}
  \StringTok{"Underlying swap fixed{-}leg schedule"}\NormalTok{,}
  \AttributeTok{n =} \DecValTok{20}
\NormalTok{)}
\CommentTok{\#\textgreater{} \# A tibble: 3 x 3}
\CommentTok{\#\textgreater{}   payment\_time accrual\_fraction discount\_factor}
\CommentTok{\#\textgreater{}          \textless{}dbl\textgreater{}            \textless{}dbl\textgreater{}           \textless{}dbl\textgreater{}}
\CommentTok{\#\textgreater{} 1            2                1           0.932}
\CommentTok{\#\textgreater{} 2            3                1           0.900}
\CommentTok{\#\textgreater{} 3            4                1           0.869}

\NormalTok{GaussQuant}\SpecialCharTok{::}\FunctionTok{show\_tbl\_GQH}\NormalTok{(}
\NormalTok{  forward\_swap\_tbl,}
  \StringTok{"Forward swap rate and annuity"}\NormalTok{,}
  \AttributeTok{n =} \DecValTok{20}
\NormalTok{)}
\CommentTok{\#\textgreater{} \# A tibble: 4 x 2}
\CommentTok{\#\textgreater{}   metric                      value}
\CommentTok{\#\textgreater{}   \textless{}chr\textgreater{}                       \textless{}dbl\textgreater{}}
\CommentTok{\#\textgreater{} 1 swap\_start\_discount\_factor 0.966 }
\CommentTok{\#\textgreater{} 2 swap\_end\_discount\_factor   0.869 }
\CommentTok{\#\textgreater{} 3 swap\_annuity               2.70  }
\CommentTok{\#\textgreater{} 4 forward\_swap\_rate          0.0356}
\end{Highlighting}
\end{Shaded}

\subsection{5.5.2 Bachelier Swaption}\label{bachelier-swaption}

payer
swaptionは、将来の固定金利支払いスワップへ入る権利である。フォワードスワップレートが行使金利を上回ると価値を持つため、フォワードスワップレートに対するcall
optionとして評価できる。

swap
measureの下でフォワードスワップレートをマルチンゲールとし、Normalモデルを仮定する。

\[
dS_t
=
\sigma_S\,dW_t^A
\]

payer swaptionの価格は次式となる。

\[
V_{\mathrm{payer}}
=
NA(0)
\left[
(S_0-K)N(d)
+
\sigma_S\sqrt{T_e}\phi(d)
\right]
\]

ここで、

\[
d
=
\frac{
S_0-K
}{
\sigma_S\sqrt{T_e}
}
\]

である。receiver swaptionはput optionとして評価する。

\begin{Shaded}
\begin{Highlighting}[]
\NormalTok{swaption\_notional }\OtherTok{\textless{}{-}} \DecValTok{10000000}
\NormalTok{swaption\_strike }\OtherTok{\textless{}{-}} \FloatTok{0.034}
\NormalTok{swaption\_normal\_volatility }\OtherTok{\textless{}{-}} \DecValTok{100} \SpecialCharTok{/} \DecValTok{10000}

\NormalTok{payer\_option\_value\_per\_annuity }\OtherTok{\textless{}{-}}\NormalTok{ GaussQuant}\SpecialCharTok{::}\FunctionTok{normal\_option\_price\_GQL}\NormalTok{(}
  \AttributeTok{option\_sign =} \DecValTok{1}\NormalTok{,}
  \AttributeTok{forward =}\NormalTok{ forward\_swap\_rate,}
  \AttributeTok{strike =}\NormalTok{ swaption\_strike,}
  \AttributeTok{vol =}\NormalTok{ swaption\_normal\_volatility,}
  \AttributeTok{maturity =}\NormalTok{ swap\_exercise\_time,}
  \AttributeTok{discount\_factor =} \DecValTok{1}
\NormalTok{)}

\NormalTok{receiver\_option\_value\_per\_annuity }\OtherTok{\textless{}{-}}\NormalTok{ GaussQuant}\SpecialCharTok{::}\FunctionTok{normal\_option\_price\_GQL}\NormalTok{(}
  \AttributeTok{option\_sign =} \SpecialCharTok{{-}}\DecValTok{1}\NormalTok{,}
  \AttributeTok{forward =}\NormalTok{ forward\_swap\_rate,}
  \AttributeTok{strike =}\NormalTok{ swaption\_strike,}
  \AttributeTok{vol =}\NormalTok{ swaption\_normal\_volatility,}
  \AttributeTok{maturity =}\NormalTok{ swap\_exercise\_time,}
  \AttributeTok{discount\_factor =} \DecValTok{1}
\NormalTok{)}

\NormalTok{swaption\_tbl }\OtherTok{\textless{}{-}}\NormalTok{ tibble}\SpecialCharTok{::}\FunctionTok{tribble}\NormalTok{(}
  \SpecialCharTok{\textasciitilde{}}\NormalTok{option\_type, }\SpecialCharTok{\textasciitilde{}}\NormalTok{option\_sign, }\SpecialCharTok{\textasciitilde{}}\NormalTok{npv,}
  \StringTok{"payer"}\NormalTok{,       }\DecValTok{1}\NormalTok{,           swaption\_notional }\SpecialCharTok{*}
\NormalTok{                               swap\_annuity }\SpecialCharTok{*}
\NormalTok{                               payer\_option\_value\_per\_annuity,}
  \StringTok{"receiver"}\NormalTok{,   }\SpecialCharTok{{-}}\DecValTok{1}\NormalTok{,           swaption\_notional }\SpecialCharTok{*}
\NormalTok{                               swap\_annuity }\SpecialCharTok{*}
\NormalTok{                               receiver\_option\_value\_per\_annuity}
\NormalTok{) }\SpecialCharTok{|\textgreater{}}
\NormalTok{  dplyr}\SpecialCharTok{::}\FunctionTok{mutate}\NormalTok{(}
    \AttributeTok{forward\_swap\_rate =}\NormalTok{ forward\_swap\_rate,}
    \AttributeTok{strike =}\NormalTok{ swaption\_strike,}
    \AttributeTok{annuity =}\NormalTok{ swap\_annuity,}
    \AttributeTok{normal\_volatility =}\NormalTok{ swaption\_normal\_volatility}
\NormalTok{  ) }\SpecialCharTok{|\textgreater{}}
\NormalTok{  dplyr}\SpecialCharTok{::}\FunctionTok{relocate}\NormalTok{(}
\NormalTok{    .data}\SpecialCharTok{$}\NormalTok{option\_type,}
\NormalTok{    .data}\SpecialCharTok{$}\NormalTok{forward\_swap\_rate,}
\NormalTok{    .data}\SpecialCharTok{$}\NormalTok{strike,}
\NormalTok{    .data}\SpecialCharTok{$}\NormalTok{annuity,}
\NormalTok{    .data}\SpecialCharTok{$}\NormalTok{normal\_volatility,}
\NormalTok{    .data}\SpecialCharTok{$}\NormalTok{npv}
\NormalTok{  )}

\NormalTok{GaussQuant}\SpecialCharTok{::}\FunctionTok{show\_tbl\_GQH}\NormalTok{(}
\NormalTok{  swaption\_tbl,}
  \StringTok{"Bachelier swaption summary"}\NormalTok{,}
  \AttributeTok{n =} \DecValTok{20}
\NormalTok{)}
\CommentTok{\#\textgreater{} \# A tibble: 2 x 7}
\CommentTok{\#\textgreater{}   option\_type forward\_swap\_rate strike annuity normal\_volatility     npv}
\CommentTok{\#\textgreater{}   \textless{}chr\textgreater{}                   \textless{}dbl\textgreater{}  \textless{}dbl\textgreater{}   \textless{}dbl\textgreater{}             \textless{}dbl\textgreater{}   \textless{}dbl\textgreater{}}
\CommentTok{\#\textgreater{} 1 payer                  0.0356  0.034    2.70              0.01 131091.}
\CommentTok{\#\textgreater{} 2 receiver               0.0356  0.034    2.70              0.01  87325.}
\CommentTok{\#\textgreater{} \# i 1 more variable: option\_sign \textless{}dbl\textgreater{}}
\end{Highlighting}
\end{Shaded}

アニュイティ \(A(0)\)
はすでに将来の固定レグ支払いを現在価値へ割り引いた量である。そのため、Normal公式側では割引係数を1とし、結果へ元本とアニュイティを掛けている。

\section{5.6 Black volatilityとNormal
volatility}\label{black-volatilityux3068normal-volatility}

Blackモデルでは、ATM付近のオプション価値は概ね次式に比例する。

\[
V_{\mathrm{Black,ATM}}
\approx
D(0,T)
F_0\sigma_B\sqrt{T}\phi(0)
\]

Normalモデルでは次式となる。

\[
V_{\mathrm{Normal,ATM}}
\approx
D(0,T)
\sigma_N\sqrt{T}\phi(0)
\]

したがって、ATM付近では両モデルのボラティリティに次の近似関係がある。

\[
\sigma_N
\approx
F_0\sigma_B
\]

Normal volatilityは金利水準と同じ単位を持つため、同じBlack
volatilityであっても、金利水準が低下すると対応するNormal
volatilityは小さくなる。

\begin{Shaded}
\begin{Highlighting}[]
\NormalTok{black\_volatility }\OtherTok{\textless{}{-}} \FloatTok{0.20}

\NormalTok{black\_normal\_comparison\_tbl }\OtherTok{\textless{}{-}}\NormalTok{ tibble}\SpecialCharTok{::}\FunctionTok{tribble}\NormalTok{(}
  \SpecialCharTok{\textasciitilde{}}\NormalTok{forward\_rate, }\SpecialCharTok{\textasciitilde{}}\NormalTok{black\_volatility,}
  \FloatTok{0.0500}\NormalTok{,         black\_volatility,}
  \FloatTok{0.0200}\NormalTok{,         black\_volatility,}
  \FloatTok{0.0050}\NormalTok{,         black\_volatility}
\NormalTok{) }\SpecialCharTok{|\textgreater{}}
\NormalTok{  dplyr}\SpecialCharTok{::}\FunctionTok{mutate}\NormalTok{(}
    \AttributeTok{approximate\_normal\_volatility =}
\NormalTok{      .data}\SpecialCharTok{$}\NormalTok{forward\_rate }\SpecialCharTok{*}
\NormalTok{      .data}\SpecialCharTok{$}\NormalTok{black\_volatility,}
    \AttributeTok{approximate\_normal\_volatility\_bp =}
\NormalTok{      .data}\SpecialCharTok{$}\NormalTok{approximate\_normal\_volatility }\SpecialCharTok{*}
      \DecValTok{10000}
\NormalTok{  )}

\NormalTok{GaussQuant}\SpecialCharTok{::}\FunctionTok{show\_tbl\_GQH}\NormalTok{(}
\NormalTok{  black\_normal\_comparison\_tbl,}
  \StringTok{"Approximate ATM Black{-}to{-}Normal volatility conversion"}\NormalTok{,}
  \AttributeTok{n =} \DecValTok{20}
\NormalTok{)}
\CommentTok{\#\textgreater{} \# A tibble: 3 x 4}
\CommentTok{\#\textgreater{}   forward\_rate black\_volatility approximate\_normal\_vola\textasciitilde{}1 approximate\_normal\_v\textasciitilde{}2}
\CommentTok{\#\textgreater{}          \textless{}dbl\textgreater{}            \textless{}dbl\textgreater{}                     \textless{}dbl\textgreater{}                  \textless{}dbl\textgreater{}}
\CommentTok{\#\textgreater{} 1        0.05               0.2                     0.01                     100}
\CommentTok{\#\textgreater{} 2        0.02               0.2                     0.004                     40}
\CommentTok{\#\textgreater{} 3        0.005              0.2                     0.001                     10}
\CommentTok{\#\textgreater{} \# i abbreviated names: 1: approximate\_normal\_volatility,}
\CommentTok{\#\textgreater{} \#   2: approximate\_normal\_volatility\_bp}
\end{Highlighting}
\end{Shaded}

この関係はATM近似であり、StrikeがATMから離れる場合やボラティリティスマイルが存在する場合には、単純な換算だけでは十分でない。

\section{5.7 まとめ}\label{ux307eux3068ux3081-9}

本章では、金利期間構造全体の確率モデルを与えず、各商品の評価に必要なフォワード金利またはフォワードスワップレートを、適切な測度の下でマルチンゲールとして直接モデル化した。

Normalモデルでは、フォワード金利の変化額が正規分布に従う。

\[
dF_t
=
\sigma_N\,dW_t
\]

このため、フォワード金利がゼロまたはマイナスとなる場合でも、価格公式を変更せずに評価できる。

Normal
volatilityは相対変化率ではなく、金利と同じ単位を持つ。金利市場では通常、ベーシスポイント単位で表示される。Normal
volatility 50bpは小数表示で0.0050である。

金利先物オプションでは、先物価格から金利へ変換するとPayoffの向きが反転する。先物価格のcallは、金利を原資産とした場合にはputとして評価される。

Capは複数のCapletの合計であり、各Capletは対応するフォワード金利を原資産とする。Swaptionはフォワードスワップレートを原資産とし、スワップアニュイティを価格倍率として使用する。

Normalモデルはマイナス金利を扱える一方、金利水準にかかわらず同じ絶対変動幅を仮定し、極端な負の金利も許容する。実際の市場では、Black、Shifted
Black、NormalおよびSABRなどを、金利水準とボラティリティ市場の慣行に応じて使い分ける。

\begin{Shaded}
\begin{Highlighting}[]
\NormalTok{chapter05\_summary\_tbl }\OtherTok{\textless{}{-}}\NormalTok{ tibble}\SpecialCharTok{::}\FunctionTok{tribble}\NormalTok{(}
  \SpecialCharTok{\textasciitilde{}}\NormalTok{section,                    }\SpecialCharTok{\textasciitilde{}}\NormalTok{main\_result,}
  \StringTok{"Normal dynamics"}\NormalTok{,           }\StringTok{"Additive forward{-}rate process"}\NormalTok{,}
  \StringTok{"Futures option"}\NormalTok{,            }\StringTok{"Rate conversion, price, and Greeks"}\NormalTok{,}
  \StringTok{"Implied normal volatility"}\NormalTok{, }\StringTok{"Volatility inversion and repricing"}\NormalTok{,}
  \StringTok{"Negative rates"}\NormalTok{,            }\StringTok{"Pricing without a logarithmic restriction"}\NormalTok{,}
  \StringTok{"Cap"}\NormalTok{,                       }\StringTok{"Sum of Bachelier caplets"}\NormalTok{,}
  \StringTok{"Swaption"}\NormalTok{,                  }\StringTok{"Forward swap rate times annuity"}
\NormalTok{)}

\NormalTok{GaussQuant}\SpecialCharTok{::}\FunctionTok{show\_tbl\_GQH}\NormalTok{(}
\NormalTok{  chapter05\_summary\_tbl,}
  \StringTok{"Chapter II{-}5 summary"}\NormalTok{,}
  \AttributeTok{n =} \DecValTok{20}
\NormalTok{)}
\CommentTok{\#\textgreater{} \# A tibble: 6 x 2}
\CommentTok{\#\textgreater{}   section                   main\_result                              }
\CommentTok{\#\textgreater{}   \textless{}chr\textgreater{}                     \textless{}chr\textgreater{}                                    }
\CommentTok{\#\textgreater{} 1 Normal dynamics           Additive forward{-}rate process            }
\CommentTok{\#\textgreater{} 2 Futures option            Rate conversion, price, and Greeks       }
\CommentTok{\#\textgreater{} 3 Implied normal volatility Volatility inversion and repricing       }
\CommentTok{\#\textgreater{} 4 Negative rates            Pricing without a logarithmic restriction}
\CommentTok{\#\textgreater{} 5 Cap                       Sum of Bachelier caplets                 }
\CommentTok{\#\textgreater{} 6 Swaption                  Forward swap rate times annuity}

\FunctionTok{cat}\NormalTok{(}
  \StringTok{"}\SpecialCharTok{\textbackslash{}n}\StringTok{Chapter II{-}5 Normal Bachelier model completed successfully.}\SpecialCharTok{\textbackslash{}n}\StringTok{"}
\NormalTok{)}
\CommentTok{\#\textgreater{} }
\CommentTok{\#\textgreater{} Chapter II{-}5 Normal Bachelier model completed successfully.}
\end{Highlighting}
\end{Shaded}

\chapter{第II-6章
SABRモデル}\label{ux7b2cii-6ux7ae0-sabrux30e2ux30c7ux30eb}

前章までのBlackモデルとNormalモデルでは、満期とStrikeが同じオプションに対して、一つのボラティリティを与えた。しかし、実際のオプション市場では、同じ満期であってもStrikeごとに異なるインプライド・ボラティリティが観測される。このStrike方向の形状をボラティリティ・スマイルまたはスキューと呼ぶ。

SABRモデルは、フォワード価格またはフォワード金利と、そのボラティリティを同時に確率変数として扱う。モデル名は、Stochastic
Alpha Beta Rhoの頭文字に由来する。

フォワード \(F_t\) とボラティリティ係数 \(\alpha_t\)
の動学を、概念的に次のように表す。

\[
dF_t
=
\alpha_t F_t^\beta\,dW_t^{(1)}
\]

\[
d\alpha_t
=
\nu\alpha_t\,dW_t^{(2)}
\]

\[
dW_t^{(1)}dW_t^{(2)}
=
\rho\,dt
\]

SABRモデルには、\(\alpha\)、\(\beta\)、\(\nu\)、\(\rho\)
の4パラメータがある。

\begin{itemize}
\tightlist
\item
  \(\alpha\) はボラティリティ水準を主に決定する。
\item
  \(\beta\) はフォワード水準に対する変動幅の依存度を表す。
\item
  \(\nu\)
  はボラティリティ自体の変動性であり、スマイルの曲率に強く影響する。
\item
  \(\rho\)
  はフォワードとボラティリティの相関であり、スマイルの左右非対称性に影響する。
\end{itemize}

\(\beta=1\) は対数正規型に近く、\(\beta=0\)
はNormal型に近い。実務では4パラメータを同時推定すると不安定になりやすいため、\(\beta\)
をあらかじめ固定し、\(\alpha\)、\(\nu\)、\(\rho\)
を市場ボラティリティへカリブレーションすることが多い。

本章では、まずBlack
volatilityのスマイルへSABRモデルを適合させる。次に各パラメータがスマイルの形状へ与える影響を確認する。最後に、負のStrikeを含むNormal
volatilityの市場データに対してShifted SABRを適用する。

\begin{Shaded}
\begin{Highlighting}[]
\NormalTok{knitr}\SpecialCharTok{::}\NormalTok{opts\_chunk}\SpecialCharTok{$}\FunctionTok{set}\NormalTok{(}
  \AttributeTok{collapse =} \ConstantTok{TRUE}\NormalTok{,}
  \AttributeTok{comment =} \StringTok{"\#\textgreater{}"}\NormalTok{,}
  \AttributeTok{warning =} \ConstantTok{FALSE}\NormalTok{,}
  \AttributeTok{message =} \ConstantTok{FALSE}
\NormalTok{)}


\FunctionTok{suppressPackageStartupMessages}\NormalTok{(\{}
  \FunctionTok{library}\NormalTok{(tidyverse)}
  \FunctionTok{library}\NormalTok{(ggthemes)}
  \FunctionTok{library}\NormalTok{(QuantLib)}
  \FunctionTok{library}\NormalTok{(GaussQuant)}
  \FunctionTok{library}\NormalTok{(tidyverse)}
  \FunctionTok{library}\NormalTok{(broom)}
\NormalTok{\})}
\end{Highlighting}
\end{Shaded}

\section{6.1
SABRモデルとパラメータ}\label{sabrux30e2ux30c7ux30ebux3068ux30d1ux30e9ux30e1ux30fcux30bf}

\subsection{6.1.1
SABRパラメータの役割}\label{sabrux30d1ux30e9ux30e1ux30fcux30bfux306eux5f79ux5272}

SABRの各パラメータは完全に独立してスマイルを動かすわけではないが、主な役割は次のように整理できる。

\subsection{6.1.2 Alpha}\label{alpha}

\(\alpha\)
は初期ボラティリティ水準である。他の条件を固定して\(\alpha\)を上げると、スマイル全体が概ね上方へ移動する。

\subsection{6.1.3 Beta}\label{beta}

\(\beta\) はフォワード水準に対する弾性を表す。

\[
\text{instantaneous volatility}
=
\alpha_tF_t^\beta
\]

\(\beta=1\)
では変動幅がフォワード水準に比例し、Blackモデルに近い。\(\beta=0\)
では変動幅がフォワード水準に依存せず、Normalモデルに近い。

\subsection{6.1.4 Volatility of
volatility}\label{volatility-of-volatility}

\(\nu\) はvolatility of
volatility、すなわちvol-of-volである。\(\nu\)が大きいほど、ボラティリティが将来大きく変化する可能性が高まり、一般にスマイルの曲率が強くなる。

本章のコードでは、変数名として \texttt{volvol} を使用する。

\subsection{6.1.5 Rho}\label{rho}

\(\rho\)
はフォワード変化とボラティリティ変化の相関である。負の\(\rho\)は、フォワード低下時にボラティリティが上昇しやすい構造を表し、低Strike側のボラティリティを相対的に高くする。

\subsection{6.1.6
Hagan近似とカリブレーション}\label{haganux8fd1ux4f3cux3068ux30abux30eaux30d6ux30ecux30fcux30b7ux30e7ux30f3}

SABRモデルの市場利用では、Haganらによる近似式を用いて、SABRパラメータからBlack
volatilityまたはNormal volatilityを計算することが多い。

市場Strikeを \(K_i\)、市場ボラティリティを \(\sigma_i^{\mathrm{mkt}}\)、
SABR近似値を \(\sigma_i^{\mathrm{SABR}}\)
とする。本章では、適合度を次のRMSEで測定する。

\[
\mathrm{RMSE}
=
\sqrt{
\frac{1}{n}
\sum_{i=1}^{n}
\left(
\sigma_i^{\mathrm{mkt}}
-
\sigma_i^{\mathrm{SABR}}
\right)^2
}
\]

\(\beta\)を固定し、RMSEが最小となる \(\alpha\)、\(\nu\)、\(\rho\)
を数値的に求める。

\section{6.2 Black
SABRのカリブレーション}\label{black-sabrux306eux30abux30eaux30d6ux30ecux30fcux30b7ux30e7ux30f3}

\subsection{6.2.1 市場データ}\label{ux5e02ux5834ux30c7ux30fcux30bf}

フォワードを0.56\%、満期を2年、\(\beta=0.5\)とする。市場Black
volatilityはStrikeごとに異なり、明確なスマイルを形成している。

\begin{Shaded}
\begin{Highlighting}[]
\NormalTok{black\_market\_data }\OtherTok{\textless{}{-}}\NormalTok{ tibble}\SpecialCharTok{::}\FunctionTok{tribble}\NormalTok{(}
  \SpecialCharTok{\textasciitilde{}}\NormalTok{strike\_pct, }\SpecialCharTok{\textasciitilde{}}\NormalTok{market\_vol\_pct,}
  \FloatTok{0.06}\NormalTok{,         }\FloatTok{90.78}\NormalTok{,}
  \FloatTok{0.31}\NormalTok{,         }\FloatTok{46.09}\NormalTok{,}
  \FloatTok{0.56}\NormalTok{,         }\FloatTok{45.30}\NormalTok{,}
  \FloatTok{0.81}\NormalTok{,         }\FloatTok{50.17}\NormalTok{,}
  \FloatTok{1.06}\NormalTok{,         }\FloatTok{53.85}\NormalTok{,}
  \FloatTok{1.56}\NormalTok{,         }\FloatTok{58.42}\NormalTok{,}
  \FloatTok{2.56}\NormalTok{,         }\FloatTok{62.72}
\NormalTok{) }\SpecialCharTok{|\textgreater{}}
\NormalTok{  dplyr}\SpecialCharTok{::}\FunctionTok{mutate}\NormalTok{(}
    \AttributeTok{strike =}\NormalTok{ .data}\SpecialCharTok{$}\NormalTok{strike\_pct }\SpecialCharTok{/} \DecValTok{100}\NormalTok{,}
    \AttributeTok{market\_vol =}\NormalTok{ .data}\SpecialCharTok{$}\NormalTok{market\_vol\_pct }\SpecialCharTok{/} \DecValTok{100}
\NormalTok{  )}

\NormalTok{forward\_black }\OtherTok{\textless{}{-}} \FloatTok{0.0056}
\NormalTok{maturity\_black }\OtherTok{\textless{}{-}} \DecValTok{2}
\NormalTok{beta\_black }\OtherTok{\textless{}{-}} \FloatTok{0.5}

\NormalTok{black\_model\_inputs }\OtherTok{\textless{}{-}}\NormalTok{ tibble}\SpecialCharTok{::}\FunctionTok{tribble}\NormalTok{(}
  \SpecialCharTok{\textasciitilde{}}\NormalTok{input,       }\SpecialCharTok{\textasciitilde{}}\NormalTok{value,}
  \StringTok{"forward"}\NormalTok{,     forward\_black,}
  \StringTok{"maturity"}\NormalTok{,    maturity\_black,}
  \StringTok{"fixed\_beta"}\NormalTok{,  beta\_black}
\NormalTok{)}

\NormalTok{initial\_parameters\_black }\OtherTok{\textless{}{-}} \FunctionTok{c}\NormalTok{(}
  \AttributeTok{alpha =} \FloatTok{0.1}\NormalTok{,}
  \AttributeTok{volvol =} \FloatTok{0.1}\NormalTok{,}
  \AttributeTok{rho =} \FloatTok{0.1}
\NormalTok{)}

\NormalTok{GaussQuant}\SpecialCharTok{::}\FunctionTok{show\_tbl\_GQH}\NormalTok{(}
\NormalTok{  black\_market\_data }\SpecialCharTok{|\textgreater{}}
\NormalTok{    dplyr}\SpecialCharTok{::}\FunctionTok{select}\NormalTok{(}
\NormalTok{      .data}\SpecialCharTok{$}\NormalTok{strike,}
\NormalTok{      .data}\SpecialCharTok{$}\NormalTok{market\_vol}
\NormalTok{    ),}
  \StringTok{"Black SABR market data"}\NormalTok{,}
  \AttributeTok{n =} \DecValTok{20}
\NormalTok{)}
\CommentTok{\#\textgreater{} \# A tibble: 7 x 2}
\CommentTok{\#\textgreater{}   strike market\_vol}
\CommentTok{\#\textgreater{}    \textless{}dbl\textgreater{}      \textless{}dbl\textgreater{}}
\CommentTok{\#\textgreater{} 1 0.0006      0.908}
\CommentTok{\#\textgreater{} 2 0.0031      0.461}
\CommentTok{\#\textgreater{} 3 0.0056      0.453}
\CommentTok{\#\textgreater{} 4 0.0081      0.502}
\CommentTok{\#\textgreater{} 5 0.0106      0.538}
\CommentTok{\#\textgreater{} 6 0.0156      0.584}
\CommentTok{\#\textgreater{} 7 0.0256      0.627}

\NormalTok{GaussQuant}\SpecialCharTok{::}\FunctionTok{show\_tbl\_GQH}\NormalTok{(}
\NormalTok{  black\_model\_inputs,}
  \StringTok{"Black SABR model inputs"}\NormalTok{,}
  \AttributeTok{n =} \DecValTok{20}
\NormalTok{)}
\CommentTok{\#\textgreater{} \# A tibble: 3 x 2}
\CommentTok{\#\textgreater{}   input       value}
\CommentTok{\#\textgreater{}   \textless{}chr\textgreater{}       \textless{}dbl\textgreater{}}
\CommentTok{\#\textgreater{} 1 forward    0.0056}
\CommentTok{\#\textgreater{} 2 maturity   2     }
\CommentTok{\#\textgreater{} 3 fixed\_beta 0.5}
\end{Highlighting}
\end{Shaded}

ATMに近いStrikeは0.56\%であり、フォワード0.56\%と一致する。ATM付近の水準は主に\(\alpha\)に影響し、両翼の形状が\(\nu\)と\(\rho\)の推定に影響する。

\subsection{6.2.2
パラメータ・カリブレーション}\label{ux30d1ux30e9ux30e1ux30fcux30bfux30abux30eaux30d6ux30ecux30fcux30b7ux30e7ux30f3}

目的関数では、各StrikeにおけるQuantLibのSABR Black volatilityと市場Black
volatilityのRMSEを計算する。

\begin{Shaded}
\begin{Highlighting}[]
\NormalTok{objective\_sabr\_black }\OtherTok{\textless{}{-}} \ControlFlowTok{function}\NormalTok{(parameters) \{}
\NormalTok{  alpha }\OtherTok{\textless{}{-}}\NormalTok{ parameters[[}\DecValTok{1}\NormalTok{]]}
\NormalTok{  volvol }\OtherTok{\textless{}{-}}\NormalTok{ parameters[[}\DecValTok{2}\NormalTok{]]}
\NormalTok{  rho }\OtherTok{\textless{}{-}}\NormalTok{ parameters[[}\DecValTok{3}\NormalTok{]]}

  \ControlFlowTok{if}\NormalTok{ (}
\NormalTok{    alpha }\SpecialCharTok{\textless{}=} \DecValTok{0} \SpecialCharTok{||}
\NormalTok{    volvol }\SpecialCharTok{\textless{}} \DecValTok{0} \SpecialCharTok{||}
\NormalTok{    rho }\SpecialCharTok{\textless{}=} \SpecialCharTok{{-}}\FloatTok{0.9999} \SpecialCharTok{||}
\NormalTok{    rho }\SpecialCharTok{\textgreater{}=} \FloatTok{0.9999}
\NormalTok{  ) \{}
    \FunctionTok{return}\NormalTok{(}\FloatTok{1e6}\NormalTok{)}
\NormalTok{  \}}

\NormalTok{  calculated\_volatility }\OtherTok{\textless{}{-}}\NormalTok{ purrr}\SpecialCharTok{::}\FunctionTok{map\_dbl}\NormalTok{(}
\NormalTok{    black\_market\_data}\SpecialCharTok{$}\NormalTok{strike,}
    \ControlFlowTok{function}\NormalTok{(strike\_value) \{}
\NormalTok{      GaussQuant}\SpecialCharTok{::}\FunctionTok{safe\_sabr\_vol\_GQL}\NormalTok{(}
        \AttributeTok{strike =}\NormalTok{ strike\_value,}
        \AttributeTok{forward =}\NormalTok{ forward\_black,}
        \AttributeTok{maturity =}\NormalTok{ maturity\_black,}
        \AttributeTok{alpha =}\NormalTok{ alpha,}
        \AttributeTok{beta =}\NormalTok{ beta\_black,}
        \AttributeTok{volvol =}\NormalTok{ volvol,}
        \AttributeTok{rho =}\NormalTok{ rho}
\NormalTok{      )}
\NormalTok{    \}}
\NormalTok{  )}

  \ControlFlowTok{if}\NormalTok{ (}
    \FunctionTok{any}\NormalTok{(}\FunctionTok{is.na}\NormalTok{(calculated\_volatility)) }\SpecialCharTok{||}
    \FunctionTok{any}\NormalTok{(}\SpecialCharTok{!}\FunctionTok{is.finite}\NormalTok{(calculated\_volatility))}
\NormalTok{  ) \{}
    \FunctionTok{return}\NormalTok{(}\FloatTok{1e6}\NormalTok{)}
\NormalTok{  \}}

\NormalTok{  GaussQuant}\SpecialCharTok{::}\FunctionTok{rmse\_GQH}\NormalTok{(}
    \AttributeTok{actual =}\NormalTok{ black\_market\_data}\SpecialCharTok{$}\NormalTok{market\_vol,}
    \AttributeTok{fitted =}\NormalTok{ calculated\_volatility}
\NormalTok{  )}
\NormalTok{\}}

\NormalTok{black\_fit }\OtherTok{\textless{}{-}}\NormalTok{ stats}\SpecialCharTok{::}\FunctionTok{optim}\NormalTok{(}
  \AttributeTok{par =}\NormalTok{ initial\_parameters\_black,}
  \AttributeTok{fn =}\NormalTok{ objective\_sabr\_black,}
  \AttributeTok{method =} \StringTok{"L{-}BFGS{-}B"}\NormalTok{,}
  \AttributeTok{lower =} \FunctionTok{c}\NormalTok{(}
    \FloatTok{0.0001}\NormalTok{,}
    \DecValTok{0}\NormalTok{,}
    \SpecialCharTok{{-}}\FloatTok{0.9999}
\NormalTok{  ),}
  \AttributeTok{upper =} \FunctionTok{c}\NormalTok{(}
    \ConstantTok{Inf}\NormalTok{,}
    \ConstantTok{Inf}\NormalTok{,}
    \FloatTok{0.9999}
\NormalTok{  )}
\NormalTok{)}

\NormalTok{black\_parameters }\OtherTok{\textless{}{-}}\NormalTok{ black\_fit}\SpecialCharTok{$}\NormalTok{par}

\NormalTok{black\_calculated\_volatility }\OtherTok{\textless{}{-}}\NormalTok{ purrr}\SpecialCharTok{::}\FunctionTok{map\_dbl}\NormalTok{(}
\NormalTok{  black\_market\_data}\SpecialCharTok{$}\NormalTok{strike,}
  \ControlFlowTok{function}\NormalTok{(strike\_value) \{}
\NormalTok{    GaussQuant}\SpecialCharTok{::}\FunctionTok{safe\_sabr\_vol\_GQL}\NormalTok{(}
      \AttributeTok{strike =}\NormalTok{ strike\_value,}
      \AttributeTok{forward =}\NormalTok{ forward\_black,}
      \AttributeTok{maturity =}\NormalTok{ maturity\_black,}
      \AttributeTok{alpha =}\NormalTok{ black\_parameters[[}\DecValTok{1}\NormalTok{]],}
      \AttributeTok{beta =}\NormalTok{ beta\_black,}
      \AttributeTok{volvol =}\NormalTok{ black\_parameters[[}\DecValTok{2}\NormalTok{]],}
      \AttributeTok{rho =}\NormalTok{ black\_parameters[[}\DecValTok{3}\NormalTok{]]}
\NormalTok{    )}
\NormalTok{  \}}
\NormalTok{)}

\NormalTok{black\_calibration\_tbl }\OtherTok{\textless{}{-}}\NormalTok{ tibble}\SpecialCharTok{::}\FunctionTok{tribble}\NormalTok{(}
  \SpecialCharTok{\textasciitilde{}}\NormalTok{parameter,          }\SpecialCharTok{\textasciitilde{}}\NormalTok{value,}
  \StringTok{"alpha"}\NormalTok{,               black\_parameters[[}\DecValTok{1}\NormalTok{]],}
  \StringTok{"beta"}\NormalTok{,                beta\_black,}
  \StringTok{"volvol"}\NormalTok{,              black\_parameters[[}\DecValTok{2}\NormalTok{]],}
  \StringTok{"rho"}\NormalTok{,                 black\_parameters[[}\DecValTok{3}\NormalTok{]],}
  \StringTok{"rmse"}\NormalTok{,                black\_fit}\SpecialCharTok{$}\NormalTok{value,}
  \StringTok{"convergence\_code"}\NormalTok{,    black\_fit}\SpecialCharTok{$}\NormalTok{convergence}
\NormalTok{)}

\NormalTok{GaussQuant}\SpecialCharTok{::}\FunctionTok{show\_tbl\_GQH}\NormalTok{(}
\NormalTok{  black\_calibration\_tbl,}
  \StringTok{"Black SABR calibration summary"}\NormalTok{,}
  \AttributeTok{n =} \DecValTok{20}
\NormalTok{)}
\CommentTok{\#\textgreater{} \# A tibble: 6 x 2}
\CommentTok{\#\textgreater{}   parameter         value}
\CommentTok{\#\textgreater{}   \textless{}chr\textgreater{}             \textless{}dbl\textgreater{}}
\CommentTok{\#\textgreater{} 1 alpha            0.0306}
\CommentTok{\#\textgreater{} 2 beta             0.5   }
\CommentTok{\#\textgreater{} 3 volvol           0.728 }
\CommentTok{\#\textgreater{} 4 rho              0.449 }
\CommentTok{\#\textgreater{} 5 rmse             0.0139}
\CommentTok{\#\textgreater{} 6 convergence\_code 0}
\end{Highlighting}
\end{Shaded}

\texttt{convergence\_code}が0であれば、\texttt{optim()}は通常の収束条件を満たして終了している。RMSEは市場ボラティリティとモデル値の平均的な誤差を表す。

\subsection{6.2.3
ボラティリティ・スマイル}\label{ux30dcux30e9ux30c6ux30a3ux30eaux30c6ux30a3ux30b9ux30deux30a4ux30eb}

\begin{Shaded}
\begin{Highlighting}[]
\NormalTok{black\_smile\_tbl }\OtherTok{\textless{}{-}}\NormalTok{ black\_market\_data }\SpecialCharTok{|\textgreater{}}
\NormalTok{  dplyr}\SpecialCharTok{::}\FunctionTok{transmute}\NormalTok{(}
    \AttributeTok{strike =}\NormalTok{ .data}\SpecialCharTok{$}\NormalTok{strike,}
    \AttributeTok{market\_vol =}\NormalTok{ .data}\SpecialCharTok{$}\NormalTok{market\_vol,}
    \AttributeTok{sabr\_vol =}\NormalTok{ black\_calculated\_volatility,}
    \AttributeTok{residual =}\NormalTok{ .data}\SpecialCharTok{$}\NormalTok{market\_vol }\SpecialCharTok{{-}}
\NormalTok{      black\_calculated\_volatility}
\NormalTok{  )}

\NormalTok{black\_smile\_plot\_tbl }\OtherTok{\textless{}{-}}\NormalTok{ black\_smile\_tbl }\SpecialCharTok{|\textgreater{}}
\NormalTok{  dplyr}\SpecialCharTok{::}\FunctionTok{select}\NormalTok{(}
\NormalTok{    .data}\SpecialCharTok{$}\NormalTok{strike,}
\NormalTok{    .data}\SpecialCharTok{$}\NormalTok{market\_vol,}
\NormalTok{    .data}\SpecialCharTok{$}\NormalTok{sabr\_vol}
\NormalTok{  ) }\SpecialCharTok{|\textgreater{}}
\NormalTok{  tidyr}\SpecialCharTok{::}\FunctionTok{pivot\_longer}\NormalTok{(}
    \AttributeTok{cols =} \FunctionTok{c}\NormalTok{(}
\NormalTok{      market\_vol,}
\NormalTok{      sabr\_vol}
\NormalTok{    ),}
    \AttributeTok{names\_to =} \StringTok{"series"}\NormalTok{,}
    \AttributeTok{values\_to =} \StringTok{"volatility"}
\NormalTok{  )}

\NormalTok{black\_smile\_plot }\OtherTok{\textless{}{-}}\NormalTok{ ggplot2}\SpecialCharTok{::}\FunctionTok{ggplot}\NormalTok{(}
\NormalTok{  black\_smile\_plot\_tbl,}
\NormalTok{  ggplot2}\SpecialCharTok{::}\FunctionTok{aes}\NormalTok{(}
    \AttributeTok{x =}\NormalTok{ .data}\SpecialCharTok{$}\NormalTok{strike,}
    \AttributeTok{y =}\NormalTok{ .data}\SpecialCharTok{$}\NormalTok{volatility,}
    \AttributeTok{group =}\NormalTok{ .data}\SpecialCharTok{$}\NormalTok{series,}
    \AttributeTok{linetype =}\NormalTok{ .data}\SpecialCharTok{$}\NormalTok{series,}
    \AttributeTok{shape =}\NormalTok{ .data}\SpecialCharTok{$}\NormalTok{series}
\NormalTok{  )}
\NormalTok{) }\SpecialCharTok{+}
\NormalTok{  ggplot2}\SpecialCharTok{::}\FunctionTok{geom\_line}\NormalTok{(}
    \AttributeTok{linewidth =} \FloatTok{0.6}
\NormalTok{  ) }\SpecialCharTok{+}
\NormalTok{  ggplot2}\SpecialCharTok{::}\FunctionTok{geom\_point}\NormalTok{(}
    \AttributeTok{size =} \FloatTok{2.2}
\NormalTok{  ) }\SpecialCharTok{+}
\NormalTok{  ggplot2}\SpecialCharTok{::}\FunctionTok{scale\_x\_continuous}\NormalTok{(}
    \AttributeTok{labels =} \ControlFlowTok{function}\NormalTok{(x) }\FunctionTok{sprintf}\NormalTok{(}\StringTok{"\%.2f\%\%"}\NormalTok{, }\DecValTok{100} \SpecialCharTok{*}\NormalTok{ x)}
\NormalTok{  ) }\SpecialCharTok{+}
\NormalTok{  ggplot2}\SpecialCharTok{::}\FunctionTok{scale\_y\_continuous}\NormalTok{(}
    \AttributeTok{labels =} \ControlFlowTok{function}\NormalTok{(x) }\FunctionTok{sprintf}\NormalTok{(}\StringTok{"\%.0f\%\%"}\NormalTok{, }\DecValTok{100} \SpecialCharTok{*}\NormalTok{ x)}
\NormalTok{  ) }\SpecialCharTok{+}
\NormalTok{  ggplot2}\SpecialCharTok{::}\FunctionTok{labs}\NormalTok{(}
    \AttributeTok{title =} \StringTok{"Black SABR calibration"}\NormalTok{,}
    \AttributeTok{subtitle =} \FunctionTok{paste0}\NormalTok{(}
      \StringTok{"alpha="}\NormalTok{,}
      \FunctionTok{sprintf}\NormalTok{(}\StringTok{"\%.5f"}\NormalTok{, black\_parameters[[}\DecValTok{1}\NormalTok{]]),}
      \StringTok{", beta="}\NormalTok{,}
      \FunctionTok{sprintf}\NormalTok{(}\StringTok{"\%.3f"}\NormalTok{, beta\_black),}
      \StringTok{", volvol="}\NormalTok{,}
      \FunctionTok{sprintf}\NormalTok{(}\StringTok{"\%.5f"}\NormalTok{, black\_parameters[[}\DecValTok{2}\NormalTok{]]),}
      \StringTok{", rho="}\NormalTok{,}
      \FunctionTok{sprintf}\NormalTok{(}\StringTok{"\%.5f"}\NormalTok{, black\_parameters[[}\DecValTok{3}\NormalTok{]])}
\NormalTok{    ),}
    \AttributeTok{x =} \StringTok{"Strike"}\NormalTok{,}
    \AttributeTok{y =} \StringTok{"Black volatility"}\NormalTok{,}
    \AttributeTok{linetype =} \StringTok{"Series"}\NormalTok{,}
    \AttributeTok{shape =} \StringTok{"Series"}
\NormalTok{  ) }\SpecialCharTok{+}
\NormalTok{  ggplot2}\SpecialCharTok{::}\FunctionTok{theme\_minimal}\NormalTok{()}

\NormalTok{GaussQuant}\SpecialCharTok{::}\FunctionTok{show\_tbl\_GQH}\NormalTok{(}
\NormalTok{  black\_smile\_tbl,}
  \StringTok{"Black SABR smile table"}\NormalTok{,}
  \AttributeTok{n =} \DecValTok{20}
\NormalTok{)}
\CommentTok{\#\textgreater{} \# A tibble: 7 x 4}
\CommentTok{\#\textgreater{}   strike market\_vol sabr\_vol  residual}
\CommentTok{\#\textgreater{}    \textless{}dbl\textgreater{}      \textless{}dbl\textgreater{}    \textless{}dbl\textgreater{}     \textless{}dbl\textgreater{}}
\CommentTok{\#\textgreater{} 1 0.0006      0.908    0.897  0.0108  }
\CommentTok{\#\textgreater{} 2 0.0031      0.461    0.485 {-}0.0240  }
\CommentTok{\#\textgreater{} 3 0.0056      0.453    0.449  0.00402 }
\CommentTok{\#\textgreater{} 4 0.0081      0.502    0.489  0.0125  }
\CommentTok{\#\textgreater{} 5 0.0106      0.538    0.528  0.0106  }
\CommentTok{\#\textgreater{} 6 0.0156      0.584    0.583  0.000924}
\CommentTok{\#\textgreater{} 7 0.0256      0.627    0.646 {-}0.0193}

\FunctionTok{print}\NormalTok{(black\_smile\_plot)}
\end{Highlighting}
\end{Shaded}

\pandocbounded{\includegraphics[keepaspectratio]{part2_quantive_riemann/ch06_sabr_model_files/figure-pdf/chapter06-black-smile-1.pdf}}

SABRはすべての市場点を完全に通過する補間法ではなく、少数のパラメータでスマイル全体の形状を表現するモデルである。残差が完全にゼロでなくても、Strike方向に滑らかなボラティリティを与えられる点に意味がある。

\section{6.3
SABRパラメータの感応度}\label{sabrux30d1ux30e9ux30e1ux30fcux30bfux306eux611fux5fdcux5ea6}

推定結果を基準として、一つのパラメータだけを上下させた場合のスマイルを確認する。パラメータ間には相互作用があるため、これは厳密なリスク分解ではなく、形状を理解するための比較である。

\begin{Shaded}
\begin{Highlighting}[]
\NormalTok{base\_parameter\_tbl }\OtherTok{\textless{}{-}}\NormalTok{ tibble}\SpecialCharTok{::}\FunctionTok{tribble}\NormalTok{(}
  \SpecialCharTok{\textasciitilde{}}\NormalTok{parameter, }\SpecialCharTok{\textasciitilde{}}\NormalTok{base\_value,                 }\SpecialCharTok{\textasciitilde{}}\NormalTok{shift,}
  \StringTok{"alpha"}\NormalTok{,     black\_parameters[[}\DecValTok{1}\NormalTok{]],        }\FloatTok{0.02}\NormalTok{,}
  \StringTok{"beta"}\NormalTok{,      beta\_black,                   }\FloatTok{0.20}\NormalTok{,}
  \StringTok{"volvol"}\NormalTok{,    black\_parameters[[}\DecValTok{2}\NormalTok{]],        }\FloatTok{0.40}\NormalTok{,}
  \StringTok{"rho"}\NormalTok{,       black\_parameters[[}\DecValTok{3}\NormalTok{]],        }\FloatTok{0.44}
\NormalTok{)}

\NormalTok{base\_parameters }\OtherTok{\textless{}{-}}\NormalTok{ stats}\SpecialCharTok{::}\FunctionTok{setNames}\NormalTok{(}
\NormalTok{  base\_parameter\_tbl}\SpecialCharTok{$}\NormalTok{base\_value,}
\NormalTok{  base\_parameter\_tbl}\SpecialCharTok{$}\NormalTok{parameter}
\NormalTok{)}

\NormalTok{parameter\_shifts }\OtherTok{\textless{}{-}}\NormalTok{ stats}\SpecialCharTok{::}\FunctionTok{setNames}\NormalTok{(}
\NormalTok{  base\_parameter\_tbl}\SpecialCharTok{$}\NormalTok{shift,}
\NormalTok{  base\_parameter\_tbl}\SpecialCharTok{$}\NormalTok{parameter}
\NormalTok{)}

\NormalTok{sabr\_sensitivity\_tbl }\OtherTok{\textless{}{-}}\NormalTok{ purrr}\SpecialCharTok{::}\FunctionTok{map\_dfr}\NormalTok{(}
  \FunctionTok{names}\NormalTok{(base\_parameters),}
  \ControlFlowTok{function}\NormalTok{(parameter\_name) \{}
\NormalTok{    parameters\_up }\OtherTok{\textless{}{-}}\NormalTok{ base\_parameters}
\NormalTok{    parameters\_down }\OtherTok{\textless{}{-}}\NormalTok{ base\_parameters}

\NormalTok{    parameters\_up[[parameter\_name]] }\OtherTok{\textless{}{-}}
\NormalTok{      parameters\_up[[parameter\_name]] }\SpecialCharTok{+}
\NormalTok{      parameter\_shifts[[parameter\_name]]}

\NormalTok{    parameters\_down[[parameter\_name]] }\OtherTok{\textless{}{-}}
\NormalTok{      parameters\_down[[parameter\_name]] }\SpecialCharTok{{-}}
\NormalTok{      parameter\_shifts[[parameter\_name]]}

\NormalTok{    parameters\_up[[}\StringTok{"alpha"}\NormalTok{]] }\OtherTok{\textless{}{-}} \FunctionTok{max}\NormalTok{(}
\NormalTok{      parameters\_up[[}\StringTok{"alpha"}\NormalTok{]],}
      \FloatTok{1e{-}6}
\NormalTok{    )}
\NormalTok{    parameters\_down[[}\StringTok{"alpha"}\NormalTok{]] }\OtherTok{\textless{}{-}} \FunctionTok{max}\NormalTok{(}
\NormalTok{      parameters\_down[[}\StringTok{"alpha"}\NormalTok{]],}
      \FloatTok{1e{-}6}
\NormalTok{    )}
\NormalTok{    parameters\_up[[}\StringTok{"beta"}\NormalTok{]] }\OtherTok{\textless{}{-}} \FunctionTok{min}\NormalTok{(}
      \FunctionTok{max}\NormalTok{(parameters\_up[[}\StringTok{"beta"}\NormalTok{]], }\DecValTok{0}\NormalTok{),}
      \DecValTok{1}
\NormalTok{    )}
\NormalTok{    parameters\_down[[}\StringTok{"beta"}\NormalTok{]] }\OtherTok{\textless{}{-}} \FunctionTok{min}\NormalTok{(}
      \FunctionTok{max}\NormalTok{(parameters\_down[[}\StringTok{"beta"}\NormalTok{]], }\DecValTok{0}\NormalTok{),}
      \DecValTok{1}
\NormalTok{    )}
\NormalTok{    parameters\_up[[}\StringTok{"volvol"}\NormalTok{]] }\OtherTok{\textless{}{-}} \FunctionTok{max}\NormalTok{(}
\NormalTok{      parameters\_up[[}\StringTok{"volvol"}\NormalTok{]],}
      \DecValTok{0}
\NormalTok{    )}
\NormalTok{    parameters\_down[[}\StringTok{"volvol"}\NormalTok{]] }\OtherTok{\textless{}{-}} \FunctionTok{max}\NormalTok{(}
\NormalTok{      parameters\_down[[}\StringTok{"volvol"}\NormalTok{]],}
      \DecValTok{0}
\NormalTok{    )}
\NormalTok{    parameters\_up[[}\StringTok{"rho"}\NormalTok{]] }\OtherTok{\textless{}{-}} \FunctionTok{min}\NormalTok{(}
      \FunctionTok{max}\NormalTok{(parameters\_up[[}\StringTok{"rho"}\NormalTok{]], }\SpecialCharTok{{-}}\FloatTok{0.95}\NormalTok{),}
      \FloatTok{0.95}
\NormalTok{    )}
\NormalTok{    parameters\_down[[}\StringTok{"rho"}\NormalTok{]] }\OtherTok{\textless{}{-}} \FunctionTok{min}\NormalTok{(}
      \FunctionTok{max}\NormalTok{(parameters\_down[[}\StringTok{"rho"}\NormalTok{]], }\SpecialCharTok{{-}}\FloatTok{0.95}\NormalTok{),}
      \FloatTok{0.95}
\NormalTok{    )}

\NormalTok{    scenario\_parameter\_tbl }\OtherTok{\textless{}{-}}\NormalTok{ tibble}\SpecialCharTok{::}\FunctionTok{tribble}\NormalTok{(}
      \SpecialCharTok{\textasciitilde{}}\NormalTok{scenario, }\SpecialCharTok{\textasciitilde{}}\NormalTok{alpha,                       }\SpecialCharTok{\textasciitilde{}}\NormalTok{beta,                       }\SpecialCharTok{\textasciitilde{}}\NormalTok{volvol,                       }\SpecialCharTok{\textasciitilde{}}\NormalTok{rho,}
      \StringTok{"base"}\NormalTok{,     base\_parameters[[}\StringTok{"alpha"}\NormalTok{]],   base\_parameters[[}\StringTok{"beta"}\NormalTok{]],  base\_parameters[[}\StringTok{"volvol"}\NormalTok{]], base\_parameters[[}\StringTok{"rho"}\NormalTok{]],}
      \StringTok{"up"}\NormalTok{,       parameters\_up[[}\StringTok{"alpha"}\NormalTok{]],     parameters\_up[[}\StringTok{"beta"}\NormalTok{]],    parameters\_up[[}\StringTok{"volvol"}\NormalTok{]],   parameters\_up[[}\StringTok{"rho"}\NormalTok{]],}
      \StringTok{"down"}\NormalTok{,     parameters\_down[[}\StringTok{"alpha"}\NormalTok{]],   parameters\_down[[}\StringTok{"beta"}\NormalTok{]],  parameters\_down[[}\StringTok{"volvol"}\NormalTok{]], parameters\_down[[}\StringTok{"rho"}\NormalTok{]]}
\NormalTok{    )}

\NormalTok{    purrr}\SpecialCharTok{::}\FunctionTok{pmap\_dfr}\NormalTok{(}
\NormalTok{      scenario\_parameter\_tbl,}
      \ControlFlowTok{function}\NormalTok{(}
\NormalTok{        scenario,}
\NormalTok{        alpha,}
\NormalTok{        beta,}
\NormalTok{        volvol,}
\NormalTok{        rho}
\NormalTok{      ) \{}
\NormalTok{        tibble}\SpecialCharTok{::}\FunctionTok{tibble}\NormalTok{(}
          \AttributeTok{strike =}\NormalTok{ black\_market\_data}\SpecialCharTok{$}\NormalTok{strike,}
          \AttributeTok{volatility =}\NormalTok{ purrr}\SpecialCharTok{::}\FunctionTok{map\_dbl}\NormalTok{(}
\NormalTok{            black\_market\_data}\SpecialCharTok{$}\NormalTok{strike,}
            \ControlFlowTok{function}\NormalTok{(strike\_value) \{}
\NormalTok{              GaussQuant}\SpecialCharTok{::}\FunctionTok{safe\_sabr\_vol\_GQL}\NormalTok{(}
                \AttributeTok{strike =}\NormalTok{ strike\_value,}
                \AttributeTok{forward =}\NormalTok{ forward\_black,}
                \AttributeTok{maturity =}\NormalTok{ maturity\_black,}
                \AttributeTok{alpha =}\NormalTok{ alpha,}
                \AttributeTok{beta =}\NormalTok{ beta,}
                \AttributeTok{volvol =}\NormalTok{ volvol,}
                \AttributeTok{rho =}\NormalTok{ rho}
\NormalTok{              )}
\NormalTok{            \}}
\NormalTok{          ),}
          \AttributeTok{scenario =}\NormalTok{ scenario,}
          \AttributeTok{parameter =}\NormalTok{ parameter\_name}
\NormalTok{        )}
\NormalTok{      \}}
\NormalTok{    )}
\NormalTok{  \}}
\NormalTok{)}

\NormalTok{sabr\_sensitivity\_plot }\OtherTok{\textless{}{-}}\NormalTok{ ggplot2}\SpecialCharTok{::}\FunctionTok{ggplot}\NormalTok{(}
\NormalTok{  sabr\_sensitivity\_tbl,}
\NormalTok{  ggplot2}\SpecialCharTok{::}\FunctionTok{aes}\NormalTok{(}
    \AttributeTok{x =}\NormalTok{ .data}\SpecialCharTok{$}\NormalTok{strike,}
    \AttributeTok{y =}\NormalTok{ .data}\SpecialCharTok{$}\NormalTok{volatility,}
    \AttributeTok{group =}\NormalTok{ .data}\SpecialCharTok{$}\NormalTok{scenario,}
    \AttributeTok{linetype =}\NormalTok{ .data}\SpecialCharTok{$}\NormalTok{scenario}
\NormalTok{  )}
\NormalTok{) }\SpecialCharTok{+}
\NormalTok{  ggplot2}\SpecialCharTok{::}\FunctionTok{geom\_line}\NormalTok{(}
    \AttributeTok{linewidth =} \FloatTok{0.6}
\NormalTok{  ) }\SpecialCharTok{+}
\NormalTok{  ggplot2}\SpecialCharTok{::}\FunctionTok{facet\_wrap}\NormalTok{(}
    \SpecialCharTok{\textasciitilde{}}\NormalTok{parameter,}
    \AttributeTok{scales =} \StringTok{"free\_y"}
\NormalTok{  ) }\SpecialCharTok{+}
\NormalTok{  ggplot2}\SpecialCharTok{::}\FunctionTok{scale\_x\_continuous}\NormalTok{(}
    \AttributeTok{labels =} \ControlFlowTok{function}\NormalTok{(x) }\FunctionTok{sprintf}\NormalTok{(}\StringTok{"\%.2f\%\%"}\NormalTok{, }\DecValTok{100} \SpecialCharTok{*}\NormalTok{ x)}
\NormalTok{  ) }\SpecialCharTok{+}
\NormalTok{  ggplot2}\SpecialCharTok{::}\FunctionTok{scale\_y\_continuous}\NormalTok{(}
    \AttributeTok{labels =} \ControlFlowTok{function}\NormalTok{(x) }\FunctionTok{sprintf}\NormalTok{(}\StringTok{"\%.0f\%\%"}\NormalTok{, }\DecValTok{100} \SpecialCharTok{*}\NormalTok{ x)}
\NormalTok{  ) }\SpecialCharTok{+}
\NormalTok{  ggplot2}\SpecialCharTok{::}\FunctionTok{labs}\NormalTok{(}
    \AttributeTok{title =} \StringTok{"SABR parameter sensitivity"}\NormalTok{,}
    \AttributeTok{x =} \StringTok{"Strike"}\NormalTok{,}
    \AttributeTok{y =} \StringTok{"Black volatility"}\NormalTok{,}
    \AttributeTok{linetype =} \StringTok{"Scenario"}
\NormalTok{  ) }\SpecialCharTok{+}
\NormalTok{  ggplot2}\SpecialCharTok{::}\FunctionTok{theme\_minimal}\NormalTok{()}

\NormalTok{GaussQuant}\SpecialCharTok{::}\FunctionTok{show\_tbl\_GQH}\NormalTok{(}
\NormalTok{  base\_parameter\_tbl,}
  \StringTok{"SABR sensitivity settings"}\NormalTok{,}
  \AttributeTok{n =} \DecValTok{20}
\NormalTok{)}
\CommentTok{\#\textgreater{} \# A tibble: 4 x 3}
\CommentTok{\#\textgreater{}   parameter base\_value shift}
\CommentTok{\#\textgreater{}   \textless{}chr\textgreater{}          \textless{}dbl\textgreater{} \textless{}dbl\textgreater{}}
\CommentTok{\#\textgreater{} 1 alpha         0.0306  0.02}
\CommentTok{\#\textgreater{} 2 beta          0.5     0.2 }
\CommentTok{\#\textgreater{} 3 volvol        0.728   0.4 }
\CommentTok{\#\textgreater{} 4 rho           0.449   0.44}

\FunctionTok{print}\NormalTok{(sabr\_sensitivity\_plot)}
\end{Highlighting}
\end{Shaded}

\pandocbounded{\includegraphics[keepaspectratio]{part2_quantive_riemann/ch06_sabr_model_files/figure-pdf/chapter06-parameter-sensitivity-1.pdf}}

一般に、\(\alpha\)はスマイルの水準、\(\nu\)は曲率、\(\rho\)は傾きと左右非対称性に強く影響する。\(\beta\)も水準と傾きの双方へ影響するため、他のパラメータとの識別が難しい。このため、\(\beta\)を市場慣行や過去推定値に基づいて固定することが多い。

\section{6.4 Shifted
SABRと負の金利}\label{shifted-sabrux3068ux8ca0ux306eux91d1ux5229}

通常のBlack
SABR近似では、フォワードとStrikeが正である必要がある。負の金利を含む市場では、フォワードとStrikeへ同じShift
\(s\) を加える。

\[
\widetilde{F}
=
F+s
\]

\[
\widetilde{K}
=
K+s
\]

Shift後の値が正であれば、対数やべき乗を含むSABR近似を適用できる。

本章では、Normal
volatilityの市場データに対して、Shiftを加えたHagan近似を使用する。出力はNormal
volatilityであるが、内部ではShift後のフォワードとStrikeを使ってスマイルを表現する。

Shiftは負の金利を消してしまう調整ではない。モデル計算上の座標を移動するだけであり、実際のStrikeとフォワードは元の水準のまま解釈する。

\section{6.5 Shifted SABR
Normalのカリブレーション}\label{shifted-sabr-normalux306eux30abux30eaux30d6ux30ecux30fcux30b7ux30e7ux30f3}

\subsection{6.5.1 市場データ}\label{ux5e02ux5834ux30c7ux30fcux30bf-1}

フォワードを0.05\%、満期を2年、\(\beta=0.5\)、Shiftを2.5\%とする。市場データにはマイナス0.95\%のStrikeが含まれる。

\begin{Shaded}
\begin{Highlighting}[]
\NormalTok{normal\_market\_data }\OtherTok{\textless{}{-}}\NormalTok{ tibble}\SpecialCharTok{::}\FunctionTok{tribble}\NormalTok{(}
  \SpecialCharTok{\textasciitilde{}}\NormalTok{strike\_pct, }\SpecialCharTok{\textasciitilde{}}\NormalTok{market\_vol\_bp,}
  \SpecialCharTok{{-}}\FloatTok{0.95}\NormalTok{,        }\FloatTok{41.94}\NormalTok{,}
  \SpecialCharTok{{-}}\FloatTok{0.45}\NormalTok{,        }\FloatTok{38.64}\NormalTok{,}
  \SpecialCharTok{{-}}\FloatTok{0.20}\NormalTok{,        }\FloatTok{37.71}\NormalTok{,}
   \FloatTok{0.05}\NormalTok{,        }\FloatTok{37.80}\NormalTok{,}
   \FloatTok{0.30}\NormalTok{,        }\FloatTok{39.05}\NormalTok{,}
   \FloatTok{0.55}\NormalTok{,        }\FloatTok{41.37}\NormalTok{,}
   \FloatTok{1.05}\NormalTok{,        }\FloatTok{47.81}
\NormalTok{) }\SpecialCharTok{|\textgreater{}}
\NormalTok{  dplyr}\SpecialCharTok{::}\FunctionTok{mutate}\NormalTok{(}
    \AttributeTok{strike =}\NormalTok{ .data}\SpecialCharTok{$}\NormalTok{strike\_pct }\SpecialCharTok{/} \DecValTok{100}\NormalTok{,}
    \AttributeTok{market\_vol =}\NormalTok{ .data}\SpecialCharTok{$}\NormalTok{market\_vol\_bp }\SpecialCharTok{/} \DecValTok{10000}
\NormalTok{  )}

\NormalTok{forward\_normal }\OtherTok{\textless{}{-}} \FloatTok{0.0005}
\NormalTok{maturity\_normal }\OtherTok{\textless{}{-}} \DecValTok{2}
\NormalTok{beta\_normal }\OtherTok{\textless{}{-}} \FloatTok{0.5}
\NormalTok{shift\_normal }\OtherTok{\textless{}{-}} \FloatTok{0.025}

\NormalTok{normal\_model\_inputs }\OtherTok{\textless{}{-}}\NormalTok{ tibble}\SpecialCharTok{::}\FunctionTok{tribble}\NormalTok{(}
  \SpecialCharTok{\textasciitilde{}}\NormalTok{input,       }\SpecialCharTok{\textasciitilde{}}\NormalTok{value,}
  \StringTok{"forward"}\NormalTok{,     forward\_normal,}
  \StringTok{"maturity"}\NormalTok{,    maturity\_normal,}
  \StringTok{"fixed\_beta"}\NormalTok{,  beta\_normal,}
  \StringTok{"shift"}\NormalTok{,       shift\_normal}
\NormalTok{)}

\NormalTok{initial\_parameters\_normal }\OtherTok{\textless{}{-}} \FunctionTok{c}\NormalTok{(}
  \AttributeTok{alpha =} \FloatTok{0.1}\NormalTok{,}
  \AttributeTok{volvol =} \FloatTok{0.1}\NormalTok{,}
  \AttributeTok{rho =} \FloatTok{0.1}
\NormalTok{)}

\NormalTok{GaussQuant}\SpecialCharTok{::}\FunctionTok{show\_tbl\_GQH}\NormalTok{(}
\NormalTok{  normal\_market\_data }\SpecialCharTok{|\textgreater{}}
\NormalTok{    dplyr}\SpecialCharTok{::}\FunctionTok{select}\NormalTok{(}
\NormalTok{      .data}\SpecialCharTok{$}\NormalTok{strike,}
\NormalTok{      .data}\SpecialCharTok{$}\NormalTok{market\_vol}
\NormalTok{    ),}
  \StringTok{"Shifted SABR normal market data"}\NormalTok{,}
  \AttributeTok{n =} \DecValTok{20}
\NormalTok{)}
\CommentTok{\#\textgreater{} \# A tibble: 7 x 2}
\CommentTok{\#\textgreater{}    strike market\_vol}
\CommentTok{\#\textgreater{}     \textless{}dbl\textgreater{}      \textless{}dbl\textgreater{}}
\CommentTok{\#\textgreater{} 1 {-}0.0095    0.00419}
\CommentTok{\#\textgreater{} 2 {-}0.0045    0.00386}
\CommentTok{\#\textgreater{} 3 {-}0.002     0.00377}
\CommentTok{\#\textgreater{} 4  0.0005    0.00378}
\CommentTok{\#\textgreater{} 5  0.003     0.00390}
\CommentTok{\#\textgreater{} 6  0.0055    0.00414}
\CommentTok{\#\textgreater{} 7  0.0105    0.00478}

\NormalTok{GaussQuant}\SpecialCharTok{::}\FunctionTok{show\_tbl\_GQH}\NormalTok{(}
\NormalTok{  normal\_model\_inputs,}
  \StringTok{"Shifted SABR normal model inputs"}\NormalTok{,}
  \AttributeTok{n =} \DecValTok{20}
\NormalTok{)}
\CommentTok{\#\textgreater{} \# A tibble: 4 x 2}
\CommentTok{\#\textgreater{}   input       value}
\CommentTok{\#\textgreater{}   \textless{}chr\textgreater{}       \textless{}dbl\textgreater{}}
\CommentTok{\#\textgreater{} 1 forward    0.0005}
\CommentTok{\#\textgreater{} 2 maturity   2     }
\CommentTok{\#\textgreater{} 3 fixed\_beta 0.5   }
\CommentTok{\#\textgreater{} 4 shift      0.025}
\end{Highlighting}
\end{Shaded}

すべての市場Strikeについて、

\[
K+s>0
\]

が必要である。本例の最小Strikeはマイナス0.95\%、Shiftは2.5\%であるため、Shift後のStrikeは正となる。

\subsection{6.5.2
パラメータ・カリブレーション}\label{ux30d1ux30e9ux30e1ux30fcux30bfux30abux30eaux30d6ux30ecux30fcux30b7ux30e7ux30f3-1}

Black
SABRと同様に\(\beta\)を固定し、\(\alpha\)、\(\nu\)、\(\rho\)を推定する。目的関数はNormal
volatility単位のRMSEである。

\begin{Shaded}
\begin{Highlighting}[]
\NormalTok{objective\_sabr\_normal }\OtherTok{\textless{}{-}} \ControlFlowTok{function}\NormalTok{(parameters) \{}
\NormalTok{  alpha }\OtherTok{\textless{}{-}}\NormalTok{ parameters[[}\DecValTok{1}\NormalTok{]]}
\NormalTok{  volvol }\OtherTok{\textless{}{-}}\NormalTok{ parameters[[}\DecValTok{2}\NormalTok{]]}
\NormalTok{  rho }\OtherTok{\textless{}{-}}\NormalTok{ parameters[[}\DecValTok{3}\NormalTok{]]}

  \ControlFlowTok{if}\NormalTok{ (}
\NormalTok{    alpha }\SpecialCharTok{\textless{}=} \FloatTok{0.001} \SpecialCharTok{||}
\NormalTok{    volvol }\SpecialCharTok{\textless{}} \DecValTok{0} \SpecialCharTok{||}
\NormalTok{    rho }\SpecialCharTok{\textless{}=} \SpecialCharTok{{-}}\FloatTok{0.999} \SpecialCharTok{||}
\NormalTok{    rho }\SpecialCharTok{\textgreater{}=} \FloatTok{0.999}
\NormalTok{  ) \{}
    \FunctionTok{return}\NormalTok{(}\FloatTok{1e6}\NormalTok{)}
\NormalTok{  \}}

\NormalTok{  calculated\_volatility }\OtherTok{\textless{}{-}}\NormalTok{ purrr}\SpecialCharTok{::}\FunctionTok{map\_dbl}\NormalTok{(}
\NormalTok{    normal\_market\_data}\SpecialCharTok{$}\NormalTok{strike,}
    \ControlFlowTok{function}\NormalTok{(strike\_value) \{}
\NormalTok{      GaussQuant}\SpecialCharTok{::}\FunctionTok{shifted\_normal\_vol\_hagan\_GQL}\NormalTok{(}
        \AttributeTok{strike =}\NormalTok{ strike\_value,}
        \AttributeTok{forward =}\NormalTok{ forward\_normal,}
        \AttributeTok{maturity =}\NormalTok{ maturity\_normal,}
        \AttributeTok{beta =}\NormalTok{ beta\_normal,}
        \AttributeTok{alpha =}\NormalTok{ alpha,}
        \AttributeTok{volvol =}\NormalTok{ volvol,}
        \AttributeTok{rho =}\NormalTok{ rho,}
        \AttributeTok{shift =}\NormalTok{ shift\_normal}
\NormalTok{      )}
\NormalTok{    \}}
\NormalTok{  )}

  \ControlFlowTok{if}\NormalTok{ (}
    \FunctionTok{any}\NormalTok{(}\FunctionTok{is.na}\NormalTok{(calculated\_volatility)) }\SpecialCharTok{||}
    \FunctionTok{any}\NormalTok{(}\SpecialCharTok{!}\FunctionTok{is.finite}\NormalTok{(calculated\_volatility))}
\NormalTok{  ) \{}
    \FunctionTok{return}\NormalTok{(}\FloatTok{1e6}\NormalTok{)}
\NormalTok{  \}}

\NormalTok{  GaussQuant}\SpecialCharTok{::}\FunctionTok{rmse\_GQH}\NormalTok{(}
    \AttributeTok{actual =}\NormalTok{ normal\_market\_data}\SpecialCharTok{$}\NormalTok{market\_vol,}
    \AttributeTok{fitted =}\NormalTok{ calculated\_volatility}
\NormalTok{  )}
\NormalTok{\}}

\NormalTok{normal\_fit }\OtherTok{\textless{}{-}}\NormalTok{ stats}\SpecialCharTok{::}\FunctionTok{optim}\NormalTok{(}
  \AttributeTok{par =}\NormalTok{ initial\_parameters\_normal,}
  \AttributeTok{fn =}\NormalTok{ objective\_sabr\_normal,}
  \AttributeTok{method =} \StringTok{"L{-}BFGS{-}B"}\NormalTok{,}
  \AttributeTok{lower =} \FunctionTok{c}\NormalTok{(}
    \FloatTok{0.001}\NormalTok{,}
    \DecValTok{0}\NormalTok{,}
    \SpecialCharTok{{-}}\FloatTok{0.999}
\NormalTok{  ),}
  \AttributeTok{upper =} \FunctionTok{c}\NormalTok{(}
    \ConstantTok{Inf}\NormalTok{,}
    \ConstantTok{Inf}\NormalTok{,}
    \FloatTok{0.999}
\NormalTok{  )}
\NormalTok{)}

\NormalTok{normal\_parameters }\OtherTok{\textless{}{-}}\NormalTok{ normal\_fit}\SpecialCharTok{$}\NormalTok{par}

\NormalTok{normal\_calculated\_volatility }\OtherTok{\textless{}{-}}\NormalTok{ purrr}\SpecialCharTok{::}\FunctionTok{map\_dbl}\NormalTok{(}
\NormalTok{  normal\_market\_data}\SpecialCharTok{$}\NormalTok{strike,}
  \ControlFlowTok{function}\NormalTok{(strike\_value) \{}
\NormalTok{    GaussQuant}\SpecialCharTok{::}\FunctionTok{shifted\_normal\_vol\_hagan\_GQL}\NormalTok{(}
      \AttributeTok{strike =}\NormalTok{ strike\_value,}
      \AttributeTok{forward =}\NormalTok{ forward\_normal,}
      \AttributeTok{maturity =}\NormalTok{ maturity\_normal,}
      \AttributeTok{beta =}\NormalTok{ beta\_normal,}
      \AttributeTok{alpha =}\NormalTok{ normal\_parameters[[}\DecValTok{1}\NormalTok{]],}
      \AttributeTok{volvol =}\NormalTok{ normal\_parameters[[}\DecValTok{2}\NormalTok{]],}
      \AttributeTok{rho =}\NormalTok{ normal\_parameters[[}\DecValTok{3}\NormalTok{]],}
      \AttributeTok{shift =}\NormalTok{ shift\_normal}
\NormalTok{    )}
\NormalTok{  \}}
\NormalTok{)}

\NormalTok{normal\_calibration\_tbl }\OtherTok{\textless{}{-}}\NormalTok{ tibble}\SpecialCharTok{::}\FunctionTok{tribble}\NormalTok{(}
  \SpecialCharTok{\textasciitilde{}}\NormalTok{parameter,          }\SpecialCharTok{\textasciitilde{}}\NormalTok{value,}
  \StringTok{"alpha"}\NormalTok{,               normal\_parameters[[}\DecValTok{1}\NormalTok{]],}
  \StringTok{"beta"}\NormalTok{,                beta\_normal,}
  \StringTok{"volvol"}\NormalTok{,              normal\_parameters[[}\DecValTok{2}\NormalTok{]],}
  \StringTok{"rho"}\NormalTok{,                 normal\_parameters[[}\DecValTok{3}\NormalTok{]],}
  \StringTok{"shift"}\NormalTok{,               shift\_normal,}
  \StringTok{"rmse"}\NormalTok{,                normal\_fit}\SpecialCharTok{$}\NormalTok{value,}
  \StringTok{"convergence\_code"}\NormalTok{,    normal\_fit}\SpecialCharTok{$}\NormalTok{convergence}
\NormalTok{)}

\NormalTok{GaussQuant}\SpecialCharTok{::}\FunctionTok{show\_tbl\_GQH}\NormalTok{(}
\NormalTok{  normal\_calibration\_tbl,}
  \StringTok{"Shifted SABR normal calibration summary"}\NormalTok{,}
  \AttributeTok{n =} \DecValTok{20}
\NormalTok{)}
\CommentTok{\#\textgreater{} \# A tibble: 7 x 2}
\CommentTok{\#\textgreater{}   parameter           value}
\CommentTok{\#\textgreater{}   \textless{}chr\textgreater{}               \textless{}dbl\textgreater{}}
\CommentTok{\#\textgreater{} 1 alpha            0.0254  }
\CommentTok{\#\textgreater{} 2 beta             0.5     }
\CommentTok{\#\textgreater{} 3 volvol           0.0994  }
\CommentTok{\#\textgreater{} 4 rho              0.0998  }
\CommentTok{\#\textgreater{} 5 shift            0.025   }
\CommentTok{\#\textgreater{} 6 rmse             0.000298}
\CommentTok{\#\textgreater{} 7 convergence\_code 0}
\end{Highlighting}
\end{Shaded}

\subsection{6.5.3
Normalボラティリティ・スマイル}\label{normalux30dcux30e9ux30c6ux30a3ux30eaux30c6ux30a3ux30b9ux30deux30a4ux30eb}

\begin{Shaded}
\begin{Highlighting}[]
\NormalTok{normal\_smile\_tbl }\OtherTok{\textless{}{-}}\NormalTok{ normal\_market\_data }\SpecialCharTok{|\textgreater{}}
\NormalTok{  dplyr}\SpecialCharTok{::}\FunctionTok{transmute}\NormalTok{(}
    \AttributeTok{strike =}\NormalTok{ .data}\SpecialCharTok{$}\NormalTok{strike,}
    \AttributeTok{market\_vol =}\NormalTok{ .data}\SpecialCharTok{$}\NormalTok{market\_vol,}
    \AttributeTok{sabr\_vol =}\NormalTok{ normal\_calculated\_volatility,}
    \AttributeTok{residual =}\NormalTok{ .data}\SpecialCharTok{$}\NormalTok{market\_vol }\SpecialCharTok{{-}}
\NormalTok{      normal\_calculated\_volatility}
\NormalTok{  )}

\NormalTok{normal\_smile\_plot\_tbl }\OtherTok{\textless{}{-}}\NormalTok{ normal\_smile\_tbl }\SpecialCharTok{|\textgreater{}}
\NormalTok{  dplyr}\SpecialCharTok{::}\FunctionTok{select}\NormalTok{(}
\NormalTok{    .data}\SpecialCharTok{$}\NormalTok{strike,}
\NormalTok{    .data}\SpecialCharTok{$}\NormalTok{market\_vol,}
\NormalTok{    .data}\SpecialCharTok{$}\NormalTok{sabr\_vol}
\NormalTok{  ) }\SpecialCharTok{|\textgreater{}}
\NormalTok{  tidyr}\SpecialCharTok{::}\FunctionTok{pivot\_longer}\NormalTok{(}
    \AttributeTok{cols =} \FunctionTok{c}\NormalTok{(}
\NormalTok{      market\_vol,}
\NormalTok{      sabr\_vol}
\NormalTok{    ),}
    \AttributeTok{names\_to =} \StringTok{"series"}\NormalTok{,}
    \AttributeTok{values\_to =} \StringTok{"volatility"}
\NormalTok{  )}

\NormalTok{normal\_smile\_plot }\OtherTok{\textless{}{-}}\NormalTok{ ggplot2}\SpecialCharTok{::}\FunctionTok{ggplot}\NormalTok{(}
\NormalTok{  normal\_smile\_plot\_tbl,}
\NormalTok{  ggplot2}\SpecialCharTok{::}\FunctionTok{aes}\NormalTok{(}
    \AttributeTok{x =}\NormalTok{ .data}\SpecialCharTok{$}\NormalTok{strike,}
    \AttributeTok{y =}\NormalTok{ .data}\SpecialCharTok{$}\NormalTok{volatility,}
    \AttributeTok{group =}\NormalTok{ .data}\SpecialCharTok{$}\NormalTok{series,}
    \AttributeTok{linetype =}\NormalTok{ .data}\SpecialCharTok{$}\NormalTok{series,}
    \AttributeTok{shape =}\NormalTok{ .data}\SpecialCharTok{$}\NormalTok{series}
\NormalTok{  )}
\NormalTok{) }\SpecialCharTok{+}
\NormalTok{  ggplot2}\SpecialCharTok{::}\FunctionTok{geom\_line}\NormalTok{(}
    \AttributeTok{linewidth =} \FloatTok{0.6}
\NormalTok{  ) }\SpecialCharTok{+}
\NormalTok{  ggplot2}\SpecialCharTok{::}\FunctionTok{geom\_point}\NormalTok{(}
    \AttributeTok{size =} \FloatTok{2.2}
\NormalTok{  ) }\SpecialCharTok{+}
\NormalTok{  ggplot2}\SpecialCharTok{::}\FunctionTok{scale\_x\_continuous}\NormalTok{(}
    \AttributeTok{labels =} \ControlFlowTok{function}\NormalTok{(x) }\FunctionTok{sprintf}\NormalTok{(}\StringTok{"\%.2f\%\%"}\NormalTok{, }\DecValTok{100} \SpecialCharTok{*}\NormalTok{ x)}
\NormalTok{  ) }\SpecialCharTok{+}
\NormalTok{  ggplot2}\SpecialCharTok{::}\FunctionTok{scale\_y\_continuous}\NormalTok{(}
    \AttributeTok{labels =} \ControlFlowTok{function}\NormalTok{(x) }\FunctionTok{sprintf}\NormalTok{(}\StringTok{"\%.1f bp"}\NormalTok{, }\DecValTok{10000} \SpecialCharTok{*}\NormalTok{ x)}
\NormalTok{  ) }\SpecialCharTok{+}
\NormalTok{  ggplot2}\SpecialCharTok{::}\FunctionTok{labs}\NormalTok{(}
    \AttributeTok{title =} \StringTok{"Shifted SABR normal calibration"}\NormalTok{,}
    \AttributeTok{subtitle =} \FunctionTok{paste0}\NormalTok{(}
      \StringTok{"alpha="}\NormalTok{,}
      \FunctionTok{sprintf}\NormalTok{(}\StringTok{"\%.5f"}\NormalTok{, normal\_parameters[[}\DecValTok{1}\NormalTok{]]),}
      \StringTok{", beta="}\NormalTok{,}
      \FunctionTok{sprintf}\NormalTok{(}\StringTok{"\%.3f"}\NormalTok{, beta\_normal),}
      \StringTok{", volvol="}\NormalTok{,}
      \FunctionTok{sprintf}\NormalTok{(}\StringTok{"\%.5f"}\NormalTok{, normal\_parameters[[}\DecValTok{2}\NormalTok{]]),}
      \StringTok{", rho="}\NormalTok{,}
      \FunctionTok{sprintf}\NormalTok{(}\StringTok{"\%.5f"}\NormalTok{, normal\_parameters[[}\DecValTok{3}\NormalTok{]]),}
      \StringTok{", shift="}\NormalTok{,}
      \FunctionTok{sprintf}\NormalTok{(}\StringTok{"\%.2f\%\%"}\NormalTok{, }\DecValTok{100} \SpecialCharTok{*}\NormalTok{ shift\_normal)}
\NormalTok{    ),}
    \AttributeTok{x =} \StringTok{"Strike"}\NormalTok{,}
    \AttributeTok{y =} \StringTok{"Normal volatility"}\NormalTok{,}
    \AttributeTok{linetype =} \StringTok{"Series"}\NormalTok{,}
    \AttributeTok{shape =} \StringTok{"Series"}
\NormalTok{  ) }\SpecialCharTok{+}
\NormalTok{  ggplot2}\SpecialCharTok{::}\FunctionTok{theme\_minimal}\NormalTok{()}

\NormalTok{GaussQuant}\SpecialCharTok{::}\FunctionTok{show\_tbl\_GQH}\NormalTok{(}
\NormalTok{  normal\_smile\_tbl }\SpecialCharTok{|\textgreater{}}
\NormalTok{    dplyr}\SpecialCharTok{::}\FunctionTok{mutate}\NormalTok{(}
      \AttributeTok{market\_vol\_bp =} \DecValTok{10000} \SpecialCharTok{*}\NormalTok{ .data}\SpecialCharTok{$}\NormalTok{market\_vol,}
      \AttributeTok{sabr\_vol\_bp =} \DecValTok{10000} \SpecialCharTok{*}\NormalTok{ .data}\SpecialCharTok{$}\NormalTok{sabr\_vol,}
      \AttributeTok{residual\_bp =} \DecValTok{10000} \SpecialCharTok{*}\NormalTok{ .data}\SpecialCharTok{$}\NormalTok{residual}
\NormalTok{    ),}
  \StringTok{"Shifted SABR normal smile table"}\NormalTok{,}
  \AttributeTok{n =} \DecValTok{20}
\NormalTok{)}
\CommentTok{\#\textgreater{} \# A tibble: 7 x 7}
\CommentTok{\#\textgreater{}    strike market\_vol sabr\_vol   residual market\_vol\_bp sabr\_vol\_bp residual\_bp}
\CommentTok{\#\textgreater{}     \textless{}dbl\textgreater{}      \textless{}dbl\textgreater{}    \textless{}dbl\textgreater{}      \textless{}dbl\textgreater{}         \textless{}dbl\textgreater{}       \textless{}dbl\textgreater{}       \textless{}dbl\textgreater{}}
\CommentTok{\#\textgreater{} 1 {-}0.0095    0.00419  0.00360  0.000593           41.9        36.0       5.93 }
\CommentTok{\#\textgreater{} 2 {-}0.0045    0.00386  0.00383  0.0000348          38.6        38.3       0.348}
\CommentTok{\#\textgreater{} 3 {-}0.002     0.00377  0.00394 {-}0.000170           37.7        39.4      {-}1.70 }
\CommentTok{\#\textgreater{} 4  0.0005    0.00378  0.00405 {-}0.000273           37.8        40.5      {-}2.73 }
\CommentTok{\#\textgreater{} 5  0.003     0.00390  0.00417 {-}0.000261           39.0        41.7      {-}2.61 }
\CommentTok{\#\textgreater{} 6  0.0055    0.00414  0.00428 {-}0.000141           41.4        42.8      {-}1.41 }
\CommentTok{\#\textgreater{} 7  0.0105    0.00478  0.00450  0.000277           47.8        45.0       2.77}

\FunctionTok{print}\NormalTok{(normal\_smile\_plot)}
\end{Highlighting}
\end{Shaded}

\pandocbounded{\includegraphics[keepaspectratio]{part2_quantive_riemann/ch06_sabr_model_files/figure-pdf/chapter06-normal-smile-1.pdf}}

Black volatilityとNormal volatilityは単位が異なるため、Black
SABRのRMSEとShifted SABR
NormalのRMSEを数値の大きさだけで直接比較してはならない。Black側は相対ボラティリティ、Normal側は金利の絶対変化幅である。

\section{6.6
外挿とモデル比較}\label{ux5916ux633fux3068ux30e2ux30c7ux30ebux6bd4ux8f03}

\subsection{6.6.1
カリブレーション範囲外のStrike}\label{ux30abux30eaux30d6ux30ecux30fcux30b7ux30e7ux30f3ux7bc4ux56f2ux5916ux306estrike}

SABRモデルは、市場点の間を滑らかにつなぐだけでなく、市場点の外側にもボラティリティを与える。ただし、外挿結果は観測データによる制約が弱く、パラメータの影響を強く受ける。

次のコードでは、市場データの外側にある二つのStrikeについてNormal
volatilityを確認する。

\begin{Shaded}
\begin{Highlighting}[]
\NormalTok{outer\_strike\_input\_tbl }\OtherTok{\textless{}{-}}\NormalTok{ tibble}\SpecialCharTok{::}\FunctionTok{tribble}\NormalTok{(}
  \SpecialCharTok{\textasciitilde{}}\NormalTok{label,       }\SpecialCharTok{\textasciitilde{}}\NormalTok{strike,}
  \StringTok{"lower\_wing"}\NormalTok{, }\SpecialCharTok{{-}}\FloatTok{0.0195}\NormalTok{,}
  \StringTok{"upper\_wing"}\NormalTok{,  }\FloatTok{0.0205}
\NormalTok{)}

\NormalTok{outer\_strike\_tbl }\OtherTok{\textless{}{-}}\NormalTok{ outer\_strike\_input\_tbl }\SpecialCharTok{|\textgreater{}}
\NormalTok{  dplyr}\SpecialCharTok{::}\FunctionTok{mutate}\NormalTok{(}
    \AttributeTok{shifted\_strike =}\NormalTok{ .data}\SpecialCharTok{$}\NormalTok{strike }\SpecialCharTok{+}
\NormalTok{      shift\_normal,}
    \AttributeTok{normal\_volatility =}\NormalTok{ purrr}\SpecialCharTok{::}\FunctionTok{map\_dbl}\NormalTok{(}
\NormalTok{      .data}\SpecialCharTok{$}\NormalTok{strike,}
      \ControlFlowTok{function}\NormalTok{(strike\_value) \{}
\NormalTok{        GaussQuant}\SpecialCharTok{::}\FunctionTok{shifted\_normal\_vol\_hagan\_GQL}\NormalTok{(}
          \AttributeTok{strike =}\NormalTok{ strike\_value,}
          \AttributeTok{forward =}\NormalTok{ forward\_normal,}
          \AttributeTok{maturity =}\NormalTok{ maturity\_normal,}
          \AttributeTok{beta =}\NormalTok{ beta\_normal,}
          \AttributeTok{alpha =}\NormalTok{ normal\_parameters[[}\DecValTok{1}\NormalTok{]],}
          \AttributeTok{volvol =}\NormalTok{ normal\_parameters[[}\DecValTok{2}\NormalTok{]],}
          \AttributeTok{rho =}\NormalTok{ normal\_parameters[[}\DecValTok{3}\NormalTok{]],}
          \AttributeTok{shift =}\NormalTok{ shift\_normal}
\NormalTok{        )}
\NormalTok{      \}}
\NormalTok{    ),}
    \AttributeTok{normal\_volatility\_bp =} \DecValTok{10000} \SpecialCharTok{*}
\NormalTok{      .data}\SpecialCharTok{$}\NormalTok{normal\_volatility}
\NormalTok{  )}

\NormalTok{GaussQuant}\SpecialCharTok{::}\FunctionTok{show\_tbl\_GQH}\NormalTok{(}
\NormalTok{  outer\_strike\_tbl,}
  \StringTok{"Shifted SABR outer{-}strike check"}\NormalTok{,}
  \AttributeTok{n =} \DecValTok{20}
\NormalTok{)}
\CommentTok{\#\textgreater{} \# A tibble: 2 x 5}
\CommentTok{\#\textgreater{}   label       strike shifted\_strike normal\_volatility normal\_volatility\_bp}
\CommentTok{\#\textgreater{}   \textless{}chr\textgreater{}        \textless{}dbl\textgreater{}          \textless{}dbl\textgreater{}             \textless{}dbl\textgreater{}                \textless{}dbl\textgreater{}}
\CommentTok{\#\textgreater{} 1 lower\_wing {-}0.0195         0.0055           0.00307                 30.7}
\CommentTok{\#\textgreater{} 2 upper\_wing  0.0205         0.0455           0.00496                 49.6}
\end{Highlighting}
\end{Shaded}

Shift後のStrikeがゼロ以下となる領域では、この近似式を使用できない。また、数式上計算できる場合でも、カリブレーション範囲から遠い外挿値には注意が必要である。

\subsection{6.6.2 Black SABRとShifted SABR
Normal}\label{black-sabrux3068shifted-sabr-normal}

二つのカリブレーション結果を一覧にする。

\begin{Shaded}
\begin{Highlighting}[]
\NormalTok{calibration\_comparison\_tbl }\OtherTok{\textless{}{-}}\NormalTok{ tibble}\SpecialCharTok{::}\FunctionTok{tribble}\NormalTok{(}
  \SpecialCharTok{\textasciitilde{}}\NormalTok{model,                }\SpecialCharTok{\textasciitilde{}}\NormalTok{alpha,                  }\SpecialCharTok{\textasciitilde{}}\NormalTok{beta,       }\SpecialCharTok{\textasciitilde{}}\NormalTok{volvol,                  }\SpecialCharTok{\textasciitilde{}}\NormalTok{rho,                   }\SpecialCharTok{\textasciitilde{}}\NormalTok{shift,        }\SpecialCharTok{\textasciitilde{}}\NormalTok{rmse,             }\SpecialCharTok{\textasciitilde{}}\NormalTok{convergence\_code,}
  \StringTok{"Black SABR"}\NormalTok{,           black\_parameters[[}\DecValTok{1}\NormalTok{]],   beta\_black,  black\_parameters[[}\DecValTok{2}\NormalTok{]],   black\_parameters[[}\DecValTok{3}\NormalTok{]],  }\DecValTok{0}\NormalTok{,             black\_fit}\SpecialCharTok{$}\NormalTok{value,   black\_fit}\SpecialCharTok{$}\NormalTok{convergence,}
  \StringTok{"Shifted SABR Normal"}\NormalTok{,  normal\_parameters[[}\DecValTok{1}\NormalTok{]],  beta\_normal, normal\_parameters[[}\DecValTok{2}\NormalTok{]],  normal\_parameters[[}\DecValTok{3}\NormalTok{]], shift\_normal,  normal\_fit}\SpecialCharTok{$}\NormalTok{value,  normal\_fit}\SpecialCharTok{$}\NormalTok{convergence}
\NormalTok{)}

\NormalTok{GaussQuant}\SpecialCharTok{::}\FunctionTok{show\_tbl\_GQH}\NormalTok{(}
\NormalTok{  calibration\_comparison\_tbl,}
  \StringTok{"SABR calibration comparison"}\NormalTok{,}
  \AttributeTok{n =} \DecValTok{20}
\NormalTok{)}
\CommentTok{\#\textgreater{} \# A tibble: 2 x 8}
\CommentTok{\#\textgreater{}   model                alpha  beta volvol    rho shift     rmse convergence\_code}
\CommentTok{\#\textgreater{}   \textless{}chr\textgreater{}                \textless{}dbl\textgreater{} \textless{}dbl\textgreater{}  \textless{}dbl\textgreater{}  \textless{}dbl\textgreater{} \textless{}dbl\textgreater{}    \textless{}dbl\textgreater{}            \textless{}int\textgreater{}}
\CommentTok{\#\textgreater{} 1 Black SABR          0.0306   0.5 0.728  0.449  0     0.0139                  0}
\CommentTok{\#\textgreater{} 2 Shifted SABR Normal 0.0254   0.5 0.0994 0.0998 0.025 0.000298                0}
\end{Highlighting}
\end{Shaded}

Black SABRでは、正のフォワードとStrikeからBlack
volatilityを計算した。Shifted SABR
Normalでは、フォワードとStrikeへShiftを加え、負のStrikeを含む市場データへNormal
volatilityを適合させた。

どちらも同じ4つのSABRパラメータを用いるが、出力するインプライド・ボラティリティの定義と単位が異なる。したがって、パラメータ値やRMSEを異なるボラティリティ表現の間で単純比較するのではなく、それぞれの市場表示におけるスマイル適合度を確認する。

\section{6.7 まとめ}\label{ux307eux3068ux3081-10}

SABRモデルは、フォワードとそのボラティリティを同時に確率変数として扱い、Strikeごとに異なるインプライド・ボラティリティを少数のパラメータで表現する。

\(\alpha\)は主にボラティリティ水準、\(\nu\)はスマイルの曲率、\(\rho\)は傾きと左右非対称性へ影響する。\(\beta\)はフォワード水準に対する変動幅の弾性を表し、他のパラメータとの識別が難しいため、本章では固定してカリブレーションした。

Black SABRでは、正のフォワードとStrikeを前提としてBlack
volatilityのスマイルを評価した。負の金利を含むNormal
volatility市場では、フォワードとStrikeへ共通のShiftを加えるShifted
SABRを使用した。

SABRは市場ボラティリティを完全に補間することだけを目的とするものではない。Strike方向に滑らかで、パラメータによる解釈が可能なスマイルを構築し、観測されていないStrikeの評価やリスク管理へ利用する。

\begin{Shaded}
\begin{Highlighting}[]
\NormalTok{chapter06\_summary\_tbl }\OtherTok{\textless{}{-}}\NormalTok{ tibble}\SpecialCharTok{::}\FunctionTok{tribble}\NormalTok{(}
  \SpecialCharTok{\textasciitilde{}}\NormalTok{section,               }\SpecialCharTok{\textasciitilde{}}\NormalTok{main\_result,}
  \StringTok{"Black SABR"}\NormalTok{,            }\StringTok{"Calibration to a Black{-}volatility smile"}\NormalTok{,}
  \StringTok{"Parameter sensitivity"}\NormalTok{, }\StringTok{"Alpha, beta, vol{-}of{-}vol, and rho effects"}\NormalTok{,}
  \StringTok{"Shifted SABR Normal"}\NormalTok{,   }\StringTok{"Normal{-}volatility smile with negative strikes"}\NormalTok{,}
  \StringTok{"Outer strikes"}\NormalTok{,         }\StringTok{"Extrapolation outside calibration points"}
\NormalTok{)}

\NormalTok{GaussQuant}\SpecialCharTok{::}\FunctionTok{show\_tbl\_GQH}\NormalTok{(}
\NormalTok{  chapter06\_summary\_tbl,}
  \StringTok{"Chapter II{-}6 summary"}\NormalTok{,}
  \AttributeTok{n =} \DecValTok{20}
\NormalTok{)}
\CommentTok{\#\textgreater{} \# A tibble: 4 x 2}
\CommentTok{\#\textgreater{}   section               main\_result                                  }
\CommentTok{\#\textgreater{}   \textless{}chr\textgreater{}                 \textless{}chr\textgreater{}                                        }
\CommentTok{\#\textgreater{} 1 Black SABR            Calibration to a Black{-}volatility smile      }
\CommentTok{\#\textgreater{} 2 Parameter sensitivity Alpha, beta, vol{-}of{-}vol, and rho effects     }
\CommentTok{\#\textgreater{} 3 Shifted SABR Normal   Normal{-}volatility smile with negative strikes}
\CommentTok{\#\textgreater{} 4 Outer strikes         Extrapolation outside calibration points}

\FunctionTok{cat}\NormalTok{(}
  \StringTok{"}\SpecialCharTok{\textbackslash{}n}\StringTok{Chapter II{-}6 SABR model completed successfully.}\SpecialCharTok{\textbackslash{}n}\StringTok{"}
\NormalTok{)}
\CommentTok{\#\textgreater{} }
\CommentTok{\#\textgreater{} Chapter II{-}6 SABR model completed successfully.}
\end{Highlighting}
\end{Shaded}

\chapter{第II-7章 金利期間構造モデル Hull--White
Callable債評価}\label{ux7b2cii-7ux7ae0-ux91d1ux5229ux671fux9593ux69cbux9020ux30e2ux30c7ux30eb-hullwhite-callableux50b5ux8a55ux4fa1}

Callable
Bondは、発行体があらかじめ定められた日と価格で期限前償還できる債券である。市場金利が低下すると、発行体は既存の高いクーポンを支払い続けるよりも、債券を期限前償還して低い金利で資金を再調達する誘因を持つ。

したがって、コール条項は発行体にとって価値を持つ一方、投資家にとっては不利な権利である。同じクーポン、満期および信用条件を持つ場合、Callable
Bondの価格はオプションを含まないBullet Bondの価格以下となる。

概念的には、次の関係で表される。

\[
V_{\mathrm{callable}}
=
V_{\mathrm{bullet}}
-
V_{\mathrm{issuer\ call}}
\]

発行体のコール権は、金利低下時に価値が上昇する。この性質は、固定金利を受け取り変動金利を支払うReceiver
Swaptionに近い。複数のコール可能日がある場合、その権利はBermudan
Receiver Swaptionに対応する。

ただし、Callable
BondとSwaptionでは、元本交換、クーポンと固定レグ、決済、経過利子、行使価格および支払日の調整が必ずしも完全には一致しない。このため、本章では、

\[
V_{\mathrm{callable}}
\approx
V_{\mathrm{bullet}}
-
V_{\mathrm{Bermudan\ receiver\ swaption}}
\]

をモデル構造の比較として確認する。

期限前償還権を評価するためには、将来の金利変動と各行使日における継続価値を扱う必要がある。本章では、初期金利期間構造へ適合できるHull--White
1-factorモデルと金利Treeを使用する。

\begin{Shaded}
\begin{Highlighting}[]
\NormalTok{knitr}\SpecialCharTok{::}\NormalTok{opts\_chunk}\SpecialCharTok{$}\FunctionTok{set}\NormalTok{(}
  \AttributeTok{collapse =} \ConstantTok{TRUE}\NormalTok{,}
  \AttributeTok{comment =} \StringTok{"\#\textgreater{}"}\NormalTok{,}
  \AttributeTok{warning =} \ConstantTok{FALSE}\NormalTok{,}
  \AttributeTok{message =} \ConstantTok{FALSE}
\NormalTok{)}


\FunctionTok{suppressPackageStartupMessages}\NormalTok{(\{}
  \FunctionTok{library}\NormalTok{(tidyverse)}
  \FunctionTok{library}\NormalTok{(ggthemes)}
  \FunctionTok{library}\NormalTok{(QuantLib)}
  \FunctionTok{library}\NormalTok{(GaussQuant)}
  \FunctionTok{library}\NormalTok{(tidyverse)}
  \FunctionTok{library}\NormalTok{(broom)}
\NormalTok{\})}
\end{Highlighting}
\end{Shaded}

\section{7.1
Hull--Whiteモデルと債券条件}\label{hullwhiteux30e2ux30c7ux30ebux3068ux50b5ux5238ux6761ux4ef6}

\subsection{7.1.1 Hull--White
1-factorモデル}\label{hullwhite-1-factorux30e2ux30c7ux30eb}

Hull--Whiteモデルでは、短期金利 \(r_t\) を次のGaussian過程で表す。

\[
dr_t
=
\left[
\theta(t)-ar_t
\right]dt
+
\sigma\,dW_t
\]

ここで、\(a\) は平均回帰速度、\(\sigma\)
は短期金利のボラティリティ、\(\theta(t)\) は時間依存のドリフトである。

\(a\) が大きいほど、金利ショックの影響は速く減衰する。\(\sigma\)
が大きいほど将来金利の分布が広がり、金利オプションの価値は一般に高くなる。

Hull--Whiteモデルの重要な特徴は、\(\theta(t)\)を調整することによって、評価時点で観測された初期イールドカーブを再現できることである。モデルの確率部分を状態変数
\(x_t\) として分離すると、次のOrnstein--Uhlenbeck過程として表せる。

\[
dx_t
=
-ax_t\,dt
+
\sigma\,dW_t
\]

短期金利は、初期期間構造を表す決定論的な部分と状態変数の和として解釈できる。

\[
r_t
=
\phi(t)+x_t
\]

Hull--WhiteモデルはGaussianモデルであるため、短期金利が負になる可能性を持つ。前章までのNormalモデルと同様、ゼロ金利近傍を扱いやすい一方、極端な負の金利も理論上許容する。

\subsection{7.1.2 債券条件}\label{ux50b5ux5238ux6761ux4ef6}

5年満期、年率3\%、半年払いの固定利付債を考える。市場利回りも3\%とし、オプションを含まないBullet
Bondの価格がおおむね額面付近となる条件を設定する。

最初の1年間をノンコール期間とし、その後の各クーポン日に額面100で期限前償還できるものとする。満期日は通常の元本償還日であるため、コール可能日から除外する。

\begin{Shaded}
\begin{Highlighting}[]
\NormalTok{trade\_date }\OtherTok{\textless{}{-}}\NormalTok{ GaussQuant}\SpecialCharTok{::}\FunctionTok{date\_GQL}\NormalTok{(}\StringTok{"2024{-}04{-}15"}\NormalTok{)}
\NormalTok{issue\_date }\OtherTok{\textless{}{-}}\NormalTok{ GaussQuant}\SpecialCharTok{::}\FunctionTok{date\_GQL}\NormalTok{(}\StringTok{"2024{-}04{-}17"}\NormalTok{)}
\NormalTok{maturity\_date }\OtherTok{\textless{}{-}}\NormalTok{ GaussQuant}\SpecialCharTok{::}\FunctionTok{date\_GQL}\NormalTok{(}\StringTok{"2029{-}04{-}17"}\NormalTok{)}

\NormalTok{GaussQuant}\SpecialCharTok{::}\FunctionTok{set\_eval\_date\_GQL}\NormalTok{(trade\_date)}

\NormalTok{coupon\_rate }\OtherTok{\textless{}{-}} \FloatTok{0.03}
\NormalTok{market\_yield }\OtherTok{\textless{}{-}} \FloatTok{0.03}
\NormalTok{settlement\_days }\OtherTok{\textless{}{-}} \DecValTok{2}\NormalTok{L}
\NormalTok{face\_amount }\OtherTok{\textless{}{-}} \DecValTok{100}
\NormalTok{call\_price\_clean }\OtherTok{\textless{}{-}} \DecValTok{100}
\NormalTok{non\_call\_periods }\OtherTok{\textless{}{-}} \DecValTok{2}\NormalTok{L}

\NormalTok{bond\_calendar }\OtherTok{\textless{}{-}}\NormalTok{ QuantLib}\SpecialCharTok{::}\FunctionTok{WeekendsOnly}\NormalTok{()}
\NormalTok{bond\_day\_counter }\OtherTok{\textless{}{-}}\NormalTok{ QuantLib}\SpecialCharTok{::}\FunctionTok{ActualActual}\NormalTok{(}\StringTok{"Bond"}\NormalTok{)}
\NormalTok{model\_day\_counter }\OtherTok{\textless{}{-}}\NormalTok{ QuantLib}\SpecialCharTok{::}\FunctionTok{Actual365Fixed}\NormalTok{()}
\NormalTok{bond\_compounding }\OtherTok{\textless{}{-}}\NormalTok{ QuantLib}\SpecialCharTok{::}\FunctionTok{Compounding\_Compounded\_get}\NormalTok{()}
\NormalTok{bond\_frequency }\OtherTok{\textless{}{-}}\NormalTok{ QuantLib}\SpecialCharTok{::}\FunctionTok{Frequency\_Semiannual\_get}\NormalTok{()}

\NormalTok{bond\_input\_tbl }\OtherTok{\textless{}{-}}\NormalTok{ tibble}\SpecialCharTok{::}\FunctionTok{tribble}\NormalTok{(}
  \SpecialCharTok{\textasciitilde{}}\NormalTok{input,             }\SpecialCharTok{\textasciitilde{}}\NormalTok{value,}
  \StringTok{"trade\_date"}\NormalTok{,        }\StringTok{"2024{-}04{-}15"}\NormalTok{,}
  \StringTok{"issue\_date"}\NormalTok{,        }\StringTok{"2024{-}04{-}17"}\NormalTok{,}
  \StringTok{"maturity\_date"}\NormalTok{,     }\StringTok{"2029{-}04{-}17"}\NormalTok{,}
  \StringTok{"coupon\_rate"}\NormalTok{,        }\FunctionTok{as.character}\NormalTok{(coupon\_rate),}
  \StringTok{"market\_yield"}\NormalTok{,       }\FunctionTok{as.character}\NormalTok{(market\_yield),}
  \StringTok{"settlement\_days"}\NormalTok{,    }\FunctionTok{as.character}\NormalTok{(settlement\_days),}
  \StringTok{"face\_amount"}\NormalTok{,        }\FunctionTok{as.character}\NormalTok{(face\_amount),}
  \StringTok{"call\_price"}\NormalTok{,         }\FunctionTok{as.character}\NormalTok{(call\_price\_clean),}
  \StringTok{"non\_call\_periods"}\NormalTok{,   }\FunctionTok{as.character}\NormalTok{(non\_call\_periods)}
\NormalTok{)}

\NormalTok{GaussQuant}\SpecialCharTok{::}\FunctionTok{show\_tbl\_GQH}\NormalTok{(}
\NormalTok{  bond\_input\_tbl,}
  \StringTok{"Callable bond input assumptions"}\NormalTok{,}
  \AttributeTok{n =} \DecValTok{20}
\NormalTok{)}
\CommentTok{\#\textgreater{} \# A tibble: 9 x 2}
\CommentTok{\#\textgreater{}   input            value     }
\CommentTok{\#\textgreater{}   \textless{}chr\textgreater{}            \textless{}chr\textgreater{}     }
\CommentTok{\#\textgreater{} 1 trade\_date       2024{-}04{-}15}
\CommentTok{\#\textgreater{} 2 issue\_date       2024{-}04{-}17}
\CommentTok{\#\textgreater{} 3 maturity\_date    2029{-}04{-}17}
\CommentTok{\#\textgreater{} 4 coupon\_rate      0.03      }
\CommentTok{\#\textgreater{} 5 market\_yield     0.03      }
\CommentTok{\#\textgreater{} 6 settlement\_days  2         }
\CommentTok{\#\textgreater{} 7 face\_amount      100       }
\CommentTok{\#\textgreater{} 8 call\_price       100       }
\CommentTok{\#\textgreater{} 9 non\_call\_periods 2}
\end{Highlighting}
\end{Shaded}

\subsection{7.1.3 Bullet Bond}\label{bullet-bond}

クーポンスケジュールを作成し、コール条項を含まない固定利付債を評価する。

\begin{Shaded}
\begin{Highlighting}[]
\NormalTok{bond\_schedule }\OtherTok{\textless{}{-}}\NormalTok{ QuantLib}\SpecialCharTok{::}\FunctionTok{Schedule}\NormalTok{(}
\NormalTok{  issue\_date,}
\NormalTok{  maturity\_date,}
\NormalTok{  GaussQuant}\SpecialCharTok{::}\FunctionTok{period\_months\_GQL}\NormalTok{(}\DecValTok{6}\NormalTok{),}
\NormalTok{  bond\_calendar,}
  \StringTok{"Unadjusted"}\NormalTok{,}
  \StringTok{"Unadjusted"}\NormalTok{,}
  \StringTok{"Backward"}\NormalTok{,}
  \ConstantTok{FALSE}
\NormalTok{)}

\NormalTok{bond\_schedule\_tbl }\OtherTok{\textless{}{-}}\NormalTok{ GaussQuant}\SpecialCharTok{::}\FunctionTok{schedule\_table\_GQL}\NormalTok{(}
\NormalTok{  bond\_schedule}
\NormalTok{)}

\NormalTok{bullet\_bond }\OtherTok{\textless{}{-}}\NormalTok{ QuantLib}\SpecialCharTok{::}\FunctionTok{FixedRateBond}\NormalTok{(}
\NormalTok{  settlement\_days,}
\NormalTok{  face\_amount,}
\NormalTok{  bond\_schedule,}
  \FunctionTok{c}\NormalTok{(coupon\_rate),}
\NormalTok{  bond\_day\_counter}
\NormalTok{)}

\NormalTok{bullet\_settlement\_date }\OtherTok{\textless{}{-}} \FunctionTok{tryCatch}\NormalTok{(}
\NormalTok{  bullet\_bond}\SpecialCharTok{$}\FunctionTok{settlementDate}\NormalTok{(),}
  \AttributeTok{error =} \ControlFlowTok{function}\NormalTok{(e) }\ConstantTok{NULL}
\NormalTok{)}

\NormalTok{bullet\_clean\_price }\OtherTok{\textless{}{-}}\NormalTok{ GaussQuant}\SpecialCharTok{::}\FunctionTok{bond\_clean\_price\_from\_yield\_safe\_GQL}\NormalTok{(}
  \AttributeTok{bond =}\NormalTok{ bullet\_bond,}
  \AttributeTok{yield\_rate =}\NormalTok{ market\_yield,}
  \AttributeTok{day\_counter =}\NormalTok{ bond\_day\_counter,}
  \AttributeTok{compounding =}\NormalTok{ bond\_compounding,}
  \AttributeTok{frequency =}\NormalTok{ bond\_frequency,}
  \AttributeTok{settlement\_date =}\NormalTok{ bullet\_settlement\_date}
\NormalTok{)}

\NormalTok{bullet\_summary\_tbl }\OtherTok{\textless{}{-}}\NormalTok{ tibble}\SpecialCharTok{::}\FunctionTok{tribble}\NormalTok{(}
  \SpecialCharTok{\textasciitilde{}}\NormalTok{metric,                        }\SpecialCharTok{\textasciitilde{}}\NormalTok{value,}
  \StringTok{"coupon\_rate"}\NormalTok{,                   coupon\_rate,}
  \StringTok{"market\_yield"}\NormalTok{,                  market\_yield,}
  \StringTok{"clean\_price\_from\_yield"}\NormalTok{,        bullet\_clean\_price,}
  \StringTok{"face\_amount"}\NormalTok{,                   face\_amount}
\NormalTok{)}

\NormalTok{GaussQuant}\SpecialCharTok{::}\FunctionTok{show\_tbl\_GQH}\NormalTok{(}
\NormalTok{  bond\_schedule\_tbl,}
  \StringTok{"Bullet bond schedule"}\NormalTok{,}
  \AttributeTok{n =} \DecValTok{20}
\NormalTok{)}
\CommentTok{\#\textgreater{} \# A tibble: 11 x 2}
\CommentTok{\#\textgreater{}    period schedule\_date}
\CommentTok{\#\textgreater{}     \textless{}int\textgreater{} \textless{}date\textgreater{}       }
\CommentTok{\#\textgreater{}  1      1 2024{-}04{-}17   }
\CommentTok{\#\textgreater{}  2      2 2024{-}10{-}17   }
\CommentTok{\#\textgreater{}  3      3 2025{-}04{-}17   }
\CommentTok{\#\textgreater{}  4      4 2025{-}10{-}17   }
\CommentTok{\#\textgreater{}  5      5 2026{-}04{-}17   }
\CommentTok{\#\textgreater{}  6      6 2026{-}10{-}17   }
\CommentTok{\#\textgreater{}  7      7 2027{-}04{-}17   }
\CommentTok{\#\textgreater{}  8      8 2027{-}10{-}17   }
\CommentTok{\#\textgreater{}  9      9 2028{-}04{-}17   }
\CommentTok{\#\textgreater{} 10     10 2028{-}10{-}17   }
\CommentTok{\#\textgreater{} 11     11 2029{-}04{-}17}

\NormalTok{GaussQuant}\SpecialCharTok{::}\FunctionTok{show\_tbl\_GQH}\NormalTok{(}
\NormalTok{  bullet\_summary\_tbl,}
  \StringTok{"Bullet bond summary"}\NormalTok{,}
  \AttributeTok{n =} \DecValTok{20}
\NormalTok{)}
\CommentTok{\#\textgreater{} \# A tibble: 4 x 2}
\CommentTok{\#\textgreater{}   metric                  value}
\CommentTok{\#\textgreater{}   \textless{}chr\textgreater{}                   \textless{}dbl\textgreater{}}
\CommentTok{\#\textgreater{} 1 coupon\_rate              0.03}
\CommentTok{\#\textgreater{} 2 market\_yield             0.03}
\CommentTok{\#\textgreater{} 3 clean\_price\_from\_yield 100.0 }
\CommentTok{\#\textgreater{} 4 face\_amount            100}
\end{Highlighting}
\end{Shaded}

\subsection{7.1.4
初期カーブとHull--Whiteモデル}\label{ux521dux671fux30abux30fcux30d6ux3068hullwhiteux30e2ux30c7ux30eb}

Callable
BondとSwaptionを同じ初期カーブとHull--Whiteパラメータで評価する。

本節では、年率3\%のフラットカーブを使用する。Callable
Bondの主目的はオプション構造とTree評価を確認することであり、複雑な市場カーブ構築は行わない。

\begin{Shaded}
\begin{Highlighting}[]
\NormalTok{hull\_white\_a }\OtherTok{\textless{}{-}} \FloatTok{0.03}
\NormalTok{hull\_white\_sigma }\OtherTok{\textless{}{-}} \FloatTok{0.01}
\NormalTok{tree\_steps }\OtherTok{\textless{}{-}} \DecValTok{120}\NormalTok{L}

\NormalTok{flat\_curve }\OtherTok{\textless{}{-}}\NormalTok{ GaussQuant}\SpecialCharTok{::}\FunctionTok{flat\_curve\_GQL}\NormalTok{(}
  \AttributeTok{reference\_date =}\NormalTok{ trade\_date,}
  \AttributeTok{rate =}\NormalTok{ market\_yield,}
  \AttributeTok{day\_counter =}\NormalTok{ model\_day\_counter,}
  \AttributeTok{compounding =} \StringTok{"Continuous"}
\NormalTok{)}

\NormalTok{flat\_curve\_handle }\OtherTok{\textless{}{-}}\NormalTok{ QuantLib}\SpecialCharTok{::}\FunctionTok{YieldTermStructureHandle}\NormalTok{(}
\NormalTok{  flat\_curve}
\NormalTok{)}

\NormalTok{hull\_white\_model }\OtherTok{\textless{}{-}}\NormalTok{ GaussQuant}\SpecialCharTok{::}\FunctionTok{hull\_white\_model\_GQL}\NormalTok{(}
  \AttributeTok{term\_structure =}\NormalTok{ flat\_curve\_handle,}
  \AttributeTok{a =}\NormalTok{ hull\_white\_a,}
  \AttributeTok{sigma =}\NormalTok{ hull\_white\_sigma}
\NormalTok{)}

\NormalTok{hull\_white\_input\_tbl }\OtherTok{\textless{}{-}}\NormalTok{ tibble}\SpecialCharTok{::}\FunctionTok{tribble}\NormalTok{(}
  \SpecialCharTok{\textasciitilde{}}\NormalTok{parameter,          }\SpecialCharTok{\textasciitilde{}}\NormalTok{value,}
  \StringTok{"mean\_reversion\_a"}\NormalTok{,   hull\_white\_a,}
  \StringTok{"sigma"}\NormalTok{,              hull\_white\_sigma,}
  \StringTok{"tree\_steps"}\NormalTok{,         tree\_steps,}
  \StringTok{"flat\_curve\_rate"}\NormalTok{,    market\_yield}
\NormalTok{)}

\NormalTok{GaussQuant}\SpecialCharTok{::}\FunctionTok{show\_tbl\_GQH}\NormalTok{(}
\NormalTok{  hull\_white\_input\_tbl,}
  \StringTok{"Hull{-}White model assumptions"}\NormalTok{,}
  \AttributeTok{n =} \DecValTok{20}
\NormalTok{)}
\CommentTok{\#\textgreater{} \# A tibble: 4 x 2}
\CommentTok{\#\textgreater{}   parameter         value}
\CommentTok{\#\textgreater{}   \textless{}chr\textgreater{}             \textless{}dbl\textgreater{}}
\CommentTok{\#\textgreater{} 1 mean\_reversion\_a   0.03}
\CommentTok{\#\textgreater{} 2 sigma              0.01}
\CommentTok{\#\textgreater{} 3 tree\_steps       120   }
\CommentTok{\#\textgreater{} 4 flat\_curve\_rate    0.03}
\end{Highlighting}
\end{Shaded}

\section{7.2 Callable
Bondの構築とTree評価}\label{callable-bondux306eux69cbux7bc9ux3068treeux8a55ux4fa1}

\subsection{7.2.1
コールスケジュール}\label{ux30b3ux30fcux30ebux30b9ux30b1ux30b8ux30e5ux30fcux30eb}

スケジュールには発行日と満期日も含まれる。最初の2クーポン期間をノンコールとし、満期日を除く残りのクーポン日をコール可能日とする。

\begin{Shaded}
\begin{Highlighting}[]
\NormalTok{bond\_schedule\_dates }\OtherTok{\textless{}{-}}\NormalTok{ GaussQuant}\SpecialCharTok{::}\FunctionTok{schedule\_date\_vector\_GQL}\NormalTok{(}
\NormalTok{  bond\_schedule}
\NormalTok{)}

\NormalTok{first\_call\_position }\OtherTok{\textless{}{-}}\NormalTok{ non\_call\_periods }\SpecialCharTok{+} \DecValTok{1}\NormalTok{L}

\NormalTok{call\_dates }\OtherTok{\textless{}{-}}\NormalTok{ bond\_schedule\_dates[}
\NormalTok{  first\_call\_position}\SpecialCharTok{:}
\NormalTok{    (}\FunctionTok{length}\NormalTok{(bond\_schedule\_dates) }\SpecialCharTok{{-}} \DecValTok{1}\NormalTok{L)}
\NormalTok{]}

\NormalTok{call\_schedule }\OtherTok{\textless{}{-}}\NormalTok{ GaussQuant}\SpecialCharTok{::}\FunctionTok{make\_callability\_schedule\_GQL}\NormalTok{(}
  \AttributeTok{call\_dates =}\NormalTok{ call\_dates,}
  \AttributeTok{call\_price\_clean =}\NormalTok{ call\_price\_clean}
\NormalTok{)}

\NormalTok{call\_schedule\_tbl }\OtherTok{\textless{}{-}}\NormalTok{ tibble}\SpecialCharTok{::}\FunctionTok{tibble}\NormalTok{(}
  \AttributeTok{call\_number =} \FunctionTok{seq\_along}\NormalTok{(call\_dates),}
  \AttributeTok{call\_date =}\NormalTok{ call\_dates,}
  \AttributeTok{call\_price =} \FunctionTok{rep}\NormalTok{(}
\NormalTok{    call\_price\_clean,}
    \FunctionTok{length}\NormalTok{(call\_dates)}
\NormalTok{  )}
\NormalTok{)}

\NormalTok{GaussQuant}\SpecialCharTok{::}\FunctionTok{show\_tbl\_GQH}\NormalTok{(}
\NormalTok{  call\_schedule\_tbl,}
  \StringTok{"Callable bond exercise schedule"}\NormalTok{,}
  \AttributeTok{n =} \DecValTok{20}
\NormalTok{)}
\CommentTok{\#\textgreater{} \# A tibble: 8 x 3}
\CommentTok{\#\textgreater{}   call\_number call\_date  call\_price}
\CommentTok{\#\textgreater{}         \textless{}int\textgreater{} \textless{}date\textgreater{}          \textless{}dbl\textgreater{}}
\CommentTok{\#\textgreater{} 1           1 2025{-}04{-}17        100}
\CommentTok{\#\textgreater{} 2           2 2025{-}10{-}17        100}
\CommentTok{\#\textgreater{} 3           3 2026{-}04{-}17        100}
\CommentTok{\#\textgreater{} 4           4 2026{-}10{-}17        100}
\CommentTok{\#\textgreater{} 5           5 2027{-}04{-}17        100}
\CommentTok{\#\textgreater{} 6           6 2027{-}10{-}17        100}
\CommentTok{\#\textgreater{} 7           7 2028{-}04{-}17        100}
\CommentTok{\#\textgreater{} 8           8 2028{-}10{-}17        100}
\end{Highlighting}
\end{Shaded}

\subsection{7.2.2 Callable
BondのTree評価}\label{callable-bondux306etreeux8a55ux4fa1}

Callable Bondを作成し、Hull--Whiteモデルに基づく金利TreeをPricing
Engineとして設定する。

各Tree
nodeでは、発行体がその時点でコールする価値と、債券を存続させる価値が比較される。発行体にとってコールが有利であれば期限前償還が選択される。

\begin{Shaded}
\begin{Highlighting}[]
\NormalTok{callable\_bond }\OtherTok{\textless{}{-}}\NormalTok{ QuantLib}\SpecialCharTok{::}\FunctionTok{CallableFixedRateBond}\NormalTok{(}
\NormalTok{  settlement\_days,}
\NormalTok{  face\_amount,}
\NormalTok{  bond\_schedule,}
  \FunctionTok{c}\NormalTok{(coupon\_rate),}
\NormalTok{  bond\_day\_counter,}
  \StringTok{"Unadjusted"}\NormalTok{,}
\NormalTok{  call\_price\_clean,}
\NormalTok{  issue\_date,}
\NormalTok{  call\_schedule}
\NormalTok{)}

\NormalTok{callable\_bond\_engine }\OtherTok{\textless{}{-}}\NormalTok{ QuantLib}\SpecialCharTok{::}\FunctionTok{TreeCallableFixedRateBondEngine}\NormalTok{(}
\NormalTok{  hull\_white\_model,}
\NormalTok{  tree\_steps}
\NormalTok{)}

\NormalTok{callable\_bond}\SpecialCharTok{$}\FunctionTok{setPricingEngine}\NormalTok{(}
\NormalTok{  callable\_bond\_engine}
\NormalTok{)}
\CommentTok{\#\textgreater{} NULL}

\NormalTok{callable\_clean\_price }\OtherTok{\textless{}{-}}\NormalTok{ GaussQuant}\SpecialCharTok{::}\FunctionTok{callable\_bond\_clean\_price\_safe\_GQL}\NormalTok{(}
\NormalTok{  callable\_bond}
\NormalTok{)}

\NormalTok{callable\_dirty\_price }\OtherTok{\textless{}{-}}\NormalTok{ GaussQuant}\SpecialCharTok{::}\FunctionTok{callable\_bond\_dirty\_price\_safe\_GQL}\NormalTok{(}
\NormalTok{  callable\_bond}
\NormalTok{)}

\NormalTok{callable\_settlement\_value }\OtherTok{\textless{}{-}}\NormalTok{ GaussQuant}\SpecialCharTok{::}\FunctionTok{callable\_bond\_settlement\_value\_safe\_GQL}\NormalTok{(}
\NormalTok{  callable\_bond}
\NormalTok{)}

\NormalTok{embedded\_call\_value }\OtherTok{\textless{}{-}}\NormalTok{ bullet\_clean\_price }\SpecialCharTok{{-}}
\NormalTok{  callable\_clean\_price}

\NormalTok{callable\_summary\_tbl }\OtherTok{\textless{}{-}}\NormalTok{ tibble}\SpecialCharTok{::}\FunctionTok{tribble}\NormalTok{(}
  \SpecialCharTok{\textasciitilde{}}\NormalTok{metric,                      }\SpecialCharTok{\textasciitilde{}}\NormalTok{value,}
  \StringTok{"bullet\_clean\_price"}\NormalTok{,          bullet\_clean\_price,}
  \StringTok{"callable\_clean\_price"}\NormalTok{,        callable\_clean\_price,}
  \StringTok{"callable\_dirty\_price"}\NormalTok{,        callable\_dirty\_price,}
  \StringTok{"callable\_settlement\_value"}\NormalTok{,   callable\_settlement\_value,}
  \StringTok{"embedded\_call\_value"}\NormalTok{,         embedded\_call\_value}
\NormalTok{)}

\NormalTok{GaussQuant}\SpecialCharTok{::}\FunctionTok{show\_tbl\_GQH}\NormalTok{(}
\NormalTok{  callable\_summary\_tbl,}
  \StringTok{"Callable bond summary"}\NormalTok{,}
  \AttributeTok{n =} \DecValTok{20}
\NormalTok{)}
\CommentTok{\#\textgreater{} \# A tibble: 5 x 2}
\CommentTok{\#\textgreater{}   metric                     value}
\CommentTok{\#\textgreater{}   \textless{}chr\textgreater{}                      \textless{}dbl\textgreater{}}
\CommentTok{\#\textgreater{} 1 bullet\_clean\_price        100.0 }
\CommentTok{\#\textgreater{} 2 callable\_clean\_price       98.0 }
\CommentTok{\#\textgreater{} 3 callable\_dirty\_price       98.0 }
\CommentTok{\#\textgreater{} 4 callable\_settlement\_value  98.0 }
\CommentTok{\#\textgreater{} 5 embedded\_call\_value         2.00}
\end{Highlighting}
\end{Shaded}

Callable Bond価格がBullet
Bond価格を下回る差額は、投資家が発行体へ与えている期限前償還権の価値と解釈できる。

\section{7.3
Hull--Whiteパラメータ感応度}\label{hullwhiteux30d1ux30e9ux30e1ux30fcux30bfux611fux5fdcux5ea6}

短期金利のボラティリティが高いほど、将来金利が大きく低下する可能性も高くなる。したがって、発行体のコール権は一般に高くなり、Callable
Bond価格は低下する。

平均回帰速度 \(a\)
は、金利ショックがどの程度の速さで消滅するかを決める。平均回帰が速い場合、長い期間にわたる金利ショックの影響は小さくなる。

\begin{Shaded}
\begin{Highlighting}[]
\NormalTok{callable\_sensitivity\_input\_tbl }\OtherTok{\textless{}{-}}\NormalTok{ tibble}\SpecialCharTok{::}\FunctionTok{tribble}\NormalTok{(}
  \SpecialCharTok{\textasciitilde{}}\NormalTok{scenario,       }\SpecialCharTok{\textasciitilde{}}\NormalTok{mean\_reversion, }\SpecialCharTok{\textasciitilde{}}\NormalTok{sigma,}
  \StringTok{"low\_sigma"}\NormalTok{,      }\FloatTok{0.03}\NormalTok{,            }\FloatTok{0.005}\NormalTok{,}
  \StringTok{"base"}\NormalTok{,           }\FloatTok{0.03}\NormalTok{,            }\FloatTok{0.010}\NormalTok{,}
  \StringTok{"high\_sigma"}\NormalTok{,     }\FloatTok{0.03}\NormalTok{,            }\FloatTok{0.020}\NormalTok{,}
  \StringTok{"slow\_reversion"}\NormalTok{, }\FloatTok{0.01}\NormalTok{,            }\FloatTok{0.010}\NormalTok{,}
  \StringTok{"fast\_reversion"}\NormalTok{, }\FloatTok{0.10}\NormalTok{,            }\FloatTok{0.010}
\NormalTok{)}

\NormalTok{callable\_sensitivity\_tbl }\OtherTok{\textless{}{-}}\NormalTok{ purrr}\SpecialCharTok{::}\FunctionTok{pmap\_dfr}\NormalTok{(}
\NormalTok{  callable\_sensitivity\_input\_tbl,}
  \ControlFlowTok{function}\NormalTok{(}
\NormalTok{    scenario,}
\NormalTok{    mean\_reversion,}
\NormalTok{    sigma}
\NormalTok{  ) \{}
\NormalTok{    model\_i }\OtherTok{\textless{}{-}}\NormalTok{ GaussQuant}\SpecialCharTok{::}\FunctionTok{hull\_white\_model\_GQL}\NormalTok{(}
      \AttributeTok{term\_structure =}\NormalTok{ flat\_curve\_handle,}
      \AttributeTok{a =}\NormalTok{ mean\_reversion,}
      \AttributeTok{sigma =}\NormalTok{ sigma}
\NormalTok{    )}

\NormalTok{    engine\_i }\OtherTok{\textless{}{-}}\NormalTok{ QuantLib}\SpecialCharTok{::}\FunctionTok{TreeCallableFixedRateBondEngine}\NormalTok{(}
\NormalTok{      model\_i,}
\NormalTok{      tree\_steps}
\NormalTok{    )}

\NormalTok{    callable\_bond}\SpecialCharTok{$}\FunctionTok{setPricingEngine}\NormalTok{(}
\NormalTok{      engine\_i}
\NormalTok{    )}

\NormalTok{    callable\_price\_i }\OtherTok{\textless{}{-}}\NormalTok{ GaussQuant}\SpecialCharTok{::}\FunctionTok{callable\_bond\_clean\_price\_safe\_GQL}\NormalTok{(}
\NormalTok{      callable\_bond}
\NormalTok{    )}

\NormalTok{    tibble}\SpecialCharTok{::}\FunctionTok{tibble}\NormalTok{(}
      \AttributeTok{scenario =}\NormalTok{ scenario,}
      \AttributeTok{mean\_reversion =}\NormalTok{ mean\_reversion,}
      \AttributeTok{sigma =}\NormalTok{ sigma,}
      \AttributeTok{callable\_clean\_price =}\NormalTok{ callable\_price\_i,}
      \AttributeTok{embedded\_call\_value =}\NormalTok{ bullet\_clean\_price }\SpecialCharTok{{-}}
\NormalTok{        callable\_price\_i}
\NormalTok{    )}
\NormalTok{  \}}
\NormalTok{)}

\NormalTok{callable\_bond}\SpecialCharTok{$}\FunctionTok{setPricingEngine}\NormalTok{(}
\NormalTok{  callable\_bond\_engine}
\NormalTok{)}
\CommentTok{\#\textgreater{} NULL}

\NormalTok{GaussQuant}\SpecialCharTok{::}\FunctionTok{show\_tbl\_GQH}\NormalTok{(}
\NormalTok{  callable\_sensitivity\_tbl,}
  \StringTok{"Callable bond Hull{-}White sensitivity"}\NormalTok{,}
  \AttributeTok{n =} \DecValTok{20}
\NormalTok{)}
\CommentTok{\#\textgreater{} \# A tibble: 5 x 5}
\CommentTok{\#\textgreater{}   scenario       mean\_reversion sigma callable\_clean\_price embedded\_call\_value}
\CommentTok{\#\textgreater{}   \textless{}chr\textgreater{}                   \textless{}dbl\textgreater{} \textless{}dbl\textgreater{}                \textless{}dbl\textgreater{}               \textless{}dbl\textgreater{}}
\CommentTok{\#\textgreater{} 1 low\_sigma                0.03 0.005                 99.0                1.04}
\CommentTok{\#\textgreater{} 2 base                     0.03 0.01                  98.0                2.00}
\CommentTok{\#\textgreater{} 3 high\_sigma               0.03 0.02                  96.1                3.91}
\CommentTok{\#\textgreater{} 4 slow\_reversion           0.01 0.01                  97.9                2.08}
\CommentTok{\#\textgreater{} 5 fast\_reversion           0.1  0.01                  98.2                1.76}
\end{Highlighting}
\end{Shaded}

パラメータ感応度は他の条件を固定した比較である。市場カリブレーション後のパラメータは互いに独立ではないため、実務上のリスク分析ではカーブとボラティリティ市場を同時に考慮する必要がある。

\section{7.4 Bermudan Receiver
Swaptionとの比較}\label{bermudan-receiver-swaptionux3068ux306eux6bd4ux8f03}

発行体が債券をコールする誘因は、市場金利がクーポン金利を下回ったときに強くなる。これは固定金利を受け取るReceiver
Swapの価値が、金利低下時に上昇する性質と対応する。

そこで、債券のクーポンスケジュールを固定レグとするReceiver
Swapを作成し、Callable Bondと同じ行使日を持つBermudan Receiver
Swaptionを評価する。

\begin{Shaded}
\begin{Highlighting}[]
\NormalTok{floating\_index }\OtherTok{\textless{}{-}}\NormalTok{ QuantLib}\SpecialCharTok{::}\FunctionTok{JPYLibor}\NormalTok{(}
\NormalTok{  GaussQuant}\SpecialCharTok{::}\FunctionTok{period\_months\_GQL}\NormalTok{(}\DecValTok{6}\NormalTok{),}
\NormalTok{  flat\_curve\_handle}
\NormalTok{)}

\FunctionTok{tryCatch}\NormalTok{(}
\NormalTok{  floating\_index}\SpecialCharTok{$}\FunctionTok{addFixing}\NormalTok{(}
\NormalTok{    trade\_date,}
\NormalTok{    market\_yield,}
    \ConstantTok{TRUE}
\NormalTok{  ),}
  \AttributeTok{error =} \ControlFlowTok{function}\NormalTok{(e) }\ConstantTok{NULL}
\NormalTok{)}
\CommentTok{\#\textgreater{} NULL}

\NormalTok{underlying\_receiver\_swap }\OtherTok{\textless{}{-}}\NormalTok{ QuantLib}\SpecialCharTok{::}\FunctionTok{VanillaSwap}\NormalTok{(}
\NormalTok{  QuantLib}\SpecialCharTok{::}\FunctionTok{Swap\_Receiver\_get}\NormalTok{(),}
\NormalTok{  face\_amount,}
\NormalTok{  bond\_schedule,}
\NormalTok{  coupon\_rate,}
\NormalTok{  bond\_day\_counter,}
\NormalTok{  bond\_schedule,}
\NormalTok{  floating\_index,}
  \DecValTok{0}\NormalTok{,}
\NormalTok{  bond\_day\_counter}
\NormalTok{)}

\NormalTok{bermudan\_receiver\_engine }\OtherTok{\textless{}{-}}\NormalTok{ GaussQuant}\SpecialCharTok{::}\FunctionTok{hull\_white\_swaption\_engine\_GQL}\NormalTok{(}
  \AttributeTok{term\_structure =}\NormalTok{ flat\_curve\_handle,}
  \AttributeTok{a =}\NormalTok{ hull\_white\_a,}
  \AttributeTok{sigma =}\NormalTok{ hull\_white\_sigma,}
  \AttributeTok{method =} \StringTok{"tree"}\NormalTok{,}
  \AttributeTok{time\_steps =}\NormalTok{ tree\_steps}
\NormalTok{)}

\NormalTok{bermudan\_receiver\_swaption }\OtherTok{\textless{}{-}}\NormalTok{ GaussQuant}\SpecialCharTok{::}\FunctionTok{make\_bermudan\_swaption\_GQL}\NormalTok{(}
  \AttributeTok{underlying\_swap =}\NormalTok{ underlying\_receiver\_swap,}
  \AttributeTok{exercise\_dates =}\NormalTok{ call\_dates,}
  \AttributeTok{pricing\_engine =}\NormalTok{ bermudan\_receiver\_engine}
\NormalTok{)}

\NormalTok{bermudan\_receiver\_npv }\OtherTok{\textless{}{-}}\NormalTok{ GaussQuant}\SpecialCharTok{::}\FunctionTok{swaption\_npv\_GQL}\NormalTok{(}
\NormalTok{  bermudan\_receiver\_swaption}
\NormalTok{)}

\NormalTok{receiver\_swaption\_value\_per\_100 }\OtherTok{\textless{}{-}}\NormalTok{ bermudan\_receiver\_npv }\SpecialCharTok{/}
\NormalTok{  face\_amount }\SpecialCharTok{*}
  \DecValTok{100}

\NormalTok{bullet\_minus\_receiver\_swaption }\OtherTok{\textless{}{-}}\NormalTok{ bullet\_clean\_price }\SpecialCharTok{{-}}
\NormalTok{  receiver\_swaption\_value\_per\_100}

\NormalTok{callable\_comparison\_tbl }\OtherTok{\textless{}{-}}\NormalTok{ tibble}\SpecialCharTok{::}\FunctionTok{tribble}\NormalTok{(}
  \SpecialCharTok{\textasciitilde{}}\NormalTok{valuation,                           }\SpecialCharTok{\textasciitilde{}}\NormalTok{value,}
  \StringTok{"bullet\_bond\_clean\_price"}\NormalTok{,             bullet\_clean\_price,}
  \StringTok{"callable\_bond\_clean\_price"}\NormalTok{,           callable\_clean\_price,}
  \StringTok{"receiver\_swaption\_value\_per\_100"}\NormalTok{,     receiver\_swaption\_value\_per\_100,}
  \StringTok{"bullet\_minus\_receiver\_swaption"}\NormalTok{,      bullet\_minus\_receiver\_swaption,}
  \StringTok{"pricing\_difference"}\NormalTok{,                  callable\_clean\_price }\SpecialCharTok{{-}}
\NormalTok{                                           bullet\_minus\_receiver\_swaption}
\NormalTok{)}

\NormalTok{GaussQuant}\SpecialCharTok{::}\FunctionTok{show\_tbl\_GQH}\NormalTok{(}
\NormalTok{  callable\_comparison\_tbl,}
  \StringTok{"Callable bond and Bermudan receiver swaption comparison"}\NormalTok{,}
  \AttributeTok{n =} \DecValTok{20}
\NormalTok{)}
\CommentTok{\#\textgreater{} \# A tibble: 5 x 2}
\CommentTok{\#\textgreater{}   valuation                         value}
\CommentTok{\#\textgreater{}   \textless{}chr\textgreater{}                             \textless{}dbl\textgreater{}}
\CommentTok{\#\textgreater{} 1 bullet\_bond\_clean\_price         100.0  }
\CommentTok{\#\textgreater{} 2 callable\_bond\_clean\_price        98.0  }
\CommentTok{\#\textgreater{} 3 receiver\_swaption\_value\_per\_100   1.89 }
\CommentTok{\#\textgreater{} 4 bullet\_minus\_receiver\_swaption   98.1  }
\CommentTok{\#\textgreater{} 5 pricing\_difference               {-}0.112}
\end{Highlighting}
\end{Shaded}

Callable BondとBermudan Receiver
Swaptionは、同じHull--Whiteモデルと行使日を使用している。しかし、次の条件は完全には一致していない。

\begin{itemize}
\tightlist
\item
  債券の元本償還とSwapの元本交換
\item
  Clean priceとSwaption NPVの決済単位
\item
  経過利子
\item
  コール価格
\item
  固定クーポンとSwap固定レグ
\item
  浮動レグのFixingおよび支払条件
\end{itemize}

したがって、\texttt{pricing\_difference}が厳密にゼロとなることは要求しない。この比較の目的は、Callable
Bondに埋め込まれた金利低下時のオプション価値を、Receiver
Swaptionとの関係から理解することである。

\section{7.5
VasicekモデルとHull--Whiteモデル}\label{vasicekux30e2ux30c7ux30ebux3068hullwhiteux30e2ux30c7ux30eb}

\subsection{7.5.1
Vasicekモデルとの違い}\label{vasicekux30e2ux30c7ux30ebux3068ux306eux9055ux3044}

Vasicekモデルは、短期金利が一定の長期平均へ回帰するGaussianモデルである。

\[
dr_t
=
a(b-r_t)dt
+
\sigma\,dW_t
\]

ここで、\(b\) は長期平均金利である。

Hull--WhiteモデルもGaussian平均回帰モデルであるが、時間依存ドリフト
\(\theta(t)\)
を導入することによって、評価時点のイールドカーブへ適合できる。

{\def\LTcaptype{none} % do not increment counter
\begin{longtable}[]{@{}lll@{}}
\toprule\noalign{}
項目 & Vasicek & Hull--White \\
\midrule\noalign{}
\endhead
\bottomrule\noalign{}
\endlastfoot
平均回帰先 & 一定の長期平均 \(b\) & 時間依存ドリフトで調整 \\
初期カーブ & 一般には完全適合しない & 初期カーブへ適合可能 \\
主な入力 & \(a,b,\sigma,r_0\) & 初期カーブ、\(a,\sigma\) \\
金利分布 & Gaussian & Gaussian \\
負の金利 & 許容 & 許容 \\
\end{longtable}
}

Vasicekモデルは金利期間構造の長期的な形状をモデルパラメータから生成する。Hull--Whiteモデルは、観測カーブを出発点として、その周りの確率的変動を表現する。

\subsection{7.5.2
VasicekとHull--Whiteのゼロ金利曲線}\label{vasicekux3068hullwhiteux306eux30bcux30edux91d1ux5229ux66f2ux7dda}

Vasicekでは初期短期金利を変えると、モデルが生成する期間構造全体が変化する。Hull--Whiteでは状態変数がゼロであれば初期カーブを再現し、状態変数のショックが平均回帰しながら期間構造へ影響する。

\begin{Shaded}
\begin{Highlighting}[]
\NormalTok{maturity\_grid }\OtherTok{\textless{}{-}} \FunctionTok{seq}\NormalTok{(}
  \FloatTok{0.25}\NormalTok{,}
  \DecValTok{10}\NormalTok{,}
  \AttributeTok{by =} \FloatTok{0.25}
\NormalTok{)}

\NormalTok{vasicek\_scenario\_tbl }\OtherTok{\textless{}{-}}\NormalTok{ tibble}\SpecialCharTok{::}\FunctionTok{tribble}\NormalTok{(}
  \SpecialCharTok{\textasciitilde{}}\NormalTok{scenario,     }\SpecialCharTok{\textasciitilde{}}\NormalTok{short\_rate,}
  \StringTok{"r0 = 1\%"}\NormalTok{,      }\FloatTok{0.01}\NormalTok{,}
  \StringTok{"r0 = 3\%"}\NormalTok{,      }\FloatTok{0.03}\NormalTok{,}
  \StringTok{"r0 = 5\%"}\NormalTok{,      }\FloatTok{0.05}
\NormalTok{)}

\NormalTok{vasicek\_curve\_tbl }\OtherTok{\textless{}{-}}\NormalTok{ purrr}\SpecialCharTok{::}\FunctionTok{pmap\_dfr}\NormalTok{(}
\NormalTok{  vasicek\_scenario\_tbl,}
  \ControlFlowTok{function}\NormalTok{(}
\NormalTok{    scenario,}
\NormalTok{    short\_rate}
\NormalTok{  ) \{}
\NormalTok{    tibble}\SpecialCharTok{::}\FunctionTok{tibble}\NormalTok{(}
      \AttributeTok{maturity =}\NormalTok{ maturity\_grid,}
      \AttributeTok{zero\_rate =}\NormalTok{ purrr}\SpecialCharTok{::}\FunctionTok{map\_dbl}\NormalTok{(}
\NormalTok{        maturity\_grid,}
        \ControlFlowTok{function}\NormalTok{(maturity\_value) \{}
\NormalTok{          GaussQuant}\SpecialCharTok{::}\FunctionTok{vasicek\_zero\_rate\_GQL}\NormalTok{(}
            \AttributeTok{time =} \DecValTok{0}\NormalTok{,}
            \AttributeTok{maturity =}\NormalTok{ maturity\_value,}
            \AttributeTok{mean\_reversion =} \FloatTok{0.30}\NormalTok{,}
            \AttributeTok{sigma =} \FloatTok{0.01}\NormalTok{,}
            \AttributeTok{long\_run\_rate =} \FloatTok{0.03}\NormalTok{,}
            \AttributeTok{short\_rate =}\NormalTok{ short\_rate}
\NormalTok{          )}
\NormalTok{        \}}
\NormalTok{      ),}
      \AttributeTok{model =} \StringTok{"Vasicek"}\NormalTok{,}
      \AttributeTok{scenario =}\NormalTok{ scenario}
\NormalTok{    )}
\NormalTok{  \}}
\NormalTok{)}

\NormalTok{hull\_white\_scenario\_tbl }\OtherTok{\textless{}{-}}\NormalTok{ tibble}\SpecialCharTok{::}\FunctionTok{tribble}\NormalTok{(}
  \SpecialCharTok{\textasciitilde{}}\NormalTok{scenario,       }\SpecialCharTok{\textasciitilde{}}\NormalTok{state\_variable,}
  \StringTok{"x0 = {-}0.3\%"}\NormalTok{,    }\SpecialCharTok{{-}}\FloatTok{0.003}\NormalTok{,}
  \StringTok{"base curve"}\NormalTok{,     }\FloatTok{0.000}\NormalTok{,}
  \StringTok{"x0 = 0.3\%"}\NormalTok{,      }\FloatTok{0.003}
\NormalTok{)}

\NormalTok{hull\_white\_curve\_tbl }\OtherTok{\textless{}{-}}\NormalTok{ purrr}\SpecialCharTok{::}\FunctionTok{pmap\_dfr}\NormalTok{(}
\NormalTok{  hull\_white\_scenario\_tbl,}
  \ControlFlowTok{function}\NormalTok{(}
\NormalTok{    scenario,}
\NormalTok{    state\_variable}
\NormalTok{  ) \{}
\NormalTok{    tibble}\SpecialCharTok{::}\FunctionTok{tibble}\NormalTok{(}
      \AttributeTok{maturity =}\NormalTok{ maturity\_grid,}
      \AttributeTok{zero\_rate =}\NormalTok{ purrr}\SpecialCharTok{::}\FunctionTok{map\_dbl}\NormalTok{(}
\NormalTok{        maturity\_grid,}
        \ControlFlowTok{function}\NormalTok{(maturity\_value) \{}
\NormalTok{          GaussQuant}\SpecialCharTok{::}\FunctionTok{hull\_white\_zero\_rate\_GQL}\NormalTok{(}
            \AttributeTok{time =} \DecValTok{0}\NormalTok{,}
            \AttributeTok{maturity =}\NormalTok{ maturity\_value,}
            \AttributeTok{curve =}\NormalTok{ flat\_curve,}
            \AttributeTok{mean\_reversion =} \FloatTok{0.30}\NormalTok{,}
            \AttributeTok{sigma =} \FloatTok{0.01}\NormalTok{,}
            \AttributeTok{state\_variable =}\NormalTok{ state\_variable}
\NormalTok{          )}
\NormalTok{        \}}
\NormalTok{      ),}
      \AttributeTok{model =} \StringTok{"Hull{-}White"}\NormalTok{,}
      \AttributeTok{scenario =}\NormalTok{ scenario}
\NormalTok{    )}
\NormalTok{  \}}
\NormalTok{)}

\NormalTok{short\_rate\_curve\_tbl }\OtherTok{\textless{}{-}}\NormalTok{ dplyr}\SpecialCharTok{::}\FunctionTok{bind\_rows}\NormalTok{(}
\NormalTok{  vasicek\_curve\_tbl,}
\NormalTok{  hull\_white\_curve\_tbl}
\NormalTok{)}

\NormalTok{short\_rate\_curve\_plot }\OtherTok{\textless{}{-}}\NormalTok{ ggplot2}\SpecialCharTok{::}\FunctionTok{ggplot}\NormalTok{(}
\NormalTok{  short\_rate\_curve\_tbl,}
\NormalTok{  ggplot2}\SpecialCharTok{::}\FunctionTok{aes}\NormalTok{(}
    \AttributeTok{x =}\NormalTok{ .data}\SpecialCharTok{$}\NormalTok{maturity,}
    \AttributeTok{y =}\NormalTok{ .data}\SpecialCharTok{$}\NormalTok{zero\_rate,}
    \AttributeTok{group =}\NormalTok{ .data}\SpecialCharTok{$}\NormalTok{scenario,}
    \AttributeTok{linetype =}\NormalTok{ .data}\SpecialCharTok{$}\NormalTok{scenario}
\NormalTok{  )}
\NormalTok{) }\SpecialCharTok{+}
\NormalTok{  ggplot2}\SpecialCharTok{::}\FunctionTok{geom\_line}\NormalTok{(}
    \AttributeTok{linewidth =} \FloatTok{0.7}
\NormalTok{  ) }\SpecialCharTok{+}
\NormalTok{  ggplot2}\SpecialCharTok{::}\FunctionTok{facet\_wrap}\NormalTok{(}
    \SpecialCharTok{\textasciitilde{}}\NormalTok{model,}
    \AttributeTok{scales =} \StringTok{"free\_y"}
\NormalTok{  ) }\SpecialCharTok{+}
\NormalTok{  ggplot2}\SpecialCharTok{::}\FunctionTok{scale\_y\_continuous}\NormalTok{(}
    \AttributeTok{labels =} \ControlFlowTok{function}\NormalTok{(x) }\FunctionTok{sprintf}\NormalTok{(}\StringTok{"\%.2f\%\%"}\NormalTok{, }\DecValTok{100} \SpecialCharTok{*}\NormalTok{ x)}
\NormalTok{  ) }\SpecialCharTok{+}
\NormalTok{  ggplot2}\SpecialCharTok{::}\FunctionTok{labs}\NormalTok{(}
    \AttributeTok{title =} \StringTok{"Vasicek and Hull{-}White zero{-}rate curves"}\NormalTok{,}
    \AttributeTok{x =} \StringTok{"Residual maturity (years)"}\NormalTok{,}
    \AttributeTok{y =} \StringTok{"Continuously compounded zero rate"}\NormalTok{,}
    \AttributeTok{linetype =} \StringTok{"Scenario"}
\NormalTok{  ) }\SpecialCharTok{+}
\NormalTok{  ggplot2}\SpecialCharTok{::}\FunctionTok{theme\_minimal}\NormalTok{()}

\NormalTok{GaussQuant}\SpecialCharTok{::}\FunctionTok{show\_tbl\_GQH}\NormalTok{(}
\NormalTok{  short\_rate\_curve\_tbl,}
  \StringTok{"Short{-}rate model zero{-}rate curves"}\NormalTok{,}
  \AttributeTok{n =} \DecValTok{30}
\NormalTok{)}
\CommentTok{\#\textgreater{} \# A tibble: 30 x 4}
\CommentTok{\#\textgreater{}    maturity zero\_rate model   scenario}
\CommentTok{\#\textgreater{}       \textless{}dbl\textgreater{}     \textless{}dbl\textgreater{} \textless{}chr\textgreater{}   \textless{}chr\textgreater{}   }
\CommentTok{\#\textgreater{}  1     0.25    0.0107 Vasicek r0 = 1\% }
\CommentTok{\#\textgreater{}  2     0.5     0.0114 Vasicek r0 = 1\% }
\CommentTok{\#\textgreater{}  3     0.75    0.0121 Vasicek r0 = 1\% }
\CommentTok{\#\textgreater{}  4     1       0.0127 Vasicek r0 = 1\% }
\CommentTok{\#\textgreater{}  5     1.25    0.0133 Vasicek r0 = 1\% }
\CommentTok{\#\textgreater{}  6     1.5     0.0139 Vasicek r0 = 1\% }
\CommentTok{\#\textgreater{}  7     1.75    0.0144 Vasicek r0 = 1\% }
\CommentTok{\#\textgreater{}  8     2       0.0149 Vasicek r0 = 1\% }
\CommentTok{\#\textgreater{}  9     2.25    0.0154 Vasicek r0 = 1\% }
\CommentTok{\#\textgreater{} 10     2.5     0.0159 Vasicek r0 = 1\% }
\CommentTok{\#\textgreater{} \# i 20 more rows}

\FunctionTok{print}\NormalTok{(short\_rate\_curve\_plot)}
\end{Highlighting}
\end{Shaded}

\pandocbounded{\includegraphics[keepaspectratio]{part2_quantive_riemann/ch07_callable_hull_white_files/figure-pdf/chapter07-short-rate-curves-1.pdf}}

Hull--Whiteの\texttt{base\ curve}が初期フラットカーブ付近に位置することが、初期期間構造へ適合する性質を表している。状態変数のショックは短期側で大きく、満期が長くなるほど平均回帰によって影響が小さくなる。

\section{7.6
カリブレーションと数値検証}\label{ux30abux30eaux30d6ux30ecux30fcux30b7ux30e7ux30f3ux3068ux6570ux5024ux691cux8a3c}

\subsection{7.6.1
Hull--Whiteモデルのカリブレーション}\label{hullwhiteux30e2ux30c7ux30ebux306eux30abux30eaux30d6ux30ecux30fcux30b7ux30e7ux30f3}

Hull--Whiteモデルのパラメータ \(a\) と \(\sigma\)
は、Swaption市場の価格またはインプライド・ボラティリティへ適合させる。

本節では、1年後に開始して5年間継続する1Y×5Y ATM
Swaptionを使用する。市場Normal volatilityを36.06bpとし、平均回帰速度
\(a=0.03\) を固定して、短期金利ボラティリティ \(\sigma\)
だけを推定する。

単一の市場商品から \(a\) と \(\sigma\)
の双方を安定して識別することはできない。複数満期のSwaptionを利用する場合には、複数のCalibration
Helperを同時に用いて両パラメータを推定する。

\subsection{7.6.2 1Y×5Y ATM Swap}\label{y5y-atm-swap}

\begin{Shaded}
\begin{Highlighting}[]
\NormalTok{GaussQuant}\SpecialCharTok{::}\FunctionTok{set\_eval\_date\_GQL}\NormalTok{(trade\_date)}

\NormalTok{calibration\_calendar }\OtherTok{\textless{}{-}}\NormalTok{ QuantLib}\SpecialCharTok{::}\FunctionTok{Japan}\NormalTok{()}
\NormalTok{calibration\_day\_counter }\OtherTok{\textless{}{-}}\NormalTok{ QuantLib}\SpecialCharTok{::}\FunctionTok{Actual365Fixed}\NormalTok{()}

\NormalTok{swaption\_expiry\_date }\OtherTok{\textless{}{-}}\NormalTok{ calibration\_calendar}\SpecialCharTok{$}\FunctionTok{advance}\NormalTok{(}
\NormalTok{  trade\_date,}
  \DecValTok{1}\NormalTok{L,}
  \StringTok{"Years"}
\NormalTok{)}

\NormalTok{swap\_effective\_date }\OtherTok{\textless{}{-}}\NormalTok{ calibration\_calendar}\SpecialCharTok{$}\FunctionTok{advance}\NormalTok{(}
\NormalTok{  swaption\_expiry\_date,}
  \DecValTok{2}\NormalTok{L,}
  \StringTok{"Days"}
\NormalTok{)}

\NormalTok{swap\_maturity\_date }\OtherTok{\textless{}{-}}\NormalTok{ calibration\_calendar}\SpecialCharTok{$}\FunctionTok{advance}\NormalTok{(}
\NormalTok{  swap\_effective\_date,}
  \DecValTok{5}\NormalTok{L,}
  \StringTok{"Years"}
\NormalTok{)}

\NormalTok{calibration\_schedule }\OtherTok{\textless{}{-}}\NormalTok{ QuantLib}\SpecialCharTok{::}\FunctionTok{Schedule}\NormalTok{(}
\NormalTok{  swap\_effective\_date,}
\NormalTok{  swap\_maturity\_date,}
\NormalTok{  GaussQuant}\SpecialCharTok{::}\FunctionTok{period\_months\_GQL}\NormalTok{(}\DecValTok{6}\NormalTok{),}
\NormalTok{  calibration\_calendar,}
  \StringTok{"ModifiedFollowing"}\NormalTok{,}
  \StringTok{"ModifiedFollowing"}\NormalTok{,}
  \StringTok{"Backward"}\NormalTok{,}
  \ConstantTok{FALSE}
\NormalTok{)}

\NormalTok{calibration\_index }\OtherTok{\textless{}{-}}\NormalTok{ QuantLib}\SpecialCharTok{::}\FunctionTok{JPYLibor}\NormalTok{(}
\NormalTok{  GaussQuant}\SpecialCharTok{::}\FunctionTok{period\_months\_GQL}\NormalTok{(}\DecValTok{6}\NormalTok{),}
\NormalTok{  flat\_curve\_handle}
\NormalTok{)}

\NormalTok{calibration\_annuity }\OtherTok{\textless{}{-}}\NormalTok{ GaussQuant}\SpecialCharTok{::}\FunctionTok{swap\_annuity\_from\_schedule\_GQL}\NormalTok{(}
  \AttributeTok{schedule =}\NormalTok{ calibration\_schedule,}
  \AttributeTok{curve =}\NormalTok{ flat\_curve,}
  \AttributeTok{day\_counter =}\NormalTok{ calibration\_day\_counter}
\NormalTok{)}

\NormalTok{calibration\_forward\_swap\_rate }\OtherTok{\textless{}{-}}\NormalTok{ GaussQuant}\SpecialCharTok{::}\FunctionTok{forward\_swap\_rate\_from\_schedule\_GQL}\NormalTok{(}
  \AttributeTok{schedule =}\NormalTok{ calibration\_schedule,}
  \AttributeTok{curve =}\NormalTok{ flat\_curve,}
  \AttributeTok{day\_counter =}\NormalTok{ calibration\_day\_counter}
\NormalTok{)}
\NormalTok{calibration\_swap\_tbl }\OtherTok{\textless{}{-}}\NormalTok{ tibble}\SpecialCharTok{::}\FunctionTok{tribble}\NormalTok{(}
  \SpecialCharTok{\textasciitilde{}}\NormalTok{metric,                }\SpecialCharTok{\textasciitilde{}}\NormalTok{value,}
  \StringTok{"expiry\_date"}\NormalTok{,           GaussQuant}\SpecialCharTok{::}\FunctionTok{iso\_GQL}\NormalTok{(swaption\_expiry\_date),}
  \StringTok{"swap\_effective\_date"}\NormalTok{,   GaussQuant}\SpecialCharTok{::}\FunctionTok{iso\_GQL}\NormalTok{(swap\_effective\_date),}
  \StringTok{"swap\_maturity\_date"}\NormalTok{,    GaussQuant}\SpecialCharTok{::}\FunctionTok{iso\_GQL}\NormalTok{(swap\_maturity\_date),}
  \StringTok{"annuity"}\NormalTok{,               }\FunctionTok{as.character}\NormalTok{(calibration\_annuity),}
  \StringTok{"forward\_swap\_rate"}\NormalTok{,     }\FunctionTok{as.character}\NormalTok{(calibration\_forward\_swap\_rate)}
\NormalTok{)}

\NormalTok{GaussQuant}\SpecialCharTok{::}\FunctionTok{show\_tbl\_GQH}\NormalTok{(}
\NormalTok{  GaussQuant}\SpecialCharTok{::}\FunctionTok{schedule\_table\_GQL}\NormalTok{(}
\NormalTok{    calibration\_schedule}
\NormalTok{  ),}
  \StringTok{"1Yx5Y fixed{-}leg schedule"}\NormalTok{,}
  \AttributeTok{n =} \DecValTok{20}
\NormalTok{)}
\CommentTok{\#\textgreater{} \# A tibble: 11 x 2}
\CommentTok{\#\textgreater{}    period schedule\_date}
\CommentTok{\#\textgreater{}     \textless{}int\textgreater{} \textless{}date\textgreater{}       }
\CommentTok{\#\textgreater{}  1      1 2025{-}04{-}17   }
\CommentTok{\#\textgreater{}  2      2 2025{-}10{-}17   }
\CommentTok{\#\textgreater{}  3      3 2026{-}04{-}17   }
\CommentTok{\#\textgreater{}  4      4 2026{-}10{-}19   }
\CommentTok{\#\textgreater{}  5      5 2027{-}04{-}19   }
\CommentTok{\#\textgreater{}  6      6 2027{-}10{-}18   }
\CommentTok{\#\textgreater{}  7      7 2028{-}04{-}17   }
\CommentTok{\#\textgreater{}  8      8 2028{-}10{-}17   }
\CommentTok{\#\textgreater{}  9      9 2029{-}04{-}17   }
\CommentTok{\#\textgreater{} 10     10 2029{-}10{-}17   }
\CommentTok{\#\textgreater{} 11     11 2030{-}04{-}17}

\NormalTok{GaussQuant}\SpecialCharTok{::}\FunctionTok{show\_tbl\_GQH}\NormalTok{(}
\NormalTok{  calibration\_swap\_tbl,}
  \StringTok{"1Yx5Y forward swap summary"}\NormalTok{,}
  \AttributeTok{n =} \DecValTok{20}
\NormalTok{)}
\CommentTok{\#\textgreater{} \# A tibble: 5 x 2}
\CommentTok{\#\textgreater{}   metric              value             }
\CommentTok{\#\textgreater{}   \textless{}chr\textgreater{}               \textless{}chr\textgreater{}             }
\CommentTok{\#\textgreater{} 1 expiry\_date         2025{-}04{-}15        }
\CommentTok{\#\textgreater{} 2 swap\_effective\_date 2025{-}04{-}17        }
\CommentTok{\#\textgreater{} 3 swap\_maturity\_date  2030{-}04{-}17        }
\CommentTok{\#\textgreater{} 4 annuity             4.47364922408499  }
\CommentTok{\#\textgreater{} 5 forward\_swap\_rate   0.0302262727429665}
\end{Highlighting}
\end{Shaded}

\subsection{7.6.3 One-helper calibration}\label{one-helper-calibration}

\begin{Shaded}
\begin{Highlighting}[]
\NormalTok{market\_normal\_volatility }\OtherTok{\textless{}{-}} \FloatTok{36.06} \SpecialCharTok{/} \DecValTok{10000}
\NormalTok{fixed\_mean\_reversion }\OtherTok{\textless{}{-}} \FloatTok{0.03}
\NormalTok{initial\_sigma }\OtherTok{\textless{}{-}} \FloatTok{0.0001}

\NormalTok{calibration\_model }\OtherTok{\textless{}{-}}\NormalTok{ GaussQuant}\SpecialCharTok{::}\FunctionTok{hull\_white\_model\_GQL}\NormalTok{(}
  \AttributeTok{term\_structure =}\NormalTok{ flat\_curve\_handle,}
  \AttributeTok{a =}\NormalTok{ fixed\_mean\_reversion,}
  \AttributeTok{sigma =}\NormalTok{ initial\_sigma}
\NormalTok{)}

\NormalTok{calibration\_jamshidian\_engine }\OtherTok{\textless{}{-}}\NormalTok{ QuantLib}\SpecialCharTok{::}\FunctionTok{JamshidianSwaptionEngine}\NormalTok{(}
\NormalTok{  calibration\_model}
\NormalTok{)}

\NormalTok{swaption\_helper }\OtherTok{\textless{}{-}}\NormalTok{ QuantLib}\SpecialCharTok{::}\FunctionTok{SwaptionHelper}\NormalTok{(}
\NormalTok{  swaption\_expiry\_date,}
\NormalTok{  GaussQuant}\SpecialCharTok{::}\FunctionTok{period\_years\_GQL}\NormalTok{(}\DecValTok{5}\NormalTok{),}
\NormalTok{  QuantLib}\SpecialCharTok{::}\FunctionTok{QuoteHandle}\NormalTok{(}
\NormalTok{    QuantLib}\SpecialCharTok{::}\FunctionTok{SimpleQuote}\NormalTok{(}
\NormalTok{      market\_normal\_volatility}
\NormalTok{    )}
\NormalTok{  ),}
\NormalTok{  calibration\_index,}
\NormalTok{  GaussQuant}\SpecialCharTok{::}\FunctionTok{period\_months\_GQL}\NormalTok{(}\DecValTok{6}\NormalTok{),}
\NormalTok{  calibration\_day\_counter,}
\NormalTok{  calibration\_day\_counter,}
\NormalTok{  flat\_curve\_handle,}
\NormalTok{  QuantLib}\SpecialCharTok{::}\FunctionTok{BlackCalibrationHelper\_RelativePriceError\_get}\NormalTok{(),}
\NormalTok{  calibration\_forward\_swap\_rate,}
  \DecValTok{1}\NormalTok{,}
\NormalTok{  QuantLib}\SpecialCharTok{::}\FunctionTok{VolatilityType\_Normal\_get}\NormalTok{()}
\NormalTok{)}

\NormalTok{swaption\_helper}\SpecialCharTok{$}\FunctionTok{setPricingEngine}\NormalTok{(}
\NormalTok{  calibration\_jamshidian\_engine}
\NormalTok{)}
\CommentTok{\#\textgreater{} NULL}

\NormalTok{calibration\_helper\_vector }\OtherTok{\textless{}{-}}\NormalTok{ GaussQuant}\SpecialCharTok{::}\FunctionTok{push\_calibration\_helpers\_GQL}\NormalTok{(}
  \FunctionTok{list}\NormalTok{(swaption\_helper)}
\NormalTok{)}

\NormalTok{calibration\_end\_criteria }\OtherTok{\textless{}{-}}\NormalTok{ QuantLib}\SpecialCharTok{::}\FunctionTok{EndCriteria}\NormalTok{(}
  \DecValTok{10000}\NormalTok{L,}
  \DecValTok{100}\NormalTok{L,}
  \FloatTok{1e{-}6}\NormalTok{,}
  \FloatTok{1e{-}8}\NormalTok{,}
  \FloatTok{1e{-}8}
\NormalTok{)}

\NormalTok{calibration\_model}\SpecialCharTok{$}\FunctionTok{calibrate}\NormalTok{(}
\NormalTok{  calibration\_helper\_vector,}
\NormalTok{  QuantLib}\SpecialCharTok{::}\FunctionTok{LevenbergMarquardt}\NormalTok{(),}
\NormalTok{  calibration\_end\_criteria,}
\NormalTok{  QuantLib}\SpecialCharTok{::}\FunctionTok{NoConstraint}\NormalTok{(),}
  \FunctionTok{numeric}\NormalTok{(),}
  \FunctionTok{c}\NormalTok{(}
    \ConstantTok{TRUE}\NormalTok{,}
    \ConstantTok{FALSE}
\NormalTok{  )}
\NormalTok{)}
\CommentTok{\#\textgreater{} NULL}

\NormalTok{calibration\_parameter\_array }\OtherTok{\textless{}{-}}\NormalTok{ calibration\_model}\SpecialCharTok{$}\FunctionTok{params}\NormalTok{()}

\NormalTok{calibrated\_parameters }\OtherTok{\textless{}{-}} \FunctionTok{c}\NormalTok{(}
  \FunctionTok{as.numeric}\NormalTok{(calibration\_parameter\_array[}\DecValTok{1}\NormalTok{]),}
  \FunctionTok{as.numeric}\NormalTok{(calibration\_parameter\_array[}\DecValTok{2}\NormalTok{])}
\NormalTok{)}

\NormalTok{calibration\_result\_tbl }\OtherTok{\textless{}{-}}\NormalTok{ tibble}\SpecialCharTok{::}\FunctionTok{tribble}\NormalTok{(}
  \SpecialCharTok{\textasciitilde{}}\NormalTok{parameter,            }\SpecialCharTok{\textasciitilde{}}\NormalTok{value,}
  \StringTok{"a\_fixed"}\NormalTok{,              calibrated\_parameters[[}\DecValTok{1}\NormalTok{]],}
  \StringTok{"sigma\_calibrated"}\NormalTok{,     calibrated\_parameters[[}\DecValTok{2}\NormalTok{]],}
  \StringTok{"market\_normal\_vol"}\NormalTok{,    market\_normal\_volatility}
\NormalTok{)}

\NormalTok{GaussQuant}\SpecialCharTok{::}\FunctionTok{show\_tbl\_GQH}\NormalTok{(}
\NormalTok{  calibration\_result\_tbl,}
  \StringTok{"Hull{-}White one{-}helper calibration"}\NormalTok{,}
  \AttributeTok{n =} \DecValTok{20}
\NormalTok{)}
\CommentTok{\#\textgreater{} \# A tibble: 3 x 2}
\CommentTok{\#\textgreater{}   parameter           value}
\CommentTok{\#\textgreater{}   \textless{}chr\textgreater{}               \textless{}dbl\textgreater{}}
\CommentTok{\#\textgreater{} 1 a\_fixed           0.03   }
\CommentTok{\#\textgreater{} 2 sigma\_calibrated  0.00334}
\CommentTok{\#\textgreater{} 3 market\_normal\_vol 0.00361}
\NormalTok{calibration\_parameter\_array }\OtherTok{\textless{}{-}}\NormalTok{ calibration\_model}\SpecialCharTok{$}\FunctionTok{params}\NormalTok{()}

\NormalTok{calibration\_result\_tbl }\OtherTok{\textless{}{-}}\NormalTok{ tibble}\SpecialCharTok{::}\FunctionTok{tribble}\NormalTok{(}
  \SpecialCharTok{\textasciitilde{}}\NormalTok{parameter,             }\SpecialCharTok{\textasciitilde{}}\NormalTok{value,}
  \StringTok{"a\_fixed"}\NormalTok{,               calibrated\_parameters[[}\DecValTok{1}\NormalTok{]],}
  \StringTok{"sigma\_calibrated"}\NormalTok{,      calibrated\_parameters[[}\DecValTok{2}\NormalTok{]],}
  \StringTok{"market\_normal\_vol"}\NormalTok{,     market\_normal\_volatility}
\NormalTok{)}

\NormalTok{GaussQuant}\SpecialCharTok{::}\FunctionTok{show\_tbl\_GQH}\NormalTok{(}
\NormalTok{  calibration\_result\_tbl,}
  \StringTok{"Hull{-}White one{-}helper calibration"}\NormalTok{,}
  \AttributeTok{n =} \DecValTok{20}
\NormalTok{)}
\CommentTok{\#\textgreater{} \# A tibble: 3 x 2}
\CommentTok{\#\textgreater{}   parameter           value}
\CommentTok{\#\textgreater{}   \textless{}chr\textgreater{}               \textless{}dbl\textgreater{}}
\CommentTok{\#\textgreater{} 1 a\_fixed           0.03   }
\CommentTok{\#\textgreater{} 2 sigma\_calibrated  0.00334}
\CommentTok{\#\textgreater{} 3 market\_normal\_vol 0.00361}
\end{Highlighting}
\end{Shaded}

\texttt{c(TRUE,\ FALSE)}は、最初のパラメータ \(a\)
を固定し、2番目のパラメータ \(\sigma\) だけを最適化する指定である。

\subsection{7.6.4 European
SwaptionのEngine比較}\label{european-swaptionux306eengineux6bd4ux8f03}

カリブレーションに用いた1Y×5Y ATM Swapを原資産とするEuropean Payer
Swaptionを作成し、次の3つの方法で評価する。

\begin{itemize}
\tightlist
\item
  Jamshidian decomposition
\item
  Hull--White Tree
\item
  Bachelier analytic engine
\end{itemize}

JamshidianとTreeは同じカリブレーション後のHull--Whiteモデルを使用する。Bachelier
Engineは市場Normal volatilityを直接使用する。

\begin{Shaded}
\begin{Highlighting}[]
\NormalTok{calibration\_underlying\_swap }\OtherTok{\textless{}{-}}\NormalTok{ QuantLib}\SpecialCharTok{::}\FunctionTok{VanillaSwap}\NormalTok{(}
\NormalTok{  QuantLib}\SpecialCharTok{::}\FunctionTok{Swap\_Payer\_get}\NormalTok{(),}
  \DecValTok{1}\NormalTok{,}
\NormalTok{  calibration\_schedule,}
\NormalTok{  calibration\_forward\_swap\_rate,}
\NormalTok{  calibration\_day\_counter,}
\NormalTok{  calibration\_schedule,}
\NormalTok{  calibration\_index,}
  \DecValTok{0}\NormalTok{,}
\NormalTok{  calibration\_day\_counter}
\NormalTok{)}

\NormalTok{calibration\_european\_swaption }\OtherTok{\textless{}{-}}\NormalTok{ QuantLib}\SpecialCharTok{::}\FunctionTok{Swaption}\NormalTok{(}
\NormalTok{  calibration\_underlying\_swap,}
\NormalTok{  QuantLib}\SpecialCharTok{::}\FunctionTok{EuropeanExercise}\NormalTok{(}
\NormalTok{    swaption\_expiry\_date}
\NormalTok{  )}
\NormalTok{)}

\NormalTok{jamshidian\_engine }\OtherTok{\textless{}{-}}\NormalTok{ QuantLib}\SpecialCharTok{::}\FunctionTok{JamshidianSwaptionEngine}\NormalTok{(}
\NormalTok{  calibration\_model}
\NormalTok{)}

\NormalTok{calibration\_european\_swaption}\SpecialCharTok{$}\FunctionTok{setPricingEngine}\NormalTok{(}
\NormalTok{  jamshidian\_engine}
\NormalTok{)}
\CommentTok{\#\textgreater{} NULL}

\NormalTok{npv\_jamshidian }\OtherTok{\textless{}{-}}\NormalTok{ GaussQuant}\SpecialCharTok{::}\FunctionTok{instrument\_npv\_safe\_GQL}\NormalTok{(}
\NormalTok{  calibration\_european\_swaption}
\NormalTok{)}

\NormalTok{tree\_engine\_500 }\OtherTok{\textless{}{-}}\NormalTok{ QuantLib}\SpecialCharTok{::}\FunctionTok{TreeSwaptionEngine}\NormalTok{(}
\NormalTok{  calibration\_model,}
  \DecValTok{500}\NormalTok{L}
\NormalTok{)}

\NormalTok{calibration\_european\_swaption}\SpecialCharTok{$}\FunctionTok{setPricingEngine}\NormalTok{(}
\NormalTok{  tree\_engine\_500}
\NormalTok{)}
\CommentTok{\#\textgreater{} NULL}

\NormalTok{npv\_tree\_500 }\OtherTok{\textless{}{-}}\NormalTok{ GaussQuant}\SpecialCharTok{::}\FunctionTok{instrument\_npv\_safe\_GQL}\NormalTok{(}
\NormalTok{  calibration\_european\_swaption}
\NormalTok{)}

\NormalTok{bachelier\_engine }\OtherTok{\textless{}{-}}\NormalTok{ QuantLib}\SpecialCharTok{::}\FunctionTok{BachelierSwaptionEngine}\NormalTok{(}
\NormalTok{  flat\_curve\_handle,}
\NormalTok{  QuantLib}\SpecialCharTok{::}\FunctionTok{QuoteHandle}\NormalTok{(}
\NormalTok{    QuantLib}\SpecialCharTok{::}\FunctionTok{SimpleQuote}\NormalTok{(}
\NormalTok{      market\_normal\_volatility}
\NormalTok{    )}
\NormalTok{  )}
\NormalTok{)}

\NormalTok{calibration\_european\_swaption}\SpecialCharTok{$}\FunctionTok{setPricingEngine}\NormalTok{(}
\NormalTok{  bachelier\_engine}
\NormalTok{)}
\CommentTok{\#\textgreater{} NULL}

\NormalTok{npv\_bachelier }\OtherTok{\textless{}{-}}\NormalTok{ GaussQuant}\SpecialCharTok{::}\FunctionTok{instrument\_npv\_safe\_GQL}\NormalTok{(}
\NormalTok{  calibration\_european\_swaption}
\NormalTok{)}

\NormalTok{engine\_comparison\_tbl }\OtherTok{\textless{}{-}}\NormalTok{ tibble}\SpecialCharTok{::}\FunctionTok{tribble}\NormalTok{(}
  \SpecialCharTok{\textasciitilde{}}\NormalTok{engine,                  }\SpecialCharTok{\textasciitilde{}}\NormalTok{npv,}
  \StringTok{"Jamshidian"}\NormalTok{,              npv\_jamshidian,}
  \StringTok{"Hull{-}White tree 500"}\NormalTok{,     npv\_tree\_500,}
  \StringTok{"Bachelier"}\NormalTok{,               npv\_bachelier}
\NormalTok{) }\SpecialCharTok{|\textgreater{}}
\NormalTok{  dplyr}\SpecialCharTok{::}\FunctionTok{mutate}\NormalTok{(}
    \AttributeTok{difference\_vs\_bachelier =}\NormalTok{ .data}\SpecialCharTok{$}\NormalTok{npv }\SpecialCharTok{{-}}
\NormalTok{      npv\_bachelier,}
    \AttributeTok{difference\_vs\_jamshidian =}\NormalTok{ .data}\SpecialCharTok{$}\NormalTok{npv }\SpecialCharTok{{-}}
\NormalTok{      npv\_jamshidian}
\NormalTok{  )}

\NormalTok{GaussQuant}\SpecialCharTok{::}\FunctionTok{show\_tbl\_GQH}\NormalTok{(}
\NormalTok{  engine\_comparison\_tbl,}
  \StringTok{"1Yx5Y swaption engine comparison"}\NormalTok{,}
  \AttributeTok{n =} \DecValTok{20}
\NormalTok{)}
\CommentTok{\#\textgreater{} \# A tibble: 3 x 4}
\CommentTok{\#\textgreater{}   engine                  npv difference\_vs\_bachelier difference\_vs\_jamshidian}
\CommentTok{\#\textgreater{}   \textless{}chr\textgreater{}                 \textless{}dbl\textgreater{}                   \textless{}dbl\textgreater{}                    \textless{}dbl\textgreater{}}
\CommentTok{\#\textgreater{} 1 Jamshidian          0.00555           {-}0.0000000580             0           }
\CommentTok{\#\textgreater{} 2 Hull{-}White tree 500 0.00555           {-}0.00000395              {-}0.00000390  }
\CommentTok{\#\textgreater{} 3 Bachelier           0.00555            0                        0.0000000580}
\end{Highlighting}
\end{Shaded}

カリブレーション対象と同じATM
Swaptionであるため、Hull--White価格とBachelier市場価格は近い値になることが期待される。Jamshidianと十分に細かいTreeの差は、主にTree離散化誤差を表す。

\subsection{7.6.5 Tree stepの収束}\label{tree-stepux306eux53ceux675f}

Hull--White Treeでは、時間軸と金利状態を有限個のnodeへ離散化する。Tree
stepが少ない場合、金利分布とオプションPayoffの非線形性を粗くしか表現できない。

同じEuropean Swaptionを異なるTree
stepで評価し、Jamshidian価格への収束を確認する。

\begin{Shaded}
\begin{Highlighting}[]
\NormalTok{tree\_step\_grid }\OtherTok{\textless{}{-}} \FunctionTok{c}\NormalTok{(}
  \DecValTok{6}\NormalTok{L,}
  \DecValTok{12}\NormalTok{L,}
  \DecValTok{24}\NormalTok{L,}
  \DecValTok{60}\NormalTok{L,}
  \DecValTok{120}\NormalTok{L,}
  \DecValTok{240}\NormalTok{L,}
  \DecValTok{500}\NormalTok{L}
\NormalTok{)}

\NormalTok{tree\_convergence\_tbl }\OtherTok{\textless{}{-}}\NormalTok{ purrr}\SpecialCharTok{::}\FunctionTok{map\_dfr}\NormalTok{(}
\NormalTok{  tree\_step\_grid,}
  \ControlFlowTok{function}\NormalTok{(tree\_step) \{}
\NormalTok{    tree\_engine\_i }\OtherTok{\textless{}{-}}\NormalTok{ QuantLib}\SpecialCharTok{::}\FunctionTok{TreeSwaptionEngine}\NormalTok{(}
\NormalTok{      calibration\_model,}
\NormalTok{      tree\_step}
\NormalTok{    )}

\NormalTok{    calibration\_european\_swaption}\SpecialCharTok{$}\FunctionTok{setPricingEngine}\NormalTok{(}
\NormalTok{      tree\_engine\_i}
\NormalTok{    )}

\NormalTok{    tree\_npv\_i }\OtherTok{\textless{}{-}}\NormalTok{ GaussQuant}\SpecialCharTok{::}\FunctionTok{instrument\_npv\_safe\_GQL}\NormalTok{(}
\NormalTok{      calibration\_european\_swaption}
\NormalTok{    )}

\NormalTok{    tibble}\SpecialCharTok{::}\FunctionTok{tibble}\NormalTok{(}
      \AttributeTok{tree\_steps =}\NormalTok{ tree\_step,}
      \AttributeTok{npv =}\NormalTok{ tree\_npv\_i,}
      \AttributeTok{difference\_vs\_jamshidian =}\NormalTok{ tree\_npv\_i }\SpecialCharTok{{-}}
\NormalTok{        npv\_jamshidian,}
      \AttributeTok{difference\_vs\_bachelier =}\NormalTok{ tree\_npv\_i }\SpecialCharTok{{-}}
\NormalTok{        npv\_bachelier}
\NormalTok{    )}
\NormalTok{  \}}
\NormalTok{)}

\NormalTok{tree\_convergence\_plot }\OtherTok{\textless{}{-}}\NormalTok{ ggplot2}\SpecialCharTok{::}\FunctionTok{ggplot}\NormalTok{(}
\NormalTok{  tree\_convergence\_tbl,}
\NormalTok{  ggplot2}\SpecialCharTok{::}\FunctionTok{aes}\NormalTok{(}
    \AttributeTok{x =}\NormalTok{ .data}\SpecialCharTok{$}\NormalTok{tree\_steps,}
    \AttributeTok{y =}\NormalTok{ .data}\SpecialCharTok{$}\NormalTok{npv}
\NormalTok{  )}
\NormalTok{) }\SpecialCharTok{+}
\NormalTok{  ggplot2}\SpecialCharTok{::}\FunctionTok{geom\_line}\NormalTok{(}
    \AttributeTok{linewidth =} \FloatTok{0.7}
\NormalTok{  ) }\SpecialCharTok{+}
\NormalTok{  ggplot2}\SpecialCharTok{::}\FunctionTok{geom\_point}\NormalTok{(}
    \AttributeTok{size =} \FloatTok{2.2}
\NormalTok{  ) }\SpecialCharTok{+}
\NormalTok{  ggplot2}\SpecialCharTok{::}\FunctionTok{geom\_hline}\NormalTok{(}
    \AttributeTok{yintercept =}\NormalTok{ npv\_jamshidian,}
    \AttributeTok{linetype =} \StringTok{"dashed"}
\NormalTok{  ) }\SpecialCharTok{+}
\NormalTok{  ggplot2}\SpecialCharTok{::}\FunctionTok{labs}\NormalTok{(}
    \AttributeTok{title =} \StringTok{"Hull{-}White tree{-}step convergence"}\NormalTok{,}
    \AttributeTok{subtitle =} \StringTok{"Dashed line: Jamshidian price"}\NormalTok{,}
    \AttributeTok{x =} \StringTok{"Tree steps"}\NormalTok{,}
    \AttributeTok{y =} \StringTok{"Swaption NPV"}
\NormalTok{  ) }\SpecialCharTok{+}
\NormalTok{  ggplot2}\SpecialCharTok{::}\FunctionTok{theme\_minimal}\NormalTok{()}

\NormalTok{GaussQuant}\SpecialCharTok{::}\FunctionTok{show\_tbl\_GQH}\NormalTok{(}
\NormalTok{  tree\_convergence\_tbl,}
  \StringTok{"Hull{-}White tree convergence"}\NormalTok{,}
  \AttributeTok{n =} \DecValTok{20}
\NormalTok{)}
\CommentTok{\#\textgreater{} \# A tibble: 7 x 4}
\CommentTok{\#\textgreater{}   tree\_steps     npv difference\_vs\_jamshidian difference\_vs\_bachelier}
\CommentTok{\#\textgreater{}        \textless{}int\textgreater{}   \textless{}dbl\textgreater{}                    \textless{}dbl\textgreater{}                   \textless{}dbl\textgreater{}}
\CommentTok{\#\textgreater{} 1          6 0.00409              {-}0.00146                {-}0.00146   }
\CommentTok{\#\textgreater{} 2         12 0.00480              {-}0.000753               {-}0.000753  }
\CommentTok{\#\textgreater{} 3         24 0.00521              {-}0.000340               {-}0.000340  }
\CommentTok{\#\textgreater{} 4         60 0.00543              {-}0.000117               {-}0.000117  }
\CommentTok{\#\textgreater{} 5        120 0.00550              {-}0.0000508              {-}0.0000509 }
\CommentTok{\#\textgreater{} 6        240 0.00553              {-}0.0000194              {-}0.0000194 }
\CommentTok{\#\textgreater{} 7        500 0.00555              {-}0.00000390             {-}0.00000395}

\FunctionTok{print}\NormalTok{(tree\_convergence\_plot)}
\end{Highlighting}
\end{Shaded}

\pandocbounded{\includegraphics[keepaspectratio]{part2_quantive_riemann/ch07_callable_hull_white_files/figure-pdf/chapter07-tree-convergence-1.pdf}}

Tree
stepを増やすほど、Tree価格がJamshidian価格へ近づくことを確認する。Bermudan
SwaptionやCallable
Bondには一般的なJamshidian解析解を利用できないため、European
SwaptionでTree精度を検証してから、早期行使商品へ同じTree設定を適用することが重要である。

\section{7.7 まとめ}\label{ux307eux3068ux3081-11}

Callable Bondは、Bullet
Bondから発行体の期限前償還権を差し引いた商品として理解できる。金利が低下すると発行体のコール権は価値を持ち、投資家が保有するCallable
Bondの価格上昇は抑制される。

発行体のコール権は、金利低下時に価値が上昇するReceiver
Swaptionと同じ方向の金利オプション性を持つ。複数のコール日がある場合にはBermudan
Receiver
Swaptionに近い。ただし、債券とSwapでは元本、決済およびキャッシュフロー条件が完全には一致しないため、価格関係は近似として解釈する必要がある。

Hull--WhiteモデルはGaussian平均回帰型の短期金利モデルであり、時間依存ドリフトによって初期イールドカーブへ適合できる。Vasicekモデルがパラメータから期間構造を生成するのに対し、Hull--Whiteモデルは観測カーブを出発点として、その周りの確率変動を表現する。

Hull--Whiteの平均回帰速度 \(a\) は金利ショックの持続性を、\(\sigma\)
は短期金利変動の大きさを表す。本章では、1Y×5Y ATM SwaptionのNormal
volatilityを使い、\(a\)を固定して\(\sigma\)をカリブレーションした。

European Swaptionでは、Jamshidian decomposition、Hull--White
TreeおよびBachelier Engineを比較した。Tree
stepを増やすとTree価格はJamshidian価格へ近づく。Callable BondやBermudan
Swaptionのような早期行使商品では、このTree精度を確認したうえで評価する必要がある。

\begin{Shaded}
\begin{Highlighting}[]
\NormalTok{chapter07\_summary\_tbl }\OtherTok{\textless{}{-}}\NormalTok{ tibble}\SpecialCharTok{::}\FunctionTok{tribble}\NormalTok{(}
  \SpecialCharTok{\textasciitilde{}}\NormalTok{section,                  }\SpecialCharTok{\textasciitilde{}}\NormalTok{main\_result,}
  \StringTok{"Callable Bond"}\NormalTok{,           }\StringTok{"Bullet bond minus issuer call option"}\NormalTok{,}
  \StringTok{"Receiver Swaption"}\NormalTok{,       }\StringTok{"Interest{-}rate direction of the issuer call"}\NormalTok{,}
  \StringTok{"Vasicek and Hull{-}White"}\NormalTok{,  }\StringTok{"Initial{-}curve fit and mean reversion"}\NormalTok{,}
  \StringTok{"Calibration"}\NormalTok{,             }\StringTok{"Fixed a and calibrated sigma"}\NormalTok{,}
  \StringTok{"Tree convergence"}\NormalTok{,        }\StringTok{"Convergence toward Jamshidian pricing"}
\NormalTok{)}

\NormalTok{GaussQuant}\SpecialCharTok{::}\FunctionTok{show\_tbl\_GQH}\NormalTok{(}
\NormalTok{  chapter07\_summary\_tbl,}
  \StringTok{"Chapter II{-}7 summary"}\NormalTok{,}
  \AttributeTok{n =} \DecValTok{20}
\NormalTok{)}
\CommentTok{\#\textgreater{} \# A tibble: 5 x 2}
\CommentTok{\#\textgreater{}   section                main\_result                               }
\CommentTok{\#\textgreater{}   \textless{}chr\textgreater{}                  \textless{}chr\textgreater{}                                     }
\CommentTok{\#\textgreater{} 1 Callable Bond          Bullet bond minus issuer call option      }
\CommentTok{\#\textgreater{} 2 Receiver Swaption      Interest{-}rate direction of the issuer call}
\CommentTok{\#\textgreater{} 3 Vasicek and Hull{-}White Initial{-}curve fit and mean reversion      }
\CommentTok{\#\textgreater{} 4 Calibration            Fixed a and calibrated sigma              }
\CommentTok{\#\textgreater{} 5 Tree convergence       Convergence toward Jamshidian pricing}

\FunctionTok{cat}\NormalTok{(}
  \StringTok{"}\SpecialCharTok{\textbackslash{}n}\StringTok{Chapter II{-}7 Callable Bond and Hull{-}White model completed successfully.}\SpecialCharTok{\textbackslash{}n}\StringTok{"}
\NormalTok{)}
\CommentTok{\#\textgreater{} }
\CommentTok{\#\textgreater{} Chapter II{-}7 Callable Bond and Hull{-}White model completed successfully.}
\end{Highlighting}
\end{Shaded}

\chapter{第II-8章
信用リスクとCDS}\label{ux7b2cii-8ux7ae0-ux4fe1ux7528ux30eaux30b9ux30afux3068cds}

前章までの評価では、将来の市場価格、金利またはボラティリティの変動を扱った。本章では、契約相手または参照企業が債務を履行できなくなる信用リスクを検討する。

Credit Default
Swap（CDS）は、参照企業などに信用事由が発生した場合に、プロテクションの売り手が買い手へ損失補償を支払う契約である。プロテクションの買い手は、その対価として契約期間中に定期的なPremiumを支払う。

したがって、CDSは二つのレグから構成される。

\begin{itemize}
\tightlist
\item
  Premium Leg：プロテクション買い手が定期的なSpreadを支払う。
\item
  Protection
  Leg：信用事由が発生した場合に、プロテクション売り手が損失を補償する。
\end{itemize}

プロテクション買い手から見たCDS価値は、次のように表される。

\[
V_{\mathrm{CDS}}
=
V_{\mathrm{protection}}
+
V_{\mathrm{premium}}
\]

Premium Legは支払いであるため通常は負、Protection
Legは受取りであるため正となる。

CDS評価には、割引カーブだけでなく、将来のデフォルト確率を表すDefault
Probability Curveが必要である。本章では、Hazard
Rateから生存確率とデフォルト確率を求め、Premium LegとProtection
Legを評価する。

まず、フラット割引率とフラットHazard
Rateを使用して、MidPoint法によるCDS評価と手計算を比較する。次に、標準CDS2015契約をISDA
CDS Engineで評価し、Fair spread、Fair upfrontおよびAccrual
rebateを確認する。

\begin{Shaded}
\begin{Highlighting}[]
\NormalTok{knitr}\SpecialCharTok{::}\NormalTok{opts\_chunk}\SpecialCharTok{$}\FunctionTok{set}\NormalTok{(}
  \AttributeTok{collapse =} \ConstantTok{TRUE}\NormalTok{,}
  \AttributeTok{comment =} \StringTok{"\#\textgreater{}"}\NormalTok{,}
  \AttributeTok{warning =} \ConstantTok{FALSE}\NormalTok{,}
  \AttributeTok{message =} \ConstantTok{FALSE}
\NormalTok{)}


\FunctionTok{suppressPackageStartupMessages}\NormalTok{(\{}
  \FunctionTok{library}\NormalTok{(tidyverse)}
  \FunctionTok{library}\NormalTok{(ggthemes)}
  \FunctionTok{library}\NormalTok{(QuantLib)}
  \FunctionTok{library}\NormalTok{(GaussQuant)}
  \FunctionTok{library}\NormalTok{(tidyverse)}
  \FunctionTok{library}\NormalTok{(broom)}
\NormalTok{\})}
\end{Highlighting}
\end{Shaded}

\section{8.1
信用リスクとCDSの基礎}\label{ux4fe1ux7528ux30eaux30b9ux30afux3068cdsux306eux57faux790e}

\subsection{8.1.1
デフォルト時刻と生存確率}\label{ux30c7ux30d5ux30a9ux30ebux30c8ux6642ux523bux3068ux751fux5b58ux78baux7387}

デフォルトが発生する確率時刻を \(\tau\) とする。時点 \(t\)
まで企業がデフォルトせずに存続する確率を、生存確率と呼ぶ。

\[
Q(t)
=
P(\tau>t)
\]

時点0から \(t\) までの累積デフォルト確率は、

\[
P(\tau\leq t)
=
1-Q(t)
\]

である。

Hazard Rate \(\lambda(t)\) は、時点 \(t\)
まで存続しているという条件の下で、その直後にデフォルトする瞬間的な強度を表す。

\[
\lambda(t)
=
\lim_{\Delta t\to 0}
\frac{
P(t<\tau\leq t+\Delta t\mid \tau>t)
}{
\Delta t
}
\]

Hazard Rateが一定のフラットモデルでは、生存確率は次式となる。

\[
Q(t)
=
\exp(-\lambda t)
\]

累積デフォルト確率は、

\[
1-Q(t)
=
1-\exp(-\lambda t)
\]

である。

Hazard
Rateは、通常の意味での「一年間のデフォルト確率」そのものではない。たとえばHazard
Rateが2\%であっても、1年累積デフォルト確率は、

\[
1-e^{-0.02}
\]

であり、厳密には2\%と一致しない。

\subsection{8.1.2 Recovery RateとLoss Given
Default}\label{recovery-rateux3068loss-given-default}

デフォルトが発生した場合でも、債権価値のすべてが失われるとは限らない。デフォルト後に回収できる元本比率をRecovery
Rate \(R\) とする。

損失率Loss Given Defaultは、

\[
LGD
=
1-R
\]

である。

Notionalを \(N\) とすると、Protection Legが補償する損失額は概ね、

\[
N(1-R)
\]

となる。

Recovery
Rateが高いほどデフォルト時の損失は小さくなり、同じデフォルト確率であればCDS
Spreadは低くなる。

\subsection{8.1.3 CDS SpreadとHazard
Rateの近似}\label{cds-spreadux3068hazard-rateux306eux8fd1ux4f3c}

連続的なPremium支払い、一定Hazard Rate、一定Recovery
Rateおよび単純な割引を仮定すると、CDS Spread \(s\) とHazard Rate
\(\lambda\) には次の近似関係がある。

\[
s
\approx
(1-R)\lambda
\]

したがって、

\[
\lambda
\approx
\frac{s}{1-R}
\]

となる。

この近似はCDSの信用水準を理解するために有用である。ただし、実際の評価では、四半期ごとの支払い、Accrual
on default、割引、標準日付およびカーブ補間が影響するため、厳密なHazard
Rateとは一致しない。

\section{8.2
フラットカーブによるCDS評価}\label{ux30d5ux30e9ux30c3ux30c8ux30abux30fcux30d6ux306bux3088ux308bcdsux8a55ux4fa1}

\subsection{8.2.1
フラットカーブによる1年物CDS}\label{ux30d5ux30e9ux30c3ux30c8ux30abux30fcux30d6ux306bux3088ux308b1ux5e74ux7269cds}

1年物CDSを考える。

\begin{itemize}
\tightlist
\item
  評価日：2022年9月19日
\item
  満期：2023年9月20日
\item
  Notional：10,000,000
\item
  割引率：10\%
\item
  市場Spread：130bp
\item
  契約Coupon：100bp
\item
  Recovery Rate：40\%
\end{itemize}

市場Spreadは130bpであるが、契約Couponは100bpである。このため、契約時点のCDS価値はゼロにならない。

\begin{Shaded}
\begin{Highlighting}[]
\NormalTok{trade\_date\_flat }\OtherTok{\textless{}{-}}\NormalTok{ GaussQuant}\SpecialCharTok{::}\FunctionTok{date\_GQL}\NormalTok{(}\StringTok{"2022{-}09{-}19"}\NormalTok{)}
\NormalTok{maturity\_date\_flat }\OtherTok{\textless{}{-}}\NormalTok{ GaussQuant}\SpecialCharTok{::}\FunctionTok{date\_GQL}\NormalTok{(}\StringTok{"2023{-}09{-}20"}\NormalTok{)}

\NormalTok{GaussQuant}\SpecialCharTok{::}\FunctionTok{set\_eval\_date\_GQL}\NormalTok{(trade\_date\_flat)}

\NormalTok{notional\_flat }\OtherTok{\textless{}{-}} \DecValTok{10000000}
\NormalTok{risk\_free\_rate\_flat }\OtherTok{\textless{}{-}} \FloatTok{0.10}
\NormalTok{market\_spread\_flat }\OtherTok{\textless{}{-}} \DecValTok{130} \SpecialCharTok{/} \DecValTok{10000}
\NormalTok{contract\_coupon\_flat }\OtherTok{\textless{}{-}} \DecValTok{100} \SpecialCharTok{/} \DecValTok{10000}
\NormalTok{recovery\_rate\_flat }\OtherTok{\textless{}{-}} \FloatTok{0.40}

\NormalTok{actual\_365 }\OtherTok{\textless{}{-}}\NormalTok{ QuantLib}\SpecialCharTok{::}\FunctionTok{Actual365Fixed}\NormalTok{()}
\NormalTok{actual\_360 }\OtherTok{\textless{}{-}}\NormalTok{ QuantLib}\SpecialCharTok{::}\FunctionTok{Actual360}\NormalTok{()}
\NormalTok{weekend\_calendar }\OtherTok{\textless{}{-}}\NormalTok{ QuantLib}\SpecialCharTok{::}\FunctionTok{WeekendsOnly}\NormalTok{()}

\NormalTok{flat\_cds\_input\_tbl }\OtherTok{\textless{}{-}}\NormalTok{ tibble}\SpecialCharTok{::}\FunctionTok{tribble}\NormalTok{(}
  \SpecialCharTok{\textasciitilde{}}\NormalTok{input,             }\SpecialCharTok{\textasciitilde{}}\NormalTok{value,}
  \StringTok{"trade\_date"}\NormalTok{,        }\StringTok{"2022{-}09{-}19"}\NormalTok{,}
  \StringTok{"maturity\_date"}\NormalTok{,     }\StringTok{"2023{-}09{-}20"}\NormalTok{,}
  \StringTok{"notional"}\NormalTok{,          }\FunctionTok{as.character}\NormalTok{(notional\_flat),}
  \StringTok{"risk\_free\_rate"}\NormalTok{,    }\FunctionTok{as.character}\NormalTok{(risk\_free\_rate\_flat),}
  \StringTok{"market\_spread"}\NormalTok{,     }\FunctionTok{as.character}\NormalTok{(market\_spread\_flat),}
  \StringTok{"contract\_coupon"}\NormalTok{,   }\FunctionTok{as.character}\NormalTok{(contract\_coupon\_flat),}
  \StringTok{"recovery\_rate"}\NormalTok{,     }\FunctionTok{as.character}\NormalTok{(recovery\_rate\_flat)}
\NormalTok{)}

\NormalTok{GaussQuant}\SpecialCharTok{::}\FunctionTok{show\_tbl\_GQH}\NormalTok{(}
\NormalTok{  flat\_cds\_input\_tbl,}
  \StringTok{"Flat CDS input assumptions"}\NormalTok{,}
  \AttributeTok{n =} \DecValTok{20}
\NormalTok{)}
\CommentTok{\#\textgreater{} \# A tibble: 7 x 2}
\CommentTok{\#\textgreater{}   input           value     }
\CommentTok{\#\textgreater{}   \textless{}chr\textgreater{}           \textless{}chr\textgreater{}     }
\CommentTok{\#\textgreater{} 1 trade\_date      2022{-}09{-}19}
\CommentTok{\#\textgreater{} 2 maturity\_date   2023{-}09{-}20}
\CommentTok{\#\textgreater{} 3 notional        1e+07     }
\CommentTok{\#\textgreater{} 4 risk\_free\_rate  0.1       }
\CommentTok{\#\textgreater{} 5 market\_spread   0.013     }
\CommentTok{\#\textgreater{} 6 contract\_coupon 0.01      }
\CommentTok{\#\textgreater{} 7 recovery\_rate   0.4}
\end{Highlighting}
\end{Shaded}

\subsection{8.2.2 Hazard
Rate近似と確率曲線}\label{hazard-rateux8fd1ux4f3cux3068ux78baux7387ux66f2ux7dda}

市場Spread 130bpとRecovery Rate 40\%から、フラットHazard
Rateを近似する。

\begin{Shaded}
\begin{Highlighting}[]
\NormalTok{hazard\_rate\_approximation }\OtherTok{\textless{}{-}}\NormalTok{ market\_spread\_flat }\SpecialCharTok{/}
\NormalTok{  (}\DecValTok{1} \SpecialCharTok{{-}}\NormalTok{ recovery\_rate\_flat)}

\NormalTok{hazard\_approximation\_tbl }\OtherTok{\textless{}{-}}\NormalTok{ tibble}\SpecialCharTok{::}\FunctionTok{tribble}\NormalTok{(}
  \SpecialCharTok{\textasciitilde{}}\NormalTok{metric,                 }\SpecialCharTok{\textasciitilde{}}\NormalTok{value,}
  \StringTok{"market\_spread"}\NormalTok{,          market\_spread\_flat,}
  \StringTok{"recovery\_rate"}\NormalTok{,          recovery\_rate\_flat,}
  \StringTok{"loss\_given\_default"}\NormalTok{,     }\DecValTok{1} \SpecialCharTok{{-}}\NormalTok{ recovery\_rate\_flat,}
  \StringTok{"hazard\_rate\_approx"}\NormalTok{,     hazard\_rate\_approximation}
\NormalTok{)}

\NormalTok{probability\_time\_grid }\OtherTok{\textless{}{-}} \FunctionTok{seq}\NormalTok{(}
  \DecValTok{0}\NormalTok{,}
  \DecValTok{5}\NormalTok{,}
  \AttributeTok{by =} \FloatTok{0.25}
\NormalTok{)}

\NormalTok{flat\_probability\_tbl }\OtherTok{\textless{}{-}}\NormalTok{ tibble}\SpecialCharTok{::}\FunctionTok{tibble}\NormalTok{(}
  \AttributeTok{time =}\NormalTok{ probability\_time\_grid}
\NormalTok{) }\SpecialCharTok{|\textgreater{}}
\NormalTok{  dplyr}\SpecialCharTok{::}\FunctionTok{mutate}\NormalTok{(}
    \AttributeTok{survival\_probability =} \FunctionTok{exp}\NormalTok{(}
      \SpecialCharTok{{-}}\NormalTok{hazard\_rate\_approximation }\SpecialCharTok{*}\NormalTok{ .data}\SpecialCharTok{$}\NormalTok{time}
\NormalTok{    ),}
    \AttributeTok{cumulative\_default\_probability =} \DecValTok{1} \SpecialCharTok{{-}}
\NormalTok{      .data}\SpecialCharTok{$}\NormalTok{survival\_probability}
\NormalTok{  )}

\NormalTok{flat\_probability\_plot\_tbl }\OtherTok{\textless{}{-}}\NormalTok{ flat\_probability\_tbl }\SpecialCharTok{|\textgreater{}}
\NormalTok{  tidyr}\SpecialCharTok{::}\FunctionTok{pivot\_longer}\NormalTok{(}
    \AttributeTok{cols =} \FunctionTok{c}\NormalTok{(}
\NormalTok{      survival\_probability,}
\NormalTok{      cumulative\_default\_probability}
\NormalTok{    ),}
    \AttributeTok{names\_to =} \StringTok{"probability\_type"}\NormalTok{,}
    \AttributeTok{values\_to =} \StringTok{"probability"}
\NormalTok{  )}

\NormalTok{flat\_probability\_plot }\OtherTok{\textless{}{-}}\NormalTok{ ggplot2}\SpecialCharTok{::}\FunctionTok{ggplot}\NormalTok{(}
\NormalTok{  flat\_probability\_plot\_tbl,}
\NormalTok{  ggplot2}\SpecialCharTok{::}\FunctionTok{aes}\NormalTok{(}
    \AttributeTok{x =}\NormalTok{ .data}\SpecialCharTok{$}\NormalTok{time,}
    \AttributeTok{y =}\NormalTok{ .data}\SpecialCharTok{$}\NormalTok{probability,}
    \AttributeTok{linetype =}\NormalTok{ .data}\SpecialCharTok{$}\NormalTok{probability\_type}
\NormalTok{  )}
\NormalTok{) }\SpecialCharTok{+}
\NormalTok{  ggplot2}\SpecialCharTok{::}\FunctionTok{geom\_line}\NormalTok{(}
    \AttributeTok{linewidth =} \FloatTok{0.7}
\NormalTok{  ) }\SpecialCharTok{+}
\NormalTok{  ggplot2}\SpecialCharTok{::}\FunctionTok{scale\_y\_continuous}\NormalTok{(}
    \AttributeTok{labels =}\NormalTok{ scales}\SpecialCharTok{::}\FunctionTok{percent\_format}\NormalTok{(}
      \AttributeTok{accuracy =} \FloatTok{0.1}
\NormalTok{    )}
\NormalTok{  ) }\SpecialCharTok{+}
\NormalTok{  ggplot2}\SpecialCharTok{::}\FunctionTok{labs}\NormalTok{(}
    \AttributeTok{title =} \StringTok{"Flat{-}hazard survival and default probabilities"}\NormalTok{,}
    \AttributeTok{x =} \StringTok{"Time (years)"}\NormalTok{,}
    \AttributeTok{y =} \StringTok{"Probability"}\NormalTok{,}
    \AttributeTok{linetype =} \StringTok{"Probability"}
\NormalTok{  ) }\SpecialCharTok{+}
\NormalTok{  ggplot2}\SpecialCharTok{::}\FunctionTok{theme\_minimal}\NormalTok{()}

\NormalTok{GaussQuant}\SpecialCharTok{::}\FunctionTok{show\_tbl\_GQH}\NormalTok{(}
\NormalTok{  hazard\_approximation\_tbl,}
  \StringTok{"Flat hazard{-}rate approximation"}\NormalTok{,}
  \AttributeTok{n =} \DecValTok{20}
\NormalTok{)}
\CommentTok{\#\textgreater{} \# A tibble: 4 x 2}
\CommentTok{\#\textgreater{}   metric              value}
\CommentTok{\#\textgreater{}   \textless{}chr\textgreater{}               \textless{}dbl\textgreater{}}
\CommentTok{\#\textgreater{} 1 market\_spread      0.013 }
\CommentTok{\#\textgreater{} 2 recovery\_rate      0.4   }
\CommentTok{\#\textgreater{} 3 loss\_given\_default 0.6   }
\CommentTok{\#\textgreater{} 4 hazard\_rate\_approx 0.0217}

\NormalTok{GaussQuant}\SpecialCharTok{::}\FunctionTok{show\_tbl\_GQH}\NormalTok{(}
\NormalTok{  flat\_probability\_tbl,}
  \StringTok{"Flat{-}hazard probability curve"}\NormalTok{,}
  \AttributeTok{n =} \DecValTok{30}
\NormalTok{)}
\CommentTok{\#\textgreater{} \# A tibble: 21 x 3}
\CommentTok{\#\textgreater{}     time survival\_probability cumulative\_default\_probability}
\CommentTok{\#\textgreater{}    \textless{}dbl\textgreater{}                \textless{}dbl\textgreater{}                          \textless{}dbl\textgreater{}}
\CommentTok{\#\textgreater{}  1  0                   1                            0      }
\CommentTok{\#\textgreater{}  2  0.25                0.995                        0.00540}
\CommentTok{\#\textgreater{}  3  0.5                 0.989                        0.0108 }
\CommentTok{\#\textgreater{}  4  0.75                0.984                        0.0161 }
\CommentTok{\#\textgreater{}  5  1                   0.979                        0.0214 }
\CommentTok{\#\textgreater{}  6  1.25                0.973                        0.0267 }
\CommentTok{\#\textgreater{}  7  1.5                 0.968                        0.0320 }
\CommentTok{\#\textgreater{}  8  1.75                0.963                        0.0372 }
\CommentTok{\#\textgreater{}  9  2                   0.958                        0.0424 }
\CommentTok{\#\textgreater{} 10  2.25                0.952                        0.0476 }
\CommentTok{\#\textgreater{} \# i 11 more rows}

\FunctionTok{print}\NormalTok{(flat\_probability\_plot)}
\end{Highlighting}
\end{Shaded}

\pandocbounded{\includegraphics[keepaspectratio]{part2_quantive_riemann/ch08_credit_risk_cds_files/figure-pdf/chapter08-flat-hazard-1.pdf}}

\subsection{8.2.3 CDS
Scheduleとカーブ}\label{cds-scheduleux3068ux30abux30fcux30d6}

Premiumは四半期ごとに支払われるものとする。単純な教育用例として、評価日の翌日から満期までのForward
Scheduleを作成する。

\begin{Shaded}
\begin{Highlighting}[]
\NormalTok{cds\_schedule\_flat }\OtherTok{\textless{}{-}}\NormalTok{ QuantLib}\SpecialCharTok{::}\FunctionTok{Schedule}\NormalTok{(}
\NormalTok{  trade\_date\_flat }\SpecialCharTok{+} \DecValTok{1}\NormalTok{L,}
\NormalTok{  maturity\_date\_flat,}
\NormalTok{  GaussQuant}\SpecialCharTok{::}\FunctionTok{period\_months\_GQL}\NormalTok{(}\DecValTok{3}\NormalTok{),}
\NormalTok{  weekend\_calendar,}
  \StringTok{"Following"}\NormalTok{,}
  \StringTok{"Unadjusted"}\NormalTok{,}
  \StringTok{"Forward"}\NormalTok{,}
  \ConstantTok{FALSE}
\NormalTok{)}

\NormalTok{discount\_curve\_flat }\OtherTok{\textless{}{-}}\NormalTok{ GaussQuant}\SpecialCharTok{::}\FunctionTok{flat\_curve\_GQL}\NormalTok{(}
  \AttributeTok{reference\_date =}\NormalTok{ trade\_date\_flat,}
  \AttributeTok{rate =}\NormalTok{ risk\_free\_rate\_flat,}
  \AttributeTok{day\_counter =}\NormalTok{ actual\_365,}
  \AttributeTok{compounding =} \StringTok{"Continuous"}
\NormalTok{)}

\NormalTok{discount\_curve\_handle\_flat }\OtherTok{\textless{}{-}}\NormalTok{ QuantLib}\SpecialCharTok{::}\FunctionTok{YieldTermStructureHandle}\NormalTok{(}
\NormalTok{  discount\_curve\_flat}
\NormalTok{)}

\NormalTok{hazard\_curve\_flat }\OtherTok{\textless{}{-}}\NormalTok{ QuantLib}\SpecialCharTok{::}\FunctionTok{FlatHazardRate}\NormalTok{(}
\NormalTok{  trade\_date\_flat,}
\NormalTok{  GaussQuant}\SpecialCharTok{::}\FunctionTok{quote\_handle\_GQL}\NormalTok{(}
\NormalTok{    hazard\_rate\_approximation}
\NormalTok{  ),}
\NormalTok{  actual\_365}
\NormalTok{)}

\NormalTok{hazard\_curve\_handle\_flat }\OtherTok{\textless{}{-}}\NormalTok{ QuantLib}\SpecialCharTok{::}\FunctionTok{DefaultProbabilityTermStructureHandle}\NormalTok{(}
\NormalTok{  hazard\_curve\_flat}
\NormalTok{)}

\NormalTok{flat\_curve\_check\_tbl }\OtherTok{\textless{}{-}}\NormalTok{ tibble}\SpecialCharTok{::}\FunctionTok{tribble}\NormalTok{(}
  \SpecialCharTok{\textasciitilde{}}\NormalTok{metric,                            }\SpecialCharTok{\textasciitilde{}}\NormalTok{value,}
  \StringTok{"discount\_factor\_at\_maturity"}\NormalTok{,       discount\_curve\_flat}\SpecialCharTok{$}\FunctionTok{discount}\NormalTok{(}
\NormalTok{                                          maturity\_date\_flat}
\NormalTok{                                        ),}
  \StringTok{"survival\_probability\_at\_maturity"}\NormalTok{, hazard\_curve\_flat}\SpecialCharTok{$}\FunctionTok{survivalProbability}\NormalTok{(}
\NormalTok{                                          maturity\_date\_flat}
\NormalTok{                                        ),}
  \StringTok{"default\_probability\_to\_maturity"}\NormalTok{,  }\DecValTok{1} \SpecialCharTok{{-}}
\NormalTok{                                        hazard\_curve\_flat}\SpecialCharTok{$}\FunctionTok{survivalProbability}\NormalTok{(}
\NormalTok{                                          maturity\_date\_flat}
\NormalTok{                                        ),}
  \StringTok{"hazard\_rate\_at\_maturity"}\NormalTok{,          hazard\_curve\_flat}\SpecialCharTok{$}\FunctionTok{hazardRate}\NormalTok{(}
\NormalTok{                                          maturity\_date\_flat}
\NormalTok{                                        )}
\NormalTok{)}

\NormalTok{GaussQuant}\SpecialCharTok{::}\FunctionTok{show\_tbl\_GQH}\NormalTok{(}
\NormalTok{  GaussQuant}\SpecialCharTok{::}\FunctionTok{schedule\_table\_GQL}\NormalTok{(}
\NormalTok{    cds\_schedule\_flat}
\NormalTok{  ),}
  \StringTok{"CDS schedule: flat{-}curve example"}\NormalTok{,}
  \AttributeTok{n =} \DecValTok{20}
\NormalTok{)}
\CommentTok{\#\textgreater{} \# A tibble: 5 x 2}
\CommentTok{\#\textgreater{}   period schedule\_date}
\CommentTok{\#\textgreater{}    \textless{}int\textgreater{} \textless{}date\textgreater{}       }
\CommentTok{\#\textgreater{} 1      1 2022{-}09{-}20   }
\CommentTok{\#\textgreater{} 2      2 2022{-}12{-}20   }
\CommentTok{\#\textgreater{} 3      3 2023{-}03{-}20   }
\CommentTok{\#\textgreater{} 4      4 2023{-}06{-}20   }
\CommentTok{\#\textgreater{} 5      5 2023{-}09{-}20}

\NormalTok{GaussQuant}\SpecialCharTok{::}\FunctionTok{show\_tbl\_GQH}\NormalTok{(}
\NormalTok{  flat\_curve\_check\_tbl,}
  \StringTok{"Flat discount and hazard{-}curve check"}\NormalTok{,}
  \AttributeTok{n =} \DecValTok{20}
\NormalTok{)}
\CommentTok{\#\textgreater{} \# A tibble: 4 x 2}
\CommentTok{\#\textgreater{}   metric                            value}
\CommentTok{\#\textgreater{}   \textless{}chr\textgreater{}                             \textless{}dbl\textgreater{}}
\CommentTok{\#\textgreater{} 1 discount\_factor\_at\_maturity      0.905 }
\CommentTok{\#\textgreater{} 2 survival\_probability\_at\_maturity 0.979 }
\CommentTok{\#\textgreater{} 3 default\_probability\_to\_maturity  0.0215}
\CommentTok{\#\textgreater{} 4 hazard\_rate\_at\_maturity          0.0217}
\end{Highlighting}
\end{Shaded}

\subsection{8.2.4 MidPoint CDS Engine}\label{midpoint-cds-engine}

MidPoint法では、各Premium期間中に発生するデフォルトを、その期間の中点で発生したものとして近似する。

100bp CouponのCDSを作成し、GaussQuantのMidPoint CDS Engineで評価する。

\begin{Shaded}
\begin{Highlighting}[]
\NormalTok{cds\_flat }\OtherTok{\textless{}{-}}\NormalTok{ QuantLib}\SpecialCharTok{::}\FunctionTok{CreditDefaultSwap}\NormalTok{(}
\NormalTok{  GaussQuant}\SpecialCharTok{::}\FunctionTok{cds\_buyer\_side\_GQL}\NormalTok{(),}
\NormalTok{  notional\_flat,}
\NormalTok{  contract\_coupon\_flat,}
\NormalTok{  cds\_schedule\_flat,}
  \StringTok{"Following"}\NormalTok{,}
\NormalTok{  actual\_360}
\NormalTok{)}

\NormalTok{midpoint\_engine\_flat }\OtherTok{\textless{}{-}}\NormalTok{ GaussQuant}\SpecialCharTok{::}\FunctionTok{cds\_midpoint\_engine\_GQL}\NormalTok{(}
  \AttributeTok{probability\_handle =}\NormalTok{ hazard\_curve\_handle\_flat,}
  \AttributeTok{recovery\_rate =}\NormalTok{ recovery\_rate\_flat,}
  \AttributeTok{discount\_handle =}\NormalTok{ discount\_curve\_handle\_flat}
\NormalTok{)}

\NormalTok{cds\_flat}\SpecialCharTok{$}\FunctionTok{setPricingEngine}\NormalTok{(}
\NormalTok{  midpoint\_engine\_flat}
\NormalTok{)}
\CommentTok{\#\textgreater{} NULL}

\NormalTok{cds\_flat\_summary\_tbl }\OtherTok{\textless{}{-}}\NormalTok{ GaussQuant}\SpecialCharTok{::}\FunctionTok{cds\_summary\_GQL}\NormalTok{(}
\NormalTok{  cds\_flat}
\NormalTok{)}

\NormalTok{cds\_flat\_leg\_check\_tbl }\OtherTok{\textless{}{-}}\NormalTok{ tibble}\SpecialCharTok{::}\FunctionTok{tribble}\NormalTok{(}
  \SpecialCharTok{\textasciitilde{}}\NormalTok{metric,             }\SpecialCharTok{\textasciitilde{}}\NormalTok{value,}
  \StringTok{"coupon\_leg\_npv"}\NormalTok{,     cds\_flat}\SpecialCharTok{$}\FunctionTok{couponLegNPV}\NormalTok{(),}
  \StringTok{"default\_leg\_npv"}\NormalTok{,    cds\_flat}\SpecialCharTok{$}\FunctionTok{defaultLegNPV}\NormalTok{(),}
  \StringTok{"sum\_of\_two\_legs"}\NormalTok{,    cds\_flat}\SpecialCharTok{$}\FunctionTok{couponLegNPV}\NormalTok{() }\SpecialCharTok{+}
\NormalTok{                          cds\_flat}\SpecialCharTok{$}\FunctionTok{defaultLegNPV}\NormalTok{(),}
  \StringTok{"reported\_npv"}\NormalTok{,       GaussQuant}\SpecialCharTok{::}\FunctionTok{cds\_npv\_GQL}\NormalTok{(}
\NormalTok{                          cds\_flat}
\NormalTok{                        ),}
  \StringTok{"difference"}\NormalTok{,         GaussQuant}\SpecialCharTok{::}\FunctionTok{cds\_npv\_GQL}\NormalTok{(}
\NormalTok{                          cds\_flat}
\NormalTok{                        ) }\SpecialCharTok{{-}}
\NormalTok{                          (}
\NormalTok{                            cds\_flat}\SpecialCharTok{$}\FunctionTok{couponLegNPV}\NormalTok{() }\SpecialCharTok{+}
\NormalTok{                              cds\_flat}\SpecialCharTok{$}\FunctionTok{defaultLegNPV}\NormalTok{()}
\NormalTok{                          )}
\NormalTok{)}

\NormalTok{GaussQuant}\SpecialCharTok{::}\FunctionTok{show\_tbl\_GQH}\NormalTok{(}
\NormalTok{  cds\_flat\_summary\_tbl,}
  \StringTok{"CDS summary: MidPoint engine with flat curves"}\NormalTok{,}
  \AttributeTok{n =} \DecValTok{20}
\NormalTok{)}
\CommentTok{\#\textgreater{} \# A tibble: 4 x 2}
\CommentTok{\#\textgreater{}   metric                value}
\CommentTok{\#\textgreater{}   \textless{}chr\textgreater{}                 \textless{}dbl\textgreater{}}
\CommentTok{\#\textgreater{} 1 npv              28116.    }
\CommentTok{\#\textgreater{} 2 fair\_spread          0.0130}
\CommentTok{\#\textgreater{} 3 coupon\_leg\_npv  {-}94246.    }
\CommentTok{\#\textgreater{} 4 default\_leg\_npv 122362.}

\NormalTok{GaussQuant}\SpecialCharTok{::}\FunctionTok{show\_tbl\_GQH}\NormalTok{(}
\NormalTok{  cds\_flat\_leg\_check\_tbl,}
  \StringTok{"CDS two{-}leg NPV check"}\NormalTok{,}
  \AttributeTok{n =} \DecValTok{20}
\NormalTok{)}
\CommentTok{\#\textgreater{} \# A tibble: 5 x 2}
\CommentTok{\#\textgreater{}   metric            value}
\CommentTok{\#\textgreater{}   \textless{}chr\textgreater{}             \textless{}dbl\textgreater{}}
\CommentTok{\#\textgreater{} 1 coupon\_leg\_npv  {-}94246.}
\CommentTok{\#\textgreater{} 2 default\_leg\_npv 122362.}
\CommentTok{\#\textgreater{} 3 sum\_of\_two\_legs  28116.}
\CommentTok{\#\textgreater{} 4 reported\_npv     28116.}
\CommentTok{\#\textgreater{} 5 difference           0}
\end{Highlighting}
\end{Shaded}

プロテクション買い手から見て、Premium Legは負、Protection
Legは正となる。契約CouponがFair
Spreadより低ければ、プロテクション買い手に有利となり、CDS
NPVは一般に正となる。

\subsection{8.2.5
CDSキャッシュフローと確率表}\label{cdsux30adux30e3ux30c3ux30b7ux30e5ux30d5ux30edux30fcux3068ux78baux7387ux8868}

GaussQuantの\texttt{cds\_cashflow\_table\_GQL()}を利用し、各Premium期間の割引係数、生存確率、デフォルト確率および期間中点の情報を確認する。

\begin{Shaded}
\begin{Highlighting}[]
\NormalTok{cds\_cashflow\_flat\_tbl }\OtherTok{\textless{}{-}}\NormalTok{ GaussQuant}\SpecialCharTok{::}\FunctionTok{cds\_cashflow\_table\_GQL}\NormalTok{(}
  \AttributeTok{cds\_schedule =}\NormalTok{ cds\_schedule\_flat,}
  \AttributeTok{protection\_start\_date =}\NormalTok{ cds\_flat}\SpecialCharTok{$}\FunctionTok{protectionStartDate}\NormalTok{(),}
  \AttributeTok{coupon\_rate =}\NormalTok{ contract\_coupon\_flat,}
  \AttributeTok{hazard\_curve =}\NormalTok{ hazard\_curve\_flat,}
  \AttributeTok{discount\_curve =}\NormalTok{ discount\_curve\_flat,}
  \AttributeTok{trade\_date =}\NormalTok{ trade\_date\_flat,}
  \AttributeTok{notional =}\NormalTok{ notional\_flat}
\NormalTok{)}

\NormalTok{GaussQuant}\SpecialCharTok{::}\FunctionTok{show\_tbl\_GQH}\NormalTok{(}
\NormalTok{  cds\_cashflow\_flat\_tbl,}
  \StringTok{"CDS cash{-}flow and probability table"}\NormalTok{,}
  \AttributeTok{n =} \DecValTok{20}
\NormalTok{)}
\CommentTok{\#\textgreater{} \# A tibble: 5 x 13}
\CommentTok{\#\textgreater{}   payment\_date coupon\_rate accrual\_start accrual\_end accrual\_days}
\CommentTok{\#\textgreater{}   \textless{}chr\textgreater{}              \textless{}dbl\textgreater{} \textless{}chr\textgreater{}         \textless{}chr\textgreater{}              \textless{}dbl\textgreater{}}
\CommentTok{\#\textgreater{} 1 2022{-}09{-}20         NA    \textless{}NA\textgreater{}          2022{-}09{-}20            NA}
\CommentTok{\#\textgreater{} 2 2022{-}12{-}20          0.01 2022{-}09{-}20    2022{-}12{-}20            91}
\CommentTok{\#\textgreater{} 3 2023{-}03{-}20          0.01 2022{-}12{-}20    2023{-}03{-}20            90}
\CommentTok{\#\textgreater{} 4 2023{-}06{-}20          0.01 2023{-}03{-}20    2023{-}06{-}20            92}
\CommentTok{\#\textgreater{} 5 2023{-}09{-}20          0.01 2023{-}06{-}20    2023{-}09{-}20            92}
\CommentTok{\#\textgreater{} \# i 8 more variables: accrual\_year\_fraction\_365 \textless{}dbl\textgreater{}, coupon\_amount \textless{}dbl\textgreater{},}
\CommentTok{\#\textgreater{} \#   discount\_factor \textless{}dbl\textgreater{}, survival\_probability \textless{}dbl\textgreater{},}
\CommentTok{\#\textgreater{} \#   default\_probability \textless{}dbl\textgreater{}, midpoint\_date \textless{}chr\textgreater{},}
\CommentTok{\#\textgreater{} \#   midpoint\_discount\_factor \textless{}dbl\textgreater{}, midpoint\_survival\_probability \textless{}dbl\textgreater{}}
\end{Highlighting}
\end{Shaded}

\section{8.3 Premium LegとProtection
Leg}\label{premium-legux3068protection-leg}

\subsection{8.3.1 Premium
Legの手計算}\label{premium-legux306eux624bux8a08ux7b97}

Premium Legは、参照企業が各支払日まで存続した場合に支払われる。

単純化した現在価値は次式で表される。

\[
V_{\mathrm{premium}}
\approx
-Ns
\sum_i
\alpha_i
D(0,t_i)
Q(t_i)
\]

ここで、\(N\)はNotional、\(s\)は契約Coupon、\(\alpha_i\)は利息計算期間、\(D(0,t_i)\)は割引係数、\(Q(t_i)\)は生存確率である。

\begin{Shaded}
\begin{Highlighting}[]
\NormalTok{premium\_leg\_hand\_flat }\OtherTok{\textless{}{-}} \SpecialCharTok{{-}}\FunctionTok{sum}\NormalTok{(}
\NormalTok{  cds\_cashflow\_flat\_tbl}\SpecialCharTok{$}\NormalTok{coupon\_amount }\SpecialCharTok{*}
\NormalTok{    cds\_cashflow\_flat\_tbl}\SpecialCharTok{$}\NormalTok{discount\_factor }\SpecialCharTok{*}
\NormalTok{    cds\_cashflow\_flat\_tbl}\SpecialCharTok{$}\NormalTok{midpoint\_survival\_probability,}
  \AttributeTok{na.rm =} \ConstantTok{TRUE}
\NormalTok{)}
\end{Highlighting}
\end{Shaded}

この近似では期間中点の生存確率を使用している。実際のCDSでは、期間中にデフォルトした場合の経過PremiumであるAccrual
on defaultも考慮される。

\subsection{8.3.2 Protection
Legの手計算}\label{protection-legux306eux624bux8a08ux7b97}

期間 \((t_{i-1},t_i]\) にデフォルトする確率を、

\[
\Delta PD_i
=
Q(t_{i-1})-Q(t_i)
\]

とする。

Protection Legは、各期間の中点 \(m_i\)
で損失補償が支払われると近似する。

\[
V_{\mathrm{protection}}
\approx
N(1-R)
\sum_i
D(0,m_i)
\Delta PD_i
\]

\begin{Shaded}
\begin{Highlighting}[]
\NormalTok{protection\_leg\_per\_unit\_hand\_flat }\OtherTok{\textless{}{-}}\NormalTok{ (}
  \DecValTok{1} \SpecialCharTok{{-}}\NormalTok{ recovery\_rate\_flat}
\NormalTok{) }\SpecialCharTok{*}
  \FunctionTok{sum}\NormalTok{(}
\NormalTok{    cds\_cashflow\_flat\_tbl}\SpecialCharTok{$}\NormalTok{midpoint\_discount\_factor }\SpecialCharTok{*}
\NormalTok{      cds\_cashflow\_flat\_tbl}\SpecialCharTok{$}\NormalTok{default\_probability,}
    \AttributeTok{na.rm =} \ConstantTok{TRUE}
\NormalTok{  )}

\NormalTok{protection\_leg\_hand\_flat }\OtherTok{\textless{}{-}}\NormalTok{ notional\_flat }\SpecialCharTok{*}
\NormalTok{  protection\_leg\_per\_unit\_hand\_flat}
\end{Highlighting}
\end{Shaded}

\subsection{8.3.3 Risky PV01とFair
Spread}\label{risky-pv01ux3068fair-spread}

Risky PV01は、CDS Spreadを1単位支払うPremium Legの現在価値に相当する。

\[
RPV01
\approx
\sum_i
\alpha_i
D(0,t_i)
Q(t_i)
\]

Fair Spreadは、Premium LegとProtection
Legの価値を等しくし、UpfrontなしのCDS NPVをゼロにするRunning
Spreadである。

\[
s_{\mathrm{fair}}
=
\frac{
V_{\mathrm{protection}}/N
}{
RPV01
}
\]

\begin{Shaded}
\begin{Highlighting}[]
\NormalTok{risky\_pv01\_hand\_flat }\OtherTok{\textless{}{-}} \FunctionTok{sum}\NormalTok{(}
\NormalTok{  (}
\NormalTok{    cds\_cashflow\_flat\_tbl}\SpecialCharTok{$}\NormalTok{accrual\_days }\SpecialCharTok{/}
      \DecValTok{360}
\NormalTok{  ) }\SpecialCharTok{*}
\NormalTok{    cds\_cashflow\_flat\_tbl}\SpecialCharTok{$}\NormalTok{discount\_factor }\SpecialCharTok{*}
\NormalTok{    cds\_cashflow\_flat\_tbl}\SpecialCharTok{$}\NormalTok{midpoint\_survival\_probability,}
  \AttributeTok{na.rm =} \ConstantTok{TRUE}
\NormalTok{)}

\NormalTok{fair\_spread\_hand\_flat }\OtherTok{\textless{}{-}}\NormalTok{ protection\_leg\_per\_unit\_hand\_flat }\SpecialCharTok{/}
\NormalTok{  risky\_pv01\_hand\_flat}

\NormalTok{cds\_npv\_hand\_flat }\OtherTok{\textless{}{-}}\NormalTok{ premium\_leg\_hand\_flat }\SpecialCharTok{+}
\NormalTok{  protection\_leg\_hand\_flat}

\NormalTok{flat\_hand\_check\_tbl }\OtherTok{\textless{}{-}}\NormalTok{ tibble}\SpecialCharTok{::}\FunctionTok{tribble}\NormalTok{(}
  \SpecialCharTok{\textasciitilde{}}\NormalTok{metric,               }\SpecialCharTok{\textasciitilde{}}\NormalTok{hand\_value,              }\SpecialCharTok{\textasciitilde{}}\NormalTok{quantlib\_value,}
  \StringTok{"coupon\_leg\_npv"}\NormalTok{,       premium\_leg\_hand\_flat,     cds\_flat}\SpecialCharTok{$}\FunctionTok{couponLegNPV}\NormalTok{(),}
  \StringTok{"default\_leg\_npv"}\NormalTok{,      protection\_leg\_hand\_flat,  cds\_flat}\SpecialCharTok{$}\FunctionTok{defaultLegNPV}\NormalTok{(),}
  \StringTok{"cds\_npv"}\NormalTok{,              cds\_npv\_hand\_flat,         GaussQuant}\SpecialCharTok{::}\FunctionTok{cds\_npv\_GQL}\NormalTok{(}
\NormalTok{                                                       cds\_flat}
\NormalTok{                                                     ),}
  \StringTok{"fair\_spread"}\NormalTok{,          fair\_spread\_hand\_flat,     GaussQuant}\SpecialCharTok{::}\FunctionTok{cds\_fair\_spread\_GQL}\NormalTok{(}
\NormalTok{                                                       cds\_flat}
\NormalTok{                                                     )}
\NormalTok{) }\SpecialCharTok{|\textgreater{}}
\NormalTok{  dplyr}\SpecialCharTok{::}\FunctionTok{mutate}\NormalTok{(}
    \AttributeTok{difference =}\NormalTok{ .data}\SpecialCharTok{$}\NormalTok{hand\_value }\SpecialCharTok{{-}}
\NormalTok{      .data}\SpecialCharTok{$}\NormalTok{quantlib\_value}
\NormalTok{  )}

\NormalTok{flat\_hand\_summary\_tbl }\OtherTok{\textless{}{-}}\NormalTok{ tibble}\SpecialCharTok{::}\FunctionTok{tribble}\NormalTok{(}
  \SpecialCharTok{\textasciitilde{}}\NormalTok{metric,                    }\SpecialCharTok{\textasciitilde{}}\NormalTok{value,}
  \StringTok{"risky\_pv01"}\NormalTok{,                risky\_pv01\_hand\_flat,}
  \StringTok{"protection\_leg\_per\_unit"}\NormalTok{,   protection\_leg\_per\_unit\_hand\_flat,}
  \StringTok{"hand\_fair\_spread"}\NormalTok{,          fair\_spread\_hand\_flat,}
  \StringTok{"quantlib\_fair\_spread"}\NormalTok{,      GaussQuant}\SpecialCharTok{::}\FunctionTok{cds\_fair\_spread\_GQL}\NormalTok{(}
\NormalTok{                                  cds\_flat}
\NormalTok{                                ),}
  \StringTok{"fair\_spread\_difference"}\NormalTok{,    fair\_spread\_hand\_flat }\SpecialCharTok{{-}}
\NormalTok{                                  GaussQuant}\SpecialCharTok{::}\FunctionTok{cds\_fair\_spread\_GQL}\NormalTok{(}
\NormalTok{                                    cds\_flat}
\NormalTok{                                  )}
\NormalTok{)}

\NormalTok{GaussQuant}\SpecialCharTok{::}\FunctionTok{show\_tbl\_GQH}\NormalTok{(}
\NormalTok{  flat\_hand\_check\_tbl,}
  \StringTok{"CDS hand calculation and MidPoint engine comparison"}\NormalTok{,}
  \AttributeTok{n =} \DecValTok{20}
\NormalTok{)}
\CommentTok{\#\textgreater{} \# A tibble: 4 x 4}
\CommentTok{\#\textgreater{}   metric           hand\_value quantlib\_value  difference}
\CommentTok{\#\textgreater{}   \textless{}chr\textgreater{}                 \textless{}dbl\textgreater{}          \textless{}dbl\textgreater{}       \textless{}dbl\textgreater{}}
\CommentTok{\#\textgreater{} 1 coupon\_leg\_npv  {-}94244.        {-}94246.     2.10       }
\CommentTok{\#\textgreater{} 2 default\_leg\_npv 122362.        122362.     0          }
\CommentTok{\#\textgreater{} 3 cds\_npv          28118.         28116.     2.10       }
\CommentTok{\#\textgreater{} 4 fair\_spread          0.0130         0.0130 0.000000289}

\NormalTok{GaussQuant}\SpecialCharTok{::}\FunctionTok{show\_tbl\_GQH}\NormalTok{(}
\NormalTok{  flat\_hand\_summary\_tbl,}
  \StringTok{"CDS hand{-}calculation summary"}\NormalTok{,}
  \AttributeTok{n =} \DecValTok{20}
\NormalTok{)}
\CommentTok{\#\textgreater{} \# A tibble: 5 x 2}
\CommentTok{\#\textgreater{}   metric                        value}
\CommentTok{\#\textgreater{}   \textless{}chr\textgreater{}                         \textless{}dbl\textgreater{}}
\CommentTok{\#\textgreater{} 1 risky\_pv01              0.942      }
\CommentTok{\#\textgreater{} 2 protection\_leg\_per\_unit 0.0122     }
\CommentTok{\#\textgreater{} 3 hand\_fair\_spread        0.0130     }
\CommentTok{\#\textgreater{} 4 quantlib\_fair\_spread    0.0130     }
\CommentTok{\#\textgreater{} 5 fair\_spread\_difference  0.000000289}
\end{Highlighting}
\end{Shaded}

手計算とMidPoint
Engineは近い値となるが、厳密には一致しない。差にはAccrual on
default、日付規則、Protection開始日およびQuantLib内部の計算規約が含まれる。

\section{8.4 Hazard
Rate推定と感応度}\label{hazard-rateux63a8ux5b9aux3068ux611fux5fdcux5ea6}

\subsection{8.4.1 Zero-NPV CDSからHazard
Rateを推定する}\label{zero-npv-cdsux304bux3089hazard-rateux3092ux63a8ux5b9aux3059ux308b}

市場Spreadを契約Couponとして設定したCDSは、契約開始時点でNPVがゼロとなるようにHazard
Rateを調整できる。

ここでは、市場Spread 130bpのCDSについて、MidPoint
EngineのNPVがゼロとなるフラットHazard
Rateを\texttt{uniroot()}で推定する。

\begin{Shaded}
\begin{Highlighting}[]
\NormalTok{cds\_market\_flat }\OtherTok{\textless{}{-}}\NormalTok{ QuantLib}\SpecialCharTok{::}\FunctionTok{CreditDefaultSwap}\NormalTok{(}
\NormalTok{  GaussQuant}\SpecialCharTok{::}\FunctionTok{cds\_buyer\_side\_GQL}\NormalTok{(),}
\NormalTok{  notional\_flat,}
\NormalTok{  market\_spread\_flat,}
\NormalTok{  cds\_schedule\_flat,}
  \StringTok{"Following"}\NormalTok{,}
\NormalTok{  actual\_360}
\NormalTok{)}

\NormalTok{cds\_market\_npv\_at\_hazard }\OtherTok{\textless{}{-}} \ControlFlowTok{function}\NormalTok{(hazard\_rate) \{}
\NormalTok{  engine }\OtherTok{\textless{}{-}}\NormalTok{ GaussQuant}\SpecialCharTok{::}\FunctionTok{cds\_midpoint\_engine\_GQL}\NormalTok{(}
    \AttributeTok{recovery\_rate =}\NormalTok{ recovery\_rate\_flat,}
    \AttributeTok{hazard\_rate =}\NormalTok{ hazard\_rate,}
    \AttributeTok{discount\_rate =}\NormalTok{ risk\_free\_rate\_flat,}
    \AttributeTok{reference\_date =}\NormalTok{ trade\_date\_flat,}
    \AttributeTok{day\_counter =} \StringTok{"Actual365Fixed"}
\NormalTok{  )}

\NormalTok{  cds\_market\_flat}\SpecialCharTok{$}\FunctionTok{setPricingEngine}\NormalTok{(}
\NormalTok{    engine}
\NormalTok{  )}

\NormalTok{  GaussQuant}\SpecialCharTok{::}\FunctionTok{cds\_npv\_GQL}\NormalTok{(}
\NormalTok{    cds\_market\_flat}
\NormalTok{  )}
\NormalTok{\}}

\NormalTok{hazard\_calibration\_result }\OtherTok{\textless{}{-}}\NormalTok{ stats}\SpecialCharTok{::}\FunctionTok{uniroot}\NormalTok{(}
\NormalTok{  cds\_market\_npv\_at\_hazard,}
  \AttributeTok{interval =} \FunctionTok{c}\NormalTok{(}
    \FloatTok{1e{-}8}\NormalTok{,}
    \FloatTok{0.50}
\NormalTok{  ),}
  \AttributeTok{tol =} \FloatTok{1e{-}10}
\NormalTok{)}

\NormalTok{hazard\_rate\_calibrated }\OtherTok{\textless{}{-}}\NormalTok{ hazard\_calibration\_result}\SpecialCharTok{$}\NormalTok{root}

\NormalTok{calibrated\_midpoint\_engine }\OtherTok{\textless{}{-}}\NormalTok{ GaussQuant}\SpecialCharTok{::}\FunctionTok{cds\_midpoint\_engine\_GQL}\NormalTok{(}
  \AttributeTok{recovery\_rate =}\NormalTok{ recovery\_rate\_flat,}
  \AttributeTok{hazard\_rate =}\NormalTok{ hazard\_rate\_calibrated,}
  \AttributeTok{discount\_rate =}\NormalTok{ risk\_free\_rate\_flat,}
  \AttributeTok{reference\_date =}\NormalTok{ trade\_date\_flat,}
  \AttributeTok{day\_counter =} \StringTok{"Actual365Fixed"}
\NormalTok{)}

\NormalTok{cds\_market\_flat}\SpecialCharTok{$}\FunctionTok{setPricingEngine}\NormalTok{(}
\NormalTok{  calibrated\_midpoint\_engine}
\NormalTok{)}
\CommentTok{\#\textgreater{} NULL}

\NormalTok{hazard\_calibration\_tbl }\OtherTok{\textless{}{-}}\NormalTok{ tibble}\SpecialCharTok{::}\FunctionTok{tribble}\NormalTok{(}
  \SpecialCharTok{\textasciitilde{}}\NormalTok{metric,                       }\SpecialCharTok{\textasciitilde{}}\NormalTok{value,}
  \StringTok{"market\_spread"}\NormalTok{,                market\_spread\_flat,}
  \StringTok{"simple\_hazard\_approximation"}\NormalTok{,  hazard\_rate\_approximation,}
  \StringTok{"calibrated\_hazard\_rate"}\NormalTok{,       hazard\_rate\_calibrated,}
  \StringTok{"zero\_npv\_check"}\NormalTok{,               GaussQuant}\SpecialCharTok{::}\FunctionTok{cds\_npv\_GQL}\NormalTok{(}
\NormalTok{                                      cds\_market\_flat}
\NormalTok{                                    ),}
  \StringTok{"fair\_spread\_check"}\NormalTok{,            GaussQuant}\SpecialCharTok{::}\FunctionTok{cds\_fair\_spread\_GQL}\NormalTok{(}
\NormalTok{                                      cds\_market\_flat}
\NormalTok{                                    )}
\NormalTok{)}

\NormalTok{GaussQuant}\SpecialCharTok{::}\FunctionTok{show\_tbl\_GQH}\NormalTok{(}
\NormalTok{  hazard\_calibration\_tbl,}
  \StringTok{"Flat hazard calibration from a zero{-}NPV CDS"}\NormalTok{,}
  \AttributeTok{n =} \DecValTok{20}
\NormalTok{)}
\CommentTok{\#\textgreater{} \# A tibble: 5 x 2}
\CommentTok{\#\textgreater{}   metric                            value}
\CommentTok{\#\textgreater{}   \textless{}chr\textgreater{}                             \textless{}dbl\textgreater{}}
\CommentTok{\#\textgreater{} 1 market\_spread               0.013      }
\CommentTok{\#\textgreater{} 2 simple\_hazard\_approximation 0.0217     }
\CommentTok{\#\textgreater{} 3 calibrated\_hazard\_rate      0.0217     }
\CommentTok{\#\textgreater{} 4 zero\_npv\_check              0.000000367}
\CommentTok{\#\textgreater{} 5 fair\_spread\_check           0.0130}
\end{Highlighting}
\end{Shaded}

単純近似 \(\lambda\approx s/(1-R)\)
とカリブレーション結果は近いが、支払頻度と割引の影響により差が生じる。

\subsection{8.4.2 SpreadとRecovery
Rateの感応度}\label{spreadux3068recovery-rateux306eux611fux5fdcux5ea6}

単純近似の下では、Spreadが高いほどHazard
Rateは高くなる。また、同じSpreadでRecovery
Rateを高く仮定すると、デフォルト時損失が小さくなるため、そのSpreadを説明するために必要なHazard
Rateは高くなる。

\begin{Shaded}
\begin{Highlighting}[]
\NormalTok{credit\_sensitivity\_input\_tbl }\OtherTok{\textless{}{-}}\NormalTok{ tidyr}\SpecialCharTok{::}\FunctionTok{crossing}\NormalTok{(}
  \AttributeTok{market\_spread\_bp =} \FunctionTok{c}\NormalTok{(}
    \DecValTok{50}\NormalTok{,}
    \DecValTok{100}\NormalTok{,}
    \DecValTok{200}
\NormalTok{  ),}
  \AttributeTok{recovery\_rate =} \FunctionTok{c}\NormalTok{(}
    \FloatTok{0.20}\NormalTok{,}
    \FloatTok{0.40}\NormalTok{,}
    \FloatTok{0.60}
\NormalTok{  )}
\NormalTok{)}

\NormalTok{credit\_sensitivity\_tbl }\OtherTok{\textless{}{-}}\NormalTok{ credit\_sensitivity\_input\_tbl }\SpecialCharTok{|\textgreater{}}
\NormalTok{  dplyr}\SpecialCharTok{::}\FunctionTok{mutate}\NormalTok{(}
    \AttributeTok{market\_spread =}\NormalTok{ .data}\SpecialCharTok{$}\NormalTok{market\_spread\_bp }\SpecialCharTok{/}
      \DecValTok{10000}\NormalTok{,}
    \AttributeTok{loss\_given\_default =} \DecValTok{1} \SpecialCharTok{{-}}
\NormalTok{      .data}\SpecialCharTok{$}\NormalTok{recovery\_rate,}
    \AttributeTok{hazard\_rate\_approx =}\NormalTok{ .data}\SpecialCharTok{$}\NormalTok{market\_spread }\SpecialCharTok{/}
\NormalTok{      .data}\SpecialCharTok{$}\NormalTok{loss\_given\_default,}
    \AttributeTok{survival\_probability\_5y =} \FunctionTok{exp}\NormalTok{(}
      \SpecialCharTok{{-}}\DecValTok{5} \SpecialCharTok{*}\NormalTok{ .data}\SpecialCharTok{$}\NormalTok{hazard\_rate\_approx}
\NormalTok{    ),}
    \AttributeTok{cumulative\_default\_probability\_5y =} \DecValTok{1} \SpecialCharTok{{-}}
\NormalTok{      .data}\SpecialCharTok{$}\NormalTok{survival\_probability\_5y}
\NormalTok{  )}

\NormalTok{GaussQuant}\SpecialCharTok{::}\FunctionTok{show\_tbl\_GQH}\NormalTok{(}
\NormalTok{  credit\_sensitivity\_tbl,}
  \StringTok{"Spread and recovery{-}rate sensitivity"}\NormalTok{,}
  \AttributeTok{n =} \DecValTok{30}
\NormalTok{)}
\CommentTok{\#\textgreater{} \# A tibble: 9 x 7}
\CommentTok{\#\textgreater{}   market\_spread\_bp recovery\_rate market\_spread loss\_given\_default}
\CommentTok{\#\textgreater{}              \textless{}dbl\textgreater{}         \textless{}dbl\textgreater{}         \textless{}dbl\textgreater{}              \textless{}dbl\textgreater{}}
\CommentTok{\#\textgreater{} 1               50           0.2         0.005                0.8}
\CommentTok{\#\textgreater{} 2               50           0.4         0.005                0.6}
\CommentTok{\#\textgreater{} 3               50           0.6         0.005                0.4}
\CommentTok{\#\textgreater{} 4              100           0.2         0.01                 0.8}
\CommentTok{\#\textgreater{} 5              100           0.4         0.01                 0.6}
\CommentTok{\#\textgreater{} 6              100           0.6         0.01                 0.4}
\CommentTok{\#\textgreater{} 7              200           0.2         0.02                 0.8}
\CommentTok{\#\textgreater{} 8              200           0.4         0.02                 0.6}
\CommentTok{\#\textgreater{} 9              200           0.6         0.02                 0.4}
\CommentTok{\#\textgreater{} \# i 3 more variables: hazard\_rate\_approx \textless{}dbl\textgreater{}, survival\_probability\_5y \textless{}dbl\textgreater{},}
\CommentTok{\#\textgreater{} \#   cumulative\_default\_probability\_5y \textless{}dbl\textgreater{}}
\end{Highlighting}
\end{Shaded}

Recovery RateとHazard Rateは、CDS
Spreadだけから完全に独立して識別できない。実務ではRecovery
Rateを外生的に仮定し、その仮定の下でDefault Probability
Curveを構築することが多い。

\section{8.5
標準CDS2015とISDAカーブ}\label{ux6a19ux6e96cds2015ux3068isdaux30abux30fcux30d6}

\subsection{8.5.1
標準CDSとCDS2015満期}\label{ux6a19ux6e96cdsux3068cds2015ux6e80ux671f}

標準CDSでは、満期を単純に取引日から5年後とするのではなく、CDS
IMM日とCDS2015ルールに従って決定する。

次の5年物CDSを考える。

\begin{itemize}
\tightlist
\item
  取引日：2022年8月19日
\item
  市場Spread：30.426bp
\item
  契約Coupon：100bp
\item
  Notional：10,000,000
\item
  Recovery Rate：40\%
\item
  外生的なOIS割引率：0.25\%
\end{itemize}

第II-3章ではOISカーブの構築を扱った。本章の目的は信用カーブとCDS評価であるため、割引カーブの動学やブートストラップを繰り返さず、フラットOIS割引率を外生的に与える。

\begin{Shaded}
\begin{Highlighting}[]
\NormalTok{trade\_date\_isda }\OtherTok{\textless{}{-}}\NormalTok{ GaussQuant}\SpecialCharTok{::}\FunctionTok{date\_GQL}\NormalTok{(}
  \StringTok{"2022{-}08{-}19"}
\NormalTok{)}

\NormalTok{GaussQuant}\SpecialCharTok{::}\FunctionTok{set\_eval\_date\_GQL}\NormalTok{(}
\NormalTok{  trade\_date\_isda}
\NormalTok{)}

\NormalTok{cds\_tenor\_isda }\OtherTok{\textless{}{-}}\NormalTok{ GaussQuant}\SpecialCharTok{::}\FunctionTok{period\_years\_GQL}\NormalTok{(}
  \DecValTok{5}\NormalTok{L}
\NormalTok{)}

\NormalTok{maturity\_date\_isda }\OtherTok{\textless{}{-}}\NormalTok{ QuantLib}\SpecialCharTok{::}\FunctionTok{cdsMaturity}\NormalTok{(}
\NormalTok{  trade\_date\_isda,}
\NormalTok{  cds\_tenor\_isda,}
  \StringTok{"CDS2015"}
\NormalTok{)}

\NormalTok{market\_spread\_isda }\OtherTok{\textless{}{-}} \FloatTok{30.426} \SpecialCharTok{/} \DecValTok{10000}
\NormalTok{contract\_coupon\_isda }\OtherTok{\textless{}{-}} \DecValTok{100} \SpecialCharTok{/} \DecValTok{10000}
\NormalTok{notional\_isda }\OtherTok{\textless{}{-}} \DecValTok{10000000}
\NormalTok{recovery\_rate\_isda }\OtherTok{\textless{}{-}} \FloatTok{0.40}
\NormalTok{ois\_discount\_rate\_isda }\OtherTok{\textless{}{-}} \FloatTok{0.0025}

\NormalTok{isda\_cds\_input\_tbl }\OtherTok{\textless{}{-}}\NormalTok{ tibble}\SpecialCharTok{::}\FunctionTok{tribble}\NormalTok{(}
  \SpecialCharTok{\textasciitilde{}}\NormalTok{input,               }\SpecialCharTok{\textasciitilde{}}\NormalTok{value,}
  \StringTok{"trade\_date"}\NormalTok{,          }\StringTok{"2022{-}08{-}19"}\NormalTok{,}
  \StringTok{"cds\_tenor"}\NormalTok{,           }\StringTok{"5Y"}\NormalTok{,}
  \StringTok{"maturity\_date"}\NormalTok{,       GaussQuant}\SpecialCharTok{::}\FunctionTok{iso\_GQL}\NormalTok{(}
\NormalTok{                           maturity\_date\_isda}
\NormalTok{                         ),}
  \StringTok{"market\_spread\_bp"}\NormalTok{,    }\FunctionTok{as.character}\NormalTok{(}
\NormalTok{                           market\_spread\_isda }\SpecialCharTok{*}
                             \DecValTok{10000}
\NormalTok{                         ),}
  \StringTok{"contract\_coupon\_bp"}\NormalTok{,  }\FunctionTok{as.character}\NormalTok{(}
\NormalTok{                           contract\_coupon\_isda }\SpecialCharTok{*}
                             \DecValTok{10000}
\NormalTok{                         ),}
  \StringTok{"notional"}\NormalTok{,            }\FunctionTok{as.character}\NormalTok{(}
\NormalTok{                           notional\_isda}
\NormalTok{                         ),}
  \StringTok{"recovery\_rate"}\NormalTok{,       }\FunctionTok{as.character}\NormalTok{(}
\NormalTok{                           recovery\_rate\_isda}
\NormalTok{                         ),}
  \StringTok{"ois\_discount\_rate"}\NormalTok{,   }\FunctionTok{as.character}\NormalTok{(}
\NormalTok{                           ois\_discount\_rate\_isda}
\NormalTok{                         )}
\NormalTok{)}

\NormalTok{GaussQuant}\SpecialCharTok{::}\FunctionTok{show\_tbl\_GQH}\NormalTok{(}
\NormalTok{  isda\_cds\_input\_tbl,}
  \StringTok{"ISDA CDS input assumptions"}\NormalTok{,}
  \AttributeTok{n =} \DecValTok{20}
\NormalTok{)}
\CommentTok{\#\textgreater{} \# A tibble: 8 x 2}
\CommentTok{\#\textgreater{}   input              value     }
\CommentTok{\#\textgreater{}   \textless{}chr\textgreater{}              \textless{}chr\textgreater{}     }
\CommentTok{\#\textgreater{} 1 trade\_date         2022{-}08{-}19}
\CommentTok{\#\textgreater{} 2 cds\_tenor          5Y        }
\CommentTok{\#\textgreater{} 3 maturity\_date      2027{-}06{-}20}
\CommentTok{\#\textgreater{} 4 market\_spread\_bp   30.426    }
\CommentTok{\#\textgreater{} 5 contract\_coupon\_bp 100       }
\CommentTok{\#\textgreater{} 6 notional           1e+07     }
\CommentTok{\#\textgreater{} 7 recovery\_rate      0.4       }
\CommentTok{\#\textgreater{} 8 ois\_discount\_rate  0.0025}
\end{Highlighting}
\end{Shaded}

\subsection{8.5.2 標準CDS2015
Schedule}\label{ux6a19ux6e96cds2015-schedule}

CDS2015ルールにより、Premium支払日は標準化された四半期日へ調整される。

\begin{Shaded}
\begin{Highlighting}[]
\NormalTok{cds\_calendar\_isda }\OtherTok{\textless{}{-}}\NormalTok{ QuantLib}\SpecialCharTok{::}\FunctionTok{WeekendsOnly}\NormalTok{()}

\NormalTok{cds\_schedule\_isda }\OtherTok{\textless{}{-}}\NormalTok{ QuantLib}\SpecialCharTok{::}\FunctionTok{Schedule}\NormalTok{(}
\NormalTok{  trade\_date\_isda,}
\NormalTok{  maturity\_date\_isda,}
\NormalTok{  GaussQuant}\SpecialCharTok{::}\FunctionTok{period\_months\_GQL}\NormalTok{(}\DecValTok{3}\NormalTok{),}
\NormalTok{  cds\_calendar\_isda,}
  \StringTok{"Following"}\NormalTok{,}
  \StringTok{"Unadjusted"}\NormalTok{,}
  \StringTok{"CDS2015"}\NormalTok{,}
  \ConstantTok{FALSE}
\NormalTok{)}

\NormalTok{GaussQuant}\SpecialCharTok{::}\FunctionTok{show\_tbl\_GQH}\NormalTok{(}
\NormalTok{  GaussQuant}\SpecialCharTok{::}\FunctionTok{schedule\_table\_GQL}\NormalTok{(}
\NormalTok{    cds\_schedule\_isda}
\NormalTok{  ),}
  \StringTok{"CDS schedule: 5Y CDS2015 example"}\NormalTok{,}
  \AttributeTok{n =} \DecValTok{30}
\NormalTok{)}
\CommentTok{\#\textgreater{} \# A tibble: 21 x 2}
\CommentTok{\#\textgreater{}    period schedule\_date}
\CommentTok{\#\textgreater{}     \textless{}int\textgreater{} \textless{}date\textgreater{}       }
\CommentTok{\#\textgreater{}  1      1 2022{-}06{-}20   }
\CommentTok{\#\textgreater{}  2      2 2022{-}09{-}20   }
\CommentTok{\#\textgreater{}  3      3 2022{-}12{-}20   }
\CommentTok{\#\textgreater{}  4      4 2023{-}03{-}20   }
\CommentTok{\#\textgreater{}  5      5 2023{-}06{-}20   }
\CommentTok{\#\textgreater{}  6      6 2023{-}09{-}20   }
\CommentTok{\#\textgreater{}  7      7 2023{-}12{-}20   }
\CommentTok{\#\textgreater{}  8      8 2024{-}03{-}20   }
\CommentTok{\#\textgreater{}  9      9 2024{-}06{-}20   }
\CommentTok{\#\textgreater{} 10     10 2024{-}09{-}20   }
\CommentTok{\#\textgreater{} \# i 11 more rows}
\end{Highlighting}
\end{Shaded}

\subsection{8.5.3
ISDA用割引カーブ}\label{isdaux7528ux5272ux5f15ux30abux30fcux30d6}

\begin{Shaded}
\begin{Highlighting}[]
\NormalTok{discount\_curve\_isda }\OtherTok{\textless{}{-}}\NormalTok{ GaussQuant}\SpecialCharTok{::}\FunctionTok{flat\_curve\_GQL}\NormalTok{(}
  \AttributeTok{reference\_date =}\NormalTok{ trade\_date\_isda,}
  \AttributeTok{rate =}\NormalTok{ ois\_discount\_rate\_isda,}
  \AttributeTok{day\_counter =}\NormalTok{ actual\_365,}
  \AttributeTok{compounding =} \StringTok{"Continuous"}
\NormalTok{)}

\NormalTok{discount\_curve\_handle\_isda }\OtherTok{\textless{}{-}}\NormalTok{ QuantLib}\SpecialCharTok{::}\FunctionTok{YieldTermStructureHandle}\NormalTok{(}
\NormalTok{  discount\_curve\_isda}
\NormalTok{)}

\NormalTok{isda\_discount\_check\_tbl }\OtherTok{\textless{}{-}}\NormalTok{ tibble}\SpecialCharTok{::}\FunctionTok{tribble}\NormalTok{(}
  \SpecialCharTok{\textasciitilde{}}\NormalTok{metric,                         }\SpecialCharTok{\textasciitilde{}}\NormalTok{value,}
  \StringTok{"ois\_discount\_rate"}\NormalTok{,              ois\_discount\_rate\_isda,}
  \StringTok{"discount\_factor\_at\_maturity"}\NormalTok{,    discount\_curve\_isda}\SpecialCharTok{$}\FunctionTok{discount}\NormalTok{(}
\NormalTok{                                       maturity\_date\_isda}
\NormalTok{                                     )}
\NormalTok{)}

\NormalTok{GaussQuant}\SpecialCharTok{::}\FunctionTok{show\_tbl\_GQH}\NormalTok{(}
\NormalTok{  isda\_discount\_check\_tbl,}
  \StringTok{"ISDA discount{-}curve check"}\NormalTok{,}
  \AttributeTok{n =} \DecValTok{20}
\NormalTok{)}
\CommentTok{\#\textgreater{} \# A tibble: 2 x 2}
\CommentTok{\#\textgreater{}   metric                       value}
\CommentTok{\#\textgreater{}   \textless{}chr\textgreater{}                        \textless{}dbl\textgreater{}}
\CommentTok{\#\textgreater{} 1 ois\_discount\_rate           0.0025}
\CommentTok{\#\textgreater{} 2 discount\_factor\_at\_maturity 0.988}
\end{Highlighting}
\end{Shaded}

\subsection{8.5.4 ISDA整合Hazard
Rateの推定}\label{isdaux6574ux5408hazard-rateux306eux63a8ux5b9a}

市場Spread 30.426bpをRunning Spreadとする標準CDSを作成し、ISDA
EngineのNPVがゼロとなるフラットHazard Rateを推定する。

\begin{Shaded}
\begin{Highlighting}[]
\NormalTok{market\_cds\_isda }\OtherTok{\textless{}{-}}\NormalTok{ GaussQuant}\SpecialCharTok{::}\FunctionTok{make\_cds\_GQL}\NormalTok{(}
  \AttributeTok{maturity\_date =}\NormalTok{ maturity\_date\_isda,}
  \AttributeTok{running\_spread =}\NormalTok{ market\_spread\_isda,}
  \AttributeTok{notional =}\NormalTok{ notional\_isda,}
  \AttributeTok{side =} \StringTok{"buyer"}\NormalTok{,}
  \AttributeTok{trade\_date =}\NormalTok{ trade\_date\_isda,}
  \AttributeTok{coupon\_tenor =} \StringTok{"3M"}\NormalTok{,}
  \AttributeTok{day\_counter =} \StringTok{"Actual360"}\NormalTok{,}
  \AttributeTok{date\_generation\_rule =} \StringTok{"CDS2015"}
\NormalTok{)}

\NormalTok{isda\_market\_npv\_at\_hazard }\OtherTok{\textless{}{-}} \ControlFlowTok{function}\NormalTok{(hazard\_rate) \{}
\NormalTok{  probability\_handle }\OtherTok{\textless{}{-}}\NormalTok{ GaussQuant}\SpecialCharTok{::}\FunctionTok{flat\_hazard\_rate\_GQL}\NormalTok{(}
    \AttributeTok{hazard\_rate =}\NormalTok{ hazard\_rate,}
    \AttributeTok{reference\_date =}\NormalTok{ trade\_date\_isda,}
    \AttributeTok{day\_counter =} \StringTok{"Actual365Fixed"}
\NormalTok{  )}

\NormalTok{  engine }\OtherTok{\textless{}{-}}\NormalTok{ GaussQuant}\SpecialCharTok{::}\FunctionTok{cds\_isda\_engine\_GQL}\NormalTok{(}
    \AttributeTok{hazard\_curve\_handle =}\NormalTok{ probability\_handle,}
    \AttributeTok{recovery\_rate =}\NormalTok{ recovery\_rate\_isda,}
    \AttributeTok{discount\_curve\_handle =}\NormalTok{ discount\_curve\_handle\_isda}
\NormalTok{  )}

\NormalTok{  market\_cds\_isda}\SpecialCharTok{$}\FunctionTok{setPricingEngine}\NormalTok{(}
\NormalTok{    engine}
\NormalTok{  )}

\NormalTok{  GaussQuant}\SpecialCharTok{::}\FunctionTok{cds\_npv\_GQL}\NormalTok{(}
\NormalTok{    market\_cds\_isda}
\NormalTok{  )}
\NormalTok{\}}

\NormalTok{isda\_hazard\_calibration\_result }\OtherTok{\textless{}{-}}\NormalTok{ stats}\SpecialCharTok{::}\FunctionTok{uniroot}\NormalTok{(}
\NormalTok{  isda\_market\_npv\_at\_hazard,}
  \AttributeTok{interval =} \FunctionTok{c}\NormalTok{(}
    \FloatTok{1e{-}8}\NormalTok{,}
    \FloatTok{0.50}
\NormalTok{  ),}
  \AttributeTok{tol =} \FloatTok{1e{-}10}
\NormalTok{)}

\NormalTok{hazard\_rate\_isda }\OtherTok{\textless{}{-}}\NormalTok{ isda\_hazard\_calibration\_result}\SpecialCharTok{$}\NormalTok{root}

\NormalTok{hazard\_curve\_isda }\OtherTok{\textless{}{-}}\NormalTok{ QuantLib}\SpecialCharTok{::}\FunctionTok{FlatHazardRate}\NormalTok{(}
\NormalTok{  trade\_date\_isda,}
\NormalTok{  GaussQuant}\SpecialCharTok{::}\FunctionTok{quote\_handle\_GQL}\NormalTok{(}
\NormalTok{    hazard\_rate\_isda}
\NormalTok{  ),}
\NormalTok{  actual\_365}
\NormalTok{)}

\NormalTok{hazard\_curve\_handle\_isda }\OtherTok{\textless{}{-}}\NormalTok{ QuantLib}\SpecialCharTok{::}\FunctionTok{DefaultProbabilityTermStructureHandle}\NormalTok{(}
\NormalTok{  hazard\_curve\_isda}
\NormalTok{)}

\NormalTok{isda\_hazard\_check\_tbl }\OtherTok{\textless{}{-}}\NormalTok{ tibble}\SpecialCharTok{::}\FunctionTok{tribble}\NormalTok{(}
  \SpecialCharTok{\textasciitilde{}}\NormalTok{metric,                         }\SpecialCharTok{\textasciitilde{}}\NormalTok{value,}
  \StringTok{"market\_spread"}\NormalTok{,                  market\_spread\_isda,}
  \StringTok{"recovery\_rate"}\NormalTok{,                  recovery\_rate\_isda,}
  \StringTok{"implied\_hazard\_rate"}\NormalTok{,            hazard\_rate\_isda,}
  \StringTok{"simple\_hazard\_approximation"}\NormalTok{,    market\_spread\_isda }\SpecialCharTok{/}
\NormalTok{                                      (}
                                        \DecValTok{1} \SpecialCharTok{{-}}
\NormalTok{                                          recovery\_rate\_isda}
\NormalTok{                                      ),}
  \StringTok{"survival\_probability\_at\_maturity"}\NormalTok{,}
\NormalTok{  hazard\_curve\_isda}\SpecialCharTok{$}\FunctionTok{survivalProbability}\NormalTok{(}
\NormalTok{    maturity\_date\_isda}
\NormalTok{  ),}
  \StringTok{"default\_probability\_to\_maturity"}\NormalTok{,}
  \DecValTok{1} \SpecialCharTok{{-}}
\NormalTok{    hazard\_curve\_isda}\SpecialCharTok{$}\FunctionTok{survivalProbability}\NormalTok{(}
\NormalTok{      maturity\_date\_isda}
\NormalTok{    )}
\NormalTok{)}

\NormalTok{GaussQuant}\SpecialCharTok{::}\FunctionTok{show\_tbl\_GQH}\NormalTok{(}
\NormalTok{  isda\_hazard\_check\_tbl,}
  \StringTok{"ISDA implied hazard{-}rate check"}\NormalTok{,}
  \AttributeTok{n =} \DecValTok{20}
\NormalTok{)}
\CommentTok{\#\textgreater{} \# A tibble: 6 x 2}
\CommentTok{\#\textgreater{}   metric                             value}
\CommentTok{\#\textgreater{}   \textless{}chr\textgreater{}                              \textless{}dbl\textgreater{}}
\CommentTok{\#\textgreater{} 1 market\_spread                    0.00304}
\CommentTok{\#\textgreater{} 2 recovery\_rate                    0.4    }
\CommentTok{\#\textgreater{} 3 implied\_hazard\_rate              0.00514}
\CommentTok{\#\textgreater{} 4 simple\_hazard\_approximation      0.00507}
\CommentTok{\#\textgreater{} 5 survival\_probability\_at\_maturity 0.975  }
\CommentTok{\#\textgreater{} 6 default\_probability\_to\_maturity  0.0246}
\end{Highlighting}
\end{Shaded}

市場SpreadとRecovery Rateから得られる単純近似と、ISDA
EngineでゼロNPVとなるHazard Rateには差が生じる。標準CDSの日付、Accrual
on default、割引および数値積分が評価へ反映されるためである。

\section{8.6 ISDA
Engineと評価比較}\label{isda-engineux3068ux8a55ux4fa1ux6bd4ux8f03}

\subsection{8.6.1 ISDA CDS Engine}\label{isda-cds-engine}

契約Couponを標準Coupon 100bpへ固定した5年物CDSを評価する。

市場Spreadは30.426bpであり、契約Coupon
100bpより低い。プロテクション買い手は市場水準より高いPremiumを支払うため、Running
Couponだけで評価すると一般に負のNPVとなる。

\begin{Shaded}
\begin{Highlighting}[]
\NormalTok{isda\_engine }\OtherTok{\textless{}{-}}\NormalTok{ GaussQuant}\SpecialCharTok{::}\FunctionTok{cds\_isda\_engine\_GQL}\NormalTok{(}
  \AttributeTok{hazard\_curve\_handle =}\NormalTok{ hazard\_curve\_handle\_isda,}
  \AttributeTok{recovery\_rate =}\NormalTok{ recovery\_rate\_isda,}
  \AttributeTok{discount\_curve\_handle =}\NormalTok{ discount\_curve\_handle\_isda}
\NormalTok{)}

\NormalTok{cds\_isda }\OtherTok{\textless{}{-}}\NormalTok{ GaussQuant}\SpecialCharTok{::}\FunctionTok{make\_cds\_GQL}\NormalTok{(}
  \AttributeTok{maturity\_date =}\NormalTok{ maturity\_date\_isda,}
  \AttributeTok{running\_spread =}\NormalTok{ contract\_coupon\_isda,}
  \AttributeTok{notional =}\NormalTok{ notional\_isda,}
  \AttributeTok{side =} \StringTok{"buyer"}\NormalTok{,}
  \AttributeTok{trade\_date =}\NormalTok{ trade\_date\_isda,}
  \AttributeTok{coupon\_tenor =} \StringTok{"3M"}\NormalTok{,}
  \AttributeTok{day\_counter =} \StringTok{"Actual360"}\NormalTok{,}
  \AttributeTok{date\_generation\_rule =} \StringTok{"CDS2015"}\NormalTok{,}
  \AttributeTok{pricing\_engine =}\NormalTok{ isda\_engine}
\NormalTok{)}

\NormalTok{cds\_isda\_summary\_base\_tbl }\OtherTok{\textless{}{-}}\NormalTok{ GaussQuant}\SpecialCharTok{::}\FunctionTok{cds\_summary\_GQL}\NormalTok{(}
\NormalTok{  cds\_isda}
\NormalTok{)}

\NormalTok{GaussQuant}\SpecialCharTok{::}\FunctionTok{show\_tbl\_GQH}\NormalTok{(}
\NormalTok{  cds\_isda\_summary\_base\_tbl,}
  \StringTok{"ISDA CDS base summary"}\NormalTok{,}
  \AttributeTok{n =} \DecValTok{20}
\NormalTok{)}
\CommentTok{\#\textgreater{} \# A tibble: 4 x 2}
\CommentTok{\#\textgreater{}   metric                  value}
\CommentTok{\#\textgreater{}   \textless{}chr\textgreater{}                   \textless{}dbl\textgreater{}}
\CommentTok{\#\textgreater{} 1 npv             {-}334963.     }
\CommentTok{\#\textgreater{} 2 fair\_spread           0.00304}
\CommentTok{\#\textgreater{} 3 coupon\_leg\_npv  {-}498392.     }
\CommentTok{\#\textgreater{} 4 default\_leg\_npv  146485.}
\end{Highlighting}
\end{Shaded}

\subsection{8.6.2 Accrual
rebateとNPV分解}\label{accrual-rebateux3068npvux5206ux89e3}

標準CDSでは、取引日が前回の標準Coupon日と次回Coupon日の間にある場合、契約開始時にAccrual
rebateが発生することがある。

この場合、プロテクション買い手から見たNPVは、

\[
V_{\mathrm{CDS}}
=
V_{\mathrm{coupon}}
+
V_{\mathrm{default}}
+
V_{\mathrm{accrual\ rebate}}
\]

となる。

\begin{Shaded}
\begin{Highlighting}[]
\NormalTok{coupon\_leg\_npv\_isda }\OtherTok{\textless{}{-}}\NormalTok{ cds\_isda}\SpecialCharTok{$}\FunctionTok{couponLegNPV}\NormalTok{()}
\NormalTok{default\_leg\_npv\_isda }\OtherTok{\textless{}{-}}\NormalTok{ cds\_isda}\SpecialCharTok{$}\FunctionTok{defaultLegNPV}\NormalTok{()}

\NormalTok{accrual\_rebate\_npv\_isda }\OtherTok{\textless{}{-}} \FunctionTok{tryCatch}\NormalTok{(}
  \FunctionTok{as.numeric}\NormalTok{(}
\NormalTok{    cds\_isda}\SpecialCharTok{$}\FunctionTok{accrualRebateNPV}\NormalTok{()}
\NormalTok{  ),}
  \AttributeTok{error =} \ControlFlowTok{function}\NormalTok{(e) }\ConstantTok{NA\_real\_}
\NormalTok{)}

\NormalTok{three\_component\_sum\_isda }\OtherTok{\textless{}{-}}\NormalTok{ coupon\_leg\_npv\_isda }\SpecialCharTok{+}
\NormalTok{  default\_leg\_npv\_isda }\SpecialCharTok{+}
\NormalTok{  accrual\_rebate\_npv\_isda}

\NormalTok{isda\_leg\_check\_tbl }\OtherTok{\textless{}{-}}\NormalTok{ tibble}\SpecialCharTok{::}\FunctionTok{tribble}\NormalTok{(}
  \SpecialCharTok{\textasciitilde{}}\NormalTok{metric,               }\SpecialCharTok{\textasciitilde{}}\NormalTok{value,}
  \StringTok{"coupon\_leg\_npv"}\NormalTok{,       coupon\_leg\_npv\_isda,}
  \StringTok{"default\_leg\_npv"}\NormalTok{,      default\_leg\_npv\_isda,}
  \StringTok{"accrual\_rebate\_npv"}\NormalTok{,   accrual\_rebate\_npv\_isda,}
  \StringTok{"three\_component\_sum"}\NormalTok{,  three\_component\_sum\_isda,}
  \StringTok{"reported\_npv"}\NormalTok{,         GaussQuant}\SpecialCharTok{::}\FunctionTok{cds\_npv\_GQL}\NormalTok{(}
\NormalTok{                            cds\_isda}
\NormalTok{                          ),}
  \StringTok{"difference"}\NormalTok{,           GaussQuant}\SpecialCharTok{::}\FunctionTok{cds\_npv\_GQL}\NormalTok{(}
\NormalTok{                            cds\_isda}
\NormalTok{                          ) }\SpecialCharTok{{-}}
\NormalTok{                            three\_component\_sum\_isda}
\NormalTok{)}

\NormalTok{GaussQuant}\SpecialCharTok{::}\FunctionTok{show\_tbl\_GQH}\NormalTok{(}
\NormalTok{  isda\_leg\_check\_tbl,}
  \StringTok{"ISDA CDS NPV decomposition"}\NormalTok{,}
  \AttributeTok{n =} \DecValTok{20}
\NormalTok{)}
\CommentTok{\#\textgreater{} \# A tibble: 6 x 2}
\CommentTok{\#\textgreater{}   metric                 value}
\CommentTok{\#\textgreater{}   \textless{}chr\textgreater{}                  \textless{}dbl\textgreater{}}
\CommentTok{\#\textgreater{} 1 coupon\_leg\_npv      {-}498392.}
\CommentTok{\#\textgreater{} 2 default\_leg\_npv      146485.}
\CommentTok{\#\textgreater{} 3 accrual\_rebate\_npv    16944.}
\CommentTok{\#\textgreater{} 4 three\_component\_sum {-}334963.}
\CommentTok{\#\textgreater{} 5 reported\_npv        {-}334963.}
\CommentTok{\#\textgreater{} 6 difference                0}
\end{Highlighting}
\end{Shaded}

\subsection{8.6.3 Fair SpreadとFair
Upfront}\label{fair-spreadux3068fair-upfront}

Fair Spreadは、Upfrontを授受せずにCDS NPVをゼロとするRunning
Spreadである。

Fair
Upfrontは、契約Couponを標準Couponへ固定したまま、契約価値を調整するために開始時点で授受するNotional比率である。

\begin{Shaded}
\begin{Highlighting}[]
\NormalTok{fair\_spread\_isda }\OtherTok{\textless{}{-}}\NormalTok{ GaussQuant}\SpecialCharTok{::}\FunctionTok{cds\_fair\_spread\_GQL}\NormalTok{(}
\NormalTok{  cds\_isda}
\NormalTok{)}

\NormalTok{fair\_upfront\_isda }\OtherTok{\textless{}{-}} \FunctionTok{tryCatch}\NormalTok{(}
  \FunctionTok{as.numeric}\NormalTok{(}
\NormalTok{    cds\_isda}\SpecialCharTok{$}\FunctionTok{fairUpfront}\NormalTok{()}
\NormalTok{  ),}
  \AttributeTok{error =} \ControlFlowTok{function}\NormalTok{(e) }\ConstantTok{NA\_real\_}
\NormalTok{)}

\NormalTok{fair\_upfront\_amount\_isda }\OtherTok{\textless{}{-}}\NormalTok{ fair\_upfront\_isda }\SpecialCharTok{*}
\NormalTok{  notional\_isda}

\NormalTok{isda\_cds\_summary\_tbl }\OtherTok{\textless{}{-}}\NormalTok{ tibble}\SpecialCharTok{::}\FunctionTok{tribble}\NormalTok{(}
  \SpecialCharTok{\textasciitilde{}}\NormalTok{metric,                  }\SpecialCharTok{\textasciitilde{}}\NormalTok{value,}
  \StringTok{"coupon\_leg\_npv"}\NormalTok{,          coupon\_leg\_npv\_isda,}
  \StringTok{"default\_leg\_npv"}\NormalTok{,         default\_leg\_npv\_isda,}
  \StringTok{"accrual\_rebate\_npv"}\NormalTok{,      accrual\_rebate\_npv\_isda,}
  \StringTok{"cds\_npv"}\NormalTok{,                 GaussQuant}\SpecialCharTok{::}\FunctionTok{cds\_npv\_GQL}\NormalTok{(}
\NormalTok{                               cds\_isda}
\NormalTok{                             ),}
  \StringTok{"market\_spread"}\NormalTok{,           market\_spread\_isda,}
  \StringTok{"contract\_coupon"}\NormalTok{,         contract\_coupon\_isda,}
  \StringTok{"fair\_spread"}\NormalTok{,             fair\_spread\_isda,}
  \StringTok{"fair\_upfront"}\NormalTok{,            fair\_upfront\_isda,}
  \StringTok{"fair\_upfront\_amount"}\NormalTok{,     fair\_upfront\_amount\_isda,}
  \StringTok{"hazard\_rate"}\NormalTok{,             hazard\_rate\_isda}
\NormalTok{)}

\NormalTok{GaussQuant}\SpecialCharTok{::}\FunctionTok{show\_tbl\_GQH}\NormalTok{(}
\NormalTok{  isda\_cds\_summary\_tbl,}
  \StringTok{"CDS summary: ISDA engine"}\NormalTok{,}
  \AttributeTok{n =} \DecValTok{30}
\NormalTok{)}
\CommentTok{\#\textgreater{} \# A tibble: 10 x 2}
\CommentTok{\#\textgreater{}    metric                      value}
\CommentTok{\#\textgreater{}    \textless{}chr\textgreater{}                       \textless{}dbl\textgreater{}}
\CommentTok{\#\textgreater{}  1 coupon\_leg\_npv      {-}498392.     }
\CommentTok{\#\textgreater{}  2 default\_leg\_npv      146485.     }
\CommentTok{\#\textgreater{}  3 accrual\_rebate\_npv    16944.     }
\CommentTok{\#\textgreater{}  4 cds\_npv             {-}334963.     }
\CommentTok{\#\textgreater{}  5 market\_spread             0.00304}
\CommentTok{\#\textgreater{}  6 contract\_coupon           0.01   }
\CommentTok{\#\textgreater{}  7 fair\_spread               0.00304}
\CommentTok{\#\textgreater{}  8 fair\_upfront             {-}0.0335 }
\CommentTok{\#\textgreater{}  9 fair\_upfront\_amount {-}334974.     }
\CommentTok{\#\textgreater{} 10 hazard\_rate               0.00514}
\end{Highlighting}
\end{Shaded}

この例では、契約Couponが市場Spreadより高いため、プロテクション買い手は市場より多くのPremiumを支払う。したがって、契約開始時の価値を調整するUpfrontが必要となる。

Fair
Upfrontの正負は、QuantLibにおけるBuyer／Sellerの符号規約と契約設定に依存する。符号だけで判断せず、CDS
NPV、契約Coupon、Fair
SpreadおよびUpfront金額を一緒に確認する必要がある。

\subsection{8.6.4
ISDA用キャッシュフロー表}\label{isdaux7528ux30adux30e3ux30c3ux30b7ux30e5ux30d5ux30edux30fcux8868}

標準CDS2015 Schedule、ISDA整合Hazard
Curveおよび割引カーブから、Premium期間ごとの確率情報を作成する。

\begin{Shaded}
\begin{Highlighting}[]
\NormalTok{cds\_cashflow\_isda\_tbl }\OtherTok{\textless{}{-}}\NormalTok{ GaussQuant}\SpecialCharTok{::}\FunctionTok{cds\_cashflow\_table\_GQL}\NormalTok{(}
  \AttributeTok{cds\_schedule =}\NormalTok{ cds\_schedule\_isda,}
  \AttributeTok{protection\_start\_date =}\NormalTok{ cds\_isda}\SpecialCharTok{$}\FunctionTok{protectionStartDate}\NormalTok{(),}
  \AttributeTok{coupon\_rate =}\NormalTok{ contract\_coupon\_isda,}
  \AttributeTok{hazard\_curve =}\NormalTok{ hazard\_curve\_isda,}
  \AttributeTok{discount\_curve =}\NormalTok{ discount\_curve\_isda,}
  \AttributeTok{trade\_date =}\NormalTok{ trade\_date\_isda,}
  \AttributeTok{notional =}\NormalTok{ notional\_isda}
\NormalTok{)}

\NormalTok{GaussQuant}\SpecialCharTok{::}\FunctionTok{show\_tbl\_GQH}\NormalTok{(}
\NormalTok{  cds\_cashflow\_isda\_tbl,}
  \StringTok{"CDS cash{-}flow table: CDS2015 example"}\NormalTok{,}
  \AttributeTok{n =} \DecValTok{30}
\NormalTok{)}
\CommentTok{\#\textgreater{} \# A tibble: 21 x 13}
\CommentTok{\#\textgreater{}    payment\_date coupon\_rate accrual\_start accrual\_end accrual\_days}
\CommentTok{\#\textgreater{}    \textless{}chr\textgreater{}              \textless{}dbl\textgreater{} \textless{}chr\textgreater{}         \textless{}chr\textgreater{}              \textless{}dbl\textgreater{}}
\CommentTok{\#\textgreater{}  1 2022{-}08{-}19         NA    \textless{}NA\textgreater{}          2022{-}08{-}19            NA}
\CommentTok{\#\textgreater{}  2 2022{-}09{-}20          0.01 2022{-}06{-}20    2022{-}09{-}20            92}
\CommentTok{\#\textgreater{}  3 2022{-}12{-}20          0.01 2022{-}09{-}20    2022{-}12{-}20            91}
\CommentTok{\#\textgreater{}  4 2023{-}03{-}20          0.01 2022{-}12{-}20    2023{-}03{-}20            90}
\CommentTok{\#\textgreater{}  5 2023{-}06{-}20          0.01 2023{-}03{-}20    2023{-}06{-}20            92}
\CommentTok{\#\textgreater{}  6 2023{-}09{-}20          0.01 2023{-}06{-}20    2023{-}09{-}20            92}
\CommentTok{\#\textgreater{}  7 2023{-}12{-}20          0.01 2023{-}09{-}20    2023{-}12{-}20            91}
\CommentTok{\#\textgreater{}  8 2024{-}03{-}20          0.01 2023{-}12{-}20    2024{-}03{-}20            91}
\CommentTok{\#\textgreater{}  9 2024{-}06{-}20          0.01 2024{-}03{-}20    2024{-}06{-}20            92}
\CommentTok{\#\textgreater{} 10 2024{-}09{-}20          0.01 2024{-}06{-}20    2024{-}09{-}20            92}
\CommentTok{\#\textgreater{} \# i 11 more rows}
\CommentTok{\#\textgreater{} \# i 8 more variables: accrual\_year\_fraction\_365 \textless{}dbl\textgreater{}, coupon\_amount \textless{}dbl\textgreater{},}
\CommentTok{\#\textgreater{} \#   discount\_factor \textless{}dbl\textgreater{}, survival\_probability \textless{}dbl\textgreater{},}
\CommentTok{\#\textgreater{} \#   default\_probability \textless{}dbl\textgreater{}, midpoint\_date \textless{}chr\textgreater{},}
\CommentTok{\#\textgreater{} \#   midpoint\_discount\_factor \textless{}dbl\textgreater{}, midpoint\_survival\_probability \textless{}dbl\textgreater{}}
\end{Highlighting}
\end{Shaded}

この表はCDSの経済構造を理解するためのMidPoint型分解である。ISDA
Engineは、標準日付、Accrual on
default、数値積分および曲線補間をより厳密に扱うため、この表から得られる単純近似と完全には一致しない。

\subsection{8.6.5 MidPoint型近似とISDA
Engine}\label{midpointux578bux8fd1ux4f3cux3068isda-engine}

ISDA用キャッシュフロー表から、Premium Leg、Protection
Leg、RPV01およびFair Spreadを簡易計算する。

\begin{Shaded}
\begin{Highlighting}[]
\NormalTok{premium\_leg\_hand\_isda }\OtherTok{\textless{}{-}} \SpecialCharTok{{-}}\FunctionTok{sum}\NormalTok{(}
\NormalTok{  cds\_cashflow\_isda\_tbl}\SpecialCharTok{$}\NormalTok{coupon\_amount }\SpecialCharTok{*}
\NormalTok{    cds\_cashflow\_isda\_tbl}\SpecialCharTok{$}\NormalTok{discount\_factor }\SpecialCharTok{*}
\NormalTok{    cds\_cashflow\_isda\_tbl}\SpecialCharTok{$}\NormalTok{midpoint\_survival\_probability,}
  \AttributeTok{na.rm =} \ConstantTok{TRUE}
\NormalTok{)}

\NormalTok{protection\_leg\_per\_unit\_hand\_isda }\OtherTok{\textless{}{-}}\NormalTok{ (}
  \DecValTok{1} \SpecialCharTok{{-}}
\NormalTok{    recovery\_rate\_isda}
\NormalTok{) }\SpecialCharTok{*}
  \FunctionTok{sum}\NormalTok{(}
\NormalTok{    cds\_cashflow\_isda\_tbl}\SpecialCharTok{$}\NormalTok{midpoint\_discount\_factor }\SpecialCharTok{*}
\NormalTok{      cds\_cashflow\_isda\_tbl}\SpecialCharTok{$}\NormalTok{default\_probability,}
    \AttributeTok{na.rm =} \ConstantTok{TRUE}
\NormalTok{  )}

\NormalTok{protection\_leg\_hand\_isda }\OtherTok{\textless{}{-}}\NormalTok{ notional\_isda }\SpecialCharTok{*}
\NormalTok{  protection\_leg\_per\_unit\_hand\_isda}

\NormalTok{risky\_pv01\_hand\_isda }\OtherTok{\textless{}{-}} \FunctionTok{sum}\NormalTok{(}
\NormalTok{  (}
\NormalTok{    cds\_cashflow\_isda\_tbl}\SpecialCharTok{$}\NormalTok{accrual\_days }\SpecialCharTok{/}
      \DecValTok{360}
\NormalTok{  ) }\SpecialCharTok{*}
\NormalTok{    cds\_cashflow\_isda\_tbl}\SpecialCharTok{$}\NormalTok{discount\_factor }\SpecialCharTok{*}
\NormalTok{    cds\_cashflow\_isda\_tbl}\SpecialCharTok{$}\NormalTok{midpoint\_survival\_probability,}
  \AttributeTok{na.rm =} \ConstantTok{TRUE}
\NormalTok{)}

\NormalTok{fair\_spread\_hand\_isda }\OtherTok{\textless{}{-}}\NormalTok{ protection\_leg\_per\_unit\_hand\_isda }\SpecialCharTok{/}
\NormalTok{  risky\_pv01\_hand\_isda}

\NormalTok{cds\_npv\_hand\_isda }\OtherTok{\textless{}{-}}\NormalTok{ premium\_leg\_hand\_isda }\SpecialCharTok{+}
\NormalTok{  protection\_leg\_hand\_isda }\SpecialCharTok{+}
\NormalTok{  accrual\_rebate\_npv\_isda}

\NormalTok{isda\_hand\_check\_tbl }\OtherTok{\textless{}{-}}\NormalTok{ tibble}\SpecialCharTok{::}\FunctionTok{tribble}\NormalTok{(}
  \SpecialCharTok{\textasciitilde{}}\NormalTok{metric,              }\SpecialCharTok{\textasciitilde{}}\NormalTok{hand\_value,                 }\SpecialCharTok{\textasciitilde{}}\NormalTok{isda\_value,}
  \StringTok{"coupon\_leg\_npv"}\NormalTok{,      premium\_leg\_hand\_isda,        coupon\_leg\_npv\_isda,}
  \StringTok{"default\_leg\_npv"}\NormalTok{,     protection\_leg\_hand\_isda,     default\_leg\_npv\_isda,}
  \StringTok{"cds\_npv"}\NormalTok{,             cds\_npv\_hand\_isda,            GaussQuant}\SpecialCharTok{::}\FunctionTok{cds\_npv\_GQL}\NormalTok{(}
\NormalTok{                                                        cds\_isda}
\NormalTok{                                                      ),}
  \StringTok{"fair\_spread"}\NormalTok{,         fair\_spread\_hand\_isda,        fair\_spread\_isda}
\NormalTok{) }\SpecialCharTok{|\textgreater{}}
\NormalTok{  dplyr}\SpecialCharTok{::}\FunctionTok{mutate}\NormalTok{(}
    \AttributeTok{difference =}\NormalTok{ .data}\SpecialCharTok{$}\NormalTok{hand\_value }\SpecialCharTok{{-}}
\NormalTok{      .data}\SpecialCharTok{$}\NormalTok{isda\_value}
\NormalTok{  )}

\NormalTok{GaussQuant}\SpecialCharTok{::}\FunctionTok{show\_tbl\_GQH}\NormalTok{(}
\NormalTok{  isda\_hand\_check\_tbl,}
  \StringTok{"MidPoint{-}style approximation and ISDA comparison"}\NormalTok{,}
  \AttributeTok{n =} \DecValTok{20}
\NormalTok{)}
\CommentTok{\#\textgreater{} \# A tibble: 4 x 4}
\CommentTok{\#\textgreater{}   metric             hand\_value    isda\_value difference}
\CommentTok{\#\textgreater{}   \textless{}chr\textgreater{}                   \textless{}dbl\textgreater{}         \textless{}dbl\textgreater{}      \textless{}dbl\textgreater{}}
\CommentTok{\#\textgreater{} 1 coupon\_leg\_npv  {-}498112.      {-}498392.      280.      }
\CommentTok{\#\textgreater{} 2 default\_leg\_npv  146486.       146485.        0.194   }
\CommentTok{\#\textgreater{} 3 cds\_npv         {-}334682.      {-}334963.      280.      }
\CommentTok{\#\textgreater{} 4 fair\_spread           0.00294       0.00304  {-}0.000102}
\end{Highlighting}
\end{Shaded}

\subsection{8.6.6 MidPoint EngineとISDA
Engineの比較}\label{midpoint-engineux3068isda-engineux306eux6bd4ux8f03}

同じHazard Curve、Recovery
Rate、割引カーブおよび標準Couponを使用し、MidPoint EngineとISDA
Engineの評価差を確認する。

\begin{Shaded}
\begin{Highlighting}[]
\NormalTok{midpoint\_engine\_isda\_inputs }\OtherTok{\textless{}{-}}\NormalTok{ GaussQuant}\SpecialCharTok{::}\FunctionTok{cds\_midpoint\_engine\_GQL}\NormalTok{(}
  \AttributeTok{probability\_handle =}\NormalTok{ hazard\_curve\_handle\_isda,}
  \AttributeTok{recovery\_rate =}\NormalTok{ recovery\_rate\_isda,}
  \AttributeTok{discount\_handle =}\NormalTok{ discount\_curve\_handle\_isda}
\NormalTok{)}

\NormalTok{cds\_midpoint\_standard }\OtherTok{\textless{}{-}}\NormalTok{ GaussQuant}\SpecialCharTok{::}\FunctionTok{make\_cds\_GQL}\NormalTok{(}
  \AttributeTok{maturity\_date =}\NormalTok{ maturity\_date\_isda,}
  \AttributeTok{running\_spread =}\NormalTok{ contract\_coupon\_isda,}
  \AttributeTok{notional =}\NormalTok{ notional\_isda,}
  \AttributeTok{side =} \StringTok{"buyer"}\NormalTok{,}
  \AttributeTok{trade\_date =}\NormalTok{ trade\_date\_isda,}
  \AttributeTok{coupon\_tenor =} \StringTok{"3M"}\NormalTok{,}
  \AttributeTok{day\_counter =} \StringTok{"Actual360"}\NormalTok{,}
  \AttributeTok{date\_generation\_rule =} \StringTok{"CDS2015"}\NormalTok{,}
  \AttributeTok{pricing\_engine =}\NormalTok{ midpoint\_engine\_isda\_inputs}
\NormalTok{)}

\NormalTok{midpoint\_standard\_summary\_tbl }\OtherTok{\textless{}{-}}\NormalTok{ GaussQuant}\SpecialCharTok{::}\FunctionTok{cds\_summary\_GQL}\NormalTok{(}
\NormalTok{  cds\_midpoint\_standard}
\NormalTok{) }\SpecialCharTok{|\textgreater{}}
\NormalTok{  dplyr}\SpecialCharTok{::}\FunctionTok{mutate}\NormalTok{(}
    \AttributeTok{engine =} \StringTok{"MidPoint"}
\NormalTok{  )}

\NormalTok{isda\_standard\_summary\_tbl }\OtherTok{\textless{}{-}}\NormalTok{ GaussQuant}\SpecialCharTok{::}\FunctionTok{cds\_summary\_GQL}\NormalTok{(}
\NormalTok{  cds\_isda}
\NormalTok{) }\SpecialCharTok{|\textgreater{}}
\NormalTok{  dplyr}\SpecialCharTok{::}\FunctionTok{mutate}\NormalTok{(}
    \AttributeTok{engine =} \StringTok{"ISDA"}
\NormalTok{  )}

\NormalTok{engine\_comparison\_tbl }\OtherTok{\textless{}{-}}\NormalTok{ dplyr}\SpecialCharTok{::}\FunctionTok{bind\_rows}\NormalTok{(}
\NormalTok{  midpoint\_standard\_summary\_tbl,}
\NormalTok{  isda\_standard\_summary\_tbl}
\NormalTok{) }\SpecialCharTok{|\textgreater{}}
\NormalTok{  dplyr}\SpecialCharTok{::}\FunctionTok{relocate}\NormalTok{(}
\NormalTok{    .data}\SpecialCharTok{$}\NormalTok{engine}
\NormalTok{  ) }\SpecialCharTok{|\textgreater{}}
\NormalTok{  tidyr}\SpecialCharTok{::}\FunctionTok{pivot\_wider}\NormalTok{(}
    \AttributeTok{names\_from =}\NormalTok{ engine,}
    \AttributeTok{values\_from =}\NormalTok{ value}
\NormalTok{  ) }\SpecialCharTok{|\textgreater{}}
\NormalTok{  dplyr}\SpecialCharTok{::}\FunctionTok{mutate}\NormalTok{(}
    \AttributeTok{difference\_isda\_minus\_midpoint =}\NormalTok{ .data}\SpecialCharTok{$}\NormalTok{ISDA }\SpecialCharTok{{-}}
\NormalTok{      .data}\SpecialCharTok{$}\NormalTok{MidPoint}
\NormalTok{  )}

\NormalTok{GaussQuant}\SpecialCharTok{::}\FunctionTok{show\_tbl\_GQH}\NormalTok{(}
\NormalTok{  engine\_comparison\_tbl,}
  \StringTok{"MidPoint and ISDA CDS engine comparison"}\NormalTok{,}
  \AttributeTok{n =} \DecValTok{20}
\NormalTok{)}
\CommentTok{\#\textgreater{} \# A tibble: 4 x 4}
\CommentTok{\#\textgreater{}   metric               MidPoint          ISDA difference\_isda\_minus\_midpoint}
\CommentTok{\#\textgreater{}   \textless{}chr\textgreater{}                   \textless{}dbl\textgreater{}         \textless{}dbl\textgreater{}                          \textless{}dbl\textgreater{}}
\CommentTok{\#\textgreater{} 1 npv             {-}334951.      {-}334963.                            {-}1.17e+1}
\CommentTok{\#\textgreater{} 2 fair\_spread           0.00304       0.00304                       {-}7.66e{-}8}
\CommentTok{\#\textgreater{} 3 coupon\_leg\_npv  {-}498380.      {-}498392.                            {-}1.15e+1}
\CommentTok{\#\textgreater{} 4 default\_leg\_npv  146486.       146485.                            {-}1.94e{-}1}
\end{Highlighting}
\end{Shaded}

MidPoint法は、各Premium期間内のデフォルト時点を期間中点へ集約する。ISDA
Engineは標準市場規約に合わせて、より詳細な日付処理と積分を行う。そのため、同じカーブを用いてもLeg
NPVとFair Spreadに差が生じる。

\section{8.7 まとめ}\label{ux307eux3068ux3081-12}

CDSは、Premium LegとProtection
Legから構成される。プロテクション買い手はPremiumを支払い、信用事由が発生した場合に損失補償を受け取る。

Hazard
Rateは、企業がその時点まで存続しているという条件の下での瞬間的なデフォルト強度である。フラットHazard
Rateでは、生存確率は次式で表される。

\[
Q(t)
=
e^{-\lambda t}
\]

Recovery Rate \(R\) を仮定すると、CDS SpreadとHazard Rateには、

\[
\lambda
\approx
\frac{s}{1-R}
\]

という近似関係がある。ただし、厳密なHazard
Rateは、Premium支払日、割引、Accrual on
defaultおよび標準CDS規約を反映して推定する。

Premium Legの単位Spread当たり現在価値がRisky PV01であり、Fair
SpreadはProtection Legの単位Notional価値をRisky PV01で割って求められる。

MidPoint
Engineは、デフォルト時点を各期間の中点で近似するため、CDS評価の構造を理解しやすい。ISDA
Engineは、標準CDS2015の日付、Accrual rebate、Accrual on
defaultおよび市場規約をより厳密に扱う。

標準Couponと市場Spreadが異なる場合、Running
Couponだけでは契約価値がゼロにならない。この差を契約開始時に調整するものがFair
Upfrontである。

\begin{Shaded}
\begin{Highlighting}[]
\NormalTok{chapter08\_summary\_tbl }\OtherTok{\textless{}{-}}\NormalTok{ tibble}\SpecialCharTok{::}\FunctionTok{tribble}\NormalTok{(}
  \SpecialCharTok{\textasciitilde{}}\NormalTok{section,                     }\SpecialCharTok{\textasciitilde{}}\NormalTok{main\_result,}
  \StringTok{"Hazard Rate"}\NormalTok{,                }\StringTok{"Survival and cumulative default probabilities"}\NormalTok{,}
  \StringTok{"MidPoint CDS"}\NormalTok{,               }\StringTok{"Premium and protection leg decomposition"}\NormalTok{,}
  \StringTok{"Hazard calibration"}\NormalTok{,         }\StringTok{"Zero{-}NPV market CDS"}\NormalTok{,}
  \StringTok{"CDS2015"}\NormalTok{,                    }\StringTok{"Standard maturity and coupon schedule"}\NormalTok{,}
  \StringTok{"ISDA CDS"}\NormalTok{,                   }\StringTok{"Fair spread, upfront, and accrual rebate"}\NormalTok{,}
  \StringTok{"Engine comparison"}\NormalTok{,          }\StringTok{"MidPoint approximation versus ISDA pricing"}
\NormalTok{)}

\NormalTok{GaussQuant}\SpecialCharTok{::}\FunctionTok{show\_tbl\_GQH}\NormalTok{(}
\NormalTok{  chapter08\_summary\_tbl,}
  \StringTok{"Chapter II{-}8 summary"}\NormalTok{,}
  \AttributeTok{n =} \DecValTok{20}
\NormalTok{)}
\CommentTok{\#\textgreater{} \# A tibble: 6 x 2}
\CommentTok{\#\textgreater{}   section            main\_result                                  }
\CommentTok{\#\textgreater{}   \textless{}chr\textgreater{}              \textless{}chr\textgreater{}                                        }
\CommentTok{\#\textgreater{} 1 Hazard Rate        Survival and cumulative default probabilities}
\CommentTok{\#\textgreater{} 2 MidPoint CDS       Premium and protection leg decomposition     }
\CommentTok{\#\textgreater{} 3 Hazard calibration Zero{-}NPV market CDS                          }
\CommentTok{\#\textgreater{} 4 CDS2015            Standard maturity and coupon schedule        }
\CommentTok{\#\textgreater{} 5 ISDA CDS           Fair spread, upfront, and accrual rebate     }
\CommentTok{\#\textgreater{} 6 Engine comparison  MidPoint approximation versus ISDA pricing}

\FunctionTok{cat}\NormalTok{(}
  \StringTok{"}\SpecialCharTok{\textbackslash{}n}\StringTok{Chapter II{-}8 Credit risk and CDS completed successfully.}\SpecialCharTok{\textbackslash{}n}\StringTok{"}
\NormalTok{)}
\CommentTok{\#\textgreater{} }
\CommentTok{\#\textgreater{} Chapter II{-}8 Credit risk and CDS completed successfully.}
\end{Highlighting}
\end{Shaded}

\chapter{第II-9章
エキゾチック・デリバティブとPRDC}\label{ux7b2cii-9ux7ae0-ux30a8ux30adux30beux30c1ux30c3ux30afux30c7ux30eaux30d0ux30c6ux30a3ux30d6ux3068prdc}

これまでの章では、債券、金利先物、Swap、Option、Callable
BondおよびCDSを、それぞれの主要なリスク要因に基づいて評価した。本章では、第II部の締めくくりとして、金利と為替を組み合わせたエキゾチック・デリバティブであるPRDCを検討する。

PRDC（Power Reverse Dual Currency
Note）は、為替レートに連動するCouponを支払う長期の仕組債である。典型的なCoupon
rateは、次の形で表される。

\[
C(S_t)
=
\min
\left[
C_{\max},
\max
\left(
a\frac{S_t}{S_0}-b,
C_{\min}
\right)
\right]
\]

ここで、

\begin{itemize}
\item
  \(S_t\) はCoupon決定日における為替レート
\item
  \(S_0\) はCoupon計算上の基準為替レート
\item
  \(a\) は為替連動係数
\item
  \(b\) は固定控除率
\item
  \(C_{\min}\) はCoupon Floor
\item
  \(C_{\max}\) はCoupon Cap
\end{itemize}

である。

為替レートが上昇するとCoupon
rateは増加するが、FloorとCapによって下限と上限が制限される。したがって、PRDC
CouponはFX
Spotに対して区分線形である一方、将来価値は為替分布に依存する非線形商品となる。

本章では、Domestic金利、Foreign金利およびFX
volatilityを一定とするGarman--Kohlhagenモデルを用いて、将来のCoupon日ごとにFXパスを生成する。各パスでCouponを計算し、その平均値からCoupon
LegとRedemption Legの現在価値を求める。

実際のPRDCには、発行体の期限前償還権を含むことが多い。本章では商品の基本構造を明確にするため、\textbf{発行体コール権を含まない非Callable
PRDC}を評価する。Callable
PRDCには、金利・為替の複数因子モデル、相関構造およびLeast-Squares Monte
Carloなどが追加で必要となる。

\begin{Shaded}
\begin{Highlighting}[]
\NormalTok{knitr}\SpecialCharTok{::}\NormalTok{opts\_chunk}\SpecialCharTok{$}\FunctionTok{set}\NormalTok{(}
  \AttributeTok{collapse =} \ConstantTok{TRUE}\NormalTok{,}
  \AttributeTok{comment =} \StringTok{"\#\textgreater{}"}\NormalTok{,}
  \AttributeTok{warning =} \ConstantTok{FALSE}\NormalTok{,}
  \AttributeTok{message =} \ConstantTok{FALSE}
\NormalTok{)}


\FunctionTok{suppressPackageStartupMessages}\NormalTok{(\{}
  \FunctionTok{library}\NormalTok{(tidyverse)}
  \FunctionTok{library}\NormalTok{(ggthemes)}
  \FunctionTok{library}\NormalTok{(QuantLib)}
  \FunctionTok{library}\NormalTok{(GaussQuant)}
  \FunctionTok{library}\NormalTok{(tidyverse)}
  \FunctionTok{library}\NormalTok{(broom)}
\NormalTok{\})}
\end{Highlighting}
\end{Shaded}

\section{9.1
PRDCの商品構造と市場条件}\label{prdcux306eux5546ux54c1ux69cbux9020ux3068ux5e02ux5834ux6761ux4ef6}

\subsection{9.1.1 商品条件}\label{ux5546ux54c1ux6761ux4ef6}

2015年9月3日から2041年9月3日まで、半年ごとにCouponを支払うPRDCを考える。最初の1年間は固定Couponとし、その後はFX連動Couponへ移行する。

評価日は2023年4月11日であるため、固定Intro
Coupon期間はすでに終了している。ただし、GaussQuantのPRDC評価関数がIntro
Couponを扱えることを確認するため、契約条件として残す。

本章では、計算時間を抑えながらMonte
Carlo評価を確認するため、基準評価のパス数を3,000本、感応度評価を1,000本とする。

\begin{Shaded}
\begin{Highlighting}[]
\NormalTok{valuation\_date\_prdc }\OtherTok{\textless{}{-}} \FunctionTok{as.Date}\NormalTok{(}\StringTok{"2023{-}04{-}11"}\NormalTok{)}
\NormalTok{effective\_date\_prdc }\OtherTok{\textless{}{-}} \FunctionTok{as.Date}\NormalTok{(}\StringTok{"2015{-}09{-}03"}\NormalTok{)}
\NormalTok{termination\_date\_prdc }\OtherTok{\textless{}{-}} \FunctionTok{as.Date}\NormalTok{(}\StringTok{"2041{-}09{-}03"}\NormalTok{)}
\NormalTok{intro\_coupon\_end\_date\_prdc }\OtherTok{\textless{}{-}} \FunctionTok{as.Date}\NormalTok{(}\StringTok{"2016{-}09{-}03"}\NormalTok{)}

\NormalTok{notional\_prdc }\OtherTok{\textless{}{-}} \DecValTok{300000000}
\NormalTok{intro\_coupon\_rate\_prdc }\OtherTok{\textless{}{-}} \FloatTok{0.022}
\NormalTok{coupon\_frequency\_months\_prdc }\OtherTok{\textless{}{-}} \DecValTok{6}\NormalTok{L}

\NormalTok{domestic\_rate\_prdc }\OtherTok{\textless{}{-}} \FloatTok{0.01}
\NormalTok{foreign\_rate\_prdc }\OtherTok{\textless{}{-}} \FloatTok{0.03}
\NormalTok{discount\_rate\_prdc }\OtherTok{\textless{}{-}} \FloatTok{0.005}

\NormalTok{fx\_spot\_prdc }\OtherTok{\textless{}{-}} \FloatTok{133.2681}
\NormalTok{fx\_volatility\_prdc }\OtherTok{\textless{}{-}} \FloatTok{0.10}

\NormalTok{coupon\_fx\_base\_prdc }\OtherTok{\textless{}{-}} \DecValTok{120}
\NormalTok{coupon\_leverage\_prdc }\OtherTok{\textless{}{-}} \FloatTok{0.122}
\NormalTok{coupon\_subtraction\_prdc }\OtherTok{\textless{}{-}} \FloatTok{0.10}
\NormalTok{coupon\_floor\_prdc }\OtherTok{\textless{}{-}} \DecValTok{0}
\NormalTok{coupon\_cap\_prdc }\OtherTok{\textless{}{-}} \FloatTok{0.022}

\NormalTok{number\_of\_paths\_prdc }\OtherTok{\textless{}{-}} \DecValTok{3000}\NormalTok{L}
\NormalTok{sensitivity\_paths\_prdc }\OtherTok{\textless{}{-}} \DecValTok{1000}\NormalTok{L}
\NormalTok{random\_seed\_prdc }\OtherTok{\textless{}{-}} \DecValTok{123}\NormalTok{L}

\NormalTok{GaussQuant}\SpecialCharTok{::}\FunctionTok{set\_eval\_date\_GQL}\NormalTok{(}
\NormalTok{  valuation\_date\_prdc}
\NormalTok{)}

\NormalTok{prdc\_input\_tbl }\OtherTok{\textless{}{-}}\NormalTok{ tibble}\SpecialCharTok{::}\FunctionTok{tribble}\NormalTok{(}
  \SpecialCharTok{\textasciitilde{}}\NormalTok{input,                    }\SpecialCharTok{\textasciitilde{}}\NormalTok{value,}
  \StringTok{"valuation\_date"}\NormalTok{,           }\FunctionTok{as.character}\NormalTok{(valuation\_date\_prdc),}
  \StringTok{"effective\_date"}\NormalTok{,           }\FunctionTok{as.character}\NormalTok{(effective\_date\_prdc),}
  \StringTok{"termination\_date"}\NormalTok{,         }\FunctionTok{as.character}\NormalTok{(termination\_date\_prdc),}
  \StringTok{"intro\_coupon\_end\_date"}\NormalTok{,    }\FunctionTok{as.character}\NormalTok{(intro\_coupon\_end\_date\_prdc),}
  \StringTok{"notional\_jpy"}\NormalTok{,             }\FunctionTok{as.character}\NormalTok{(notional\_prdc),}
  \StringTok{"intro\_coupon\_rate"}\NormalTok{,        }\FunctionTok{as.character}\NormalTok{(intro\_coupon\_rate\_prdc),}
  \StringTok{"coupon\_frequency\_months"}\NormalTok{,  }\FunctionTok{as.character}\NormalTok{(coupon\_frequency\_months\_prdc),}
  \StringTok{"domestic\_rate"}\NormalTok{,            }\FunctionTok{as.character}\NormalTok{(domestic\_rate\_prdc),}
  \StringTok{"foreign\_rate"}\NormalTok{,             }\FunctionTok{as.character}\NormalTok{(foreign\_rate\_prdc),}
  \StringTok{"discount\_rate"}\NormalTok{,            }\FunctionTok{as.character}\NormalTok{(discount\_rate\_prdc),}
  \StringTok{"fx\_spot"}\NormalTok{,                  }\FunctionTok{as.character}\NormalTok{(fx\_spot\_prdc),}
  \StringTok{"fx\_volatility"}\NormalTok{,            }\FunctionTok{as.character}\NormalTok{(fx\_volatility\_prdc),}
  \StringTok{"coupon\_fx\_base"}\NormalTok{,           }\FunctionTok{as.character}\NormalTok{(coupon\_fx\_base\_prdc),}
  \StringTok{"coupon\_leverage"}\NormalTok{,          }\FunctionTok{as.character}\NormalTok{(coupon\_leverage\_prdc),}
  \StringTok{"coupon\_subtraction"}\NormalTok{,       }\FunctionTok{as.character}\NormalTok{(coupon\_subtraction\_prdc),}
  \StringTok{"coupon\_floor"}\NormalTok{,             }\FunctionTok{as.character}\NormalTok{(coupon\_floor\_prdc),}
  \StringTok{"coupon\_cap"}\NormalTok{,               }\FunctionTok{as.character}\NormalTok{(coupon\_cap\_prdc),}
  \StringTok{"valuation\_paths"}\NormalTok{,          }\FunctionTok{as.character}\NormalTok{(number\_of\_paths\_prdc),}
  \StringTok{"sensitivity\_paths"}\NormalTok{,        }\FunctionTok{as.character}\NormalTok{(sensitivity\_paths\_prdc),}
  \StringTok{"random\_seed"}\NormalTok{,              }\FunctionTok{as.character}\NormalTok{(random\_seed\_prdc)}
\NormalTok{)}

\NormalTok{GaussQuant}\SpecialCharTok{::}\FunctionTok{show\_tbl\_GQH}\NormalTok{(}
\NormalTok{  prdc\_input\_tbl,}
  \StringTok{"PRDC product and model assumptions"}\NormalTok{,}
  \AttributeTok{n =} \DecValTok{30}
\NormalTok{)}
\CommentTok{\#\textgreater{} \# A tibble: 20 x 2}
\CommentTok{\#\textgreater{}    input                   value     }
\CommentTok{\#\textgreater{}    \textless{}chr\textgreater{}                   \textless{}chr\textgreater{}     }
\CommentTok{\#\textgreater{}  1 valuation\_date          2023{-}04{-}11}
\CommentTok{\#\textgreater{}  2 effective\_date          2015{-}09{-}03}
\CommentTok{\#\textgreater{}  3 termination\_date        2041{-}09{-}03}
\CommentTok{\#\textgreater{}  4 intro\_coupon\_end\_date   2016{-}09{-}03}
\CommentTok{\#\textgreater{}  5 notional\_jpy            3e+08     }
\CommentTok{\#\textgreater{}  6 intro\_coupon\_rate       0.022     }
\CommentTok{\#\textgreater{}  7 coupon\_frequency\_months 6         }
\CommentTok{\#\textgreater{}  8 domestic\_rate           0.01      }
\CommentTok{\#\textgreater{}  9 foreign\_rate            0.03      }
\CommentTok{\#\textgreater{} 10 discount\_rate           0.005     }
\CommentTok{\#\textgreater{} 11 fx\_spot                 133.2681  }
\CommentTok{\#\textgreater{} 12 fx\_volatility           0.1       }
\CommentTok{\#\textgreater{} 13 coupon\_fx\_base          120       }
\CommentTok{\#\textgreater{} 14 coupon\_leverage         0.122     }
\CommentTok{\#\textgreater{} 15 coupon\_subtraction      0.1       }
\CommentTok{\#\textgreater{} 16 coupon\_floor            0         }
\CommentTok{\#\textgreater{} 17 coupon\_cap              0.022     }
\CommentTok{\#\textgreater{} 18 valuation\_paths         3000      }
\CommentTok{\#\textgreater{} 19 sensitivity\_paths       1000      }
\CommentTok{\#\textgreater{} 20 random\_seed             123}
\end{Highlighting}
\end{Shaded}

\subsection{9.1.2 Coupon Schedule}\label{coupon-schedule}

契約開始日から満期までの全Coupon Scheduleと、固定Intro
Coupon期間のScheduleを作成する。

\begin{Shaded}
\begin{Highlighting}[]
\NormalTok{valuation\_date\_ql\_prdc }\OtherTok{\textless{}{-}}\NormalTok{ GaussQuant}\SpecialCharTok{::}\FunctionTok{date\_GQL}\NormalTok{(}
\NormalTok{  valuation\_date\_prdc}
\NormalTok{)}

\NormalTok{effective\_date\_ql\_prdc }\OtherTok{\textless{}{-}}\NormalTok{ GaussQuant}\SpecialCharTok{::}\FunctionTok{date\_GQL}\NormalTok{(}
\NormalTok{  effective\_date\_prdc}
\NormalTok{)}

\NormalTok{termination\_date\_ql\_prdc }\OtherTok{\textless{}{-}}\NormalTok{ GaussQuant}\SpecialCharTok{::}\FunctionTok{date\_GQL}\NormalTok{(}
\NormalTok{  termination\_date\_prdc}
\NormalTok{)}

\NormalTok{intro\_coupon\_end\_date\_ql\_prdc }\OtherTok{\textless{}{-}}\NormalTok{ GaussQuant}\SpecialCharTok{::}\FunctionTok{date\_GQL}\NormalTok{(}
\NormalTok{  intro\_coupon\_end\_date\_prdc}
\NormalTok{)}

\NormalTok{coupon\_schedule\_prdc }\OtherTok{\textless{}{-}}\NormalTok{ QuantLib}\SpecialCharTok{::}\FunctionTok{Schedule}\NormalTok{(}
\NormalTok{  effective\_date\_ql\_prdc,}
\NormalTok{  termination\_date\_ql\_prdc,}
\NormalTok{  GaussQuant}\SpecialCharTok{::}\FunctionTok{period\_months\_GQL}\NormalTok{(}
\NormalTok{    coupon\_frequency\_months\_prdc}
\NormalTok{  ),}
\NormalTok{  QuantLib}\SpecialCharTok{::}\FunctionTok{TARGET}\NormalTok{(),}
  \StringTok{"ModifiedFollowing"}\NormalTok{,}
  \StringTok{"ModifiedFollowing"}\NormalTok{,}
  \StringTok{"Backward"}\NormalTok{,}
  \ConstantTok{FALSE}
\NormalTok{)}

\NormalTok{intro\_coupon\_schedule\_prdc }\OtherTok{\textless{}{-}}\NormalTok{ QuantLib}\SpecialCharTok{::}\FunctionTok{Schedule}\NormalTok{(}
\NormalTok{  effective\_date\_ql\_prdc,}
\NormalTok{  intro\_coupon\_end\_date\_ql\_prdc,}
\NormalTok{  GaussQuant}\SpecialCharTok{::}\FunctionTok{period\_months\_GQL}\NormalTok{(}
\NormalTok{    coupon\_frequency\_months\_prdc}
\NormalTok{  ),}
\NormalTok{  QuantLib}\SpecialCharTok{::}\FunctionTok{TARGET}\NormalTok{(),}
  \StringTok{"ModifiedFollowing"}\NormalTok{,}
  \StringTok{"ModifiedFollowing"}\NormalTok{,}
  \StringTok{"Backward"}\NormalTok{,}
  \ConstantTok{FALSE}
\NormalTok{)}

\NormalTok{coupon\_schedule\_tbl }\OtherTok{\textless{}{-}}\NormalTok{ GaussQuant}\SpecialCharTok{::}\FunctionTok{schedule\_table\_GQL}\NormalTok{(}
\NormalTok{  coupon\_schedule\_prdc}
\NormalTok{)}

\NormalTok{intro\_coupon\_dates\_prdc }\OtherTok{\textless{}{-}}\NormalTok{ GaussQuant}\SpecialCharTok{::}\FunctionTok{schedule\_date\_vector\_GQL}\NormalTok{(}
\NormalTok{  intro\_coupon\_schedule\_prdc}
\NormalTok{)}

\NormalTok{remaining\_intro\_coupon\_dates\_prdc }\OtherTok{\textless{}{-}}\NormalTok{ intro\_coupon\_dates\_prdc[}
\NormalTok{  intro\_coupon\_dates\_prdc }\SpecialCharTok{\textgreater{}}
\NormalTok{    valuation\_date\_prdc}
\NormalTok{]}

\NormalTok{intro\_coupon\_status\_tbl }\OtherTok{\textless{}{-}}\NormalTok{ tibble}\SpecialCharTok{::}\FunctionTok{tribble}\NormalTok{(}
  \SpecialCharTok{\textasciitilde{}}\NormalTok{metric,                        }\SpecialCharTok{\textasciitilde{}}\NormalTok{value,}
  \StringTok{"intro\_coupon\_rate"}\NormalTok{,             intro\_coupon\_rate\_prdc,}
  \StringTok{"total\_intro\_schedule\_dates"}\NormalTok{,    }\FunctionTok{length}\NormalTok{(intro\_coupon\_dates\_prdc),}
  \StringTok{"remaining\_intro\_coupon\_dates"}\NormalTok{,  }\FunctionTok{length}\NormalTok{(remaining\_intro\_coupon\_dates\_prdc),}
  \StringTok{"has\_remaining\_intro\_coupon"}\NormalTok{,    }\FunctionTok{as.numeric}\NormalTok{(}
                                      \FunctionTok{length}\NormalTok{(}
\NormalTok{                                        remaining\_intro\_coupon\_dates\_prdc}
\NormalTok{                                      ) }\SpecialCharTok{\textgreater{}} \DecValTok{0}
\NormalTok{                                    )}
\NormalTok{)}

\NormalTok{GaussQuant}\SpecialCharTok{::}\FunctionTok{show\_tbl\_GQH}\NormalTok{(}
\NormalTok{  coupon\_schedule\_tbl,}
  \StringTok{"PRDC full coupon schedule"}\NormalTok{,}
  \AttributeTok{n =} \DecValTok{60}
\NormalTok{)}
\CommentTok{\#\textgreater{} \# A tibble: 53 x 2}
\CommentTok{\#\textgreater{}    period schedule\_date}
\CommentTok{\#\textgreater{}     \textless{}int\textgreater{} \textless{}date\textgreater{}       }
\CommentTok{\#\textgreater{}  1      1 2015{-}09{-}03   }
\CommentTok{\#\textgreater{}  2      2 2016{-}03{-}03   }
\CommentTok{\#\textgreater{}  3      3 2016{-}09{-}05   }
\CommentTok{\#\textgreater{}  4      4 2017{-}03{-}03   }
\CommentTok{\#\textgreater{}  5      5 2017{-}09{-}04   }
\CommentTok{\#\textgreater{}  6      6 2018{-}03{-}05   }
\CommentTok{\#\textgreater{}  7      7 2018{-}09{-}03   }
\CommentTok{\#\textgreater{}  8      8 2019{-}03{-}04   }
\CommentTok{\#\textgreater{}  9      9 2019{-}09{-}03   }
\CommentTok{\#\textgreater{} 10     10 2020{-}03{-}03   }
\CommentTok{\#\textgreater{} \# i 43 more rows}

\NormalTok{GaussQuant}\SpecialCharTok{::}\FunctionTok{show\_tbl\_GQH}\NormalTok{(}
\NormalTok{  intro\_coupon\_status\_tbl,}
  \StringTok{"PRDC introductory coupon status"}\NormalTok{,}
  \AttributeTok{n =} \DecValTok{20}
\NormalTok{)}
\CommentTok{\#\textgreater{} \# A tibble: 4 x 2}
\CommentTok{\#\textgreater{}   metric                       value}
\CommentTok{\#\textgreater{}   \textless{}chr\textgreater{}                        \textless{}dbl\textgreater{}}
\CommentTok{\#\textgreater{} 1 intro\_coupon\_rate            0.022}
\CommentTok{\#\textgreater{} 2 total\_intro\_schedule\_dates   3    }
\CommentTok{\#\textgreater{} 3 remaining\_intro\_coupon\_dates 0    }
\CommentTok{\#\textgreater{} 4 has\_remaining\_intro\_coupon   0}
\end{Highlighting}
\end{Shaded}

\subsection{9.1.3 Domestic、ForeignおよびDiscount
Curve}\label{domesticforeignux304aux3088ux3073discount-curve}

FX Quoteを、外国通貨1単位当たりの円貨額として表す。

\[
S_t
=
\text{JPY per unit of foreign currency}
\]

このとき、

\begin{itemize}
\tightlist
\item
  Domestic CurveはJPY金利
\item
  Foreign Curveは外国通貨金利
\end{itemize}

である。

Garman--Kohlhagenモデルのリスク中立ドリフトには、Domestic金利とForeign金利の差を使用する。

\[
\frac{dS_t}{S_t}
=
(r_d-r_f)dt
+
\sigma_{\mathrm{FX}}dW_t
\]

本例では、Domestic金利を1\%、Foreign金利を3\%とする。そのため、リスク中立測度の下でのFXドリフトは年率マイナス2\%となる。

Couponと元本の現在価値には、担保またはファンディング条件を表す外生的な割引率0.5\%を使用する。これは簡略化した設定であり、本格的な評価ではDomestic
OIS CurveとCollateral条件を整合させる必要がある。

\begin{Shaded}
\begin{Highlighting}[]
\NormalTok{fx\_day\_counter\_prdc }\OtherTok{\textless{}{-}}\NormalTok{ QuantLib}\SpecialCharTok{::}\FunctionTok{Actual360}\NormalTok{()}

\NormalTok{domestic\_curve\_prdc }\OtherTok{\textless{}{-}}\NormalTok{ GaussQuant}\SpecialCharTok{::}\FunctionTok{flat\_curve\_GQL}\NormalTok{(}
  \AttributeTok{reference\_date =}\NormalTok{ valuation\_date\_ql\_prdc,}
  \AttributeTok{rate =}\NormalTok{ domestic\_rate\_prdc,}
  \AttributeTok{day\_counter =}\NormalTok{ fx\_day\_counter\_prdc,}
  \AttributeTok{compounding =} \StringTok{"Continuous"}
\NormalTok{)}

\NormalTok{foreign\_curve\_prdc }\OtherTok{\textless{}{-}}\NormalTok{ GaussQuant}\SpecialCharTok{::}\FunctionTok{flat\_curve\_GQL}\NormalTok{(}
  \AttributeTok{reference\_date =}\NormalTok{ valuation\_date\_ql\_prdc,}
  \AttributeTok{rate =}\NormalTok{ foreign\_rate\_prdc,}
  \AttributeTok{day\_counter =}\NormalTok{ fx\_day\_counter\_prdc,}
  \AttributeTok{compounding =} \StringTok{"Continuous"}
\NormalTok{)}

\NormalTok{discount\_curve\_prdc }\OtherTok{\textless{}{-}}\NormalTok{ GaussQuant}\SpecialCharTok{::}\FunctionTok{flat\_curve\_GQL}\NormalTok{(}
  \AttributeTok{reference\_date =}\NormalTok{ valuation\_date\_ql\_prdc,}
  \AttributeTok{rate =}\NormalTok{ discount\_rate\_prdc,}
  \AttributeTok{day\_counter =}\NormalTok{ fx\_day\_counter\_prdc,}
  \AttributeTok{compounding =} \StringTok{"Continuous"}
\NormalTok{)}

\NormalTok{domestic\_curve\_handle\_prdc }\OtherTok{\textless{}{-}}\NormalTok{ QuantLib}\SpecialCharTok{::}\FunctionTok{YieldTermStructureHandle}\NormalTok{(}
\NormalTok{  domestic\_curve\_prdc}
\NormalTok{)}

\NormalTok{foreign\_curve\_handle\_prdc }\OtherTok{\textless{}{-}}\NormalTok{ QuantLib}\SpecialCharTok{::}\FunctionTok{YieldTermStructureHandle}\NormalTok{(}
\NormalTok{  foreign\_curve\_prdc}
\NormalTok{)}

\NormalTok{discount\_curve\_handle\_prdc }\OtherTok{\textless{}{-}}\NormalTok{ QuantLib}\SpecialCharTok{::}\FunctionTok{YieldTermStructureHandle}\NormalTok{(}
\NormalTok{  discount\_curve\_prdc}
\NormalTok{)}

\NormalTok{curve\_summary\_tbl }\OtherTok{\textless{}{-}}\NormalTok{ tibble}\SpecialCharTok{::}\FunctionTok{tribble}\NormalTok{(}
  \SpecialCharTok{\textasciitilde{}}\NormalTok{curve,                 }\SpecialCharTok{\textasciitilde{}}\NormalTok{flat\_rate,}
  \StringTok{"Domestic FX curve"}\NormalTok{,     domestic\_rate\_prdc,}
  \StringTok{"Foreign FX curve"}\NormalTok{,      foreign\_rate\_prdc,}
  \StringTok{"PRDC discount curve"}\NormalTok{,   discount\_rate\_prdc,}
  \StringTok{"FX risk{-}neutral drift"}\NormalTok{, domestic\_rate\_prdc }\SpecialCharTok{{-}}
\NormalTok{                             foreign\_rate\_prdc}
\NormalTok{)}

\NormalTok{GaussQuant}\SpecialCharTok{::}\FunctionTok{show\_tbl\_GQH}\NormalTok{(}
\NormalTok{  curve\_summary\_tbl,}
  \StringTok{"PRDC flat{-}curve assumptions"}\NormalTok{,}
  \AttributeTok{n =} \DecValTok{20}
\NormalTok{)}
\CommentTok{\#\textgreater{} \# A tibble: 4 x 2}
\CommentTok{\#\textgreater{}   curve                 flat\_rate}
\CommentTok{\#\textgreater{}   \textless{}chr\textgreater{}                     \textless{}dbl\textgreater{}}
\CommentTok{\#\textgreater{} 1 Domestic FX curve         0.01 }
\CommentTok{\#\textgreater{} 2 Foreign FX curve          0.03 }
\CommentTok{\#\textgreater{} 3 PRDC discount curve       0.005}
\CommentTok{\#\textgreater{} 4 FX risk{-}neutral drift    {-}0.02}
\end{Highlighting}
\end{Shaded}

\section{9.2 FXプロセスとCoupon
Payoff}\label{fxux30d7ux30edux30bbux30b9ux3068coupon-payoff}

\subsection{9.2.1 FX Process}\label{fx-process}

FX volatilityを10\%とし、Garman--Kohlhagen Processを作成する。

\begin{Shaded}
\begin{Highlighting}[]
\NormalTok{fx\_volatility\_curve\_prdc }\OtherTok{\textless{}{-}}\NormalTok{ QuantLib}\SpecialCharTok{::}\FunctionTok{BlackConstantVol}\NormalTok{(}
\NormalTok{  valuation\_date\_ql\_prdc,}
\NormalTok{  QuantLib}\SpecialCharTok{::}\FunctionTok{NullCalendar}\NormalTok{(),}
\NormalTok{  GaussQuant}\SpecialCharTok{::}\FunctionTok{quote\_handle\_GQL}\NormalTok{(}
\NormalTok{    fx\_volatility\_prdc}
\NormalTok{  ),}
\NormalTok{  fx\_day\_counter\_prdc}
\NormalTok{)}

\NormalTok{fx\_volatility\_handle\_prdc }\OtherTok{\textless{}{-}}\NormalTok{ QuantLib}\SpecialCharTok{::}\FunctionTok{BlackVolTermStructureHandle}\NormalTok{(}
\NormalTok{  fx\_volatility\_curve\_prdc}
\NormalTok{)}

\NormalTok{fx\_process\_prdc }\OtherTok{\textless{}{-}}\NormalTok{ GaussQuant}\SpecialCharTok{::}\FunctionTok{fx\_process\_GQL}\NormalTok{(}
  \AttributeTok{spot =}\NormalTok{ fx\_spot\_prdc,}
  \AttributeTok{foreign\_curve\_handle =}\NormalTok{ foreign\_curve\_handle\_prdc,}
  \AttributeTok{domestic\_curve\_handle =}\NormalTok{ domestic\_curve\_handle\_prdc,}
  \AttributeTok{volatility\_handle =}\NormalTok{ fx\_volatility\_handle\_prdc}
\NormalTok{)}

\NormalTok{fx\_market\_summary\_tbl }\OtherTok{\textless{}{-}}\NormalTok{ tibble}\SpecialCharTok{::}\FunctionTok{tribble}\NormalTok{(}
  \SpecialCharTok{\textasciitilde{}}\NormalTok{metric,               }\SpecialCharTok{\textasciitilde{}}\NormalTok{value,}
  \StringTok{"fx\_spot"}\NormalTok{,              fx\_spot\_prdc,}
  \StringTok{"fx\_volatility"}\NormalTok{,        fx\_volatility\_prdc,}
  \StringTok{"domestic\_rate"}\NormalTok{,        domestic\_rate\_prdc,}
  \StringTok{"foreign\_rate"}\NormalTok{,         foreign\_rate\_prdc,}
  \StringTok{"risk\_neutral\_drift"}\NormalTok{,   domestic\_rate\_prdc }\SpecialCharTok{{-}}
\NormalTok{                            foreign\_rate\_prdc}
\NormalTok{)}

\NormalTok{GaussQuant}\SpecialCharTok{::}\FunctionTok{show\_tbl\_GQH}\NormalTok{(}
\NormalTok{  fx\_market\_summary\_tbl,}
  \StringTok{"PRDC FX market assumptions"}\NormalTok{,}
  \AttributeTok{n =} \DecValTok{20}
\NormalTok{)}
\CommentTok{\#\textgreater{} \# A tibble: 5 x 2}
\CommentTok{\#\textgreater{}   metric              value}
\CommentTok{\#\textgreater{}   \textless{}chr\textgreater{}               \textless{}dbl\textgreater{}}
\CommentTok{\#\textgreater{} 1 fx\_spot            133.  }
\CommentTok{\#\textgreater{} 2 fx\_volatility        0.1 }
\CommentTok{\#\textgreater{} 3 domestic\_rate        0.01}
\CommentTok{\#\textgreater{} 4 foreign\_rate         0.03}
\CommentTok{\#\textgreater{} 5 risk\_neutral\_drift  {-}0.02}
\end{Highlighting}
\end{Shaded}

\subsection{9.2.2 Coupon Payoff}\label{coupon-payoff}

本例のFX連動Couponは次式である。

\[
C(S)
=
\min
\left[
2.2\%,
\max
\left(
12.2\%
\frac{S}{120}
-
10\%,
0
\right)
\right]
\]

Floorに到達するFX水準は、

\[
S_{\mathrm{floor}}
=
S_0
\frac{
b+C_{\min}
}{
a
}
\]

Capに到達するFX水準は、

\[
S_{\mathrm{cap}}
=
S_0
\frac{
b+C_{\max}
}{
a
}
\]

である。

\begin{Shaded}
\begin{Highlighting}[]
\NormalTok{fx\_floor\_threshold\_prdc }\OtherTok{\textless{}{-}}\NormalTok{ coupon\_fx\_base\_prdc }\SpecialCharTok{*}
\NormalTok{  (}
\NormalTok{    coupon\_subtraction\_prdc }\SpecialCharTok{+}
\NormalTok{      coupon\_floor\_prdc}
\NormalTok{  ) }\SpecialCharTok{/}
\NormalTok{  coupon\_leverage\_prdc}

\NormalTok{fx\_cap\_threshold\_prdc }\OtherTok{\textless{}{-}}\NormalTok{ coupon\_fx\_base\_prdc }\SpecialCharTok{*}
\NormalTok{  (}
\NormalTok{    coupon\_subtraction\_prdc }\SpecialCharTok{+}
\NormalTok{      coupon\_cap\_prdc}
\NormalTok{  ) }\SpecialCharTok{/}
\NormalTok{  coupon\_leverage\_prdc}

\NormalTok{fx\_payoff\_grid\_tbl }\OtherTok{\textless{}{-}}\NormalTok{ tibble}\SpecialCharTok{::}\FunctionTok{tibble}\NormalTok{(}
  \AttributeTok{fx\_rate =} \FunctionTok{seq}\NormalTok{(}
    \DecValTok{70}\NormalTok{,}
    \DecValTok{180}\NormalTok{,}
    \AttributeTok{by =} \DecValTok{2}
\NormalTok{  )}
\NormalTok{) }\SpecialCharTok{|\textgreater{}}
\NormalTok{  dplyr}\SpecialCharTok{::}\FunctionTok{mutate}\NormalTok{(}
    \AttributeTok{uncapped\_coupon =}\NormalTok{ coupon\_leverage\_prdc }\SpecialCharTok{*}
\NormalTok{      .data}\SpecialCharTok{$}\NormalTok{fx\_rate }\SpecialCharTok{/}
\NormalTok{      coupon\_fx\_base\_prdc }\SpecialCharTok{{-}}
\NormalTok{      coupon\_subtraction\_prdc,}
    \AttributeTok{coupon\_rate =}\NormalTok{ GaussQuant}\SpecialCharTok{::}\FunctionTok{prdc\_coupon\_GQL}\NormalTok{(}
      \AttributeTok{fx\_rate =}\NormalTok{ .data}\SpecialCharTok{$}\NormalTok{fx\_rate,}
      \AttributeTok{fx\_base =}\NormalTok{ coupon\_fx\_base\_prdc,}
      \AttributeTok{leverage =}\NormalTok{ coupon\_leverage\_prdc,}
      \AttributeTok{subtraction =}\NormalTok{ coupon\_subtraction\_prdc,}
      \AttributeTok{floor\_rate =}\NormalTok{ coupon\_floor\_prdc,}
      \AttributeTok{cap\_rate =}\NormalTok{ coupon\_cap\_prdc}
\NormalTok{    )}
\NormalTok{  )}

\NormalTok{coupon\_threshold\_tbl }\OtherTok{\textless{}{-}}\NormalTok{ tibble}\SpecialCharTok{::}\FunctionTok{tribble}\NormalTok{(}
  \SpecialCharTok{\textasciitilde{}}\NormalTok{threshold,       }\SpecialCharTok{\textasciitilde{}}\NormalTok{fx\_rate,}
  \StringTok{"floor\_threshold"}\NormalTok{, fx\_floor\_threshold\_prdc,}
  \StringTok{"cap\_threshold"}\NormalTok{,   fx\_cap\_threshold\_prdc,}
  \StringTok{"current\_spot"}\NormalTok{,    fx\_spot\_prdc}
\NormalTok{)}

\NormalTok{prdc\_payoff\_plot }\OtherTok{\textless{}{-}}\NormalTok{ ggplot2}\SpecialCharTok{::}\FunctionTok{ggplot}\NormalTok{(}
\NormalTok{  fx\_payoff\_grid\_tbl,}
\NormalTok{  ggplot2}\SpecialCharTok{::}\FunctionTok{aes}\NormalTok{(}
    \AttributeTok{x =}\NormalTok{ .data}\SpecialCharTok{$}\NormalTok{fx\_rate,}
    \AttributeTok{y =}\NormalTok{ .data}\SpecialCharTok{$}\NormalTok{coupon\_rate}
\NormalTok{  )}
\NormalTok{) }\SpecialCharTok{+}
\NormalTok{  ggplot2}\SpecialCharTok{::}\FunctionTok{geom\_line}\NormalTok{(}
    \AttributeTok{linewidth =} \FloatTok{0.8}
\NormalTok{  ) }\SpecialCharTok{+}
\NormalTok{  ggplot2}\SpecialCharTok{::}\FunctionTok{geom\_vline}\NormalTok{(}
    \AttributeTok{xintercept =} \FunctionTok{c}\NormalTok{(}
\NormalTok{      fx\_floor\_threshold\_prdc,}
\NormalTok{      fx\_cap\_threshold\_prdc}
\NormalTok{    ),}
    \AttributeTok{linetype =} \StringTok{"dashed"}
\NormalTok{  ) }\SpecialCharTok{+}
\NormalTok{  ggplot2}\SpecialCharTok{::}\FunctionTok{scale\_y\_continuous}\NormalTok{(}
    \AttributeTok{labels =}\NormalTok{ scales}\SpecialCharTok{::}\FunctionTok{percent\_format}\NormalTok{(}
      \AttributeTok{accuracy =} \FloatTok{0.1}
\NormalTok{    )}
\NormalTok{  ) }\SpecialCharTok{+}
\NormalTok{  ggplot2}\SpecialCharTok{::}\FunctionTok{labs}\NormalTok{(}
    \AttributeTok{title =} \StringTok{"PRDC coupon payoff"}\NormalTok{,}
    \AttributeTok{subtitle =} \StringTok{"Floored and capped FX{-}linked coupon"}\NormalTok{,}
    \AttributeTok{x =} \StringTok{"FX rate"}\NormalTok{,}
    \AttributeTok{y =} \StringTok{"Coupon rate"}
\NormalTok{  ) }\SpecialCharTok{+}
\NormalTok{  ggplot2}\SpecialCharTok{::}\FunctionTok{theme\_minimal}\NormalTok{()}

\NormalTok{GaussQuant}\SpecialCharTok{::}\FunctionTok{show\_tbl\_GQH}\NormalTok{(}
\NormalTok{  coupon\_threshold\_tbl,}
  \StringTok{"PRDC coupon FX thresholds"}\NormalTok{,}
  \AttributeTok{n =} \DecValTok{20}
\NormalTok{)}
\CommentTok{\#\textgreater{} \# A tibble: 3 x 2}
\CommentTok{\#\textgreater{}   threshold       fx\_rate}
\CommentTok{\#\textgreater{}   \textless{}chr\textgreater{}             \textless{}dbl\textgreater{}}
\CommentTok{\#\textgreater{} 1 floor\_threshold    98.4}
\CommentTok{\#\textgreater{} 2 cap\_threshold     120  }
\CommentTok{\#\textgreater{} 3 current\_spot      133.}

\NormalTok{GaussQuant}\SpecialCharTok{::}\FunctionTok{show\_tbl\_GQH}\NormalTok{(}
\NormalTok{  fx\_payoff\_grid\_tbl,}
  \StringTok{"PRDC coupon payoff table"}\NormalTok{,}
  \AttributeTok{n =} \DecValTok{60}
\NormalTok{)}
\CommentTok{\#\textgreater{} \# A tibble: 56 x 3}
\CommentTok{\#\textgreater{}    fx\_rate uncapped\_coupon coupon\_rate}
\CommentTok{\#\textgreater{}      \textless{}dbl\textgreater{}           \textless{}dbl\textgreater{}       \textless{}dbl\textgreater{}}
\CommentTok{\#\textgreater{}  1      70         {-}0.0288           0}
\CommentTok{\#\textgreater{}  2      72         {-}0.0268           0}
\CommentTok{\#\textgreater{}  3      74         {-}0.0248           0}
\CommentTok{\#\textgreater{}  4      76         {-}0.0227           0}
\CommentTok{\#\textgreater{}  5      78         {-}0.0207           0}
\CommentTok{\#\textgreater{}  6      80         {-}0.0187           0}
\CommentTok{\#\textgreater{}  7      82         {-}0.0166           0}
\CommentTok{\#\textgreater{}  8      84         {-}0.0146           0}
\CommentTok{\#\textgreater{}  9      86         {-}0.0126           0}
\CommentTok{\#\textgreater{} 10      88         {-}0.0105           0}
\CommentTok{\#\textgreater{} \# i 46 more rows}

\FunctionTok{print}\NormalTok{(prdc\_payoff\_plot)}
\end{Highlighting}
\end{Shaded}

\pandocbounded{\includegraphics[keepaspectratio]{part2_quantive_riemann/ch09_exotic_derivatives_prdc_files/figure-pdf/chapter09-coupon-payoff-1.pdf}}

現在のFX
SpotはCap到達水準を上回っているため、評価時点のCoupon式だけを見ればCapに張り付いている。しかし、満期までのFX分布にはFloor、Linear領域およびCap領域のすべてが含まれ得る。

\section{9.3 Monte
CarloによるPRDC評価}\label{monte-carloux306bux3088ux308bprdcux8a55ux4fa1}

\subsection{9.3.1
Coupon日ベースのFX時間グリッド}\label{couponux65e5ux30d9ux30fcux30b9ux306efxux6642ux9593ux30b0ux30eaux30c3ux30c9}

PRDC評価に必要なのは、評価日より後に到来するCoupon日である。GaussQuantの\texttt{fx\_path\_generator\_GQL()}は、実際のCoupon日を用いて不均一な時間間隔を保持する。

\begin{Shaded}
\begin{Highlighting}[]
\NormalTok{fx\_path\_generator\_prdc }\OtherTok{\textless{}{-}}\NormalTok{ GaussQuant}\SpecialCharTok{::}\FunctionTok{fx\_path\_generator\_GQL}\NormalTok{(}
  \AttributeTok{valuation\_date =}\NormalTok{ valuation\_date\_prdc,}
  \AttributeTok{coupon\_schedule =}\NormalTok{ coupon\_schedule\_prdc,}
  \AttributeTok{day\_counter =}\NormalTok{ fx\_day\_counter\_prdc,}
  \AttributeTok{process =}\NormalTok{ fx\_process\_prdc}
\NormalTok{)}

\NormalTok{time\_grid\_tbl }\OtherTok{\textless{}{-}}\NormalTok{ tibble}\SpecialCharTok{::}\FunctionTok{tibble}\NormalTok{(}
  \AttributeTok{coupon\_number =} \FunctionTok{seq\_along}\NormalTok{(}
\NormalTok{    fx\_path\_generator\_prdc}\SpecialCharTok{$}\NormalTok{remaining\_coupon\_dates}
\NormalTok{  ),}
  \AttributeTok{coupon\_date =}\NormalTok{ fx\_path\_generator\_prdc}\SpecialCharTok{$}\NormalTok{remaining\_coupon\_dates,}
  \AttributeTok{time\_from\_valuation =}\NormalTok{ fx\_path\_generator\_prdc}\SpecialCharTok{$}\NormalTok{coupon\_times,}
  \AttributeTok{time\_step =}\NormalTok{ fx\_path\_generator\_prdc}\SpecialCharTok{$}\NormalTok{time\_steps}
\NormalTok{)}

\NormalTok{path\_generator\_summary\_tbl }\OtherTok{\textless{}{-}}\NormalTok{ tibble}\SpecialCharTok{::}\FunctionTok{tribble}\NormalTok{(}
  \SpecialCharTok{\textasciitilde{}}\NormalTok{metric,                  }\SpecialCharTok{\textasciitilde{}}\NormalTok{value,}
  \StringTok{"remaining\_coupon\_dates"}\NormalTok{,  }\FunctionTok{length}\NormalTok{(}
\NormalTok{                               fx\_path\_generator\_prdc}\SpecialCharTok{$}\NormalTok{remaining\_coupon\_dates}
\NormalTok{                             ),}
  \StringTok{"number\_of\_steps"}\NormalTok{,         fx\_path\_generator\_prdc}\SpecialCharTok{$}\NormalTok{number\_of\_steps,}
  \StringTok{"first\_coupon\_time"}\NormalTok{,       }\FunctionTok{min}\NormalTok{(}
\NormalTok{                               fx\_path\_generator\_prdc}\SpecialCharTok{$}\NormalTok{coupon\_times}
\NormalTok{                             ),}
  \StringTok{"final\_coupon\_time"}\NormalTok{,       }\FunctionTok{max}\NormalTok{(}
\NormalTok{                               fx\_path\_generator\_prdc}\SpecialCharTok{$}\NormalTok{coupon\_times}
\NormalTok{                             )}
\NormalTok{)}

\NormalTok{GaussQuant}\SpecialCharTok{::}\FunctionTok{show\_tbl\_GQH}\NormalTok{(}
\NormalTok{  path\_generator\_summary\_tbl,}
  \StringTok{"PRDC FX path{-}generator summary"}\NormalTok{,}
  \AttributeTok{n =} \DecValTok{20}
\NormalTok{)}
\CommentTok{\#\textgreater{} \# A tibble: 4 x 2}
\CommentTok{\#\textgreater{}   metric                  value}
\CommentTok{\#\textgreater{}   \textless{}chr\textgreater{}                   \textless{}dbl\textgreater{}}
\CommentTok{\#\textgreater{} 1 remaining\_coupon\_dates 37    }
\CommentTok{\#\textgreater{} 2 number\_of\_steps        37    }
\CommentTok{\#\textgreater{} 3 first\_coupon\_time       0.406}
\CommentTok{\#\textgreater{} 4 final\_coupon\_time      18.7}

\NormalTok{GaussQuant}\SpecialCharTok{::}\FunctionTok{show\_tbl\_GQH}\NormalTok{(}
\NormalTok{  time\_grid\_tbl,}
  \StringTok{"PRDC remaining coupon{-}date time grid"}\NormalTok{,}
  \AttributeTok{n =} \DecValTok{60}
\NormalTok{)}
\CommentTok{\#\textgreater{} \# A tibble: 37 x 4}
\CommentTok{\#\textgreater{}    coupon\_number coupon\_date time\_from\_valuation time\_step}
\CommentTok{\#\textgreater{}            \textless{}int\textgreater{} \textless{}date\textgreater{}                    \textless{}dbl\textgreater{}     \textless{}dbl\textgreater{}}
\CommentTok{\#\textgreater{}  1             1 2023{-}09{-}04                0.406     0.406}
\CommentTok{\#\textgreater{}  2             2 2024{-}03{-}04                0.911     0.506}
\CommentTok{\#\textgreater{}  3             3 2024{-}09{-}03                1.42      0.508}
\CommentTok{\#\textgreater{}  4             4 2025{-}03{-}03                1.92      0.503}
\CommentTok{\#\textgreater{}  5             5 2025{-}09{-}03                2.43      0.511}
\CommentTok{\#\textgreater{}  6             6 2026{-}03{-}03                2.94      0.503}
\CommentTok{\#\textgreater{}  7             7 2026{-}09{-}03                3.45      0.511}
\CommentTok{\#\textgreater{}  8             8 2027{-}03{-}03                3.95      0.503}
\CommentTok{\#\textgreater{}  9             9 2027{-}09{-}03                4.46      0.511}
\CommentTok{\#\textgreater{} 10            10 2028{-}03{-}03                4.97      0.506}
\CommentTok{\#\textgreater{} \# i 27 more rows}
\end{Highlighting}
\end{Shaded}

QuantLibの\texttt{StochasticProcess1D\$evolve()}へは標準正規変量を渡す。Brownian
incrementを外側で二重に時間調整してはならない。GaussQuantのPath
Generatorでは、標準正規乱数と各Coupon間の実際の時間差を分けて処理している。

\subsection{9.3.2
サンプルFXパス}\label{ux30b5ux30f3ux30d7ux30ebfxux30d1ux30b9}

Monte
Carlo評価へ進む前に、20本のFXパスを生成し、将来のFX分布とCouponの動きを確認する。

\begin{Shaded}
\begin{Highlighting}[]
\NormalTok{sample\_path\_count\_prdc }\OtherTok{\textless{}{-}} \DecValTok{20}\NormalTok{L}

\NormalTok{sample\_fx\_path\_matrix\_prdc }\OtherTok{\textless{}{-}}\NormalTok{ GaussQuant}\SpecialCharTok{::}\FunctionTok{prdc\_path\_matrix\_GQL}\NormalTok{(}
  \AttributeTok{path\_generator =}\NormalTok{ fx\_path\_generator\_prdc,}
  \AttributeTok{n\_paths =}\NormalTok{ sample\_path\_count\_prdc,}
  \AttributeTok{seed =}\NormalTok{ random\_seed\_prdc}
\NormalTok{)}

\NormalTok{sample\_fx\_path\_tbl }\OtherTok{\textless{}{-}}\NormalTok{ purrr}\SpecialCharTok{::}\FunctionTok{map\_dfr}\NormalTok{(}
  \FunctionTok{seq\_len}\NormalTok{(}
    \FunctionTok{ncol}\NormalTok{(sample\_fx\_path\_matrix\_prdc)}
\NormalTok{  ),}
  \ControlFlowTok{function}\NormalTok{(path\_id) \{}
\NormalTok{    path\_values }\OtherTok{\textless{}{-}}\NormalTok{ sample\_fx\_path\_matrix\_prdc[}
\NormalTok{      ,}
\NormalTok{      path\_id}
\NormalTok{    ]}

\NormalTok{    tibble}\SpecialCharTok{::}\FunctionTok{tibble}\NormalTok{(}
      \AttributeTok{path\_id =}\NormalTok{ path\_id,}
      \AttributeTok{coupon\_number =} \FunctionTok{seq\_along}\NormalTok{(}
\NormalTok{        path\_values}
\NormalTok{      ),}
      \AttributeTok{coupon\_date =}\NormalTok{ fx\_path\_generator\_prdc}\SpecialCharTok{$}\NormalTok{remaining\_coupon\_dates,}
      \AttributeTok{time =}\NormalTok{ fx\_path\_generator\_prdc}\SpecialCharTok{$}\NormalTok{coupon\_times,}
      \AttributeTok{fx\_rate =}\NormalTok{ path\_values,}
      \AttributeTok{coupon\_rate =}\NormalTok{ GaussQuant}\SpecialCharTok{::}\FunctionTok{prdc\_coupon\_GQL}\NormalTok{(}
        \AttributeTok{fx\_rate =}\NormalTok{ path\_values,}
        \AttributeTok{fx\_base =}\NormalTok{ coupon\_fx\_base\_prdc,}
        \AttributeTok{leverage =}\NormalTok{ coupon\_leverage\_prdc,}
        \AttributeTok{subtraction =}\NormalTok{ coupon\_subtraction\_prdc,}
        \AttributeTok{floor\_rate =}\NormalTok{ coupon\_floor\_prdc,}
        \AttributeTok{cap\_rate =}\NormalTok{ coupon\_cap\_prdc}
\NormalTok{      )}
\NormalTok{    )}
\NormalTok{  \}}
\NormalTok{)}

\NormalTok{sample\_fx\_path\_plot }\OtherTok{\textless{}{-}}\NormalTok{ ggplot2}\SpecialCharTok{::}\FunctionTok{ggplot}\NormalTok{(}
\NormalTok{  sample\_fx\_path\_tbl,}
\NormalTok{  ggplot2}\SpecialCharTok{::}\FunctionTok{aes}\NormalTok{(}
    \AttributeTok{x =}\NormalTok{ .data}\SpecialCharTok{$}\NormalTok{time,}
    \AttributeTok{y =}\NormalTok{ .data}\SpecialCharTok{$}\NormalTok{fx\_rate,}
    \AttributeTok{group =}\NormalTok{ .data}\SpecialCharTok{$}\NormalTok{path\_id}
\NormalTok{  )}
\NormalTok{) }\SpecialCharTok{+}
\NormalTok{  ggplot2}\SpecialCharTok{::}\FunctionTok{geom\_line}\NormalTok{(}
    \AttributeTok{linewidth =} \FloatTok{0.5}\NormalTok{,}
    \AttributeTok{alpha =} \FloatTok{0.7}
\NormalTok{  ) }\SpecialCharTok{+}
\NormalTok{  ggplot2}\SpecialCharTok{::}\FunctionTok{labs}\NormalTok{(}
    \AttributeTok{title =} \StringTok{"Sample FX paths for PRDC valuation"}\NormalTok{,}
    \AttributeTok{x =} \StringTok{"Time from valuation date (years)"}\NormalTok{,}
    \AttributeTok{y =} \StringTok{"FX rate"}
\NormalTok{  ) }\SpecialCharTok{+}
\NormalTok{  ggplot2}\SpecialCharTok{::}\FunctionTok{theme\_minimal}\NormalTok{()}

\NormalTok{GaussQuant}\SpecialCharTok{::}\FunctionTok{show\_tbl\_GQH}\NormalTok{(}
\NormalTok{  sample\_fx\_path\_tbl,}
  \StringTok{"Sample PRDC FX and coupon paths"}\NormalTok{,}
  \AttributeTok{n =} \DecValTok{40}
\NormalTok{)}
\CommentTok{\#\textgreater{} \# A tibble: 40 x 6}
\CommentTok{\#\textgreater{}    path\_id coupon\_number coupon\_date  time fx\_rate coupon\_rate}
\CommentTok{\#\textgreater{}      \textless{}int\textgreater{}         \textless{}int\textgreater{} \textless{}date\textgreater{}      \textless{}dbl\textgreater{}   \textless{}dbl\textgreater{}       \textless{}dbl\textgreater{}}
\CommentTok{\#\textgreater{}  1       1             1 2023{-}09{-}04  0.406    127.       0.022}
\CommentTok{\#\textgreater{}  2       1             2 2024{-}03{-}04  0.911    124.       0.022}
\CommentTok{\#\textgreater{}  3       1             3 2024{-}09{-}03  1.42     136.       0.022}
\CommentTok{\#\textgreater{}  4       1             4 2025{-}03{-}03  1.92     135.       0.022}
\CommentTok{\#\textgreater{}  5       1             5 2025{-}09{-}03  2.43     135.       0.022}
\CommentTok{\#\textgreater{}  6       1             6 2026{-}03{-}03  2.94     150.       0.022}
\CommentTok{\#\textgreater{}  7       1             7 2026{-}09{-}03  3.45     154.       0.022}
\CommentTok{\#\textgreater{}  8       1             8 2027{-}03{-}03  3.95     139.       0.022}
\CommentTok{\#\textgreater{}  9       1             9 2027{-}09{-}03  4.46     130.       0.022}
\CommentTok{\#\textgreater{} 10       1            10 2028{-}03{-}03  4.97     125.       0.022}
\CommentTok{\#\textgreater{} \# i 30 more rows}

\FunctionTok{print}\NormalTok{(sample\_fx\_path\_plot)}
\end{Highlighting}
\end{Shaded}

\pandocbounded{\includegraphics[keepaspectratio]{part2_quantive_riemann/ch09_exotic_derivatives_prdc_files/figure-pdf/chapter09-sample-paths-1.pdf}}

\subsection{9.3.3
PRDCの現在価値}\label{prdcux306eux73feux5728ux4fa1ux5024}

各Coupon日の期待Coupon rateは、Monte Carloによって次のように近似する。

\[
E[C(S_{t_i})]
\approx
\frac{1}{M}
\sum_{j=1}^{M}
C(S_{t_i}^{(j)})
\]

ここで、\(M\)はパス数、\(S_{t_i}^{(j)}\)はパス \(j\) におけるCoupon日
\(t_i\) のFX Rateである。

各Coupon cash flowは、

\[
CF_i
=
N_0\alpha_iE[C(S_{t_i})]
\]

であり、PRDC現在価値は、

\[
PV_{\mathrm{PRDC}}
=
\sum_i
D(0,t_i)CF_i
+
D(0,T)N_0
\]

となる。

\begin{Shaded}
\begin{Highlighting}[]
\NormalTok{prdc\_result }\OtherTok{\textless{}{-}}\NormalTok{ GaussQuant}\SpecialCharTok{::}\FunctionTok{prdc\_npv\_GQL}\NormalTok{(}
  \AttributeTok{valuation\_date =}\NormalTok{ valuation\_date\_prdc,}
  \AttributeTok{coupon\_schedule =}\NormalTok{ coupon\_schedule\_prdc,}
  \AttributeTok{process =}\NormalTok{ fx\_process\_prdc,}
  \AttributeTok{discount\_curve =}\NormalTok{ discount\_curve\_prdc,}
  \AttributeTok{notional =}\NormalTok{ notional\_prdc,}
  \AttributeTok{fx\_base =}\NormalTok{ coupon\_fx\_base\_prdc,}
  \AttributeTok{leverage =}\NormalTok{ coupon\_leverage\_prdc,}
  \AttributeTok{subtraction =}\NormalTok{ coupon\_subtraction\_prdc,}
  \AttributeTok{floor\_rate =}\NormalTok{ coupon\_floor\_prdc,}
  \AttributeTok{cap\_rate =}\NormalTok{ coupon\_cap\_prdc,}
  \AttributeTok{day\_counter =}\NormalTok{ fx\_day\_counter\_prdc,}
  \AttributeTok{intro\_coupon\_dates =}\NormalTok{ intro\_coupon\_dates\_prdc,}
  \AttributeTok{intro\_coupon\_rate =}\NormalTok{ intro\_coupon\_rate\_prdc,}
  \AttributeTok{n\_paths =}\NormalTok{ number\_of\_paths\_prdc,}
  \AttributeTok{seed =}\NormalTok{ random\_seed\_prdc,}
  \AttributeTok{redemption\_amount =}\NormalTok{ notional\_prdc,}
  \AttributeTok{return\_details =} \ConstantTok{TRUE}
\NormalTok{)}

\NormalTok{fx\_path\_matrix\_prdc }\OtherTok{\textless{}{-}}\NormalTok{ prdc\_result}\SpecialCharTok{$}\NormalTok{path\_matrix}
\NormalTok{expected\_coupon\_tbl }\OtherTok{\textless{}{-}}\NormalTok{ prdc\_result}\SpecialCharTok{$}\NormalTok{expected\_coupon}
\NormalTok{prdc\_cashflow\_tbl }\OtherTok{\textless{}{-}}\NormalTok{ prdc\_result}\SpecialCharTok{$}\NormalTok{cashflows}

\NormalTok{coupon\_leg\_present\_value\_prdc }\OtherTok{\textless{}{-}}\NormalTok{ prdc\_result}\SpecialCharTok{$}\NormalTok{summary }\SpecialCharTok{|\textgreater{}}
\NormalTok{  dplyr}\SpecialCharTok{::}\FunctionTok{filter}\NormalTok{(}
\NormalTok{    .data}\SpecialCharTok{$}\NormalTok{metric }\SpecialCharTok{==} \StringTok{"coupon\_leg\_npv"}
\NormalTok{  ) }\SpecialCharTok{|\textgreater{}}
\NormalTok{  dplyr}\SpecialCharTok{::}\FunctionTok{pull}\NormalTok{(}
\NormalTok{    value}
\NormalTok{  )}

\NormalTok{redemption\_present\_value\_prdc }\OtherTok{\textless{}{-}}\NormalTok{ prdc\_result}\SpecialCharTok{$}\NormalTok{summary }\SpecialCharTok{|\textgreater{}}
\NormalTok{  dplyr}\SpecialCharTok{::}\FunctionTok{filter}\NormalTok{(}
\NormalTok{    .data}\SpecialCharTok{$}\NormalTok{metric }\SpecialCharTok{==} \StringTok{"redemption\_leg\_npv"}
\NormalTok{  ) }\SpecialCharTok{|\textgreater{}}
\NormalTok{  dplyr}\SpecialCharTok{::}\FunctionTok{pull}\NormalTok{(}
\NormalTok{    value}
\NormalTok{  )}

\NormalTok{prdc\_present\_value\_jpy }\OtherTok{\textless{}{-}}\NormalTok{ prdc\_result}\SpecialCharTok{$}\NormalTok{npv}

\NormalTok{prdc\_valuation\_summary\_tbl }\OtherTok{\textless{}{-}}\NormalTok{ tibble}\SpecialCharTok{::}\FunctionTok{tribble}\NormalTok{(}
  \SpecialCharTok{\textasciitilde{}}\NormalTok{metric,                    }\SpecialCharTok{\textasciitilde{}}\NormalTok{value,}
  \StringTok{"coupon\_leg\_pv\_jpy"}\NormalTok{,         coupon\_leg\_present\_value\_prdc,}
  \StringTok{"redemption\_leg\_pv\_jpy"}\NormalTok{,     redemption\_present\_value\_prdc,}
  \StringTok{"total\_pv\_jpy"}\NormalTok{,              prdc\_present\_value\_jpy,}
  \StringTok{"pv\_as\_percent\_of\_notional"}\NormalTok{, prdc\_present\_value\_jpy }\SpecialCharTok{/}
\NormalTok{                                  notional\_prdc,}
  \StringTok{"number\_of\_paths"}\NormalTok{,           number\_of\_paths\_prdc,}
  \StringTok{"number\_of\_coupon\_steps"}\NormalTok{,    }\FunctionTok{nrow}\NormalTok{(}
\NormalTok{                                  fx\_path\_matrix\_prdc}
\NormalTok{                                )}
\NormalTok{)}

\NormalTok{GaussQuant}\SpecialCharTok{::}\FunctionTok{show\_tbl\_GQH}\NormalTok{(}
\NormalTok{  prdc\_valuation\_summary\_tbl,}
  \StringTok{"PRDC Monte Carlo valuation summary"}\NormalTok{,}
  \AttributeTok{n =} \DecValTok{20}
\NormalTok{)}
\CommentTok{\#\textgreater{} \# A tibble: 6 x 2}
\CommentTok{\#\textgreater{}   metric                           value}
\CommentTok{\#\textgreater{}   \textless{}chr\textgreater{}                            \textless{}dbl\textgreater{}}
\CommentTok{\#\textgreater{} 1 coupon\_leg\_pv\_jpy          61290348.  }
\CommentTok{\#\textgreater{} 2 redemption\_leg\_pv\_jpy     273266946.  }
\CommentTok{\#\textgreater{} 3 total\_pv\_jpy              334557294.  }
\CommentTok{\#\textgreater{} 4 pv\_as\_percent\_of\_notional         1.12}
\CommentTok{\#\textgreater{} 5 number\_of\_paths                3000   }
\CommentTok{\#\textgreater{} 6 number\_of\_coupon\_steps           37}
\end{Highlighting}
\end{Shaded}

本章の元本償還額は固定されているため、Redemption
Legは割引率に依存するがFXには依存しない。FX SpotとFX
volatilityの影響は、主にCoupon Legを通じて現れる。

\subsection{9.3.4
期待Couponとキャッシュフロー}\label{ux671fux5f85couponux3068ux30adux30e3ux30c3ux30b7ux30e5ux30d5ux30edux30fc}

\begin{Shaded}
\begin{Highlighting}[]
\NormalTok{expected\_coupon\_display\_tbl }\OtherTok{\textless{}{-}}\NormalTok{ expected\_coupon\_tbl }\SpecialCharTok{|\textgreater{}}
\NormalTok{  dplyr}\SpecialCharTok{::}\FunctionTok{select}\NormalTok{(}
\NormalTok{    coupon\_number,}
\NormalTok{    coupon\_date,}
\NormalTok{    time,}
\NormalTok{    expected\_fx,}
\NormalTok{    expected\_coupon\_rate,}
\NormalTok{    is\_intro\_coupon,}
\NormalTok{    final\_coupon\_rate}
\NormalTok{  )}

\NormalTok{prdc\_cashflow\_display\_tbl }\OtherTok{\textless{}{-}}\NormalTok{ prdc\_cashflow\_tbl }\SpecialCharTok{|\textgreater{}}
\NormalTok{  dplyr}\SpecialCharTok{::}\FunctionTok{select}\NormalTok{(}
\NormalTok{    coupon\_number,}
\NormalTok{    accrual\_start,}
\NormalTok{    accrual\_end,}
\NormalTok{    final\_coupon\_rate,}
\NormalTok{    accrual\_fraction,}
\NormalTok{    coupon\_amount,}
\NormalTok{    redemption\_amount,}
\NormalTok{    total\_amount,}
\NormalTok{    discount\_factor,}
\NormalTok{    total\_present\_value}
\NormalTok{  )}

\NormalTok{GaussQuant}\SpecialCharTok{::}\FunctionTok{show\_tbl\_GQH}\NormalTok{(}
\NormalTok{  expected\_coupon\_display\_tbl,}
  \StringTok{"PRDC expected FX and coupon rates"}\NormalTok{,}
  \AttributeTok{n =} \DecValTok{60}
\NormalTok{)}
\CommentTok{\#\textgreater{} \# A tibble: 37 x 7}
\CommentTok{\#\textgreater{}    coupon\_number coupon\_date  time expected\_fx expected\_coupon\_rate}
\CommentTok{\#\textgreater{}            \textless{}int\textgreater{} \textless{}date\textgreater{}      \textless{}dbl\textgreater{}       \textless{}dbl\textgreater{}                \textless{}dbl\textgreater{}}
\CommentTok{\#\textgreater{}  1             1 2023{-}09{-}04  0.406        132.               0.0218}
\CommentTok{\#\textgreater{}  2             2 2024{-}03{-}04  0.911        131.               0.0208}
\CommentTok{\#\textgreater{}  3             3 2024{-}09{-}03  1.42         130.               0.0196}
\CommentTok{\#\textgreater{}  4             4 2025{-}03{-}03  1.92         128.               0.0185}
\CommentTok{\#\textgreater{}  5             5 2025{-}09{-}03  2.43         127.               0.0176}
\CommentTok{\#\textgreater{}  6             6 2026{-}03{-}03  2.94         126.               0.0168}
\CommentTok{\#\textgreater{}  7             7 2026{-}09{-}03  3.45         124.               0.0161}
\CommentTok{\#\textgreater{}  8             8 2027{-}03{-}03  3.95         123.               0.0153}
\CommentTok{\#\textgreater{}  9             9 2027{-}09{-}03  4.46         122.               0.0147}
\CommentTok{\#\textgreater{} 10            10 2028{-}03{-}03  4.97         121.               0.0140}
\CommentTok{\#\textgreater{} \# i 27 more rows}
\CommentTok{\#\textgreater{} \# i 2 more variables: is\_intro\_coupon \textless{}lgl\textgreater{}, final\_coupon\_rate \textless{}dbl\textgreater{}}

\NormalTok{GaussQuant}\SpecialCharTok{::}\FunctionTok{show\_tbl\_GQH}\NormalTok{(}
\NormalTok{  prdc\_cashflow\_display\_tbl,}
  \StringTok{"PRDC expected cash{-}flow table"}\NormalTok{,}
  \AttributeTok{n =} \DecValTok{60}
\NormalTok{)}
\CommentTok{\#\textgreater{} \# A tibble: 37 x 10}
\CommentTok{\#\textgreater{}    coupon\_number accrual\_start accrual\_end final\_coupon\_rate accrual\_fraction}
\CommentTok{\#\textgreater{}            \textless{}int\textgreater{} \textless{}date\textgreater{}        \textless{}date\textgreater{}                  \textless{}dbl\textgreater{}            \textless{}dbl\textgreater{}}
\CommentTok{\#\textgreater{}  1             1 2023{-}03{-}03    2023{-}09{-}04             0.0218            0.514}
\CommentTok{\#\textgreater{}  2             2 2023{-}09{-}04    2024{-}03{-}04             0.0208            0.506}
\CommentTok{\#\textgreater{}  3             3 2024{-}03{-}04    2024{-}09{-}03             0.0196            0.508}
\CommentTok{\#\textgreater{}  4             4 2024{-}09{-}03    2025{-}03{-}03             0.0185            0.503}
\CommentTok{\#\textgreater{}  5             5 2025{-}03{-}03    2025{-}09{-}03             0.0176            0.511}
\CommentTok{\#\textgreater{}  6             6 2025{-}09{-}03    2026{-}03{-}03             0.0168            0.503}
\CommentTok{\#\textgreater{}  7             7 2026{-}03{-}03    2026{-}09{-}03             0.0161            0.511}
\CommentTok{\#\textgreater{}  8             8 2026{-}09{-}03    2027{-}03{-}03             0.0153            0.503}
\CommentTok{\#\textgreater{}  9             9 2027{-}03{-}03    2027{-}09{-}03             0.0147            0.511}
\CommentTok{\#\textgreater{} 10            10 2027{-}09{-}03    2028{-}03{-}03             0.0140            0.506}
\CommentTok{\#\textgreater{} \# i 27 more rows}
\CommentTok{\#\textgreater{} \# i 5 more variables: coupon\_amount \textless{}dbl\textgreater{}, redemption\_amount \textless{}dbl\textgreater{},}
\CommentTok{\#\textgreater{} \#   total\_amount \textless{}dbl\textgreater{}, discount\_factor \textless{}dbl\textgreater{}, total\_present\_value \textless{}dbl\textgreater{}}
\end{Highlighting}
\end{Shaded}

長期のFX期待値はDomestic金利とForeign金利の差の影響を受ける。本例ではForeign金利がDomestic金利より高いため、リスク中立測度の下ではFX
Rateに下向きのドリフトが働く。その結果、満期が長くなるほどCap領域からLinear領域またはFloor領域へ移る確率が高くなる。

\section{9.4
Couponの期間構造と分布}\label{couponux306eux671fux9593ux69cbux9020ux3068ux5206ux5e03}

\subsection{9.4.1
期待Couponの期間構造}\label{ux671fux5f85couponux306eux671fux9593ux69cbux9020}

\begin{Shaded}
\begin{Highlighting}[]
\NormalTok{expected\_coupon\_plot }\OtherTok{\textless{}{-}}\NormalTok{ ggplot2}\SpecialCharTok{::}\FunctionTok{ggplot}\NormalTok{(}
\NormalTok{  expected\_coupon\_tbl,}
\NormalTok{  ggplot2}\SpecialCharTok{::}\FunctionTok{aes}\NormalTok{(}
    \AttributeTok{x =}\NormalTok{ .data}\SpecialCharTok{$}\NormalTok{time,}
    \AttributeTok{y =}\NormalTok{ .data}\SpecialCharTok{$}\NormalTok{final\_coupon\_rate}
\NormalTok{  )}
\NormalTok{) }\SpecialCharTok{+}
\NormalTok{  ggplot2}\SpecialCharTok{::}\FunctionTok{geom\_line}\NormalTok{(}
    \AttributeTok{linewidth =} \FloatTok{0.8}
\NormalTok{  ) }\SpecialCharTok{+}
\NormalTok{  ggplot2}\SpecialCharTok{::}\FunctionTok{geom\_point}\NormalTok{(}
    \AttributeTok{size =} \FloatTok{1.8}
\NormalTok{  ) }\SpecialCharTok{+}
\NormalTok{  ggplot2}\SpecialCharTok{::}\FunctionTok{scale\_y\_continuous}\NormalTok{(}
    \AttributeTok{labels =}\NormalTok{ scales}\SpecialCharTok{::}\FunctionTok{percent\_format}\NormalTok{(}
      \AttributeTok{accuracy =} \FloatTok{0.1}
\NormalTok{    )}
\NormalTok{  ) }\SpecialCharTok{+}
\NormalTok{  ggplot2}\SpecialCharTok{::}\FunctionTok{labs}\NormalTok{(}
    \AttributeTok{title =} \StringTok{"Expected PRDC coupon rates"}\NormalTok{,}
    \AttributeTok{subtitle =} \StringTok{"Monte Carlo expectation at each future coupon date"}\NormalTok{,}
    \AttributeTok{x =} \StringTok{"Time from valuation date (years)"}\NormalTok{,}
    \AttributeTok{y =} \StringTok{"Expected coupon rate"}
\NormalTok{  ) }\SpecialCharTok{+}
\NormalTok{  ggplot2}\SpecialCharTok{::}\FunctionTok{theme\_minimal}\NormalTok{()}

\FunctionTok{print}\NormalTok{(expected\_coupon\_plot)}
\end{Highlighting}
\end{Shaded}

\pandocbounded{\includegraphics[keepaspectratio]{part2_quantive_riemann/ch09_exotic_derivatives_prdc_files/figure-pdf/chapter09-expected-coupon-plot-1.pdf}}

期待Couponは、単に期待FXをCoupon式へ代入した値とは一般に一致しない。

\[
E[C(S_t)]
\neq
C(E[S_t])
\]

CapとFloorを含む非線形Payoffでは、FX分布全体を考慮する必要がある。

\subsection{9.4.2 Coupon
Cap／Floor到達確率}\label{coupon-capfloorux5230ux9054ux78baux7387}

各Coupon日について、シミュレーションされたCouponがFloorまたはCapに到達する確率を計算する。

\begin{Shaded}
\begin{Highlighting}[]
\NormalTok{coupon\_cap\_floor\_probability\_tbl }\OtherTok{\textless{}{-}}\NormalTok{ GaussQuant}\SpecialCharTok{::}\FunctionTok{prdc\_cap\_floor\_probability\_GQL}\NormalTok{(}
  \AttributeTok{path\_matrix =}\NormalTok{ fx\_path\_matrix\_prdc,}
  \AttributeTok{fx\_base =}\NormalTok{ coupon\_fx\_base\_prdc,}
  \AttributeTok{leverage =}\NormalTok{ coupon\_leverage\_prdc,}
  \AttributeTok{subtraction =}\NormalTok{ coupon\_subtraction\_prdc,}
  \AttributeTok{floor\_rate =}\NormalTok{ coupon\_floor\_prdc,}
  \AttributeTok{cap\_rate =}\NormalTok{ coupon\_cap\_prdc,}
  \AttributeTok{coupon\_dates =}\NormalTok{ fx\_path\_generator\_prdc}\SpecialCharTok{$}\NormalTok{remaining\_coupon\_dates,}
  \AttributeTok{coupon\_times =}\NormalTok{ fx\_path\_generator\_prdc}\SpecialCharTok{$}\NormalTok{coupon\_times}
\NormalTok{)}

\NormalTok{GaussQuant}\SpecialCharTok{::}\FunctionTok{show\_tbl\_GQH}\NormalTok{(}
\NormalTok{  coupon\_cap\_floor\_probability\_tbl,}
  \StringTok{"PRDC coupon floor and cap probabilities"}\NormalTok{,}
  \AttributeTok{n =} \DecValTok{60}
\NormalTok{)}
\CommentTok{\#\textgreater{} \# A tibble: 37 x 6}
\CommentTok{\#\textgreater{}    coupon\_number floor\_probability cap\_probability interior\_probability}
\CommentTok{\#\textgreater{}            \textless{}int\textgreater{}             \textless{}dbl\textgreater{}           \textless{}dbl\textgreater{}                \textless{}dbl\textgreater{}}
\CommentTok{\#\textgreater{}  1             1          0                  0.931               0.0687}
\CommentTok{\#\textgreater{}  2             2          0.000667           0.803               0.197 }
\CommentTok{\#\textgreater{}  3             3          0.0137             0.722               0.265 }
\CommentTok{\#\textgreater{}  4             4          0.0347             0.653               0.312 }
\CommentTok{\#\textgreater{}  5             5          0.0633             0.615               0.322 }
\CommentTok{\#\textgreater{}  6             6          0.0917             0.581               0.327 }
\CommentTok{\#\textgreater{}  7             7          0.118              0.546               0.336 }
\CommentTok{\#\textgreater{}  8             8          0.150              0.523               0.327 }
\CommentTok{\#\textgreater{}  9             9          0.175              0.493               0.332 }
\CommentTok{\#\textgreater{} 10            10          0.205              0.465               0.330 }
\CommentTok{\#\textgreater{} \# i 27 more rows}
\CommentTok{\#\textgreater{} \# i 2 more variables: coupon\_date \textless{}date\textgreater{}, time \textless{}dbl\textgreater{}}
\end{Highlighting}
\end{Shaded}

Cap確率、Linear領域の確率およびFloor確率は、長期的なCouponの性質を直接示す。現在のSpotがCap水準を上回っていても、長期Couponが常にCapへ固定されるとは限らない。

\subsection{9.4.3 Coupon分布}\label{couponux5206ux5e03}

短期、中期、長期の代表的なCoupon日について、Coupon
rateの分布を比較する。

\begin{Shaded}
\begin{Highlighting}[]
\NormalTok{coupon\_path\_matrix\_prdc }\OtherTok{\textless{}{-}}\NormalTok{ GaussQuant}\SpecialCharTok{::}\FunctionTok{prdc\_coupon\_GQL}\NormalTok{(}
  \AttributeTok{fx\_rate =}\NormalTok{ fx\_path\_matrix\_prdc,}
  \AttributeTok{fx\_base =}\NormalTok{ coupon\_fx\_base\_prdc,}
  \AttributeTok{leverage =}\NormalTok{ coupon\_leverage\_prdc,}
  \AttributeTok{subtraction =}\NormalTok{ coupon\_subtraction\_prdc,}
  \AttributeTok{floor\_rate =}\NormalTok{ coupon\_floor\_prdc,}
  \AttributeTok{cap\_rate =}\NormalTok{ coupon\_cap\_prdc}
\NormalTok{)}

\NormalTok{number\_of\_coupon\_steps\_prdc }\OtherTok{\textless{}{-}} \FunctionTok{nrow}\NormalTok{(}
\NormalTok{  coupon\_path\_matrix\_prdc}
\NormalTok{)}

\NormalTok{distribution\_coupon\_ids }\OtherTok{\textless{}{-}} \FunctionTok{unique}\NormalTok{(}
  \FunctionTok{pmax}\NormalTok{(}
    \DecValTok{1}\NormalTok{L,}
    \FunctionTok{pmin}\NormalTok{(}
\NormalTok{      number\_of\_coupon\_steps\_prdc,}
      \FunctionTok{c}\NormalTok{(}
        \DecValTok{1}\NormalTok{L,}
        \FunctionTok{round}\NormalTok{(}
\NormalTok{          number\_of\_coupon\_steps\_prdc }\SpecialCharTok{/}
            \DecValTok{2}
\NormalTok{        ),}
\NormalTok{        number\_of\_coupon\_steps\_prdc}
\NormalTok{      )}
\NormalTok{    )}
\NormalTok{  )}
\NormalTok{)}

\NormalTok{coupon\_distribution\_tbl }\OtherTok{\textless{}{-}}\NormalTok{ purrr}\SpecialCharTok{::}\FunctionTok{map\_dfr}\NormalTok{(}
\NormalTok{  distribution\_coupon\_ids,}
  \ControlFlowTok{function}\NormalTok{(coupon\_id) \{}
\NormalTok{    tibble}\SpecialCharTok{::}\FunctionTok{tibble}\NormalTok{(}
      \AttributeTok{coupon\_number =}\NormalTok{ coupon\_id,}
      \AttributeTok{coupon\_date =} \FunctionTok{as.character}\NormalTok{(}
\NormalTok{        fx\_path\_generator\_prdc}\SpecialCharTok{$}\NormalTok{remaining\_coupon\_dates[}
\NormalTok{          coupon\_id}
\NormalTok{        ]}
\NormalTok{      ),}
      \AttributeTok{coupon\_rate =}\NormalTok{ coupon\_path\_matrix\_prdc[}
\NormalTok{        coupon\_id,}
\NormalTok{      ]}
\NormalTok{    )}
\NormalTok{  \}}
\NormalTok{)}

\NormalTok{coupon\_distribution\_plot }\OtherTok{\textless{}{-}}\NormalTok{ ggplot2}\SpecialCharTok{::}\FunctionTok{ggplot}\NormalTok{(}
\NormalTok{  coupon\_distribution\_tbl,}
\NormalTok{  ggplot2}\SpecialCharTok{::}\FunctionTok{aes}\NormalTok{(}
    \AttributeTok{x =}\NormalTok{ .data}\SpecialCharTok{$}\NormalTok{coupon\_rate}
\NormalTok{  )}
\NormalTok{) }\SpecialCharTok{+}
\NormalTok{  ggplot2}\SpecialCharTok{::}\FunctionTok{geom\_histogram}\NormalTok{(}
    \AttributeTok{bins =} \DecValTok{30}\NormalTok{,}
    \AttributeTok{linewidth =} \FloatTok{0.3}
\NormalTok{  ) }\SpecialCharTok{+}
\NormalTok{  ggplot2}\SpecialCharTok{::}\FunctionTok{facet\_wrap}\NormalTok{(}
    \SpecialCharTok{\textasciitilde{}}\NormalTok{coupon\_date,}
    \AttributeTok{scales =} \StringTok{"free\_y"}
\NormalTok{  ) }\SpecialCharTok{+}
\NormalTok{  ggplot2}\SpecialCharTok{::}\FunctionTok{scale\_x\_continuous}\NormalTok{(}
    \AttributeTok{labels =}\NormalTok{ scales}\SpecialCharTok{::}\FunctionTok{percent\_format}\NormalTok{(}
      \AttributeTok{accuracy =} \FloatTok{0.1}
\NormalTok{    )}
\NormalTok{  ) }\SpecialCharTok{+}
\NormalTok{  ggplot2}\SpecialCharTok{::}\FunctionTok{labs}\NormalTok{(}
    \AttributeTok{title =} \StringTok{"Simulated PRDC coupon distributions"}\NormalTok{,}
    \AttributeTok{x =} \StringTok{"Coupon rate"}\NormalTok{,}
    \AttributeTok{y =} \StringTok{"Path count"}
\NormalTok{  ) }\SpecialCharTok{+}
\NormalTok{  ggplot2}\SpecialCharTok{::}\FunctionTok{theme\_minimal}\NormalTok{()}

\FunctionTok{print}\NormalTok{(coupon\_distribution\_plot)}
\end{Highlighting}
\end{Shaded}

\pandocbounded{\includegraphics[keepaspectratio]{part2_quantive_riemann/ch09_exotic_derivatives_prdc_files/figure-pdf/chapter09-coupon-distribution-1.pdf}}

Coupon分布は連続的な正規分布または対数正規分布ではない。FloorとCapに確率質量が集まり、その間にLinear領域の連続分布が存在する。

\section{9.5 FX感応度とMonte
Carlo誤差}\label{fxux611fux5fdcux5ea6ux3068monte-carloux8aa4ux5dee}

\subsection{9.5.1 FX Spot感応度}\label{fx-spotux611fux5fdcux5ea6}

PRDC CouponはFX Spotへ直接依存する。Spotを上下させたシナリオを比較する。

各シナリオで同じ乱数Seedを使用することにより、Monte
Carloノイズを抑えたCommon Random Numbers型の比較となる。

\begin{Shaded}
\begin{Highlighting}[]
\NormalTok{spot\_scenario\_input\_tbl }\OtherTok{\textless{}{-}}\NormalTok{ tibble}\SpecialCharTok{::}\FunctionTok{tribble}\NormalTok{(}
  \SpecialCharTok{\textasciitilde{}}\NormalTok{scenario,       }\SpecialCharTok{\textasciitilde{}}\NormalTok{fx\_spot,}
  \StringTok{"spot\_minus\_10"}\NormalTok{,  fx\_spot\_prdc }\SpecialCharTok{*} \FloatTok{0.90}\NormalTok{,}
  \StringTok{"spot\_minus\_5"}\NormalTok{,   fx\_spot\_prdc }\SpecialCharTok{*} \FloatTok{0.95}\NormalTok{,}
  \StringTok{"base"}\NormalTok{,           fx\_spot\_prdc,}
  \StringTok{"spot\_plus\_5"}\NormalTok{,    fx\_spot\_prdc }\SpecialCharTok{*} \FloatTok{1.05}\NormalTok{,}
  \StringTok{"spot\_plus\_10"}\NormalTok{,   fx\_spot\_prdc }\SpecialCharTok{*} \FloatTok{1.10}
\NormalTok{)}

\NormalTok{spot\_sensitivity\_tbl }\OtherTok{\textless{}{-}}\NormalTok{ GaussQuant}\SpecialCharTok{::}\FunctionTok{prdc\_spot\_sensitivity\_GQL}\NormalTok{(}
  \AttributeTok{spot\_scenarios =}\NormalTok{ spot\_scenario\_input\_tbl,}
  \AttributeTok{valuation\_date =}\NormalTok{ valuation\_date\_prdc,}
  \AttributeTok{coupon\_schedule =}\NormalTok{ coupon\_schedule\_prdc,}
  \AttributeTok{foreign\_curve\_handle =}\NormalTok{ foreign\_curve\_handle\_prdc,}
  \AttributeTok{domestic\_curve\_handle =}\NormalTok{ domestic\_curve\_handle\_prdc,}
  \AttributeTok{volatility\_handle =}\NormalTok{ fx\_volatility\_handle\_prdc,}
  \AttributeTok{discount\_curve =}\NormalTok{ discount\_curve\_prdc,}
  \AttributeTok{notional =}\NormalTok{ notional\_prdc,}
  \AttributeTok{fx\_base =}\NormalTok{ coupon\_fx\_base\_prdc,}
  \AttributeTok{leverage =}\NormalTok{ coupon\_leverage\_prdc,}
  \AttributeTok{subtraction =}\NormalTok{ coupon\_subtraction\_prdc,}
  \AttributeTok{floor\_rate =}\NormalTok{ coupon\_floor\_prdc,}
  \AttributeTok{cap\_rate =}\NormalTok{ coupon\_cap\_prdc,}
  \AttributeTok{day\_counter =}\NormalTok{ fx\_day\_counter\_prdc,}
  \AttributeTok{intro\_coupon\_dates =}\NormalTok{ intro\_coupon\_dates\_prdc,}
  \AttributeTok{intro\_coupon\_rate =}\NormalTok{ intro\_coupon\_rate\_prdc,}
  \AttributeTok{n\_paths =}\NormalTok{ sensitivity\_paths\_prdc,}
  \AttributeTok{seed =}\NormalTok{ random\_seed\_prdc,}
  \AttributeTok{redemption\_amount =}\NormalTok{ notional\_prdc,}
  \AttributeTok{base\_scenario =} \StringTok{"base"}
\NormalTok{)}

\NormalTok{spot\_sensitivity\_plot }\OtherTok{\textless{}{-}}\NormalTok{ ggplot2}\SpecialCharTok{::}\FunctionTok{ggplot}\NormalTok{(}
\NormalTok{  spot\_sensitivity\_tbl,}
\NormalTok{  ggplot2}\SpecialCharTok{::}\FunctionTok{aes}\NormalTok{(}
    \AttributeTok{x =}\NormalTok{ .data}\SpecialCharTok{$}\NormalTok{fx\_spot,}
    \AttributeTok{y =}\NormalTok{ .data}\SpecialCharTok{$}\NormalTok{prdc\_npv}
\NormalTok{  )}
\NormalTok{) }\SpecialCharTok{+}
\NormalTok{  ggplot2}\SpecialCharTok{::}\FunctionTok{geom\_line}\NormalTok{(}
    \AttributeTok{linewidth =} \FloatTok{0.7}
\NormalTok{  ) }\SpecialCharTok{+}
\NormalTok{  ggplot2}\SpecialCharTok{::}\FunctionTok{geom\_point}\NormalTok{(}
    \AttributeTok{size =} \FloatTok{2.2}
\NormalTok{  ) }\SpecialCharTok{+}
\NormalTok{  ggplot2}\SpecialCharTok{::}\FunctionTok{labs}\NormalTok{(}
    \AttributeTok{title =} \StringTok{"PRDC FX spot sensitivity"}\NormalTok{,}
    \AttributeTok{x =} \StringTok{"FX spot"}\NormalTok{,}
    \AttributeTok{y =} \StringTok{"PRDC NPV"}
\NormalTok{  ) }\SpecialCharTok{+}
\NormalTok{  ggplot2}\SpecialCharTok{::}\FunctionTok{theme\_minimal}\NormalTok{()}

\NormalTok{GaussQuant}\SpecialCharTok{::}\FunctionTok{show\_tbl\_GQH}\NormalTok{(}
\NormalTok{  spot\_sensitivity\_tbl,}
  \StringTok{"PRDC FX spot sensitivity"}\NormalTok{,}
  \AttributeTok{n =} \DecValTok{20}
\NormalTok{)}
\CommentTok{\#\textgreater{} \# A tibble: 5 x 4}
\CommentTok{\#\textgreater{}   scenario      fx\_spot   prdc\_npv npv\_difference\_from\_base}
\CommentTok{\#\textgreater{}   \textless{}chr\textgreater{}           \textless{}dbl\textgreater{}      \textless{}dbl\textgreater{}                    \textless{}dbl\textgreater{}}
\CommentTok{\#\textgreater{} 1 spot\_minus\_10    120. 317016976.               {-}16325124.}
\CommentTok{\#\textgreater{} 2 spot\_minus\_5     127. 325687168.                {-}7654931.}
\CommentTok{\#\textgreater{} 3 base             133. 333342099.                       0 }
\CommentTok{\#\textgreater{} 4 spot\_plus\_5      140. 340136016.                 6793916.}
\CommentTok{\#\textgreater{} 5 spot\_plus\_10     147. 346121967.                12779868.}

\FunctionTok{print}\NormalTok{(spot\_sensitivity\_plot)}
\end{Highlighting}
\end{Shaded}

\pandocbounded{\includegraphics[keepaspectratio]{part2_quantive_riemann/ch09_exotic_derivatives_prdc_files/figure-pdf/chapter09-spot-sensitivity-1.pdf}}

Spot上昇はCouponを上昇させる方向に働くが、Capへ到達した後の限界効果は小さくなる。Spot低下時も、Floorへ到達するとCouponの低下は止まる。このため、PRDC価値のSpot感応度は全域で一定ではない。

\subsection{9.5.2 FX
Volatility感応度}\label{fx-volatilityux611fux5fdcux5ea6}

PRDC Couponは、FloorとCapを持つ非線形Payoffである。そのため、FX
volatilityの変化は将来のFX分布だけでなく、Floor、Linear領域およびCapへ入る確率を変化させる。

\begin{Shaded}
\begin{Highlighting}[]
\NormalTok{volatility\_scenario\_input\_tbl }\OtherTok{\textless{}{-}}\NormalTok{ tibble}\SpecialCharTok{::}\FunctionTok{tribble}\NormalTok{(}
  \SpecialCharTok{\textasciitilde{}}\NormalTok{scenario,     }\SpecialCharTok{\textasciitilde{}}\NormalTok{fx\_volatility,}
  \StringTok{"vol\_5pct"}\NormalTok{,     }\FloatTok{0.05}\NormalTok{,}
  \StringTok{"vol\_7\_5pct"}\NormalTok{,   }\FloatTok{0.075}\NormalTok{,}
  \StringTok{"base"}\NormalTok{,         fx\_volatility\_prdc,}
  \StringTok{"vol\_12\_5pct"}\NormalTok{,  }\FloatTok{0.125}\NormalTok{,}
  \StringTok{"vol\_15pct"}\NormalTok{,    }\FloatTok{0.15}
\NormalTok{)}

\NormalTok{volatility\_sensitivity\_tbl }\OtherTok{\textless{}{-}}\NormalTok{ GaussQuant}\SpecialCharTok{::}\FunctionTok{prdc\_volatility\_sensitivity\_GQL}\NormalTok{(}
  \AttributeTok{volatility\_scenarios =}\NormalTok{ volatility\_scenario\_input\_tbl,}
  \AttributeTok{valuation\_date =}\NormalTok{ valuation\_date\_prdc,}
  \AttributeTok{coupon\_schedule =}\NormalTok{ coupon\_schedule\_prdc,}
  \AttributeTok{spot =}\NormalTok{ fx\_spot\_prdc,}
  \AttributeTok{foreign\_curve\_handle =}\NormalTok{ foreign\_curve\_handle\_prdc,}
  \AttributeTok{domestic\_curve\_handle =}\NormalTok{ domestic\_curve\_handle\_prdc,}
  \AttributeTok{discount\_curve =}\NormalTok{ discount\_curve\_prdc,}
  \AttributeTok{notional =}\NormalTok{ notional\_prdc,}
  \AttributeTok{fx\_base =}\NormalTok{ coupon\_fx\_base\_prdc,}
  \AttributeTok{leverage =}\NormalTok{ coupon\_leverage\_prdc,}
  \AttributeTok{subtraction =}\NormalTok{ coupon\_subtraction\_prdc,}
  \AttributeTok{floor\_rate =}\NormalTok{ coupon\_floor\_prdc,}
  \AttributeTok{cap\_rate =}\NormalTok{ coupon\_cap\_prdc,}
  \AttributeTok{day\_counter =}\NormalTok{ fx\_day\_counter\_prdc,}
  \AttributeTok{calendar =}\NormalTok{ QuantLib}\SpecialCharTok{::}\FunctionTok{NullCalendar}\NormalTok{(),}
  \AttributeTok{intro\_coupon\_dates =}\NormalTok{ intro\_coupon\_dates\_prdc,}
  \AttributeTok{intro\_coupon\_rate =}\NormalTok{ intro\_coupon\_rate\_prdc,}
  \AttributeTok{n\_paths =}\NormalTok{ sensitivity\_paths\_prdc,}
  \AttributeTok{seed =}\NormalTok{ random\_seed\_prdc,}
  \AttributeTok{redemption\_amount =}\NormalTok{ notional\_prdc,}
  \AttributeTok{base\_scenario =} \StringTok{"base"}
\NormalTok{)}

\NormalTok{volatility\_sensitivity\_plot }\OtherTok{\textless{}{-}}\NormalTok{ ggplot2}\SpecialCharTok{::}\FunctionTok{ggplot}\NormalTok{(}
\NormalTok{  volatility\_sensitivity\_tbl,}
\NormalTok{  ggplot2}\SpecialCharTok{::}\FunctionTok{aes}\NormalTok{(}
    \AttributeTok{x =}\NormalTok{ .data}\SpecialCharTok{$}\NormalTok{fx\_volatility,}
    \AttributeTok{y =}\NormalTok{ .data}\SpecialCharTok{$}\NormalTok{prdc\_npv}
\NormalTok{  )}
\NormalTok{) }\SpecialCharTok{+}
\NormalTok{  ggplot2}\SpecialCharTok{::}\FunctionTok{geom\_line}\NormalTok{(}
    \AttributeTok{linewidth =} \FloatTok{0.7}
\NormalTok{  ) }\SpecialCharTok{+}
\NormalTok{  ggplot2}\SpecialCharTok{::}\FunctionTok{geom\_point}\NormalTok{(}
    \AttributeTok{size =} \FloatTok{2.2}
\NormalTok{  ) }\SpecialCharTok{+}
\NormalTok{  ggplot2}\SpecialCharTok{::}\FunctionTok{scale\_x\_continuous}\NormalTok{(}
    \AttributeTok{labels =}\NormalTok{ scales}\SpecialCharTok{::}\FunctionTok{percent\_format}\NormalTok{(}
      \AttributeTok{accuracy =} \FloatTok{0.1}
\NormalTok{    )}
\NormalTok{  ) }\SpecialCharTok{+}
\NormalTok{  ggplot2}\SpecialCharTok{::}\FunctionTok{labs}\NormalTok{(}
    \AttributeTok{title =} \StringTok{"PRDC FX volatility sensitivity"}\NormalTok{,}
    \AttributeTok{x =} \StringTok{"FX volatility"}\NormalTok{,}
    \AttributeTok{y =} \StringTok{"PRDC NPV"}
\NormalTok{  ) }\SpecialCharTok{+}
\NormalTok{  ggplot2}\SpecialCharTok{::}\FunctionTok{theme\_minimal}\NormalTok{()}

\NormalTok{GaussQuant}\SpecialCharTok{::}\FunctionTok{show\_tbl\_GQH}\NormalTok{(}
\NormalTok{  volatility\_sensitivity\_tbl,}
  \StringTok{"PRDC FX volatility sensitivity"}\NormalTok{,}
  \AttributeTok{n =} \DecValTok{20}
\NormalTok{)}
\CommentTok{\#\textgreater{} \# A tibble: 5 x 4}
\CommentTok{\#\textgreater{}   scenario    fx\_volatility   prdc\_npv npv\_difference\_from\_base}
\CommentTok{\#\textgreater{}   \textless{}chr\textgreater{}               \textless{}dbl\textgreater{}      \textless{}dbl\textgreater{}                    \textless{}dbl\textgreater{}}
\CommentTok{\#\textgreater{} 1 vol\_5pct            0.05  337897654.                 4555555.}
\CommentTok{\#\textgreater{} 2 vol\_7\_5pct          0.075 335766964.                 2424865.}
\CommentTok{\#\textgreater{} 3 base                0.1   333342099.                       0 }
\CommentTok{\#\textgreater{} 4 vol\_12\_5pct         0.125 330911366.                {-}2430733.}
\CommentTok{\#\textgreater{} 5 vol\_15pct           0.15  328542072.                {-}4800028.}

\FunctionTok{print}\NormalTok{(volatility\_sensitivity\_plot)}
\end{Highlighting}
\end{Shaded}

\pandocbounded{\includegraphics[keepaspectratio]{part2_quantive_riemann/ch09_exotic_derivatives_prdc_files/figure-pdf/chapter09-volatility-sensitivity-1.pdf}}

通常のOptionのように「Volatility上昇は必ず価値上昇」と単純には言えない。PRDC
CouponはCap付きの上昇Payoffであると同時にFloorを持ち、現在のSpot、金利差、満期および各境界への到達確率によってVolatility感応度が変わる。

\subsection{9.5.3 Monte
Carlo標準誤差}\label{monte-carloux6a19ux6e96ux8aa4ux5dee}

Monte
Carlo推定値には統計誤差が存在する。各パスについて全CouponとRedemptionの現在価値を計算し、標本平均、標準偏差、標準誤差および信頼区間を求める。

標準誤差は、パス現在価値の標準偏差を \(\widehat{\sigma}_{PV}\)、パス数を
\(M\) とすると、

\[
SE
=
\frac{
\widehat{\sigma}_{PV}
}{
\sqrt{M}
}
\]

である。

\begin{Shaded}
\begin{Highlighting}[]
\NormalTok{monte\_carlo\_error\_result }\OtherTok{\textless{}{-}}\NormalTok{ GaussQuant}\SpecialCharTok{::}\FunctionTok{prdc\_monte\_carlo\_error\_GQL}\NormalTok{(}
  \AttributeTok{path\_matrix =}\NormalTok{ fx\_path\_matrix\_prdc,}
  \AttributeTok{cashflow\_tbl =}\NormalTok{ prdc\_cashflow\_tbl,}
  \AttributeTok{notional =}\NormalTok{ notional\_prdc,}
  \AttributeTok{fx\_base =}\NormalTok{ coupon\_fx\_base\_prdc,}
  \AttributeTok{leverage =}\NormalTok{ coupon\_leverage\_prdc,}
  \AttributeTok{subtraction =}\NormalTok{ coupon\_subtraction\_prdc,}
  \AttributeTok{floor\_rate =}\NormalTok{ coupon\_floor\_prdc,}
  \AttributeTok{cap\_rate =}\NormalTok{ coupon\_cap\_prdc,}
  \AttributeTok{intro\_coupon\_dates =}\NormalTok{ intro\_coupon\_dates\_prdc,}
  \AttributeTok{intro\_coupon\_rate =}\NormalTok{ intro\_coupon\_rate\_prdc,}
  \AttributeTok{confidence\_level =} \FloatTok{0.95}
\NormalTok{)}

\NormalTok{path\_present\_value\_tbl }\OtherTok{\textless{}{-}}\NormalTok{ monte\_carlo\_error\_result}\SpecialCharTok{$}\NormalTok{path\_present\_values}
\NormalTok{monte\_carlo\_error\_tbl }\OtherTok{\textless{}{-}}\NormalTok{ monte\_carlo\_error\_result}\SpecialCharTok{$}\NormalTok{summary}

\NormalTok{GaussQuant}\SpecialCharTok{::}\FunctionTok{show\_tbl\_GQH}\NormalTok{(}
\NormalTok{  monte\_carlo\_error\_tbl,}
  \StringTok{"PRDC Monte Carlo error estimate"}\NormalTok{,}
  \AttributeTok{n =} \DecValTok{20}
\NormalTok{)}
\CommentTok{\#\textgreater{} \# A tibble: 7 x 2}
\CommentTok{\#\textgreater{}   metric                        value}
\CommentTok{\#\textgreater{}   \textless{}chr\textgreater{}                         \textless{}dbl\textgreater{}}
\CommentTok{\#\textgreater{} 1 number\_of\_paths             3000   }
\CommentTok{\#\textgreater{} 2 mean\_present\_value     334557294.  }
\CommentTok{\#\textgreater{} 3 present\_value\_sd        36143296.  }
\CommentTok{\#\textgreater{} 4 standard\_error            659883.  }
\CommentTok{\#\textgreater{} 5 lower\_confidence\_bound 333263946.  }
\CommentTok{\#\textgreater{} 6 upper\_confidence\_bound 335850641.  }
\CommentTok{\#\textgreater{} 7 confidence\_level               0.95}
\end{Highlighting}
\end{Shaded}

パス数を4倍にすると、他の条件が同じであれば標準誤差は概ね半分になる。ただし、計算量も増加する。実務では、パス数を増やすだけでなく、Antithetic
Variates、Control Variates、Quasi-Monte
Carloなどの分散削減法を併用する。

\section{9.6 モデルの限界とPart
II総括}\label{ux30e2ux30c7ux30ebux306eux9650ux754cux3068part-iiux7dcfux62ec}

\subsection{9.6.1
PRDC評価の限界}\label{prdcux8a55ux4fa1ux306eux9650ux754c}

本章のモデルは、PRDCの基本的な非線形CouponとMonte
Carlo評価を理解するための簡略モデルである。主な制約は次のとおりである。

\subsection{9.6.2
金利を決定論的とした}\label{ux91d1ux5229ux3092ux6c7aux5b9aux8ad6ux7684ux3068ux3057ux305f}

Domestic、ForeignおよびDiscount
Curveをフラットかつ決定論的とした。実際の長期PRDCでは、Domestic金利とForeign金利の変動が、FX
Forward、割引係数およびCoupon価値へ影響する。

\subsection{9.6.3 FX
volatilityを一定とした}\label{fx-volatilityux3092ux4e00ux5b9aux3068ux3057ux305f}

一定のBlack volatilityを使用した。実際のFX
Option市場には、満期方向とStrike方向のVolatility
Surfaceが存在する。長期PRDCでは、Smile、SkewおよびVolatility term
structureの影響が重要となる。

\subsection{9.6.4
相関を明示的に扱っていない}\label{ux76f8ux95a2ux3092ux660eux793aux7684ux306bux6271ux3063ux3066ux3044ux306aux3044}

本格的なPRDCモデルでは、Domestic金利、Foreign金利およびFXの間の相関を扱う必要がある。金利とFXを別々に評価して単純に合成するだけでは、長期のCross-currencyリスクを十分に表現できない。

\subsection{9.6.5
発行体コール権を含めていない}\label{ux767aux884cux4f53ux30b3ux30fcux30ebux6a29ux3092ux542bux3081ux3066ux3044ux306aux3044}

実際のPRDCには、発行体がCoupon日に商品を期限前償還できるCallable条項が含まれることが多い。

発行体は各行使日で、

\[
\text{Call value}
\quad\text{と}\quad
\text{Continuation value}
\]

を比較する。将来のContinuation
valueはClosed-formでは求めにくいため、Least-Squares Monte
CarloまたはBackward Inductionを使用する。

本章の評価結果は、非Callable PRDCのCoupon LegとRedemption
Legの価値である。Callable PRDCの投資家価値は、概念的には次式となる。

\[
V_{\mathrm{callable\ PRDC}}
=
V_{\mathrm{noncallable\ PRDC}}
-
V_{\mathrm{issuer\ call}}
\]

\subsection{9.6.6 Part
IIのまとめ}\label{part-iiux306eux307eux3068ux3081}

第II部では、金融商品のキャッシュフロー、マーケットコンベンション、確率モデルおよび数値計算を段階的に扱った。

\begin{Shaded}
\begin{Highlighting}[]
\NormalTok{part2\_chapter\_summary\_tbl }\OtherTok{\textless{}{-}}\NormalTok{ tibble}\SpecialCharTok{::}\FunctionTok{tribble}\NormalTok{(}
  \SpecialCharTok{\textasciitilde{}}\NormalTok{chapter, }\SpecialCharTok{\textasciitilde{}}\NormalTok{subject,                              }\SpecialCharTok{\textasciitilde{}}\NormalTok{main\_method,}
  \StringTok{"II{-}1"}\NormalTok{,   }\StringTok{"Bonds and market conventions"}\NormalTok{,        }\StringTok{"Yield, price, duration, and convexity"}\NormalTok{,}
  \StringTok{"II{-}2"}\NormalTok{,   }\StringTok{"Treasury futures and CTD"}\NormalTok{,             }\StringTok{"Basis, carry, implied repo, and futures BPV"}\NormalTok{,}
  \StringTok{"II{-}3"}\NormalTok{,   }\StringTok{"OIS and multi{-}curve valuation"}\NormalTok{,        }\StringTok{"Discount and forecast curve separation"}\NormalTok{,}
  \StringTok{"II{-}4"}\NormalTok{,   }\StringTok{"Black model"}\NormalTok{,                          }\StringTok{"Lognormal forward{-}option pricing"}\NormalTok{,}
  \StringTok{"II{-}5"}\NormalTok{,   }\StringTok{"Normal model"}\NormalTok{,                         }\StringTok{"Bachelier pricing for zero and negative rates"}\NormalTok{,}
  \StringTok{"II{-}6"}\NormalTok{,   }\StringTok{"SABR model"}\NormalTok{,                           }\StringTok{"Volatility{-}smile calibration"}\NormalTok{,}
  \StringTok{"II{-}7"}\NormalTok{,   }\StringTok{"Callable Bond and Hull{-}White"}\NormalTok{,          }\StringTok{"Short{-}rate tree and early exercise"}\NormalTok{,}
  \StringTok{"II{-}8"}\NormalTok{,   }\StringTok{"Credit risk and CDS"}\NormalTok{,                  }\StringTok{"Hazard curve, MidPoint, and ISDA valuation"}\NormalTok{,}
  \StringTok{"II{-}9"}\NormalTok{,   }\StringTok{"Exotic derivatives and PRDC"}\NormalTok{,          }\StringTok{"FX{-}linked coupon Monte Carlo valuation"}
\NormalTok{)}

\NormalTok{GaussQuant}\SpecialCharTok{::}\FunctionTok{show\_tbl\_GQH}\NormalTok{(}
\NormalTok{  part2\_chapter\_summary\_tbl,}
  \StringTok{"Part II chapter summary"}\NormalTok{,}
  \AttributeTok{n =} \DecValTok{20}
\NormalTok{)}
\CommentTok{\#\textgreater{} \# A tibble: 9 x 3}
\CommentTok{\#\textgreater{}   chapter subject                       main\_method                             }
\CommentTok{\#\textgreater{}   \textless{}chr\textgreater{}   \textless{}chr\textgreater{}                         \textless{}chr\textgreater{}                                   }
\CommentTok{\#\textgreater{} 1 II{-}1    Bonds and market conventions  Yield, price, duration, and convexity   }
\CommentTok{\#\textgreater{} 2 II{-}2    Treasury futures and CTD      Basis, carry, implied repo, and futures\textasciitilde{}}
\CommentTok{\#\textgreater{} 3 II{-}3    OIS and multi{-}curve valuation Discount and forecast curve separation  }
\CommentTok{\#\textgreater{} 4 II{-}4    Black model                   Lognormal forward{-}option pricing        }
\CommentTok{\#\textgreater{} 5 II{-}5    Normal model                  Bachelier pricing for zero and negative\textasciitilde{}}
\CommentTok{\#\textgreater{} 6 II{-}6    SABR model                    Volatility{-}smile calibration            }
\CommentTok{\#\textgreater{} 7 II{-}7    Callable Bond and Hull{-}White  Short{-}rate tree and early exercise      }
\CommentTok{\#\textgreater{} 8 II{-}8    Credit risk and CDS           Hazard curve, MidPoint, and ISDA valuat\textasciitilde{}}
\CommentTok{\#\textgreater{} 9 II{-}9    Exotic derivatives and PRDC   FX{-}linked coupon Monte Carlo valuation}
\end{Highlighting}
\end{Shaded}

債券評価から始まり、金利先物、OIS、Option、Volatility
Smile、短期金利モデル、信用リスク、そして複数の市場要因を持つエキゾチック商品へ進んだ。

Part
IIを通じた共通原則は、商品ごとに異なる数式を暗記することではない。次の要素を分離し、再利用可能な計算層へ落とすことである。

\begin{itemize}
\tightlist
\item
  契約条件とSchedule
\item
  市場データとCurve
\item
  リスク中立測度と確率過程
\item
  Cash flowとPayoff
\item
  Pricing Engineまたは数値計算法
\item
  NPV、RiskおよびCalibration
\item
  表示と検証
\end{itemize}

GaussQuantは、QuantLib
SWIGの低水準オブジェクトを金融商品の計算層へまとめる。LagrangeFinanceの章側は、商品構造、数式、入力条件、表、グラフおよび金融的解釈に集中する。

\section{9.7 まとめ}\label{ux307eux3068ux3081-13}

PRDCは、将来のFX Rateに連動するCouponを持つ長期の仕組債である。

Coupon rateは、

\[
C(S)
=
\min
\left[
C_{\max},
\max
\left(
a\frac{S}{S_0}-b,
C_{\min}
\right)
\right]
\]

というFloor付き・Cap付きの非線形Payoffである。

Garman--Kohlhagenモデルでは、FX
Rateのリスク中立ドリフトはDomestic金利とForeign金利の差によって決まる。

\[
\frac{dS_t}{S_t}
=
(r_d-r_f)dt
+
\sigma_{\mathrm{FX}}dW_t
\]

Coupon日のFX分布をMonte Carloで生成し、各Coupon日の期待Coupon、Coupon
Leg、Redemption LegおよびPRDC NPVを計算した。

FX SpotとFX volatilityの感応度は、FloorとCapによって非線形となる。Monte
Carlo評価には統計誤差があるため、NPVだけでなく標準誤差と信頼区間を確認する必要がある。

本章は非Callable PRDCを対象とした。実際のCallable
PRDCには、Domestic金利、Foreign金利、FX、Volatility
Surface、相関および発行体の早期行使判断を同時に扱う、より高度なモデルが必要となる。

\begin{Shaded}
\begin{Highlighting}[]
\NormalTok{chapter09\_summary\_tbl }\OtherTok{\textless{}{-}}\NormalTok{ tibble}\SpecialCharTok{::}\FunctionTok{tribble}\NormalTok{(}
  \SpecialCharTok{\textasciitilde{}}\NormalTok{section,                 }\SpecialCharTok{\textasciitilde{}}\NormalTok{main\_result,}
  \StringTok{"Coupon payoff"}\NormalTok{,           }\StringTok{"Floored and capped FX{-}linked coupon"}\NormalTok{,}
  \StringTok{"FX process"}\NormalTok{,              }\StringTok{"Garman{-}Kohlhagen risk{-}neutral dynamics"}\NormalTok{,}
  \StringTok{"Monte Carlo valuation"}\NormalTok{,   }\StringTok{"Coupon leg and redemption leg NPV"}\NormalTok{,}
  \StringTok{"Spot sensitivity"}\NormalTok{,        }\StringTok{"Nonlinear response caused by cap and floor"}\NormalTok{,}
  \StringTok{"Volatility sensitivity"}\NormalTok{,  }\StringTok{"Distributional effect on coupon regions"}\NormalTok{,}
  \StringTok{"Monte Carlo error"}\NormalTok{,       }\StringTok{"Standard error and confidence interval"}\NormalTok{,}
  \StringTok{"Model boundary"}\NormalTok{,          }\StringTok{"Non{-}callable deterministic{-}rate PRDC"}
\NormalTok{)}

\NormalTok{GaussQuant}\SpecialCharTok{::}\FunctionTok{show\_tbl\_GQH}\NormalTok{(}
\NormalTok{  chapter09\_summary\_tbl,}
  \StringTok{"Chapter II{-}9 summary"}\NormalTok{,}
  \AttributeTok{n =} \DecValTok{20}
\NormalTok{)}
\CommentTok{\#\textgreater{} \# A tibble: 7 x 2}
\CommentTok{\#\textgreater{}   section                main\_result                               }
\CommentTok{\#\textgreater{}   \textless{}chr\textgreater{}                  \textless{}chr\textgreater{}                                     }
\CommentTok{\#\textgreater{} 1 Coupon payoff          Floored and capped FX{-}linked coupon       }
\CommentTok{\#\textgreater{} 2 FX process             Garman{-}Kohlhagen risk{-}neutral dynamics    }
\CommentTok{\#\textgreater{} 3 Monte Carlo valuation  Coupon leg and redemption leg NPV         }
\CommentTok{\#\textgreater{} 4 Spot sensitivity       Nonlinear response caused by cap and floor}
\CommentTok{\#\textgreater{} 5 Volatility sensitivity Distributional effect on coupon regions   }
\CommentTok{\#\textgreater{} 6 Monte Carlo error      Standard error and confidence interval    }
\CommentTok{\#\textgreater{} 7 Model boundary         Non{-}callable deterministic{-}rate PRDC}

\FunctionTok{cat}\NormalTok{(}
  \StringTok{"}\SpecialCharTok{\textbackslash{}n}\StringTok{Chapter II{-}9 Exotic derivatives and PRDC completed successfully.}\SpecialCharTok{\textbackslash{}n}\StringTok{"}
\NormalTok{)}
\CommentTok{\#\textgreater{} }
\CommentTok{\#\textgreater{} Chapter II{-}9 Exotic derivatives and PRDC completed successfully.}
\end{Highlighting}
\end{Shaded}





\end{document}
